% !TeX program = pdflatex
\documentclass[11pt,a4paper,oneside]{book}
\usepackage[T1]{fontenc}
\usepackage[utf8]{inputenc}
\usepackage[polish]{babel}
\usepackage{lmodern}
\usepackage{amsmath,amssymb,amsthm,mathtools}
\usepackage{microtype}
\usepackage{longtable}
\usepackage{geometry}
\usepackage{enumitem}
\usepackage{xcolor}
\usepackage{tikz}
\usepackage[unicode]{hyperref}
\hypersetup{colorlinks=true,linkcolor=blue!45!black,urlcolor=blue!45!black}
\geometry{margin=2.6cm}
\setlength{\parindent}{0pt}
\setlength{\parskip}{4pt}
\setlength{\emergencystretch}{2em}
\setlist[enumerate]{topsep=4pt,itemsep=2pt}
\numberwithin{equation}{chapter}
\theoremstyle{plain}
\newtheorem{twierdzenie}{Twierdzenie}[chapter]
\newtheorem{lemat}[twierdzenie]{Lemat}
\newtheorem{wniosek}[twierdzenie]{Wniosek}
\theoremstyle{definition}
\newtheorem{definicja}[twierdzenie]{Definicja}
\newtheorem{przyklad}[twierdzenie]{Przykład}
\theoremstyle{remark}
\newtheorem{uwaga}[twierdzenie]{Uwaga}
\renewcommand{\proofname}{Dowód}
\newcommand{\dd}{\,\mathrm d}
% Funkcje charakterystyczne oznaczamy wyłącznie przez \chi_A.
\title{Teoria miary\\\large Rozdziały 1--31: od konstrukcji miary do dżetów przybliżonych}
\author{}
\date{2026}
\begin{document}
\frontmatter
\hypersetup{pageanchor=false}
\maketitle
\hypersetup{pageanchor=true}
\tableofcontents

% BEGIN chapters/pl/measure_theory/supplements/spis_oznaczen.tex
% Redakcyjny spis najczęstszych symboli (bez powtarzania definicji z rozdziałów).
\chapter*{Spis oznaczeń}
\markboth{SPIS OZNACZEŃ}{SPIS OZNACZEŃ}
\phantomsection\addcontentsline{toc}{chapter}{Spis oznaczeń}
\noindent Symbole mogą przyjmować dodatkowe założenia w poszczególnych twierdzeniach.
Pierwszeństwo mają zawsze definicje podane we wskazanym rozdziale.

\section*{Zbiory i mierzalność}
\begin{longtable}{@{}p{0.285\textwidth}p{0.555\textwidth}p{0.095\textwidth}@{}}
$X,\,\mathcal P(X)$ & Przestrzeń bazowa i zbiór jej wszystkich podzbiorów. & rozdz.~\ref{ch:przestrzenie-mierzalne} \\
$\mathcal A,\,\mathcal F$ & Rodziny zbiorów; w szczególności $\sigma$-algebry. & rozdz.~\ref{ch:przestrzenie-mierzalne} \\
$\sigma(\mathcal G)$ & Najmniejsza $\sigma$-algebra zawierająca rodzinę $\mathcal G$. & rozdz.~\ref{ch:przestrzenie-mierzalne} \\
$\mathcal B(X)$ & $\sigma$-algebra borelowska przestrzeni topologicznej $X$. & rozdz.~\ref{ch:przestrzenie-mierzalne} \\
$\chi_A$ & Funkcja charakterystyczna zbioru $A$; jedyna stosowana w książce notacja tego pojęcia. & rozdz.~\ref{ch:funkcje-mierzalne} \\
$\mu^*$ & Miara zewnętrzna (niekoniecznie zdefiniowana na $\sigma$-algebrze). & rozdz.~\ref{ch:caratheodory} \\
$\lambda_d$ & $d$-wymiarowa miara Lebesgue'a na $\mathbb R^d$. & rozdz.~\ref{ch:lebesgue} \\
$\mathrm d\mu(x)$ & Element notacji całkowania wskazujący miarę i zmienną całkowania. & rozdz.~\ref{ch:calkowanie-lebesgue} \\
\end{longtable}

\section*{Zbieżność i całkowanie}
\begin{longtable}{@{}p{0.285\textwidth}p{0.555\textwidth}p{0.095\textwidth}@{}}
$\mu\otimes\nu$ & Miara produktowa na mierzalnej przestrzeni produktowej. & rozdz.~\ref{ch:produkty-miary} \\
$L^p(\mu)$ & Przestrzeń klas funkcji $p$-całkowalnych; dla $p=\infty$ funkcji istotnie ograniczonych. & rozdz.~\ref{ch:przestrzenie-lp} \\
$\|f\|_{L^p}$ & Norma w przestrzeni $L^p(\mu)$. & rozdz.~\ref{ch:przestrzenie-lp} \\
$C_c(X)$ & Przestrzeń ciągłych funkcji o zwartym nośniku. & rozdz.~\ref{ch:radon-measures-20} \\
$C_0(X)$ & Przestrzeń ciągłych funkcji zanikających w nieskończoności. & rozdz.~\ref{ch:riesz-markov-21} \\
$F_*\mu$ & Miara obrazowa: $(F_*\mu)(A)=\mu(F^{-1}(A))$. & rozdz.~\ref{ch:rn-applications-19} \\
\end{longtable}

\section*{Miary i ich rozkłady}
\begin{longtable}{@{}p{0.285\textwidth}p{0.555\textwidth}p{0.095\textwidth}@{}}
$\nu\ll\mu$ & Bezwzględna ciągłość miary $\nu$ względem $\mu$. & rozdz.~\ref{ch:ac-singular-17} \\
$\nu\perp\mu$ & Wzajemna osobliwość miar. & rozdz.~\ref{ch:miary-ze-znakiem} \\
$\nu^+,\,\nu^-$ & Dodatnia i ujemna część rozkładu Jordana miary ze znakiem. & rozdz.~\ref{ch:miary-ze-znakiem} \\
$|\nu|$ & Wariacja całkowita miary ze znakiem, zespolonej lub wektorowej; zależnie od kontekstu. & rozdz.~\ref{ch:miary-ze-znakiem} \\
$\frac{\mathrm d\nu}{\mathrm d\mu}$ & Pochodna Radona--Nikodyma, a nie symbol całkowania. & rozdz.~\ref{ch:radon-nikodym-18} \\
$\operatorname{supp}\mu$ & Nośnik miary, z uwzględnieniem odpowiednich założeń topologicznych. & rozdz.~\ref{ch:radon-measures-20} \\
\end{longtable}

\section*{Miary geometryczne}
\begin{longtable}{@{}p{0.285\textwidth}p{0.555\textwidth}p{0.095\textwidth}@{}}
$\mathcal H^s_\delta,\,\mathcal H^s$ & Miary Hausdorffa na skali $\delta$ oraz po przejściu do granicy; z ustaloną normalizacją. & rozdz.~\ref{ch:hausdorff-26} \\
$\dim_{\mathrm H}(E)$ & Wymiar Hausdorffa zbioru $E$. & rozdz.~\ref{ch:hausdorff-26} \\
$\mu_g$ & Miara objętości związana z metryką riemannowską $g$. & rozdz.~\ref{ch:miary-rozmaitosci-27} \\
$J_k(A)$ & $k$-wymiarowy jakobian odwzorowania liniowego $A$. & rozdz.~\ref{ch:area-coarea-28} \\
$\operatorname{Jac}_m(B)$ & $m$-wymiarowy jakobian koobszaru odwzorowania liniowego $B$. & rozdz.~\ref{ch:area-coarea-28} \\
$BV(\Omega)$ & Przestrzeń funkcji o ograniczonej wariacji w $\Omega$. & rozdz.~\ref{ch:geometric-measure-29} \\
$BV_{\mathrm{loc}}(\Omega)$ & Przestrzeń funkcji o wariacji lokalnie skończonej. & rozdz.~\ref{ch:geometric-measure-29} \\
$Du,\,|Du|$ & Miara pochodnej dystrybucyjnej i jej wariacja. & rozdz.~\ref{ch:geometric-measure-29} \\
$P(E;\Omega)$ & Obwód zbioru $E$ względem otwartego zbioru $\Omega$. & rozdz.~\ref{ch:geometric-measure-29} \\
$\partial^*E,\,\nu_E$ & Brzeg zredukowany i zewnętrzna normalna miarowo-teoretyczna. & rozdz.~\ref{ch:geometric-measure-29} \\
\end{longtable}

\section*{Miary wektorowe}
\begin{longtable}{@{}p{0.285\textwidth}p{0.555\textwidth}p{0.095\textwidth}@{}}
$\nu:\mathcal A\to X$ & Miara o wartościach w przestrzeni Banacha $X$. & rozdz.~\ref{ch:vector-measures-30} \\
$\|\nu\|_{\mathrm{sv}}$ & Półwariacja miary wektorowej. & rozdz.~\ref{ch:vector-measures-30} \\
$\int_A g(x)\,\mathrm d\mu(x)$ & Całka funkcji wektorowej w sensie Bochnera, gdy $g$ spełnia odpowiednie założenia. & rozdz.~\ref{ch:vector-measures-30} \\
$\mathrm{RNP}$ & Własność Radona--Nikodyma przestrzeni Banacha. & rozdz.~\ref{ch:vector-measures-30} \\
\end{longtable}

\noindent\emph{Uwaga.} Notacja $\chi_A$ jest stosowana konsekwentnie
w całym skrypcie; miary geometryczne stosują wspólną kalibrację
ustaloną w rozdziale~\ref{ch:hausdorff-26}.


\section*{Dżety i przedłużanie danych}
\begin{longtable}{@{}p{0.285\textwidth}p{0.555\textwidth}p{0.095\textwidth}@{}}
$j_p^kf$ & Dżet $k$-tego rzędu funkcji w punkcie $p$. & rozdz.~\ref{ch:dzety} \\
$J^kE$ & Wiązka dżetów $k$-tego rzędu przekrojów $E\to M$. & rozdz.~\ref{ch:dzety} \\
$D_{\mathrm{ap}}f$ & Pochodna przybliżona, określona w odpowiednich punktach. & rozdz.~\ref{ch:dzety} \\
\end{longtable}


% END chapters/pl/measure_theory/supplements/spis_oznaczen.tex

\mainmatter

% BEGIN chapters/pl/measure_theory/01-rozdzial.tex
% ROZDZIAŁ 1 — plik włączany przez main.tex
\chapter{Rodziny zbiorów i przestrzenie mierzalne}
% Konwencja obowiązująca w całym skrypcie: funkcję charakterystyczną zbioru A oznaczamy przez $\chi_A$.
\label{ch:przestrzenie-mierzalne}

\section{Dlaczego potrzebujemy \texorpdfstring{$\sigma$}{sigma}-algebr?}

W teorii miary chcemy przypisywać rozmiar zbiorom, a następnie całkować funkcje.
Aksjomaty teorii mają zapewniać zgodność tego rozmiaru z rozkładaniem zbioru
na przeliczalnie wiele rozłącznych części. Nie wystarczy zatem umieć mierzyć
pojedynczych zbiorów: potrzebujemy rodziny zbiorów zamkniętej na operacje,
które pojawiają się w granicach ciągów i szeregach. Właśnie tę rolę pełnią
$\sigma$-algebry. Istotne jest, że nie zakładamy mierzalności wszystkich
podzbiorów rozważanej przestrzeni.

W całym rozdziale $X$ i $Y$ oznaczają zbiory, $\mathcal P(X)$ --- zbiór
wszystkich podzbiorów $X$, natomiast $A^c=X\setminus A$ oznacza dopełnienie
zbioru $A\subseteq X$ względem ustalonej przestrzeni $X$. Symbol $\mathbb N$
oznacza zbiór liczb naturalnych dodatnich. Rodziny zbiorów oznaczamy
$\mathcal A,\mathcal F,\mathcal G,\mathcal M$.

\begin{definicja}[Pierścień, algebra, $\sigma$-algebra]
\label{def:algebry}
Niech $X$ będzie zbiorem, a $\mathcal R\subseteq\mathcal P(X)$.
\begin{enumerate}
\item $\mathcal R$ nazywamy \emph{pierścieniem zbiorów}, jeżeli
$\varnothing\in\mathcal R$ i dla dowolnych $A,B\in\mathcal R$ zachodzą
$A\cup B\in\mathcal R$ oraz $A\setminus B\in\mathcal R$.
\item $\mathcal R$ nazywamy \emph{algebrą zbiorów} na $X$, jeżeli jest
pierścieniem i $X\in\mathcal R$.
\item $\mathcal R$ nazywamy \emph{$\sigma$-algebrą} na $X$, jeżeli
$X\in\mathcal R$, $A^c\in\mathcal R$ dla każdego $A\in\mathcal R$ oraz
$\bigcup_{n=1}^{\infty}A_n\in\mathcal R$ dla każdego ciągu
$(A_n)_{n\geq1}$ elementów $\mathcal R$.
\end{enumerate}
Parę $(X,\mathcal F)$, gdzie $\mathcal F$ jest $\sigma$-algebrą na $X$,
nazywamy \emph{przestrzenią mierzalną}; jej elementy to \emph{zbiory
mierzalne}.
\end{definicja}

\begin{uwaga}
Algebra jest zamknięta na skończone sumy, przekroje i różnice zbiorów.
$\sigma$-algebra jest ponadto zamknięta na przeliczalne przekroje, gdyż
z praw de Morgana
$\bigcap_{n\geq1}A_n=(\bigcup_{n\geq1}A_n^c)^c$.
Każda $\sigma$-algebra jest algebrą, ale nie odwrotnie.
\end{uwaga}

\begin{przyklad}[Najmniejsza i największa $\sigma$-algebra]
Dla dowolnego $X$ rodziny $\{\varnothing,X\}$ oraz $\mathcal P(X)$ są
$\sigma$-algebrami. Pierwszą nazywamy \emph{trywialną}, drugą ---
\emph{dyskretną}. Gdy $X=\varnothing$, są one równe.
\end{przyklad}

\begin{przyklad}[Zbiory skończone i kofinitywne]
\label{ex:finite-cofinite}
Niech $X$ będzie nieskończonym zbiorem. Rodzina
\[
\mathcal A=\{A\subseteq X:A\ \text{jest skończony lub }X\setminus A
\ \text{jest skończony}\}
\]
jest algebrą. Jeśli $X=\mathbb N$, nie jest $\sigma$-algebrą:
każdy singleton $\{2n\}$ należy do $\mathcal A$, lecz
$\bigcup_{n\geq1}\{2n\}$ ani nie jest skończony, ani nie ma skończonego
dopełnienia.
\end{przyklad}

\begin{przyklad}[Algebra przeliczalno-koprzeliczalna]
\label{ex:countable-cocountable}
Niech $X$ będzie dowolnym zbiorem. Rodzina
\[
\mathcal C_X=\{A\subseteq X:A\text{ jest przeliczalny lub }X\setminus A
\text{ jest przeliczalny}\}
\]
jest $\sigma$-algebrą. Istotnie, dopełnienia zachowują podaną własność.
Jeżeli w przeliczalnej rodzinie $(A_n)$ występuje zbiór o przeliczalnym
dopełnieniu, to dopełnienie $\bigcup_nA_n$ jest przeliczalne; jeśli zaś
wszystkie $A_n$ są przeliczalne, ich suma również jest przeliczalna.
Jeżeli $X$ jest nieprzeliczalny, to $\mathcal C_X\subsetneq\mathcal P(X)$.
Rzeczywiście, możemy wybrać dwa rozłączne nieprzeliczalne podzbiory
$A,B$ o sumie $X$ (wybór takiego podziału jest możliwy w zwykłych
założeniach teorii mnogości z aksjomatem wyboru). Wtedy ani $A$,
ani $X\setminus A=B$ nie jest przeliczalny, więc
$A\notin\mathcal C_X$.
\end{przyklad}

\begin{przyklad}[Zbiory określone przez podział]
Niech $(E_i)_{i\in I}$ będzie rozbiciem zbioru $X$ na niepuste, parami
rozłączne zbiory. Rodzina wszystkich sum $\bigcup_{i\in J}E_i$ dla
$J\subseteq I$ jest $\sigma$-algebrą. W szczególności dla
$X=\{a,b,c\}$ podział $\{\{a,b\},\{c\}\}$ wyznacza $\sigma$-algebrę
$\{\varnothing,\{a,b\},\{c\},X\}$.
\end{przyklad}

\section{Generowanie \texorpdfstring{$\sigma$}{sigma}-algebr}

\begin{twierdzenie}[Przekroje $\sigma$-algebr]
\label{thm:intersection-sigma}
Niech $X$ będzie zbiorem i niech $(\mathcal F_i)_{i\in I}$ będzie
niepustą rodziną $\sigma$-algebr na $X$. Wówczas
$\bigcap_{i\in I}\mathcal F_i$ jest $\sigma$-algebrą na $X$.
\end{twierdzenie}
\begin{proof}
Wszystkie $\mathcal F_i$ zawierają $X$. Jeśli $A$ należy do każdego
$\mathcal F_i$, to jego dopełnienie również należy do każdego z nich.
Analogicznie, jeśli wszystkie $A_n$ należą do przekroju tych rodzin,
to $\bigcup_nA_n$ należy do każdej $\mathcal F_i$, a więc do ich przekroju.
\end{proof}

\begin{definicja}[$\sigma$-algebra generowana]
\label{def:sigma-generated}
Dla zbioru $X$ i rodziny $\mathcal E\subseteq\mathcal P(X)$ definiujemy
\[
\sigma_X(\mathcal E)=\bigcap\{\mathcal F:\mathcal F\text{ jest
$\sigma$-algebrą na }X\text{ i }\mathcal E\subseteq\mathcal F\}.
\]
Jest to \emph{$\sigma$-algebra generowana przez $\mathcal E$}.
Jeśli przestrzeń $X$ jest ustalona, piszemy po prostu $\sigma(\mathcal E)$.
\end{definicja}

Rodzina pod przekrojem jest niepusta, ponieważ zawiera $\mathcal P(X)$.
Z twierdzenia \ref{thm:intersection-sigma} wynika, że $\sigma(\mathcal E)$
jest $\sigma$-algebrą. Ponadto jest najmniejszą $\sigma$-algebrą
zawierającą $\mathcal E$, w sensie inkluzji.

\begin{przyklad}
Dla $E\subseteq X$ mamy
$\sigma(\{E\})=\{\varnothing,E,E^c,X\}$, z ewentualnymi powtórzeniami,
gdy $E=\varnothing$ lub $E=X$.
\end{przyklad}

\begin{uwaga}[Pułapka dotycząca sumy rodzin]
Nawet rosnąca suma $\sigma$-algebr nie musi być $\sigma$-algebrą.
Na $X=\mathbb N$ niech $\mathcal F_n$ oznacza $\sigma$-algebrę
generowaną przez singletony $\{1\},\ldots,\{n\}$.
Wówczas $\mathcal F_n\subseteq\mathcal F_{n+1}$, ale
$\bigcup_n\mathcal F_n$ jest tylko algebrą zbiorów skończonych i
kofinitywnych, znaną z przykładu \ref{ex:finite-cofinite}.
\end{uwaga}

\begin{definicja}[$\sigma$-algebra śladowa]\label{idx:01:definicja:04}

Niech $(X,\mathcal F)$ będzie przestrzenią mierzalną i niech
$E\subseteq X$ będzie dowolnym podzbiorem. Definiujemy
$\mathcal F|_E=\{A\cap E:A\in\mathcal F\}$.
Rodzina ta jest $\sigma$-algebrą na $E$, nazywaną \emph{śladem}
$\mathcal F$ na $E$.
\end{definicja}
\begin{proof}[Uzasadnienie]
Do rodziny należy $E=X\cap E$. Dopełnienie w $E$ zbioru $A\cap E$
wynosi $A^c\cap E$, a przeliczalna suma zbiorów $A_n\cap E$ jest równa
$(\bigcup_n A_n)\cap E$.
\end{proof}

\section{Układy \texorpdfstring{$\pi$}{pi} i układy Dynkina}

\begin{definicja}[Układ $\pi$ i układ Dynkina]
\label{def:pi-lambda}
Niech $X$ będzie zbiorem. Rodzinę $\mathcal P\subseteq\mathcal P(X)$
nazywamy \emph{układem $\pi$}, jeżeli jest niepusta i
$A\cap B\in\mathcal P$ dla wszystkich $A,B\in\mathcal P$.
Rodzinę $\mathcal D\subseteq\mathcal P(X)$ nazywamy \emph{układem
Dynkina} (lub \emph{układem $\lambda$}), jeżeli:
\begin{enumerate}
\item $X\in\mathcal D$;
\item jeśli $A,B\in\mathcal D$ i $A\subseteq B$, to $B\setminus A\in\mathcal D$;
\item dla dowolnych parami rozłącznych $A_1,A_2,\ldots\in\mathcal D$
zachodzi $\bigcup_nA_n\in\mathcal D$.
\end{enumerate}
\end{definicja}

Każda $\sigma$-algebra jest układem Dynkina i układem $\pi$.
Samo bycie układem Dynkina nie gwarantuje zamknięcia na przekroje.
Przykładowo, jeśli $X=\{1,2,3,4\}$, to rodzina
\[
\mathcal D=\{\varnothing,X,\{1,2\},\{3,4\},\{1,3\},\{2,4\}\}
\]
jest układem Dynkina (jedynymi właściwymi, niepustymi zbiorami są tu
zbiory dwuelementowe), ale nie jest $\sigma$-algebrą, gdyż
$\{1,2\}\cap\{1,3\}=\{1\}\notin\mathcal D$.

\begin{lemat}
\label{lem:dynkin-intersection}
Układ Dynkina $\mathcal D$ na zbiorze $X$ jest $\sigma$-algebrą wtedy
i tylko wtedy, gdy jest zamknięty na skończone przekroje.
\end{lemat}
\begin{proof}
Konieczność jest oczywista. Załóżmy, że $\mathcal D$ jest układem
Dynkina zamkniętym na przekroje. Z warunku drugiego otrzymujemy
$A^c=X\setminus A\in\mathcal D$. Dla $A,B\in\mathcal D$ zachodzi
$B\setminus A=B\setminus(A\cap B)\in\mathcal D$, więc
$A\cup B=A\,\dot\cup\,(B\setminus A)\in\mathcal D$; symbol
$\dot\cup$ oznacza sumę rozłączną. Powtarzając ten krok, uzyskujemy
zamknięcie na sumy skończone. Dla dowolnego ciągu $(A_n)$ elementów
$\mathcal D$ zastępujemy $A_n$ przez
$B_n=A_n\setminus\bigcup_{k<n}A_k$. Zbiory $B_n$ są mierzalne względem
$\mathcal D$, rozłączne, a ich suma jest równa $\bigcup_nA_n$.
Trzeci aksjomat układu Dynkina kończy dowód.
\end{proof}

\begin{twierdzenie}[Twierdzenie $\pi$--$\lambda$ Dynkina]
\label{thm:pi-lambda}
Niech $X$ będzie zbiorem, $\mathcal P$ układem $\pi$ na $X$,
a $\mathcal D$ układem Dynkina na $X$. Jeśli
$\mathcal P\subseteq\mathcal D$, to
$\sigma(\mathcal P)\subseteq\mathcal D$.
\end{twierdzenie}
\begin{proof}
Przekrój dowolnej rodziny układów Dynkina na $X$ jest układem Dynkina.
Niech $\mathcal L$ będzie najmniejszym układem Dynkina zawierającym
$\mathcal P$. Wystarczy pokazać, że $\mathcal L$ jest $\sigma$-algebrą.

Dla ustalonego $A\in\mathcal P$ definiujemy
$\mathcal L_A=\{B\in\mathcal L:A\cap B\in\mathcal L\}$.
Rodzina $\mathcal L_A$ jest układem Dynkina: należy do niej $X$, gdyż
$A\cap X=A\in\mathcal L$; własność różnicy wynika z równości
$A\cap(B\setminus C)=(A\cap B)\setminus(A\cap C)$ dla $C\subseteq B$;
własność rozłącznych sum wynika z rozdzielności przekroju względem sumy.
Ponieważ $\mathcal P$ jest układem $\pi$, mamy
$\mathcal P\subseteq\mathcal L_A$. Z minimalności $\mathcal L$ wynika
$\mathcal L_A=\mathcal L$. Oznacza to, że
$A\cap B\in\mathcal L$ dla wszystkich $A\in\mathcal P$, $B\in\mathcal L$.

Teraz dla ustalonego $B\in\mathcal L$ rozważamy
$\mathcal K_B=\{A\in\mathcal L:A\cap B\in\mathcal L\}$.
Ponownie jest to układ Dynkina. Wykazana własność daje
$\mathcal P\subseteq\mathcal K_B$, skąd
$\mathcal K_B=\mathcal L$. Zatem $\mathcal L$ jest zamknięty
na skończone przekroje. Na mocy lematu \ref{lem:dynkin-intersection}
jest $\sigma$-algebrą, a więc $\sigma(\mathcal P)\subseteq\mathcal L$.
Odwrotna inkluzja wynika z minimalności $\mathcal L$, ponieważ
$\sigma(\mathcal P)$ jest układem Dynkina zawierającym $\mathcal P$.
Ostatecznie $\mathcal L=\sigma(\mathcal P)\subseteq\mathcal D$.
\end{proof}

\begin{uwaga}
Twierdzenie $\pi$--$\lambda$ pozwala dowodzić własności wszystkich
zbiorów w $\sigma(\mathcal P)$, badając je najpierw na układzie
$\pi$. W rozdziale \ref{ch:miary-wlasnosci} zastosujemy je do
jednoznaczności miar; później również do jednoznaczności miar produktowych.
\end{uwaga}

\section{Klasy monotoniczne}

\begin{definicja}[Klasa monotoniczna]
\label{def:monotone-class}
Rodzinę $\mathcal M\subseteq\mathcal P(X)$ nazywamy \emph{klasą
monotoniczną}, jeżeli jest zamknięta na granice monotoniczne zbiorów:
\begin{align*}
A_n\in\mathcal M,\ A_n\uparrow A
&\Longrightarrow A\in\mathcal M,\\
B_n\in\mathcal M,\ B_n\downarrow B
&\Longrightarrow B\in\mathcal M.
\end{align*}
Zapis $A_n\uparrow A$ oznacza $A_n\subseteq A_{n+1}$ oraz
$A=\bigcup_nA_n$, a $B_n\downarrow B$ oznacza
$B_{n+1}\subseteq B_n$ oraz $B=\bigcap_nB_n$.
\end{definicja}

Przekrój klas monotonicznych jest klasą monotoniczną. Dlatego każda
rodzina $\mathcal E\subseteq\mathcal P(X)$ wyznacza najmniejszą klasę
monotoniczną $\mathcal M(\mathcal E)$, która ją zawiera.

\begin{twierdzenie}[Twierdzenie o klasie monotonicznej]
\label{thm:monotone-class}
Niech $X$ będzie zbiorem, a $\mathcal A$ algebrą podzbiorów $X$.
Wówczas $\mathcal M(\mathcal A)=\sigma(\mathcal A)$.
\end{twierdzenie}
\begin{proof}
Każda $\sigma$-algebra jest klasą monotoniczną, stąd
$\mathcal M(\mathcal A)\subseteq\sigma(\mathcal A)$.
Niech $\mathcal M=\mathcal M(\mathcal A)$. Pokażemy, że $\mathcal M$
jest $\sigma$-algebrą.

Najpierw rozważmy rodzinę
$\mathcal C=\{E\in\mathcal M:E^c\in\mathcal M\}$.
Jest ona klasą monotoniczną: dopełnienie zamienia ciągi rosnące
na malejące i odwrotnie. Ponieważ $\mathcal A$ jest algebrą,
$\mathcal A\subseteq\mathcal C$. Z minimalności $\mathcal M$
wynika $\mathcal C=\mathcal M$, czyli $\mathcal M$ jest zamknięta
na dopełnienia.

Ustalmy $A\in\mathcal A$ i połóżmy
$\mathcal M_A=\{E\in\mathcal M:E\cup A\in\mathcal M\}$.
Suma z ustalonym $A$ zachowuje granice monotoniczne, więc
$\mathcal M_A$ jest klasą monotoniczną. Algebraiczność $\mathcal A$
daje $\mathcal A\subseteq\mathcal M_A$, a zatem
$\mathcal M_A=\mathcal M$. Otrzymaliśmy
$E\cup A\in\mathcal M$ dla każdego $E\in\mathcal M$ i $A\in\mathcal A$.
Dla ustalonego $E\in\mathcal M$ rodzina
$\mathcal N_E=\{B\in\mathcal M:E\cup B\in\mathcal M\}$
jest więc klasą monotoniczną zawierającą $\mathcal A$, zatem
$\mathcal N_E=\mathcal M$. W konsekwencji $\mathcal M$ jest zamknięta
na sumy skończone.

Jeśli $E_1,E_2,\ldots\in\mathcal M$, to ich częściowe sumy
$F_n=\bigcup_{k=1}^nE_k$ należą do $\mathcal M$ i tworzą ciąg
rosnący. Stąd $\bigcup_nE_n\in\mathcal M$. Ponieważ
$X\in\mathcal A\subseteq\mathcal M$, otrzymaliśmy $\sigma$-algebrę.
Minimalność $\sigma(\mathcal A)$ daje drugą inkluzję.
\end{proof}

\begin{uwaga}[Dwa narzędzia rozszerzania własności]
Twierdzenia \ref{thm:pi-lambda} i \ref{thm:monotone-class} są podobne,
ale ich założenia nie są zamienne: twierdzenie Dynkina rozpoczyna się
od układu $\pi$, a twierdzenie o klasie monotonicznej --- od algebry.
Później użyjemy również wersji funkcyjnej twierdzenia o klasie
monotonicznej, po wprowadzeniu całki.
\end{uwaga}

\section{Zbiory borelowskie}

\begin{definicja}[$\sigma$-algebra borelowska]
\label{def:borel}
Niech $(X,\tau)$ będzie przestrzenią topologiczną, gdzie $\tau$
stanowi rodzinę jej zbiorów otwartych. Definiujemy
$\mathcal B(X)=\sigma_X(\tau)$.
Jest to \emph{$\sigma$-algebra borelowska}; jej elementy nazywamy
\emph{zbiorami borelowskimi}. Gdy $X=\mathbb R^d$, używamy standardowej
topologii euklidesowej.
\end{definicja}

Każdy zbiór otwarty i domknięty jest borelowski, podobnie jak każdy
przeliczalny przekrój zbiorów otwartych i przeliczalna suma zbiorów
domkniętych. Zbiory tych ostatnich dwóch typów oznacza się odpowiednio
$G_\delta$ i $F_\sigma$. Nie każdy zbiór borelowski ma jedną z tych
dwóch postaci.

\begin{twierdzenie}[Przeliczalne generatory na $\mathbb R$]
\label{thm:borel-generators-R}
Dla standardowej topologii na $\mathbb R$ zachodzą równości
\[
\mathcal B(\mathbb R)
=\sigma\bigl(\{(a,b):a,b\in\mathbb Q,\ a<b\}\bigr)
=\sigma\bigl(\{(-\infty,q):q\in\mathbb Q\}\bigr)
=\sigma\bigl(\{(-\infty,q]:q\in\mathbb Q\}\bigr).
\]
\end{twierdzenie}
\begin{proof}
Wszystkie wymienione przedziały są borelowskie, więc każda z
wygenerowanych rodzin jest zawarta w $\mathcal B(\mathbb R)$.
Każdy otwarty $U\subseteq\mathbb R$ jest sumą tych przedziałów
$(a,b)$ o końcach wymiernych, które zawierają się w $U$. Takich
przedziałów jest przeliczalnie wiele; daje to pierwszą równość.

Dalej, dla $a<b$ mamy $(a,b)=(-\infty,b)\setminus(-\infty,a]$ oraz
$(-\infty,a]=\bigcap_{n\geq1}(-\infty,a+1/n)$.
Wystarczy brać wymierne $a,b$ w pierwszym wzorze; w drugim
$a+1/n$ jest wymierne dla wymiernego $a$. To dowodzi drugiej
równości. Wreszcie dla $q\in\mathbb Q$ zachodzą
$(-\infty,q]=\bigcap_{n\geq1}(-\infty,q+1/n)$ oraz
$(-\infty,q)=\bigcup_{n\geq1}(-\infty,q-1/n]$,
co daje trzecią równość.
\end{proof}

\begin{twierdzenie}[Borelowskie prostokąty]
\label{thm:borel-Rd}
Dla $d\in\mathbb N$ $\sigma$-algebra $\mathcal B(\mathbb R^d)$
jest generowana przez otwarte prostokąty
$\prod_{j=1}^d(a_j,b_j)$ o wymiernych końcach. Ponadto
\[
\mathcal B(\mathbb R^d)=\mathcal B(\mathbb R)^{\otimes d},
\]
gdzie prawa strona oznacza $\sigma$-algebrę generowaną przez
prostokąty $B_1\times\cdots\times B_d$ z $B_j\in\mathcal B(\mathbb R)$.
\end{twierdzenie}
\begin{proof}
Otwarte prostokąty o wymiernych końcach tworzą przeliczalną bazę
topologii $\mathbb R^d$. Zatem dowolny zbiór otwarty jest sumą
przeliczalnej podrodziny tej bazy, co daje pierwsze twierdzenie.
Oznaczmy przez $\mathcal F$ prawą stronę deklarowanej równości.
Każdy prostokąt bazowy należy do $\mathcal F$, skąd
$\mathcal B(\mathbb R^d)\subseteq\mathcal F$.
Z drugiej strony rzuty współrzędnych $p_j:\mathbb R^d\to\mathbb R$
są ciągłe, więc $p_j^{-1}(B_j)$ jest borelowski dla
$B_j\in\mathcal B(\mathbb R)$: rodzina zbiorów $B$ o tej
własności jest $\sigma$-algebrą zawierającą otwarte podzbiory
$\mathbb R$. Ponieważ
$\prod_{j=1}^dB_j=\bigcap_{j=1}^dp_j^{-1}(B_j)$,
każdy prostokąt borelowski należy do $\mathcal B(\mathbb R^d)$.
Stąd $\mathcal F\subseteq\mathcal B(\mathbb R^d)$.
\end{proof}

Dla dowolnego zbioru $A\subseteq X$ przez $\chi_A$ oznaczamy
funkcję charakterystyczną: $\chi_A(x)=1$ dla $x\in A$
oraz $\chi_A(x)=0$ dla $x\notin A$.

\begin{przyklad}[Racjonalne i irracjonalne]
$\mathbb Q$ jest podzbiorem borelowskim $\mathbb R$, bo jest
przeliczalną sumą singletonów, a singletony są domknięte.
Zatem $\mathbb R\setminus\mathbb Q$ również jest borelowski.
Jednak funkcja charakterystyczna $\chi_{\mathbb Q}$ jest nieciągła
w każdym punkcie $\mathbb R$. Mierzalność nie jest więc bliska
ciągłości w sensie punktowym.
\end{przyklad}

\section{Odwzorowania mierzalne i ich charakteryzacje}

\begin{definicja}[Odwzorowanie mierzalne]
\label{def:measurable-map}
Niech $(X,\mathcal F)$ i $(Y,\mathcal G)$ będą przestrzeniami
mierzalnymi. Odwzorowanie $f:X\to Y$ nazywamy
\emph{$(\mathcal F,\mathcal G)$-mierzalnym}, jeżeli
$f^{-1}(B)\in\mathcal F$ dla każdego $B\in\mathcal G$.
Jeśli $Y$ jest przestrzenią topologiczną z
$\mathcal G=\mathcal B(Y)$, mówimy krótko, że $f$ jest
\emph{borelowsko mierzalne} względem $\mathcal F$.
\end{definicja}

\begin{twierdzenie}[Kryterium przez generatory]
\label{thm:measurable-generator}
Niech $(X,\mathcal F)$ będzie przestrzenią mierzalną, $Y$ zbiorem,
$\mathcal E\subseteq\mathcal P(Y)$ rodziną zbiorów i
$\mathcal G=\sigma_Y(\mathcal E)$. Wówczas odwzorowanie
$f:X\to Y$ jest $(\mathcal F,\mathcal G)$-mierzalne wtedy i tylko
wtedy, gdy $f^{-1}(E)\in\mathcal F$ dla wszystkich $E\in\mathcal E$.
\end{twierdzenie}
\begin{proof}
Konieczność wynika z inkluzji $\mathcal E\subseteq\mathcal G$.
Dla dostateczności niech
$\mathcal H=\{B\subseteq Y:f^{-1}(B)\in\mathcal F\}$.
To $\sigma$-algebra na $Y$, ponieważ obrazy odwrotne zachowują
dopełnienia i przeliczalne sumy. Z założenia
$\mathcal E\subseteq\mathcal H$, zatem
$\mathcal G=\sigma_Y(\mathcal E)\subseteq\mathcal H$.
\end{proof}

\begin{twierdzenie}[Złożenia i funkcje ciągłe]
\label{thm:composition-measurable}
Niech $(X,\mathcal F)$, $(Y,\mathcal G)$ i $(Z,\mathcal H)$
będą przestrzeniami mierzalnymi. Jeżeli $f:X\to Y$ jest
$(\mathcal F,\mathcal G)$-mierzalne, a $g:Y\to Z$ jest
$(\mathcal G,\mathcal H)$-mierzalne, to $g\circ f$ jest
$(\mathcal F,\mathcal H)$-mierzalne.
Ponadto każde odwzorowanie ciągłe $g:Y\to Z$ pomiędzy
przestrzeniami topologicznymi jest
$(\mathcal B(Y),\mathcal B(Z))$-mierzalne.
\end{twierdzenie}
\begin{proof}
Dla $C\in\mathcal H$ mamy
$(g\circ f)^{-1}(C)=f^{-1}(g^{-1}(C))\in\mathcal F$.
Jeśli $g$ jest ciągłe, obrazy odwrotne zbiorów otwartych są otwarte;
z kryterium \ref{thm:measurable-generator} wynika jego mierzalność
borelowska.
\end{proof}

\begin{twierdzenie}[Funkcje rzeczywiste i działania]
\label{thm:measurable-real}
Niech $(X,\mathcal F)$ będzie przestrzenią mierzalną.
Dla funkcji $f:X\to\mathbb R$ następujące warunki są równoważne:
\begin{enumerate}
\item $f$ jest $(\mathcal F,\mathcal B(\mathbb R))$-mierzalna;
\item $\{x\in X:f(x)<q\}\in\mathcal F$ dla każdego $q\in\mathbb Q$;
\item $\{x\in X:f(x)\leq q\}\in\mathcal F$ dla każdego $q\in\mathbb Q$.
\end{enumerate}
Jeżeli $f,g:X\to\mathbb R$ są mierzalne, to funkcje
$f+g$, $fg$, $\max(f,g)$ i $\min(f,g)$ są mierzalne.
\end{twierdzenie}
\begin{proof}
Równoważność trzech warunków wynika z twierdzeń
\ref{thm:borel-generators-R} i \ref{thm:measurable-generator}.
Odwzorowanie $x\mapsto(f(x),g(x))$ do $\mathbb R^2$ jest mierzalne:
obrazy odwrotne prostokątów $B_1\times B_2$, gdzie
$B_1,B_2\in\mathcal B(\mathbb R)$, są przekrojami
$f^{-1}(B_1)\cap g^{-1}(B_2)$. Prostokąty te generują
$\mathcal B(\mathbb R^2)$ na mocy twierdzenia
\ref{thm:borel-Rd}. Funkcje na $\mathbb R^2$ określone przez
$(u,v)\mapsto u+v$, $(u,v)\mapsto uv$,
$(u,v)\mapsto\max(u,v)$ i $(u,v)\mapsto\min(u,v)$ są ciągłe.
Teza wynika więc z twierdzenia \ref{thm:composition-measurable}.
\end{proof}

\begin{przyklad}[Funkcje charakterystyczne zbiorów]
Dla $A\subseteq X$ definiujemy funkcję charakterystyczną
$\chi_A:X\to\mathbb R$ wzorem
$\chi_A(x)=1$ dla $x\in A$ i $\chi_A(x)=0$ dla $x\notin A$.
Wtedy $\chi_A$ jest mierzalna względem $\mathcal F$
dokładnie wtedy, gdy $A\in\mathcal F$.
Dostateczność wynika z tego, że funkcja przyjmuje tylko dwie
wartości; konieczność z równości
$A=\chi_A^{-1}((1/2,\infty))$.
\end{przyklad}

\begin{przyklad}[Mierzalność nie implikuje regularności geometrycznej]
Na $X=\mathbb R$ z $\mathcal F=\mathcal B(\mathbb R)$ funkcja
$\chi_{\mathbb Q}$ jest mierzalna, choć nie jest ciągła
w żadnym punkcie. Z drugiej strony, jeśli na nieprzeliczalnym
zbiorze $X$ przyjmiemy $\mathcal F=\{\varnothing,X\}$, to
każda mierzalna funkcja $X\to\mathbb R$ jest stała. Rzeczywiście,
gdyby przyjmowała wartości $a<b$, to dla pewnej liczby wymiernej
$q\in(a,b)$ zbiór $\{f<q\}$ byłby niepusty i różny od $X$.
\end{przyklad}

\begin{uwaga}[Co zostało, a co nie zostało wykazane]
Wprowadziliśmy pojęcie mierzalności bez definiowania jakiejkolwiek
miary. To rozróżnienie jest zasadnicze: przestrzeń mierzalna określa,
\emph{którym} zbiorom wolno przypisać rozmiar; miara określa
\emph{jaki} jest ten rozmiar. Miary pojawią się w rozdziale
\ref{ch:miary-wlasnosci}, a mierzalność funkcji o wartościach
w $[-\infty,\infty]$ omówimy szczegółowo przed zdefiniowaniem całki.
\end{uwaga}

% END chapters/pl/measure_theory/01-rozdzial.tex


% BEGIN chapters/pl/measure_theory/02-rozdzial.tex
% ROZDZIAŁ 2 — plik włączany przez main.tex
\chapter{Miary i ich podstawowe własności}
\label{ch:miary-wlasnosci}

\section{Definicja i pierwsze przykłady}

W rozdziale \ref{ch:przestrzenie-mierzalne} ustaliliśmy, które podzbiory
przestrzeni mają być mierzalne. Teraz przypiszemy im liczby z przedziału
$[0,\infty]$. Dopuszczenie wartości $\infty$ jest konieczne już dla
naturalnego pojęcia liczby elementów zbioru, a następnie dla miary
objętości na przestrzeniach nieograniczonych.

\begin{definicja}[Miara]
\label{def:measure}
Niech $(X,\mathcal F)$ będzie przestrzenią mierzalną.
\emph{Miarą} na $(X,\mathcal F)$ nazywamy funkcję
$\mu:\mathcal F\to[0,\infty]$ spełniającą:
\begin{enumerate}
\item $\mu(\varnothing)=0$;
\item dla każdego ciągu $(A_n)_{n\geq1}$ parami rozłącznych zbiorów
z $\mathcal F$ zachodzi \emph{przeliczalna addytywność}
\begin{equation}
\label{eq:countable-additivity}
\mu\left(\bigcup_{n=1}^{\infty}A_n\right)
=\sum_{n=1}^{\infty}\mu(A_n).
\end{equation}
\end{enumerate}
Trójkę $(X,\mathcal F,\mu)$ nazywamy \emph{przestrzenią z miarą}.
Wszystkie sumy liczb nieujemnych rozumiemy jako kres górny sum
częściowych, dopuszczając wartość $\infty$.
\end{definicja}

\begin{uwaga}
Nie jest wymagane, aby $\mu(X)$ było skończone. Aksjomat
$\mu(\varnothing)=0$ zapisujemy oddzielnie, ponieważ bez niego
funkcja stale równa $\infty$ spełniałaby równość
\eqref{eq:countable-additivity} dla każdej rodziny rozłącznej,
ale nie miałaby zamierzonych własności miary.
\end{uwaga}

\begin{przyklad}[Miara zerowa i miara Diraca]
\label{ex:dirac}
Na dowolnej przestrzeni mierzalnej $(X,\mathcal F)$ funkcja
$\mu(A)=0$ dla $A\in\mathcal F$ jest miarą zerową.
Dla ustalonego $x\in X$ definiujemy \emph{miarę Diraca} wzorem
\[
\delta_x(A)=\begin{cases}1,&x\in A,\\0,&x\notin A,\end{cases}
\qquad A\in\mathcal F.
\]
Jeśli $(A_n)$ są parami rozłączne, punkt $x$ należy co najwyżej
do jednego z tych zbiorów. Daje to przeliczalną addytywność.
Nie zakładamy przy tym, że singleton $\{x\}$ należy do $\mathcal F$.
\end{przyklad}

\begin{przyklad}[Miara licząca]
\label{ex:counting}
Na $(X,\mathcal P(X))$ definiujemy
\[
\#(A)=\begin{cases}
|A|,&A\text{ jest skończony},\\
\infty,&A\text{ jest nieskończony}.
\end{cases}
\]
Dla parami rozłącznych zbiorów $A_n$ liczność ich sumy jest sumą
liczności w podanym sensie: jeśli suma jest skończona, tylko
skończenie wiele $A_n$ jest niepustych i każde jest skończone;
jeśli jest nieskończona, po prawej stronie także otrzymujemy
$\infty$. Jest to \emph{miara licząca}.
\end{przyklad}

\begin{przyklad}[Miara na $\sigma$-algebrze przeliczalno-koprzeliczalnej]
\label{ex:cocountable-measure}
Niech $X$ będzie nieprzeliczalnym zbiorem i niech
$\mathcal C_X$ będzie $\sigma$-algebrą z przykładu
\ref{ex:countable-cocountable}. Określmy
\[
\mu(A)=\begin{cases}
0,&A\text{ jest przeliczalny},\\
1,&X\setminus A\text{ jest przeliczalny}.
\end{cases}
\qquad A\in\mathcal C_X.
\]
Definicja jest jednoznaczna, ponieważ nieprzeliczalny zbiór
$X$ nie może być jednocześnie przeliczalny i mieć przeliczalne
dopełnienie. Dwie rozłączne części $X$ nie mogą obie mieć
przeliczalnego dopełnienia. Dla przeliczalnej rodziny
rozłącznych zbiorów z $\mathcal C_X$ albo wszystkie mają miarę
zero, a ich suma jest przeliczalna, albo dokładnie jeden ma
miarę jeden, a suma ma miarę jeden. Jest to zatem miara.
Ten przykład pokaże w sekcji \ref{sec:atoms}, że atomy
nie muszą być singletonami.
\end{przyklad}

\begin{przyklad}[Ważone miary dyskretne]
\label{ex:weighted-discrete}
Niech $X$ będzie zbiorem, $\mathcal F=\mathcal P(X)$,
$(x_n)_{n\geq1}$ ciągiem punktów $X$, a $(w_n)_{n\geq1}$
ciągiem liczb z $[0,\infty)$. Wtedy
\[
\mu(A)=\sum_{n=1}^{\infty}w_n\chi_A(x_n),
\qquad A\subseteq X,
\]
jest miarą. Dla rozłącznych $A_k$ równość
$\chi_{\cup_kA_k}(x_n)=\sum_k\chi_{A_k}(x_n)$
pozwala zmienić kolejność dwóch szeregów o wyrazach
nieujemnych. Nie wymagamy, aby punkty $x_n$ były różne ani by
$\sum_nw_n$ było skończone.
\end{przyklad}

\section{Monotoniczność i przeliczalna podaddytywność}

\begin{twierdzenie}[Podstawowe własności miary]
\label{thm:basic-measure}
Niech $(X,\mathcal F,\mu)$ będzie przestrzenią z miarą.
Wówczas:
\begin{enumerate}
\item dla dowolnych $A,B\in\mathcal F$ z $A\subseteq B$ zachodzi
$\mu(A)\leq\mu(B)$;
\item dla dowolnych $A,B\in\mathcal F$ zachodzi
$\mu(A\cup B)\leq\mu(A)+\mu(B)$;
\item dla dowolnych $A_n\in\mathcal F$, $n\geq1$, zachodzi
\begin{equation}
\label{eq:subadditivity}
\mu\left(\bigcup_{n=1}^{\infty}A_n\right)
\leq\sum_{n=1}^{\infty}\mu(A_n);
\end{equation}
\item jeśli $A,B\in\mathcal F$, $A\subseteq B$ i $\mu(A)<\infty$,
to $\mu(B\setminus A)=\mu(B)-\mu(A)$.
\end{enumerate}
\end{twierdzenie}
\begin{proof}
Zbiory $A$ i $B\setminus A$ są rozłączne, a ich suma jest równa $B$.
Przeliczalna addytywność, zastosowana do rodziny uzupełnionej
zbiorami pustymi, daje
$\mu(B)=\mu(A)+\mu(B\setminus A)$. Wynika stąd monotoniczność,
a jeżeli $\mu(A)<\infty$, również wzór na różnicę.

Ponadto $A\cup B=A\,\dot\cup\,(B\setminus A)$, więc
$\mu(A\cup B)\leq\mu(A)+\mu(B)$.
Dla dowolnego ciągu $(A_n)$ definiujemy
$B_1=A_1$ oraz
$B_n=A_n\setminus\bigcup_{k<n}A_k$ dla $n\geq2$.
Zbiory $B_n$ są parami rozłączne, $B_n\subseteq A_n$
i $\bigcup_nB_n=\bigcup_nA_n$. Stąd
$\mu(\bigcup_nA_n)=\sum_n\mu(B_n)\leq\sum_n\mu(A_n)$.
\end{proof}

\begin{uwaga}[Arytmetyka z nieskończonością]
Warunku $\mu(A)<\infty$ w ostatnim punkcie twierdzenia
\ref{thm:basic-measure} nie wolno usuwać. Jeśli
$\mu(A)=\mu(B)=\infty$, wyrażenie $\infty-\infty$
nie ma określonej wartości. Prawidłowym wzorem pozostaje
$\mu(B)=\mu(A)+\mu(B\setminus A)$, który nie wymaga odejmowania.
\end{uwaga}

\begin{wniosek}[Przeliczalne sumy zbiorów miary zero]
\label{cor:null-union}
Jeżeli $(X,\mathcal F,\mu)$ jest przestrzenią z miarą oraz
$N_n\in\mathcal F$ i $\mu(N_n)=0$ dla każdego $n\in\mathbb N$,
to $\mu(\bigcup_nN_n)=0$.
\end{wniosek}
\begin{proof}
Z podaddytywności mamy
$0\leq\mu(\bigcup_nN_n)\leq\sum_n\mu(N_n)=0$.
\end{proof}

\section{Ciągłość miary względem granic zbiorów}

\begin{twierdzenie}[Ciągłość z dołu i z góry]
\label{thm:measure-continuity}
Niech $(X,\mathcal F,\mu)$ będzie przestrzenią z miarą.
\begin{enumerate}
\item Jeśli $A_n\in\mathcal F$ oraz $A_n\uparrow A$, to
$A\in\mathcal F$ i
\begin{equation}
\label{eq:continuity-below}
\mu(A)=\lim_{n\to\infty}\mu(A_n).
\end{equation}
\item Jeśli $B_n\in\mathcal F$, $B_n\downarrow B$ oraz
$\mu(B_1)<\infty$, to $B\in\mathcal F$ i
\begin{equation}
\label{eq:continuity-above}
\mu(B)=\lim_{n\to\infty}\mu(B_n).
\end{equation}
\end{enumerate}
\end{twierdzenie}
\begin{proof}
Dla ciągu rosnącego przyjmujemy $E_1=A_1$ i
$E_n=A_n\setminus A_{n-1}$ dla $n\geq2$.
Zbiory $E_n$ są parami rozłączne oraz
$A=\bigcup_nE_n$ i $A_k=\bigcup_{n=1}^kE_n$.
Stąd
$\mu(A)=\sum_n\mu(E_n)=\lim_{k\to\infty}\sum_{n=1}^k\mu(E_n)
=\lim_k\mu(A_k)$.

Dla ciągu malejącego zbiory $C_n=B_1\setminus B_n$ rosną
do $B_1\setminus B$. Z pierwszej części mamy
$\mu(C_n)\to\mu(B_1\setminus B)$. Ponieważ
$\mu(B_1)<\infty$, wolno zastosować wzór na różnicę:
$\mu(C_n)=\mu(B_1)-\mu(B_n)$ oraz
$\mu(B_1\setminus B)=\mu(B_1)-\mu(B)$.
Po odjęciu od $\mu(B_1)$ otrzymujemy żądaną równość.
\end{proof}

\begin{uwaga}[Wystarczy skończoność miary jednego wyrazu]
W drugim punkcie twierdzenia~\ref{thm:measure-continuity} można
zastąpić założenie $\mu(B_1)<\infty$ słabszym
$\mu(B_m)<\infty$ dla pewnego $m\ge1$: wystarczy
zastosować twierdzenie do ogona ciągu
$(B_{m+k-1})_{k\ge1}$, którego przekrój jest nadal równy $B$.
Nie wystarcza natomiast sama skończoność $\mu(B)$.
\end{uwaga}

\begin{przyklad}[Dlaczego założenie skończoności jest niezbędne]
\label{ex:continuity-counterexample}
Na $X=\mathbb N$ z miarą liczącą zbiory
$B_n=\{n,n+1,n+2,\ldots\}$ tworzą ciąg malejący o pustym
przekroju. Każdy $B_n$ ma miarę $\infty$, ale
$\mu(\bigcap_nB_n)=\mu(\varnothing)=0$.
Tak więc ciągłość z góry nie zachodzi bez odpowiedniej
hipotezy skończoności.
\end{przyklad}

\begin{definicja}[Granica górna i dolna zbiorów]
\label{def:limsup-sets}
Dla ciągu $(A_n)_{n\geq1}$ podzbiorów $X$ definiujemy
\[
\liminf_{n\to\infty}A_n
=\bigcup_{m=1}^{\infty}\bigcap_{n\geq m}A_n,
\qquad
\limsup_{n\to\infty}A_n
=\bigcap_{m=1}^{\infty}\bigcup_{n\geq m}A_n.
\]
Pierwszy zbiór składa się z punktów należących do wszystkich
$A_n$ od pewnego miejsca, drugi --- z punktów należących do
nieskończenie wielu $A_n$. Jeśli oba zbiory są równe,
ich wspólną wartość oznaczamy przez $\lim_nA_n$.
\end{definicja}

\begin{twierdzenie}[Nierówności dla granic zbiorów]
\label{thm:fatou-for-sets}
Niech $(X,\mathcal F,\mu)$ będzie przestrzenią z miarą
oraz $A_n\in\mathcal F$ dla wszystkich $n\ge1$.
Wówczas
\[
 \mu(\liminf_{n\to\infty}A_n)
 \le \liminf_{n\to\infty}\mu(A_n).
\]
Jeżeli dodatkowo $\mu(X)<\infty$, to
\[
 \limsup_{n\to\infty}\mu(A_n)
 \le \mu(\limsup_{n\to\infty}A_n).
\]
\end{twierdzenie}
\begin{proof}
Niech $C_m=\bigcap_{n\ge m}A_n$.
Zbiory $C_m$ tworzą ciąg rosnący,
$\bigcup_m C_m=\liminf_n A_n$ oraz
$\mu(C_m)\le\inf_{n\ge m}\mu(A_n)$.
Ciągłość miary z dołu daje
$\mu(\liminf_n A_n)=\lim_m\mu(C_m)
\le\lim_m\inf_{n\ge m}\mu(A_n)$.
Dla drugiej nierówności kładziemy
$D_m=\bigcup_{n\ge m}A_n$.
Zbiory $D_m$ maleją do $\limsup_n A_n$,
a $\mu(D_m)\ge\sup_{n\ge m}\mu(A_n)$.
Ponieważ $\mu(D_1)\le\mu(X)<\infty$,
ciągłość miary z góry pozwala przejść do granicy:
$\mu(\limsup_n A_n)=\lim_m\mu(D_m)
\ge\lim_m\sup_{n\ge m}\mu(A_n)$.
\end{proof}
\begin{uwaga}
Założenie skończoności jest istotne w drugiej nierówności.
Dla miary liczącej na $\mathbb N$ i
$A_n=\{n,n+1,\ldots\}$ mamy
$\limsup_n\mu(A_n)=\infty$, lecz
$\limsup_n A_n=\varnothing$.
\end{uwaga}

\begin{twierdzenie}[Lemat Borela--Cantellego, część pierwsza]
\label{thm:borel-cantelli-1}
Niech $(X,\mathcal F,\mu)$ będzie przestrzenią z miarą
oraz $(A_n)_{n\geq1}$ ciągiem zbiorów z $\mathcal F$.
Jeżeli $\sum_{n=1}^{\infty}\mu(A_n)<\infty$, to
$\mu(\limsup_{n\to\infty}A_n)=0$.
\end{twierdzenie}
\begin{proof}
Zbiór $\limsup_nA_n$ jest mierzalny i dla każdego
$m\in\mathbb N$ zawiera się w $\bigcup_{n\geq m}A_n$.
Z monotoniczności i podaddytywności wynika
\[
0\leq\mu(\limsup_nA_n)
\leq\sum_{n\geq m}\mu(A_n).
\]
Ponieważ ogólny szereg jest zbieżny, jego ogon dąży do zera,
gdy $m\to\infty$. To kończy dowód. Nie użyliśmy ani
skończoności $\mu(X)$, ani niezależności zbiorów.
\end{proof}

\begin{przyklad}[Interpretacja]
Jeśli $\mu(A_n)\leq2^{-n}$ dla wszystkich $n\geq1$, to
$\mu(\limsup_nA_n)=0$. Prawie każdy punkt (w sensie, który
zaraz zdefiniujemy) może należeć tylko do skończenie wielu $A_n$.
W rachunku prawdopodobieństwa po wprowadzeniu niezależności
uzyskuje się drugą część lematu Borela--Cantellego, z innymi
założeniami i przeciwnym wnioskiem.
\end{przyklad}

\section{Skończoność, \texorpdfstring{$\sigma$}{sigma}-skończoność i półskończoność}

\begin{definicja}[Klasy miar]
\label{def:measure-finiteness}
Niech $(X,\mathcal F,\mu)$ będzie przestrzenią z miarą.
\begin{enumerate}
\item Miara $\mu$ jest \emph{skończona}, gdy $\mu(X)<\infty$.
\item Miara $\mu$ jest \emph{probabilistyczna}, gdy $\mu(X)=1$.
\item Miara $\mu$ jest \emph{$\sigma$-skończona}, gdy istnieją
$E_n\in\mathcal F$, $n\geq1$, takie że
$X=\bigcup_nE_n$ i $\mu(E_n)<\infty$ dla wszystkich $n$.
\item Miara $\mu$ jest \emph{półskończona}, gdy dla każdego
$A\in\mathcal F$ z $\mu(A)>0$ istnieje $B\in\mathcal F$
takie, że $B\subseteq A$ oraz $0<\mu(B)<\infty$.
\end{enumerate}
Jeśli $(X,\tau)$ jest przestrzenią topologiczną, a
$\mathcal F\supseteq\mathcal B(X)$, nazywamy $\mu$
\emph{lokalnie skończoną}, gdy każdy $x\in X$ posiada
otwarte otoczenie $U$ o $\mu(U)<\infty$.
\end{definicja}

\begin{twierdzenie}[Implikacje]
\label{thm:finiteness-implications}
Każda miara skończona na przestrzeni mierzalnej jest
$\sigma$-skończona, a każda miara $\sigma$-skończona
jest półskończona. Implikacje odwrotne są na ogół fałszywe.
\end{twierdzenie}
\begin{proof}
Jeśli $\mu(X)<\infty$, można w definicji $\sigma$-skończoności
przyjąć $E_n=X$. Niech teraz $X=\bigcup_nE_n$ z
$\mu(E_n)<\infty$. Dla $A\in\mathcal F$ o $\mu(A)>0$
nie wszystkie przekroje $A\cap E_n$ mogą mieć miarę zero:
w przeciwnym razie podaddytywność dałaby
$\mu(A)\leq\sum_n\mu(A\cap E_n)=0$.
Istnieje więc $n$ z $0<\mu(A\cap E_n)<\infty$.

Miara licząca na $\mathbb N$ jest $\sigma$-skończona,
ale nieskończona. Miara licząca na nieprzeliczalnym
zbiorze $X$ jest półskończona (każdy niepusty zbiór
zawiera singleton miary jeden), ale nie jest
$\sigma$-skończona, gdyż zbiory o skończonej mierze
są skończone, a przeliczalna suma zbiorów skończonych
jest przeliczalna.
\end{proof}

\begin{przyklad}[Miara, która nie jest półskończona]
Niech $X$ będzie niepustym zbiorem z $\mathcal F=\mathcal P(X)$
i niech $\mu(\varnothing)=0$ oraz $\mu(A)=\infty$
dla $A\ne\varnothing$. Jest to miara: dla rodziny
rozłącznej z niepustą sumą co najmniej jeden składnik
ma miarę $\infty$. Nie ma jednak żadnego podzbioru
o mierze dodatniej i skończonej.
\end{przyklad}

\begin{uwaga}[Lokalna skończoność a $\sigma$-skończoność]
Na przestrzeni topologicznej mającej przeliczalną bazę
lokalnie skończona miara borelowska jest $\sigma$-skończona.
Istotnie, otwarte otoczenia o skończonej mierze tworzą
pokrycie; własność Lindelöfa (wynikająca z istnienia
przeliczalnej bazy) pozwala wybrać pokrycie przeliczalne.
Bez założenia topologicznego takiej implikacji nie ma:
miara licząca na nieprzeliczalnej przestrzeni dyskretnej
jest lokalnie skończona, lecz nie jest $\sigma$-skończona.
\end{uwaga}

\section{Zbiory miary zero, własności prawie wszędzie i zupełność}

\begin{definicja}[Własność prawie wszędzie]
\label{def:almost-everywhere}
Niech $(X,\mathcal F,\mu)$ będzie przestrzenią z miarą.
Mówimy, że własność $P(x)$ zachodzi \emph{$\mu$-prawie wszędzie}
(skrót: $\mu$-p.w.), jeśli istnieje $N\in\mathcal F$
o $\mu(N)=0$ takie, że $P(x)$ zachodzi dla każdego
$x\in X\setminus N$. Nie wymagamy mierzalności zbioru
$\{x\in X:P(x)\text{ nie zachodzi}\}$.
Jeśli jest on mierzalny, jest to równoważne temu,
że ma miarę zero.
\end{definicja}

\begin{uwaga}
Dzięki wnioskowi \ref{cor:null-union} przeliczalnie wiele
własności zachodzących $\mu$-p.w. zachodzi jednocześnie
$\mu$-p.w. Bez przeliczalności nie można wyciągać tego
wniosku: przy miarze z przykładu \ref{ex:cocountable-measure}
każdy singleton ma miarę zero, lecz suma wszystkich singletonów
jest równa $X$ i ma miarę jeden.
\end{uwaga}

\begin{definicja}[Miara zupełna]
\label{def:complete-measure}
Miara $\mu$ na $(X,\mathcal F)$ jest \emph{zupełna}, jeżeli
dla każdego $N\in\mathcal F$ spełniającego $\mu(N)=0$
każdy podzbiór $A\subseteq N$ należy do $\mathcal F$.
Z monotoniczności wynika wówczas $\mu(A)=0$.
\end{definicja}

\begin{przyklad}[Miara niezupełna na trzech punktach]
\label{ex:incomplete-measure}
Niech $X=\{a,b,c\}$ oraz
$\mathcal F=\{\varnothing,\{a,b\},\{c\},X\}$.
Określmy $\mu(\{a,b\})=0$ i $\mu(\{c\})=1$,
a pozostałe wartości wyznaczmy przez addytywność:
$\mu(\varnothing)=0$, $\mu(X)=1$.
To dobrze określona miara na $\mathcal F$.
Zbiór $N=\{a,b\}$ ma miarę zero, ale jego podzbiór
$\{a\}$ nie jest mierzalny. Zatem $\mu$ nie jest zupełna.
\end{przyklad}

\begin{twierdzenie}[Uzupełnienie przestrzeni z miarą]
\label{thm:completion}
Niech $(X,\mathcal F,\mu)$ będzie przestrzenią z miarą.
Zdefiniujmy
\[
\overline{\mathcal F}
=\{E\subseteq X:\text{istnieją }A,B\in\mathcal F
\text{ takie, że }A\subseteq E\subseteq B
\text{ oraz }\mu(B\setminus A)=0\}.
\]
Wówczas $\overline{\mathcal F}$ jest $\sigma$-algebrą zawierającą
$\mathcal F$, a wzór
\[
\overline\mu(E)=\mu(A),\qquad
A\subseteq E\subseteq B,\quad A,B\in\mathcal F,
\quad\mu(B\setminus A)=0,
\]
definiuje na niej miarę zupełną. Miara $\overline\mu$
przedłuża $\mu$, a $\overline{\mathcal F}$ jest najmniejszą
$\sigma$-algebrą zawierającą $\mathcal F$ i wszystkie
podzbiory zbiorów z $\mathcal F$ o mierze zero.
\end{twierdzenie}
\begin{proof}
Najpierw sprawdzamy jednoznaczność definicji $\overline\mu$.
Jeśli $A_i\subseteq E\subseteq B_i$ i
$\mu(B_i\setminus A_i)=0$ dla $i=1,2$, to
$A_1\setminus A_2\subseteq B_2\setminus A_2$.
Z monotoniczności mamy $\mu(A_1\setminus A_2)=0$,
a stąd $\mu(A_1)\leq\mu(A_2)$, nawet jeśli
jedna z miar jest nieskończona. Zamiana indeksów
daje $\mu(A_1)=\mu(A_2)$. Ponadto
$\mu(A_i)=\mu(B_i)$, gdyż
$B_i=A_i\,\dot\cup\,(B_i\setminus A_i)$.

Jeśli $E\in\overline{\mathcal F}$ i
$A\subseteq E\subseteq B$ są jak w tezie, to
$X\setminus B\subseteq X\setminus E\subseteq X\setminus A$,
a różnica zbiorów skrajnych wynosi $B\setminus A$
i ma miarę zero. To dowodzi zamknięcia na dopełnienia.
Jeśli $E_n\in\overline{\mathcal F}$ i
$A_n\subseteq E_n\subseteq B_n$ z
$\mu(B_n\setminus A_n)=0$, to
\[
\bigcup_nA_n\subseteq\bigcup_nE_n\subseteq\bigcup_nB_n,
\qquad
\left(\bigcup_nB_n\right)\setminus\left(\bigcup_nA_n\right)
\subseteq\bigcup_n(B_n\setminus A_n).
\]
Ostatni zbiór ma miarę zero na mocy wniosku
\ref{cor:null-union}, więc przeliczalna suma należy do
$\overline{\mathcal F}$. Zawieranie $\mathcal F$ wynika
z wyboru $A=B=E\in\mathcal F$.

Dla parami rozłącznych $E_n\in\overline{\mathcal F}$
wybieramy $A_n\subseteq E_n$ jak wyżej. Zbiory $A_n$
są parami rozłączne i
$\overline\mu(\bigcup_nE_n)=\mu(\bigcup_nA_n)
=\sum_n\mu(A_n)=\sum_n\overline\mu(E_n)$.
Zatem $\overline\mu$ jest miarą i przedłuża $\mu$.

Aby wykazać zupełność, niech $E\in\overline{\mathcal F}$
i $\overline\mu(E)=0$. Wybierzmy $A,B\in\mathcal F$
takie, że $A\subseteq E\subseteq B$ i
$\mu(B\setminus A)=0$. Mamy $\mu(B)=\mu(A)=0$.
Dla dowolnego $T\subseteq E$ zachodzi
$\varnothing\subseteq T\subseteq B$, przy $\mu(B)=0$,
więc $T\in\overline{\mathcal F}$.

Wreszcie, jeśli $\mathcal G$ jest $\sigma$-algebrą zawierającą
$\mathcal F$ i wszystkie podzbiory jej zbiorów miary zero,
to dla $E\in\overline{\mathcal F}$ jak wyżej zbiór
$E\setminus A$ jest podzbiorem mierzalnego zbioru
$B\setminus A$ miary zero. Zatem $E\setminus A\in\mathcal G$
i $E=A\cup(E\setminus A)\in\mathcal G$. Dowodzi to minimalności.
\end{proof}

\begin{uwaga}[Jednoznaczność uzupełnienia]
Konstrukcja z twierdzenia \ref{thm:completion} nie polega
na arbitralnym wyborze wartości dla nowych podzbiorów.
Zbiory nowo dołączone różnią się od starych zbiorów
mierzalnych jedynie o części zbiorów miary zero.
Rozróżnienie między miarą borelowską a jej uzupełnieniem
będzie ważne przy konstrukcji miary Lebesgue'a.
\end{uwaga}

\section{Operacje na miarach}

\begin{twierdzenie}[Sumy i mnożenie przez stałą]
\label{thm:measure-operations}
Niech $(X,\mathcal F)$ będzie przestrzenią mierzalną,
$(\mu_n)_{n\geq1}$ ciągiem miar na $\mathcal F$,
a $(c_n)_{n\geq1}$ ciągiem liczb nieujemnych i skończonych.
Wówczas
$\mu(A)=\sum_{n=1}^{\infty}c_n\mu_n(A)$ dla $A\in\mathcal F$
jest miarą, przy konwencji, że $0\mu_n$ oznacza miarę
zerową również wtedy, gdy $\mu_n(A)=\infty$.
\end{twierdzenie}
\begin{proof}
Wszystkie składniki są nieujemne i wartość na zbiorze pustym
wynosi zero. Dla parami rozłącznych $A_k$ otrzymujemy
\[
\mu\left(\bigcup_kA_k\right)
=\sum_n c_n\sum_k\mu_n(A_k)
=\sum_k\sum_n c_n\mu_n(A_k)
=\sum_k\mu(A_k).
\]
Wolno przestawić kolejność sumowania, ponieważ obie strony
środkowej równości są kresami górnymi tych samych sum
skończonych wyrazów nieujemnych.
\end{proof}

\begin{definicja}[Miara śladowa i ograniczenie]
\label{def:restriction}
Niech $(X,\mathcal F,\mu)$ będzie przestrzenią z miarą
i $E\in\mathcal F$. Na przestrzeni mierzalnej
$(E,\mathcal F|_E)$ definiujemy \emph{miarę śladową}
$\mu|_E$ wzorem
$(\mu|_E)(B)=\mu(B)$ dla $B\in\mathcal F|_E$.
Ponieważ $E\in\mathcal F$, każdy zbiór śladowy
$B=A\cap E$ jest elementem $\mathcal F$.
Można też określić na $(X,\mathcal F)$ miarę
$\mu_E(A)=\mu(A\cap E)$; obie konstrukcje mają
różne dziedziny i nie należy ich utożsamiać.
\end{definicja}

\begin{proof}[Sprawdzenie aksjomatów]
Miary obu typów przypisują zbiorowi pustemu wartość zero.
Jeśli $B_n\subseteq E$ są parami rozłączne, to ich
suma jest rozłączna również w $X$, a przeliczalna
addytywność $\mu$ daje addytywność $\mu|_E$.
Dla $\mu_E$ korzystamy z tożsamości
$(\bigcup_nA_n)\cap E=\bigcup_n(A_n\cap E)$.
\end{proof}

\begin{definicja}[Miara obrazowa]
\label{def:pushforward}
Niech $(X,\mathcal F,\mu)$ będzie przestrzenią z miarą,
$(Y,\mathcal G)$ przestrzenią mierzalną, a
$T:X\to Y$ odwzorowaniem $(\mathcal F,\mathcal G)$-mierzalnym.
\emph{Miarą obrazową} $T_\#\mu$ na $(Y,\mathcal G)$ nazywamy
miarę określoną wzorem
\[
(T_\#\mu)(B)=\mu(T^{-1}(B)),\qquad B\in\mathcal G.
\]
\end{definicja}

\begin{proof}[Sprawdzenie aksjomatów]
Z mierzalności $T$ wszystkie obrazy odwrotne występujące
we wzorze są mierzalne. Obraz odwrotny zbioru pustego
jest pusty. Dla parami rozłącznych $B_n$ zbiory
$T^{-1}(B_n)$ są parami rozłączne i
$T^{-1}(\bigcup_nB_n)=\bigcup_nT^{-1}(B_n)$.
Równość przeliczalnej addytywności jest zatem konsekwencją
tego samego aksjomatu dla $\mu$.
\end{proof}

\begin{przyklad}[Obraz miary Diraca]
Jeżeli $x\in X$ i $T:X\to Y$ jest mierzalne, to
$T_\#\delta_x=\delta_{T(x)}$.
W szczególności obraz miary dyskretnej
$\sum_nw_n\delta_{x_n}$ jest równy
$\sum_nw_n\delta_{T(x_n)}$; punkty o tym samym obrazie
mogą powodować sumowanie wag.
\end{przyklad}

\begin{przyklad}[Normalizacja]
Jeżeli miara $\mu$ na $(X,\mathcal F)$ spełnia
$0<\mu(X)<\infty$, to $\mathbb P(A)=\mu(A)/\mu(X)$
jest miarą probabilistyczną. Nie można tego zrobić
w ten sposób dla miary o całkowitej masie nieskończonej.
\end{przyklad}

\section{Jednoznaczność miar na układach generujących}

\begin{twierdzenie}[Jednoznaczność miar skończonych]
\label{thm:uniqueness-finite}
Niech $X$ będzie zbiorem, $\mathcal P$ układem $\pi$
podzbiorów $X$, a $\mathcal F=\sigma(\mathcal P)$.
Jeżeli $\mu,\nu$ są miarami skończonymi na $(X,\mathcal F)$
spełniającymi $\mu(X)=\nu(X)$ i
$\mu(P)=\nu(P)$ dla każdego $P\in\mathcal P$,
to $\mu=\nu$ na całej $\mathcal F$.
\end{twierdzenie}
\begin{proof}
Niech $\mathcal D=\{A\in\mathcal F:\mu(A)=\nu(A)\}$.
Z założenia $X\in\mathcal D$ i $\mathcal P\subseteq\mathcal D$.
Jeżeli $A,B\in\mathcal D$ i $A\subseteq B$, to
$\mu(B\setminus A)=\mu(B)-\mu(A)
=\nu(B)-\nu(A)=\nu(B\setminus A)$.
Odejmowanie jest dopuszczalne, gdyż obie miary są skończone.
Dla parami rozłącznych $A_n\in\mathcal D$ mamy
$\mu(\bigcup_nA_n)=\sum_n\mu(A_n)
=\sum_n\nu(A_n)=\nu(\bigcup_nA_n)$.
Zatem $\mathcal D$ jest układem Dynkina zawierającym
$\mathcal P$. Twierdzenie \ref{thm:pi-lambda} daje
$\mathcal F\subseteq\mathcal D$.
\end{proof}

\begin{twierdzenie}[Jednoznaczność przy wspólnym wyczerpaniu]
\label{thm:uniqueness-sigma-finite}
Niech $X$ będzie zbiorem, $\mathcal P$ układem $\pi$ na $X$,
$\mathcal F=\sigma(\mathcal P)$, a $\mu,\nu$ miarami
na $(X,\mathcal F)$ zgodnymi na $\mathcal P$.
Załóżmy, że istnieją zbiory $E_n\in\mathcal P$ takie, że
$E_n\subseteq E_{n+1}$, $\bigcup_nE_n=X$ oraz
$\mu(E_n)=\nu(E_n)<\infty$ dla każdego $n$.
Wtedy $\mu=\nu$ na $\mathcal F$.
\end{twierdzenie}
\begin{proof}
Dla każdego $n$ definiujemy miary skończone na $\mathcal F$:
$\mu_n(A)=\mu(A\cap E_n)$ i
$\nu_n(A)=\nu(A\cap E_n)$.
Ponieważ $E_n\in\mathcal P$ i $\mathcal P$ jest układem $\pi$,
dla każdego $P\in\mathcal P$ zachodzi
$P\cap E_n\in\mathcal P$, a więc
$\mu_n(P)=\nu_n(P)$. Ponadto
$\mu_n(X)=\nu_n(X)<\infty$. Z twierdzenia
\ref{thm:uniqueness-finite} wynika $\mu_n=\nu_n$.
Dla dowolnego $A\in\mathcal F$ zbiory
$A\cap E_n$ rosną do $A$. Ciągłość miary z dołu
daje
$\mu(A)=\lim_n\mu(A\cap E_n)
=\lim_n\nu(A\cap E_n)=\nu(A)$.
\end{proof}

\begin{przyklad}[Miara określona wartościami dystrybuanty]
\label{ex:unique-cdf}
Niech $\mu$ i $\nu$ będą miarami probabilistycznymi
na $(\mathbb R,\mathcal B(\mathbb R))$.
Jeśli dla każdego $t\in\mathbb R$ zachodzi
$\mu((-\infty,t])=\nu((-\infty,t])$, to $\mu=\nu$.
Wystarczy bowiem zastosować twierdzenie
\ref{thm:uniqueness-finite} do układu $\pi$
$\{(-\infty,t]:t\in\mathbb R\}$.
Układ ten generuje $\mathcal B(\mathbb R)$,
a równość całkowitych mas wynika z probabilistyczności.
To wyjaśnia, dlaczego dystrybuanta jednoznacznie
wyznacza rozkład zmiennej losowej.
\end{przyklad}

\begin{przyklad}[Dlaczego wartości $\infty$ nie wystarczają]
\label{ex:uniqueness-fails}
Na $X=\mathbb R$ wyposażonym w $\mathcal B(\mathbb R)$
niech $\mu$ będzie miarą liczącą, a $\nu=2\mu$.
Obie miary przypisują wartość $\infty$ każdemu niepustemu
otwartemu przedziałowi o końcach wymiernych oraz wartość
zero zbiorowi pustemu. Zgadzają się zatem na układzie
$\pi$ tych przedziałów (razem ze zbiorem pustym),
który generuje $\mathcal B(\mathbb R)$.
Jednak $\mu(\{0\})=1$, a $\nu(\{0\})=2$.
Nie ma tu wspólnego wyczerpania zbiorami z układu $\pi$
o skończonych miarach. Przykład dowodzi, że zgodność
na generatorach bez odpowiednich warunków skończoności
nie wyznacza miary jednoznacznie.
\end{przyklad}

\section{Atomy miar}
\label{sec:atoms}

\begin{definicja}[Atom i miara bezatomowa]
\label{def:atom}
Niech $(X,\mathcal F,\mu)$ będzie przestrzenią z miarą.
Zbiór $A\in\mathcal F$ nazywamy \emph{atomem miary $\mu$},
jeżeli $\mu(A)>0$ oraz dla każdego $B\in\mathcal F$
zawartego w $A$ zachodzi $\mu(B)=0$ lub
$\mu(A\setminus B)=0$.
Miarę $\mu$ nazywamy \emph{bezatomową}, jeśli nie ma atomów.
\end{definicja}

\begin{przyklad}[Atomy miary dyskretnej]
Dla miary Diraca $\delta_x$ każdy mierzalny zbiór
zawierający $x$ jest atomem; wszystkie takie zbiory
są równoważne z dokładnością do zbiorów miary zero.
Dla miary liczącej na $\mathcal P(X)$ atomami są
dokładnie singletony: zbiór liczący co najmniej dwa punkty
można rozłożyć na dwie niepuste mierzalne części
o dodatniej mierze.
\end{przyklad}

\begin{przyklad}[Atom, który nie jest singletonem]
W przestrzeni z przykładu \ref{ex:cocountable-measure}
cały zbiór $X$ jest atomem: jego mierzalne podzbiory
mają miarę zero albo ich dopełnienia mają miarę zero.
Jednocześnie każdy singleton ma miarę zero.
Zatem nawet miara probabilistyczna może być atomowa,
chociaż żaden punkt nie ma dodatniej miary. Wniosek
``miara bez atomów wtedy i tylko wtedy, gdy wszystkie
singletony mają miarę zero'' jest bez dodatkowych
założeń fałszywy.
\end{przyklad}

\begin{twierdzenie}[Atomy miar $\sigma$-skończonych mają skończoną masę]
\label{thm:atom-sigma-finite}
Jeżeli $\mu$ jest miarą $\sigma$-skończoną na
$(X,\mathcal F)$ i $A\in\mathcal F$ jest atomem $\mu$,
to $0<\mu(A)<\infty$.
\end{twierdzenie}
\begin{proof}
Niech $X=\bigcup_nE_n$ z $\mu(E_n)<\infty$.
Gdyby wszystkie $A\cap E_n$ miały miarę zero,
to $\mu(A)=0$ na mocy przeliczalnej podaddytywności.
Istnieje więc $n$ z $\mu(A\cap E_n)>0$.
Atomowość $A$ implikuje $\mu(A\setminus E_n)=0$.
Otrzymujemy
$\mu(A)=\mu(A\cap E_n)+\mu(A\setminus E_n)
=\mu(A\cap E_n)\leq\mu(E_n)<\infty$.
\end{proof}

\begin{wniosek}[Przeliczalność rozłącznych atomów]
\label{cor:countably-many-disjoint-atoms}
Niech $\mu$ będzie miarą $\sigma$-skończoną na
$(X,\mathcal F)$, a $(A_i)_{i\in I}$ rodziną parami rozłącznych
atomów miary $\mu$. Wtedy zbiór indeksów $I$ jest co najwyżej
przeliczalny.
\end{wniosek}
\begin{proof}
Wybierzmy $X=\bigcup_{n\ge1}E_n$ z $\mu(E_n)<\infty$.
Z twierdzenia~\ref{thm:atom-sigma-finite} każdy $A_i$ ma
miarę dodatnią i skończoną. Istnieje więc $n$ takie, że
$\mu(A_i\cap E_n)>0$; z atomowości wynika wówczas
$\mu(A_i\cap E_n)=\mu(A_i)$. Dla ustalonych $n,k\ge1$
rodzina tych indeksów $i$, dla których
$\mu(A_i\cap E_n)\ge1/k$, jest skończona:
w przeciwnym przypadku suma miar dowolnie wielu
rozłącznych fragmentów przekroczyłaby $\mu(E_n)$.
Każdy indeks należy do przynajmniej jednej takiej
rodziny, więc $I$ jest przeliczalną sumą zbiorów
skończonych.
\end{proof}

\begin{uwaga}[Co daje $\sigma$-skończoność?]
$\sigma$-skończoność nie jest jedynie wygodnym
wzmocnieniem skończoności. Stanowi istotne założenie
jednoznaczności rozszerzenia miary, wielu wersji
twierdzenia Fubiniego oraz twierdzenia Radona--Nikodyma.
Nie należy jednak automatycznie zastępować
$\sigma$-skończoności półskończonością: są to różne hipotezy.
\end{uwaga}

\section{Synteza rozdziału i perspektywa dalszych konstrukcji}

Z rozdziałów \ref{ch:przestrzenie-mierzalne} i
\ref{ch:miary-wlasnosci} otrzymaliśmy następujący obraz.
$\sigma$-algebra kontroluje stabilność zbiorów mierzalnych
przy granicach przeliczalnych, zaś miara przekształca
rozłączne sumy w szeregi liczb nieujemnych. Te dwa
aksjomaty wystarczają do otrzymania monotoniczności,
podaddytywności, ciągłości z dołu i z góry oraz pierwszych
argumentów typu ``prawie wszędzie''.

Pozostaje jednak zasadnicze pytanie konstrukcyjne:
jak otrzymać miarę o zadanych wartościach na prostych
zbiorach, na przykład na prostokątach w $\mathbb R^d$?
Twierdzenia o jednoznaczności mówią, że pewne dane
wyznaczają miarę co najwyżej jedną; nie zapewniają
jej istnienia. W następnym rozdziale przejdziemy
do miar zewnętrznych i konstrukcji Carathéodory'ego,
która pozwoli przejść od prostych danych do prawdziwej
przeliczalnie addytywnej miary.

% END chapters/pl/measure_theory/02-rozdzial.tex


% BEGIN chapters/pl/measure_theory/03-rozdzial.tex
% ROZDZIAŁ 3 — plik włączany przez main.tex
\chapter{Miary zewnętrzne i konstrukcja Carathéodory'ego}
\label{ch:caratheodory}

\section{Problem konstrukcji miary}

W rozdziale \ref{ch:miary-wlasnosci} badaliśmy miarę już określoną na
$\sigma$-algebrze. Teraz odwrócimy pytanie: jeśli znamy ``rozmiary''
prostych zbiorów, to jak zbudować miarę na znacznie większej rodzinie?
Przykładowo na prostej chcemy przypisać przedziałowi $[a,b)$ długość
$b-a$, a następnie mierzyć zbiory otrzymane przez przeliczalne operacje.
Sam warunek skończonej addytywności jest za słaby. Kluczową rolę odgrywa
przeliczalna addytywność oraz możliwość pokrywania dowolnego zbioru
zbiorami prostszymi.

W całym rozdziale $X$ oznacza ustalony zbiór, a wszystkie funkcje
zbiorów przyjmują wartości w $[0,\infty]$, jeśli wyraźnie nie
zaznaczono inaczej. Przyjmujemy $\inf\varnothing=\infty$ oraz
$\sum_{n=1}^\infty a_n=\sup_{N\ge1}\sum_{n=1}^N a_n$ dla $a_n\ge0$.

\section{Miara zewnętrzna: definicja i własności}

\begin{definicja}[Miara zewnętrzna]
\label{def:outer-measure}
Funkcję $\mu^*:\mathcal P(X)\to[0,\infty]$ nazywamy
\emph{miarą zewnętrzną}, jeśli:
\begin{enumerate}[label=\textup{(\roman*)}]
\item $\mu^*(\varnothing)=0$;
\item $A\subseteq B$ implikuje $\mu^*(A)\le\mu^*(B)$;
\item dla dowolnego ciągu $(A_n)_{n\ge1}$ podzbiorów $X$ zachodzi
\begin{equation}\label{eq:outer-subadd}
\mu^*\Big(\bigcup_{n=1}^\infty A_n\Big)
\le\sum_{n=1}^\infty\mu^*(A_n).
\end{equation}
\end{enumerate}
Nie zakładamy, że $\mu^*$ jest miarą na $\mathcal P(X)$:
na ogół nie jest addytywna nawet dla rozłącznych zbiorów.
\end{definicja}

\begin{przyklad}[Najprostsze miary zewnętrzne]
Dla każdego zbioru $X$ funkcja $\mu^*(A)=0$ jest miarą zewnętrzną.
Miara licząca $\#(A)$ z przykładu \ref{ex:counting}, traktowana
jako funkcja na $\mathcal P(X)$, także jest miarą zewnętrzną.
Bardziej pouczający przykład otrzymujemy na $X=\{a,b\}$, gdy
$\mu^*(\varnothing)=0$ i $\mu^*(A)=1$ dla każdego $A\ne\varnothing$.
Wtedy $\mu^*(\{a\}\cup\{b\})=1\ne2=\mu^*(\{a\})+\mu^*(\{b\})$.
Wkrótce zobaczymy, że jej zbiory mierzalne w sensie Carathéodory'ego
to jedynie $\varnothing$ i $X$.
\end{przyklad}

\subsection{Miara zewnętrzna z kosztu pokryć}

\begin{twierdzenie}[Konstrukcja przez pokrycia]
\label{thm:outer-from-cover}
Niech $\mathcal C\subseteq\mathcal P(X)$ będzie rodziną zawierającą
$\varnothing$ oraz taką, że $X$ jest sumą przeliczalnie wielu zbiorów
z $\mathcal C$. Niech $c:\mathcal C\to[0,\infty]$ spełnia
$c(\varnothing)=0$. Dla każdego $E\subseteq X$ określmy
\begin{equation}\label{eq:outer-cover}
\mu_c^*(E):=\inf\Big\{\sum_{n=1}^\infty c(C_n):
C_n\in\mathcal C,\quad E\subseteq\bigcup_{n=1}^\infty C_n\Big\}.
\end{equation}
Wówczas $\mu_c^*$ jest miarą zewnętrzną.
\end{twierdzenie}

\begin{proof}
Pusty zbiór można pokryć ciągiem zbiorów pustych, więc
$0\le\mu_c^*(\varnothing)\le0$. Jeżeli $E\subseteq F$, to każde
pokrycie $F$ dopuszczalne w \eqref{eq:outer-cover} jest pokryciem $E$;
przy infimum po większej rodzinie nie można uzyskać większej wartości.

Pozostaje podaddytywność. Niech $E=\bigcup_{j\ge1}E_j$.
Jeżeli $\sum_j\mu_c^*(E_j)=\infty$, teza jest automatyczna.
W przeciwnym przypadku, przy ustalonym $\varepsilon>0$, dla każdego
$j$ wybieramy pokrycie $E_j\subseteq\bigcup_{k\ge1}C_{j,k}$, dla
którego
$\sum_{k\ge1}c(C_{j,k})<\mu_c^*(E_j)+\varepsilon2^{-j}$.
Taki wybór jest możliwy z definicji infimum; sumy skończone też
można uzupełnić pustymi zbiorami. Rodzina
$\{C_{j,k}:j,k\ge1\}$ jest przeliczalna i pokrywa $E$. Zatem
\[
\mu_c^*(E)\le \sum_{j,k\ge1}c(C_{j,k})
\le\sum_{j\ge1}\mu_c^*(E_j)+\varepsilon.
\]
Ponieważ $\varepsilon>0$ było dowolne, otrzymujemy
\eqref{eq:outer-subadd}. Przestawianie sum jest tu dozwolone,
gdyż wszystkie składniki są nieujemne: obie sumy to kresy górne
sum po skończonych podzbiorach $\mathbb N^2$.
\end{proof}

\begin{uwaga}[Rola założeń]
W konstrukcji wystarczy, aby każde $E\subseteq X$ miało jakieś
przeliczalne pokrycie zbiorami z $\mathcal C$; przyjęte założenie
zapewnia to jednocześnie wszystkim $E$. Nie wymagamy addytywności $c$.
Z drugiej strony samo pokrywanie nie gwarantuje, że
$\mu_c^*(C)=c(C)$ dla $C\in\mathcal C$.
Równość uzyskamy dopiero przy założeniu, że $c$ jest premiarą.
\end{uwaga}

\section{Kryterium mierzalności Carathéodory'ego}

\begin{definicja}[Zbiór mierzalny w sensie Carathéodory'ego]
\label{def:caratheodory-measurable}
Niech $\mu^*$ będzie miarą zewnętrzną na $X$. Zbiór $E\subseteq X$
nazywamy $\mu^*$-\emph{mierzalnym}, jeśli dla każdego $T\subseteq X$
zachodzi
\begin{equation}\label{eq:caratheodory-criterion}
\mu^*(T)=\mu^*(T\cap E)+\mu^*(T\setminus E).
\end{equation}
Rodzinę takich zbiorów oznaczamy przez $\mathcal M_{\mu^*}$.
\end{definicja}

Z nierówności \eqref{eq:outer-subadd} zawsze wynika
$\mu^*(T)\le\mu^*(T\cap E)+\mu^*(T\setminus E)$.
Weryfikacja \eqref{eq:caratheodory-criterion} polega więc na
udowodnieniu \emph{odwrotnej} nierówności, także wtedy, gdy
$\mu^*(T)=\infty$.

\begin{twierdzenie}[Twierdzenie Carathéodory'ego]
\label{thm:caratheodory}
Jeśli $\mu^*$ jest miarą zewnętrzną na $X$, to
$\mathcal M_{\mu^*}$ jest $\sigma$-algebrą na $X$, zaś
$\mu:=\mu^*|_{\mathcal M_{\mu^*}}$ jest miarą. Ponadto miara
$\mu$ jest zupełna: każdy podzbiór zbioru $\mu$-miary zero jest
$\mu^*$-mierzalny i ma miarę zero.
\end{twierdzenie}

\begin{proof}
\emph{Krok 1: dopełnienia i zbiory trywialne.}
Dla $E=\varnothing$ równość
\eqref{eq:caratheodory-criterion} ma postać
$\mu^*(T)=0+\mu^*(T)$. Warunek jest symetryczny względem
$E$ i $X\setminus E$, toteż mierzalność $E$ pociąga mierzalność
jego dopełnienia. W szczególności $X\in\mathcal M_{\mu^*}$.

\emph{Krok 2: sumy dwóch zbiorów.}
Dla $E,F\in\mathcal M_{\mu^*}$ i dowolnego $T\subseteq X$
rozkładamy najpierw $T$ według $E$, a potem $T\setminus E$
według $F$:
\begin{align*}
\mu^*(T)
&=\mu^*(T\cap E)+\mu^*(T\setminus E)\\
&=\mu^*(T\cap E)+\mu^*((T\setminus E)\cap F)
  +\mu^*(T\setminus(E\cup F))\\
&\ge\mu^*(T\cap(E\cup F))
  +\mu^*(T\setminus(E\cup F)).
\end{align*}
Ostatnia nierówność wynika z podaddytywności, gdyż
$T\cap(E\cup F)=(T\cap E)\cup((T\setminus E)\cap F)$.
Nierówność w drugą stronę obowiązuje dla każdego zbioru, a zatem
$E\cup F$ jest mierzalny. Mamy też zamkniętość na skończone
przekroje, dzięki prawom de Morgana.

\emph{Krok 3: przeliczalne sumy rozłączne.}
Niech $(E_n)_{n\ge1}$ będzie parami rozłącznym ciągiem
zbiorów mierzalnych, a $E=\bigcup_nE_n$.
Wielokrotnie stosując kryterium do $E_1,\ldots,E_N$,
otrzymujemy dla dowolnego $T\subseteq X$
\[
\mu^*(T)=\sum_{n=1}^N\mu^*(T\cap E_n)
+\mu^*\Big(T\setminus\bigcup_{n=1}^N E_n\Big).
\]
Ponieważ $T\setminus E\subseteq T\setminus\bigcup_{n=1}^N E_n$,
monotoniczność daje
$\mu^*(T)\ge\sum_{n=1}^N\mu^*(T\cap E_n)
+\mu^*(T\setminus E)$.
Przechodząc do kresu górnego po $N$, otrzymujemy
\[
\mu^*(T)\ge\sum_{n=1}^\infty\mu^*(T\cap E_n)
+\mu^*(T\setminus E)
\ge\mu^*(T\cap E)+\mu^*(T\setminus E).
\]
Ostatnie przejście korzysta z podaddytywności dla
$T\cap E=\bigcup_n(T\cap E_n)$. To dowodzi mierzalności $E$.

\emph{Krok 4: dowolne sumy przeliczalne.}
Dla mierzalnych $A_n$ stosujemy rozłącznienie
$E_1=A_1$ oraz $E_n=A_n\setminus\bigcup_{k<n}A_k$ dla $n\ge2$.
Na mocy poprzednich kroków wszystkie $E_n$ są mierzalne i
$\bigcup_n A_n=\bigcup_n E_n$; zatem
$\mathcal M_{\mu^*}$ jest $\sigma$-algebrą.

\emph{Krok 5: przeliczalna addytywność.}
Dla rozłącznych mierzalnych $E_n$ podstawiamy $T=E=\bigcup_nE_n$
do nierówności z kroku 3. Otrzymujemy
$\mu^*(E)\ge\sum_n\mu^*(E_n)$, a odwrotną nierówność zapewnia
\eqref{eq:outer-subadd}. Ponadto $\mu^*(\varnothing)=0$,
więc ograniczenie $\mu$ jest miarą.

\emph{Krok 6: zupełność.}
Niech $N\in\mathcal M_{\mu^*}$, $\mu^*(N)=0$ i $S\subseteq N$.
Dla dowolnego $T$ z monotoniczności mamy
$\mu^*(T\cap S)=0$ i
$\mu^*(T\setminus S)\le\mu^*(T)$.
Stąd
$\mu^*(T\cap S)+\mu^*(T\setminus S)\le\mu^*(T)$.
Podaddytywność daje nierówność przeciwną. Zatem $S$ jest
mierzalny oraz $\mu^*(S)=0$.
\end{proof}

\begin{przyklad}[Miara zewnętrzna, która nie mierzy wszystkich zbiorów]
Niech $X=\{a,b\}$ oraz $\mu^*(A)=1$ dla $A\ne\varnothing$.
Gdyby $E=\{a\}$ było mierzalne, podstawienie $T=X$ dałoby
$1=\mu^*(X)=\mu^*(\{a\})+\mu^*(\{b\})=2$,
sprzeczność. Stąd $\mathcal M_{\mu^*}=\{\varnothing,X\}$.
To pokazuje, że kryterium Carathéodory'ego ma rzeczywistą treść.
\end{przyklad}

\section{Premiary na algebrach zbiorów}

\begin{definicja}[Premiara]
\label{def:premeasure}
Niech $\mathcal A$ będzie algebrą podzbiorów $X$.
Funkcję $\mu_0:\mathcal A\to[0,\infty]$ nazywamy
\emph{premiarą}, jeżeli $\mu_0(\varnothing)=0$ i dla każdego
ciągu parami rozłącznych $A_n\in\mathcal A$ takiego, że
$A=\bigcup_{n\ge1}A_n\in\mathcal A$, zachodzi
\begin{equation}\label{eq:premeasure}
\mu_0(A)=\sum_{n=1}^\infty\mu_0(A_n).
\end{equation}
Warunek nie zakłada, że $\bigcup_nA_n$ zawsze należy do
$\mathcal A$; algebra nie musi być $\sigma$-algebrą.
\end{definicja}

\begin{lemat}[Podaddytywność premiary na pokryciach]
\label{lem:premeasure-cover}
Jeżeli $\mu_0$ jest premiarą na $\mathcal A$, zaś
$A,A_n\in\mathcal A$ i $A\subseteq\bigcup_{n\ge1}A_n$,
to $\mu_0(A)\le\sum_{n\ge1}\mu_0(A_n)$.
\end{lemat}

\begin{proof}
Tworzymy rozłączne zbiory
$B_1=A\cap A_1$ i
$B_n=(A\cap A_n)\setminus\bigcup_{k<n}A_k$ dla $n\ge2$.
Wszystkie należą do algebry, a $A=\bigsqcup_nB_n$.
Z \eqref{eq:premeasure} oraz monotoniczności premiary
(wynikającej ze skończonej addytywności) dostajemy
$\mu_0(A)=\sum_n\mu_0(B_n)\le\sum_n\mu_0(A_n)$.
\end{proof}

\begin{przyklad}[Skończona addytywność nie wystarcza]
Na $X=\mathbb N$ niech $\mathcal A$ będzie algebrą
zbiorów skończonych i kofinitywnych. Określmy
$\nu_0(A)=0$, gdy $A$ jest skończony, oraz
$\nu_0(A)=1$, gdy $\mathbb N\setminus A$ jest skończony.
Jest to funkcja skończenie addytywna: dwa rozłączne
zbiory kofinitywne nie istnieją, a dołączenie zbioru
skończonego nie zmienia wartości. Nie jest premiarą, gdyż
$\mathbb N=\bigsqcup_{n\ge1}\{n\}$, natomiast
$\nu_0(\mathbb N)=1\ne\sum_n\nu_0(\{n\})=0$.
Każda miara rozszerzająca $\nu_0$ musiałaby spełniać
sprzeczną równość, toteż takie rozszerzenie nie istnieje.
\end{przyklad}

\begin{twierdzenie}[Kryterium ciągłości premiary skończonej]
\label{thm:premeasure-continuity}
Niech $\mathcal A$ będzie algebrą podzbiorów $X$, a
$\mu_0:\mathcal A\to[0,\infty)$ funkcją skończenie addytywną
spełniającą $\mu_0(X)<\infty$ i $\mu_0(\varnothing)=0$.
Wówczas $\mu_0$ jest premiarą wtedy i tylko wtedy, gdy
$A_n\in\mathcal A$, $A_n\downarrow\varnothing$ implikuje
$\mu_0(A_n)\to0$.
\end{twierdzenie}

\begin{proof}
Jeśli $\mu_0$ jest premiarą, to dla
$A_n\downarrow\varnothing$ zbiory $D_n=A_n\setminus A_{n+1}$
są rozłączne i $A_1=\bigsqcup_nD_n$.
W konsekwencji
$\mu_0(A_n)=\sum_{k\ge n}\mu_0(D_k)\to0$,
gdyż szereg ma skończoną sumę $\mu_0(A_1)\le\mu_0(X)$.

Odwrotnie, niech $A=\bigsqcup_n A_n$ i wszystkie te zbiory
należą do $\mathcal A$. Dla
$R_N=A\setminus\bigcup_{n=1}^N A_n$ mamy
$R_N\downarrow\varnothing$, a skończona addytywność daje
$\mu_0(A)=\sum_{n=1}^N\mu_0(A_n)+\mu_0(R_N)$.
Po przejściu do granicy otrzymujemy \eqref{eq:premeasure}.
\end{proof}

\begin{uwaga}
Nie wolno bez zmian stosować twierdzenia
\ref{thm:premeasure-continuity} przy $\mu_0(X)=\infty$.
Na przykład miara licząca na $\mathbb N$ jest premiarą,
lecz dla $A_n=\{n,n+1,\ldots\}$ mamy $A_n\downarrow\varnothing$
i $\mu_0(A_n)=\infty$ dla wszystkich $n$.
\end{uwaga}

\section{Twierdzenie o rozszerzaniu miar}

\begin{twierdzenie}[Rozszerzenie Carathéodory'ego]
\label{thm:extension-caratheodory}
Niech $\mathcal A$ będzie algebrą podzbiorów $X$, a
$\mu_0$ premiarą na $\mathcal A$. Zdefiniujmy $\mu^*$
przez \eqref{eq:outer-cover}, biorąc $\mathcal C=\mathcal A$
i $c=\mu_0$. Wówczas:
\begin{enumerate}[label=\textup{(\roman*)}]
\item $\mu^*$ jest miarą zewnętrzną;
\item każdy $A\in\mathcal A$ jest $\mu^*$-mierzalny;
\item $\mu^*(A)=\mu_0(A)$ dla wszystkich $A\in\mathcal A$;
\item $\mu=\mu^*|_{\mathcal M_{\mu^*}}$ jest zupełną miarą
rozszerzającą $\mu_0$, a jej ograniczenie do
$\sigma(\mathcal A)$ jest miarą na $\sigma(\mathcal A)$.
Miara na $\sigma(\mathcal A)$ nie musi jednak być zupełna.
\end{enumerate}
\end{twierdzenie}

\begin{proof}
Punkt (i) wynika z twierdzenia \ref{thm:outer-from-cover},
gdyż $X\in\mathcal A$.
Ustalmy $A\in\mathcal A$ i dowolny $T\subseteq X$.
Dla każdego pokrycia $T\subseteq\bigcup_nC_n$ przez
$C_n\in\mathcal A$ zbiory $C_n\cap A$ pokrywają $T\cap A$,
a $C_n\setminus A$ pokrywają $T\setminus A$. Dlatego
\begin{align*}
\mu^*(T\cap A)+\mu^*(T\setminus A)
&\le\sum_n\mu_0(C_n\cap A)+\sum_n\mu_0(C_n\setminus A)\\
&=\sum_n\mu_0(C_n).
\end{align*}
Równość korzysta ze skończonej addytywności $\mu_0$.
Biorąc infimum po pokryciach, otrzymujemy
$\mu^*(T\cap A)+\mu^*(T\setminus A)\le\mu^*(T)$.
W połączeniu z podaddytywnością jest to warunek
\eqref{eq:caratheodory-criterion}; dowodzi (ii).

Dla każdego $A\in\mathcal A$ pokrycie $A$ samym $A$
(i pozostałymi zbiorami pustymi) daje
$\mu^*(A)\le\mu_0(A)$. Z drugiej strony
lemat \ref{lem:premeasure-cover} zapewnia
$\mu_0(A)\le\sum_n\mu_0(C_n)$ dla dowolnego pokrycia
$A\subseteq\bigcup_nC_n$. Infimum daje
$\mu_0(A)\le\mu^*(A)$, a więc (iii).

Punkt (iv) jest już konsekwencją twierdzenia
\ref{thm:caratheodory}: $\mathcal M_{\mu^*}$ jest
$\sigma$-algebrą zawierającą $\mathcal A$, zatem
zawiera $\sigma(\mathcal A)$; ograniczenie $\mu^*$
jest na niej miarą, zaś na $\mathcal M_{\mu^*}$ miarą zupełną.
\end{proof}

\begin{definicja}[\texorpdfstring{$\sigma$}{sigma}-skończona premiara]\label{idx:03:definicja:09}

Premiarę $\mu_0$ na algebrze $\mathcal A$ nazywamy
$\sigma$-\emph{skończoną}, jeżeli istnieją
$A_n\in\mathcal A$ takie, że $X=\bigcup_nA_n$ oraz
$\mu_0(A_n)<\infty$ dla każdego $n$.
Możemy zawsze zastąpić $A_n$ sumami skończonymi
$A_1\cup\cdots\cup A_n$, otrzymując ciąg rosnący.
\end{definicja}

\begin{twierdzenie}[Jednoznaczność rozszerzenia]
\label{thm:extension-uniqueness}
Jeśli premiara $\mu_0$ na algebrze $\mathcal A$ jest
$\sigma$-skończona, to istnieje dokładnie jedna miara
na $\sigma(\mathcal A)$, której ograniczenie do
$\mathcal A$ jest równe $\mu_0$.
\end{twierdzenie}

\begin{proof}
Istnienie daje twierdzenie \ref{thm:extension-caratheodory}.
Niech $\mu$ i $\nu$ będą dwoma rozszerzeniami.
Wybierzmy rosnący ciąg $A_n\in\mathcal A$ o
$\bigcup_n A_n=X$ i $\mu_0(A_n)<\infty$.
Dla ustalonego $n$ funkcje
$\mu_n(E)=\mu(E\cap A_n)$ oraz $\nu_n(E)=\nu(E\cap A_n)$,
$E\in\sigma(\mathcal A)$, są miarami skończonymi.
Zgadzają się na algebrze $\mathcal A$, ponieważ
$E\cap A_n\in\mathcal A$ dla $E\in\mathcal A$.
Algebra jest układem $\pi$ i generuje
$\sigma(\mathcal A)$. Z twierdzenia
\ref{thm:uniqueness-finite} otrzymujemy
$\mu_n(E)=\nu_n(E)$ dla każdego $E\in\sigma(\mathcal A)$.
Ponieważ $E\cap A_n\uparrow E$, ciągłość miary z dołu
(twierdzenie \ref{thm:measure-continuity}) daje
$\mu(E)=\lim_n\mu_n(E)=\lim_n\nu_n(E)=\nu(E)$.
\end{proof}

\begin{przyklad}[Dlaczego potrzebna jest $\sigma$-skończoność]
\label{ex:nonunique-extension}
Na $X=\mathbb R$ niech $\mathcal A$ będzie algebrą
generowaną przez przedziały $[a,b)$, gdzie $a<b$ są
rzeczywiste. Każdy niepusty element $\mathcal A$ zawiera
niezdegenerowany przedział; ponadto
$\sigma(\mathcal A)=\mathcal B(\mathbb R)$.
Na zbiorach borelowskich rozważmy dwie miary:
\[
\mu(E)=\#E,\qquad
\nu(E)=\begin{cases}0,&E\text{ jest przeliczalny},\\
\infty,&E\text{ jest nieprzeliczalny}.
\end{cases}
\]
Funkcja $\nu$ jest miarą, gdyż przeliczalna suma zbiorów
przeliczalnych jest przeliczalna, a suma zawierająca
zbiór nieprzeliczalny jest nieprzeliczalna. Na każdym
niepustym elemencie $\mathcal A$ obie miary mają wartość
$\infty$, zaś na pustym obie są równe zero. Są to więc
różne rozszerzenia tej samej premiary:
$\mu(\{0\})=1$, natomiast $\nu(\{0\})=0$.
Premiara nie jest $\sigma$-skończona: jej jedynym
zbiorem o skończonej mierze jest $\varnothing$.
\end{przyklad}

\begin{uwaga}[Różne role założeń]
Premiara daje istnienie rozszerzenia na $\sigma(\mathcal A)$
bez założenia $\sigma$-skończoności. To założenie pojawia się
dopiero przy jednoznaczności rozszerzenia oraz przy
identyfikacji całej rodziny zbiorów Carathéodory'ego
z uzupełnieniem $\sigma(\mathcal A)$. W szczególności
stwierdzenie o istnieniu i stwierdzenie o jednoznaczności
nie są jednym i tym samym twierdzeniem.
\end{uwaga}

\section{Kiedy mierzalność Carathéodory'ego jest uzupełnieniem rozszerzenia?}

Rozszerzenie otrzymane z twierdzenia
\ref{thm:extension-caratheodory} jest określone na
$\mathcal M_{\mu^*}$, zwykle większej niż $\sigma(\mathcal A)$.
Nie należy utożsamiać tych rodzin bez dowodu.
Poniższy rezultat określa ich związek przy
$\sigma$-skończoności.

\begin{lemat}[Aproksymacja od góry]
\label{lem:outer-approx-algebra}
Niech $\mu_0$ będzie premiarą na algebrze $\mathcal A$ i
$\mu^*$ jej miarą zewnętrzną. Jeśli $E\subseteq X$ oraz
$\mu^*(E)<\infty$, to dla każdego $\varepsilon>0$
istnieje $B\in\sigma(\mathcal A)$, $E\subseteq B$, takie, że
$\mu^*(B)\le\mu^*(E)+\varepsilon$.
Jeżeli dodatkowo $E\in\mathcal M_{\mu^*}$, to
$\mu(B\setminus E)<\varepsilon$.
\end{lemat}

\begin{proof}
Wybieramy pokrycie $E\subseteq\bigcup_nA_n$ przez
$A_n\in\mathcal A$ z $\sum_n\mu_0(A_n)<\mu^*(E)+\varepsilon$.
Wystarczy wziąć $B=\bigcup_nA_n$. Z podaddytywności i
równości $\mu^*(A_n)=\mu_0(A_n)$ dostajemy
$\mu^*(B)\le\sum_n\mu_0(A_n)<\mu^*(E)+\varepsilon$.
Jeśli $E$ jest mierzalny, to $B\setminus E$ jest mierzalny
i $\mu(B)=\mu(E)+\mu(B\setminus E)$.
Obie pierwsze wartości są skończone, więc
$\mu(B\setminus E)<\varepsilon$.
\end{proof}

\begin{twierdzenie}[Uzupełnienie rozszerzenia]
\label{thm:caratheodory-completion}
Załóżmy, że premiara $\mu_0$ na algebrze $\mathcal A$
jest $\sigma$-skończona. Niech $\mu^*$ będzie wyznaczoną
przez nią miarą zewnętrzną, a $\mu_B=\mu^*|_{\sigma(\mathcal A)}$.
Wówczas przestrzeń
$(X,\mathcal M_{\mu^*},\mu^*|_{\mathcal M_{\mu^*}})$
jest uzupełnieniem przestrzeni
$(X,\sigma(\mathcal A),\mu_B)$ w sensie twierdzenia
\ref{thm:completion}. Dokładniej: dla każdego
$E\in\mathcal M_{\mu^*}$ istnieją
$C,D\in\sigma(\mathcal A)$ takie, że
$C\subseteq E\subseteq D$ i $\mu_B(D\setminus C)=0$.
\end{twierdzenie}

\begin{proof}
Wybierzmy rosnący ciąg $A_n\in\mathcal A$ pokrywający $X$,
przy czym $\mu_0(A_n)<\infty$. Dla ustalonego mierzalnego
$E$ zbiór $E_n=E\cap A_n$ jest mierzalny i ma skończoną miarę.
Dla każdego $k,n\ge1$ lemat \ref{lem:outer-approx-algebra}
daje $B_{n,k}\in\sigma(\mathcal A)$ z
$E_n\subseteq B_{n,k}$ oraz
$\mu(B_{n,k}\setminus E_n)<2^{-n}/k$.
Połóżmy $D_k=\bigcup_{n\ge1}B_{n,k}$.
Wtedy $D_k\in\sigma(\mathcal A)$,
$E\subseteq D_k$, gdyż $E=\bigcup_nE_n$, i
\[
\mu(D_k\setminus E)
\le\sum_n\mu(B_{n,k}\setminus E)
\le\sum_n\mu(B_{n,k}\setminus E_n)<\frac1k.
\]
Zbiór $D=\bigcap_{k\ge1}D_k$ należy do
$\sigma(\mathcal A)$, zawiera $E$ i spełnia
$\mu(D\setminus E)=0$, ponieważ
$D\setminus E\subseteq D_k\setminus E$ dla każdego $k$.
Powtarzamy konstrukcję dla mierzalnego zbioru
$X\setminus E$, otrzymując zbiór należący do
$\sigma(\mathcal A)$, mianowicie $H\supseteq X\setminus E$,
taki że $\mu(H\setminus(X\setminus E))=0$.
Wtedy $C=X\setminus H\in\sigma(\mathcal A)$,
$C\subseteq E$ i $\mu(E\setminus C)=0$.
Wobec $C\subseteq E\subseteq D$ mamy
$D\setminus C=(D\setminus E)\sqcup(E\setminus C)$,
zatem $\mu_B(D\setminus C)=0$.

Każdy zbiór z uzupełnienia $\sigma(\mathcal A)$ różni
się od zbioru z $\sigma(\mathcal A)$ o podzbiór
zbioru $\mu_B$-zerowego, a więc jest
$\mu^*$-mierzalny z zupełności twierdzenia
\ref{thm:caratheodory}. Przed chwilą udowodniliśmy
zawieranie przeciwne, kończąc dowód.
\end{proof}

\begin{uwaga}[Dlaczego to rozróżnienie jest ważne?]
W rozdziale \ref{ch:lebesgue} rozróżnimy zbiory borelowskie
i zbiory mierzalne w sensie Lebesgue'a. Konstrukcja
Carathéodory'ego od razu daje miarę zupełną; ograniczenie
do borelowskiej $\sigma$-algebry nie musi być zupełne.
\end{uwaga}

\section{Jeszcze jeden sposób interpretacji konstrukcji}

Dla zbioru $E\subseteq X$ liczba $\mu^*(E)$ jest
najmniejszym kosztem \emph{przeliczalnego pokrycia} $E$
przez zbiory z wyjściowej algebry. Mierzalność w sensie
Carathéodory'ego oznacza, że przecięcie z $E$ oraz z jego
dopełnieniem nie powoduje sztucznej utraty kosztu,
niezależnie od badanego zbioru $T$.

Następujące trzy kroki to logiczny szkielet rozdziału:
\begin{enumerate}
\item z premiary $\mu_0$ tworzymy $\mu^*$ na całej $\mathcal P(X)$;
\item warunek \eqref{eq:caratheodory-criterion} wyodrębnia
$\sigma$-algebrę, na której $\mu^*$ jest miarą;
\item przy $\sigma$-skończoności dostajemy jedyne
rozszerzenie na $\sigma(\mathcal A)$, a po uzupełnieniu
pełną dziedzinę mierzalności Carathéodory'ego.
\end{enumerate}
W następnym rozdziale zrealizujemy ten program dla
objętości prostopadłościanów w $\mathbb R^d$.

% END chapters/pl/measure_theory/03-rozdzial.tex


% BEGIN chapters/pl/measure_theory/04-rozdzial.tex
% ROZDZIAŁ 4 — plik włączany przez main.tex
\chapter{Miara Lebesgue'a na przestrzeni euklidesowej}
\label{ch:lebesgue}

\section{Co właściwie chcemy skonstruować?}

Długość przedziału $[a,b)$ to $b-a$, a objętość
$d$-wymiarowego prostopadłościanu
$Q=\prod_{j=1}^d[a_j,b_j)$ wynosi
$|Q|=\prod_{j=1}^d(b_j-a_j)$.
Chcemy rozciągnąć tę funkcję na bogatą rodzinę
podzbiorów $\mathbb R^d$, zachowując przeliczalną
addytywność i naturalną zgodność z geometrią.

Przy każdym $d\in\mathbb N$, $d\ge1$, przez
$\mathcal B(\mathbb R^d)$ oznaczamy $\sigma$-algebrę
borelowską z definicji \ref{def:borel}. Odróżniamy
\emph{miarę borelowską Lebesgue'a} $\lambda_d^B$
od jej uzupełnienia $\lambda_d$. Ich dziedziny nie są
jednakowe. Dla prostopadłościanów osiowych z dowolnym wyborem
otwartości i domknięcia ścian przez $|R|$ oznaczamy
\emph{elementarną objętość geometryczną}, czyli iloczyn
długości boków; nie zakładamy jeszcze, że jest ona miarą.
W szczególności dla $Q$ jest to wyjściowa objętość,
a dla zwartego $K$ lub otwartego $O$ ta sama liczba.
Symbol $|R|$ nie jest na tym etapie miarą dowolnego zbioru.

\section{Objętość elementarna i lemat o pokryciach}

\begin{definicja}[Prostopadłościan półotwarty]
\label{def:halfopen-box}
Przez prostopadłościan półotwarty rozumiemy zbiór
$Q=\prod_{j=1}^d[a_j,b_j)$, gdzie $a_j,b_j\in\mathbb R$,
$a_j<b_j$. Rodzinę takich \emph{ograniczonych} zbiorów
oraz zbioru pustego oznaczamy $\mathcal Q_d$.
Kładziemy $|\varnothing|=0$ i
$|Q|=\prod_{j=1}^d(b_j-a_j)$ dla niepustego $Q$.
\end{definicja}

Użycie przedziałów półotwartych nie jest kosmetyczne.
Dla $a<c<b$ mamy rozkład rozłączny
$[a,b)=[a,c)\sqcup[c,b)$. Analogicznie przecięcie
dwóch prostopadłościanów półotwartych jest
prostopadłościanem tego samego typu lub zbiorem pustym,
a różnica takich zbiorów jest skończoną sumą rozłączną
prostopadłościanów półotwartych.

\begin{lemat}[Skończony podział na siatkę]
\label{lem:rect-grid}
Niech $Q_1,\ldots,Q_N\in\mathcal Q_d$.
Istnieje skończona rodzina parami rozłącznych
prostopadłościanów półotwartych $P_1,\ldots,P_M$, taka że
każdy z $Q_i$ jest sumą pewnych $P_j$. W konsekwencji
można jednoznacznie określić na rodzinie $\mathcal R_d$
wszystkich skończonych sum zbiorów z $\mathcal Q_d$
funkcję $v$, przyjmując
$v(\bigsqcup_{i=1}^NQ_i)=\sum_{i=1}^N|Q_i|$
dla dowolnego rozłącznego przedstawienia.
Funkcja $v$ jest nieujemna i skończenie addytywna.
\end{lemat}

\begin{proof}
W każdej współrzędnej wypisujemy w porządku rosnącym
wszystkie końce przedziałów występujących w $Q_i$.
Sąsiednie końce wyznaczają przedziały półotwarte;
ich iloczyny kartezjańskie dają skończoną siatkę.
Każdy $Q_i$ składa się z całych komórek siatki.
Dla każdej komórki $P=\prod_j[c_j,e_j)$ jej objętość
to $\prod_j(e_j-c_j)$. Dodawanie objętości przy
podziale w jednej współrzędnej jest rozdzielnością
mnożenia względem dodawania; iterując po współrzędnych,
otrzymujemy $|Q_i|$ jako sumę objętości jej komórek.
Jeżeli ten sam zbiór ma dwa rozłączne przedstawienia,
wspólna siatka dla obu przedstawień pokazuje, że oba
sumowania dają tę samą liczbę. Tak samo dowodzi się
skończonej addytywności $v$.
\end{proof}

\begin{uwaga}
$\mathcal R_d$ jest pierścieniem zbiorów, ale nie
algebrą podzbiorów $\mathbb R^d$, ponieważ
$\mathbb R^d\notin\mathcal R_d$. Na ograniczonym
prostopadłościanie $H$ rodzina
$\mathcal R_d|_H=\{A\cap H:A\in\mathcal R_d\}$
jest natomiast algebrą podzbiorów $H$.
To lokalne rozróżnienie pozwoli nam zastosować
rozdział \ref{ch:caratheodory} bez udawania, że
nieograniczona przestrzeń jest ograniczonym prostokątem.
\end{uwaga}

\begin{lemat}[Nierówność objętościowa dla pokryć]
\label{lem:rect-cover-ineq}
Jeżeli $Q,Q_1,Q_2,\ldots\in\mathcal Q_d$ i
$Q\subseteq\bigcup_{n\ge1}Q_n$, to
\begin{equation}\label{eq:rect-cover-volume}
|Q|\le\sum_{n=1}^\infty|Q_n|.
\end{equation}
\end{lemat}

\begin{proof}
Dla $Q=\varnothing$ twierdzenie jest oczywiste.
Niech $Q=\prod_j[a_j,b_j)$ będzie niepusty.
Wystarczy rozważyć przypadek, w którym suma po prawej
stronie jest skończona. Ustalmy $\varepsilon>0$ oraz
$0<\delta<\frac12\min_j(b_j-a_j)$ i połóżmy
$K_\delta=\prod_j[a_j+\delta,b_j-\delta]$.
Jest to zwarty podzbiór $Q$.

Każdy $Q_n$ można zawrzeć w otwartym prostopadłościanie
$O_n$ o objętości elementarnej mniejszej niż
$|Q_n|+\varepsilon2^{-n}$: wystarczy nieznacznie
przesunąć wszystkie jego ściany na zewnątrz.
Zbiory $O_n$ tworzą otwarte pokrycie $K_\delta$.
Ze zwartości istnieją $n_1,\ldots,n_N$ takie, że
$K_\delta\subseteq\bigcup_{i=1}^N O_{n_i}$.

Wykażmy osobno elementarną nierówność dla skończonego
pokrycia zwartego prostopadłościanu
$K=\prod_{j=1}^d[\alpha_j,\beta_j]$ przez otwarte
prostopadłościany $O_1,\ldots,O_m$:
$|K|\le\sum_{i=1}^m|O_i|$.
Dla dostatecznie małego $\eta>0$ wprowadzamy
$P_\eta=\prod_j[\alpha_j+\eta,\beta_j-\eta)$.
W każdej współrzędnej wybieramy wszystkie końce
przedziałów tworzących $O_i$ i $P_\eta$.
Skończona siatka złożona z tych punktów rozkłada
$P_\eta$ na parami rozłączne półotwarte komórki $S_l$.
Wnętrze komórki nie przecina żadnej hiperpłaszczyzny,
która jest ścianą $O_i$. Wybierzmy
$x_l\in\operatorname{int}S_l\subseteq K$.
Punkt $x_l$ należy do pewnego $O_{i(l)}$.
Ponieważ granice współrzędnych $O_{i(l)}$ należą do
siatki, wynika stąd
$\operatorname{int}S_l\subseteq O_{i(l)}$, a więc
$S_l\subseteq\overline{O_{i(l)}}$.

Dla ustalonego $i$ oznaczmy przez $\mathcal S_i$ rodzinę
komórek przypisanych do $O_i$. Niech $H_i^{(\gamma)}$
będzie półotwartym prostopadłościanem zawierającym
$\overline O_i$, którego długości boków przewyższają
odpowiednie długości $O_i$ o $2\gamma$.
Każda komórka $S_l\in\mathcal S_i$ jest zawarta w
$H_i^{(\gamma)}$. Na mocy skończonej addytywności
i monotoniczności elementarnej objętości na pierścieniu
$\mathcal R_d$ mamy
$\sum_{S_l\in\mathcal S_i}|S_l|
\le |H_i^{(\gamma)}|$.
Gdy $\gamma\downarrow0$, prawa strona dąży do $|O_i|$.
Komórki $S_l$ rozkładają $P_\eta$ i każdą przypisaliśmy
jednemu z $O_i$, zatem
\[
 |P_\eta|=\sum_l|S_l|
 =\sum_{i=1}^m\sum_{S_l\in\mathcal S_i}|S_l|
 \le\sum_{i=1}^m|O_i|.
\]
Wreszcie $|P_\eta|\to|K|$ przy $\eta\downarrow0$.
Argument ten wykorzystuje wyłącznie skończone podziały
prostopadłościanów i nie zakłada wcześniej istnienia miary
Lebesgue'a.
Stąd
\[
|K_\delta|\le\sum_{i=1}^N|O_{n_i}|
\le\sum_{n=1}^\infty|Q_n|+\varepsilon.
\]
Przy $\delta\downarrow0$ mamy
$|K_\delta|=\prod_j(b_j-a_j-2\delta)\to|Q|$.
Następnie $\varepsilon\downarrow0$, co dowodzi
\eqref{eq:rect-cover-volume}.
\end{proof}

\begin{uwaga}[W którym miejscu używamy zwartości?]
Bez zwartości nie mielibyśmy prawa zastąpić
przeliczalnego pokrycia skończonym i skorzystać
z rozkładu na siatkę. To zasadniczy krok geometryczny
w dowodzie, że elementarna objętość przechodzi
z addytywności skończonej do przeliczalnej.
\end{uwaga}

\begin{twierdzenie}[Premiara na ograniczonym prostopadłościanie]
\label{thm:volume-premeasure}
Dla ustalonego $H=\prod_{j=1}^d[a_j,b_j)\in\mathcal Q_d$
funkcja $v$ z lematu \ref{lem:rect-grid}, ograniczona do
algebry $\mathcal R_d|_H$, jest skończoną premiarą.
Jej przedłużenie na $\sigma(\mathcal R_d|_H)$ istnieje
i jest jednoznaczne.
\end{twierdzenie}

\begin{proof}
Wszystkie zbiory z $\mathcal R_d|_H$ są skończonymi
rozłącznymi sumami prostopadłościanów półotwartych
zawartych w $H$; wspólna siatka zapewnia skończoną
addytywność. W szczególności $v(H)=|H|<\infty$.

Niech $A=\bigsqcup_{n\ge1}A_n$, gdzie
$A,A_n\in\mathcal R_d|_H$. Z addytywności skończonej
i monotoniczności dostajemy
$\sum_{n=1}^Nv(A_n)\le v(A)$ dla każdego $N$, więc
$\sum_nv(A_n)\le v(A)$.
Rozbijmy $A$ i każdy $A_n$ na skończenie wiele
rozłącznych prostopadłościanów. Każdy składnik rozkładu
$A$ jest pokryty składnikami wszystkich $A_n$.
Lemat \ref{lem:rect-cover-ineq} daje po zsumowaniu
$v(A)\le\sum_nv(A_n)$. Zachodzi więc addytywność
przeliczalna w sensie definicji \ref{def:premeasure}.
Istnienie przedłużenia wynika z twierdzenia
\ref{thm:extension-caratheodory}, a jednoznaczność
z twierdzenia \ref{thm:extension-uniqueness}.
\end{proof}

\section{Konstrukcja miary Lebesgue'a}

\begin{definicja}[Miara zewnętrzna Lebesgue'a]
\label{def:lebesgue-outer}
Dla $E\subseteq\mathbb R^d$ określamy
\begin{equation}\label{eq:lebesgue-outer}
\lambda_d^*(E)=\inf\Big\{\sum_{n=1}^\infty|Q_n|:
Q_n\in\mathcal Q_d,\quad E\subseteq\bigcup_{n=1}^\infty Q_n\Big\}.
\end{equation}
Rodzinę zbiorów $\lambda_d^*$-mierzalnych w sensie
Carathéodory'ego oznaczamy $\mathcal L_d$.
Miarę $\lambda_d=\lambda_d^*|_{\mathcal L_d}$
nazywamy \emph{$d$-wymiarową miarą Lebesgue'a}.
\end{definicja}

\begin{twierdzenie}[Istnienie i podstawowe własności]
\label{thm:lebesgue-exists}
Dla każdego $d\ge1$ zachodzą następujące fakty:
\begin{enumerate}[label=\textup{(\roman*)}]
\item $\lambda_d^*$ jest miarą zewnętrzną;
\item każdy $Q\in\mathcal Q_d$ należy do $\mathcal L_d$ i
$\lambda_d(Q)=|Q|$;
\item $\mathcal B(\mathbb R^d)\subseteq\mathcal L_d$;
\item $\lambda_d$ jest zupełną i $\sigma$-skończoną miarą;
\item jej ograniczenie
$\lambda_d^B=\lambda_d|_{\mathcal B(\mathbb R^d)}$
jest jedyną miarą borelowską, która przypisuje każdemu
$Q\in\mathcal Q_d$ liczbę $|Q|$.
\end{enumerate}
\end{twierdzenie}

\begin{proof}
\emph{(i)} Rodzina $\mathcal Q_d$ zawiera $\varnothing$,
a $\mathbb R^d$ ma przeliczalne pokrycie prostopadłościanami
$[-n,n)^d$. Twierdzenie \ref{thm:outer-from-cover}
gwarantuje zatem, że \eqref{eq:lebesgue-outer}
określa miarę zewnętrzną.

\emph{(ii), równość objętości.}
Pokrycie $Q$ jednym $Q$ daje $\lambda_d^*(Q)\le|Q|$.
Lemat \ref{lem:rect-cover-ineq} daje nierówność
przeciwną dla każdego pokrycia dopuszczalnego w
\eqref{eq:lebesgue-outer}; zatem $\lambda_d^*(Q)=|Q|$.

\emph{(ii), mierzalność.}
Ustalmy $Q\in\mathcal Q_d$ i $T\subseteq\mathbb R^d$.
Dla każdego pokrycia $T\subseteq\bigcup_nR_n$, gdzie
$R_n\in\mathcal Q_d$, dzielimy $R_n$ wzdłuż ścian $Q$.
Przecięcie $R_n\cap Q$ jest jednym prostopadłościanem
półotwartym lub zbiorem pustym, natomiast
$R_n\setminus Q$ jest skończoną rozłączną sumą
takich prostopadłościanów. Na mocy lematu
\ref{lem:rect-grid} łączna objętość obu części wynosi
$|R_n|$. Części pierwszego rodzaju pokrywają
$T\cap Q$, a części drugiego $T\setminus Q$, więc
\[
\lambda_d^*(T\cap Q)+\lambda_d^*(T\setminus Q)
\le\sum_{n=1}^\infty |R_n|.
\]
Biorąc infimum po pokryciach, uzyskujemy nierówność
$\le\lambda_d^*(T)$; przeciwna jest zawsze prawdziwa.
Wynika stąd $Q\in\mathcal L_d$.

\emph{(iii)} Twierdzenie \ref{thm:caratheodory} daje,
że $\mathcal L_d$ jest $\sigma$-algebrą. Ponieważ zawiera
wszystkie prostopadłościany półotwarte, zawiera również
ich $\sigma$-algebrę, czyli $\mathcal B(\mathbb R^d)$
(twierdzenie \ref{thm:borel-Rd}).

\emph{(iv)} Zupełność daje twierdzenie
\ref{thm:caratheodory}. Dla $H_n=[-n,n)^d$ mamy
$\mathbb R^d=\bigcup_nH_n$ i
$\lambda_d(H_n)=(2n)^d<\infty$, więc miara jest
$\sigma$-skończona.

\emph{(v)} Niech $\nu$ będzie inną miarą borelowską
o zadanych wartościach na $\mathcal Q_d$.
Rodzina $\mathcal Q_d$ wraz z $\varnothing$ jest
układem $\pi$, generuje $\mathcal B(\mathbb R^d)$,
a rosnące prostopadłościany $H_n$ pokrywają całą
przestrzeń i mają tę samą skończoną miarę dla
$\nu$ oraz $\lambda_d^B$. Twierdzenie
\ref{thm:uniqueness-sigma-finite} daje
$\nu=\lambda_d^B$.
\end{proof}

\begin{uwaga}[Zgodność obu konstrukcji]
Dla każdego ograniczonego prostopadłościanu $H$
miara z twierdzenia \ref{thm:volume-premeasure} jest
równa ograniczeniu $\lambda_d^B$ do
$\mathcal B(H)$. Obie są skończonymi miarami na
tej samej $\sigma$-algebrze i zgadzają się na
półotwartych prostopadłościanach zawartych w $H$.
W ten sposób globalna konstrukcja przez pokrycia
zgadza się z lokalnym rozszerzaniem premiary.
\end{uwaga}

\section{Niezmienniczość względem przesunięć i skalowania}

\begin{lemat}[Zbiory zawarte w hiperpłaszczyźnie]
\label{lem:hyperplane-null}
Dla każdego $c\in\mathbb R$ i $1\le j\le d$
hiperpłaszczyzna $\{x\in\mathbb R^d:x_j=c\}$ ma
miarę Lebesgue'a zero. W szczególności każda skończona
lub przeliczalna rodzina punktów ma miarę zero.
\end{lemat}

\begin{proof}
Przecięcie hiperpłaszczyzny z $[-N,N)^d$ zawiera się
w półotwartym prostopadłościanie, którego $j$-ty bok
ma długość $2\delta$, a pozostałe długość $2N$
(po ewentualnym nieznacznym powiększeniu boków).
Z definicji miary zewnętrznej
$\lambda_d^*(\{x:x_j=c\}\cap[-N,N)^d)
\le (2N)^{d-1}2\delta$ dla dowolnego $\delta>0$.
Przechodząc z $\delta\downarrow0$, dostajemy wartość zero.
Hiperpłaszczyzna jest przeliczalną sumą takich przecięć,
a więc ma miarę zero. Singleton jest podzbiorem
hiperpłaszczyzny, a przeliczalna suma zbiorów zerowych
ma miarę zero na mocy \eqref{eq:subadditivity}.
\end{proof}

\begin{twierdzenie}[Translacje i jednolite skalowanie]
\label{thm:invariance-lebesgue}
Dla każdego $E\in\mathcal L_d$, każdego
$a\in\mathbb R^d$ i każdego $t\in\mathbb R\setminus\{0\}$
zbiory $a+E$ i $tE$ należą do $\mathcal L_d$ oraz
\begin{equation}\label{eq:lebesgue-invariance}
\lambda_d(a+E)=\lambda_d(E),\qquad
\lambda_d(tE)=|t|^d\lambda_d(E).
\end{equation}
Przyjmujemy tu $|t|^d\cdot\infty=\infty$.
\end{twierdzenie}

\begin{proof}
Udowodnimy najpierw mocniejszy fakt dotyczący miary
\emph{zewnętrznej}: dla każdego $E\subseteq\mathbb R^d$
zachodzi
$\lambda_d^*(a+E)=\lambda_d^*(E)$ oraz
$\lambda_d^*(tE)=|t|^d\lambda_d^*(E)$.

Dla translacji argument jest bezpośredni: jeśli
$E\subseteq\bigcup_nQ_n$, to
$a+E\subseteq\bigcup_n(a+Q_n)$, a
$|a+Q_n|=|Q_n|$. Z definicji infimum otrzymujemy
$\lambda_d^*(a+E)\le\lambda_d^*(E)$.
Stosując tę samą nierówność do przesunięcia o $-a$,
dostajemy równość.

Jeżeli $t>0$, to dla każdego $Q\in\mathcal Q_d$
zbiór $tQ$ jest prostopadłościanem półotwartym
objętości $t^d|Q|$. Przeskalowanie dowolnego
pokrycia $E$ daje
$\lambda_d^*(tE)\le t^d\lambda_d^*(E)$;
powtarzając argument ze skalą $1/t$, otrzymujemy
równość. Jeśli $t<0$, obraz $tQ$ jest iloczynem
przedziałów typu $(tb_j,ta_j]$. Dla każdego
$\eta>0$ można go pokryć prostopadłościanem
półotwartym $R\in\mathcal Q_d$ o
$|R|<|t|^d|Q|+\eta$. Przeskalowanie pokrycia
$E\subseteq\bigcup_nQ_n$ i wybór błędów
$\eta_n=\varepsilon2^{-n}$ daje
$\lambda_d^*(tE)\le |t|^d\sum_n|Q_n|+\varepsilon$.
Po przejściu do infimum i $\varepsilon\downarrow0$
uzyskujemy nierówność $\le$; zastosowanie tego
argumentu do $1/t$ daje nierówność przeciwną.

Wreszcie niech $E\in\mathcal L_d$ i niech
$S(x)=a+tx$, gdzie $t\ne0$. Dla dowolnego
$T\subseteq\mathbb R^d$ możemy podstawić
$S^{-1}(T)$ w kryterium Carathéodory'ego
\eqref{eq:caratheodory-criterion}. Ponieważ $S$ jest
bijekcją oraz
$S(A\cap B)=S(A)\cap S(B)$,
$S(A\setminus B)=S(A)\setminus S(B)$,
a miara zewnętrzna skaluje się przez dodatnią
stałą $|t|^d$, otrzymujemy
\[
\lambda_d^*(T)=\lambda_d^*(T\cap S(E))
+\lambda_d^*(T\setminus S(E)).
\]
Zatem $S(E)\in\mathcal L_d$ i
$\lambda_d(S(E))=|t|^d\lambda_d(E)$.
Dla samej translacji postępujemy identycznie,
używając zachowania miary zewnętrznej bez czynnika.
\end{proof}

\begin{uwaga}[Ogólne przekształcenia liniowe]
Dla odwracalnego przekształcenia liniowego
$T:\mathbb R^d\to\mathbb R^d$ prawdziwy jest
silniejszy wzór $\lambda_d(TE)=|\det T|\lambda_d(E)$.
Jego pełny dowód, obejmujący rozkład przekształcenia
na operacje elementarne lub twierdzenie o zamianie
zmiennych, umieścimy w rozdziale o przekształcaniu całek.
Nie wykorzystujemy tego ogólnego wzoru w żadnym
dowodzie bieżącego rozdziału.
\end{uwaga}

\section{Regularność miary Lebesgue'a}

Przypomnijmy oznaczenia opisujące złożoność topologiczną:
zbiór typu $G_\delta$ jest przeliczalnym przekrojem zbiorów
otwartych, zaś zbiór typu $F_\sigma$ jest przeliczalną
sumą zbiorów domkniętych. Oba rodzaje zbiorów są borelowskie.
Wyniki poniżej wskazują, jak blisko nich leżą dowolne
zbiory mierzalne Lebesgue'a.

\begin{twierdzenie}[Aproksymacja otwartymi zbiorami]
\label{thm:lebesgue-outer-regular}
Dla każdego $E\subseteq\mathbb R^d$ zachodzi
\begin{equation}\label{eq:outer-open-regular}
\lambda_d^*(E)=\inf\{\lambda_d(U):
U\subseteq\mathbb R^d\ \text{otwarty},\ E\subseteq U\}.
\end{equation}
Jeśli ponadto $E\in\mathcal L_d$ i
$\lambda_d(E)<\infty$, to dla każdego
$\varepsilon>0$ można wybrać otwarty $U\supseteq E$
taki, że $\lambda_d(U\setminus E)<\varepsilon$.
\end{twierdzenie}

\begin{proof}
Każdy zbiór otwarty jest borelowski, a więc
mierzalny. Z monotoniczności
$\lambda_d^*(E)\le\lambda_d(U)$, skąd nierówność
$\le$ w \eqref{eq:outer-open-regular}.

Przypuśćmy najpierw, że $\lambda_d^*(E)<\infty$.
Dla $\varepsilon>0$ wybieramy pokrycie
$E\subseteq\bigcup_nQ_n$ przez półotwarte
prostopadłościany tak, że
$\sum_n|Q_n|<\lambda_d^*(E)+\varepsilon/2$.
Dla każdego $n$ wybieramy otwarty prostopadłościan
$O_n\supseteq Q_n$ taki, że
$|O_n|<|Q_n|+\varepsilon2^{-n-1}$.
Zbiór $U=\bigcup_nO_n$ jest otwarty, zawiera $E$ i
\[
\lambda_d(U)\le\sum_n\lambda_d(O_n)
=\sum_n|O_n|<\lambda_d^*(E)+\varepsilon.
\]
Objętość otwartego prostopadłościanu jest równa
iloczynowi długości boków: różni się on od
półotwartego tylko skończoną sumą fragmentów ścian,
które są zerowe na mocy lematu \ref{lem:hyperplane-null}.
Dla $\lambda_d^*(E)=\infty$ wszystkie otwarte
nadzbiory mają miarę nieskończoną, więc równość
\eqref{eq:outer-open-regular} również zachodzi.

Jeśli $E$ jest mierzalny i ma miarę skończoną,
$U\setminus E$ jest mierzalny, zaś
$\lambda_d(U)=\lambda_d(E)+\lambda_d(U\setminus E)$.
Wybór $U$ z poprzedniej części daje tezę.
\end{proof}

\begin{twierdzenie}[Aproksymacja zbiorami zwartymi]
\label{thm:lebesgue-inner-regular}
Dla każdego $E\in\mathcal L_d$ zachodzi
\begin{equation}\label{eq:inner-compact-regular}
\lambda_d(E)=\sup\{\lambda_d(K):K\subseteq E,
\ K\subseteq\mathbb R^d\ \text{zwarty}\}.
\end{equation}
Jeżeli $\lambda_d(E)<\infty$, to dla dowolnego
$\varepsilon>0$ istnieje zwarty $K\subseteq E$
spełniający $\lambda_d(E\setminus K)<\varepsilon$.
\end{twierdzenie}

\begin{proof}
Najpierw załóżmy $\lambda_d(E)<\infty$.
Z ciągłości miary z dołu dla
$E\cap[-N,N]^d\uparrow E$ wybieramy $N$ tak duże, że
$\lambda_d(E\setminus H)<\varepsilon/2$, gdzie
$H=[-N,N]^d$ jest zwarty.
Niech $F=H\setminus E$; jest to zbiór mierzalny
i $\lambda_d(F)\le\lambda_d(H)<\infty$.
Z twierdzenia \ref{thm:lebesgue-outer-regular}
wybieramy zbiór otwarty $U\supseteq F$ o
$\lambda_d(U\setminus F)<\varepsilon/2$.
Zbiór $K=H\setminus U$ jest zwarty i zawiera się
w $E$. Ponadto
$E\setminus K\subseteq(E\setminus H)
\cup((E\cap H)\cap U)$, a drugi składnik zawiera się
w $U\setminus F$. Z podaddytywności
$\lambda_d(E\setminus K)<\varepsilon$.

Jeśli $\lambda_d(E)=\infty$, to z
$E\cap[-N,N]^d\uparrow E$ wynika
$\sup_N\lambda_d(E\cap[-N,N]^d)=\infty$.
Każdy z tych przecięć ma miarę skończoną, więc
poprzedni argument daje zawarty w nim zwarty podzbiór
o mierze dowolnie bliskiej mierze przecięcia. Miary
zwartych podzbiorów $E$ nie są zatem ograniczone.
Z drugiej strony $\lambda_d(K)\le\lambda_d(E)$ dla
każdego $K\subseteq E$, co kończy dowód.
\end{proof}

\begin{twierdzenie}[Borelowskie aproksymacje i uzupełnienie]
\label{thm:lebesgue-completion}
Dla każdego $E\in\mathcal L_d$ istnieją
zbiór $G$ typu $G_\delta$ i zbiór $F$ typu $F_\sigma$
takie, że
\begin{equation}\label{eq:lebesgue-Borel-approx}
F\subseteq E\subseteq G,\qquad
\lambda_d(G\setminus F)=0.
\end{equation}
W konsekwencji $\mathcal L_d$ jest dokładnie
uzupełnieniem $\mathcal B(\mathbb R^d)$ względem
$\lambda_d^B$.
\end{twierdzenie}

\begin{proof}
Niech $H_n=[-n,n)^d$ i $E_n=E\cap H_n$.
Każdy $E_n$ ma miarę skończoną.
Dla $n,k\ge1$ z twierdzenia
\ref{thm:lebesgue-outer-regular} wybieramy
otwarty $U_{n,k}\supseteq E_n$ taki, że
$\lambda_d(U_{n,k}\setminus E_n)<2^{-n}/k$.
Kładziemy $U_k=\bigcup_{n\ge1}U_{n,k}$.
Zbiory $U_k$ są otwarte, zawierają $E$ i
\[
\lambda_d(U_k\setminus E)
\le\sum_{n\ge1}\lambda_d(U_{n,k}\setminus E)
\le\sum_{n\ge1}\lambda_d(U_{n,k}\setminus E_n)
<\frac1k.
\]
Zbiór $G=\bigcap_{k\ge1}U_k$ jest typu $G_\delta$,
zawiera $E$ i spełnia $\lambda_d(G\setminus E)=0$,
ponieważ $G\setminus E\subseteq U_k\setminus E$
dla każdego $k$.
Powtarzamy konstrukcję dla $\mathbb R^d\setminus E$,
otrzymując $G'\supseteq\mathbb R^d\setminus E$ typu
$G_\delta$ oraz $\lambda_d(G'\cap E)=0$.
Wówczas $F=\mathbb R^d\setminus G'$ jest typu
$F_\sigma$, $F\subseteq E$ i
$\lambda_d(E\setminus F)=0$. Zatem
$\lambda_d(G\setminus F)=0$.

Równość dziedziny $\mathcal L_d$ z uzupełnieniem
wynika z \eqref{eq:lebesgue-Borel-approx} i
zupełności $\lambda_d$: dowolny $E\in\mathcal L_d$
różni się od borelowskiego $F$ o podzbiór borelowskiego
zbioru zerowego $G\setminus F$. Odwrotnie, każdy
podzbiór borelowskiego zbioru zerowego jest
mierzalny w sensie Carathéodory'ego.
\end{proof}

\begin{wniosek}[Geometryczna charakterystyka mierzalności]
\label{cor:lebesgue-regularity-characterization}
Dla $E\subseteq\mathbb R^d$ następujące warunki
są równoważne:
\begin{enumerate}[label=\textup{(\roman*)}]
\item $E\in\mathcal L_d$;
\item dla każdego $\varepsilon>0$ istnieje zbiór
otwarty $U\supseteq E$, dla którego
$\lambda_d^*(U\setminus E)<\varepsilon$;
\item istnieje zbiór borelowski $B\supseteq E$
taki, że $\lambda_d^*(B\setminus E)=0$.
\end{enumerate}
\end{wniosek}

\begin{proof}
Implikacja (i)$\Rightarrow$(ii) dla zbiorów o miarze
nieskończonej wynika z zastosowania konstrukcji
$U_k$ z dowodu twierdzenia
\ref{thm:lebesgue-completion}; dla wszystkich $E$
zachodzi $\lambda_d(U_k\setminus E)<1/k$.
Implikacja (ii)$\Rightarrow$(iii) wymaga ostrożności:
wybieramy otwarte $U_k\supseteq E$ z
$\lambda_d^*(U_k\setminus E)<1/k$ i kładziemy
$B=\bigcap_{k\ge1}U_k$. Wtedy $B$ jest borelowski,
$E\subseteq B$ oraz
$\lambda_d^*(B\setminus E)\le1/k$ dla każdego $k$,
czyli wartość wynosi zero.
Dla (iii)$\Rightarrow$(i) zbiór $N=B\setminus E$
ma miarę zewnętrzną zero, zatem jest mierzalny
z twierdzenia \ref{thm:caratheodory}; wobec
$E=B\setminus N$ i mierzalności $B$ dostajemy tezę.
\end{proof}

\section{Pierwsze konsekwencje i przykłady}

\begin{przyklad}[Zbiory gęste miary zero]
Zbiór $\mathbb Q^d$ jest przeliczalny, a więc
$\lambda_d(\mathbb Q^d)=0$ mimo jego gęstości
w $\mathbb R^d$. Jednocześnie zbiór
$[0,1]^d\setminus\mathbb Q^d$ ma miarę $1$.
Gęstość topologiczna nie pozwala zatem wnioskować
o dodatniości miary ani o jej wielkości.
\end{przyklad}

\begin{przyklad}[Wykres funkcji Lipschitzowskiej]
\label{ex:lipschitz-graph-null}
Niech $f:[a,b]\to\mathbb R$ spełnia warunek
$|f(x)-f(y)|\le L|x-y|$.
Pokażemy, że wykres
$\Gamma_f=\{(x,f(x)):x\in[a,b]\}\subseteq\mathbb R^2$
ma miarę $\lambda_2$ równą zero.
Dla $a<b$ dzielimy $[a,b]$ na $N$ przedziałów
$[x_{j-1},x_j]$ długości $\delta=(b-a)/N$.
Na $j$-tym przedziale zachodzi
$|f(x)-f(x_{j-1})|\le L\delta$, więc część wykresu
leży w prostokącie o poziomej szerokości $\delta$
i pionowej wysokości $2L\delta$.
Dla każdego $\eta>0$ pokrywamy tę część wykresu
półotwartym prostopadłościanem o bokach długości
$\delta+2\eta$ i $2L\delta+2\eta$.
Biorąc $\eta=\delta^2$, otrzymujemy pokrycie o
łącznej powierzchni co najwyżej
\[
N(\delta+2\delta^2)(2L\delta+2\delta^2)
=(b-a)(1+2\delta)(2L\delta+2\delta^2)
\xrightarrow[N\to\infty]{}0.
\]
Zatem $\lambda_2^*(\Gamma_f)=0$, a z zupełności
$\Gamma_f$ jest mierzalny i ma miarę zero.
Ten argument nie potrzebuje jeszcze twierdzenia Fubiniego.
\end{przyklad}

\begin{przyklad}[Miara kul]
Niech $B(x,r)=\{y\in\mathbb R^d:\|y-x\|_2<r\}$
dla $x\in\mathbb R^d$ i $r>0$.
Z twierdzenia \ref{thm:invariance-lebesgue}
wynika $\lambda_d(B(x,r))=\omega_dr^d$,
gdzie $\omega_d=\lambda_d(B(0,1))$.
Liczba $\omega_d$ jest dodatnia i skończona, ponieważ
$(-1/\sqrt d,1/\sqrt d)^d\subseteq B(0,1)
\subseteq[-1,1]^d$. Wzór
$\omega_d=\pi^{d/2}/\Gamma(d/2+1)$ wymaga dodatkowego
rachunku całkowego i zostanie wyprowadzony później.
\end{przyklad}

\section{Zbiór Cantora: miara zero bez przeliczalności}

\begin{definicja}[Środkowotrzeci zbiór Cantora]
\label{def:cantor-set}
Niech $C_0=[0,1]$. Z $C_{n-1}$ otrzymujemy $C_n$,
usuwając z każdego z $2^{n-1}$ jego domkniętych
przedziałów składowych otwartą środkową trzecią część.
Definiujemy zbiór Cantora przez
$C=\bigcap_{n=0}^\infty C_n$.
\end{definicja}

\begin{center}
\begin{tikzpicture}[x=10.5cm,y=0.73cm]
\foreach \y/\lab in {0/{$C_0$},1/{$C_1$},2/{$C_2$},3/{$C_3$}} {
  \node[anchor=east] at (-0.045,-\y) {\lab};
}
\fill[blue!45] (0,-0.15) rectangle (1,0.15);
\foreach \a/\b in {0/0.333333,0.666667/1} {
  \fill[blue!45] (\a,-1.15) rectangle (\b,-0.85);
}
\foreach \a/\b in {0/0.111111,0.222222/0.333333,0.666667/0.777778,0.888889/1} {
  \fill[blue!45] (\a,-2.15) rectangle (\b,-1.85);
}
\foreach \a/\b in {0/0.037037,0.074074/0.111111,0.222222/0.259259,0.296296/0.333333,0.666667/0.703704,0.740741/0.777778,0.888889/0.925926,0.962963/1} {
  \fill[blue!45] (\a,-3.15) rectangle (\b,-2.85);
}
\node[anchor=north] at (0,-3.42) {$0$};
\node[anchor=north] at (1,-3.42) {$1$};
\end{tikzpicture}

\smallskip
\emph{Pierwsze cztery etapy konstrukcji środkowotrzeciego zbioru Cantora.
W każdym etapie pozostaje $2^n$ przedziałów długości $3^{-n}$.}
\end{center}

\begin{twierdzenie}[Podstawowe własności zbioru Cantora]
\label{thm:cantor-properties}
Zbiór Cantora $C\subseteq[0,1]$ jest niepusty, zwarty,
doskonały (zamknięty i bez punktów izolowanych),
nigdziegęsty (jego domknięcie ma puste wnętrze),
nieprzeliczalny i ma miarę Lebesgue'a zero.
Co więcej, jest w naturalny sposób homeomorficzny
z przestrzenią $\{0,1\}^{\mathbb N}$ wyposażoną
w topologię produktową. W tym przypadku topologia ta
jest generowana przez cylindry, czyli zbiory wszystkich
ciągów zero-jedynkowych o ustalonych pierwszych
$N$ wyrazach, dla dowolnego $N\ge1$.
\end{twierdzenie}

\begin{proof}
Każdy $C_n$ jest niepusty i zwarty, a zbiory są
malejące. Ich przekrój jest więc niepusty i zwarty.
Zbiór $C_n$ jest sumą $2^n$ domkniętych przedziałów
długości $3^{-n}$. Końce tych przedziałów należą do
$C$, ponieważ w dalszych etapach usuwamy wyłącznie
otwarte przedziały środkowe. Dla dowolnego $x\in C$
możemy znaleźć w dowolnie małym otoczeniu $x$
inny koniec przedziału składowego $C_n$.
Stąd $C$ nie ma punktów izolowanych.

Każdy przedział składowy $C_n$ ma długość $3^{-n}$.
Gdyby $C$ zawierał przedział otwarty dodatniej długości,
to musiałby on być zawarty w jednym z przedziałów
składowych każdego $C_n$, co jest niemożliwe dla
dostatecznie dużego $n$. Zatem wnętrze $C$ jest puste;
ponieważ $C$ jest zamknięty, jest nigdziegęsty.

Z monotoniczności miary i skończonej addytywności
otrzymujemy
$0\le\lambda_1(C)\le\lambda_1(C_n)=2^n3^{-n}=(2/3)^n$.
Przechodząc do granicy, dostajemy $\lambda_1(C)=0$.

Na każdym etapie wybrana przez punkt $x\in C$ lewa lub prawa
trzecia część przedziału wyznacza odpowiednio cyfrę $0$ lub $2$.
Długości przedziałów maleją do zera, a ich zagnieżdżenie
ma dokładnie jeden punkt wspólny. Zatem punkt $x\in C$
ma przedstawienie trójkowe
$x=\sum_{n\ge1}2\varepsilon_n3^{-n}$, gdzie
$\varepsilon_n\in\{0,1\}$. W obrębie cyfr $0$ i $2$
to przedstawienie jest jednoznaczne: jeżeli ciągi
różnią się po raz pierwszy na miejscu $k$, to różnica
składników z tego miejsca to $2\cdot3^{-k}$,
a możliwa suma różnic z późniejszych miejsc ma wartość
bezwzględną co najwyżej
$\sum_{n>k}2\cdot3^{-n}=3^{-k}$.
Odwzorowanie
$\Phi:\{0,1\}^{\mathbb N}\to C$ dane przez
$\Phi((\varepsilon_n))=\sum_n2\varepsilon_n3^{-n}$
jest więc bijekcją. Jeśli dwa ciągi mają pierwszych
$N$ cyfr jednakowych, ich obrazy różnią się o co
najwyżej $\sum_{n>N}2\cdot3^{-n}=3^{-N}$;
to dowodzi ciągłości $\Phi$ dla topologii produktowej.
Dziedzina jest zwarta. Można to stwierdzić bez użycia
twierdzenia Tychonoffa: metryka
$\rho(u,v)=\sum_{n\ge1}2^{-n}|u_n-v_n|$
wyznacza tę samą topologię co cylindry; z każdego ciągu
punktów można wybrać podciąg o stałej pierwszej
współrzędnej, następnie jego podciąg o stałej drugiej
współrzędnej itd. Przekątniowy podciąg jest zbieżny
w metryce $\rho$, co dowodzi zwartości metrycznej.
Ponieważ $C$ jest przestrzenią Hausdorffa,
ciągła bijekcja ze zwartej dziedziny jest homeomorfizmem.
Ponieważ $\{0,1\}^{\mathbb N}$ jest nieprzeliczalna,
tak samo jest z $C$.
\end{proof}

\begin{wniosek}[Lebesgue'owskie zbiory nieborelowskie]
\label{cor:nonborel-lebesgue}
Istnieją podzbiory $\mathbb R$ mierzalne w sensie
Lebesgue'a, które nie są borelowskie. W szczególności
$\mathcal B(\mathbb R)\subsetneq\mathcal L_1$.
\end{wniosek}

\begin{proof}
Zbiór Cantora $C$ jest borelowski, ma miarę zero
i ma moc continuum $\mathfrak c$.
Każdy jego podzbiór jest mierzalny w sensie Lebesgue'a,
ponieważ miara $\lambda_1$ jest zupełna.
Rodzina wszystkich podzbiorów $C$ ma moc
$2^{\mathfrak c}>\mathfrak c$ (twierdzenie Cantora),
podczas gdy rodzina zbiorów borelowskich prostej ma
moc co najwyżej $\mathfrak c$. Uzasadnijmy ostatnie
stwierdzenie. Rodzina przedziałów o końcach wymiernych
jest przeliczalną bazą topologii $\mathbb R$. Każdy
zbiór borelowski można opisać za pomocą
\emph{przeliczalnego kodu borelowskiego}: dobrze
ufundowanego drzewa, którego węzły są oznaczone
operacją dopełnienia, sumy przeliczalnej albo
jednym z przedziałów bazowych. Takie opisy otrzymuje
się indukcyjnie po etapach tworzenia najmniejszej
rodziny zamkniętej na dopełnienia i sumy przeliczalne.
Każde drzewo o przeliczalnie wielu węzłach jest
kodowane przez podzbiór zbioru wszystkich skończonych
ciągów liczb naturalnych wraz z przeliczalnym
etykietowaniem. Liczba takich kodów wynosi co najwyżej
$2^{\aleph_0}=\mathfrak c$. Wynika stąd, że zbiorów
borelowskich jest najwyżej $\mathfrak c$, zaś
podzbiorów $C$ jest $2^{\mathfrak c}$; na mocy
twierdzenia Cantora $2^{\mathfrak c}>\mathfrak c$.
Zatem nie każdy podzbiór $C$ jest borelowski.
\end{proof}

\section{Zbiór dodatniej miary bez wnętrza}

\begin{przyklad}[Zbiór Smitha--Volterry--Cantora]
\label{ex:fat-cantor}
Rozpoczynamy od $F_0=[0,1]$. Na etapie $n\ge1$
z każdego z $2^{n-1}$ przedziałów pozostawionych
w poprzednim etapie usuwamy otwarty przedział środkowy
długości $4^{-n}$; pozostały zbiór oznaczamy $F_n$.
Można to wykonać na każdym etapie: długość każdego
z dwóch potomnych przedziałów jest dodatnia.
W rzeczy samej, jeśli $\ell_n$ oznacza długość jednego
przedziału składowego $F_n$, to
$\ell_0=1$ i $\ell_n=(\ell_{n-1}-4^{-n})/2$.
Indukcja daje dokładny wzór
$\ell_n=2^{-n-1}+2^{-2n-1}>0$; co więcej,
$\ell_{n-1}>4^{-n}$, a zatem każde kolejne
usunięcie jest geometrycznie możliwe.

Łączna długość usuniętych przedziałów na etapie $n$
wynosi $2^{n-1}4^{-n}=2^{-n-1}$. Ponieważ otwarte
przedziały usuwane w różnych etapach są rozłączne,
dla $F=\bigcap_nF_n$ mamy
\[
\lambda_1(F)=1-\sum_{n=1}^\infty2^{-n-1}=\frac12.
\]
Zbiór $F$ jest zwarty. Ponieważ długości przedziałów
składowych $F_n$ dążą do zera, $F$ nie zawiera
niepustego przedziału otwartego, jest więc nigdziegęsty.
Końce każdego przedziału składowego pozostają w $F$,
a zatem, podobnie jak dla zbioru Cantora, żaden punkt
$F$ nie jest izolowany. Zbiór $F$ jest więc doskonały
i ma dodatnią miarę. A zatem nawet jednoczesna
zwartość, brak wnętrza i nigdziegęstość nie
wymuszają miary zero.
\end{przyklad}

\section{Funkcja Cantora i miara osobliwa: zapowiedź}

\begin{przyklad}[Funkcja Cantora]
\label{ex:cantor-function}
Dla $x\in C$ o przedstawieniu
$x=\sum_{n\ge1}2\varepsilon_n3^{-n}$ określamy
$F(x)=\sum_{n\ge1}\varepsilon_n2^{-n}$.
Na każdej składowej otwartej $[0,1]\setminus C$
przedłużamy $F$ jako funkcję stałą, przyjmując
wspólną wartość na końcach tej składowej.
Wartości końcowe są równe: końce przedziału
usuniętego na etapie $n$ mają przedstawienia
trójkowe, których pierwsze $n-1$ cyfr są jednakowe,
a pozostałe cyfry są odpowiednio $0,2,2,\ldots$
oraz $2,0,0,\ldots$, co w dwójkowym zapisie
prowadzi do tej samej liczby.

Tak uzyskana funkcja $F:[0,1]\to[0,1]$
jest ciągła i niemalejąca. Można to sprawdzić za
pomocą kolejnych ciągłych przybliżeń $F_n$:
na każdym przedziale składowym $C_n$ funkcja $F_n$
jest liniowa, a na przedziałach już usuniętych
stała, przy czym $F_n(0)=0$ oraz $F_n(1)=1$.
Dokładniej: na każdym z $2^n$ domkniętych przedziałów
składowych $C_n$ wartości na końcach różnią się o
$2^{-n}$, a na przedziałach już usuniętych $F_n$
pozostaje stała. Z konstrukcji wynika, że
$\|F_{n+1}-F_n\|_\infty\le 2^{-n}$.
Dla $m>n$ dostajemy
$\|F_m-F_n\|_\infty\le\sum_{k=n}^{m-1}2^{-k}$,
więc $(F_n)$ jest jednostajnym ciągiem Cauchy'ego.
Jego granica jest ciągła i niemalejąca. Na końcach
każdego usuniętego przedziału granica przyjmuje
wspólną, raz ustaloną wartość, a na $C$ zgadza się
z szeregiem definiującym $F$, ponieważ na przedziale
składowym $C_n$ zarówno $F_n$, jak i wartość tego
szeregu leżą między tymi samymi wartościami końcowymi,
oddalonymi o $2^{-n}$. To identyfikuje granicę
jednostajną z funkcją $F$ z definicji.

Ponadto $F$ jest stała na każdym przedziale
składowym $[0,1]\setminus C$, więc $F'(x)=0$
dla wszystkich $x\in(0,1)\setminus C$.
Ponieważ $\lambda_1(C)=0$, mamy $F'=0$
prawie wszędzie, ale $F(1)-F(0)=1$.
Nie jest to sprzeczność: funkcja Cantora nie jest
absolutnie ciągła. Przypomnijmy, że funkcja
$h:[a,b]\to\mathbb R$ jest absolutnie ciągła, jeśli
dla każdego $\varepsilon>0$ istnieje $\delta>0$
takie, że dla każdej skończonej rodziny parami
rozłącznych przedziałów $(a_k,b_k)\subseteq[a,b]$
warunek $\sum_k(b_k-a_k)<\delta$ implikuje
$\sum_k|h(b_k)-h(a_k)|<\varepsilon$.
Dla funkcji Cantora bierzemy $2^n$ przedziałów
składowych $C_n$. Suma ich długości wynosi
$(2/3)^n\to0$, podczas gdy suma przyrostów $F$
na tych przedziałach wynosi
$2^n\cdot2^{-n}=1$. Sprzeczność z powyższą definicją
dla $\varepsilon=1/2$ dowodzi braku absolutnej ciągłości.
W rozdziałach o całce Lebesgue'a i rozkładzie miar
zobaczymy, że $F$ wyznacza bezatomową miarę
osobliwą względem miary Lebesgue'a, to znaczy
skoncentrowaną na zbiorze o mierze Lebesgue'a zero.
\end{przyklad}

\section{Zbiór Vitaliego i ograniczenia mierzalności}

\begin{twierdzenie}[Istnienie zbioru niemierzalnego]
\label{thm:vitali}
Przy aksjomacie wyboru istnieje podzbiór
$V\subseteq[0,1]$, który nie jest mierzalny
w sensie Lebesgue'a.
\end{twierdzenie}

\begin{proof}
Na $[0,1]$ wprowadzamy relację równoważności
$x\sim y$ wtedy i tylko wtedy, gdy $x-y\in\mathbb Q$.
Z każdej klasy wybieramy dokładnie jeden element,
tworząc zbiór reprezentantów $V\subseteq[0,1]$.
Ten krok wykorzystuje aksjomat wyboru.

Niech $Q=\mathbb Q\cap[-1,1]$. Zbiory $V+q$ dla
$q\in Q$ są parami rozłączne. Rzeczywiście, gdyby
$v_1+q_1=v_2+q_2$, to
$v_1-v_2=q_2-q_1\in\mathbb Q$, więc $v_1,v_2$
należałyby do tej samej klasy; wybór reprezentantów
daje $v_1=v_2$ i następnie $q_1=q_2$.

Mamy też
\begin{equation}\label{eq:vitali-sandwich}
[0,1]\subseteq\bigcup_{q\in Q}(V+q)\subseteq[-1,2].
\end{equation}
Pierwsze zawieranie wynika stąd, że dla dowolnego
$x\in[0,1]$ istnieje $v\in V$ z $x-v\in\mathbb Q$,
a jednocześnie $|x-v|\le1$.

Przypuśćmy, że $V$ jest mierzalny. Z niezmienniczości
translacyjnej $\lambda_1(V+q)=\lambda_1(V)$.
Zbiory $V+q$ są rozłączne, $Q$ jest przeliczalny i
nieskończony, a więc
\[
\lambda_1\Big(\bigcup_{q\in Q}(V+q)\Big)
=\sum_{q\in Q}\lambda_1(V)
=\begin{cases}0,&\lambda_1(V)=0,\\
\infty,&\lambda_1(V)>0.
\end{cases}
\]
Z \eqref{eq:vitali-sandwich} wynika jednak, że lewa
strona znajduje się między $1$ a $3$. To sprzeczność.
Zatem $V\notin\mathcal L_1$.
\end{proof}

\begin{wniosek}[Nie istnieje miara długości wszystkich podzbiorów]
\label{cor:no-all-subsets-invariant}
Nie istnieje przeliczalnie addytywna miara
$\nu:\mathcal P(\mathbb R)\to[0,\infty]$, która
jest niezmiennicza względem translacji i spełnia
$\nu([0,1])=1$.
\end{wniosek}

\begin{proof}
Najpierw musimy ustalić $\nu(\{0\})=0$.
Gdyby $\nu(\{0\})>0$, przeliczalnie wiele
rozłącznych translacji singletonu zawartych w $[0,1]$
miałoby miarę nieskończoną, sprzecznie z
$\nu([0,1])=1$. Zatem $\nu(\{0\})=0$,
a translacyjna niezmienniczość daje
$\nu(\{x\})=0$ dla każdego $x\in\mathbb R$.
Wobec tego $\nu([0,1))=\nu([0,1])=1$.
Trzy rozłączne przedziały $[-1,0)$, $[0,1)$
i $[1,2)$ mają po jednej jednostce miary;
po dodaniu punktu $2$ otrzymujemy
$\nu([-1,2])=3$. Teraz dowód Vitaliego stosuje się
bez zmian do $\nu$, ponieważ wszystkie podzbiory
w nim występujące należałyby do dziedziny $\nu$.
Otrzymujemy sprzeczność.
\end{proof}

\begin{uwaga}[Precyzyjny sens niemożliwości]
Powyższe twierdzenie nie zabrania istnienia miar na
$\mathcal P(\mathbb R)$, takich jak miara licząca
czy miara zerowa. Stwierdza niemożliwość jednoczesnego
zachowania \emph{przeliczalnej addytywności},
\emph{niezmienniczości translacyjnej} i
\emph{normalizacji długości przedziału}.
\end{uwaga}

\section{Podsumowanie i przygotowanie do całki}

Wykazaliśmy, że objętość prostopadłościanów
półotwartych wyznacza jedną miarę borelowską,
której uzupełnieniem jest miara Lebesgue'a. Zbiory
mierzalne mają aproksymacje otwarte i zwarte.
Gęstość, nigdziegęstość czy nieprzeliczalność
nie przesądzają o wielkości miary zbioru.

W następnych rozdziałach zdefiniujemy funkcje
mierzalne i całkę oraz wyjaśnimy rolę absolutnej
ciągłości w związku między całką a pochodną.

% END chapters/pl/measure_theory/04-rozdzial.tex


% BEGIN chapters/pl/measure_theory/05-rozdzial.tex
% ROZDZIAŁ 5 — plik włączany przez main.tex
\chapter{Struktura przestrzeni mierzalnych}
\label{ch:struktura-przestrzeni}

\section{Dlaczego sama definicja \texorpdfstring{$\sigma$}{sigma}-algebry nie wystarcza?}

Znamy już konstrukcję $\sigma$-algebr, miar i miary Lebesgue'a.
W dalszej teorii będziemy jednak przenosić miary przez odwzorowania,
budować przestrzenie produktowe i określać zbiory na podstawie
współrzędnych. Chcemy więc wiedzieć, jak z danych o prostych
zbiorach odzyskać mierzalność dowolnego odwzorowania.

W tym rozdziale przez $(X,\mathcal F)$ i $(Y,\mathcal G)$ rozumiemy
przestrzenie mierzalne. Nie zakładamy istnienia na nich miary,
jeśli nie podajemy jej wyraźnie. Zbiór $\mathbb N$ składa się z
dodatnich liczb całkowitych. Używamy oznaczeń
$\mathcal B(T)$ dla $\sigma$-algebry borelowskiej przestrzeni
topologicznej $T$ oraz $\mathcal F|_A=\{E\cap A:E\in\mathcal F\}$
dla śladu $\sigma$-algebry na $A\subseteq X$.

\section{Ślad \texorpdfstring{$\sigma$}{sigma}-algebry i podprzestrzenie}

Ślad $\mathcal F|_A$ zdefiniowaliśmy już w rozdziale~\ref{ch:przestrzenie-mierzalne}.
Teraz badamy jego zgodność z generowaniem $\sigma$-algebr
i topologią podprzestrzeni.


\begin{lemat}[Ślad i generowanie]
\label{lem:trace-generation}
Niech $X$ będzie zbiorem, $\mathcal E\subseteq\mathcal P(X)$,
$A\subseteq X$ i $\mathcal F=\sigma_X(\mathcal E)$.
Wtedy
\[
\mathcal F|_A=\sigma_A\bigl(\{A\cap E:E\in\mathcal E\}\bigr).
\]
\end{lemat}
\begin{proof}
Oznaczmy prawą stronę przez $\mathcal H$. Zbiory
$A\cap E$ dla $E\in\mathcal E$ należą do $\mathcal F|_A$;
ponieważ $\mathcal F|_A$ jest $\sigma$-algebrą na $A$,
mamy $\mathcal H\subseteq\mathcal F|_A$.
W drugą stronę rozważmy rodzinę
$\mathcal D=\{E\subseteq X:A\cap E\in\mathcal H\}$.
Zawiera ona $X$, gdyż $A\in\mathcal H$. Jeżeli
$E\in\mathcal D$, to
$A\cap(X\setminus E)=A\setminus(A\cap E)\in\mathcal H$.
Dla $E_n\in\mathcal D$ mamy
$A\cap\bigcup_nE_n=\bigcup_n(A\cap E_n)\in\mathcal H$.
Zatem $\mathcal D$ jest $\sigma$-algebrą na $X$, zawierającą
$\mathcal E$. Wynika stąd
$\mathcal F=\sigma_X(\mathcal E)\subseteq\mathcal D$, czyli
$\mathcal F|_A\subseteq\mathcal H$.
\end{proof}

\begin{twierdzenie}[Zbiory borelowskie podprzestrzeni]
\label{thm:borel-trace}
Niech $T$ będzie przestrzenią topologiczną i $A\subseteq T$
wyposażonym w topologię podprzestrzeni. Wtedy
$\mathcal B(A)=\mathcal B(T)|_A$.
\end{twierdzenie}
\begin{proof}
Z definicji topologii podprzestrzeni zbiory otwarte w $A$
mają postać $A\cap U$, gdzie $U$ jest otwarty w $T$.
Zastosowanie lematu \ref{lem:trace-generation} do rodziny
wszystkich zbiorów otwartych przestrzeni $T$ daje tezę.
\end{proof}

\begin{przyklad}[Ślad na zbiorze nieborelowskim]
Istnieje zbiór $A\subseteq\mathbb R$ nieborelowski
(wniosek \ref{cor:nonborel-lebesgue}). Mimo to możemy
rozpatrywać $(A,\mathcal B(\mathbb R)|_A)$.
To jest przestrzeń mierzalna; brak borelowskości $A$
w przestrzeni otaczającej nie uniemożliwia
zdefiniowania śladu. Nie wolno jednak utożsamiać
$\mathcal B(\mathbb R)|_A$ z rodziną
$\{E\in\mathcal B(\mathbb R):E\subseteq A\}$:
ta druga rodzina może nawet nie zawierać całego $A$.
\end{przyklad}

\section{Mierzalność odwzorowania i rodziny generujące}

Korzystamy z definicji mierzalności~\ref{def:measurable-map} i kryterium generatorów~\ref{thm:measurable-generator}, udowodnionych w rozdziale~\ref{ch:przestrzenie-mierzalne}. Ich podstawą są tożsamości dla przeciwobrazów: $f^{-1}(Y\setminus B)=X\setminus f^{-1}(B)$ oraz $f^{-1}(\bigcup_n B_n)=\bigcup_n f^{-1}(B_n)$. Nie wymagają one różnowartościowości ani surjektywności.

\begin{wniosek}[Prawostronny test mierzalności]
\label{cor:real-borel-tests}
Niech $(X,\mathcal F)$ będzie przestrzenią mierzalną,
a $f:X\to\mathbb R$ odwzorowaniem. Jest ono
$(\mathcal F,\mathcal B(\mathbb R))$-mierzalne wtedy i tylko
wtedy, gdy $\{x\in X:f(x)>q\}\in\mathcal F$
dla każdego $q\in\mathbb Q$.
\end{wniosek}
\begin{proof}
Rodzina prawych półprostych $(q,\infty)$, $q\in\mathbb Q$,
generuje $\mathcal B(\mathbb R)$: przez dopełnienia
otrzymujemy lewe półproste domknięte, a następnie
wszystkie otwarte przedziały. Stosujemy więc
kryterium generatorów~\ref{thm:measurable-generator}.
Warunki z lewymi półprostymi zostały już udowodnione
w twierdzeniu~\ref{thm:measurable-real}.
\end{proof}

\begin{twierdzenie}[Mierzalność ograniczenia]
\label{thm:restriction-meas-5}
Niech $(X,\mathcal F)$ i $(Y,\mathcal G)$ będą przestrzeniami
mierzalnymi, $f:X\to Y$ odwzorowaniem
$(\mathcal F,\mathcal G)$-mierzalnym oraz $A\subseteq X$.
Wówczas $f|_A:(A,\mathcal F|_A)\to(Y,\mathcal G)$ jest mierzalne.
\end{twierdzenie}
\begin{proof}
Dla każdego $B\in\mathcal G$ mamy
$(f|_A)^{-1}(B)=A\cap f^{-1}(B)\in\mathcal F|_A$.
Złożenia odwzorowań mierzalnych oraz mierzalność funkcji
ciągłych zostały już wykazane w twierdzeniu~\ref{thm:composition-measurable}.
\end{proof}

\section{Najmniejsza \texorpdfstring{$\sigma$}{sigma}-algebra, względem której funkcja jest mierzalna}

\begin{definicja}[$\sigma$-algebra generowana przez odwzorowanie]
\label{def:sigma-of-function}
Dla odwzorowania $f:X\to(Y,\mathcal G)$ określamy
\[
\sigma(f):=\{f^{-1}(B):B\in\mathcal G\}.
\]
Jeżeli $(f_i)_{i\in I}$ jest rodziną odwzorowań
$f_i:X\to(Y_i,\mathcal G_i)$, to
$\sigma(f_i:i\in I):=\sigma_X\bigl(\bigcup_{i\in I}\sigma(f_i)\bigr)$.
\end{definicja}

\begin{twierdzenie}[Minimalna struktura mierzalna]
\label{thm:initial-sigma}
Dla każdego $f:X\to(Y,\mathcal G)$ rodzina $\sigma(f)$
jest $\sigma$-algebrą na $X$ i jest najmniejszą
$\sigma$-algebrą, względem której $f$ jest mierzalne.
Dla rodziny $(f_i)_{i\in I}$ analogiczną własność
ma $\sigma(f_i:i\in I)$.
\end{twierdzenie}
\begin{proof}
Przeciwobraz całego $Y$ to $X$, przeciwobrazy zachowują
dopełnienia i przeliczalne sumy, więc $\sigma(f)$
jest $\sigma$-algebrą. Jeśli $f$ jest mierzalne
względem pewnej $\mathcal F$, to z definicji
$f^{-1}(B)\in\mathcal F$ dla każdego $B\in\mathcal G$,
toteż $\sigma(f)\subseteq\mathcal F$.
Dla rodziny odwzorowań powtarzamy argument dla
$\sigma$-algebry generowanej sumą wszystkich
rodzin przeciwobrazów.
\end{proof}

\begin{przyklad}[Rozróżnialność punktów]
Dla $X=\{a,b,c,d\}$ i funkcji $f:X\to\mathbb R$ o
$f(a)=f(b)=0$ oraz $f(c)=f(d)=1$ otrzymujemy
$\sigma(f)=\{\varnothing,\{a,b\},\{c,d\},X\}$.
Funkcja nie potrafi odróżnić $a$ od $b$ ani $c$ od $d$.
To wyjaśnia, dlaczego $\sigma(f)$ może być ściśle mniejsza
od $\mathcal P(X)$.
\end{przyklad}

\section{Iloczyn przestrzeni mierzalnych}

\begin{definicja}[$\sigma$-algebra produktowa]
\label{def:product-sigma}
Dla przestrzeni mierzalnych $(X,\mathcal F)$ i
$(Y,\mathcal G)$ definiujemy
\[
\mathcal F\otimes\mathcal G
:=\sigma_{X\times Y}\bigl(\{A\times B:A\in\mathcal F,\ B\in\mathcal G\}\bigr).
\]
Parę $(X\times Y,\mathcal F\otimes\mathcal G)$
nazywamy \emph{iloczynem przestrzeni mierzalnych}.
Nie wprowadzamy w tym miejscu miary produktowej;
jej konstrukcja należy do planowanego rozdziału 12.
\end{definicja}

\begin{twierdzenie}[Charakteryzacja przez współrzędne]
\label{thm:product-projections}
Niech $(X,\mathcal F)$ i $(Y,\mathcal G)$ będą
przestrzeniami mierzalnymi, a
$\pi_X:X\times Y\to X$, $\pi_Y:X\times Y\to Y$
rzutowaniami. Wówczas
$\mathcal F\otimes\mathcal G=\sigma(\pi_X,\pi_Y)$.
Ponadto odwzorowanie $h:Z\to X\times Y$ z przestrzeni
mierzalnej $(Z,\mathcal H)$ jest
$(\mathcal H,\mathcal F\otimes\mathcal G)$-mierzalne
wtedy i tylko wtedy, gdy oba odwzorowania
$\pi_X\circ h:Z\to X$ i $\pi_Y\circ h:Z\to Y$ są mierzalne.
\end{twierdzenie}
\begin{proof}
Ponieważ $\pi_X^{-1}(A)=A\times Y$ i
$\pi_Y^{-1}(B)=X\times B$, każda z rodzin
$\sigma(\pi_X)$, $\sigma(\pi_Y)$ zawiera się w
$\mathcal F\otimes\mathcal G$. Z drugiej strony
$A\times B=\pi_X^{-1}(A)\cap\pi_Y^{-1}(B)$,
zatem każdy generator $\mathcal F\otimes\mathcal G$
należy do $\sigma(\pi_X,\pi_Y)$. To dowodzi równości.
Drugą część otrzymujemy, sprawdzając przeciwobrazy
generatorów $A\times B$:
$h^{-1}(A\times B)=(\pi_X\circ h)^{-1}(A)
\cap(\pi_Y\circ h)^{-1}(B)$.
Konieczność wynika z mierzalności rzutowań i złożenia.
\end{proof}

\begin{twierdzenie}[Produkt borelowski przestrzeni drugoprzeliczalnych]
\label{thm:borel-product}
Niech $S,T$ będą przestrzeniami topologicznymi, z których
każda ma przeliczalną bazę topologii. Wtedy
\[
\mathcal B(S\times T)=\mathcal B(S)\otimes\mathcal B(T),
\]
gdzie na $S\times T$ rozważamy topologię produktową.
\end{twierdzenie}
\begin{proof}
Rzutowania $\pi_S,\pi_T$ są ciągłe, więc
$A\times T=\pi_S^{-1}(A)$ i $S\times B=\pi_T^{-1}(B)$
są borelowskie dla borelowskich $A,B$.
Stąd $\mathcal B(S)\otimes\mathcal B(T)\subseteq
\mathcal B(S\times T)$.
Wybierzmy przeliczalne bazy $\mathcal U,\mathcal V$.
Wszystkie $U\times V$, gdzie $U\in\mathcal U$ i
$V\in\mathcal V$, należą do
$\mathcal B(S)\otimes\mathcal B(T)$ i tworzą
przeliczalną bazę produktu. Każdy zbiór otwarty
w $S\times T$ jest więc przeliczalną sumą takich
prostokątów, czyli należy do $\mathcal B(S)\otimes\mathcal B(T)$.
Generując $\sigma$-algebrę, otrzymujemy zawieranie odwrotne.
\end{proof}

\begin{wniosek}[Borelowskie odwzorowania wektorowe]
\label{cor:coordinate-measurable}
Niech $(X,\mathcal F)$ będzie przestrzenią mierzalną,
$f=(f_1,\ldots,f_d):X\to\mathbb R^d$ i $d\ge1$.
Odwzorowanie $f$ jest $(\mathcal F,\mathcal B(\mathbb R^d))$-mierzalne
wtedy i tylko wtedy, gdy każda $f_j:X\to\mathbb R$
jest $(\mathcal F,\mathcal B(\mathbb R))$-mierzalna.
\end{wniosek}
\begin{proof}
Z twierdzenia \ref{thm:borel-product}, zastosowanego
indukcyjnie, wynika
$\mathcal B(\mathbb R^d)=\bigotimes_{j=1}^d\mathcal B(\mathbb R)$.
Wykorzystujemy teraz twierdzenie \ref{thm:product-projections}.
\end{proof}

\begin{uwaga}[Dlaczego potrzebna jest przeliczalna baza?]
Dla dowolnych przestrzeni topologicznych $S,T$ zachodzi
$\mathcal B(S)\otimes\mathcal B(T)\subseteq\mathcal B(S\times T)$,
ponieważ rzutowania są ciągłe. Równość może jednak zawieść.
Dla przykładu niech $D$ będzie zbiorem o mocy większej niż
$\mathfrak c=|\mathbb R|$, wyposażonym w topologię dyskretną.
Wtedy przekątna $\Delta=\{(x,x):x\in D\}$ jest zamknięta
w $D\times D$, ale nie należy do
$\mathcal P(D)\otimes\mathcal P(D)$.
Istotnie, każdy element $\sigma$-algebry generowanej
przez dowolną rodzinę zbiorów należy do $\sigma$-algebry
generowanej przez pewną jej przeliczalną podrodzinę:
rodzina zbiorów o tej własności sama jest $\sigma$-algebrą.
Gdyby $\Delta$ była mierzalna produktowo, istniałyby
prostokąty $A_n\times B_n$, $n\ge1$, generujące
$\sigma$-algebrę zawierającą $\Delta$.
Dla różnych $x,y\in D$ punkty $(x,x)$ i $(x,y)$
muszą być rozróżnione przez któryś z tych prostokątów,
inaczej wszystkie zbiory z wygenerowanej $\sigma$-algebry
zawierałyby oba albo żaden z nich. Ponieważ pierwsza
współrzędna obu punktów jest taka sama, pewien $B_n$
musi zawierać dokładnie jeden z punktów $x,y$.
Odwzorowanie $x\mapsto(\chi_{B_n}(x))_{n\ge1}$
byłoby więc iniekcją $D\to\{0,1\}^{\mathbb N}$, co
przeczy $|D|>\mathfrak c$.
\end{uwaga}

\section{Przestrzenie standardowe borelowskie}

\begin{definicja}[Przestrzeń polska i standardowa borelowska]
\label{def:standard-borel}
Przestrzeń topologiczną nazywamy \emph{polską}, jeśli
jej topologię wyznacza pewna metryka zupełna i ośrodek
przeliczalny, czyli istnieje w niej przeliczalny zbiór
gęsty. Przestrzeń mierzalną $(X,\mathcal F)$ nazywamy
\emph{standardową borelowską}, jeżeli istnieje przestrzeń
polska $P$ i bijekcja $b:X\to P$, dla której
$b$ oraz $b^{-1}$ są mierzalne względem
$\mathcal F$ i $\mathcal B(P)$. Taką bijekcję
nazywamy \emph{izomorfizmem przestrzeni mierzalnych}.
\end{definicja}

\begin{uwaga}[Uwaga o definicji]
Równoważnie można zażądać, aby $\mathcal F$ była
$\sigma$-algebrą borelowską pewnej polskiej topologii na $X$.
Istotny jest charakter \emph{mierzalny}, a nie zgodność
z jakąkolwiek wcześniej wybraną topologią. W definicji
żądaliśmy jawnie mierzalności odwrotności; dla ogólnych
przestrzeni mierzalnych sama mierzalna bijekcja nie wystarcza.
\end{uwaga}

\begin{przyklad}[Klasyczne przestrzenie standardowe]
Przestrzenie $(\mathbb R^d,\mathcal B(\mathbb R^d))$
($d\ge1$), $(\mathbb N,\mathcal P(\mathbb N))$ oraz
$(\{0,1\}^{\mathbb N},\mathcal B(\{0,1\}^{\mathbb N}))$
są standardowe borelowskie. W pierwszych dwóch przypadkach
korzystamy ze zwykłej metryki euklidesowej lub metryki
dyskretnej. W trzecim stosujemy metrykę
$\rho(u,v)=\sum_{n\ge1}2^{-n}|u_n-v_n|$ z rozdziału
\ref{ch:lebesgue}. Zupełność wynika stąd, że w ciągu
Cauchy'ego każda współrzędna musi od pewnego miejsca
być stała; przeliczalnym zbiorem gęstym są ciągi
ostatecznie równe zeru. Metryka $\rho$ wyznacza topologię
produktową, więc mamy przestrzeń polską.
\end{przyklad}

\begin{twierdzenie}[Przeliczalne generowanie standardowych przestrzeni]
\label{thm:standard-countably-generated}
Jeżeli $(X,\mathcal F)$ jest standardową przestrzenią
borelowską, to istnieje przeliczalna rodzina
$\mathcal E\subseteq\mathcal F$ taka, że
$\mathcal F=\sigma_X(\mathcal E)$. Ponadto dla
każdej pary różnych punktów $x,y\in X$ istnieje
$E\in\mathcal E$ zawierający dokładnie jeden z nich.
\end{twierdzenie}
\begin{proof}
Niech $b:X\to P$ będzie izomorfizmem na przestrzeń polską.
Każda przestrzeń metryczna o przeliczalnym zbiorze gęstym
ma przeliczalną bazę: bierzemy kule o środkach w tym zbiorze
i dodatnich promieniach wymiernych. Niech
$\mathcal U=\{U_n:n\ge1\}$ będzie taką bazą $P$.
Rodzina $\mathcal E=\{b^{-1}(U_n):n\ge1\}$
jest przeliczalna i generuje $\mathcal F$, gdyż
$\mathcal B(P)=\sigma_P(\mathcal U)$, a przeciwobrazy
zachowują dopełnienia i sumy przeliczalne.
Jeśli $x\ne y$, to $b(x)\ne b(y)$; w przestrzeni
metrycznej istnieje zbiór otwarty zawierający $b(x)$,
lecz nie $b(y)$. Z własności bazy można wybrać
$U_n$ o tej samej własności, a $b^{-1}(U_n)$
oddziela $x$ od $y$.
\end{proof}

\begin{przyklad}[Uzupełnienie miary zmienia typ przestrzeni]
\label{ex:completion-not-standard}
Przestrzeń $(\mathbb R,\mathcal B(\mathbb R))$ jest
standardowa borelowska, ale
$(\mathbb R,\mathcal L_1)$ nie jest standardowa borelowska.
Z twierdzenia \ref{thm:standard-countably-generated}
każda standardowa $\sigma$-algebra jest przeliczalnie
generowana. Każda $\sigma$-algebra generowana przez
przeliczalną rodzinę ma co najwyżej $\mathfrak c$
elementów (kodowanie zbiorów borelowskich drzewami,
jak w dowodzie wniosku \ref{cor:nonborel-lebesgue}).
Tymczasem $\mathcal L_1$ zawiera wszystkie podzbiory
zbioru Cantora, a więc co najmniej $2^{\mathfrak c}>
\mathfrak c$ elementów. Nie może być przeliczalnie
generowana. Przykład pokazuje, że \emph{zupełność miary}
i \emph{standardowość borelowska} są innymi własnościami.
\end{przyklad}

\section{Struktura mierzalna a miara obrazowa}

Przypomnijmy, że jeśli $\mu$ jest miarą na $(X,\mathcal F)$,
a $f:(X,\mathcal F)\to(Y,\mathcal G)$ jest mierzalne,
to \emph{miara obrazowa} $f_\#\mu$ określona jest wzorem
$(f_\#\mu)(B)=\mu(f^{-1}(B))$, $B\in\mathcal G$
(definicja z rozdziału \ref{ch:miary-wlasnosci}).
To poprawne określenie miary, ponieważ przeciwobrazy
zachowują rozłączność i przeliczalne sumy.

\begin{twierdzenie}[Funktorialność miar obrazowych]
\label{thm:pushforward-composition}
Niech $\mu$ będzie miarą na $(X,\mathcal F)$,
$f:(X,\mathcal F)\to(Y,\mathcal G)$ i
$g:(Y,\mathcal G)\to(Z,\mathcal H)$ niech będą mierzalne.
Wtedy
\[
(g\circ f)_\#\mu=g_\#(f_\#\mu).
\]
\end{twierdzenie}
\begin{proof}
Dla każdego $C\in\mathcal H$ z definicji i
rachunku przeciwobrazów
$(g\circ f)_\#\mu(C)=\mu(f^{-1}(g^{-1}(C)))
=(f_\#\mu)(g^{-1}(C))=g_\#(f_\#\mu)(C)$.
Równość miar oznacza tu równość ich wartości
na wszystkich zbiorach mierzalnych.
\end{proof}

\begin{przyklad}[Miara obrazowa przez funkcję stałą]
Niech $f:X\to Y$ będzie funkcją stałą o wartości
$y_0\in Y$ i niech $\mu$ będzie miarą na $(X,\mathcal F)$.
Wtedy $(f_\#\mu)(B)=\mu(X)$, gdy $y_0\in B$, oraz
$(f_\#\mu)(B)=0$, gdy $y_0\notin B$. Jeśli $\mu(X)<\infty$, można zapisać
$f_\#\mu=\mu(X)\delta_{y_0}$, gdzie $\delta_{y_0}$
wyraża miarę Diraca w punkcie $y_0$.
Gdy $\mu(X)=\infty$, pozostajemy przy podanym wzorze
zbiorowym: wyrażenie $\infty\cdot0$ nie jest w zwykłej
arytmetyce rozszerzonej określone.
\end{przyklad}

\section{Podsumowanie i powiązania}

Rodziny generujące upraszczają testy mierzalności.
Ślad opisuje przestrzenie indukowane, a produkt
sprowadza mierzalność wektorową do współrzędnych.
Ważnym rozróżnieniem jest różnica między przestrzenią
standardową borelowską a uzupełnioną przestrzenią
z miarą Lebesgue'a.

W następnym rozdziale wykorzystamy te narzędzia do
mierzalności funkcji rzeczywistych, ich aproksymacji
funkcjami prostymi i badania granic ciągów funkcji.

% END chapters/pl/measure_theory/05-rozdzial.tex


% BEGIN chapters/pl/measure_theory/06-rozdzial.tex
% ROZDZIAŁ 6 — plik włączany przez main.tex
\chapter{Funkcje mierzalne i ich aproksymacja}
\label{ch:funkcje-mierzalne}

\section{Funkcje jako podstawowy obiekt teorii miary}

Dotychczas mierzyliśmy zbiory. W praktyce chcemy mierzyć
również zbiory postaci $\{x:f(x)>a\}$, przybliżać
funkcje funkcjami o skończenie wielu wartościach, a następnie
określić całkę. Właśnie dlatego potrzebujemy pojęcia
funkcji mierzalnej. Ważne, by odróżniać \emph{mierzalność}
od \emph{całkowalności}: pierwsza pozwala dopiero
postawić drugie pytanie.

W całym rozdziale $(X,\mathcal F)$ oznacza przestrzeń mierzalną.
Jeśli potrzebna jest miara, zapisujemy $(X,\mathcal F,\mu)$.
Przez $\overline{\mathbb R}=[-\infty,\infty]$ rozumiemy
rozszerzoną prostą rzeczywistą wyposażoną w topologię
porządkową; jej $\sigma$-algebrą jest
$\mathcal B(\overline{\mathbb R})$. Każda funkcja
$\mathbb R$-wartościowa może być traktowana również jako
$\overline{\mathbb R}$-wartościowa.

\section{Funkcje rzeczywiste i rozszerzone rzeczywiste}

\begin{definicja}[Mierzalność funkcji rozszerzonej]
\label{def:extended-measurable}
Funkcję $f:X\to\overline{\mathbb R}$ nazywamy
\emph{mierzalną}, jeśli jest
$(\mathcal F,\mathcal B(\overline{\mathbb R}))$-mierzalna,
czyli $f^{-1}(B)\in\mathcal F$ dla każdego
$B\in\mathcal B(\overline{\mathbb R})$.
\end{definicja}

\begin{twierdzenie}[Testy mierzalności na półprostych]
\label{thm:extended-tests}
Niech $(X,\mathcal F)$ będzie przestrzenią mierzalną,
a $f:X\to\overline{\mathbb R}$. Następujące warunki są równoważne:
\begin{enumerate}[label=\textup{(\roman*)}]
\item $f$ jest mierzalna;
\item $\{f<a\}\in\mathcal F$ dla każdego $a\in\mathbb R$;
\item $\{f\le a\}\in\mathcal F$ dla każdego $a\in\mathbb R$;
\item $\{f>a\}\in\mathcal F$ dla każdego $a\in\mathbb R$;
\item $\{f\ge a\}\in\mathcal F$ dla każdego $a\in\mathbb R$.
\end{enumerate}
Wystarczy w każdym z warunków (ii)--(v) badać liczby wymierne.
\end{twierdzenie}
\begin{proof}
Zbiory postaci $[-\infty,a)$, $[-\infty,a]$,
$(a,\infty]$ oraz $[a,\infty]$ dla $a\in\mathbb Q$
generują tę samą $\sigma$-algebrę
$\mathcal B(\overline{\mathbb R})$. Istotnie, każda
z tych rodzin jest rozdzielona od pozostałych za pomocą
dopełnień i przeliczalnych granic. Na przykład
$[-\infty,a]=\bigcap_{n\ge1}[-\infty,a+1/n)$ oraz
$(a,\infty]=\overline{\mathbb R}\setminus[-\infty,a]$.
Każdy otwarty przedział w rozszerzonej prostej
jest przecięciem dwóch odpowiednich półprostych,
a każdy zbiór otwarty jest przeliczalną sumą takich
przedziałów, ponieważ $\overline{\mathbb R}$ ma bazę
przeliczalną. Teraz stosujemy twierdzenie
\ref{thm:measurable-generator}.
\end{proof}

\begin{uwaga}[Wartości nieskończone]
Dla mierzalnej funkcji $f:X\to\overline{\mathbb R}$
zbiory $\{f=+\infty\}=\bigcap_{n\ge1}\{f>n\}$ i
$\{f=-\infty\}=\bigcap_{n\ge1}\{f<-n\}$ są mierzalne.
Przy działaniach arytmetycznych musimy unikać
niezdefiniowanych wyrażeń $+\infty-\infty$,
$\infty/\infty$ lub $0\cdot\infty$, chyba że wcześniej
jawnie określimy przyjętą konwencję.
\end{uwaga}

\begin{przyklad}[Funkcja charakterystyczna zbioru]
\label{ex:indicator-measurable}
Dla zbioru $A\subseteq X$ definiujemy
$\chi_A(x)=1$ dla $x\in A$ oraz
$\chi_A(x)=0$ dla $x\notin A$.
Funkcja $\chi_A:X\to\mathbb R$ jest mierzalna
wtedy i tylko wtedy, gdy $A\in\mathcal F$.
Kierunek konieczny otrzymujemy z
$A=\chi_A^{-1}((1/2,\infty))$, a dostateczność
z faktu, że każdy przeciwobraz zbioru borelowskiego
przez $\chi_A$ jest jednym z czterech zbiorów:
$\varnothing$, $A$, $X\setminus A$, $X$.
\end{przyklad}

\section{Działania na funkcjach mierzalnych}

\begin{twierdzenie}[Działania algebraiczne]
\label{thm:measurable-algebra}
Niech $(X,\mathcal F)$ będzie przestrzenią mierzalną,
a $f,g:X\to\mathbb R$ funkcjami mierzalnymi.
Wtedy funkcje $f+g$, $f-g$, $fg$, $|f|$,
$\max(f,g)$ i $\min(f,g)$ są mierzalne.
Zbiór $\{g\ne0\}$ należy do $\mathcal F$, a
funkcja $f/g:\{g\ne0\}\to\mathbb R$ jest mierzalna
względem śladu $\sigma$-algebry. Jeżeli $g(x)\ne0$
dla każdego $x\in X$, to $f/g$ jest mierzalna na całym $X$.
\end{twierdzenie}
\begin{proof}
Odwzorowanie $(f,g):X\to\mathbb R^2$ jest mierzalne
na mocy wniosku \ref{cor:coordinate-measurable}.
Dodawanie $S(u,v)=u+v$, odejmowanie $D(u,v)=u-v$,
mnożenie $M(u,v)=uv$, maksimum i minimum są ciągłymi
funkcjami na $\mathbb R^2$. Moduł jest funkcją ciągłą
na $\mathbb R$. Ich złożenia z funkcją $(f,g)$ są
więc mierzalne, na mocy wniosku \ref{thm:composition-measurable}
i twierdzenia \ref{thm:composition-measurable}.
Zbiór $\{g\ne0\}=g^{-1}(\mathbb R\setminus\{0\})$
jest mierzalny; ograniczenie $(f,g)$ do niego jest
mierzalne. Funkcja $(u,v)\mapsto u/v$ jest ciągła
na otwartym zbiorze $\{(u,v):v\ne0\}$, co kończy dowód.
\end{proof}

\begin{twierdzenie}[Operacje na funkcjach rozszerzonych]
\label{thm:extended-operations}
Niech $f,g:X\to\overline{\mathbb R}$ będą mierzalne.
Wtedy $-f$, $|f|$, $\max(f,g)$ oraz $\min(f,g)$
są mierzalne. Jeśli $f,g$ są nieujemne, ich suma
$f+g$ z naturalną regułą $a+\infty=\infty$
jest mierzalna. Jeżeli nie zachodzi w żadnym punkcie
jednocześnie $f=+\infty$, $g=-\infty$ ani
$f=-\infty$, $g=+\infty$, to również $f+g$ jest mierzalna.
\end{twierdzenie}
\begin{proof}
Negacja i moduł są ciągłe w topologii rozszerzonej
prostej. Funkcje maksimum i minimum są ciągłe na
$\overline{\mathbb R}^2$ z topologią produktową:
na przykład $\max(u,v)<a$ jest równoważne
$u<a$ i $v<a$, zaś $\max(u,v)>a$ zachodzi,
gdy $u>a$ lub $v>a$. Te testy wystarczają do
mierzalności z twierdzenia \ref{thm:extended-tests}.
Dla nieujemnych $f,g$ i dowolnego $a\in\mathbb R$
zachodzi
\[
\{f+g>a\}=\bigcup_{q\in\mathbb Q}
\bigl(\{f>q\}\cap\{g>a-q\}\bigr).
\]
Jeżeli $f(x)+g(x)>a$, można wybrać wymierne $q$
ściśle pomiędzy $a-g(x)$ i $f(x)$; gdy którakolwiek
wartość wynosi $+\infty$, odpowiedni przedział
wymiernych nadal jest niepusty. Implikacja odwrotna
wynika przez dodanie nierówności. Prawa strona
jest mierzalna, a test na półprostych kończy dowód.

W ogólnym przypadku rozdzielamy dziedzinę na mierzalne
zbiory
\[
\begin{aligned}
D_0&=\{|f|<\infty\}\cap\{|g|<\infty\},\\
D_+&=\{f=+\infty,\ g>-\infty\}
      \cup\{g=+\infty,\ f>-\infty\},\\
D_-&=\{f=-\infty,\ g<+\infty\}
      \cup\{g=-\infty,\ f<+\infty\}.
\end{aligned}
\]
W zapisach $|f|<\infty$ i $|g|<\infty$ chodzi o
skończoność wartości; dla punktów z $D_+$ lub $D_-$
suma jest odpowiednio równa $+\infty$ lub $-\infty$.
Założenie wykluczające przeciwne nieskończoności
sprawia, że $D_0,D_+,D_-$ tworzą rozłączny podział
$X$. Na $D_0$ sumę otrzymujemy przez złożenie
mierzalnego odwzorowania $(f,g)|_{D_0}$ z ciągłym
dodawaniem na $\mathbb R^2$. Na dwóch pozostałych
zbiorach funkcja jest stała. Przeciwobrazy
zbiorów borelowskich są zatem mierzalne, bez
odwoływania się do twierdzenia o granicach,
które dopiero zostanie udowodnione.
\end{proof}

\begin{uwaga}[Niezdefiniowane działania]
Dla $f(x)=+\infty$ i $g(x)=-\infty$ formalny zapis
$f(x)+g(x)$ nie określa liczby w $\overline{\mathbb R}$.
Nie można więc pytać o mierzalność takiej \emph{funkcji},
zanim nie ustali się jej wartości na zbiorze punktów
problematycznych. Ustalenie wartości stałej na tym
mierzalnym zbiorze daje jedną z możliwych mierzalnych
wersji rozszerzenia działania.
\end{uwaga}

\section{Granice punktowe funkcji mierzalnych}

\begin{twierdzenie}[Kresy i granice]
\label{thm:pointwise-limits}
Niech $(X,\mathcal F)$ będzie przestrzenią mierzalną,
a $(f_n)_{n\ge1}$ ciągiem mierzalnych funkcji
$f_n:X\to\overline{\mathbb R}$. Wówczas funkcje
\[
\sup_{n\ge1}f_n,\qquad \inf_{n\ge1}f_n,\qquad
\limsup_{n\to\infty}f_n,\qquad
\liminf_{n\to\infty}f_n
\]
są mierzalne. Zbiór
$D=\{x:\lim_{n\to\infty}f_n(x)\text{ istnieje w }
\overline{\mathbb R}\}$ należy do $\mathcal F$;
granica $f(x)=\lim_nf_n(x)$ określona na $D$ jest
mierzalna względem $\mathcal F|_D$.
\end{twierdzenie}
\begin{proof}
Dla dowolnego $a\in\mathbb R$ mamy
$\{\sup_nf_n\le a\}=\bigcap_n\{f_n\le a\}\in\mathcal F$,
a także $\{\inf_nf_n<a\}=\bigcup_n\{f_n<a\}$.
To dowodzi mierzalności supremum i infimum.
Ponieważ
$\limsup_nf_n=\inf_{m\ge1}\sup_{n\ge m}f_n$ i
$\liminf_nf_n=\sup_{m\ge1}\inf_{n\ge m}f_n$,
obie funkcje są mierzalne.
W rozszerzonej prostej granica istnieje wtedy
i tylko wtedy, gdy $\liminf_nf_n=\limsup_nf_n$.
Zbiór równości dwóch mierzalnych funkcji rozszerzonych
jest mierzalny: można na przykład zastosować
ciągłe, ściśle rosnące przekształcenie
$\theta:\overline{\mathbb R}\to[-1,1]$,
$\theta(t)=2\arctan(t)/\pi$ dla $t\in\mathbb R$,
$\theta(\pm\infty)=\pm1$, i następnie badać zbiór,
na którym $|\theta(\limsup_nf_n)-
\theta(\liminf_nf_n)|=0$.
Na $D$ granica jest ograniczeniem mierzalnej funkcji
$\limsup_nf_n$, więc jest mierzalna względem śladu.
\end{proof}

\begin{wniosek}[Monotoniczne ciągi funkcji]\label{idx:06:wniosek:06}

Jeżeli $(f_n)$ jest niemalejącym ciągiem mierzalnych
funkcji $f_n:X\to[0,\infty]$, to $f=\sup_nf_n$
jest mierzalna i $f_n(x)\uparrow f(x)$ dla każdego $x$.
Analogicznie, granica malejącego ciągu funkcji
$\overline{\mathbb R}$-wartościowych jest mierzalna.
\end{wniosek}

\begin{przyklad}[Granica punktowa funkcji ciągłych]
Funkcje $f_n(x)=x^n$ na $[0,1]$ są ciągłe,
a więc borelowsko mierzalne. Zbiegają punktowo do
$\chi_{\{1\}}$, która jest mierzalna, choć
nie jest ciągła w punkcie $1$ (dla topologii podprzestrzeni).
Granice punktowe zachowują mierzalność, lecz nie muszą
zachowywać ciągłości.
\end{przyklad}

\section{Funkcje proste}

\begin{definicja}[Funkcja prosta]
\label{def:simple-function}
Funkcję $s:X\to\mathbb R$ nazywamy
\emph{funkcją prostą mierzalną}, jeśli przyjmuje
skończenie wiele wartości i jest mierzalna.
Równoważnie istnieją $m\in\mathbb N$,
liczby $a_1,\ldots,a_m\in\mathbb R$ i zbiory
$A_1,\ldots,A_m\in\mathcal F$ tworzące podział $X$
takie, że
\[
s=\sum_{j=1}^m a_j\chi_{A_j}.
\]
Jeśli wszystkie $a_j\ge0$, funkcję prostą nazywamy
\emph{nieujemną}.
\end{definicja}

\begin{proof}[Dowód równoważności w definicji]
Jeżeli funkcja $s$ przyjmuje skończenie wiele różnych
wartości $a_1,\ldots,a_m$, to jej włókna
$A_j=\{x:s(x)=a_j\}$ są mierzalne i tworzą
podział $X$. Na $A_j$ funkcja ma wartość $a_j$,
co daje podaną reprezentację. Odwrotnie, podana suma
jest mierzalna jako suma skończona mierzalnych funkcji
charakterystycznych i ma zbiór wartości zawarty w
$\{a_1,\ldots,a_m\}$.
\end{proof}

\begin{uwaga}[Po co wymagamy podziału?]
Reprezentacja jako skończona kombinacja funkcji
charakterystycznych nie musi być jednoznaczna, jeśli
zbiory się nakładają. Zawsze można ją zastąpić
reprezentacją kanoniczną, w której $A_j$ są rozłączne
i odpowiadają dokładnie różnym wartościom $s$.
Wówczas wzór na całkę funkcji prostej będzie
niezależny od sposobu zapisania tej funkcji.
\end{uwaga}

\begin{twierdzenie}[Monotoniczna aproksymacja funkcji nieujemnej]
\label{thm:simple-up}
Niech $(X,\mathcal F)$ będzie przestrzenią mierzalną,
a $f:X\to[0,\infty]$ funkcją mierzalną. Dla
$n\ge1$ określmy
\begin{equation}\label{eq:canonical-simple}
s_n(x)=2^{-n}\Big\lfloor 2^n\min\{f(x),n\}\Big\rfloor,
\qquad x\in X.
\end{equation}
Wtedy każda $s_n$ jest nieujemną funkcją prostą mierzalną,
$0\le s_n\le s_{n+1}\le f$ oraz
$s_n(x)\to f(x)$ dla każdego $x\in X$,
przy czym granica może wynosić $+\infty$.
Dla $f(x)<\infty$ i $n>f(x)$ zachodzi
$0\le f(x)-s_n(x)<2^{-n}$.
\end{twierdzenie}
\begin{proof}
Wartość $\min(f,n)$ należy do $[0,n]$.
Po pomnożeniu przez $2^n$ i zastosowaniu części
całkowitej otrzymujemy jedną z liczb całkowitych
$0,1,\ldots,n2^n$. Funkcja $s_n$ przyjmuje zatem
skończenie wiele wartości. Dla $k=0,\ldots,n2^n-1$
zbiór $\{s_n=k2^{-n}\}$ jest opisany nierównościami
$k2^{-n}\le\min(f,n)<(k+1)2^{-n}$;
zbiór odpowiadający ostatniej wartości jest
$\{f\ge n\}$. Wszystkie te zbiory są mierzalne,
więc $s_n$ jest funkcją prostą mierzalną.

Wprowadźmy $u_n(x)=\min(f(x),n)$. Mamy
$u_n\le u_{n+1}$; funkcja
$t\mapsto 2^{-n}\lfloor2^nt\rfloor$ jest niemalejąca,
a ponadto dla każdego $t\ge0$
$2^{-(n+1)}\lfloor2^{n+1}t\rfloor\ge
2^{-n}\lfloor2^nt\rfloor$.
Ostatnia nierówność wynika z tego, że
$\lfloor2z\rfloor\ge2\lfloor z\rfloor$.
Wynika stąd $s_n\le s_{n+1}$. Nierówności
$0\le s_n\le u_n\le f$ wynikają z własności
części całkowitej. Gdy $f(x)$ jest skończona,
przy $n>f(x)$ mamy $u_n(x)=f(x)$ i
$0\le f(x)-s_n(x)<2^{-n}$.
Jeżeli $f(x)=\infty$, to $u_n(x)=n$ i
$s_n(x)=n\to\infty$. To kończy dowód.
\end{proof}

\begin{przyklad}[Aproksymacja z podziałem na poziomy]
Niech $X=[0,1]$, $f(x)=x$ i $n=2$.
Wtedy $s_2(x)=\frac14\lfloor4x\rfloor$ ma wartości
$0,1/4,1/2,3/4,1$ na kolejnych przedziałach
$[0,1/4)$, $[1/4,1/2)$, $[1/2,3/4)$,
$[3/4,1)$ i singletonie $\{1\}$.
Dalsze przybliżenia są coraz dokładniejszymi
schodkami leżącymi poniżej wykresu $f$.
\end{przyklad}

\begin{center}
\begin{tikzpicture}[x=9cm,y=4.0cm]
  \draw[->] (-0.015,0) -- (1.08,0) node[right] {$x$};
  \draw[->] (0,-0.025) -- (0,1.1) node[above] {$y$};
  \draw[blue!70!black,thick] (0,0)--(1,1) node[above left] {$f(x)=x$};
  \draw[red!75!black,thick] (0,0)--(.25,0) (.25,.25)--(.5,.25) (.5,.5)--(.75,.5) (.75,.75)--(1,.75);
  \foreach \x/\y in {.25/0,.5/.25,.75/.5,1/.75} {
    \draw[red!75!black,fill=white] (\x,\y) circle[radius=1.2pt];
  }
  \foreach \x/\y in {.25/.25,.5/.5,.75/.75,1/1} {
    \fill[red!75!black] (\x,\y) circle[radius=1.2pt];
  }
  \draw[dashed,gray!60] (1,0)--(1,1);
  \node[anchor=north] at (0,0) {$0$};
  \node[anchor=north] at (1,0) {$1$};
  \node[red!75!black,anchor=west] at (.42,.13) {$s_2$};
\end{tikzpicture}

\smallskip
\emph{Aproksymacja funkcji $f(x)=x$ na $[0,1]$ przez $s_2$.
Wartości w punktach skoku definiuje wzór \eqref{eq:canonical-simple}.}
\end{center}

\begin{twierdzenie}[Aproksymacja funkcji o wartościach rzeczywistych]
\label{thm:simple-real-approx}
Niech $(X,\mathcal F)$ będzie przestrzenią mierzalną,
a $f:X\to\mathbb R$ mierzalną funkcją skończoną.
Wówczas istnieje ciąg funkcji prostych mierzalnych
$(r_n)$, dla którego $r_n(x)\to f(x)$ dla każdego $x\in X$.
Jeśli $|f(x)|\le M$ dla wszystkich $x\in X$ i pewnego
$M<\infty$, można wybrać $r_n$ tak, by
$\|r_n-f\|_\infty\le2^{-n}$, gdzie
$\|h\|_\infty:=\sup_{x\in X}|h(x)|$.
\end{twierdzenie}
\begin{proof}
Niech $f^+=\max(f,0)$ i $f^-=\max(-f,0)$.
Obie funkcje są mierzalne i nieujemne.
Na mocy twierdzenia \ref{thm:simple-up} istnieją
niemalejące ciągi nieujemnych funkcji prostych
$s_n^+\uparrow f^+$ i $s_n^-\uparrow f^-$.
Funkcje $r_n=s_n^+-s_n^-$ są proste i mierzalne:
przyjmują skończenie wiele wartości i są różnicami
funkcji mierzalnych skończonych. Ponadto
$r_n(x)\to f^+(x)-f^-(x)=f(x)$ dla każdego $x$.
Jeśli $f$ jest ograniczona przez $M$, dla
$n>M$ oba błędy aproksymacji są mniejsze niż
$2^{-n}$, a ponieważ w każdym punkcie co najwyżej
jedna z liczb $f^+(x),f^-(x)$ jest dodatnia,
$|r_n(x)-f(x)|<2^{-n}$. Dla początkowych indeksów
można wybrać dowolne funkcje proste, gdyż nie
wpływa to na istnienie późniejszych przybliżeń;
aby uzyskać podaną nierówność dla wszystkich $n$,
bezpośrednio stosujemy obcięcie i kwantyzację
$r_n(x)=2^{-n}\lfloor2^nf(x)\rfloor$.
Przy ograniczonym $f$ funkcja ta jest prosta i
mierzalna, a dla każdego $x$ zachodzi
$0\le f(x)-r_n(x)<2^{-n}$. Dowodzi to ostatniej tezy.
\end{proof}

\begin{wniosek}[Charakteryzacja mierzalności przez funkcje proste]
\label{cor:meas-simple-characterization}
Niech $(X,\mathcal F)$ będzie przestrzenią mierzalną,
a $f:X\to\mathbb R$ funkcją. Wówczas $f$ jest
mierzalna wtedy i tylko wtedy, gdy istnieje
ciąg mierzalnych funkcji prostych $r_n:X\to\mathbb R$
zbieżny punktowo do $f$.
\end{wniosek}
\begin{proof}
Konieczność daje twierdzenie \ref{thm:simple-real-approx}.
Dostateczność wynika z twierdzenia
\ref{thm:pointwise-limits}, ponieważ punktowa granica
mierzalnych funkcji rzeczywistych jest mierzalna.
\end{proof}

\section{Mierzalność funkcji zespolonych i wektorowych}

\begin{definicja}[Funkcja zespolona mierzalna]\label{idx:06:definicja:11}

Funkcję $f:X\to\mathbb C$ nazywamy mierzalną,
jeśli jest $(\mathcal F,\mathcal B(\mathbb C))$-mierzalna,
gdzie $\mathbb C$ utożsamiamy z $\mathbb R^2$
przez $z\mapsto(\operatorname{Re}z,\operatorname{Im}z)$.
\end{definicja}

\begin{twierdzenie}[Test na częściach rzeczywistej i urojonej]
\label{thm:complex-test}
Niech $(X,\mathcal F)$ będzie przestrzenią mierzalną,
a $f:X\to\mathbb C$. Funkcja $f$ jest mierzalna
wtedy i tylko wtedy, gdy funkcje
$\operatorname{Re}f$ i $\operatorname{Im}f$ są
mierzalne jako funkcje rzeczywiste. W takim przypadku
$\overline f$ i $|f|$ także są mierzalne.
\end{twierdzenie}
\begin{proof}
Pierwsza część jest szczególnym przypadkiem
wniosku \ref{cor:coordinate-measurable} dla $d=2$.
Sprzężenie zespolone i moduł są ciągłe na $\mathbb C$,
a złożenie z funkcją mierzalną pozostaje mierzalne.
\end{proof}

\begin{przyklad}[Funkcje parametryczne]
Jeśli $u,v:X\to\mathbb R$ są mierzalne, to
$f(x)=u(x)+iv(x)$ jest mierzalna.
Nie potrzebujemy do tego żadnej dodatkowej miary
na $X$. Gdy zaś $u,v$ są funkcjami mierzalnymi
względem $\mathcal L_d$, ale niekoniecznie
borelowskimi, również $f$ jest mierzalna względem
$\mathcal L_d$.
\end{przyklad}

\section{Własności prawie wszędzie i uzupełnianie miar}

\begin{definicja}[Równość i zbieżność prawie wszędzie]
\label{def:ae-functions}
Niech $(X,\mathcal F,\mu)$ będzie przestrzenią z miarą.
Funkcje mierzalne $f,g:X\to\mathbb R$ są równe
\emph{$\mu$-prawie wszędzie}, jeśli istnieje
$N\in\mathcal F$ o $\mu(N)=0$ takie, że
$f(x)=g(x)$ dla każdego $x\in X\setminus N$.
Ciąg funkcji mierzalnych $f_n:X\to\mathbb R$
\emph{zbiega $\mu$-prawie wszędzie} do $f:X\to\mathbb R$,
co zapisujemy $f_n\to f$ $\mu$-p.w., jeśli istnieje
taki zbiór $N$, że $f_n(x)\to f(x)$ dla każdego
$x\in X\setminus N$.
\end{definicja}

\begin{uwaga}[Ostrożnie z niezupełną miarą]
Jeśli $\mu$ nie jest zupełna, podzbiór zbioru
mierzalnego miary zero nie musi być mierzalny.
Dlatego definicję wyraziliśmy przez istnienie
\emph{mierzalnego} zbioru zerowego zawierającego
zbiór punktów wyjątkowych. Dla funkcji mierzalnych
sam zbiór $\{f\ne g\}$ jest mierzalny, więc można
równoważnie wymagać $\mu(\{f\ne g\})=0$.
Podobnie dla mierzalnych $f_n,f$ zbiór braku
zbieżności jest mierzalny na mocy twierdzenia
\ref{thm:pointwise-limits}.
\end{uwaga}

\begin{twierdzenie}[Modyfikacje na zbiorach zerowych]
\label{thm:null-modifications}
Niech $(X,\mathcal F,\mu)$ będzie \emph{zupełną}
przestrzenią z miarą. Jeżeli $f:X\to\mathbb R$
jest mierzalna, $g:X\to\mathbb R$ dowolna,
a $\{x:f(x)\ne g(x)\}$ jest zawarty w pewnym
$N\in\mathcal F$ o $\mu(N)=0$, to $g$ jest mierzalna.
Bez założenia zupełności wniosek może być fałszywy.
\end{twierdzenie}
\begin{proof}
W przestrzeni zupełnej każdy podzbiór $N$ należy
do $\mathcal F$. Dla zbioru borelowskiego
$B\subseteq\mathbb R$ zbiory
$g^{-1}(B)\cap(X\setminus N)=f^{-1}(B)\cap(X\setminus N)$
i $g^{-1}(B)\cap N\subseteq N$ są mierzalne.
Ich suma jest równa $g^{-1}(B)$, zatem $g$ jest mierzalna.
Dla kontrprzykładu bierzemy nieuzupełnioną borelowską
miarę Lebesgue'a na $\mathbb R$ oraz podzbiór
$A$ zbioru Cantora, który nie jest borelowski
(wniosek \ref{cor:nonborel-lebesgue}).
Funkcja $f\equiv0$ jest mierzalna, ale
$g=\chi_A$ nie jest borelowsko mierzalna.
Jednocześnie $f=g$ poza borelowskim zbiorem
Cantora, którego miara wynosi zero.
\end{proof}

\begin{wniosek}[Wersje mierzalne funkcji określonych p.w.]\label{idx:06:wniosek:15}

Niech $(X,\mathcal F,\mu)$ będzie przestrzenią z miarą,
a $f_n:X\to\mathbb R$ mierzalnymi funkcjami.
Jeżeli ciąg $(f_n(x))$ jest zbieżny dla
$\mu$-prawie wszystkich $x$, to istnieje
mierzalna funkcja $f:X\to\mathbb R$, taka że
$f_n\to f$ $\mu$-p.w.
\end{wniosek}
\begin{proof}
Z twierdzenia \ref{thm:pointwise-limits} zbiór $D$
punktów, w których granica jest rzeczywista i
skończona, jest mierzalny, a granica na nim jest mierzalna.
Z założenia $X\setminus D$ ma miarę zero.
Definiujemy $f(x)=\lim_nf_n(x)$ na $D$ oraz
$f(x)=0$ na $X\setminus D$. Funkcja jest mierzalna,
gdyż dla borelowskiego $B$ jej przeciwobraz jest
sumą mierzalnego przeciwobrazu na $D$ i,
ewentualnie, całego $X\setminus D$.
\end{proof}

\begin{uwaga}[Zbieżność prawie wszędzie nie zapewnia mierzalności dowolnej wersji granicy]
\label{rem:noncomplete-ae-limit-6}
Jeśli nie wymagamy mierzalności funkcji granicznej,
z założenia $f_n\to f$ poza zbiorem zerowym nie można
wnioskować, że $f$ jest mierzalna. Na przestrzeni
$(\mathbb R,\mathcal B(\mathbb R),\lambda_1|_{\mathcal B})$
wybierzmy nieborelowski podzbiór $A$ zbioru Cantora.
Funkcje $f_n\equiv0$ są borelowskie i zbiegają do
$f=\chi_A$ poza borelowskim zbiorem Cantora miary zero,
lecz $f$ nie jest mierzalna względem $\mathcal B(\mathbb R)$.
Wniosek o mierzalnej wersji granicy gwarantuje istnienie
\emph{pewnego} mierzalnego przedstawiciela; nie oznacza,
że każdy wybór wartości na zbiorze wyjątkowym jest
mierzalny. W przestrzeni zupełnej takie modyfikacje są
mierzalne na mocy twierdzenia~\ref{thm:null-modifications}.
\end{uwaga}

\section{Zbieżność według miary: narzędzia potrzebne do całkowania}
Minimalna definicja i twierdzenie o podciągu pojawiają się tutaj,
ponieważ wykorzystamy je już w rozdziale 8; systematyczne porównanie
wszystkich rodzajów zbieżności pozostawiamy rozdziałowi 9.


\begin{definicja}[Zbieżność według miary]
\label{def:convergence-measure}
Niech $(X,\mathcal F,\mu)$ będzie przestrzenią z miarą,
a $f_n,f:X\to\mathbb R$ mierzalnymi funkcjami.
Mówimy, że $f_n$ \emph{zbiega do $f$ według miary $\mu$},
i piszemy $f_n\xrightarrow{\mu}f$, jeśli
\[
\forall\varepsilon>0\qquad
\mu\bigl(\{x\in X:|f_n(x)-f(x)|>\varepsilon\}\bigr)
\xrightarrow[n\to\infty]{}0.
\]
Definicja nie wymaga $\mu(X)<\infty$ ani
$\sigma$-skończoności.
\end{definicja}

\begin{twierdzenie}[Prawie wszędzie implikuje zbieżność według miary przy mierze skończonej]
\label{thm:ae-to-measure-finite}
Jeśli $(X,\mathcal F,\mu)$ jest przestrzenią
z $\mu(X)<\infty$, a $f_n,f:X\to\mathbb R$ są
mierzalne i $f_n\to f$ $\mu$-prawie wszędzie,
to $f_n\xrightarrow{\mu}f$.
\end{twierdzenie}
\begin{proof}
Ustalmy $\varepsilon>0$ i niech
$E_n=\bigcup_{k\ge n}\{|f_k-f|>\varepsilon\}$.
Wtedy $E_n\in\mathcal F$, $E_{n+1}\subseteq E_n$,
a $\bigcap_nE_n$ składa się z punktów, gdzie
nierówność $|f_k-f|>\varepsilon$ zachodzi
nieskończenie wiele razy. Zbieżność prawie wszędzie
gwarantuje $\mu(\bigcap_nE_n)=0$.
Ponieważ $\mu(E_1)\le\mu(X)<\infty$,
ciągłość miary z góry (twierdzenie
\ref{thm:measure-continuity}) daje $\mu(E_n)\to0$.
Wreszcie $\{|f_n-f|>\varepsilon\}\subseteq E_n$,
więc z monotoniczności $\mu(\{|f_n-f|>\varepsilon\})\to0$.
\end{proof}

\begin{przyklad}[Dlaczego skończoność miary jest ważna]
\label{ex:aenotmeasure}
Na $(\mathbb R,\mathcal L_1,\lambda_1)$ niech
$f_n=\chi_{[n,n+1]}$. Dla każdego ustalonego
$x\in\mathbb R$ mamy $f_n(x)=0$ dla wszystkich
dostatecznie dużych $n$, a więc $f_n\to0$ wszędzie.
Jednak dla $\varepsilon=1/2$ zachodzi
$\lambda_1(\{|f_n|>1/2\})=1$ dla wszystkich $n$.
Nie ma więc zbieżności według miary na całej $\mathbb R$.
\end{przyklad}

\begin{twierdzenie}[Podciąg zbieżny prawie wszędzie]
\label{thm:measure-subsequence}
Niech $(X,\mathcal F,\mu)$ będzie dowolną przestrzenią
z miarą, a $f_n,f:X\to\mathbb R$ mierzalnymi funkcjami.
Jeśli $f_n\xrightarrow{\mu}f$, to istnieje ściśle
rosnący ciąg indeksów $(n_k)_{k\ge1}$ taki, że
$f_{n_k}\to f$ $\mu$-prawie wszędzie.
\end{twierdzenie}
\begin{proof}
Z definicji zbieżności według miary wybieramy
indukcyjnie $n_k>n_{k-1}$ tak duże, aby
\[
\mu(A_k)<2^{-k},\qquad
A_k=\{|f_{n_k}-f|>2^{-k}\}.
\]
Wtedy $\sum_{k\ge1}\mu(A_k)\le1<\infty$.
Pierwsza część lematu Borela--Cantellego
(twierdzenie \ref{thm:borel-cantelli-1})
daje $\mu(\limsup_k A_k)=0$.
Poza tym zbiorem punkt $x$ należy do zaledwie
skończenie wielu $A_k$; dla dostatecznie dużych $k$
mamy więc $|f_{n_k}(x)-f(x)|\le2^{-k}$.
Prawa strona zbiega do zera, a zatem
$f_{n_k}(x)\to f(x)$.
\end{proof}

\begin{uwaga}[Dlaczego nie potrzebujemy skończoności miary?]
Wybór podciągu opiera się wyłącznie na sumowalności
$\sum_k\mu(A_k)$, a lemat Borela--Cantellego nie
zakłada $\mu(X)<\infty$. Jest to zasadniczo inny
argument niż ciągłość miary z góry zastosowana w
twierdzeniu~\ref{thm:ae-to-measure-finite}.
\end{uwaga}

\begin{wniosek}[Jednoznaczność granicy według miary]
\label{cor:measure-limit-unique}
Jeśli $f_n\xrightarrow{\mu}f$ i
$f_n\xrightarrow{\mu}g$ dla mierzalnych funkcji
$f,g:X\to\mathbb R$, to $f=g$ $\mu$-prawie wszędzie.
\end{wniosek}
\begin{proof}
Dla $\varepsilon>0$ z nierówności trójkąta
$\{|f-g|>\varepsilon\}\subseteq
\{|f-f_n|>\varepsilon/2\}\cup
\{|f_n-g|>\varepsilon/2\}$.
Z podaddytywności i przejścia do granicy
otrzymujemy $\mu(\{|f-g|>\varepsilon\})=0$.
Ponieważ $\{f\ne g\}=\bigcup_{m\ge1}\{|f-g|>1/m\}$,
zbiór ten ma miarę zero.
\end{proof}

\begin{przyklad}[Zbieżność według miary bez zbieżności p.w.]
\label{ex:typewriter}
Na $X=[0,1)$ z miarą Lebesgue'a, dla każdego $k\ge0$
i $j=0,\ldots,2^k-1$, kładziemy
$f_{2^k+j}=\chi_{[j2^{-k},(j+1)2^{-k})}$.
Każdy indeks $n\ge1$ występuje w tym wzorze dokładnie
raz (blok indeksów $2^k,\ldots,2^{k+1}-1$).
Dla $n=2^k+j$ mamy
$\lambda_1(\{|f_n|>1/2\})=2^{-k}\to0$,
więc $f_n\to0$ według miary.
Z drugiej strony dla każdego $x\in[0,1)$ i każdego
$k$ dokładnie jeden z przedziałów poziomu $k$
zawiera $x$. Stąd w każdym bloku istnieje jeden indeks
$n$ z $f_n(x)=1$ oraz (dla $k\ge1$) inny indeks
z $f_n(x)=0$. Ciąg $f_n(x)$ nie ma granicy w żadnym
punkcie $[0,1)$.
\end{przyklad}

\begin{twierdzenie}[Lokalna zbieżność według miary]
\label{thm:local-convergence-measure}
Niech $(X,\mathcal F,\mu)$ będzie przestrzenią z miarą,
a $f_n,f:X\to\mathbb R$ funkcjami mierzalnymi,
przy czym $f_n\to f$ $\mu$-prawie wszędzie.
Wówczas dla każdego $A\in\mathcal F$ z $\mu(A)<\infty$
i każdego $\varepsilon>0$ zachodzi
\[
\mu\bigl(A\cap\{|f_n-f|>\varepsilon\}\bigr)\longrightarrow0.
\]
\end{twierdzenie}
\begin{proof}
Rozpatrujemy przestrzeń
$(A,\mathcal F|_A,\mu|_A)$ o skończonej mierze,
a następnie stosujemy twierdzenie
\ref{thm:ae-to-measure-finite} do funkcji ograniczonych
do $A$.
\end{proof}

\begin{uwaga}[Co jest, a co nie jest implikacją?]
Zbieżność prawie wszędzie implikuje zbieżność
według miary przy $\mu(X)<\infty$, lecz nie musi
implikować jej dla $\mu(X)=\infty$.
Zbieżność według miary zawsze pozwala wybrać podciąg
zbieżny prawie wszędzie, ale nie wymusza zbieżności
prawie wszędzie całego ciągu. Kontrprzykłady powyżej
pokazują niezależność obu problemów. Twierdzenia
Jegorowa i Riesza o podciągach oraz pełne zestawienie
rodzajów zbieżności rozwinie rozdział 9.
\end{uwaga}

\section{Przygotowanie do definicji całki Lebesgue'a}

Skończone kombinacje funkcji charakterystycznych są
pomostem pomiędzy mierzeniem zbiorów a całkowaniem.
Dla nieujemnej funkcji prostej $s=\sum_{j=1}^m
a_j\chi_{A_j}$, gdzie $A_j$ tworzą mierzalny
podział $X$, naturalny kandydat na całkę to
\[
\sum_{j=1}^m a_j\mu(A_j).
\]
Jeśli $a_j=0$ i $\mu(A_j)=\infty$, przyjmujemy
w tym nieujemnym wyrażeniu konwencję $0\cdot\infty=0$.
W następnym rozdziale dowiedziemy, że liczba ta
nie zależy od reprezentacji $s$, po czym określimy
całkę dowolnej funkcji nieujemnej jako supremum
po wszystkich prostych funkcjach mierzalnych
mniejszych od niej.

W ten sposób twierdzenie \ref{thm:simple-up}
ma fundamentalną rolę: zapewnia, że zbiór takich
przybliżeń jest dostatecznie bogaty, aby odtworzyć
wartości dowolnej funkcji mierzalnej.

% END chapters/pl/measure_theory/06-rozdzial.tex


% BEGIN chapters/pl/measure_theory/07-rozdzial.tex
% ROZDZIAŁ 7 — do włączenia w main.tex
\chapter{Konstrukcja całki Lebesgue'a}
\label{ch:calkowanie-lebesgue}

\section{Od mierzenia zbiorów do mierzenia wartości funkcji}

W klasycznej sumie Riemanna dzielimy \emph{dziedzinę} na małe
przedziały. Konstrukcja Lebesgue'a rozpoczyna się od podziału
\emph{zbioru wartości}: dla funkcji prostej $s$ mierzymy zbiory,
na których $s$ przyjmuje kolejne wartości, i sumujemy odpowiednie
iloczyny. Pozwala to całkować na dowolnych przestrzeniach z miarą,
bez porządku na dziedzinie i bez pojęcia przedziału.

W całym rozdziale $(X,\mathcal F,\mu)$ jest dowolną przestrzenią
z miarą. Nie zakładamy, że $\mu(X)<\infty$, ani nawet
$\sigma$-skończoności. Jeżeli nie zaznaczamy inaczej,
funkcje są mierzalne względem $\mathcal F$.
Oznaczenia $\overline{\mathbb R}=[-\infty,\infty]$,
$f^+=\max(f,0)$ i $f^-=\max(-f,0)$ mają zwykły sens;
$f=f^+-f^-$ w punktach, gdzie $f$ jest skończona.
Dla nieujemnych wielkości z $[0,\infty]$ umawiamy się,
że $0\cdot\infty=0$, również przy wyznaczaniu całki
nieujemnej funkcji prostej. Nie dopuszczamy odejmowania
$\infty-\infty$.

\section{Całka nieujemnej funkcji prostej}

\begin{definicja}[Całka funkcji prostej]
\label{def:integral-simple}
Niech $s:X\to[0,\infty)$ będzie mierzalną funkcją prostą.
Jeżeli $a_1,\ldots,a_m$ są wszystkimi różnymi wartościami $s$,
a $A_j=\{x\in X:s(x)=a_j\}$, określamy
\begin{equation}\label{eq:simple-integral}
\int_X s(x)\,\mathrm d\mu(x)
:=\sum_{j=1}^{m}a_j\mu(A_j)\in[0,\infty].
\end{equation}
Zbiory $A_j$ są mierzalne i tworzą rozłączny podział $X$.
Przy $a_j=0$ składnik jest równy zeru, niezależnie od
$\mu(A_j)$. Wartość całki może być nieskończona.
\end{definicja}

\begin{lemat}[Niezależność od reprezentacji]
\label{lem:simple-refinement}
Jeżeli $s=\sum_{i=1}^p a_i\chi_{A_i}$, gdzie
$a_i\ge0$, a zbiory $A_i\in\mathcal F$ tworzą
rozłączny podział $X$ (dopuszczamy powtarzające się
współczynniki), to
$\int_Xs\,\mathrm d\mu(x)=\sum_{i=1}^pa_i\mu(A_i)$.
Wartość ta nie zależy od wybranego podziału ani od
powtarzania współczynników.
\end{lemat}
\begin{proof}
Niech $a_1',\ldots,a_r'$ będą różnymi wartościami
funkcji $s$, a $C_j=\{s=a_j'\}$. Dla każdego $j$
zbiór $C_j$ jest rozłączną sumą tych $A_i$, dla których
$a_i=a_j'$. Zatem, na mocy skończonej addytywności miary,
$a_j'\mu(C_j)=\sum_{i:a_i=a_j'}a_i\mu(A_i)$.
Gdy $a_j'=0$, obie strony są zerem według przyjętej
konwencji; dla $a_j'>0$ równość wynika z rozdzielności
mnożenia w $[0,\infty]$. Sumując po $j$, otrzymujemy
tezę. Dla dwóch podziałów można również zastosować
wspólne rozdrobnienie $A_i\cap B_j$; dowodzi to
niezależności bez zakładania skończoności miary.
\end{proof}

\begin{twierdzenie}[Rachunek dla funkcji prostych]
\label{thm:simple-properties}
Niech $s,t:X\to[0,\infty)$ będą mierzalnymi funkcjami
prostymi, a $c\in[0,\infty)$. Wtedy:
\begin{enumerate}[label=\textup{(\roman*)}]
\item $\int_X0\,\mathrm d\mu(x)=0$ oraz
$\int_X\chi_A\,\mathrm d\mu(x)=\mu(A)$ dla każdego
$A\in\mathcal F$;
\item jeżeli $s\le t$, to
$\int_Xs\,\mathrm d\mu(x)\le\int_Xt\,\mathrm d\mu(x)$;
\item $\int_X(cs)\,\mathrm d\mu(x)=c\int_Xs\,\mathrm d\mu(x)$;
\item $\int_X(s+t)\,\mathrm d\mu(x)=
\int_Xs\,\mathrm d\mu(x)+\int_Xt\,\mathrm d\mu(x)$.
\end{enumerate}
W punktach (iii) i (iv) wszystkie działania odbywają się
w $[0,\infty]$ z konwencją $0\cdot\infty=0$.
\end{twierdzenie}
\begin{proof}
Pierwsza część wynika wprost z definicji. Niech
$s=\sum_i a_i\chi_{A_i}$ i
$t=\sum_j b_j\chi_{B_j}$ będą reprezentacjami
na rozłącznych podziałach $X$. Zbiory
$D_{ij}=A_i\cap B_j$ tworzą wspólne rozdrobnienie.
Na każdym $D_{ij}$ obie funkcje są stałe, dlatego
\[
\int_Xs\,\mathrm d\mu(x)=\sum_{i,j}a_i\mu(D_{ij}),\qquad
\int_Xt\,\mathrm d\mu(x)=\sum_{i,j}b_j\mu(D_{ij}).
\]
Jeśli $s\le t$, to dla każdego niepustego $D_{ij}$
mamy $a_i\le b_j$, więc każdy składnik pierwszej sumy
jest nie większy od odpowiedniego składnika drugiej.
Skalowanie i addytywność wynikają z równości
$(ca_i)\mu(D_{ij})=c(a_i\mu(D_{ij}))$ oraz
$(a_i+b_j)\mu(D_{ij})=a_i\mu(D_{ij})+b_j\mu(D_{ij})$.
Żadna operacja nie wymaga odejmowania nieskończoności.
\end{proof}

\begin{przyklad}[Mierzalna funkcja prosta na przestrzeni nieskończonej]
Na $(\mathbb R,\mathcal L_1,\lambda_1)$ funkcja
$s=3\chi_{[0,2]}+5\chi_{[4,5]}$ ma całkę
$3\cdot2+5\cdot1=11$. Natomiast
$t=\chi_{[0,\infty)}$ ma całkę $+\infty$.
Obie funkcje są proste i mierzalne, chociaż druga
nie jest całkowalna w sensie skończonej całki.
\end{przyklad}

\section{Definicja całki dowolnej funkcji nieujemnej}

\begin{definicja}[Całka Lebesgue'a funkcji nieujemnej]
\label{def:integral-positive}
Niech $f:X\to[0,\infty]$ będzie mierzalna. Definiujemy
\begin{equation}\label{eq:positive-integral}
\int_X f(x)\,\mathrm d\mu(x)
:=\sup\left\{\int_X s(x)\,\mathrm d\mu(x):
\begin{array}{l}
 s:X\to[0,\infty)\text{ jest mierzalna i prosta},\\[-2pt]
 0\le s(x)\le f(x)\text{ dla wszystkich }x\in X
\end{array}\right\}.
\end{equation}
Supremum należy do $[0,\infty]$. Rodzina funkcji
prostych w nawiasie jest niepusta, ponieważ zawiera zero.
\end{definicja}

\begin{uwaga}[Co dokładnie znaczy ta konstrukcja?]
Mierzalność funkcji $f$ nie wymusza skończoności całki.
Nie wybieramy też żadnego konkretnego przybliżenia:
rozpatrujemy \emph{wszystkie} nieujemne funkcje proste
nieprzekraczające $f$. Twierdzenie
\ref{thm:simple-up} zapewnia, że istnieją również
jawne, punktowo niemalejące przybliżenia $s_n\uparrow f$.
Dopiero twierdzenie Beppo Leviego w rozdziale
\ref{ch:twierdzenia-graniczne} wykaże, że całkę można
obliczać jako granicę całek takiej sekwencji.
\end{uwaga}

\begin{twierdzenie}[Pierwsze własności całki nieujemnej]
\label{thm:nonnegative-first-properties}
Niech $f,g:X\to[0,\infty]$ będą mierzalne,
$c\in[0,\infty)$ i $A\in\mathcal F$. Wtedy:
\begin{enumerate}[label=\textup{(\roman*)}]
\item gdy $f\le g$, zachodzi
$\int_X f\,\mathrm d\mu(x)\le\int_Xg\,\mathrm d\mu(x)$;
\item $\int_X cf\,\mathrm d\mu(x)=c\int_Xf\,\mathrm d\mu(x)$;
\item dla funkcji prostej $f$ definicje
\ref{def:integral-simple} i \ref{def:integral-positive}
są zgodne;
\item $\int_X f\chi_A\,\mathrm d\mu(x)$ zależy
wyłącznie od wartości $f$ na $A$ i dla
$A\subseteq B$, $A,B\in\mathcal F$, zachodzi
$\int_X f\chi_A\,\mathrm d\mu(x)\le
\int_Xf\chi_B\,\mathrm d\mu(x)$;
\item jeśli $\mu(A)=0$, to
$\int_X f\chi_A\,\mathrm d\mu(x)=0$,
nawet gdy $f(x)=+\infty$ w punktach $A$.
\end{enumerate}
\end{twierdzenie}
\begin{proof}
W (i) każda funkcja prosta $s\le f$ spełnia
również $s\le g$, więc supremum dla $f$
nie przekracza supremum dla $g$.
Dla $c>0$ podstawienie $t=cs$ daje bijekcję
między nieujemnymi prostymi minorantami $f$ i $cf$;
na mocy twierdzenia \ref{thm:simple-properties}
całki skalują się przez $c$. Przy $c=0$ obie
strony wynoszą zero według umowy $0\cdot\infty=0$.

Jeżeli $f$ jest prosta, sama należy do rodziny
minorant w \eqref{eq:positive-integral}. Dla każdego
$s\le f$ monotoniczność całki prostej daje
$\int_Xs\,\mathrm d\mu(x)\le\int_Xf\,\mathrm d\mu(x)$;
stąd zgodność obu definicji.
Część (iv) wynika z punktowej równości funkcji
$f\chi_A$ przy ustalonych wartościach na $A$
oraz z (i). Aby wykazać (v), ustalmy dowolną
nieujemną funkcję prostą $s\le f\chi_A$.
Poza $A$ mamy $s=0$; skoro $\mu(A)=0$,
każdy jej dodatni poziom ma miarę zero, więc
$\int_Xs\,\mathrm d\mu(x)=0$.
Biorąc supremum, otrzymujemy zero.
\end{proof}

\begin{definicja}[Całka po podzbiorze]
\label{def:integral-over-set}
Dla $A\in\mathcal F$ i mierzalnej $f:X\to[0,\infty]$
określamy
$\int_A f(x)\,\mathrm d\mu(x):=
\int_X f(x)\chi_A(x)\,\mathrm d\mu(x)$,
gdzie iloczyn $f\chi_A$ rozumiemy jako funkcję
równą $f$ na $A$ i zero na $X\setminus A$.
To ostatnie określenie usuwa dwuznaczność zapisu
$0\cdot\infty$. Równoważnie całkujemy funkcję
$f|_A$ względem ograniczenia miary $\mu|_A$
na $(A,\mathcal F|_A)$.
\end{definicja}

\begin{lemat}[Zgodność z ograniczeniem miary]
\label{lem:restriction-integral}
Dla dowolnych $A\in\mathcal F$ i mierzalnej
$f:X\to[0,\infty]$ zachodzi
\[
\int_A f(x)\,\mathrm d\mu(x)
=\int_A f|_A(x)\,\mathrm d(\mu|_A)(x).
\]
\end{lemat}
\begin{proof}
Dla funkcji prostej zachodzi to z definicji i
$\mu|_A(A\cap B)=\mu(A\cap B)$.
Jeżeli $t:A\to[0,\infty)$ jest prostą mierzalną
minorantą $f|_A$, to przedłużenie jej zerem poza $A$
jest prostą mierzalną funkcją na $X$, mniejszą od
$f\chi_A$. Odwrotnie, ograniczenie każdej
prostej minoranty $f\chi_A$ do $A$ jest
prostą minorantą $f|_A$. Równość całek prostych
i supremów daje tezę.
\end{proof}

\begin{twierdzenie}[Nierówność Markowa dla całki]
\label{thm:markov-nonnegative}
Niech $f:X\to[0,\infty]$ będzie mierzalna i $a>0$.
Wówczas
\begin{equation}\label{eq:markov}
\mu(\{x\in X:f(x)\ge a\})
\le \frac1a\int_X f(x)\,\mathrm d\mu(x).
\end{equation}
Jeżeli $\int_Xf\,\mathrm d\mu(x)<\infty$, to $f$ jest
skończona $\mu$-prawie wszędzie. Ponadto
$\int_Xf\,\mathrm d\mu(x)=0$ wtedy i tylko wtedy,
gdy $f=0$ $\mu$-prawie wszędzie.
\end{twierdzenie}
\begin{proof}
Dla $E_a=\{f\ge a\}$ prosta funkcja
$a\chi_{E_a}$ nie przekracza $f$. Zatem
$a\mu(E_a)=\int_X a\chi_{E_a}\,\mathrm d\mu(x)
\le\int_Xf\,\mathrm d\mu(x)$, co daje
\eqref{eq:markov}. Jeśli ostatnia całka jest skończona,
to $\mu(\{f=\infty\})\le\mu(\{f\ge n\})
\le n^{-1}\int_Xf\,\mathrm d\mu(x)\to0$.
Jeżeli całka wynosi zero, to
$\mu(\{f\ge 1/n\})=0$ dla wszystkich $n$, a zbiór
$\{f>0\}=\bigcup_{n\ge1}\{f\ge1/n\}$ jest zerowy.
Odwrotnie, jeśli $f=0$ poza mierzalnym zbiorem zerowym,
to część (v) twierdzenia
\ref{thm:nonnegative-first-properties} daje $\int_X f=0$.
\end{proof}

\begin{przyklad}[Zero prawie wszędzie, mimo dużych wartości]
Na $\mathbb R$ z miarą Lebesgue'a funkcja
$f=+\infty\cdot\chi_{\mathbb Q}$, rozumiana jako
$f(x)=+\infty$ dla $x\in\mathbb Q$ i $0$ dla
$x\notin\mathbb Q$, jest mierzalna i ma całkę zero.
To pokazuje, dlaczego nie wystarczy traktować
całki jako zwykłego sumowania wartości punktowych.
\end{przyklad}

\section{Funkcje ze znakiem i całkowalność}

\begin{twierdzenie}[Niezmienniczość całki przy zmianie na zbiorze zerowym]
\label{thm:ae-integral-invariance}
Niech $f,g:X\to[0,\infty]$ będą mierzalne i
$f=g$ $\mu$-prawie wszędzie. Wówczas
$\int_Xf\,\mathrm d\mu(x)=\int_Xg\,\mathrm d\mu(x)$.
Jeśli $f,g:X\to\overline{\mathbb R}$ mają określone
całki ze znakiem i są równe prawie wszędzie, ich
całki również są równe.
\end{twierdzenie}
\begin{proof}
Wybierzmy $N\in\mathcal F$ o $\mu(N)=0$ takie,
że $f=g$ poza $N$. Dla dowolnej prostej minoranty
$0\le s\le f$ funkcja $s\chi_{X\setminus N}$
jest prostą minorantą $g$, a
\[
\int_Xs\,\mathrm d\mu(x)
=\int_Xs\chi_{X\setminus N}\,\mathrm d\mu(x),
\]
ponieważ składniki odpowiadające częściom poziomic
zawartym w $N$ mają miarę zero. Stąd
$\int_X f\le\int_X g$ po przejściu do supremum.
Zamiana ról $f,g$ daje nierówność odwrotną.
W przypadku funkcji ze znakiem stosujemy część
pierwszą osobno do ich części dodatnich i ujemnych.
Równość wartości wyklucza możliwość, by tylko jedna
z odpowiadających całek była nieokreślona.
\end{proof}

\begin{definicja}[Części dodatnia i ujemna]
\label{def:positive-negative-parts}
Dla mierzalnej $f:X\to\overline{\mathbb R}$ definiujemy
$f^+(x)=\max\{f(x),0\}$ i
$f^-(x)=\max\{-f(x),0\}$. Są to mierzalne funkcje
$[0,\infty]$-wartościowe, a $|f|=f^++f^-$.
Jeśli co najmniej jedna z liczb
$\int_Xf^+\,\mathrm d\mu(x)$,
$\int_Xf^-\,\mathrm d\mu(x)$ jest skończona, definiujemy
\begin{equation}\label{eq:signed-integral}
\int_X f(x)\,\mathrm d\mu(x)
:=\int_Xf^+(x)\,\mathrm d\mu(x)
-\int_Xf^-(x)\,\mathrm d\mu(x).
\end{equation}
Jeśli obie liczby są nieskończone, całka funkcji $f$
\emph{nie jest określona}. Funkcję $f$ nazywamy
\emph{(absolutnie) całkowalną}, jeżeli
$\int_X|f|\,\mathrm d\mu(x)<\infty$.
\end{definicja}

\begin{uwaga}[Dlaczego nie używamy odejmowania nieskończoności?]
Niech $X=\mathbb R$ z miarą Lebesgue'a i
$f(x)=x$. Wtedy zarówno całka $f^+$, jak i $f^-$
wynosi $\infty$, więc całka $f$ na całej $\mathbb R$
nie jest określona. Symetria i formalne równanie
"$\infty-\infty=0$" nie zastępują definicji.
Wartość główna Cauchy'ego, jeśli jest rozważana,
stanowi odrębne pojęcie wymagające osobnej definicji.
\end{uwaga}

\begin{twierdzenie}[Kryterium absolutnej całkowalności]
\label{thm:integrability-criterion}
Dla mierzalnej $f:X\to\overline{\mathbb R}$ następujące
warunki są równoważne:
\begin{enumerate}[label=\textup{(\roman*)}]
\item $\int_X|f|\,\mathrm d\mu(x)<\infty$;
\item $\int_Xf^+\,\mathrm d\mu(x)<\infty$ oraz
$\int_Xf^-\,\mathrm d\mu(x)<\infty$.
\end{enumerate}
W takim przypadku $f$ jest skończona prawie wszędzie,
całka $\int_Xf\,\mathrm d\mu(x)$ jest rzeczywista i
$|\int_Xf\,\mathrm d\mu(x)|\le\int_X|f|\,\mathrm d\mu(x)$.
\end{twierdzenie}
\begin{proof}
Ponieważ $0\le f^\pm\le|f|$, monotoniczność całki
daje (i)$\Rightarrow$(ii). Dla implikacji odwrotnej
nie możemy jeszcze korzystać bez dowodu z addytywności
całki dla dowolnych funkcji nieujemnych. Do pełnej addytywności wrócimy w rozdziale
\ref{ch:twierdzenia-graniczne}. Tu wykorzystamy
wyłącznie rozdzielność nośników części dodatniej i ujemnej.
Bezpośredni argument jest następujący: jeśli
$s\le f^++f^-$ jest nieujemna i prosta, to
dzielimy dziedzinę na mierzalne zbiory
$P=\{f\ge0\}$ i $M=\{f<0\}$. Funkcje
$s\chi_P\le f^+$ i $s\chi_M\le f^-$
są proste, a
$\int_Xs\,\mathrm d\mu(x)=
\int_Xs\chi_P\,\mathrm d\mu(x)+
\int_Xs\chi_M\,\mathrm d\mu(x)$
na mocy twierdzenia \ref{thm:simple-properties}.
Przechodząc do supremum po $s\le|f|=f^++f^-$,
otrzymujemy
$\int_X|f|\,\mathrm d\mu(x)\le
\int_Xf^+\,\mathrm d\mu(x)+
\int_Xf^-\,\mathrm d\mu(x)<\infty$.
Z twierdzenia \ref{thm:markov-nonnegative} otrzymujemy
skończoność $|f|$ prawie wszędzie. Wreszcie, dla dowolnych prostych minorant
$0\le u\le f^+$ i $0\le v\le f^-$ zachodzi
$u+v\le|f|$. Stąd
$\int_Xu\,\mathrm d\mu(x)+\int_Xv\,\mathrm d\mu(x)
=\int_X(u+v)\,\mathrm d\mu(x)
\le\int_X|f|\,\mathrm d\mu(x)$.
Przechodząc niezależnie do supremum po $u$ i $v$,
uzyskujemy także nierówność
$\int_Xf^+\,\mathrm d\mu(x)+\int_Xf^-\,\mathrm d\mu(x)
\le\int_X|f|\,\mathrm d\mu(x)$.
Z obydwu nierówności wynika równość tych liczb.
W szczególności $|\int_Xf\,\mathrm d\mu(x)|
\le\int_Xf^+\,\mathrm d\mu(x)+
\int_Xf^-\,\mathrm d\mu(x)=\int_X|f|\,\mathrm d\mu(x)$.
\end{proof}

\begin{definicja}[Całka funkcji całkowalnej po zbiorze]
\label{def:signed-integral-over-set-7}
Niech $A\in\mathcal F$, a $f:X\to\mathbb R$ będzie
mierzalna oraz całkowalna, tj.
$\int_X|f(x)|\,\mathrm d\mu(x)<\infty$.
Definiujemy
\[
\int_A f(x)\,\mathrm d\mu(x)
:=\int_X\widetilde f_A(x)\,\mathrm d\mu(x),
\qquad
\widetilde f_A(x)=
\begin{cases}f(x),&x\in A,\\0,&x\notin A.\end{cases}
\]
Funkcja $\widetilde f_A$ jest mierzalna oraz
$|\widetilde f_A|\le|f|$, więc całka po prawej
jest skończona. Zgodnie z definicją całki ze znakiem
również
$\int_A f=\int_A f^+-\int_A f^-$;
obie całki są skończone. Definicja nie wymaga
żadnej konwencji dla iloczynu $0\cdot\infty$.
\end{definicja}

\begin{uwaga}[Szczegół przy wartości $+\infty$]
W dowodzie wprowadziliśmy rozbicie
$P=\{f\ge0\}$ i $M=\{f<0\}$.
Dla $f(x)=+\infty$ obowiązuje $x\in P$,
a dla $f(x)=-\infty$ obowiązuje $x\in M$.
Nie wykonujemy w tych punktach odejmowania
$\infty-\infty$: obie części dodatnia i ujemna
nie są jednocześnie dodatnie w żadnym punkcie.
\end{uwaga}

\begin{przyklad}[Miary dyskretne i sumy szeregów]
Niech $X=\mathbb N$, $\mathcal F=\mathcal P(\mathbb N)$
i $\mu(A)=\sum_{n\in A}w_n$, gdzie $w_n\ge0$.
Dla funkcji $f:\mathbb N\to[0,\infty]$ całka
jest równa zwykłej sumie nieujemnego szeregu
$\int_{\mathbb N} f(n)\,\mathrm d\mu(n)=
\sum_{n=1}^{\infty}f(n)w_n$.
Istotnie, suma po pierwszych $N$ punktach jest
całką prostej funkcji nieprzekraczającej $f$,
więc prawa strona nie przekracza całki.
Z drugiej strony, dla dowolnej prostej $s\le f$
przeliczalna addytywność miary pozwala zapisać
$\int_X s\,\mathrm d\mu(x)=
\sum_{n=1}^{\infty}s(n)w_n\le\sum_{n=1}^{\infty}f(n)w_n$.
Dla rzeczywistej $f$ warunek absolutnej całkowalności
jest równoważny
$\sum_n|f(n)|w_n<\infty$.
W szczególności miara licząca odpowiada $w_n=1$.
\end{przyklad}

\begin{przyklad}[Całka może istnieć jako $+\infty$]
Na $(0,1]$ z miarą Lebesgue'a połóżmy $f(x)=1/x$.
Dla $j=0,\ldots,n-1$ na przedziale
$I_j=(2^{-(j+1)},2^{-j}]$ mamy $f(x)\ge2^j$.
Funkcja prosta
$t_n=\sum_{j=0}^{n-1}2^j\chi_{I_j}$
spełnia $0\le t_n\le f$, a jej całka wynosi
\[
\int_{(0,1]}t_n(x)\,\mathrm d\lambda_1(x)
=\sum_{j=0}^{n-1}2^j\lambda_1(I_j)
=\sum_{j=0}^{n-1}2^j2^{-(j+1)}=n/2.
\]
Z definicji całki nieujemnej
$\int_{(0,1]}f(x)\,\mathrm d\lambda_1(x)\ge n/2$
dla każdego $n$, a więc całka jest równa $+\infty$.
Nie trzeba do tego rachunku twierdzenia Beppo Leviego.
\end{przyklad}

\begin{uwaga}[Istnienie całki ze znakiem nie jest tym samym co całkowalność]
Funkcja nieujemna $f(x)=1/x$ na $(0,1]$ ma
określoną całkę równą $+\infty$, ale nie jest
całkowalna w sensie $L^1$. Z kolei dla $f(x)=x$
na $\mathbb R$ obie całki części dodatniej i ujemnej
są nieskończone, a więc całka ze znakiem nie jest
nawet określona. Nie należy utożsamiać trzech sytuacji:
całki skończonej, całki nieskończonej i wyrażenia
nieokreślonego.
\end{uwaga}

\section{Wzajemny stosunek całki i mierzalności}

\begin{przyklad}[Funkcja Dirichleta]
Na przedziale $[0,1]$ funkcja
$f=\chi_{\mathbb Q\cap[0,1]}$ jest borelowsko
mierzalna i nieujemna. Ponieważ zbiór liczb
wymiernych jest przeliczalny, ma miarę Lebesgue'a
zero, a więc $\int_{[0,1]} f\,\mathrm d\lambda_1(x)=0$.
Jednocześnie funkcja jest nieciągła w każdym punkcie,
więc nie jest całkowalna w sensie Riemanna.
\end{przyklad}

\begin{przyklad}[Całkowalność na przestrzeni o nieskończonej mierze]
Na $\mathbb R$ z miarą Lebesgue'a funkcja
$f(x)=1/(1+x^2)$ jest nieujemna i całkowalna, choć
$\lambda_1(\mathbb R)=\infty$. Istotnie, na $[-1,1]$
mamy $f\le1$. Dla $n\ge1$ i $x\in[n,n+1)$ mamy
$f(x)\le n^{-2}$; analogicznie na ujemnej półosi.
Dla dowolnej funkcji prostej $0\le s\le f$,
przeliczalna addytywność miary na poziomicach $s$
pozwala rozłożyć jej całkę na rozłączne przedziały
$[-1,1]$, $[n,n+1)$ i $(-(n+1),-n]$, $n\ge1$.
Na każdym przedziale całka $s$ jest nie większa
od iloczynu jego długości i odpowiednio $1$ albo
$n^{-2}$, więc
$\int_{\mathbb R}s\,\mathrm d\lambda_1(x)
\le2+2\sum_{n\ge1}n^{-2}<\infty$.
Ta nierówność jest niezależna od wyboru $s$;
biorąc supremum w definicji całki, otrzymujemy
$\int_{\mathbb R} f(x)\,\mathrm d\lambda_1(x)<\infty$. Dokładną wartość $\pi$ wyprowadzimy
po omówieniu twierdzenia o zamianie zmiennych.
\end{przyklad}

\begin{uwaga}[Jeszcze nie wszystko zostało dowiedzione]
Na tym etapie znamy monotoniczność, jednorodność,
całki funkcji prostych i kryteria całkowalności.
Nie używaliśmy bez dowodu równości
$\int_X(f+g)=\int_X f+\int_X g$ dla \emph{dowolnych}
nieujemnych funkcji mierzalnych. Tę fundamentalną
własność oraz równość $\int_X\lim f_n=\lim\int_X f_n$
w odpowiednich warunkach wyprowadzimy wprost
z definicji w kolejnym rozdziale.
\end{uwaga}

% END chapters/pl/measure_theory/07-rozdzial.tex


% BEGIN chapters/pl/measure_theory/08-rozdzial.tex
% ROZDZIAŁ 8 — do włączenia w main.tex
\chapter{Twierdzenia graniczne dla całki Lebesgue'a}
\label{ch:twierdzenia-graniczne}

\section{Dlaczego nie wolno bezwarunkowo przestawiać granicy i całki?}

Punktem wyjścia jest pytanie, kiedy dla $f_n\to f$ wolno
zamienić kolejność granicy i całkowania:
\[
\int_X f(x)\,\mathrm d\mu(x)
=\lim_{n\to\infty}\int_X f_n(x)\,\mathrm d\mu(x).
\]
Zbieżność punktowa sama nie wystarcza. Na przykład na
$(0,1)$ z miarą Lebesgue'a funkcje
$f_n=n\chi_{(0,1/n)}$ zbiegają do zera w każdym
punkcie, ale ich całki wynoszą $1$.
Twierdzenia tego rozdziału podają precyzyjne warunki,
w których przejście do granicy jest uprawnione.

W całym rozdziale $(X,\mathcal F,\mu)$ jest przestrzenią
z miarą. Każde występujące całkowanie oznacza całkowanie
w sensie rozdziału \ref{ch:calkowanie-lebesgue}.
Nie zakładamy, że $\mu$ jest skończona ani zupełna,
jeśli nie zaznaczamy tego wyraźnie. Symbol $f_n\uparrow f$
oznacza niemalejącą zbieżność punktową do $f$.

\section{Fundamentalny lemat aproksymacyjny}

\begin{lemat}[Zbieżność całek prostej funkcji na zbiorach rosnących]
\label{lem:simple-exhaustion}
Niech $s:X\to[0,\infty)$ będzie mierzalną funkcją
prostą, a $(E_n)_{n\ge1}\subseteq\mathcal F$
niemalejącym ciągiem zbiorów. Jeśli
$E=\bigcup_{n\ge1}E_n$, to
\[
\lim_{n\to\infty}\int_{E_n}s(x)\,\mathrm d\mu(x)
=\int_Es(x)\,\mathrm d\mu(x).
\]
Równość zachodzi również wtedy, gdy prawa strona
wynosi $+\infty$.
\end{lemat}
\begin{proof}
Niech $s=\sum_{j=1}^m a_j\chi_{A_j}$
będzie reprezentacją na mierzalnym podziale, gdzie
$a_j\ge0$. Z definicji całki po zbiorze
$\int_{E_n}s\,\mathrm d\mu(x)=
\sum_{j=1}^m a_j\mu(A_j\cap E_n)$.
Zbiory $A_j\cap E_n$ rosną do $A_j\cap E$, więc
ciągłość miary z dołu (twierdzenie
\ref{thm:measure-continuity}) daje
$\mu(A_j\cap E_n)\uparrow\mu(A_j\cap E)$.
Ponieważ suma zawiera skończenie wiele nieujemnych
składników, możemy przejść do granicy w każdym z nich;
przy współczynniku $a_j=0$ obie strony są zerem.
Otrzymujemy wzór dla $E$.
\end{proof}

\section{Twierdzenie Beppo Leviego}

\begin{twierdzenie}[O zbieżności monotonicznej; Beppo Levi]
\label{thm:mct}
Niech $(f_n)_{n\ge1}$ będzie ciągiem mierzalnych funkcji
$f_n:X\to[0,\infty]$, takim że dla każdego $x\in X$
$f_n(x)\le f_{n+1}(x)$. Niech
$f(x)=\lim_{n\to\infty}f_n(x)=\sup_nf_n(x)$.
Wtedy $f$ jest mierzalna i
\begin{equation}\label{eq:mct}
\int_Xf(x)\,\mathrm d\mu(x)
=\lim_{n\to\infty}\int_Xf_n(x)\,\mathrm d\mu(x)
=\sup_{n\ge1}\int_Xf_n(x)\,\mathrm d\mu(x).
\end{equation}
Dopuszczamy wartość $+\infty$ i nie wymagamy
żadnego założenia o $\mu(X)$.
\end{twierdzenie}
\begin{proof}
Mierzalność $f$ wykazano w twierdzeniu
\ref{thm:pointwise-limits}. Z monotoniczności całki
(rozdział 7) mamy
$L:=\sup_n\int_Xf_n\,\mathrm d\mu(x)\le
\int_Xf\,\mathrm d\mu(x)$.
Trzeba wykazać trudniejszą nierówność odwrotną.

Ustalmy dowolną nieujemną mierzalną funkcję prostą
$s\le f$ oraz liczbę $\theta\in(0,1)$.
Definiujemy
$E_n=\{x\in X:f_n(x)\ge\theta s(x)\}$.
Zbiory $E_n$ są mierzalne i rosną.
Ponadto $\bigcup_nE_n=X$: jeżeli $s(x)=0$,
to $x\in E_n$ dla każdego $n$; jeżeli $s(x)>0$,
to $f(x)\ge s(x)>\theta s(x)$, więc z
$f_n(x)\uparrow f(x)$ wynika $f_n(x)\ge\theta s(x)$
dla dostatecznie dużego $n$. Argument działa także
gdy $f(x)=+\infty$, ponieważ $s(x)$ jest skończona.

Z nierówności $f_n\ge\theta s\chi_{E_n}$,
monotoniczności i jednorodności całki wynika
\[
\int_Xf_n\,\mathrm d\mu(x)\ge
\theta\int_{E_n}s(x)\,\mathrm d\mu(x).
\]
Przechodząc do granicy, stosując lemat
\ref{lem:simple-exhaustion} i $E_n\uparrow X$,
dostajemy $L\ge\theta\int_Xs\,\mathrm d\mu(x)$.
Niech $\theta\uparrow1$; dla całki skończonej
wynika $L\ge\int_Xs\,\mathrm d\mu(x)$, a jeśli
$\int_Xs=\infty$, to z powyższej nierówności
już dla jednego $\theta>0$ otrzymujemy $L=\infty$.
Na koniec bierzemy supremum po wszystkich
nieujemnych prostych $s\le f$. Na mocy definicji
\ref{def:integral-positive} daje to
$L\ge\int_Xf\,\mathrm d\mu(x)$. Obie nierówności
kończą dowód.
\end{proof}

\begin{wniosek}[Konstruktywne obliczanie całki]
\label{cor:simple-approx-integral}
Jeżeli $f:X\to[0,\infty]$ jest mierzalna, a
$(s_n)$ jest \emph{dowolnym} ciągiem nieujemnych
funkcji prostych mierzalnych takim, że $s_n\uparrow f$
wszędzie, to
$\int_Xf\,\mathrm d\mu(x)=
\lim_{n\to\infty}\int_Xs_n\,\mathrm d\mu(x)$.
W szczególności można użyć funkcji z
\eqref{eq:canonical-simple}.
\end{wniosek}
\begin{proof}
To bezpośrednie zastosowanie twierdzenia
\ref{thm:mct} do ciągu $s_n$; całki funkcji prostych
znamy już z definicji \ref{def:integral-simple}.
\end{proof}

\begin{przyklad}[Ciąg monotoniczny o nieskończonej granicy całek]
Na $(0,1)$ z miarą Lebesgue'a funkcje
$f_n(x)=\min\{n,1/x\}$ są niemalejące i zbiegają
do $1/x$. Twierdzenie \ref{thm:mct} daje
$\int_{(0,1)}x^{-1}\,\mathrm d\lambda_1(x)=
\lim_n\int_{(0,1)}f_n\,\mathrm d\lambda_1(x)=+\infty$,
co jest zgodne z bezpośrednim rachunkiem z
rozdziału \ref{ch:calkowanie-lebesgue}.
Nie wolno wywnioskować, że same $f_n$ mają
nieograniczone wartości punktowe: każda z nich
jest ograniczona przez $n$.
\end{przyklad}

\begin{twierdzenie}[Zbieżność monotoniczna prawie wszędzie]
\label{thm:mct-ae}
Niech $f_n,f:X\to[0,\infty]$ będą mierzalne.
Jeśli $f_n\le f_{n+1}$ i $f_n\to f$
$\mu$-prawie wszędzie, to zachodzi równość
\eqref{eq:mct}.
\end{twierdzenie}
\begin{proof}
Niech $N\in\mathcal F$ będzie zbiorem miary zero,
poza którym monotoniczność i zbieżność zachodzą.
Funkcje $g_n=f_n\chi_{X\setminus N}$
są mierzalne, nieujemne i niemalejące wszędzie,
a ich granica jest równa
$g=f\chi_{X\setminus N}$.
Stosujemy twierdzenie \ref{thm:mct} do $g_n$.
Równość całek $g_n$ i $f_n$ oraz całek $g$ i $f$
wynika z twierdzenia \ref{thm:ae-integral-invariance}.
Dzięki temu zastosowanie wersji punktowej
Beppo Leviego daje dokładnie żądaną równość.
\end{proof}

\section{Addytywność i liniowość całki}

\begin{twierdzenie}[Addytywność całki nieujemnej]
\label{thm:positive-additivity}
Dla dowolnych mierzalnych $f,g:X\to[0,\infty]$
zachodzi
\begin{equation}\label{eq:positive-additivity}
\int_X(f+g)(x)\,\mathrm d\mu(x)
=\int_Xf(x)\,\mathrm d\mu(x)+
\int_Xg(x)\,\mathrm d\mu(x).
\end{equation}
Obie strony należą do $[0,\infty]$; dopuszczamy
$+\infty$ po którejkolwiek stronie.
\end{twierdzenie}
\begin{proof}
Z twierdzenia \ref{thm:simple-up} wybierzmy
nieujemne funkcje proste $s_n\uparrow f$ i
$t_n\uparrow g$. Wtedy $s_n+t_n\uparrow f+g$.
Twierdzenie \ref{thm:mct}, a następnie addytywność
całki prostej, dają
\[
\int_X(f+g)\,\mathrm d\mu(x)
=\lim_n\int_X(s_n+t_n)\,\mathrm d\mu(x)
=\lim_n\left(\int_Xs_n\,\mathrm d\mu(x)+
\int_Xt_n\,\mathrm d\mu(x)\right).
\]
Ponowne zastosowanie twierdzenia \ref{thm:mct}
do $s_n$ i $t_n$ daje tezę; suma dwóch ciągów
niemalejących liczb nieujemnych ma granicę równą
sumie ich granic także w $[0,\infty]$.
\end{proof}

\begin{wniosek}[Rozkład całki na rozłączne zbiory]
\label{cor:integral-countable-partition}
Niech $f:X\to[0,\infty]$ będzie mierzalna, a
$(A_n)_{n\ge1}$ przeliczalną rodziną parami rozłącznych
zbiorów z $\mathcal F$. Wtedy
\begin{equation}\label{eq:integral-partition}
\int_{\bigcup_{n\ge1}A_n}f(x)\,\mathrm d\mu(x)
=\sum_{n=1}^{\infty}\int_{A_n}f(x)\,\mathrm d\mu(x).
\end{equation}
\end{wniosek}
\begin{proof}
Dla każdego $N$ zachodzi punktowo
$f\chi_{\cup_{n=1}^NA_n}=
\sum_{n=1}^Nf\chi_{A_n}$.
Addytywność całki daje wzór dla sumy skończonej.
Gdy $N\to\infty$, lewe funkcje rosną do
$f\chi_{\cup_{n\ge1}A_n}$. Zastosowanie
Beppo Leviego kończy dowód.
\end{proof}

\begin{twierdzenie}[Liniowość całki funkcji całkowalnych]
\label{thm:integral-linear-real}
Jeżeli $f,g:X\to\mathbb R$ są mierzalne i
$\int_X|f|\,\mathrm d\mu(x)<\infty$,
$\int_X|g|\,\mathrm d\mu(x)<\infty$, to dla
$a,b\in\mathbb R$ funkcja $af+bg$ jest całkowalna i
\[
\int_X(af+bg)(x)\,\mathrm d\mu(x)
=a\int_Xf(x)\,\mathrm d\mu(x)+
b\int_Xg(x)\,\mathrm d\mu(x).
\]
Ponadto $\bigl|\int_Xf\,\mathrm d\mu(x)\bigr|
\le\int_X|f|\,\mathrm d\mu(x)$.
\end{twierdzenie}
\begin{proof}
Dla $a\ge0$ rozkłady $(af)^+=af^+$ i
$(af)^-=af^-$ oraz jednorodność całki nieujemnej
dają $\int_X af=a\int_X f$. Dla $a<0$ korzystamy z
$(-f)^+=f^-$ i $(-f)^-=f^+$.

Z nierówności $|f+g|\le|f|+|g|$ i addytywności
nieujemnej wynika całkowalność sumy.
W każdym punkcie zachodzi tożsamość
\[
(f+g)^++f^-+g^-=(f+g)^-+f^++g^+.
\]
Wszystkie wyrazy mają skończone całki, bo są
ograniczone przez sumy $|f|+|g|$.
Całkujemy tę równość, stosujemy addytywność
nieujemną i przenosimy skończone składniki na
drugą stronę. Otrzymujemy
$\int_X(f+g)=\int_X f+\int_X g$.
Połączenie obu własności daje pełną liniowość.
Na koniec z $0\le f^\pm\le|f|$ i
$|\int_X f^+-\int_X f^-|\le\int_X f^++\int_X f^-
=\int_X|f|$ wynika oszacowanie.
\end{proof}

\begin{wniosek}[Nierówność trójkąta i struktura $L^1$]
\label{cor:integral-triangle}
Jeśli $f,g:X\to\mathbb R$ są całkowalne, to
$\int_X|f+g|\,\mathrm d\mu(x)\le
\int_X|f|\,\mathrm d\mu(x)+
\int_X|g|\,\mathrm d\mu(x)$.
Funkcjonał $f\mapsto\int_Xf\,\mathrm d\mu(x)$
jest liniowy na przestrzeni rzeczywistych
funkcji całkowalnych, a nierówność
$|\int_X f|\le\int_X|f|$ wyznacza jego ograniczenie.
Po utożsamieniu funkcji równych prawie wszędzie
otrzymamy przestrzeń normowaną $L^1$ w rozdziale 14.
\end{wniosek}
\begin{proof}
Całkujemy punktową nierówność $|f+g|\le|f|+|g|$
i stosujemy addytywność całki nieujemnej. Liniowość
i oszacowanie funkcjonału zostały udowodnione
w twierdzeniu \ref{thm:integral-linear-real}.
\end{proof}

\begin{definicja}[Całka funkcji zespolonej]
\label{def:complex-integral}
Funkcję mierzalną $f:X\to\mathbb C$ nazywamy
całkowalną, jeżeli $\int_X|f|\,\mathrm d\mu(x)<\infty$.
Wówczas określamy
$\int_Xf\,\mathrm d\mu(x)=
\int_X\operatorname{Re}f\,\mathrm d\mu(x)+
 i\int_X\operatorname{Im}f\,\mathrm d\mu(x)$.
\end{definicja}

\begin{twierdzenie}[Liniowość zespolona i oszacowanie modułu]
\label{thm:integral-complex-linear}
Jeśli $f,g:X\to\mathbb C$ są całkowalne i
$\alpha,\beta\in\mathbb C$, to $\alpha f+\beta g$
jest całkowalna,
$\int_X(\alpha f+\beta g)\,\mathrm d\mu(x)=
\alpha\int_Xf\,\mathrm d\mu(x)+
\beta\int_Xg\,\mathrm d\mu(x)$, a także
\[
\left|\int_Xf(x)\,\mathrm d\mu(x)\right|
\le\int_X|f(x)|\,\mathrm d\mu(x).
\]
\end{twierdzenie}
\begin{proof}
Z $|\alpha f+\beta g|\le|\alpha||f|+|\beta||g|$
wynika całkowalność. Liniowość otrzymujemy,
rozwijając części rzeczywiste i urojone i stosując
twierdzenie \ref{thm:integral-linear-real}.
Jeśli $I=\int_Xf\,\mathrm d\mu(x)\ne0$, wybierzmy
$\zeta=\overline I/|I|$, więc $|\zeta|=1$ i
$\zeta I=|I|$. Wtedy
$|I|=\int_X\operatorname{Re}(\zeta f)\,\mathrm d\mu(x)
\le\int_X|\operatorname{Re}(\zeta f)|\,\mathrm d\mu(x)
\le\int_X|f|\,\mathrm d\mu(x)$.
Dla $I=0$ nierówność jest oczywista.
\end{proof}

\section{Lemat Fatou}

\begin{twierdzenie}[Lemat Fatou]
\label{thm:fatou}
Niech $f_n:X\to[0,\infty]$ będą funkcjami mierzalnymi.
Wtedy
\begin{equation}\label{eq:fatou}
\int_X\liminf_{n\to\infty}f_n(x)\,\mathrm d\mu(x)
\le\liminf_{n\to\infty}\int_Xf_n(x)\,\mathrm d\mu(x).
\end{equation}
Nie wymagamy monotoniczności ciągu ani skończoności miary.
\end{twierdzenie}
\begin{proof}
Dla $m\ge1$ połóżmy
$h_m(x)=\inf_{n\ge m}f_n(x)$.
Na mocy twierdzenia \ref{thm:pointwise-limits}
każda $h_m$ jest mierzalna. Ponieważ z rosnącym $m$
infimum obejmuje coraz mniej wyrazów, funkcje $h_m$
są niemalejące. Z definicji granicy dolnej
$h_m(x)\uparrow\liminf_nf_n(x)$.
Twierdzenie \ref{thm:mct} daje
\[
\int_X\liminf_nf_n\,\mathrm d\mu(x)
=\lim_{m\to\infty}\int_Xh_m\,\mathrm d\mu(x).
\]
Dla wszystkich $n\ge m$ mamy $h_m\le f_n$, więc
$\int_Xh_m\,\mathrm d\mu(x)\le
\inf_{n\ge m}\int_Xf_n\,\mathrm d\mu(x)$.
Przechodzimy do granicy po $m$ i otrzymujemy
prawą stronę \eqref{eq:fatou}.
\end{proof}

\begin{przyklad}[Nierówność Fatou bywa ostra]
Na $\mathbb R$ z miarą Lebesgue'a niech
$f_n=\chi_{[n,n+1]}$. Dla każdego $x$ zachodzi
$f_n(x)\to0$, ale $\int_{\mathbb R} f_n(x)\,\mathrm d\lambda_1(x)=1$.
Zatem $\int_{\mathbb R}\liminf_nf_n(x)\,\mathrm d\lambda_1(x)=0$
i $\liminf_n\int_{\mathbb R} f_n(x)\,\mathrm d\lambda_1(x)=1$.
Masa może \emph{uciekać do nieskończoności}.
\end{przyklad}

\begin{wniosek}[Odwrotny lemat Fatou przy dominacji]
\label{cor:reverse-fatou}
Niech $f_n:X\to\mathbb R$ będą mierzalne i
$0\le f_n\le g$ $\mu$-prawie wszędzie, gdzie
$g:X\to[0,\infty]$ jest całkowalna. Wtedy
\[
\limsup_n\int_Xf_n\,\mathrm d\mu(x)
\le\int_X\limsup_nf_n\,\mathrm d\mu(x).
\]
\end{wniosek}
\begin{proof}
Ponieważ $g$ jest całkowalna, $g<\infty$ prawie
wszędzie. Usuwamy więc jeden wspólny mierzalny
zbiór zerowy, na którym może wystąpić $g=\infty$
lub naruszenie nierówności $0\le f_n\le g$.
Na tym zbiorze zastępujemy $f_n$ i $g$ zerem;
całki nie zmieniają się. Możemy teraz założyć,
że $g$ jest skończona wszędzie i $0\le f_n\le g$.
Funkcje $g-f_n$ są nieujemne i mierzalne, a
$\liminf_n(g-f_n)=g-\limsup_nf_n$ w punktach,
gdzie $g$ jest skończona. Ponieważ $g<\infty$
prawie wszędzie, lemat Fatou i liniowość całki
funkcji całkowalnych dają
$\int_X g-\int_X\limsup f_n\le
\int_X g-\limsup\int_X f_n$.
Odejmując skończoną liczbę $\int_X g$,
otrzymujemy tezę.
\end{proof}

\begin{twierdzenie}[Lemat Fatou z całkowalnym ograniczeniem od dołu]
\label{thm:fatou-integrable-lower-8}
Niech $f_n:X\to\overline{\mathbb R}$ będą mierzalne,
a $g:X\to[0,\infty]$ całkowalna. Załóżmy,
że $f_n\ge-g$ prawie wszędzie dla każdego $n$.
Wówczas każda całka $\int_Xf_n\,\mathrm d\mu(x)$
jest określona w $(-\infty,\infty]$, podobnie całka
funkcji $\liminf_nf_n$, oraz
\[
\int_X\liminf_{n\to\infty}f_n(x)\,\mathrm d\mu(x)
\le \liminf_{n\to\infty}\int_Xf_n(x)\,\mathrm d\mu(x).
\]
Nie jest wymagana skończoność miary $\mu(X)$.
\end{twierdzenie}
\begin{proof}
Bierzemy jeden wspólny zbiór wyjątkowy miary zero,
na którym mogą zawodzić nierówności $f_n\ge-g$
lub skończoność $g$. Zmieniamy tam wszystkie
funkcje na zero; ich całki pozostają niezmienione.
Możemy więc zakładać, że $g<\infty$ wszędzie oraz
$f_n+g\ge0$ dla każdego $n$. Lemat Fatou daje
\[
\int_X(\liminf_nf_n+g)\,\mathrm d\mu(x)
=\int_X\liminf_n(f_n+g)\,\mathrm d\mu(x)
\le\liminf_n\int_X(f_n+g)\,\mathrm d\mu(x).
\]
Aby uzasadnić rachunek całek bez odejmowania
nieskończoności, zauważmy, że $f_n^-\le g$,
a więc $\int_X f_n^-<\infty$. Punktowa tożsamość
$(f_n+g)+f_n^-=f_n^++g$ i addytywność całki
nieujemnej dają
$\int_X(f_n+g)+\int_Xf_n^-=
\int_Xf_n^++\int_Xg$.
Wolno odjąć skończoną całkę $f_n^-$ i otrzymać
$\int_X(f_n+g)=\int_Xf_n+\int_Xg$.
Analogicznie postępujemy z $\liminf_nf_n$.
Odejmujemy tę samą skończoną liczbę $\int_Xg$
i otrzymujemy tezę. Żadne wyrażenie typu
$\infty-\infty$ nie występuje.
\end{proof}

\begin{przyklad}[Brak całkowalnego ograniczenia od dołu]
Na $(0,1)$ połóżmy $f_n=-n\chi_{(0,1/n)}$.
Mamy $f_n(x)\to0$ dla każdego $x\in(0,1)$,
ale $\int_{(0,1)}f_n(x)\,\mathrm d\lambda_1(x)=-1$.
Zatem nierówność z twierdzenia
\ref{thm:fatou-integrable-lower-8} byłaby fałszywa,
gdyby usunąć warunek dominacji od dołu.
\end{przyklad}

\section{Twierdzenie Lebesgue'a o zbieżności zmajoryzowanej}

\begin{twierdzenie}[O zbieżności zmajoryzowanej; Lebesgue]
\label{thm:dct}
Niech $f_n:X\to\mathbb R$ będą mierzalne i
$f_n(x)\to f(x)$ dla $\mu$-prawie wszystkich $x$, gdzie
$f:X\to\mathbb R$ jest mierzalna. Załóżmy, że
istnieje nieujemna całkowalna funkcja mierzalna
$g:X\to[0,\infty]$ taka, że
$|f_n(x)|\le g(x)$ dla wszystkich $n$ i
$\mu$-prawie wszystkich $x$, przy czym zbiór
wyjątkowy można wybrać wspólny dla całego ciągu.
Wówczas $f_n$ i $f$ są całkowalne oraz
\begin{equation}\label{eq:dct}
\lim_{n\to\infty}\int_Xf_n(x)\,\mathrm d\mu(x)
=\int_Xf(x)\,\mathrm d\mu(x).
\end{equation}
\end{twierdzenie}
\begin{proof}
Ponieważ $g$ jest całkowalna, zbiór $\{g=\infty\}$
jest zerowy. Po dołączeniu go do wspólnego zbioru
wyjątkowego możemy zastąpić na nim wszystkie
$f_n$, funkcję $f$ i $g$ zerem. Całki pozostają
niezmienione, a teraz dla każdego $x$ zachodzi
$f_n(x)\to f(x)$, $|f_n(x)|\le g(x)<\infty$.
W szczególności operacje $g\pm f_n$ są wszędzie
określone i nieujemne.
Z ciągłości wartości bezwzględnej wynika $|f|\le g$.
Monotoniczność daje
$\int_X|f_n|\le\int_X g<\infty$ oraz
$\int_X|f|\le\int_X g<\infty$.

Funkcje $g+f_n$ i $g-f_n$ są nieujemne.
Stosując do nich lemat Fatou, otrzymujemy
\[
\int_X(g+f)\,\mathrm d\mu(x)
\le\liminf_n\int_X(g+f_n)\,\mathrm d\mu(x)
=\int_Xg\,\mathrm d\mu(x)+
\liminf_n\int_Xf_n\,\mathrm d\mu(x),
\]
\[
\int_X(g-f)\,\mathrm d\mu(x)
\le\liminf_n\int_X(g-f_n)\,\mathrm d\mu(x)
=\int_Xg\,\mathrm d\mu(x)-
\limsup_n\int_Xf_n\,\mathrm d\mu(x).
\]
W drugim przejściu wszystkie całki $f_n$ są
skończone i jednostajnie ograniczone przez $\int_X g$.
Po skróceniu skończonej liczby $\int_X g$ pierwsza
nierówność daje
$\int_X f\le\liminf_n\int_X f_n$, a druga
$\limsup_n\int_X f_n\le\int_X f$.
Ponieważ $\liminf\le\limsup$, obie granice
są równe $\int_X f$.
\end{proof}

\begin{wniosek}[Zbieżność zmajoryzowana w $L^1$]
\label{cor:dct-l1}
Przy założeniach twierdzenia \ref{thm:dct} zachodzi
\[
\lim_{n\to\infty}\int_X|f_n(x)-f(x)|\,\mathrm d\mu(x)=0.
\]
Dla funkcji zespolonych wniosek i twierdzenie
\ref{thm:dct} także zachodzą, jeśli zastąpić
wartość bezwzględną modułem.
\end{wniosek}
\begin{proof}
Mamy $|f_n-f|\to0$ prawie wszędzie oraz
$|f_n-f|\le2g$ prawie wszędzie.
Stosujemy twierdzenie \ref{thm:dct} do nieujemnych
funkcji $|f_n-f|$. Dla funkcji zespolonych
stosujemy ten sam argument do modułu, a z
$\bigl|\int_X(f_n-f)\bigr|\le\int_X|f_n-f|$
wynika zbieżność całek.
\end{proof}

\begin{wniosek}[Twierdzenie o zbieżności ograniczonej]
\label{cor:bct}
Niech $\mu(X)<\infty$ oraz $f_n,f:X\to\mathbb R$
będą mierzalne. Jeśli $f_n\to f$ prawie wszędzie
i istnieje $M<\infty$ takie, że $|f_n|\le M$
prawie wszędzie dla wszystkich $n$, to
$\int_X f_n\to\int_X f$ i $\int_X|f_n-f|\to0$.
\end{wniosek}
\begin{proof}
Funkcja stała $g\equiv M$ jest całkowalna,
gdyż $\int_X g\,\mathrm d\mu(x)=M\mu(X)<\infty$.
Stosujemy twierdzenie \ref{thm:dct} i wniosek
\ref{cor:dct-l1}.
\end{proof}

\begin{przyklad}[Dlaczego potrzebna jest dominacja]
Na $(0,1)$ weźmy
$f_n=n\chi_{(0,1/n)}$.
Dla każdego ustalonego $x>0$ zachodzi $f_n(x)\to0$,
lecz $\int_{(0,1)}f_n\,\mathrm d\lambda_1(x)=1$.
Gdyby istniała całkowalna $g\ge f_n$ dla wszystkich
$n$, wówczas twierdzenie o zbieżności zmajoryzowanej
dawałoby sprzeczność. Problem nie leży zatem
w zbieżności punktowej, lecz w braku wspólnej
całkowalnej funkcji dominującej.
\end{przyklad}

\begin{przyklad}[Dlaczego skończoność miary ma znaczenie]
Na $\mathbb R$ funkcje $f_n=\chi_{[n,n+1]}$
są ograniczone wspólną stałą $1$, a mimo to
$\int_{\mathbb R}f_n(x)\,\mathrm d\lambda_1(x)=1\not\to0$, chociaż $f_n\to0$ wszędzie.
Stała $1$ nie jest całkowalna na $\mathbb R$,
więc twierdzenie \ref{cor:bct} nie ma zastosowania.
\end{przyklad}

\section{Szeregi funkcji i całek}

\begin{twierdzenie}[Szeregi nieujemnych funkcji mierzalnych]
\label{thm:sum-nonnegative-integrals}
Niech $f_n:X\to[0,\infty]$ będą mierzalne. Wtedy
\begin{equation}\label{eq:sum-nonnegative-integrals}
\int_X\sum_{n=1}^{\infty}f_n(x)\,\mathrm d\mu(x)
=\sum_{n=1}^{\infty}\int_Xf_n(x)\,\mathrm d\mu(x),
\end{equation}
gdzie obie strony mogą być nieskończone.
\end{twierdzenie}
\begin{proof}
Sumy częściowe $s_N=\sum_{n=1}^N f_n$ są
mierzalne, nieujemne i rosną do
$s=\sum_{n\ge1}f_n$. Z twierdzenia \ref{thm:mct}
$\int_Xs\,\mathrm d\mu(x)=\lim_N\int_Xs_N\,\mathrm d\mu(x)$.
Dla każdego $N$ stosujemy skończoną addytywność
całki nieujemnej z twierdzenia
\ref{thm:positive-additivity}; granica jest
z definicji sumą szeregu nieujemnego.
\end{proof}

\begin{twierdzenie}[Całkowanie szeregów absolutnie sumowalnych]
\label{thm:sum-absolutely-integrable}
Niech $f_n:X\to\mathbb C$ będą mierzalne i
całkowalne oraz niech
$\sum_{n\ge1}\int_X|f_n(x)|\,\mathrm d\mu(x)<\infty$.
Wówczas $\sum_nf_n(x)$ jest bezwzględnie zbieżny
$\mu$-prawie wszędzie. Jeśli $F$ oznacza jego sumę
na zbiorze zbieżności i zero na pozostałym zbiorze,
to $F$ jest całkowalna i
\begin{equation}\label{eq:interchange-series}
\int_XF(x)\,\mathrm d\mu(x)
=\sum_{n=1}^{\infty}\int_Xf_n(x)\,\mathrm d\mu(x).
\end{equation}
\end{twierdzenie}
\begin{proof}
Niech $G=\sum_{n\ge1}|f_n|$, mierzalna suma
nieujemna. Z twierdzenia
\ref{thm:sum-nonnegative-integrals}
$\int_XG\,\mathrm d\mu(x)=
\sum_n\int_X|f_n|\,\mathrm d\mu(x)<\infty$.
Na mocy twierdzenia \ref{thm:markov-nonnegative}
$G(x)<\infty$ prawie wszędzie. Poza zbiorem
zerowym szereg $\sum_nf_n(x)$ zbiega więc
bezwzględnie. Zbiór zbieżności jest mierzalny,
ponieważ zarówno części rzeczywiste, jak i
urojone sum częściowych są mierzalne, a warunek
Cauchy'ego można zapisać jako przeliczalną
kombinację mierzalnych nierówności.
Dla sum częściowych $S_N=\sum_{n=1}^Nf_n$
mamy $|S_N|\le G$ oraz $S_N\to F$ prawie wszędzie.
Zespolona wersja twierdzenia o zbieżności zmajoryzowanej
daje $\int_XS_N\to\int_XF$. Liniowość całki
zespolonej pozwala zapisać
$\int_XS_N=\sum_{n=1}^N\int_Xf_n$, co daje tezę.
\end{proof}

\begin{twierdzenie}[Wzór warstwowy na całkę nieujemną]
\label{thm:layer-cake-8}
Niech $f:X\to[0,\infty]$ będzie mierzalna.
Wtedy funkcja $t\mapsto\mu(\{x:f(x)>t\})$
jest borelowska na $[0,\infty)$ i zachodzi
\begin{equation}\label{eq:layer-cake-8}
\int_Xf(x)\,\mathrm d\mu(x)
=\int_{[0,\infty)}\mu(\{x:f(x)>t\})\,
\mathrm d\lambda_1(t).
\end{equation}
Obie strony mogą być nieskończone. Dowód nie korzysta
z twierdzenia Tonellego, które pojawi się później.
\end{twierdzenie}
\begin{proof}
Funkcja $t\mapsto\mu(\{f>t\})$ jest \emph{nierosnąca},
a każda funkcja monotoniczna
o wartościach w $[0,\infty]$ jest borelowsko mierzalna.
Najpierw niech $s=\sum_{j=1}^m a_j\chi_{A_j}$
będzie nieujemną funkcją prostą, gdzie zbiory $A_j$
tworzą mierzalny podział $X$. Dla $t\ge0$
\[
\mu(\{s>t\})=\sum_{j=1}^m
\mu(A_j)\chi_{[0,a_j)}(t).
\]
Całkowanie skończonej sumy po $t$ daje
$\sum_ja_j\mu(A_j)=\int_Xs\,\mathrm d\mu(x)$;
przy $a_j=0$ stosujemy umowę $0\cdot\infty=0$.
Dla ogólnej $f$ weźmy $s_n\uparrow f$ z
twierdzenia~\ref{thm:simple-up}. Dla każdego $t\ge0$
mamy $\{s_n>t\}\uparrow\{f>t\}$. Ciągłość miary
z dołu daje
$\mu(\{s_n>t\})\uparrow\mu(\{f>t\})$.
Stosujemy Beppo Leviego na $[0,\infty)$ i na $X$:
\[
\int_{[0,\infty)}\mu(\{f>t\})\,
\mathrm d\lambda_1(t)
=\lim_n\int_Xs_n\,\mathrm d\mu(x)
=\int_Xf\,\mathrm d\mu(x).
\]
\end{proof}

\begin{wniosek}[Uogólniona formuła rozkładu na poziomice]\label{idx:08:wniosek:20}

Dla każdego $p>0$ i mierzalnej $f:X\to\mathbb R$
zachodzi
\[
\int_X|f(x)|^p\,\mathrm d\mu(x)
=p\int_{[0,\infty)}t^{p-1}
\mu(\{x:|f(x)|>t\})\,\mathrm d\lambda_1(t).
\]
\end{wniosek}
\begin{proof}
Niech najpierw $s=\sum_{j=1}^m a_j\chi_{A_j}$,
gdzie $a_j\ge0$, a $(A_j)$ jest mierzalnym podziałem.
Dla $t\ge0$ zachodzi
$\mu(\{s>t\})=\sum_j\mu(A_j)\chi_{[0,a_j)}(t)$.
Z addytywności całki nieujemnej i elementarnego wzoru
$\int_0^{a}pt^{p-1}\,\mathrm dt=a^p$ dla $a\ge0$
(w sensie całki Riemanna na przedziałach
$[\delta,a]$, a następnie przy $\delta\downarrow0$)
otrzymujemy
\[
p\int_{[0,\infty)}t^{p-1}\mu(\{s>t\})\,
\mathrm d\lambda_1(t)=\sum_j a_j^p\mu(A_j)
=\int_Xs(x)^p\,\mathrm d\mu(x).
\]
Przy $t=0$ i $p<1$ kładziemy wartość czynnika
$t^{p-1}$ równą zero; nie zmienia to całki.
Dla ogólnej funkcji $|f|$ wybieramy $s_n\uparrow|f|$
ze zbioru nieujemnych funkcji prostych.
Wtedy $s_n^p\uparrow|f|^p$ i
$\mu(\{s_n>t\})\uparrow\mu(\{|f|>t\})$.
Dwukrotne zastosowanie twierdzenia Beppo Leviego
daje żądaną równość bez używania twierdzenia Tonellego.
\end{proof}

\section{Przykłady obliczeń i kryteria całkowalności}

\begin{przyklad}[Całka potęgi blisko zera]
\label{ex:power-integral}
Niech $\alpha\in\mathbb R$ i
$f(x)=x^{-\alpha}$ na $(0,1]$.
Jeżeli $\alpha<1$, to funkcja jest całkowalna oraz
\begin{equation}\label{eq:power-integral}
\int_{(0,1]}x^{-\alpha}\,\mathrm d\lambda_1(x)
=\frac1{1-\alpha}.
\end{equation}
Jeśli $\alpha\ge1$, całka wynosi $+\infty$.
\end{przyklad}
\begin{proof}
Dla $\alpha\ge1$ mamy $x^{-\alpha}\ge x^{-1}$
na $(0,1]$, a rozbieżność drugiej całki
wykazaliśmy w rozdziale 7. Dla $\alpha<1$
rozważmy $f_n=f\chi_{[1/n,1]}$.
Wtedy $f_n\uparrow f$ na $(0,1]$.
Funkcja $f_n$ jest ograniczona i ciągła na
$[1/n,1]$, więc jej całka Lebesgue'a zgadza się
z całką Riemanna na tym przedziale (zob.
twierdzenie \ref{thm:riemann-to-lebesgue} poniżej).
Z rachunku elementarnego
$\int_{[1/n,1]}x^{-\alpha}\,\mathrm d\lambda_1(x)
=(1-n^{\alpha-1})/(1-\alpha)$.
Twierdzenie \ref{thm:mct} daje po przejściu do
granicy \eqref{eq:power-integral}.
\end{proof}

\begin{przyklad}[Kryterium logarytmiczne]
Dla $p\in\mathbb R$ rozważmy na $(0,e^{-1}]$
funkcję $f_p(x)=\bigl(x(\log(1/x))^p\bigr)^{-1}$.
Całka z $f_p$ jest skończona dokładnie wtedy, gdy
$p>1$. Rzeczywiście, podstawienie klasycznej
całki Riemanna $u=\log(1/x)$ na
$[e^{-R},e^{-1}]$ (a następnie twierdzenie
\ref{thm:mct}) daje
$\int_{[e^{-R},e^{-1}]}f_p\,\mathrm d\lambda_1(x)
=\int_1^R u^{-p}\,\mathrm d\lambda_1(u)$.
Całka po prawej ma skończoną granicę przy
$R\to\infty$ dokładnie wtedy, gdy $p>1$.
Samo użycie podstawienia w tym przykładzie
można też zastąpić oszacowaniami dyadycznymi;
pełne twierdzenie o zamianie zmiennych
udowodnimy później.
\end{przyklad}

\section{Całki zależne od parametru}

\begin{twierdzenie}[Ciągłość pod znakiem całki]
\label{thm:parameter-continuity}
Niech $T$ będzie przestrzenią metryczną, $t_0\in T$
i $F:T\times X\to\mathbb C$ funkcją taką, że
$x\mapsto F(t,x)$ jest mierzalna dla każdego $t\in T$.
Załóżmy, że istnieje otoczenie $V$ punktu $t_0$
oraz całkowalna $g:X\to[0,\infty]$ takie, że
$|F(t,x)|\le g(x)$ dla wszystkich $t\in V$ i
prawie wszystkich $x$, poza jednym wspólnym zbiorem
zerowym, oraz że $F(t,x)\to F(t_0,x)$, gdy
$t\to t_0$, dla prawie wszystkich $x$.
Wówczas funkcja
$\Phi(t)=\int_XF(t,x)\,\mathrm d\mu(x)$
jest określona dla $t\in V$ i ciągła w $t_0$.
\end{twierdzenie}
\begin{proof}
Ustalmy dowolny ciąg $t_n\to t_0$ w $T$.
Od pewnego indeksu $t_n\in V$. Funkcje
$F(t_n,\cdot)$ są zdominowane przez $g$ poza
wspólnym zbiorem zerowym i zbiegają prawie wszędzie
do $F(t_0,\cdot)$. Zespolona wersja twierdzenia
\ref{thm:dct} daje $\Phi(t_n)\to\Phi(t_0)$.
Ponieważ na przestrzeni metrycznej ciągowa
ciągłość jest równoważna ciągłości, otrzymujemy tezę.
\end{proof}

\begin{twierdzenie}[Różniczkowanie pod znakiem całki]
\label{thm:differentiating-integrals}
Niech $I\subseteq\mathbb R$ będzie przedziałem otwartym,
$t_0\in I$, a $F:I\times X\to\mathbb R$ funkcją,
dla której $x\mapsto F(t,x)$ jest mierzalna
dla każdego $t\in I$. Załóżmy, że:
\begin{enumerate}[label=\textup{(\roman*)}]
\item $F(t_0,\cdot)$ jest całkowalna;
\item istnieją $\delta>0$ o
$[t_0-\delta,t_0+\delta]\subset I$ i całkowalna
$g:X\to[0,\infty]$ takie, że dla prawie wszystkich
$x$ funkcja $t\mapsto F(t,x)$ jest różniczkowalna
na $(t_0-\delta,t_0+\delta)$ i
$|\partial_tF(t,x)|\le g(x)$ dla wszystkich $t$ z
tego przedziału (zbiór wyjątkowy jest wspólny).
\end{enumerate}
Wówczas $F(t,\cdot)$ jest całkowalna dla
$t$ dostatecznie bliskich $t_0$, funkcja
$\Phi(t)=\int_XF(t,x)\,\mathrm d\mu(x)$ jest
różniczkowalna w $t_0$ i
\begin{equation}\label{eq:differentiation-under-integral}
\Phi'(t_0)=\int_X\partial_tF(t_0,x)\,\mathrm d\mu(x).
\end{equation}
\end{twierdzenie}
\begin{proof}
Poza wspólnym zbiorem zerowym zastosujmy twierdzenie
o wartości średniej do funkcji $t\mapsto F(t,x)$.
Dla $|h|<\delta$ mamy
$|F(t_0+h,x)-F(t_0,x)|\le |h|g(x)$.
Zatem $|F(t_0+h,x)|\le|F(t_0,x)|+|h|g(x)$,
a prawa strona jest całkowalna. To dowodzi
istnienia całek w otoczeniu punktu $t_0$.
Funkcje mierzalne
$Q_h(x)=(F(t_0+h,x)-F(t_0,x))/h$ spełniają
$|Q_h(x)|\le g(x)$ dla wszystkich małych $h\ne0$
i zbiegają dla prawie wszystkich $x$ do
$\partial_tF(t_0,x)$. Granica punktowa ilorazów $Q_{1/n}$ jest mierzalna,
więc funkcja $x\mapsto\partial_tF(t_0,x)$,
po zdefiniowaniu jej jako zera na wspólnym zbiorze
wyjątkowym, jest mierzalna. Zbieżność zmajoryzowana,
zastosowana do dowolnego ciągu $h_n\to0$,
$h_n\ne0$, daje
$\lim_n\int_XQ_{h_n}\,\mathrm d\mu(x)
=\int_X\partial_tF(t_0,x)\,\mathrm d\mu(x)$.
Z liniowości całki lewa strona jest granicą
ilorazów różnicowych funkcji $\Phi$; niezależność
od wybranego ciągu dowodzi istnienia pochodnej.
\end{proof}

\begin{przyklad}[Transformata Laplace'a funkcji ograniczonej]
Niech $u:[0,\infty)\to\mathbb R$ będzie mierzalna,
$|u(x)|\le M<\infty$, i dla $t>0$ określmy
$L(t)=\int_{[0,\infty)}e^{-tx}u(x)\,\mathrm d\lambda_1(x)$.
Całka jest skończona, gdyż $|e^{-tx}u(x)|\le Me^{-tx}$,
a całka z $e^{-tx}$ wynosi $1/t$ (z granicy
całek Riemanna na $[0,R]$, używając Beppo Leviego).
Dla ustalonego $t_0>0$ i $t\in(t_0/2,3t_0/2)$
mamy
$|\partial_t(e^{-tx}u(x))|=xe^{-tx}|u(x)|
\le Mxe^{-t_0x/2}$.
Funkcja dominująca jest całkowalna, np. z oszacowania
$x e^{-t_0x/2}\le C e^{-t_0x/4}$, gdzie $C<\infty$.
Twierdzenie \ref{thm:differentiating-integrals} daje
$L'(t)=-\int_{[0,\infty)}xe^{-tx}u(x)\,\mathrm d\lambda_1(x)$.
Analogicznie, przez indukcję, otrzymamy
$L^{(k)}(t)=(-1)^k\int_{[0,\infty)}x^ke^{-tx}u(x)
\,\mathrm d\lambda_1(x)$ dla wszystkich $k\ge1$.
Istotnie, dla ustalonych $t_0>0$ i $k\ge0$
wszystkie potrzebne pochodne są lokalnie
zdominowane przez $Mx^{k+1}e^{-t_0x/2}$;
funkcja ta jest całkowalna, ponieważ dla
$d=t_0/4$ iloczyn $x^{k+1}e^{-dx}$ jest
ograniczony na $[0,\infty)$.
\end{przyklad}

\section{Jednostajna całkowalność i zbieżność według miary}

\begin{definicja}[Jednostajna całkowalność]
\label{def:uniform-integrability}
Niech $\mu(X)<\infty$. Niepustą rodzinę $\mathcal U$
mierzalnych, całkowalnych funkcji rzeczywistych
na $X$ nazywamy \emph{jednostajnie całkowalną},
jeśli
\[
\lim_{K\to\infty}\sup_{u\in\mathcal U}
\int_{\{|u|>K\}}|u(x)|\,\mathrm d\mu(x)=0.
\]
Warunek kontroluje masę części funkcji o dużych
wartościach, a nie jedynie wielkość miary poziomic.
\end{definicja}

\begin{lemat}[Kontrola całek na małych zbiorach]
\label{lem:uniform-absolute-continuity}
Jeśli $\mu(X)<\infty$ i $\mathcal U$ jest jednostajnie
całkowalną rodziną funkcji rzeczywistych, to:
\begin{enumerate}[label=\textup{(\roman*)}]
\item $\sup_{u\in\mathcal U}\int_X|u|\,\mathrm d\mu(x)<\infty$;
\item dla każdego $\varepsilon>0$ istnieje $\delta>0$
takie, że dla każdego $A\in\mathcal F$ z
$\mu(A)<\delta$ zachodzi
$\sup_{u\in\mathcal U}\int_A|u|\,\mathrm d\mu(x)<\varepsilon$.
\end{enumerate}
\end{lemat}
\begin{proof}
Ustalmy $\varepsilon>0$ i wybierzmy $K>0$ takie, że
$\sup_{u\in\mathcal U}\int_{\{|u|>K\}}|u|<\varepsilon/2$.
Dla każdego $u\in\mathcal U$ i $A\in\mathcal F$
z podziału na $\{|u|\le K\}$ i $\{|u|>K\}$
wynika
$\int_A|u|\,\mathrm d\mu(x)\le K\mu(A)+\varepsilon/2$.
Biorąc $A=X$, dowodzimy (i) przy skończonej
mierze $\mu(X)$. Wybór
$\delta=\varepsilon/(2K)$ dla $K>0$ daje (ii);
przy $K=0$ oszacowanie jest nawet prostsze.
\end{proof}

\begin{twierdzenie}[Dwa równoważne opisy jednostajnej całkowalności]
\label{thm:ui-two-criteria-8}
Niech $\mu(X)<\infty$ i niech $\mathcal U$ będzie
rodziną rzeczywistych funkcji mierzalnych $u:X\to\mathbb R$
spełniających $\int_X|u|\,\mathrm d\mu(x)<\infty$.
Równość funkcji prawie wszędzie nie ma wpływu
na żadną z rozważanych własności.
Wówczas $\mathcal U$ jest jednostajnie całkowalna
w sensie definicji~\ref{def:uniform-integrability}
wtedy i tylko wtedy, gdy jednocześnie:
\begin{enumerate}[label=\textup{(\roman*)}]
\item $M:=\sup_{u\in\mathcal U}\int_X|u|\,\mathrm d\mu(x)<\infty$;
\item dla każdego $\varepsilon>0$ istnieje $\delta>0$
takie, że $\mu(A)<\delta$ implikuje
$\sup_{u\in\mathcal U}\int_A|u|\,\mathrm d\mu(x)<\varepsilon$.
\end{enumerate}
\end{twierdzenie}
\begin{proof}
Konieczność to lemat~\ref{lem:uniform-absolute-continuity}.
Odwrotnie, przy $M=0$ wszystkie funkcje są zerowe
prawie wszędzie, więc teza jest natychmiastowa.
Dla $M>0$ i zadanego $\varepsilon>0$ wybieramy
$\delta$ z (ii). Nierówność Markowa daje
$\mu(\{|u|>K\})\le M/K$ dla każdego $u\in\mathcal U$.
Gdy $K>M/\delta$, zbiór $\{|u|>K\}$ ma miarę
mniejszą niż $\delta$, zatem (ii) daje
$\int_{\{|u|>K\}}|u|<\varepsilon$ jednostajnie
względem $u$. Jest to dokładnie warunek z
definicji~\ref{def:uniform-integrability}.
\end{proof}

\begin{twierdzenie}[Zbieżność Vitaliego na przestrzeni o mierze skończonej]
\label{thm:vitali-l1}
Niech $\mu(X)<\infty$, a $f_n:X\to\mathbb R$ będą
mierzalnymi funkcjami tworzącymi rodzinę
jednostajnie całkowalną. Jeśli
$f_n\xrightarrow{\mu}f$ dla pewnej mierzalnej
$f:X\to\mathbb R$, to $f$ jest całkowalna oraz
\[
\lim_{n\to\infty}\int_X|f_n(x)-f(x)|\,\mathrm d\mu(x)=0.
\]
W szczególności $\int_Xf_n\,\mathrm d\mu(x)\to
\int_Xf\,\mathrm d\mu(x)$.
\end{twierdzenie}
\begin{proof}
Z twierdzenia \ref{thm:measure-subsequence}
można wybrać podciąg $f_{n_k}\to f$ prawie wszędzie.
Na mocy lematu Fatou i części (i) lematu
\ref{lem:uniform-absolute-continuity}
$\int_X|f|\,\mathrm d\mu(x)\le
\liminf_k\int_X|f_{n_k}|\,\mathrm d\mu(x)<\infty$.

Ustalmy $\varepsilon>0$. Z lematu
\ref{lem:uniform-absolute-continuity} wybieramy
$\delta_1>0$ takie, że
$\mu(A)<\delta_1\Rightarrow
\sup_n\int_A|f_n|\,\mathrm d\mu(x)<\varepsilon/3$.
Ponieważ $f$ jest całkowalna, również jej całka
jest absolutnie ciągła względem miary w następującym
sensie: istnieje $\delta_2>0$ takie, że
$\mu(A)<\delta_2\Rightarrow\int_A|f|<\varepsilon/3$.
Istotnie, wystarczy obciąć $|f|$ na poziomie $K$,
tak by $\int_{\{|f|>K\}}|f|<\varepsilon/6$,
i oszacować pozostałą część przez $K\mu(A)$.

Niech $\eta>0$ będzie tak małe, aby
$\eta\mu(X)<\varepsilon/3$. Ze zbieżności według
miary wynika, że dla dużych $n$ zbiór
$A_n=\{|f_n-f|>\eta\}$ ma miarę mniejszą od
$\min(\delta_1,\delta_2)$. Dzieląc całkę na
$A_n$ i dopełnienie, dostajemy
\[
\int_X|f_n-f|\,\mathrm d\mu(x)
\le\eta\mu(X)+\int_{A_n}|f_n|\,\mathrm d\mu(x)
+\int_{A_n}|f|\,\mathrm d\mu(x)<\varepsilon.
\]
To dowodzi zbieżności w całce z modułu różnicy.
Równość granic całek wynika z oszacowania
$|\int_X f_n-\int_X f|\le\int_X|f_n-f|$.
\end{proof}

\begin{przyklad}[Masę można skoncentrować na małych zbiorach]
Na $(0,1)$ niech $f_n=n\chi_{(0,1/n)}$.
Ciąg zbiega do zera według miary, ponieważ
$\lambda_1(\{|f_n|>\varepsilon\})\le1/n$
dla dużych $n$. Rodzina nie jest jednak jednostajnie
całkowalna: dla każdego $K>0$ i każdego $n>K$
zachodzi
$\int_{\{|f_n|>K\}}|f_n|\,\mathrm d\lambda_1(x)=1$.
Twierdzenie Vitaliego zatem nie ma zastosowania.
\end{przyklad}

\section{Związek z całką Riemanna}

\begin{twierdzenie}[Zgodność całki Riemanna i Lebesgue'a]
\label{thm:riemann-to-lebesgue}
Niech $a<b$ i $f:[a,b]\to\mathbb R$ będzie funkcją
ograniczoną, całkowalną w sensie Riemanna.
Wówczas $f$ jest mierzalna względem uzupełnionej
miary Lebesgue'a na $[a,b]$, jest całkowalna
w sensie Lebesgue'a, a obie całki są równe:
\[
\int_a^b f(x)\,\mathrm dx
=\int_{[a,b]}f(x)\,\mathrm d\lambda_1(x).
\]
Po lewej stronie występuje tradycyjna całka Riemanna,
a po prawej całka względem miary Lebesgue'a.
\end{twierdzenie}
\begin{proof}
Niech $M$ będzie ograniczeniem $|f|$.
Z kryterium Darboux dla całkowalności Riemanna
możemy wybrać podziały $P_n$ takie, że
$U(f,P_n)-L(f,P_n)<2^{-2n}$,
gdzie $U,L$ są odpowiednio górnymi i dolnymi
sumami Darboux. Z każdego podziału budujemy
proste funkcje borelowskie $u_n,\ell_n$:
na rozłącznych półotwartych częściach podziału
przypisujemy odpowiednio supremum i infimum $f$
na odpowiadających im \emph{domkniętych} przedziałach
podziału (ostatni przedział zawiera prawy koniec).
Dzięki temu $\ell_n\le f\le u_n$ w każdym punkcie,
a $\int_{[a,b]}u_n(x)\,\mathrm d\lambda_1(x)=U(f,P_n)$,
$\int_{[a,b]}\ell_n(x)\,\mathrm d\lambda_1(x)=L(f,P_n)$.

Ponieważ $u_n-\ell_n\ge0$ i
$\int_{[a,b]}(u_n-\ell_n)\,\mathrm d\lambda_1(x)
<2^{-2n}$, nierówność Markowa daje
$\lambda_1(\{u_n-\ell_n\ge2^{-n}\})\le2^{-n}$.
Pierwsza część lematu Borela--Cantellego implikuje,
że dla prawie wszystkich $x$ istnieje $N(x)$ takie,
że $u_n(x)-\ell_n(x)<2^{-n}$ dla $n\ge N(x)$.
Z nierówności $\ell_n\le f\le u_n$ otrzymujemy
$u_n(x)\to f(x)$ poza mierzalnym zbiorem zerowym.
Granica mierzalnych funkcji jest mierzalna poza
tym zbiorem; dzięki zupełności miary Lebesgue'a
sama $f$ jest mierzalna. Ponadto $|f|\le M$
i $\lambda_1([a,b])<\infty$, więc $f$ jest całkowalna.

Z monotoniczności całki i równości dla prostych
funkcji otrzymujemy
$L(f,P_n)\le\int_{[a,b]}f\,\mathrm d\lambda_1(x)
\le U(f,P_n)$. Obie skrajne liczby dążą do
całki Riemanna, bo ich różnica dąży do zera,
a całka Riemanna leży pomiędzy nimi.
Twierdzenie o trzech ciągach kończy dowód.
\end{proof}

\begin{twierdzenie}[Kryterium Lebesgue'a całkowalności Riemanna]
\label{thm:riemann-criterion}
Niech $f:[a,b]\to\mathbb R$ będzie ograniczona, gdzie
$a<b$. Wówczas $f$ jest całkowalna w sensie Riemanna
tedy i tylko wtedy, gdy zbiór punktów nieciągłości $f$
ma miarę Lebesgue'a zero.
\end{twierdzenie}
\begin{proof}
Dla $x\in[a,b]$ oznaczmy przez $\omega_f(x)$
oscylację $f$ w punkcie, czyli
\[
\omega_f(x):=\inf_{r>0}
\left(\sup_{y\in[a,b]\cap(x-r,x+r)}f(y)
-\inf_{y\in[a,b]\cap(x-r,x+r)}f(y)\right).
\]
Funkcja $f$ jest ciągła w $x$ wtedy i tylko wtedy,
gdy $\omega_f(x)=0$. Dla $\varepsilon>0$ zbiór
$D_\varepsilon=\{x:\omega_f(x)\ge\varepsilon\}$
jest domknięty: jeśli $\omega_f(x)<\varepsilon$,
to na pewnym otoczeniu punktu $x$ oscylacja $f$
jest mniejsza niż $\varepsilon$, a zatem także
oscylacje w punktach pobliskich są mniejsze od
$\varepsilon$. Zatem $D_\varepsilon$ jest zwarty.

Załóżmy najpierw całkowalność Riemanna i
ustalmy $\varepsilon>0$. Weźmy dowolny podział
$P=(a=x_0<\cdots<x_N=b)$. Jeśli
$x\in D_\varepsilon$ nie jest punktem podziału,
to oscylacja $f$ na domkniętym przedziale
$[x_{j-1},x_j]$ zawierającym $x$ jest co najmniej
$\varepsilon$. Dlatego $D_\varepsilon$ mieści się
w sumie skończenie wielu punktów podziału i tych
przedziałów, których oscylacja jest nie mniejsza
niż $\varepsilon$. Suma długości tych przedziałów
jest nie większa niż
$\varepsilon^{-1}(U(f,P)-L(f,P))$.
Ponieważ dla funkcji całkowalnej w sensie Riemanna
różnicę tę można uczynić dowolnie małą,
$\lambda_1(D_\varepsilon)=0$.
Zbiór nieciągłości jest sumą
$\bigcup_{m\ge1}D_{1/m}$, więc jest zerowy.

Odwrotnie, niech zbiór nieciągłości $D$ ma miarę
zero, a $|f|\le M$. Ustalmy $\varepsilon>0$.
Zbiór $D_\varepsilon$ jest zwarty i zerowy, więc
istnieje zbiór otwarty $O$ zawierający $D_\varepsilon$
o dowolnie małej mierze Lebesgue'a (regularność
zewnętrzna, rozdział 4). Zbiór
$K=[a,b]\setminus O$ jest zwarty. Dla każdego
$x\in K$ mamy $\omega_f(x)<\varepsilon$,
zatem istnieje promień $r_x>0$, na którego kuli
$B(x,r_x)\cap[a,b]$ oscylacja $f$ jest
mniejsza niż $\varepsilon$. Zwartość $K$ pozwala
wybrać skończenie wiele kul $B(x_j,r_{x_j}/2)$
pokrywających $K$. Połóżmy
$\delta=\frac12\min_j r_{x_j}>0$.
Jeżeli domknięty przedział $J\subseteq[a,b]$
o długości mniejszej niż $\delta$ przecina $K$,
wybieramy $y\in J\cap K$ i indeks $j$ taki, że
$y\in B(x_j,r_{x_j}/2)$. Dla każdego $z\in J$
mamy
$|z-x_j|\le|z-y|+|y-x_j|<\delta+r_{x_j}/2
\le r_{x_j}$.
Zatem $J\subseteq B(x_j,r_{x_j})$, gdzie
oscylacja $f$ jest mniejsza niż $\varepsilon$.

Weźmy podział $P$ o siatce mniejszej niż $\delta$.
Na przedziałach przecinających $K$ oscylacja $f$
jest mniejsza niż $\varepsilon$. Pozostałe
przedziały (z wyjątkiem ewentualnych końców)
mieszczą się w $O$, więc suma ich długości
nie przekracza $\lambda_1(O)$. Otrzymujemy
\[
U(f,P)-L(f,P)
\le\varepsilon(b-a)+2M\lambda_1(O).
\]
Wybierając najpierw $\varepsilon$, a potem
$\lambda_1(O)$ dostatecznie małe, czynimy prawą
stronę dowolnie małą. Kryterium Darboux dowodzi
całkowalności Riemanna.
\end{proof}

\begin{uwaga}[Dlaczego oba warunki mają znaczenie?]
Funkcja Dirichleta $\chi_{\mathbb Q\cap[0,1]}$
jest całkowalna w sensie Lebesgue'a, ale
nie jest całkowalna w sensie Riemanna, gdyż
zbiór jej punktów nieciągłości to cały $[0,1]$.
Z kolei każda ograniczona funkcja o przeliczalnym
zbiorze punktów nieciągłości jest całkowalna
w sensie Riemanna. Warunek ograniczoności jest
w kryterium konieczny w jego klasycznej postaci:
na przykład $x\mapsto1/x$ na $(0,1]$ nie jest
ograniczona, mimo ciągłości w całej swojej dziedzinie.
\end{uwaga}

\section{Podsumowanie logiczne}

Twierdzenie Beppo Leviego opiera się na definicji
całki jako supremum po funkcjach prostych i na
ciągłości miary z dołu. Z niego wynikają
addytywność całki nieujemnej i możliwość całkowania
szeregów nieujemnych. Lemat Fatou to zastosowanie
Beppo Leviego do ciągu infimów ogonowych, a twierdzenie
o zbieżności zmajoryzowanej uzyskujemy z lematu Fatou
przez rozważenie nieujemnych funkcji $g\pm f_n$.
Wszystkie te rezultaty pozostają ważne dla miar
nieskończonych; wyjątkiem jest twierdzenie o zbieżności
ograniczonej, w którym skończoność miary
zapewnia całkowalność stałej funkcji dominującej.

% END chapters/pl/measure_theory/08-rozdzial.tex


% BEGIN chapters/pl/measure_theory/09-rozdzial.tex
% ROZDZIAŁ 9 — samodzielny fragment do włączenia w main.tex
\chapter{Rodzaje zbieżności funkcji mierzalnych}
\label{ch:rodzaje-zbieznosci}

\section{Dlaczego potrzebujemy kilku pojęć zbieżności?}

Zbieżność punktowa opisuje zachowanie funkcji w każdym ustalonym
punkcie, lecz nie kontroluje zbiorów, na których funkcje przyjmują
duże wartości. Zbieżność według miary kontroluje rozmiar takich
zbiorów, ale dopuszcza ich przemieszczanie. Zbieżność w całce
kontroluje dodatkowo całkowitą wielkość błędu. To właśnie różnice
między tymi warunkami wyjaśniają, dlaczego nie można dowolnie
zamieniać kolejności granicy i całkowania.

Przez cały rozdział $(X,\mathcal F,\mu)$ jest przestrzenią
z miarą; nie wymagamy zupełności ani $\sigma$-skończoności, chyba
że zaznaczono inaczej. Wszystkie funkcje $f_n,f$ są mierzalne i
rzeczywiste, o ile nie wprowadzono innych założeń. Mówimy, że
własność zachodzi \emph{$\mu$-prawie wszędzie}, gdy zbiór punktów,
w których może zawodzić, zawiera się w pewnym zbiorze należącym
do $\mathcal F$ o mierze zero. Jest to ważne przy miarach
niezupełnych.

\section{Definicje i pierwsze relacje}

Zbieżność według miary zdefiniowaliśmy już w
Definicji~\ref{def:convergence-measure}, a twierdzenie o podciągu
udowodniliśmy w Twierdzeniu~\ref{thm:measure-subsequence}.
Teraz porównujemy ją z pozostałymi pojęciami zbieżności.


\begin{definicja}[Pięć rodzajów zbieżności]
\label{def:five-convergences}
Niech $f_n,f:X\to\mathbb R$ będą mierzalne.
\begin{enumerate}[label=\textup{(\roman*)}]
\item $f_n\to f$ \emph{punktowo}, jeśli
$f_n(x)\to f(x)$ dla każdego $x\in X$.
\item $f_n\to f$ \emph{$\mu$-prawie wszędzie}, w skrócie
$f_n\xrightarrow{\mathrm{p.w.}}f$, jeśli zbieżność
punktowa zachodzi poza pewnym mierzalnym zbiorem zerowym.
\item $f_n\to f$ \emph{jednostajnie}, jeśli dla każdego
$\varepsilon>0$ istnieje $N$, takie że
$|f_n(x)-f(x)|<\varepsilon$ dla wszystkich $n\ge N$
i $x\in X$.
\item $f_n\to f$ \emph{prawie jednostajnie}, jeśli dla każdego
$\eta>0$ istnieje $E\in\mathcal F$ z $\mu(E)<\eta$
takie, że $f_n\to f$ jednostajnie na $X\setminus E$.
\item $f_n\to f$ \emph{według miary}, co oznaczamy
$f_n\xrightarrow{\mu}f$, jeśli dla każdego $\varepsilon>0$
zachodzi
\begin{equation}\label{eq:conv-in-measure-9}
\lim_{n\to\infty}\mu\bigl(\{x\in X:
|f_n(x)-f(x)|>\varepsilon\}\bigr)=0.
\end{equation}
\end{enumerate}
Nie wymagamy, aby zbiory poziomicowe w
\eqref{eq:conv-in-measure-9} miały skończoną miarę dla
poszczególnych $n$; wymagana jest ich granica równa zeru.
\end{definicja}

\begin{uwaga}[Dwie kwantyfikacje, których nie wolno mylić]
W zbieżności jednostajnej kolejność kwantyfikatorów jest
następująca: dla każdego $\varepsilon>0$ istnieje jedno $N$
działające \emph{jednocześnie} dla wszystkich $x$. W zbieżności
punktowej indeks $N$ może zależeć od $x$. Zbieżność prawie
jednostajna pozwala usunąć zbiór o dowolnie małej mierze,
ale zbiór ten zależy od wymaganego dopuszczalnego błędu w mierze.
Zbiór wyjątkowy nie może natomiast zależeć od indeksu $n$
w definicji jednostajności na jego dopełnieniu.
\end{uwaga}

\begin{twierdzenie}[Podstawowe implikacje]
\label{thm:conv-implications-9}
Dla dowolnej przestrzeni z miarą $(X,\mathcal F,\mu)$
i mierzalnych funkcji $f_n,f:X\to\mathbb R$ zachodzą:
\begin{enumerate}[label=\textup{(\roman*)}]
\item zbieżność jednostajna pociąga zbieżność punktową
i według miary;
\item zbieżność prawie jednostajna pociąga zbieżność
prawie wszędzie i według miary;
\item zbieżność punktowa pociąga zbieżność prawie wszędzie.
\end{enumerate}
Jeżeli dodatkowo $\mu(X)<\infty$, to zbieżność prawie
wszędzie pociąga zbieżność według miary.
\end{twierdzenie}
\begin{proof}
W (i) dla ustalonego $\varepsilon>0$ i wszystkich dużych
$n$ zbiór $\{|f_n-f|>\varepsilon\}$ jest pusty.
Implikacja punktowa jest oczywista, podobnie (iii).

Dla (ii) ustalmy $\varepsilon,\eta>0$ i wybierzmy
$E\in\mathcal F$ o $\mu(E)<\eta$, poza którym zbieżność
jest jednostajna. Dla dużych $n$ zbiór
$\{|f_n-f|>\varepsilon\}$ zawiera się w $E$,
zatem jego miara jest nie większa od $\eta$.
Ponieważ $\eta$ jest dowolne, otrzymujemy
\eqref{eq:conv-in-measure-9}.

Niech teraz $N$ będzie zbiorem, na którym zbieżność
punktowa nie zachodzi. Ponieważ
\[
N=\bigcup_{m\ge1}\bigcap_{N_0\ge1}
\bigcup_{n\ge N_0}\{|f_n-f|\ge 1/m\},
\]
zbiór $N$ jest mierzalny. Jeśli zbieżność jest prawie
jednostajna, to dla każdego $\eta>0$ istnieje zbiór
$E$ o mierze mniejszej niż $\eta$, poza którym zbieżność
zachodzi; stąd $N\subseteq E$ i $\mu(N)\le\eta$.
Wynika $\mu(N)=0$, co dowodzi drugiej części (ii).

Ostatnia implikacja jest dokładnie twierdzeniem~\ref{thm:ae-to-measure-finite}, już udowodnionym w rozdziale~\ref{ch:funkcje-mierzalne}. To kończy dowód.
\end{proof}

\begin{przyklad}[Dlaczego w ostatniej implikacji potrzebna jest skończoność miary]
\label{ex:ae-not-measure-infinite-9}
Na $\mathbb R$ z miarą Lebesgue'a niech
$f_n=\chi_{[n,\infty)}$. Dla każdego ustalonego
$x\in\mathbb R$ mamy $f_n(x)\to0$, natomiast
dla $0<\varepsilon<1$ zachodzi
$\lambda_1(\{|f_n|>\varepsilon\})=\infty$ dla
każdego $n$. Nie wystarcza nawet zbieżność punktowa
na całej przestrzeni; problemem jest brak kontroli
masy poza zbiorami skończonej miary.
\end{przyklad}

\section{Twierdzenie Jegorowa}

\begin{twierdzenie}[Jegorow]
\label{thm:egorov-9}
Niech $(X,\mathcal F,\mu)$ będzie przestrzenią z miarą
skończoną, tj. $\mu(X)<\infty$. Jeśli
$f_n,f:X\to\mathbb R$ są mierzalne oraz
$f_n\to f$ $\mu$-prawie wszędzie, to
$f_n\to f$ prawie jednostajnie. Dokładniej: dla każdego
$\eta>0$ istnieje $E\in\mathcal F$, $\mu(E)<\eta$,
takie, że
$\sup_{x\in X\setminus E}|f_n(x)-f(x)|\to0$.
\end{twierdzenie}
\begin{proof}
Niech $N\in\mathcal F$ będzie zbiorem miary zero,
poza którym $f_n(x)\to f(x)$. Dla $m,N_0\in\mathbb N$
określmy
\[
A_{m,N_0}=\bigcup_{n\ge N_0}
\{x\in X:|f_n(x)-f(x)|\ge1/m\}.
\]
Zbiory $A_{m,N_0}$ są mierzalne i dla ustalonego $m$
maleją wraz z $N_0$. Zbieżność poza $N$ oznacza,
że $\bigcap_{N_0\ge1}A_{m,N_0}\subseteq N$.
Zatem miara tego przecięcia wynosi zero. Skończoność
miary $X$ pozwala zastosować ciągłość z góry:
$\mu(A_{m,N_0})\to0$ dla $N_0\to\infty$.

Ustalmy $\eta>0$. Dla każdego $m$ możemy wybrać
$N_m$ tak duże, by
$\mu(A_{m,N_m})<\eta 2^{-(m+1)}$.
Połóżmy $E=\bigcup_{m\ge1}A_{m,N_m}$.
Przeliczalna podaddytywność daje
$\mu(E)<\eta$.
Dla $x\notin E$ i każdego $m$ oraz $n\ge N_m$
mamy $|f_n(x)-f(x)|<1/m$. Jeśli teraz
$\varepsilon>0$, wybieramy $m$ z $1/m<\varepsilon$;
wówczas nierówność z indeksem $N_m$ zachodzi
jednocześnie dla wszystkich $x\in X\setminus E$.
To jest dokładnie zbieżność jednostajna na
$X\setminus E$.
\end{proof}

\begin{wniosek}[Wersja lokalna dla przestrzeni $\sigma$-skończonej]
\label{cor:local-egorov-9}
Niech $X=\bigcup_{j\ge1}X_j$, gdzie $X_j\in\mathcal F$,
$X_j\subseteq X_{j+1}$ i $\mu(X_j)<\infty$.
Jeżeli $f_n\to f$ prawie wszędzie, to dla każdego
$\eta>0$ istnieje $E\in\mathcal F$ o $\mu(E)<\eta$
takie, że $f_n\to f$ jednostajnie na
$X_j\setminus E$ dla każdego ustalonego $j$.
Nie twierdzimy, że zbieżność jest jednostajna
na całym $X\setminus E$.
\end{wniosek}
\begin{proof}
Dla każdego $j$ stosujemy twierdzenie
\ref{thm:egorov-9} do ograniczenia miary i funkcji
na $X_j$, z tolerancją $\eta 2^{-(j+1)}$.
Otrzymujemy $E_j\subseteq X_j$ takie, że
$\mu(E_j)<\eta 2^{-(j+1)}$ i zbieżność jest
jednostajna na $X_j\setminus E_j$.
Niech $E=\bigcup_jE_j$. Wtedy $\mu(E)<\eta$, a
$X_j\setminus E\subseteq X_j\setminus E_j$,
więc zbieżność jest jednostajna na pierwszym z tych
zbiorów.
\end{proof}

\begin{przyklad}[Nie istnieje globalna wersja Jegorowa na $\mathbb R$]
Niech $f_n=\chi_{[n,n+1)}$ na $\mathbb R$.
Funkcje zbiegają do zera wszędzie, ale jeśli
$\lambda_1(E)<\infty$, to nie może zachodzić
$f_n\to0$ jednostajnie na $\mathbb R\setminus E$.
W przeciwnym razie istniałoby $N$ takie, że
$[n,n+1)\subseteq E$ dla wszystkich $n\ge N$
(gdyż $f_n$ przyjmuje wyłącznie wartości $0,1$).
Wynikałoby $\lambda_1(E)=\infty$, sprzeczność.
\end{przyklad}

\section{Zbieżność według miary i twierdzenie Riesza o podciągu}

Jednoznaczność granicy według miary, z dokładnością do zbioru zerowego,
została wykazana w Wniosku~\ref{cor:measure-limit-unique}.
Ponadto twierdzenie Riesza o podciągu jest już w pełni
udowodnione w Twierdzeniu~\ref{thm:measure-subsequence};
w dalszych wynikach korzystamy właśnie z tej wcześniejszej wersji.

\begin{wniosek}[Kryterium za pomocą wszystkich podciągów]
\label{cor:every-subseq-9}
Niech $\mu(X)<\infty$. Dla mierzalnych
$f_n,f:X\to\mathbb R$ następujące
warunki są równoważne:
\begin{enumerate}[label=\textup{(\roman*)}]
\item $f_n\xrightarrow{\mu}f$;
\item każdy podciąg $(f_{n_k})$ ma dalszy podciąg,
który zbiega do $f$ $\mu$-prawie wszędzie.
\end{enumerate}
\end{wniosek}
\begin{proof}
Jeżeli zachodzi (i), każdy podciąg także zbiega
według miary, więc stosujemy twierdzenie
\ref{thm:measure-subsequence}.
Załóżmy (ii) i przypuśćmy, że (i) nie zachodzi.
Wówczas istnieją $\varepsilon,\delta>0$ oraz
podciąg $(f_{n_k})$ taki, że
$\mu(\{|f_{n_k}-f|>\varepsilon\})\ge\delta$
dla każdego $k$. Na mocy (ii) ma on dalszy
podciąg zbieżny do $f$ prawie wszędzie.
Ponieważ $\mu(X)<\infty$, twierdzenie
\ref{thm:conv-implications-9} daje jego zbieżność
według miary, wbrew ustalonemu oszacowaniu.
\end{proof}

\begin{uwaga}[Nie można skreślić założenia $\mu(X)<\infty$]
Na przestrzeni o mierze nieskończonej warunek (ii)
nie wystarcza: ciąg $\chi_{[n,\infty)}$ na
$\mathbb R$ zbiega wszędzie do zera, lecz nie
według miary. Właśnie dlatego równoważność
\ref{cor:every-subseq-9} została sformułowana
dla miar skończonych.
\end{uwaga}

\section{Zbieżność w $L^p$ i oszacowania Czebyszewa}

\begin{definicja}[Zbieżność w średniej rzędu $p$]
\label{def:lp-convergence-9}
Niech $1\le p<\infty$. Mierzalną funkcję
$h:X\to\mathbb R$ nazywamy \emph{$p$-całkowalną},
gdy $\int_X|h|^p\,\mathrm d\mu(x)<\infty$.
Dla $p$-całkowalnych funkcji $f_n,f:X\to\mathbb R$
mówimy, że $f_n\to f$ \emph{w $L^p(\mu)$}, jeśli
$\int_X|f_n(x)-f(x)|^p\,\mathrm d\mu(x)\to0$.
Pełną konstrukcję przestrzeni normowanej $L^p$
na klasach funkcji równych prawie wszędzie,
wraz z nierównościami H\"oldera i Minkowskiego,
przedstawimy w rozdziale 14.
\end{definicja}

\begin{twierdzenie}[Nierówność Czebyszewa i zbieżność według miary]
\label{thm:lp-implies-measure-9}
Niech $1\le p<\infty$, a $f_n,f:X\to\mathbb R$
będą mierzalne. Dla każdego $\varepsilon>0$ zachodzi
\begin{equation}\label{eq:chebyshev-lp-9}
\mu\bigl(\{|f_n-f|>\varepsilon\}\bigr)
\le\varepsilon^{-p}
\int_X|f_n(x)-f(x)|^p\,\mathrm d\mu(x).
\end{equation}
W szczególności zbieżność w $L^p(\mu)$ pociąga
zbieżność według miary.
\end{twierdzenie}
\begin{proof}
Stosujemy nierówność Markowa z rozdziału 7
do nieujemnej mierzalnej funkcji $|f_n-f|^p$
i poziomu $\varepsilon^p$. Jeśli jej całka dąży
do zera, prawa strona \eqref{eq:chebyshev-lp-9}
dąży do zera dla każdego ustalonego $\varepsilon>0$.
\end{proof}

\begin{przyklad}[Zbieżność w $L^p$ nie wymusza zbieżności prawie wszędzie]
\label{ex:typewriter-lp-9}
Dla $k\ge0$ i $0\le j<2^k$ niech
$g_{2^k+j}=\chi_{[j2^{-k},(j+1)2^{-k})}$
na $[0,1)$ z miarą Lebesgue'a. Dla każdego
$1\le p<\infty$ mamy
$\int_{[0,1)}|g_{2^k+j}|^p\,\mathrm d\lambda_1(x)=2^{-k}$,
więc $g_n\to0$ w $L^p$, a także według miary.
Jednak dla każdego $x\in[0,1)$ w każdym bloku o długości
$2^k$ dokładnie jedna funkcja przyjmuje w $x$
wartość $1$, a pozostałe wartość $0$ (dla $k\ge1$).
Stąd $g_n(x)$ nie ma granicy w żadnym punkcie $[0,1)$.
Jest to klasyczna \emph{maszyna do pisania}.
\end{przyklad}

\begin{przyklad}[Zbieżność prawie wszędzie i według miary bez zbieżności w $L^1$]
Niech $f_n=n\chi_{(0,1/n)}$ na $(0,1)$.
Wtedy $f_n\to0$ wszędzie i według miary, lecz
$\int_{(0,1)}|f_n|\,\mathrm d\lambda_1(x)=1$
dla każdego $n$. Brakuje jednostajnej całkowalności;
por. twierdzenie Vitaliego z rozdziału 8.
\end{przyklad}

\section{Ciągi Cauchy'ego według miary i zupełność}

\begin{definicja}[Ciąg Cauchy'ego według miary]
\label{def:cauchy-measure-9}
Ciąg mierzalnych funkcji rzeczywistych $(f_n)$
nazywamy \emph{ciągiem Cauchy'ego według miary}, jeśli
dla wszystkich $\varepsilon,\delta>0$ istnieje
$N\in\mathbb N$ takie, że dla $m,n\ge N$
zachodzi $\mu(\{|f_n-f_m|>\varepsilon\})<\delta$.
\end{definicja}

\begin{twierdzenie}[Zupełność zbieżności według miary]
\label{thm:cauchy-measure-complete-9}
Niech $(X,\mathcal F,\mu)$ będzie dowolną przestrzenią
z miarą. Jeżeli $(f_n)$ jest ciągiem Cauchy'ego według
miary złożonym z mierzalnych funkcji rzeczywistych,
to istnieje mierzalna funkcja $f:X\to\mathbb R$ taka,
że $f_n\xrightarrow{\mu}f$. Funkcja $f$ jest wyznaczona
jednoznacznie z dokładnością do równości prawie wszędzie.
\end{twierdzenie}
\begin{proof}
Wybierzmy ściśle rosnący ciąg indeksów $n_k$
rekurencyjnie tak, aby dla wszystkich $m,n\ge n_k$
zachodziło
$\mu(\{|f_m-f_n|>2^{-k}\})<2^{-k}$.
Możemy jednocześnie wymagać $n_{k+1}>n_k$.
Niech
$A_k=\{|f_{n_{k+1}}-f_{n_k}|>2^{-k}\}$.
Wówczas $\sum_k\mu(A_k)<\infty$ i lemat
Borela--Cantellego zapewnia, że poza mierzalnym
zbiorem zerowym $N$ każdy punkt należy tylko do
skończenie wielu $A_k$.

Dla $x\notin N$ i dużych $k$ mamy
$|f_{n_{k+1}}(x)-f_{n_k}(x)|\le2^{-k}$.
Dla $j>k$ nierówność trójkąta daje
\[
|f_{n_j}(x)-f_{n_k}(x)|
\le\sum_{\ell=k}^{j-1}2^{-\ell}\le2^{1-k}.
\]
Zatem $(f_{n_k}(x))$ jest ciągiem Cauchy'ego w
$\mathbb R$ i zbiega do pewnej skończonej liczby
$f(x)$. Na $N$ definiujemy $f(x)=0$.
Mierzalność $f$ wynika z twierdzenia o mierzalności
granic punktowych po przedłużeniu funkcji granicznej
zerem na $N$. Zbieżność $f_{n_k}\to f$ prawie
wszędzie nie gwarantuje sama przez się zbieżności
według miary na przestrzeni nieskończonej.

Tutaj jednak mamy silniejsze oszacowanie: dla
ustalonego $k$ i dowolnego $j>k$, poza zbiorem
$\bigcup_{\ell\ge k}A_\ell$, zachodzi
$|f_{n_j}-f_{n_k}|\le2^{1-k}$.
Przechodząc do granicy $j\to\infty$ poza $N$,
dostajemy
\[
\mu\bigl(\{|f-f_{n_k}|>2^{1-k}\}\bigr)
\le\sum_{\ell\ge k}\mu(A_\ell)
\le\sum_{\ell\ge k}2^{-\ell}\longrightarrow0.
\]
Zatem $f_{n_k}\to f$ według miary.
Dla ustalonych $\varepsilon,\delta>0$ warunek
Cauchy'ego daje $N_C$ takie, że dla $m,n\ge N_C$
zachodzi
$\mu(\{|f_n-f_m|>\varepsilon/2\})<\delta/2$.
Wybieramy teraz $k$ \emph{jednocześnie} tak duże,
aby $n_k\ge N_C$ i
$\mu(\{|f_{n_k}-f|>\varepsilon/2\})<\delta/2$.
Wówczas dla każdego $n\ge N_C$ mamy
$\mu(\{|f_n-f_{n_k}|>\varepsilon/2\})<\delta/2$.
Wówczas $\mu(\{|f_n-f|>\varepsilon\})<\delta$.
Jednoznaczność wynika z twierdzenia
\ref{cor:measure-limit-unique}.
\end{proof}

\section{Metryzacja zbieżności według miary na przestrzeni skończonej miary}

\begin{twierdzenie}[Metryka zbieżności według miary]
\label{thm:measure-metric-9}
Niech $\mu(X)<\infty$. Oznaczmy przez
$\mathcal M(X,\mathcal F)$ zbiór mierzalnych funkcji
rzeczywistych na $X$, a przez $\mathcal M/\!\sim$
jego iloraz przez relację
$f\sim g\Longleftrightarrow f=g$ $\mu$-prawie wszędzie.
Dla klas $[f],[g]\in\mathcal M/\!\sim$ połóżmy
\begin{equation}\label{eq:measure-metric-9}
\rho([f],[g])=
\int_X\min\{1,|f(x)-g(x)|\}\,\mathrm d\mu(x).
\end{equation}
Wówczas $\rho$ jest metryką na $\mathcal M/\!\sim$.
Ponadto $f_n\xrightarrow{\mu}f$ wtedy i tylko wtedy,
gdy $\rho([f_n],[f])\to0$.
\end{twierdzenie}
\begin{proof}
Całka w \eqref{eq:measure-metric-9} jest skończona,
bo całkujemy funkcję ograniczoną przez $1$,
a $\mu(X)<\infty$. Równość funkcji prawie
wszędzie nie zmienia wartości całki, więc wzór
zależy tylko od klas. Symetria i nieujemność
są oczywiste. Jeśli $\rho([f],[g])=0$, to
nieujemna funkcja $\min\{1,|f-g|\}$ ma całkę
zero, a więc znika prawie wszędzie; wynika
$[f]=[g]$. Nierówność trójkąta jest konsekwencją
punktowej nierówności
\[
\min\{1,|f-h|\}\le
\min\{1,|f-g|\}+\min\{1,|g-h|\},
\]
którą otrzymujemy z nierówności trójkąta dla
wartości bezwzględnej i oszacowania
$\min\{1,a+b\}\le\min\{1,a\}+\min\{1,b\}$
dla $a,b\ge0$. Całkujemy obie strony.

Załóżmy $\rho([f_n],[f])\to0$. Dla $\varepsilon>0$
zachodzi punktowo
$\min\{1,|f_n-f|\}\ge
\min\{1,\varepsilon\}\chi_{\{|f_n-f|>\varepsilon\}}$.
Zatem
\[
\mu(\{|f_n-f|>\varepsilon\})\le
\frac{\rho([f_n],[f])}{\min\{1,\varepsilon\}}
\longrightarrow0.
\]

Odwrotnie, jeśli $f_n\to f$ według miary,
to dla każdego $0<\varepsilon<1$ mamy
\[
\rho([f_n],[f])\le\varepsilon\mu(X)+
\mu(\{|f_n-f|>\varepsilon\}).
\]
Dla ustalonego $\varepsilon$ drugi składnik
dąży do zera, zatem
$\limsup_n\rho([f_n],[f])\le\varepsilon\mu(X)$.
Pozwalając $\varepsilon\downarrow0$, otrzymujemy
$\rho([f_n],[f])\to0$.
\end{proof}

\begin{wniosek}[Zupełność metryki zbieżności według miary]
\label{cor:measure-metric-complete-9}
Przy założeniu $\mu(X)<\infty$ przestrzeń
$(\mathcal M/\!\sim,\rho)$ z twierdzenia
\ref{thm:measure-metric-9} jest zupełna.
\end{wniosek}
\begin{proof}
Niech $([f_n])$ będzie ciągiem Cauchy'ego w metryce
$\rho$. Dla $\varepsilon>0$ mamy
$\mu(\{|f_n-f_m|>\varepsilon\})\le
\rho([f_n],[f_m])/\min\{1,\varepsilon\}$.
Zatem $(f_n)$ jest ciągiem Cauchy'ego według miary.
Z twierdzenia~\ref{thm:cauchy-measure-complete-9}
istnieje mierzalna $f$ taka, że $f_n\to f$
według miary. Twierdzenie~\ref{thm:measure-metric-9}
daje $\rho([f_n],[f])\to0$.
\end{proof}

\begin{uwaga}[Skończoność miary jest istotna dla tego wzoru]
Jeśli $\mu(X)=\infty$, wartość
$\int_X\min\{1,|f-g|\}\,\mathrm d\mu(x)$
może być nieskończona, więc nie musi określać
metryki o wartościach skończonych.
Na przestrzeniach $\sigma$-skończonych można budować
metryki z wagami po wyczerpaniu przestrzeni zbiorami
o skończonej mierze, ale taka metryka na ogół
opisuje zbieżność \emph{lokalną} według miary,
nie globalną zbieżność w znaczeniu
\eqref{eq:conv-in-measure-9}.
\end{uwaga}


\begin{przyklad}[Zbieżność lokalna według miary nie jest globalną]
\label{ex:local-vs-global-9}
Na $\mathbb R$ z miarą Lebesgue'a rozważmy
$f_n=\chi_{[n,n+1)}$ i $f=0$. Dla każdego
$R>0$ mamy $f_n|_{[-R,R]}=0$ dla dostatecznie
dużych $n$; funkcje zbiegają zatem według miary na
każdym zbiorze $[-R,R]$. Jednak dla $0<\varepsilon<1$
zachodzi $\lambda_1(\{|f_n-f|>\varepsilon\})=1$
dla wszystkich $n$. Nie ma globalnej zbieżności
według miary. Ten przykład objaśnia różnicę między
metryzacją zbieżności lokalnej na przestrzeni
$\sigma$-skończonej i metryką \eqref{eq:measure-metric-9}
na przestrzeni o mierze skończonej.
\end{przyklad}

\begin{twierdzenie}[Kontrola ucieczki masy przy zbieżności lokalnej]
\label{thm:local-to-global-measure-9}
Niech $X_j\in\mathcal F$, $X_j\subseteq X_{j+1}$,
$\bigcup_jX_j=X$ i $\mu(X_j)<\infty$.
Niech $f_n,f:X\to\mathbb R$ będą mierzalne.
Wówczas $f_n\to f$ według miary na całym $X$
wtedy i tylko wtedy, gdy dla każdego $\varepsilon>0$:
\begin{enumerate}[label=\textup{(\roman*)}]
\item $\mu(X_j\cap\{|f_n-f|>\varepsilon\})\to0$
dla każdego ustalonego $j$;
\item $\displaystyle \lim_{j\to\infty}\limsup_{n\to\infty}
\mu((X\setminus X_j)\cap\{|f_n-f|>\varepsilon\})=0$.
\end{enumerate}
\end{twierdzenie}
\begin{proof}
Oznaczmy $B_n=\{|f_n-f|>\varepsilon\}$.
Jeśli $\mu(B_n)\to0$, oba warunki są natychmiastowe.
Odwrotnie, dla każdego $j,n$ z podaddytywności
\[
\mu(B_n)\le\mu(X_j\cap B_n)+
\mu((X\setminus X_j)\cap B_n).
\]
Przechodząc do $\limsup$ po $n$ i korzystając
z (i), dostajemy
$\limsup_n\mu(B_n)\le
\limsup_n\mu((X\setminus X_j)\cap B_n)$.
Na mocy (ii) prawa strona dąży do zera wraz
z $j\to\infty$, co dowodzi tezy.
\end{proof}

\section{Użyteczne kryteria graniczne}

\begin{twierdzenie}[Lemat Scheffégo]
\label{thm:scheffe-9}
Niech $f_n,f:X\to[0,\infty)$ będą mierzalne,
całkowalne oraz $f_n\to f$ prawie wszędzie.
Jeżeli
$\int_Xf_n(x)\,\mathrm d\mu(x)\to
\int_Xf(x)\,\mathrm d\mu(x)<\infty$,
to
$\int_X|f_n(x)-f(x)|\,\mathrm d\mu(x)\to0$.
\end{twierdzenie}
\begin{proof}
Mamy punktową tożsamość
$|f_n-f|=f_n+f-2\min(f_n,f)$.
Ponieważ $0\le\min(f_n,f)\le f$ oraz
$\min(f_n,f)\to f$ prawie wszędzie,
twierdzenie o zbieżności zmajoryzowanej daje
$\int_X\min(f_n,f)\,\mathrm d\mu(x)\to\int_Xf\,\mathrm d\mu(x)$.
Całkując tożsamość i korzystając z założenia
o granicach całek, otrzymujemy tezę.
\end{proof}

\begin{wniosek}[Wersja dla gęstości prawdopodobieństwa]\label{idx:09:wniosek:14}

Jeżeli $f_n,f\ge0$ są mierzalne na $(X,\mathcal F,\mu)$,
$\int_X f_n\,\mathrm d\mu(x)=
\int_X f\,\mathrm d\mu(x)=1$ oraz $f_n\to f$ prawie
wszędzie, to $\int_X|f_n-f|\,\mathrm d\mu(x)\to0$.
W zastosowaniach do prawdopodobieństwa jest to
podstawowy sposób przejścia od zbieżności gęstości
prawie wszędzie do zbieżności w wariacji całkowitej.
\end{wniosek}

\begin{uwaga}[Co naprawdę wynika z poszczególnych pojęć?]
\label{rem:conv-summary-9}
Na przestrzeni skończonej miary zbieżność prawie
wszędzie pociąga zbieżność według miary; na dowolnej
przestrzeni zbieżność według miary zapewnia podciąg
zbieżny prawie wszędzie. Ani zbieżność według miary,
ani w $L^p$ nie muszą zapewniać zbieżności prawie
wszędzie całego ciągu. Sama zbieżność prawie wszędzie
nie zapewnia zbieżności całek. Dodatkową kontrolę
zapewniają dominacja, monotoniczność lub jednostajna
całkowalność.
\end{uwaga}

\begin{center}
\begin{tikzpicture}[>=stealth, every node/.style={font=\small},
box/.style={draw,rounded corners,align=center,inner sep=5pt}]
\node[box] (u) at (0,2.5) {jednostajnie};
\node[box] (au) at (4.4,2.5) {prawie jednostajnie};
\node[box] (ae) at (4.4,0.7) {prawie wszędzie};
\node[box] (lp) at (-0.1,-1) {w $L^p$, $1\le p<\infty$};
\node[box] (mu) at (4.4,-1) {według miary};
\draw[->] (u) -- (au);
\draw[->] (au) -- (ae);
\draw[->] (au.east) to[out=0,in=0] (mu.east);
\draw[->] (lp) -- (mu);
\draw[->,dashed] (ae) -- node[right,font=\scriptsize]{gdy $\mu(X)<\infty$} (mu);
\draw[->,dashed,bend left=24] (ae.west) to node[above,font=\scriptsize]{Jegorow: $\mu(X)<\infty$} (au.south west);
\end{tikzpicture}
\par\smallskip\textit{Schemat implikacji. Strzałki przerywane wymagają skończoności miary.}
\end{center}

% END chapters/pl/measure_theory/09-rozdzial.tex


% BEGIN chapters/pl/measure_theory/10-rozdzial.tex
% ROZDZIAŁ 10 — do włączenia w main.tex
\chapter[Regularność funkcji mierzalnych]{Regularność funkcji mierzalnych i aproksymacja funkcjami ciągłymi}
\label{ch:regularnosc-funkcji}

\section{Co oznacza regularność funkcji mierzalnej?}

Funkcja mierzalna może być nieciągła we wszystkich punktach,
jak funkcja Dirichleta. Mimo to mierzalność jest blisko
związana z ciągłością \emph{w sensie miary}: na zbiorze o
skończonej mierze można usunąć zbiór o dowolnie małej mierze
i otrzymać ciągłość na pozostałym odpowiednio dobranym
zbiorze zwartym. Taką treść ma twierdzenie Łuzina.

W odróżnieniu od twierdzenia Jegorowa, które dotyczy
zbieżności \emph{ciągu} funkcji, twierdzenie Łuzina
przybliża pojedynczą funkcję mierzalną przez funkcje
ciągłe. Oba twierdzenia współpracują: najpierw
aproksymujemy funkcję mierzalną funkcjami prostymi,
a następnie wymuszamy jednostajność zbieżności
poza niewielkim zbiorem wyjątkowym.

W całym rozdziale $d\ge1$ i $\lambda_d$ oznacza miarę
Lebesgue'a na $\mathbb R^d$. Przez $\mathcal L_d$
oznaczamy jej $\sigma$-algebrę zbiorów mierzalnych.
Topologię na podzbiorach $\mathbb R^d$ rozumiemy jako
topologię podprzestrzeni. W szczególności sformułowanie
``$f|_K$ jest ciągła'' \emph{nie} oznacza,
że $f$ jest ciągła w punktach $K$ jako funkcja na
całej dziedzinie.

\begin{uwaga}[Dwa różne symbole]
Funkcję charakterystyczną zbioru $A$ oznaczamy
wyłącznie przez $\chi_A$. Gładkie funkcje odcinające
oznaczamy literą $\vartheta$ lub $\psi$; nie są one
funkcjami charakterystycznymi i przyjmują na ogół
również wartości pośrednie między zerem a jedynką.
\end{uwaga}

\section{Dwa narzędzia: regularność miary i funkcje odcinające}

\begin{lemat}[Przybliżenie mierzalnego zbioru przez zbiór zwarty]
\label{lem:compact-inside-10}
Niech $A\in\mathcal L_d$ oraz $\lambda_d(A)<\infty$.
Dla każdego $\varepsilon>0$ istnieją zbiór
zwarty $K$ i zbiór otwarty $O$ w $\mathbb R^d$ takie,
że $K\subseteq A\subseteq O$ i
\begin{equation}\label{eq:compact-open-10}
\lambda_d(O\setminus K)<\varepsilon.
\end{equation}
\end{lemat}
\begin{proof}
Regularność wewnętrzna miary Lebesgue'a daje
zwarty $K\subseteq A$ z
$\lambda_d(A\setminus K)<\varepsilon/2$.
Regularność zewnętrzna daje otwarty $O\supseteq A$
z $\lambda_d(O\setminus A)<\varepsilon/2$.
Ponieważ $O\setminus K$ jest rozłączną sumą
$(O\setminus A)\cup(A\setminus K)$,
addytywność miary daje \eqref{eq:compact-open-10}.
Jeśli chcemy jedynie otrzymać otwarte otoczenie
zbioru zwartego $K$ o małej różnicy miar,
możemy je dodatkowo wybrać ograniczone:
wystarczy przeciąć $O$ z kulą zawierającą $K$.
\end{proof}

\begin{uwaga}[Niedopuszczalna zamiana kwantyfikatorów]
Zbiór zwarty $K$ zależy od $\varepsilon$.
Na ogół nie istnieje jeden zbiór zwarty zawierający
wszystkie poza mierzalnym zbiorem zerowym punkty
$A$ (w szczególności gdy $A$ nie jest ograniczony).
Nie wolno również zastępować założenia
$\lambda_d(A)<\infty$ samym założeniem
$A\in\mathcal L_d$ w żądaniu
$\lambda_d(A\setminus K)<\varepsilon$.
\end{uwaga}

\begin{lemat}[Gładka funkcja odcinająca]
\label{lem:smooth-cutoff-10}
Niech $K\subseteq\mathbb R^d$ będzie zwarty,
a $O\subseteq\mathbb R^d$ otwarty i zawierający $K$.
Istnieje funkcja $\vartheta\in C_c^\infty(\mathbb R^d)$
taka, że $0\le\vartheta\le1$, $\vartheta=1$ na $K$
oraz $\operatorname{supp}\vartheta\subseteq O$.
Przez $\operatorname{supp}\vartheta$ rozumiemy domknięcie
zbioru $\{\vartheta\ne0\}$, a $C_c^\infty$ oznacza
funkcje gładkie o zwartym nośniku.
\end{lemat}
\begin{proof}
Jeśli $K=\varnothing$, bierzemy $\vartheta=0$.
W przeciwnym razie dla każdego $a\in K$ wybieramy
$r_a>0$ tak, żeby $\overline{B(a,r_a)}\subseteq O$.
Zwartość $K$ pozwala wybrać skończenie wiele
kul $B(a_j,r_j/2)$ pokrywających $K$.

Niech $\eta:\mathbb R\to[0,\infty)$ będzie określona
wzorem $\eta(t)=e^{-1/t}$ dla $t>0$ i $\eta(t)=0$
dla $t\le0$. Wzór ten określa funkcję klasy
$C^\infty$; każda jej pochodna dla $t>0$ jest
sumą wyrazów $P(1/t)e^{-1/t}$ i ma granicę zero
przy $t\downarrow0$, co dowodzi gładkości także
w zerze. Zdefiniujmy
\[
\psi_j(x)=
\frac{\eta(r_j^2-|x-a_j|^2)}
{\eta(r_j^2-|x-a_j|^2)+
 \eta(|x-a_j|^2-r_j^2/4)}.
\]
Mianownik jest dodatni dla każdego $x$:
dla $|x-a_j|<r_j$ dodatni jest pierwszy składnik,
a dla $|x-a_j|\ge r_j$ drugi. Widzimy, że
$\psi_j=1$ na $\overline{B(a_j,r_j/2)}$,
$\psi_j=0$ poza $B(a_j,r_j)$ i
$0\le\psi_j\le1$.
Kładziemy $\vartheta=1-\prod_j(1-\psi_j)$.
Jest gładka, równa $1$ na $K$ i ma nośnik zawarty
w skończonej sumie zwartych kul
$\overline{B(a_j,r_j)}\subseteq O$.
\end{proof}

\begin{wniosek}[Przybliżanie funkcji charakterystycznej]
\label{cor:indicator-smooth-10}
Niech $A\in\mathcal L_d$ i $\lambda_d(A)<\infty$.
Dla każdego $1\le p<\infty$ i $\varepsilon>0$
istnieje $\vartheta\in C_c^\infty(\mathbb R^d)$ taka,
że $0\le\vartheta\le1$ oraz
\begin{equation}\label{eq:indicator-smooth-10}
\int_{\mathbb R^d}|\chi_A(x)-\vartheta(x)|^p
\,\mathrm d\lambda_d(x)<\varepsilon.
\end{equation}
\end{wniosek}
\begin{proof}
Z lematu \ref{lem:compact-inside-10} wybieramy
$K\subseteq A\subseteq O$ z
$\lambda_d(O\setminus K)<\varepsilon$.
Stosujemy lemat \ref{lem:smooth-cutoff-10}.
Na $K$ funkcje $\chi_A$ i $\vartheta$ są równe $1$,
a poza $O$ są równe $0$. Na pozostałym zbiorze
$O\setminus K$ ich różnica ma moduł co najwyżej $1$.
Stąd całka po $\mathbb R^d$ jest ograniczona przez
$\lambda_d(O\setminus K)<\varepsilon$.
\end{proof}

\begin{center}
\begin{tikzpicture}[scale=.9,every node/.style={font=\small}]
\draw[rounded corners=10pt,thick,fill=blue!7,draw=blue!65!black]
(0,0) .. controls (-.7,1) and (-.2,2.3) .. (1.7,2.4)
.. controls (3.6,2.9) and (6,2.15) .. (6.4,1.25)
.. controls (6.9,-.3) and (3.6,-.5) .. (0,0);
\draw[rounded corners=8pt,thick,fill=green!10,draw=green!55!black]
(.6,.35) .. controls (.4,1.2) and (1,2.1) .. (2.05,2.05)
.. controls (4.2,2.5) and (5.55,1.85) .. (5.75,1.2)
.. controls (6,.35) and (2.5,.05) .. (.6,.35);
\draw[rounded corners=8pt,thick,fill=orange!18,draw=orange!75!black]
(1.25,.7) .. controls (1.1,1.25) and (1.55,1.8) .. (2.4,1.72)
.. controls (3.6,1.98) and (4.9,1.58) .. (5.02,1.04)
.. controls (4.5,.45) and (2.2,.4) .. (1.25,.7);
\node at (3.2,1.2) {$K$};
\node at (5.45,1.75) {$A$};
\node at (6.06,2.10) {$O$};
\end{tikzpicture}
\par\smallskip
\textit{Schemat regularności miary: $K\subseteq A\subseteq O$,
przy czym $\lambda_d(O\setminus K)$ jest dowolnie mała.
Rysunek ma charakter poglądowy.}
\end{center}

\section{Funkcje proste ciągłe na dużym zbiorze zwartym}

\begin{lemat}[Łuzin dla funkcji prostych]
\label{lem:lusin-simple-10}
Niech $E\in\mathcal L_d$, $\lambda_d(E)<\infty$,
a $s:E\to\mathbb R$ będzie mierzalną funkcją prostą.
Dla każdego $\varepsilon>0$ istnieje zwarty
$K\subseteq E$ taki, że
$\lambda_d(E\setminus K)<\varepsilon$ oraz
$s|_K$ jest ciągła.
\end{lemat}
\begin{proof}
Niech $a_1,\dots,a_m$ będą różnymi wartościami $s$
i $A_j=\{x\in E:s(x)=a_j\}$.
Są to parami rozłączne zbiory mierzalne, których
suma jest równa $E$. Wszystkie mają skończoną miarę.
Dla każdego $j$ regularność wewnętrzna daje
zwarty $K_j\subseteq A_j$ taki, że
$\lambda_d(A_j\setminus K_j)<\varepsilon/m$.
Zbiór $K=\bigcup_{j=1}^mK_j$ jest zwarty,
$K\subseteq E$ i
$\lambda_d(E\setminus K)<\varepsilon$.

Ponieważ zbiory $K_j$ są parami rozłączne i zwarte,
każdy z nich jest otwarty i domknięty w podprzestrzeni
$K$: jego dopełnienie w $K$ jest skończoną sumą
zbiorów zwartych, a więc domkniętych.
Funkcja $s|_K$, stała na każdym z tych zbiorów,
jest zatem ciągła. Nie wyciągamy błędnego wniosku,
że $s$ jest ciągła na całym $E$.
\end{proof}

\section{Twierdzenie Łuzina}

\begin{twierdzenie}[Łuzin]
\label{thm:lusin-10}
Niech $E\in\mathcal L_d$ będzie zbiorem o
$\lambda_d(E)<\infty$, a $f:E\to\mathbb R$
funkcją mierzalną względem śladu
$\mathcal L_d|_E$. Wtedy dla każdego
$\varepsilon>0$ istnieje zwarty zbiór
$K\subseteq E$ taki, że
\begin{equation}\label{eq:lusin-measure-10}
\lambda_d(E\setminus K)<\varepsilon
\end{equation}
oraz $f|_K:K\to\mathbb R$ jest ciągła.
\end{twierdzenie}
\begin{proof}
\textit{Krok 1: aproksymacja funkcją prostą.}
Na mocy twierdzenia o aproksymacji funkcji
mierzalnych z rozdziału 6 istnieją mierzalne
funkcje proste $s_n:E\to\mathbb R$ takie, że
$s_n(x)\to f(x)$ dla każdego $x\in E$.
Możemy na przykład stosować obcięcia wartości
$f$ i ich kwantyzację na przedziałach długości
$2^{-n}$, najpierw do części dodatniej i ujemnej.
Nie korzystamy w tym kroku z twierdzenia Łuzina.

\textit{Krok 2: prawie jednostajność.}
Ponieważ $\lambda_d(E)<\infty$, twierdzenie
Jegorowa \ref{thm:egorov-9}, zastosowane do
przestrzeni z miarą $E$, daje mierzalny zbiór
$Z\subseteq E$ taki, że
$\lambda_d(Z)<\varepsilon/4$ i $s_n\to f$
jednostajnie na $E\setminus Z$.
Z regularności wewnętrznej wybieramy zwarty
$K_0\subseteq E\setminus Z$, dla którego
$\lambda_d((E\setminus Z)\setminus K_0)<\varepsilon/4$.
Na $K_0$ zbieżność jest jednostajna.

\textit{Krok 3: ciągłość wszystkich przybliżeń.}
Dla każdego $n\ge1$ stosujemy lemat
\ref{lem:lusin-simple-10} do $s_n$ na $E$
z tolerancją $\varepsilon 2^{-(n+2)}$.
Otrzymujemy zwarty $K_n\subseteq E$ taki, że
$\lambda_d(E\setminus K_n)<
\varepsilon2^{-(n+2)}$ i $s_n|_{K_n}$ jest ciągła.
Kładziemy
$K=K_0\cap\bigcap_{n\ge1}K_n$.
Jest to zbiór zwarty jako domknięty podzbiór
zbioru zwartego $K_0$. Ponadto
\[
\lambda_d(E\setminus K)
\le\lambda_d(Z)+
\lambda_d((E\setminus Z)\setminus K_0)+
\sum_{n\ge1}\lambda_d(E\setminus K_n)
<\frac\varepsilon4+\frac\varepsilon4+
\frac\varepsilon4<\varepsilon.
\]
Dla każdego $n$ funkcja $s_n|_K$ jest ciągła,
a ciąg $(s_n|_K)$ zbiega jednostajnie do $f|_K$.
Jednostajna granica funkcji ciągłych na przestrzeni
metrycznej jest ciągła; stąd $f|_K$ jest ciągła.
\end{proof}

\begin{wniosek}[Funkcje skończone prawie wszędzie]
\label{cor:lusin-finite-ae-10}
Twierdzenie \ref{thm:lusin-10} pozostaje prawdziwe,
gdy $f:E\to\overline{\mathbb R}$ jest mierzalna i
skończona $\lambda_d$-prawie wszędzie, pod warunkiem,
że przez ciągłość $f|_K$ rozumiemy ciągłość
funkcji \emph{rzeczywistej} na $K$.
\end{wniosek}
\begin{proof}
Zbiór $N=\{x\in E:|f(x)|=\infty\}$ jest
mierzalny i zerowy. Na $E\setminus N$ funkcja
jest rzeczywista i mierzalna. Stosujemy twierdzenie
Łuzina do $E\setminus N$, którego miara jest
równa $\lambda_d(E)$. Otrzymany zwarty
$K\subseteq E\setminus N$ ma żądane własności.
\end{proof}

\begin{wniosek}[Jednoczesne twierdzenie Łuzina dla przeliczalnej rodziny]
\label{cor:countable-lusin-10}
Niech $E\in\mathcal L_d$ ma skończoną miarę,
a $f_j:E\to\mathbb R$, $j\ge1$, będą
funkcjami mierzalnymi. Dla każdego $\varepsilon>0$
istnieje zwarty $K\subseteq E$ taki, że
$\lambda_d(E\setminus K)<\varepsilon$ i
\emph{wszystkie} ograniczenia $f_j|_K$ są ciągłe.
\end{wniosek}
\begin{proof}
Dla każdego $j$ stosujemy twierdzenie Łuzina
z tolerancją $\varepsilon2^{-(j+1)}$.
Otrzymujemy zwarte $K_j\subseteq E$ takie,
że $\lambda_d(E\setminus K_j)<\varepsilon2^{-(j+1)}$
i $f_j|_{K_j}$ jest ciągła.
Zbiór $K=\bigcap_{j\ge1}K_j$ jest zwarty
(jako domknięty podzbiór $K_1$), a
\[
\lambda_d(E\setminus K)
\le\sum_{j\ge1}\lambda_d(E\setminus K_j)
<\varepsilon.
\]
Każda $f_j|_K$ jest ciągła jako ograniczenie
funkcji ciągłej na $K_j$.
\end{proof}

\begin{uwaga}[Granice twierdzenia]
\label{rem:lusin-limitations-10}
Twierdzenie Łuzina nie orzeka, że $f$ jest
ciągła \emph{w punktach} $K$ względem całej
dziedziny $E$. Na przykład $f=\chi_{\mathbb Q}$
na $[0,1]$ jest nieciągła w każdym punkcie, ale
jej ograniczenie do dowolnego zbioru zwartego
niezawierającego liczb wymiernych jest funkcją
stałą, a więc ciągłą. Możemy znaleźć takie
zbiory zwarte o mierze dowolnie bliskiej $1$.

Skończoność miary $E$ jest istotna w podanej formie.
Dla $E=\mathbb R^d$ żaden zwarty podzbiór $K$
nie spełnia $\lambda_d(E\setminus K)<1$,
niezależnie od wyboru funkcji $f$.
\end{uwaga}

\section{Przedłużanie funkcji ciągłej ze zbioru zwartego}

\begin{lemat}[Lipschitzowska aproksymacja na zbiorze zwartym]
\label{lem:lipschitz-approx-compact-10}
Niech $K\subseteq\mathbb R^d$ będzie niepustym
zbiorem zwartym, a $h:K\to\mathbb R$ funkcją ciągłą.
Dla każdego $\varepsilon>0$ istnieje ograniczona
funkcja Lipschitzowska $H:\mathbb R^d\to\mathbb R$
taka, że
$\sup_{x\in K}|H(x)-h(x)|<\varepsilon$
i $\sup_{\mathbb R^d}|H|\le\sup_K|h|$.
\end{lemat}
\begin{proof}
Niech $M=\sup_K|h|<\infty$. Jeśli $M=0$,
wystarczy $H=0$. W przeciwnym razie z jednostajnej
ciągłości $h$ na $K$ wybieramy $\delta>0$ takie, że
$|h(x)-h(y)|<\varepsilon/2$ przy
$x,y\in K$, $|x-y|<\delta$.
Wybieramy $L>2M/\delta$ i definiujemy na
całym $\mathbb R^d$
\[
U_L(x)=\inf_{y\in K}\bigl(h(y)+L|x-y|\bigr).
\]
Dla dowolnych $x,z\in\mathbb R^d$ z nierówności
$|x-y|\le|x-z|+|z-y|$ wynika
$U_L(x)\le U_L(z)+L|x-z|$. Zamieniając $x,z$,
dostajemy $|U_L(x)-U_L(z)|\le L|x-z|$.

Dla $x\in K$ mamy $U_L(x)\le h(x)$ (wybierając
$y=x$). Jeśli $|y-x|<\delta$, to
$h(y)+L|x-y|>h(x)-\varepsilon/2$.
Jeżeli $|y-x|\ge\delta$, to
$h(y)+L|x-y|\ge-M+L\delta>M\ge h(x)$.
Wynika $h(x)-\varepsilon/2\le U_L(x)\le h(x)$.
Funkcja
$H(x)=\max(-M,\min(M,U_L(x)))$
jest ograniczona i Lipschitzowska, a obcięcie do
$[-M,M]$ nie zwiększa błędu na $K$.
\end{proof}

\begin{twierdzenie}[Przedłużenie z kompaktu w $\mathbb R^d$]
\label{thm:compact-extension-10}
Jeśli $K\subseteq\mathbb R^d$ jest zwarty,
a $h:K\to\mathbb R$ ciągła, to istnieje
$H\in C(\mathbb R^d)$ taka, że $H|_K=h$.
Można dodatkowo wymagać
$\sup_{\mathbb R^d}|H|\le\sup_K|h|$.
\end{twierdzenie}
\begin{proof}
Przy $K=\varnothing$ bierzemy $H=0$.
Niech $K\ne\varnothing$ i $M_0=\sup_K|h|$.
Konstruujemy funkcje Lipschitzowskie $u_n$ i
reszty $r_n:K\to\mathbb R$ rekurencyjnie.
Kładziemy $r_0=h$. Jeśli
$M_n=\sup_K|r_n|>0$, stosujemy lemat
\ref{lem:lipschitz-approx-compact-10} z błędem
$M_n/2$, uzyskując $u_n$ o
$\sup_{\mathbb R^d}|u_n|\le M_n$ i
$\sup_K|u_n-r_n|<M_n/2$.
Definiujemy $r_{n+1}=r_n-u_n|_K$.
Wtedy $M_{n+1}\le M_n/2$. Jeśli $M_n=0$,
bierzemy $u_n=0$ i wszystkie dalsze reszty
równe zero.

Z oszacowania
$\sup_{\mathbb R^d}|u_n|\le M_0 2^{-n}$
wynika jednostajna zbieżność szeregu
$G=\sum_{n\ge0}u_n$ na całym $\mathbb R^d$.
Jest to funkcja ciągła. Na $K$ mamy dla każdego
$N$ tożsamość
$h=\sum_{n=0}^Nu_n|_K+r_{N+1}$;
po przejściu do granicy otrzymujemy $G|_K=h$.
Na koniec funkcja $H=\max(-M_0,\min(M_0,G))$
pozostaje ciągła, zgadza się z $h$ na $K$
i ma żądane ograniczenie.
\end{proof}

\begin{wniosek}[Wersja Łuzina przez funkcję ciągłą na całej przestrzeni]
\label{cor:lusin-extension-10}
Niech $E\in\mathcal L_d$, $\lambda_d(E)<\infty$,
i $f:E\to\mathbb R$ będzie mierzalna. Dla każdego
$\varepsilon>0$ istnieją $K\subseteq E$ zwarty
i $H\in C(\mathbb R^d)$ takie, że
$\lambda_d(E\setminus K)<\varepsilon$ oraz
$H(x)=f(x)$ dla wszystkich $x\in K$.
\end{wniosek}
\begin{proof}
Stosujemy twierdzenie \ref{thm:lusin-10},
a następnie twierdzenie \ref{thm:compact-extension-10}
do funkcji ciągłej $f|_K$.
\end{proof}

\section{Aproksymacja w mierze przez funkcje ciągłe i gładkie}

\begin{twierdzenie}[Aproksymacja funkcji o nośniku mierzalnym skończonej miary]
\label{thm:cc-approx-in-measure-10}
Niech $E\in\mathcal L_d$, $\lambda_d(E)<\infty$,
i niech $f:E\to\mathbb R$ będzie mierzalna.
Przedłużmy $f$ zerem poza $E$, otrzymując
$\widetilde f:\mathbb R^d\to\mathbb R$.
Dla każdych $\varepsilon,\delta>0$ istnieje
$\varphi\in C_c(\mathbb R^d)$ takie, że
\begin{equation}\label{eq:cc-approx-measure-10}
\lambda_d\bigl(\{x\in\mathbb R^d:
|\widetilde f(x)-\varphi(x)|>\delta\}\bigr)
<\varepsilon.
\end{equation}
Można dodatkowo wybrać $\varphi$ tak, aby
$\varphi=\widetilde f$ \emph{dokładnie} na pewnym
zbiorze zwartym $K\subseteq E$ z
$\lambda_d(E\setminus K)<\varepsilon/2$.
\end{twierdzenie}
\begin{proof}
Z twierdzenia Łuzina wybieramy zwarty
$K\subseteq E$ taki, że
$\lambda_d(E\setminus K)<\varepsilon/2$
oraz $f|_K$ jest ciągła.
Z regularności zewnętrznej \emph{zbioru zwartego}
$K$ wybieramy otwarty $O$ zawierający $K$ tak,
że $\lambda_d(O\setminus K)<\varepsilon/2$.
Nie musimy wymagać, aby $O$ zawierał całe $E$.

Przedłużamy $f|_K$ do funkcji ciągłej
$H:\mathbb R^d\to\mathbb R$ na mocy
\ref{thm:compact-extension-10}. Na mocy lematu
\ref{lem:smooth-cutoff-10} istnieje
$\vartheta\in C_c^\infty(\mathbb R^d)$, $\vartheta=1$ na $K$,
z $\operatorname{supp}\vartheta\subseteq O$.
Kładziemy $\varphi=\vartheta H$. Wtedy
$\varphi\in C_c(\mathbb R^d)$ i na $K$ zachodzi
$\varphi=f$. Poza $E\cup O$ obie funkcje
$\varphi$ i $\widetilde f$ są zerowe. Stąd
zbiór, na którym funkcje mogą się różnić,
jest zawarty w $(E\setminus K)\cup(O\setminus K)$,
a jego miara jest mniejsza niż $\varepsilon$.
Implikuje to \eqref{eq:cc-approx-measure-10}
dla każdego $\delta>0$.
\end{proof}

\begin{uwaga}[Dlaczego nie istnieje takie twierdzenie dla dowolnego $f$ na $\mathbb R^d$?]
Dla funkcji $f\equiv1$ na $\mathbb R^d$ i dowolnej
$\varphi\in C_c(\mathbb R^d)$ zbiór
$\{|f-\varphi|>1/2\}$ zawiera dopełnienie
zwartego nośnika $\varphi$, a więc ma nieskończoną
miarę Lebesgue'a. Założenie o skończonej mierze
zbioru podtrzymującego funkcję nie jest czysto
technicznym dodatkiem.
\end{uwaga}

\begin{lemat}[Wygładzanie funkcji ciągłej o zwartym nośniku]
\label{lem:mollify-cc-10}
Niech $h\in C_c(\mathbb R^d)$. Istnieje ciąg
$h_n\in C_c^\infty(\mathbb R^d)$ taki, że
$\sup_{x\in\mathbb R^d}|h_n(x)-h(x)|\to0$.
\end{lemat}
\begin{proof}
Z lematu \ref{lem:smooth-cutoff-10} otrzymujemy
nieujemną, niezerową funkcję
$\rho\in C_c^\infty(\mathbb R^d)$ o nośniku
zawartym w kuli $B(0,1)$. Po podzieleniu przez
$\int_{\mathbb R^d}\rho(x)\,\mathrm d\lambda_d(x)>0$
możemy przyjąć, że ta całka jest równa $1$.
Dla $t>0$ określamy
$\rho_t(x)=t^{-d}\rho(x/t)$ oraz
\[
h_t(x)=\int_{\mathbb R^d}
 h(y)\rho_t(x-y)\,\mathrm d\lambda_d(y).
\]
Z niezmienniczości i skalowania miary Lebesgue'a
(rozdział 4) wynika
$\int_{\mathbb R^d}\rho_t(x-y)\,\mathrm d\lambda_d(y)=1$.

Funkcja $h_t$ ma nośnik zawarty w
$\operatorname{supp}h+\overline{B(0,t)}$,
który jest zwarty. Jest gładka: pochodne względem
$x$ dowolnego rzędu można przenosić pod znak
całki dzięki ograniczoności pochodnych $\rho_t$
i temu, że integrujemy po ustalonym zwartym
nośniku $h$. Formalnie stosujemy iteracyjnie
twierdzenie o różniczkowaniu pod znakiem całki
z rozdziału 8, osobno do każdej współrzędnej.

Ponadto dla każdego $x$ mamy
\[
|h_t(x)-h(x)|
\le\int_{\mathbb R^d}|h(y)-h(x)|\rho_t(x-y)
\,\mathrm d\lambda_d(y)
\le\sup_{|x-y|\le t}|h(y)-h(x)|.
\]
Funkcja $h$ jest jednostajnie ciągła na
$\mathbb R^d$, bo jest ciągła i ma zwarty nośnik.
Ostatnie supremum dąży jednostajnie do zera dla
$t\downarrow0$. Wybieramy $h_n=h_{1/n}$.
\end{proof}

\begin{wniosek}[Aproksymacja w mierze funkcjami gładkimi]
\label{cor:smooth-approx-in-measure-10}
W sytuacji twierdzenia
\ref{thm:cc-approx-in-measure-10}, dla dowolnych
$\varepsilon,\delta>0$ istnieje
$\psi\in C_c^\infty(\mathbb R^d)$ taka, że
\[
\lambda_d\bigl(\{|\widetilde f-\psi|>\delta\}\bigr)
<\varepsilon.
\]
\end{wniosek}
\begin{proof}
Z twierdzenia \ref{thm:cc-approx-in-measure-10}
wybieramy $\varphi\in C_c(\mathbb R^d)$ taką,
że $\lambda_d(\{|\widetilde f-\varphi|>
\delta/2\})<\varepsilon$.
Z lematu \ref{lem:mollify-cc-10} wybieramy
$\psi\in C_c^\infty$ taką, że
$\sup|\psi-\varphi|<\delta/2$.
Wtedy
$\{|\widetilde f-\psi|>\delta\}\subseteq
\{|\widetilde f-\varphi|>\delta/2\}$.
\end{proof}

\section{Gęstość $C_c^\infty$ w przestrzeniach $L^p$}

\begin{twierdzenie}[Gęstość funkcji gładkich]
\label{thm:cc-in-lp-10}
Niech $d\ge1$ i $1\le p<\infty$.
Dla każdej funkcji mierzalnej
$f:\mathbb R^d\to\mathbb R$ spełniającej
$\int_{\mathbb R^d}|f(x)|^p\,\mathrm d\lambda_d(x)
<\infty$ oraz dla każdego $\varepsilon>0$
istnieje $\varphi\in C_c^\infty(\mathbb R^d)$ taka, że
\begin{equation}\label{eq:lp-smooth-dense-10}
\int_{\mathbb R^d}|f(x)-\varphi(x)|^p
\,\mathrm d\lambda_d(x)<\varepsilon.
\end{equation}
W szczególności $C_c^\infty(\mathbb R^d)$ jest
gęsta w $L^p(\mathbb R^d,\lambda_d)$.
\end{twierdzenie}
\begin{proof}
\textit{Krok 1: obcięcie nośnika i wartości.}
Niech $B_R=B(0,R)$. Kładziemy
\[
f_{R,M}(x)=\max\{-M,\min\{f(x),M\}\}\chi_{B_R}(x).
\]
Przy $R,M\to\infty$ (na przykład po przekątnej
$R=M=n$) mamy $f_{R,M}(x)\to f(x)$ dla prawie
wszystkich $x$ oraz
$|f(x)-f_{R,M}(x)|^p\le 2^p|f(x)|^p$.
Twierdzenie o zbieżności zmajoryzowanej daje
$\int_{\mathbb R^d}|f-f_{R,M}|^p\,\mathrm d\lambda_d(x)\to0$. Możemy zatem przybliżać
$f$ funkcjami ograniczonymi, których nośniki są
zawarte w kulach o skończonej mierze.

\textit{Krok 2: przybliżenie funkcją prostą.}
Ustalmy ograniczoną mierzalną funkcję $g$
zerową poza kulą $B_R$. Dla $h>0$ przyjmijmy
$s_h(x)=h\lfloor g(x)/h\rfloor$ na $B_R$ i
$s_h(x)=0$ poza $B_R$. Funkcja $s_h$ jest prosta,
ponieważ $g$ jest ograniczona, a jej poziomice
są mierzalne. Na $B_R$ mamy $|s_h-g|\le h$,
więc
$\int_{\mathbb R^d}|s_h-g|^p\,\mathrm d\lambda_d(x)
\le h^p\lambda_d(B_R)\to0$.
Po ewentualnym usunięciu zerowych współczynników
zapisujemy $s_h=\sum_{j=1}^m a_j\chi_{A_j}$,
gdzie wszystkie $A_j$ są mierzalne i mają skończoną
miarę (jako podzbiory $B_R$).

\textit{Krok 3: przybliżenie funkcji prostej.}
Dla każdego $j$ stosujemy wniosek
\ref{cor:indicator-smooth-10}, aby znaleźć
$\vartheta_j\in C_c^\infty$ z dowolnie małą całką
$\int_{\mathbb R^d}|\chi_{A_j}-\vartheta_j|^p\,\mathrm d\lambda_d(x)$.
Kładziemy $\varphi=\sum_{j=1}^m a_j\vartheta_j$.
Jest to funkcja gładka o zwartym nośniku.
Z nierówności elementarnej
$|z_1+\dots+z_m|^p\le m^{p-1}
\sum_{j=1}^m|z_j|^p$ dla $p\ge1$ wynika
\[
\int_{\mathbb R^d}|s_h-\varphi|^p
\,\mathrm d\lambda_d(x)
\le m^{p-1}\sum_{j=1}^m|a_j|^p
\int_{\mathbb R^d}|\chi_{A_j}-\vartheta_j|^p
\,\mathrm d\lambda_d(x).
\]
Prawą stronę możemy uczynić dowolnie małą.

\textit{Krok 4: połączenie oszacowań.}
Używamy trzykrotnie nierówności
$|u+v|^p\le 2^{p-1}(|u|^p+|v|^p)$,
aby oszacować $|f-\varphi|^p$ przez stałą
razy sumę błędów z poprzednich kroków.
Wybierając kolejno $R,M$, następnie $h$,
a na końcu $\vartheta_j$, możemy uczynić
$\int_{\mathbb R^d}|f-\varphi|^p\,\mathrm d\lambda_d(x)$ mniejszą od zadanego
$\varepsilon$. Nie korzystamy z nierówności
Minkowskiego, której pełny dowód będzie
przedmiotem rozdziału 14.
\end{proof}

\begin{przyklad}[Dlaczego nie można zastąpić $p<\infty$ przez $p=\infty$]
Niech $f=\chi_{(0,\infty)}$ na $\mathbb R$.
Nie istnieje funkcja ciągła $g:\mathbb R\to\mathbb R$
taka, że $\operatorname*{ess\,sup}_{x\in\mathbb R}
|f(x)-g(x)|<1/2$. W przeciwnym razie dla prawie
wszystkich $x<0$ zachodziłoby $g(x)<1/2$,
a dla prawie wszystkich $x>0$ zachodziłoby
$g(x)>1/2$. Z ciągłości otrzymalibyśmy
$g(0)\le1/2$ i $g(0)\ge1/2$; przy ścisłym
oszacowaniu błędu $q<1/2$ dostalibyśmy nawet
$g(0)\le q<1-q\le g(0)$, sprzeczność.
Zatem funkcje ciągłe nie są gęste w $L^\infty$
w normie supremum istotnego.
\end{przyklad}

\section{Granice prawie wszędzie funkcji gładkich}

\begin{twierdzenie}[Prawie wszędzie gładka aproksymacja funkcji mierzalnej]
\label{thm:ae-smooth-approx-10}
Dla dowolnej funkcji mierzalnej
$f:\mathbb R^d\to\mathbb R$ istnieje ciąg
$(\varphi_n)\subseteq C_c^\infty(\mathbb R^d)$
taki, że $\varphi_n(x)\to f(x)$ dla
$\lambda_d$-prawie wszystkich $x\in\mathbb R^d$.
Nie zakładamy, że $f$ jest całkowalna.
\end{twierdzenie}
\begin{proof}
Niech $B_n=B(0,n)$. Stosujemy wniosek
\ref{cor:smooth-approx-in-measure-10} do ograniczenia
$f|_{B_n}$, tolerancji błędu $1/n$ i miary
wyjątkowego zbioru $2^{-n}$. Otrzymujemy
$\varphi_n\in C_c^\infty$ takie, że
\[
\lambda_d\bigl(B_n\cap
\{|f-\varphi_n|>1/n\}\bigr)<2^{-n}.
\]
Niech $A_n=B_n\cap\{|f-\varphi_n|>1/n\}$.
Zbiór $\limsup_n A_n$ ma miarę zero na mocy
lematu Borela--Cantellego. Dla $x$ poza nim
i dostatecznie dużego $n$ zachodzi $x\in B_n$
oraz $|f(x)-\varphi_n(x)|\le1/n$.
Stąd $\varphi_n(x)\to f(x)$.
\end{proof}

\begin{wniosek}[Borelowska reprezentacja funkcji Lebesgue'a]
\label{cor:borel-representative-10}
Jeżeli $f:\mathbb R^d\to\mathbb R$ jest
$\mathcal L_d$-mierzalna, to istnieje funkcja
$g:\mathbb R^d\to\mathbb R$ borelowsko mierzalna
taka, że $f=g$ $\lambda_d$-prawie wszędzie.
\end{wniosek}
\begin{proof}
Niech $(\varphi_n)$ będzie ciągiem z twierdzenia
\ref{thm:ae-smooth-approx-10}.
Zbiór $D$, na którym $\varphi_n(x)$ zbiega
do skończonej liczby, jest borelowski, ponieważ
istnienie granicy skończonej można wyrazić przez
warunek Cauchy'ego z przeliczalnymi kwantyfikatorami.
Definiujemy $g(x)=\lim_n\varphi_n(x)$ dla $x\in D$
i $g(x)=0$ dla $x\notin D$.
Jest to funkcja borelowska i zgadza się z $f$
poza zbiorem miary zero.
\end{proof}

\begin{uwaga}[Granica punktowa wszędzie a granica prawie wszędzie]
Nie każda funkcja mierzalna Lebesgue'a jest granicą
punktową \emph{wszędzie} funkcji ciągłych.
Istnieją funkcje Lebesgue'a, które nie są nawet
borelowsko mierzalne, natomiast każda granica
punktowa wszędzie funkcji borelowskich jest
borelowska. Wnioski
\ref{thm:ae-smooth-approx-10} i
\ref{cor:borel-representative-10} celowo mówią
o równości prawie wszędzie.
\end{uwaga}


\begin{uwaga}[Przedstawiciele borelowscy a zbiory wyjątkowe]
\label{rem:borel-null-representatives-10}
Wniosek~\ref{cor:borel-representative-10} zapewnia
istnienie funkcji borelowskiej $g$ równej $f$
$\lambda_d$-prawie wszędzie. Nie zapewnia natomiast,
że zbiór $\{f\ne g\}$ jest borelowski: jeśli $f$
jest jedynie mierzalna w sensie Lebesgue'a,
zbiór ten może być nieborelowski. Istnieje za to
\emph{borelowski zbiór zerowy} $N$ zawierający
$\{f\ne g\}$, ponieważ miara Lebesgue'a jest
uzupełnieniem miary borelowskiej i każdy mierzalny
zbiór zerowy ma borelowską otoczkę miary zero.
To rozróżnienie będzie potrzebne przy porównywaniu
miary Lebesgue'a z produktem miar w rozdziale
\ref{ch:produkty-miary}.
\end{uwaga}

\section{Podsumowanie i zastosowania}

Z regularności miary Lebesgue'a wynika, że zbiory
mierzalne o skończonej mierze można przybliżać
zbiorami zwartymi i otwartymi. W połączeniu z
twierdzeniem Jegorowa i aproksymacją funkcjami
prostymi prowadzi to do twierdzenia Łuzina.
Pozwala ono przedłużyć funkcję z dużego zbioru
zwartego do funkcji ciągłej na całej przestrzeni.
Z kolei przybliżanie funkcji charakterystycznych
przez gładkie funkcje odcinające daje gęstość
$C_c^\infty$ w $L^p$ dla $1\le p<\infty$.
Te konstrukcje przygotowują późniejsze rozdziały
o całkach produktowych, funkcjonałach na $L^p$,
miarach Radona oraz słabych pochodnych.

% END chapters/pl/measure_theory/10-rozdzial.tex


% BEGIN chapters/pl/measure_theory/11-rozdzial.tex
% ROZDZIAŁ 11 — do włączenia w main.tex
\chapter{Całki zależne od parametru}
\label{ch:parametry-calki}

\section{Problem i jego dokładne sformułowanie}

Niech $(X,\mathcal F,\mu)$ będzie przestrzenią z miarą, $T$ przestrzenią
parametrów, a $f:T\times X\to\mathbb C$ rodziną funkcji. Zajmujemy się
funkcją $F:T\to\mathbb C$ określoną wzorem
$F(t)=\int_X f(t,x)\,\mathrm d\mu(x)$, o ile $f(t,\cdot)$ jest całkowalna.
Pytamy, kiedy własności funkcji $t\mapsto f(t,x)$ przechodzą na $F$:
ciągłość, różniczkowalność, wyższa regularność i analityczność. Problem
polega na tym, że własność zachodząca osobno dla każdego $x$ nie
kontroluje automatycznie masy rozłożonej na całym $X$.

\begin{przyklad}[Parametr uciekający do nieskończoności]
\label{ex:params-escape-11}
Na $X=(0,\infty)$ z miarą Lebesgue'a połóżmy dla $t>0$
$f(t,x)=t\chi_{(0,1/t)}(x)$, a $f(0,x)=0$.
Wówczas $f(t,x)\to0$ dla każdego $x>0$ przy $t\downarrow0$,
ale $\int_X f(t,x)\,\mathrm d\lambda_1(x)=1$ dla każdego $t>0$.
Zatem $F$ nie jest ciągła w zerze, mimo ciągłości $t\mapsto f(t,x)$
w punkcie $0$ dla każdego $x$. Nie istnieje też jedna funkcja
$g\in L^1(X)$ dominująca wszystkie $f(t,\cdot)$ dla dostatecznie
małych $t>0$: w przeciwnym razie twierdzenie o zbieżności
zmajoryzowanej dawałoby sprzeczność.
\end{przyklad}

\begin{uwaga}[Zakres rozdziału]
W rozdziale 8 udowodniono podstawowe twierdzenia o całkach z parametrem.
Tutaj formułujemy ich wersje \emph{lokalne} i \emph{jednostajne},
wyprowadzamy twierdzenia wyższego rzędu i badamy kilka użytecznych
transformacji. Nie używamy twierdzenia Fubiniego, które zostanie
udowodnione dopiero w rozdziale~\ref{ch:produkty-miary}.
\end{uwaga}

\section{Dominacja lokalna i ciągłość w parametrze}

\begin{definicja}[Lokalna dominacja całkowalna]
\label{def:local-majorant-11}
Niech $(T,d)$ będzie przestrzenią metryczną, $t_0\in T$ oraz
$f:T\times X\to\mathbb C$. Mówimy, że rodzina $f$ jest
\emph{lokalnie zdominowana całkowalnie w $t_0$}, jeżeli istnieją
otoczenie $U$ punktu $t_0$ w $T$ oraz mierzalna funkcja
$g:X\to[0,\infty)$ z $\int_Xg\,\mathrm d\mu(x)<\infty$ takie, że
istnieje $N\in\mathcal F$, $\mu(N)=0$, taki, że
$|f(t,x)|\le g(x)$ dla wszystkich $t\in U$ i $x\in X\setminus N$.
Warunek ten odróżniamy od słabszego wymagania, by dla każdego $t$
istniał osobny zbiór zerowy $N_t$. W twierdzeniu o ciągłości
słabszy wariant również wystarcza: dla ustalonego ciągu
$t_n\to t_0$ zbiór $N_{t_0}\cup\bigcup_n N_{t_n}$
pozostaje zerowy, toteż nadal można użyć twierdzenia
o zbieżności zmajoryzowanej. W pozostałych twierdzeniach
sprawdzamy potrzebny wariant oddzielnie.
\end{definicja}

\begin{twierdzenie}[Lokalne kryterium ciągłości]
\label{thm:local-continuity-11}
Niech $(T,d)$ będzie przestrzenią metryczną, $t_0\in T$,
a $f:T\times X\to\mathbb C$ spełnia następujące warunki:
\begin{enumerate}[label=\textup{(\roman*)}]
\item dla każdego $t\in T$ funkcja $x\mapsto f(t,x)$ jest mierzalna;
\item dla $\mu$-prawie każdego $x$ odwzorowanie $t\mapsto f(t,x)$
      jest ciągłe w $t_0$;
\item istnieje całkowalna dominacja w pewnym otoczeniu $t_0$
      w sensie definicji~\ref{def:local-majorant-11}.
\end{enumerate}
Wtedy $F(t)=\int_X f(t,x)\,\mathrm d\mu(x)$ jest określona
w pewnym otoczeniu $t_0$ i ciągła w $t_0$.
Jeśli warunki te zachodzą w każdym $t_0\in T$, to $F$ jest
ciągła na $T$.
\end{twierdzenie}
\begin{proof}
Ustalmy otoczenie $U$ i funkcję $g$ z (iii). Dla $t\in U$
mamy $\int_X|f(t,x)|\,\mathrm d\mu(x)\le\int_Xg\,\mathrm d\mu(x)$,
więc $F(t)$ istnieje jako całka bezwzględnie zbieżna. Niech
$t_n\to t_0$. Dla dużych $n$ mamy $t_n\in U$, a funkcje
$f(t_n,\cdot)$ zbiegają do $f(t_0,\cdot)$ prawie wszędzie i
są zdominowane przez $g$. Twierdzenie o zbieżności zmajoryzowanej,
zastosowane do części rzeczywistej i urojonej lub bezpośrednio
do modułu różnicy, daje
$\int_X |f(t_n,x)-f(t_0,x)|\,\mathrm d\mu(x)\to0$.
W konsekwencji $|F(t_n)-F(t_0)|\to0$. Ciągowa charakteryzacja
ciągłości w przestrzeni metrycznej kończy dowód.
\end{proof}

\begin{wniosek}[Ciągłość w normie $L^1$]
\label{cor:l1-param-11}
Przy założeniach twierdzenia~\ref{thm:local-continuity-11}
zachodzi $\int_X|f(t,x)-f(t_0,x)|\,\mathrm d\mu(x)\to0$
dla $t\to t_0$; w szczególności nie tylko wartości całek,
lecz także całe funkcje podcałkowe zależą ciągle od parametru
w normie $L^1$.
\end{wniosek}

\begin{przyklad}[Brak potrzeby globalnej dominacji]
Niech $f(t,x)=e^{-tx}$ dla $x>0$ i $t>0$. Jeśli $t_0>0$,
to dla $t\in(t_0/2,3t_0/2)$ zachodzi
$0\le f(t,x)\le e^{-t_0x/2}$, a funkcja po prawej stronie
jest całkowalna na $(0,\infty)$. Zatem
$F(t)=\int_{(0,\infty)}e^{-tx}\,\mathrm d\lambda_1(x)=1/t$
jest ciągła na $(0,\infty)$. Nie istnieje jednak całkowalna
dominanta dla całej rodziny $\{f(t,\cdot):t>0\}$, gdyż
$\sup_{t>0} e^{-tx}=1$ dla każdego $x>0$.
\end{przyklad}

\section{Jednostajność zbieżności całek w parametrze}

\begin{twierdzenie}[Zbieżność jednostajna po scałkowaniu]
\label{thm:uniform-parameter-11}
Niech $T$ będzie niepustym zbiorem, $f_n,f:T\times X\to\mathbb C$,
a dla każdego $t\in T$ funkcje $f_n(t,\cdot)$ i $f(t,\cdot)$
niech będą mierzalne. Załóżmy, że istnieją mierzalne funkcje
$h_n:X\to[0,\infty)$ i $g\in L^1(\mu)$ takie, że
$|f_n(t,x)-f(t,x)|\le h_n(x)\le g(x)$
dla wszystkich $t\in T$, wszystkich $n$ oraz prawie wszystkich
$x$ poza wspólnym zbiorem zerowym, a także $h_n(x)\to0$
prawie wszędzie. Załóżmy dodatkowo, że dla pewnego
$t_*\in T$ funkcja $f(t_*,\cdot)$ jest całkowalna oraz że
$|f(t,x)-f(t_*,x)|\le g_0(x)$ dla $t\in T$ i pewnego
$g_0\in L^1(\mu)$, tak aby wszystkie występujące całki były
określone. Wtedy
\[
\sup_{t\in T}\left|\int_X f_n(t,x)\,\mathrm d\mu(x)
-\int_Xf(t,x)\,\mathrm d\mu(x)\right|\longrightarrow0.
\]
\end{twierdzenie}
\begin{proof}
Dla każdego $t\in T$ z nierówności dla całki wynika
$|\int_X(f_n-f)(t,x)\,\mathrm d\mu(x)|\le\int_Xh_n\,\mathrm d\mu(x)$.
Oszacowanie po prawej stronie nie zależy od $t$.
Twierdzenie Lebesgue'a o zbieżności zmajoryzowanej daje
$\int_Xh_n\,\mathrm d\mu(x)\to0$. Przechodząc do supremum
po $t$, otrzymujemy tezę.
\end{proof}

\begin{uwaga}[Mierzalność supremum nie jest automatyczna]
Jeśli chciałoby się wybrać $h_n(x)=\sup_{t\in T}|f_n(t,x)-f(t,x)|$,
trzeba najpierw sprawdzić mierzalność tej funkcji. Wystarcza
na przykład, aby $T$ było separowalną przestrzenią metryczną,
a odwzorowania $t\mapsto f_n(t,x)-f(t,x)$ były ciągłe dla
wszystkich $x$: wtedy supremum po $T$ równa się supremum
po ustalonym przeliczalnym podzbiorze gęstym. Jeśli nie
mamy tego warunku, lepiej założyć istnienie mierzalnej
majoranty $h_n$ bez odwoływania się do supremum.
\end{uwaga}

\section{Różniczkowanie pod znakiem całki}

\begin{twierdzenie}[Pierwsza pochodna — wersja lokalna]
\label{thm:local-differentiation-11}
Niech $I\subseteq\mathbb R$ będzie przedziałem otwartym,
$t_0\in I$, a $f:I\times X\to\mathbb C$. Załóżmy, że:
\begin{enumerate}[label=\textup{(\roman*)}]
\item dla każdego $t\in I$ przekrój $f(t,\cdot)$ jest mierzalny;
\item dla prawie wszystkich $x$ funkcja $t\mapsto f(t,x)$ jest
      różniczkowalna na $I$;
\item istnieje $r>0$ z $[t_0-r,t_0+r]\subset I$ i funkcja
      $g\in L^1(\mu)$, $g\ge0$, taka, że
      $|\partial_t f(t,x)|\le g(x)$ dla wszystkich
      $t\in[t_0-r,t_0+r]$ poza wspólnym zbiorem zerowym;
\item funkcja $f(t_0,\cdot)$ jest całkowalna.
\end{enumerate}
Wówczas w otoczeniu $t_0$ funkcja
$F(t)=\int_Xf(t,x)\,\mathrm d\mu(x)$ jest dobrze określona,
różniczkowalna w $t_0$ i spełnia
\begin{equation}\label{eq:diff-param-11}
F'(t_0)=\int_X\partial_t f(t_0,x)\,\mathrm d\mu(x).
\end{equation}
\end{twierdzenie}
\begin{proof}
Poza wspólnym zbiorem zerowym ustalmy $x$ oraz $t$
i zastosujmy twierdzenie o wartości średniej do
rzeczywistej funkcji
$s\mapsto\operatorname{Re}(\overline{u}f(s,x))$,
gdzie $u\in\mathbb C$, $|u|=1$. Jej pochodna
ma moduł co najwyżej $g(x)$, zatem
$|\operatorname{Re}(\overline{u}(f(t,x)-f(t_0,x)))|
\le |t-t_0|g(x)$. Biorąc supremum po
$|u|=1$, otrzymujemy dokładnie
$|f(t,x)-f(t_0,x)|\le |t-t_0|g(x)$.
Argument nie wymaga ciągłości pochodnej ani
nieuzasadnionego zastosowania wzoru Newtona--Leibniza.
Stąd $f(t,\cdot)$ jest całkowalna dla $|t-t_0|<r$.
Dla $0<|h|<r$ iloraz
$q_h(x)=(f(t_0+h,x)-f(t_0,x))/h$ jest mierzalny
jako różnica funkcji mierzalnych; poza zbiorem zerowym
$q_h(x)\to\partial_tf(t_0,x)$ przy $h\to0$,
a $|q_h(x)|\le g(x)$. Pochodna $x\mapsto\partial_tf(t_0,x)$
jest mierzalna jako granica ilorazów różnicowych, na przykład
dla ciągu $h=1/n$ dostatecznie dużych $n$.
Twierdzenie o zbieżności zmajoryzowanej daje
$\int_Xq_h(x)\,\mathrm d\mu(x)\to
\int_X\partial_tf(t_0,x)\,\mathrm d\mu(x)$,
gdy $h\to0$. Z liniowości całki
$\int_Xq_h\,\mathrm d\mu(x)=(F(t_0+h)-F(t_0))/h$,
co jest dokładnie definicją pochodnej $F$.
\end{proof}

\begin{uwaga}[Dlaczego samo istnienie pochodnej nie wystarcza]
Dla $f(t,x)=t\chi_{(0,1/t)}(x)$ z pierwszego przykładu
nie można nawet otrzymać ciągłości całki w $t=0$.
Dominacja \emph{ilorazów różnicowych} lub lokalna dominacja
pochodnej są niezależnymi, istotnymi warunkami.
\end{uwaga}

\begin{twierdzenie}[Pochodne wyższego rzędu]
\label{thm:higher-differentiation-11}
Niech $I\subseteq\mathbb R$ będzie otwartym przedziałem,
$m\ge1$ liczbą naturalną, a $f:I\times X\to\mathbb C$.
Załóżmy, że dla każdego $t\in I$ funkcja $f(t,\cdot)$ jest mierzalna,
a dla prawie wszystkich $x$ odwzorowanie $t\mapsto f(t,x)$
ma pochodne do rzędu $m$ włącznie, wszystkie ciągłe na $I$.
Załóżmy ponadto, że dla każdego zwartego przedziału
$J\subset I$ oraz każdego $1\le k\le m$
istnieje $g_{J,k}\in L^1(\mu)$ taka, że
$|\partial_t^k f(t,x)|\le g_{J,k}(x)$ dla $t\in J$
poza wspólnym zbiorem zerowym, a $f(t_*,\cdot)$ jest
całkowalna dla pewnego $t_*\in I$.
Wtedy $F(t)=\int_Xf(t,x)\,\mathrm d\mu(x)$ należy
do $C^m(I)$ oraz dla wszystkich $t\in I$ i $1\le k\le m$
\begin{equation}\label{eq:higher-param-11}
F^{(k)}(t)=\int_X\partial_t^k f(t,x)\,\mathrm d\mu(x).
\end{equation}
\end{twierdzenie}
\begin{proof}
Dla $t\in I$ odcinek o końcach $t_*$ i $t$ jest zwarty
i zawiera się w $I$. Oszacowanie pierwszej pochodnej
zapewnia $|f(t,x)|\le|f(t_*,x)|+|t-t_*|g_{J,1}(x)$,
więc $f(t,\cdot)$ jest całkowalna. Dla ustalonego
$t\in I$ zastosujmy twierdzenie~\ref{thm:local-differentiation-11}
w małym otoczeniu $t$: otrzymujemy istnienie $F'$ i wzór
z $k=1$. Gdy $1\le k<m$, funkcja $\partial_t^kf(t,\cdot)$
jest mierzalna jako punktowa granica ilorazów różnicowych.
Dzięki majorantom $g_{J,k}$ jest całkowalna. Stosując
twierdzenie~\ref{thm:local-differentiation-11} do
$\partial_t^kf$, otrzymujemy pochodną $(k+1)$-szego rzędu.
Ciągłość każdej funkcji $t\mapsto\int_X\partial_t^kf(t,x)\,\mathrm d\mu(x)$
wynika z twierdzenia~\ref{thm:local-continuity-11},
gdyż $\partial_t^kf(t,x)$ jest punktowo ciągła w $t$
i lokalnie zdominowana. Zatem $F\in C^m(I)$.
\end{proof}

\begin{przyklad}[Całka Laplace'a i pochodne]
\label{ex:laplace-11}
Dla $t>0$ kładziemy $F(t)=\int_{(0,\infty)}e^{-tx}\,\mathrm d\lambda_1(x)$.
Mamy $F(t)=1/t$. Ponieważ
$\partial_t^k e^{-tx}=(-x)^k e^{-tx}$, dla
$t$ w zwartym podprzedziale $(0,\infty)$ funkcja
$x^k e^{-tx}$ jest zdominowana przez
$x^ke^{-ax}$, gdzie $a>0$ zależy od podprzedziału.
Ta ostatnia funkcja jest całkowalna, więc
\[
\int_{(0,\infty)}x^ke^{-tx}\,\mathrm d\lambda_1(x)
=\frac{k!}{t^{k+1}}\qquad(k\in\mathbb N_0).
\]
W szczególności wzór na momenty rozkładu wykładniczego
wynika z różniczkowania pod znakiem całki, a nie z
wielokrotnego całkowania przez części.
\end{przyklad}

\section{Zbieżność według miary zamiast dominacji punktowej}

\begin{twierdzenie}[Kryterium Vitaliego dla parametru]
\label{thm:param-vitali-11}
Niech $(X,\mathcal F,\mu)$ będzie przestrzenią z miarą
skończoną, $(T,d)$ przestrzenią metryczną, $t_0\in T$,
a $f:T\times X\to\mathbb C$ rodziną funkcji o mierzalnych
przekrojach. Załóżmy, że istnieje otoczenie $U$ punktu
$t_0$ takie, że rodzina modułów
$\{|f(t,\cdot)|:t\in U\}$ jest jednostajnie całkowalna
w sensie Definicji~\ref{def:uniform-integrability}
oraz dla każdego $\varepsilon>0$ zachodzi
\[
\mu(\{x\in X:|f(t,x)-f(t_0,x)|>\varepsilon\})
\longrightarrow0\qquad(t\to t_0).
\]
Jest to naturalne rozszerzenie zbieżności według miary
z Definicji~\ref{def:convergence-measure} na funkcje
zespolone. Wówczas
$\|f(t,\cdot)-f(t_0,\cdot)\|_{L^1(\mu)}\to0$,
a więc funkcja $t\mapsto\int_X f(t,x)\,\mathrm d\mu(x)$
jest ciągła w $t_0$.
\end{twierdzenie}
\begin{proof}
Załóżmy przeciwnie, że granica norm nie jest zerem.
Istnieją wówczas $\varepsilon>0$ i ciąg $t_n\to t_0$,
$t_n\in U$, dla którego
$\|f(t_n,\cdot)-f(t_0,\cdot)\|_1\ge\varepsilon$.
Funkcje $f(t_n,\cdot)$ zbiegają według miary do
$f(t_0,\cdot)$, a rodziny ich części rzeczywistych i urojonych
są jednostajnie całkowalne. Twierdzenie Vitaliego
z rozdziału~\ref{ch:twierdzenia-graniczne}
(dowodzone osobno dla części rzeczywistych i urojonych)
implikuje zbieżność w $L^1$, sprzeczną z wybraną
nierównością. W szczególności wszystkie całki
są określone i ich ciągłość wynika z nierówności
$|\int_X(f(t,x)-f(t_0,x))\,\mathrm d\mu(x)|
\le\|f(t,\cdot)-f(t_0,\cdot)\|_1$.
\end{proof}

\begin{uwaga}
Jednostajna całkowalność nie wymaga wspólnej
całkowalnej dominanty. W przeciwieństwie do twierdzenia
o dominacji lokalnej założenie o zbieżności jest tutaj
formułowane według miary, a nie punktowo. Dla $\mu(X)=\infty$
potrzebna jest dodatkowa kontrola ogonów; sam warunek
jednostajnej całkowalności na małych zbiorach nie wystarcza.
\end{uwaga}

\section{Zmienna dziedzina całkowania}

\begin{twierdzenie}[Całkowanie po zależnej od parametru dziedzinie]
\label{thm:moving-domains-11}
Niech $(T,d)$ będzie przestrzenią metryczną, $t_0\in T$,
a $A_t\in\mathcal F$ dla $t\in T$. Niech $u:T\times X\to\mathbb C$
ma mierzalne przekroje $u(t,\cdot)$ i załóżmy, że dla prawie
wszystkich $x$ zachodzi
$\chi_{A_t}(x)u(t,x)\to\chi_{A_{t_0}}(x)u(t_0,x)$
dla $t\to t_0$. Jeżeli w pewnym otoczeniu $U$ punktu $t_0$
istnieje nieujemna funkcja $g\in L^1(\mu)$ taka, że
\emph{dla każdego} $t\in U$ nierówność
$|\chi_{A_t}(x)u(t,x)|\le g(x)$
zachodzi dla $\mu$-prawie każdego $x$ (zbiór wyjątkowy
może zależeć od $t$), to
$t\mapsto\int_{A_t}u(t,x)\,\mathrm d\mu(x)$
jest ciągła w $t_0$.
\end{twierdzenie}
\begin{proof}
Definiujemy na całym $X$ funkcję
$f(t,x)=u(t,x)$ dla $x\in A_t$ i $f(t,x)=0$
dla $x\notin A_t$. Jest ona mierzalna w $x$ dla każdego $t$;
wyrażenia $\chi_{A_t}u$ należy tak interpretować,
aby uniknąć nieoznaczonych iloczynów $0\cdot\infty$.
Ustalmy ciąg $t_n\to t_0$. Po pominięciu skończenie
wielu wyrazów możemy założyć $t_n\in U$. Dla każdego
$n\ge0$ wybierzmy zbiór zerowy $N_n$, poza którym
$|f(t_n,x)|\le g(x)$; w tym zapisie $t_0$ odpowiada
$n=0$. Zbiór $N=\bigcup_{n\ge0}N_n$ ma miarę zero.
Poza $N$ oraz zbiorem wyjątkowym z założenia o zbieżności
punktowej zachodzi jednocześnie $f(t_n,x)\to f(t_0,x)$
i $|f(t_n,x)|\le g(x)$ dla wszystkich $n\ge0$.
Twierdzenie Lebesgue'a o zbieżności zmajoryzowanej
zapewnia więc zbieżność całek wzdłuż każdego ciągu
$t_n\to t_0$. Ciągowa charakteryzacja ciągłości
w przestrzeni metrycznej kończy dowód. W ten sposób
nie użyliśmy silniejszego, niewymaganego założenia
o wspólnym zbiorze wyjątkowym dla całej rodziny parametrów.
\end{proof}

\begin{przyklad}[Przedziały ruchome i brak dominacji]
Niech $u(t,x)=1$ na $(0,\infty)$ i $A_t=(0,t)$ dla $t>0$.
Dla $t$ w otoczeniu $t_0>0$ wszystkie $A_t$ mieszczą się
w $(0,2t_0)$; funkcja $g=\chi_{(0,2t_0)}$ jest
całkowalną dominantą, a dla każdego $x\ne t_0$
zachodzi $\chi_{A_t}(x)\to\chi_{A_{t_0}}(x)$.
Stąd $\int_{A_t}1\,\mathrm d\lambda_1(x)=t$ zależy ciągle od $t$.
Ten sam argument przestaje działać, jeśli dziedziny rozrastają
się tak, że masa ucieka poza każdy zbiór skończonej miary.
\end{przyklad}

\section{Całki niewłaściwe i jednostajna kontrola ogonów}

\begin{twierdzenie}[Kryterium jednostajnej kontroli ogonów]
\label{thm:tails-11}
Niech $t_0\in T$, gdzie $(T,d)$ jest przestrzenią metryczną,
a $(X,\mathcal F,\mu)$ przestrzenią z miarą. Załóżmy, że
istnieje wyczerpanie $X_1\subseteq X_2\subseteq\dots$,
$\bigcup_{n\ge1}X_n=X$, zbiory $X_n\in\mathcal F$
o skończonej mierze, oraz $f:T\times X\to\mathbb C$,
dla której $f(t,\cdot)$ jest całkowalna na każdym $X_n$.
Załóżmy ponadto, że $t\mapsto\int_{X_n}f(t,x)\,\mathrm d\mu(x)$
jest ciągła w $t_0$ dla każdego $n$, oraz że istnieje
otoczenie $U$ punktu $t_0$ takie, że
\begin{equation}\label{eq:uniform-tails-11}
\lim_{n\to\infty}\sup_{t\in U}
\int_{X\setminus X_n}|f(t,x)|\,\mathrm d\mu(x)=0.
\end{equation}
Wtedy wszystkie $f(t,\cdot)$ dla $t\in U$ są całkowalne,
a $t\mapsto\int_Xf(t,x)\,\mathrm d\mu(x)$ jest ciągła w $t_0$.
\end{twierdzenie}
\begin{proof}
Warunek \eqref{eq:uniform-tails-11} gwarantuje, że dla
pewnego $n$ i wszystkich $t\in U$ całki modułów po
$X\setminus X_n$ są skończone; ich suma z całką po $X_n$
jest więc skończona. Ustalmy $\varepsilon>0$ i wybierzmy
$n$ tak duże, aby supremum w
\eqref{eq:uniform-tails-11} było mniejsze niż $\varepsilon/3$.
Dla $t\in U$ mamy
\begin{align*}
\left|\int_Xf(t,x)\,\mathrm d\mu(x)-\int_Xf(t_0,x)\,\mathrm d\mu(x)\right|
&\le\left|\int_{X_n}(f(t,x)-f(t_0,x))\,\mathrm d\mu(x)\right|\\
&\quad+\int_{X\setminus X_n}|f(t,x)|\,\mathrm d\mu(x)
+\int_{X\setminus X_n}|f(t_0,x)|\,\mathrm d\mu(x).
\end{align*}
Dwa ostatnie składniki są mniejsze niż $2\varepsilon/3$.
Pierwszy jest mniejszy niż $\varepsilon/3$, gdy $t$
jest dostatecznie blisko $t_0$ dzięki ciągłości
całki na $X_n$. To kończy dowód.
\end{proof}

\begin{uwaga}[Ograniczenie kryterium]
Warunek \eqref{eq:uniform-tails-11} jest wystarczający, lecz
niekonieczny dla ciągłości całki: przy funkcjach zmieniających
znak mogą występować znoszące się wkłady z ogona.
Kryterium jest użyteczne zwłaszcza dla całek niewłaściwych,
gdzie trudno znaleźć jedną majorantę całkowalną na całym $X$.
\end{uwaga}

\section{Transformata Laplace'a i funkcje holomorficzne}

\begin{twierdzenie}[Holomorficzność transformaty Laplace'a]
\label{thm:laplace-holo-11}
Niech $a\in\mathbb R$ i niech $h:(0,\infty)\to\mathbb C$
będzie mierzalna oraz
$\int_{(0,\infty)}e^{-\sigma x}|h(x)|\,\mathrm d\lambda_1(x)<\infty$
dla każdego $\sigma>a$. Dla $z\in\mathbb C$,
$\operatorname{Re}z>a$, określmy
$L_h(z)=\int_{(0,\infty)}e^{-zx}h(x)\,\mathrm d\lambda_1(x)$.
Wtedy $L_h$ jest holomorficzna w półpłaszczyźnie
$\{\operatorname{Re}z>a\}$ i dla każdego $k\in\mathbb N_0$
\begin{equation}\label{eq:laplace-derivative-11}
L_h^{(k)}(z)=\int_{(0,\infty)}(-x)^k e^{-zx}h(x)
\,\mathrm d\lambda_1(x).
\end{equation}
\end{twierdzenie}
\begin{proof}
Ustalmy $z_0$ z $\operatorname{Re}z_0>a$ i liczbę
$\sigma$ spełniającą $a<\sigma<\operatorname{Re}z_0$.
W pewnym otoczeniu $z_0$ zachodzi $\operatorname{Re}z>\sigma$.
Dla $k\ge0$ istnieje $C_{k,\sigma,z_0}<\infty$ taka, że
$x^k e^{-\operatorname{Re}z\,x}\le
C_{k,\sigma,z_0}e^{-\sigma x}$ w nieco mniejszym
otoczeniu $z_0$: wybieramy pośrednią liczbę
$\sigma'>\sigma$ z $\operatorname{Re}z>\sigma'$
i korzystamy z ograniczoności $x^ke^{-(\sigma'-\sigma)x}$.
Zatem całki po prawej stronie
\eqref{eq:laplace-derivative-11} są zbieżne bezwzględnie.

Aby udowodnić różniczkowalność zespoloną, dla małego
$w\in\mathbb C\setminus\{0\}$ rozważamy
\[
\frac{e^{-(z_0+w)x}-e^{-z_0x}}{w}
=-x\int_{[0,1]}e^{-(z_0+sw)x}\,\mathrm d\lambda_1(s).
\]
Tożsamość z prawą stroną jest klasyczną
całką po rzeczywistym odcinku $s\in[0,1]$,
a nie zamianą kolejności dwóch całkowań.
Moduł ilorazu jest ograniczony przez
$x e^{-\sigma'x}$ dla odpowiednio małych $w$,
co po pomnożeniu przez $|h(x)|$ jest zdominowane
przez stałą razy $e^{-\sigma x}|h(x)|$.
Dla każdego $x$ iloraz dąży do $-xe^{-z_0x}$.
Twierdzenie o zbieżności zmajoryzowanej daje
$L_h'(z_0)=-\int_{(0,\infty)}xe^{-z_0x}h(x)
\,\mathrm d\lambda_1(x)$. Powtarzając argument dla
$(-x)^kh(x)$, otrzymujemy wszystkie pochodne.
Ponieważ pochodna zespolona istnieje w każdym punkcie
półpłaszczyzny, $L_h$ jest holomorficzna.
\end{proof}

\begin{przyklad}[Funkcja Gamma]
\label{ex:gamma-param-11}
Dla $s>0$ definiujemy
$\Gamma(s)=\int_{(0,\infty)}x^{s-1}e^{-x}\,\mathrm d\lambda_1(x)$.
Całka jest skończona: przy zerze $x^{s-1}$ jest
całkowalne dla $s>0$, a w nieskończoności czynnik
$e^{-x}$ dominuje każdą potęgę $x$.
Dla każdego $s_0>0$ wybieramy $0<a<s_0<b$;
na $s\in[a,b]$ pochodne rzędu $k$ funkcji podcałkowej
wynoszą $x^{s-1}(\log x)^k e^{-x}$ i mają całkowalną
majorantę
\[
\chi_{(0,1]}(x)x^{a-1}|\log x|^k
+\chi_{(1,\infty)}(x)x^{b-1}(1+\log x)^k e^{-x}.
\]
W konsekwencji $\Gamma\in C^\infty((0,\infty))$ i
$\Gamma^{(k)}(s)=\int_{(0,\infty)}
 x^{s-1}(\log x)^k e^{-x}\,\mathrm d\lambda_1(x)$.
Całkując przez części na $[\varepsilon,R]$ i przechodząc
do granic $\varepsilon\downarrow0$, $R\uparrow\infty$,
dostajemy $\Gamma(s+1)=s\Gamma(s)$.
Ponadto $\Gamma(1)=1$, więc $\Gamma(n+1)=n!$ dla $n\ge0$.
\end{przyklad}

\section{Podsumowanie}
Dominacja lokalna pozwala przenosić ciągłość i
różniczkowalność pod znak całki bez narzucania
jednej globalnej majoranty. Jednostajna kontrola
błędów oraz ogonów pozwala badać rodziny funkcji,
dla których pojedyncza dominanta bywa niedostępna.
Wersja zespolona prowadzi do holomorficzności
transformaty Laplace'a i umożliwia obliczanie jej
pochodnych. W kolejnym rozdziale skonstruujemy
miary produktowe, dzięki którym można zamieniać
kolejność całkowania względem dwóch niezależnych
miar.

% END chapters/pl/measure_theory/11-rozdzial.tex


% BEGIN chapters/pl/measure_theory/12-rozdzial.tex
% ROZDZIAŁ 12 — do włączenia w main.tex
\chapter{Miary produktowe. Twierdzenia Tonellego i Fubiniego}
\label{ch:produkty-miary}

\section{Po co konstruujemy miarę produktową?}

Niech $(X,\mathcal F,\mu)$ i $(Y,\mathcal G,\nu)$ będą
przestrzeniami z miarami. Chcemy przypisać mierzalnemu
zbiorowi $E\subseteq X\times Y$ wielkość, która dla
prostokąta mierzalnego $A\times B$ wynosi $\mu(A)\nu(B)$.
Potrzebujemy następnie twierdzenia, które pozwoli obliczyć
całkę po $X\times Y$ jako dwie kolejne całki.

Już na etapie formułowania pojawiają się trzy pytania:
\emph{(i)} jaka jest odpowiednia $\sigma$-algebra
na $X\times Y$; \emph{(ii)} czy przypisanie
$A\times B\mapsto\mu(A)\nu(B)$ wyznacza miarę;
\emph{(iii)} kiedy można zmieniać kolejność całkowania.
Będziemy rozróżniać trzy odpowiedzi: konstrukcję miary
produktowej, twierdzenie Tonellego dla funkcji nieujemnych
i twierdzenie Fubiniego dla funkcji całkowalnych.

\section{Przestrzeń produktowa i przekroje}

Iloczyn $\sigma$-algebr $\mathcal F\otimes\mathcal G$
został zdefiniowany w Definicji~\ref{def:product-sigma}.
Zbiory $A\times B$ dla $A\in\mathcal F$ i $B\in\mathcal G$
nazywamy \emph{prostokątami mierzalnymi}.
Dla $E\subseteq X\times Y$ \emph{przekrojem pionowym}
przez $x\in X$ nazywamy $E_x=\{y\in Y:(x,y)\in E\}$,
a \emph{przekrojem poziomym} przez $y\in Y$ zbiór
$E^y=\{x\in X:(x,y)\in E\}$.
Dla funkcji $f:X\times Y\to\overline{\mathbb R}$ piszemy
$f_x(y)=f(x,y)$ i $f^y(x)=f(x,y)$.

\begin{lemat}[Mierzalność wszystkich przekrojów]
\label{lem:sections-measurable-12}
Niech $E\in\mathcal F\otimes\mathcal G$.
Wtedy $E_x\in\mathcal G$ dla każdego $x\in X$
oraz $E^y\in\mathcal F$ dla każdego $y\in Y$.
Jeżeli $f:(X\times Y,\mathcal F\otimes\mathcal G)
\to(\overline{\mathbb R},\mathcal B(\overline{\mathbb R}))$
jest mierzalna, to wszystkie funkcje $f_x$ oraz $f^y$
są mierzalne.
\end{lemat}
\begin{proof}
Ustalmy $x\in X$. Niech $\mathcal C_x$ oznacza rodzinę
tych zbiorów $E\subseteq X\times Y$, których przekrój
$E_x$ należy do $\mathcal G$. Ponieważ
$(X\times Y)_x=Y$, $(E^c)_x=Y\setminus E_x$ oraz
$(\bigcup_nE_n)_x=\bigcup_n(E_n)_x$, rodzina
$\mathcal C_x$ jest $\sigma$-algebrą. Dla prostokąta
$A\times B$ przekrój ma postać $B$, gdy $x\in A$,
oraz $\varnothing$, gdy $x\notin A$. Zatem
$\mathcal C_x$ zawiera wszystkie prostokąty i
$\mathcal F\otimes\mathcal G\subseteq\mathcal C_x$.
Argument dla $E^y$ jest symetryczny. Wreszcie
$\{y:f_x(y)\in H\}=(f^{-1}(H))_x$ dla każdego
zbioru borelowskiego $H\subseteq\overline{\mathbb R}$;
poprzednia część dowodu zapewnia jego mierzalność.
Analogicznie postępujemy z $f^y$.
\end{proof}

\begin{uwaga}[Nie odwracamy twierdzenia o przekrojach]
Z mierzalności wszystkich przekrojów nie wynika
mierzalność zbioru w $\mathcal F\otimes\mathcal G$.
Kontrprzykład z podzbiorem przekątnej podamy w dalszej
części rozdziału. Mierzalność przekrojów jest
\emph{konsekwencją} mierzalności produktowej, a nie
jej ogólnym kryterium.
\end{uwaga}

\section{Lemat o podwójnym liczeniu miary zbioru}

W tym miejscu nie mamy jeszcze zdefiniowanej miary na
$X\times Y$. Najpierw udowodnimy, że wielkości
$\int_X\nu(E_x)\,\mathrm d\mu(x)$ oraz
$\int_Y\mu(E^y)\,\mathrm d\nu(y)$ mają sens
i są równe, gdy obie miary są $\sigma$-skończone.
Argument jest podstawą konstrukcji miary produktowej.

\begin{lemat}[Przekroje na prostokącie skończonej miary]
\label{lem:finite-rectangle-sections-12}
Niech $C\in\mathcal F$, $D\in\mathcal G$ oraz
$\mu(C)<\infty$, $\nu(D)<\infty$. Jeżeli
$E\in\mathcal F\otimes\mathcal G$ i $E\subseteq C\times D$,
to funkcje $x\mapsto\nu(E_x)$ oraz
$y\mapsto\mu(E^y)$ są mierzalne i
\begin{equation}\label{eq:set-fubini-finite-12}
\int_C\nu(E_x)\,\mathrm d\mu(x)
=\int_D\mu(E^y)\,\mathrm d\nu(y).
\end{equation}
\end{lemat}
\begin{proof}
Rozpatrujemy przestrzeń mierzalną
$R=C\times D$ z $\sigma$-algebrą
$\mathcal H=(\mathcal F\otimes\mathcal G)|_R$.
Jest ona generowana przez prostokąty
$A\times B$ z $A\in\mathcal F|_C$ i
$B\in\mathcal G|_D$. Niech $\mathcal M$ będzie
rodziną tych $E\in\mathcal H$, dla których
obie funkcje przekrojowe są mierzalne oraz zachodzi
\eqref{eq:set-fubini-finite-12}.

\emph{Krok 1: algebra prostokątów.}
Dla $E=A\times B\subseteq R$ mamy
$\nu(E_x)=\nu(B)\chi_A(x)$ i
$\mu(E^y)=\mu(A)\chi_B(y)$. Obie całki
są równe $\mu(A)\nu(B)$.
Ta sama własność obowiązuje dla skończonej
\emph{rozłącznej} sumy prostokątów, z addytywności
miar i liniowości całki. Każdą skończoną sumę
prostokątów można przedstawić jako skończoną
rozłączną sumę prostokątów: wystarczy rozbić $C$
i $D$ na skończone rozłączne podziały wyznaczone
przez występujące pierwsze i drugie czynniki.
Zatem $\mathcal M$ zawiera algebrę $\mathcal A$
skończonych sum prostokątów.

\emph{Krok 2: ciągi rosnące.}
Jeżeli $E_n\in\mathcal M$ i $E_n\uparrow E$,
to dla każdego $x$ zachodzi $(E_n)_x\uparrow E_x$.
Ciągłość miary $\nu$ z dołu daje
$\nu(E_x)=\lim_n\nu((E_n)_x)$; zatem funkcja
$x\mapsto\nu(E_x)$ jest mierzalna jako granica
funkcji mierzalnych. Podobnie dla drugiego przekroju.
Twierdzenie Beppo Leviego zastosowane po obu stronach
\eqref{eq:set-fubini-finite-12} daje tę samą
równość dla $E$.

\emph{Krok 3: ciągi malejące.}
Jeżeli $E_n\in\mathcal M$ i $E_n\downarrow E$,
to $\nu((E_n)_x)\downarrow\nu(E_x)$, ponieważ
$\nu((E_1)_x)\le\nu(D)<\infty$. Ponadto funkcje
$\nu((E_n)_x)$ są zdominowane przez
$\nu(D)\chi_C(x)$, która ma skończoną
całkę $\mu(C)\nu(D)$ (w przypadku wartości zero
stosujemy zwykłą konwencję dla całki funkcji zerowej).
Twierdzenie o zbieżności zmajoryzowanej daje zbieżność
całek; analogicznie z drugiej strony, z dominantą
$\mu(C)\chi_D$. Stąd równość dla $E$.

Z kroków 2--3 wynika, że $\mathcal M$ jest klasą
monotoniczną. Zawiera algebrę $\mathcal A$, więc
na mocy twierdzenia o klasie monotonicznej
$\mathcal H=\sigma(\mathcal A)\subseteq\mathcal M$.
To kończy dowód bez użycia Fubiniego lub Tonellego.
\end{proof}

\begin{twierdzenie}[Tożsamość przekrojowa przy $\sigma$-skończoności]
\label{thm:section-identity-12}
Niech $\mu$ i $\nu$ będą $\sigma$-skończone.
Dla każdego $E\in\mathcal F\otimes\mathcal G$
funkcje $x\mapsto\nu(E_x)$ i $y\mapsto\mu(E^y)$
są mierzalne oraz
\begin{equation}\label{eq:section-identity-12}
\int_X\nu(E_x)\,\mathrm d\mu(x)
=\int_Y\mu(E^y)\,\mathrm d\nu(y)
\quad\text{w }[0,\infty].
\end{equation}
\end{twierdzenie}
\begin{proof}
Wybieramy rosnące wyczerpania
$X_n\uparrow X$ i $Y_n\uparrow Y$ zbiorami mierzalnymi
o skończonych miarach. Można je otrzymać, zastępując
zbiory z dowolnych wyczerpań ich skończonymi sumami.
Niech $R_n=X_n\times Y_n$ oraz $E_n=E\cap R_n$.
Lemat~\ref{lem:finite-rectangle-sections-12}
zastosowany do $E_n$ zapewnia mierzalność obu
funkcji przekrojowych dla $E_n$ i równość odpowiednich
całek. Ponieważ $E_n\uparrow E$, dla każdego $x$
zachodzi $(E_n)_x\uparrow E_x$, a dla każdego $y$
zachodzi $(E_n)^y\uparrow E^y$. Z ciągłości miary
z dołu wynika
$\nu(E_x)=\lim_n\nu((E_n)_x)$ i
$\mu(E^y)=\lim_n\mu((E_n)^y)$. Granice są
mierzalne. Stosując dwukrotnie twierdzenie
Beppo Leviego do całek po $X$ i $Y$, uzyskujemy
\eqref{eq:section-identity-12}.
\end{proof}

\section{Istnienie i jednoznaczność miary produktowej}

\begin{twierdzenie}[Konstrukcja miary produktowej]
\label{thm:product-measure-12}
Niech $(X,\mathcal F,\mu)$ i $(Y,\mathcal G,\nu)$
będą przestrzeniami z miarami $\sigma$-skończonymi.
Istnieje dokładnie jedna miara $\pi$ na
$(X\times Y,\mathcal F\otimes\mathcal G)$ taka, że
\begin{equation}\label{eq:rect-product-12}
\pi(A\times B)=\mu(A)\nu(B),
\qquad A\in\mathcal F,\quad B\in\mathcal G,
\end{equation}
przy umowie $0\cdot\infty=\infty\cdot0=0$.
Oznaczamy ją przez $\mu\otimes\nu$.
Dla każdego $E\in\mathcal F\otimes\mathcal G$ zachodzi
\begin{equation}\label{eq:measure-section-12}
(\mu\otimes\nu)(E)=
\int_X\nu(E_x)\,\mathrm d\mu(x)
=\int_Y\mu(E^y)\,\mathrm d\nu(y).
\end{equation}
\end{twierdzenie}
\begin{proof}
\emph{Istnienie.}
Definiujemy $\pi(E)$ jako pierwszą całkę
w \eqref{eq:measure-section-12}; na mocy
\ref{thm:section-identity-12} całkowana funkcja jest
mierzalna i można też użyć drugiej całki.
Oczywiście $\pi(\varnothing)=0$.
Jeśli $(E_j)_{j\ge1}$ jest ciągiem parami
rozłącznych zbiorów w $\mathcal F\otimes\mathcal G$,
to dla każdego $x$ przekroje $(E_j)_x$ są parami
rozłączne i
\[
\nu\left(\left(\bigcup_{j\ge1}E_j\right)_x\right)
=\sum_{j\ge1}\nu((E_j)_x).
\]
Stosując twierdzenie Beppo Leviego do sum częściowych,
ostatecznie otrzymujemy
$\pi(\bigcup_jE_j)=\sum_j\pi(E_j)$.
Zatem $\pi$ jest miarą.
Dla $E=A\times B$ funkcja przekrojowa jest równa
$\nu(B)\chi_A$ (w rozumieniu funkcji przyjmującej
wartość zero poza $A$), więc
$\pi(A\times B)=\mu(A)\nu(B)$ także w przypadkach
$\mu(A)=0$ lub $\nu(B)=0$.

\emph{Jednoznaczność.}
Niech $\rho$ będzie inną miarą spełniającą
\eqref{eq:rect-product-12}. Z wyczerpań użytych
w poprzednim twierdzeniu wybieramy prostokąty
$R_n=X_n\times Y_n\uparrow X\times Y$, gdzie
$\mu(X_n)\nu(Y_n)<\infty$.
Na $\sigma$-algebrze śladu na $R_n$ miary
$\rho$ i $\pi$ są skończone i zgadzają się na
prostokątach $A\times B\subseteq R_n$.
Przecięcia dwóch prostokątów są prostokątami,
a prostokąty generują $\sigma$-algebrę śladu;
twierdzenie $\pi$--$\lambda$ z rozdziału 1 daje
równość miar na całym śladzie.
W szczególności dla dowolnego
$E\in\mathcal F\otimes\mathcal G$ mamy
$\rho(E\cap R_n)=\pi(E\cap R_n)$.
Przechodząc do granicy $n\to\infty$ przez
ciągłość miar z dołu, dostajemy $\rho(E)=\pi(E)$.
\end{proof}

\begin{uwaga}[Co daje $\sigma$-skończoność?]
W powyższym dowodzie jest potrzebna do
wyczerpania przestrzeni prostokątami skończonej
miary, a w konsekwencji do jednoznaczności.
Bez $\sigma$-skończoności równość iterowanych całek
i jednoznaczność rozszerzenia nie muszą zachodzić.
Nie należy traktować jej jako dekoracyjnego założenia.
\end{uwaga}

\section{Twierdzenie Tonellego}

\begin{twierdzenie}[Tonelli]
\label{thm:tonelli-12}
Niech $\mu$ i $\nu$ będą miarami $\sigma$-skończonymi,
a $f:X\times Y\to[0,\infty]$ funkcją
$\mathcal F\otimes\mathcal G$-mierzalną.
Wtedy funkcje
\[
x\longmapsto\int_Y f(x,y)\,\mathrm d\nu(y),
\qquad
y\longmapsto\int_X f(x,y)\,\mathrm d\mu(x)
\]
są mierzalne, a trzy całki są równe
(wartość $+\infty$ jest dozwolona):
\begin{equation}\label{eq:tonelli-12}
\int_{X\times Y}f(x,y)\,\mathrm d(\mu\otimes\nu)(x,y)
=\int_X\left(\int_Yf(x,y)\,\mathrm d\nu(y)\right)
\mathrm d\mu(x)
=\int_Y\left(\int_Xf(x,y)\,\mathrm d\mu(x)\right)
\mathrm d\nu(y).
\end{equation}
\end{twierdzenie}
\begin{proof}
\emph{Krok 1: funkcje charakterystyczne.}
Dla $f=\chi_E$, $E\in\mathcal F\otimes\mathcal G$,
mamy $\int_Yf(x,y)\,\mathrm d\nu(y)=\nu(E_x)$.
Obie równości w \eqref{eq:tonelli-12} są dokładnie
wzorem \eqref{eq:measure-section-12}.

\emph{Krok 2: funkcje proste nieujemne.}
Jeśli $s=\sum_{j=1}^m a_j\chi_{E_j}$,
gdzie $a_j\ge0$, $E_j$ są mierzalne i parami
rozłączne, to liniowość całki skończonej sumy
funkcji nieujemnych daje
$\int_Y s(x,y)\,\mathrm d\nu(y)
=\sum_j a_j\nu((E_j)_x)$.
Wyrażenie po prawej jest mierzalne w $x$.
Całkując po $X$ i używając kroku 1,
otrzymujemy pierwszą równość dla $s$;
analogicznie postępujemy z odwrotną kolejnością.

\emph{Krok 3: dowolne funkcje nieujemne.}
Na mocy konstrukcji funkcji prostych z rozdziału 6
istnieje ciąg $s_n\uparrow f$ nieujemnych,
prostych i mierzalnych na $X\times Y$.
Dla każdego $x$ mamy $(s_n)_x\uparrow f_x$,
więc twierdzenie Beppo Leviego względem $\nu$
daje $\int_Y(s_n)_x\,\mathrm d\nu(y)\uparrow
\int_Yf_x\,\mathrm d\nu(y)$. Granica funkcji
mierzalnych jest mierzalna. Ponownie stosujemy
Beppo Leviego, tym razem względem $\mu$:
\[
\int_X\left(\int_Y f_x\,\mathrm d\nu(y)\right)\mathrm d\mu(x)
=\lim_n\int_X\left(\int_Y (s_n)_x\,\mathrm d\nu(y)\right)\mathrm d\mu(x)
=\lim_n\int_{X\times Y}s_n\,\mathrm d(\mu\otimes\nu)(x,y).
\]
Ostatnia granica jest równa
$\int_{X\times Y}f\,\mathrm d(\mu\otimes\nu)(x,y)$,
jeszcze raz z twierdzenia Beppo Leviego.
Argument z całkowaniem najpierw względem $\mu$
jest identyczny.
\end{proof}

\begin{wniosek}[Kryterium nieujemnej całkowalności]
\label{cor:tonelli-integrability-12}
W sytuacji twierdzenia~\ref{thm:tonelli-12}
niech $f\ge0$. Funkcja $f$ ma skończoną całkę
względem $\mu\otimes\nu$ wtedy i tylko wtedy,
gdy dowolna (a zatem każda) z iterowanych całek
w \eqref{eq:tonelli-12} jest skończona.
\end{wniosek}
\begin{proof}
Wszystkie trzy całki są równe jako elementy
$[0,\infty]$. Skończoność jednej jest zatem
równoważna skończoności pozostałych.
\end{proof}

\section{Twierdzenie Fubiniego dla funkcji całkowalnych}

\begin{twierdzenie}[Fubini]
\label{thm:fubini-12}
Niech $\mu$ i $\nu$ będą miarami $\sigma$-skończonymi,
a $f:X\times Y\to\mathbb R$ lub $\mathbb C$ funkcją
$\mathcal F\otimes\mathcal G$-mierzalną taką, że
\begin{equation}\label{eq:abs-fubini-12}
\int_{X\times Y}|f(x,y)|\,\mathrm d(\mu\otimes\nu)(x,y)<\infty.
\end{equation}
Wtedy:
\begin{enumerate}[label=\textup{(\roman*)}]
\item dla $\mu$-prawie każdego $x$ przekrój $f_x$ jest
      całkowalny względem $\nu$; analogicznie dla
      $\nu$-prawie każdego $y$ przekrój $f^y$ jest
      całkowalny względem $\mu$;
\item przyjmując wartość zero na odpowiednich
      mierzalnych zbiorach wyjątkowych otrzymujemy
      całkowalne funkcje
      $F(x)=\int_Y f(x,y)\,\mathrm d\nu(y)$ i
      $G(y)=\int_X f(x,y)\,\mathrm d\mu(x)$;
\item zachodzi
\begin{equation}\label{eq:fubini-12}
\int_{X\times Y}f(x,y)\,\mathrm d(\mu\otimes\nu)(x,y)
=\int_XF(x)\,\mathrm d\mu(x)
=\int_YG(y)\,\mathrm d\nu(y).
\end{equation}
\end{enumerate}
\end{twierdzenie}
\begin{proof}
\emph{Krok 1: całkowalność przekrojów prawie wszędzie.}
Stosując twierdzenie Tonellego do $|f|$,
otrzymujemy
\[
\int_X\left(\int_Y|f(x,y)|\,\mathrm d\nu(y)\right)
\mathrm d\mu(x)
=\int_{X\times Y}|f|\,\mathrm d(\mu\otimes\nu)(x,y)<\infty.
\]
Funkcja nieujemna $H(x)=\int_Y|f(x,y)|\,\mathrm d\nu(y)$
jest mierzalna i ma skończoną całkę. Zatem
$N=\{x:H(x)=\infty\}$ jest mierzalnym zbiorem
miary zero. Dla $x\notin N$ przekrój $f_x$
jest całkowalny; analogicznie otrzymujemy zbiór
wyjątkowy $M\subseteq Y$ miary zero.

\emph{Krok 2: mierzalność całek przekrojowych.}
Dla funkcji rzeczywistej zapisujemy $f=f^+-f^-$.
Funkcje $x\mapsto\int_Y f^+(x,y)\,\mathrm d\nu(y)$
i $x\mapsto\int_Y f^-(x,y)\,\mathrm d\nu(y)$
są mierzalne z Tonellego. Dla $x\notin N$
obie wartości są skończone, więc ich różnica
jest dobrze określona i mierzalna. Przedłużając
różnicę zerem na $N$, otrzymujemy mierzalną
funkcję $F$. Dla funkcji zespolonej powtarzamy
tę konstrukcję osobno dla części rzeczywistej
i urojonej. Analogicznie postępujemy z $G$.

\emph{Krok 3: całkowalność i równość całek.}
Dla $x\notin N$ nierówność trójkąta dla całki
daje $|F(x)|\le H(x)$, a na $N$ funkcja $F$
jest równa zero. Zatem $F\in L^1(\mu)$;
podobnie $G\in L^1(\nu)$.
Do funkcji $f^+$ i $f^-$ stosujemy Tonellego,
a następnie odejmujemy skończone liczby;
uzyskujemy pierwszą i drugą równość
\eqref{eq:fubini-12}. Dla funkcji zespolonej
stosujemy ten argument do obu części składowych.
\end{proof}

\begin{wniosek}[Kryterium praktyczne]
\label{cor:practical-fubini-12}
Jeśli dla mierzalnej funkcji $f:X\times Y\to\mathbb C$
możemy wykazać
$\int_X\bigl(\int_Y |f(x,y)|\,\mathrm d\nu(y)\bigr)
\mathrm d\mu(x)<\infty$, to $f$ jest całkowalna
na przestrzeni produktowej i wolno zamienić kolejność
całkowania w \eqref{eq:fubini-12}.
\end{wniosek}
\begin{proof}
Twierdzenie Tonellego zastosowane do $|f|$ zapewnia
równość podanej iterowanej całki i całki
$\int_{X\times Y}|f|\,\mathrm d(\mu\otimes\nu)(x,y)$.
Skończoność tej ostatniej oznacza spełnienie
warunków Fubiniego.
\end{proof}

\section{Przykłady rachunkowe}

\begin{przyklad}[Pole trójkąta bez całkowania po powierzchni]
Niech $T=\{(x,y)\in[0,1]^2:0\le y\le x\}$.
Przekrój pionowy $T_x$ jest przedziałem $[0,x]$
dla $x\in[0,1]$, a poza tym zbiorem jest pusty.
Dlatego Tonelli daje
\[
(\lambda_1\otimes\lambda_1)(T)=\int_{[0,1]}\left(
\int_{[0,1]}\chi_{\{y\le x\}}\,\mathrm d\lambda_1(y)
\right)\mathrm d\lambda_1(x)
=\int_{[0,1]}x\,\mathrm d\lambda_1(x)=\frac12.
\]
Z twierdzenia~\ref{thm:lebesgue-products-12} poniżej wynika także
$\lambda_2(T)=(\lambda_1\otimes\lambda_1)(T)=1/2$.
W odwrotnej kolejności przekrój poziomy ma długość
$1-y$, a całka wynosi
$\int_{[0,1]}(1-y)\,\mathrm d\lambda_1(y)=1/2$.
Nawet w tym elementarnym przykładzie widać dwie
różne parametryzacje tej samej miary.
\end{przyklad}

\begin{przyklad}[Całka z nieograniczoną funkcją nieujemną]
Na $(0,1)^2$ określmy $f(x,y)=\frac{1}{x+y}$.
Funkcja jest mierzalna i nieujemna; dlatego wolno
najpierw całkować po $y$, mimo że jeszcze nie
wiemy, czy podwójna całka jest skończona:
\[
\int_{(0,1)^2}\frac{1}{x+y}\,\mathrm d\lambda_2(x,y)
=\int_{(0,1)}\bigl(\log(1+x)-\log x\bigr)
\,\mathrm d\lambda_1(x)=2\log2.
\]
Ostatnia równość wynika z elementarnego
całkowania funkcji $\log(1+x)$ i $\log x$
na $(0,1)$ oraz przejścia do granicy przy zerze.
\end{przyklad}

\begin{przyklad}[Funkcja zmieniająca znak]
Niech $f(x,y)=e^{-x-y}(x-y)$ na $(0,\infty)^2$.
Ponieważ $|x-y|\le x+y$, Tonelli oraz wzory
$\int_{(0,\infty)}e^{-x}\,\mathrm d\lambda_1(x)=1$ i
$\int_{(0,\infty)}xe^{-x}\,\mathrm d\lambda_1(x)=1$
z rozdziału~\ref{ch:parametry-calki} dają
\begin{align*}
\int_{(0,\infty)^2}|f|\,
\mathrm d(\lambda_1\otimes\lambda_1)
&\le\int_{(0,\infty)^2}(x+y)e^{-x-y}
\,\mathrm d(\lambda_1\otimes\lambda_1)\\
&=1\cdot1+1\cdot1=2.
\end{align*}
Zatem wolno użyć Fubiniego. Zamiana $x$ z $y$
zmienia znak całkowanej funkcji i zachowuje
miarę produktową, a więc
$\int_{(0,\infty)^2}f\,\mathrm d(\lambda_1\otimes\lambda_1)=0$.
\end{przyklad}

\section{Dlaczego założenia twierdzeń są ważne?}

\begin{przyklad}[Dwie iterowane całki, dwa różne wyniki]
\label{ex:nonabs-fubini-12}
Niech $X=Y=\mathbb N$ z $\sigma$-algebrą potęgową
i miarą liczącą $\kappa$. Definiujemy
\[
a_{m,n}=\begin{cases}
1,&n=m,\\
-1,&n=m+1,\\
0,&\text{w pozostałych przypadkach}.
\end{cases}
\]
Dla każdego $m$ szereg $\sum_{n=1}^\infty a_{m,n}$
ma tylko dwa niezerowe wyrazy i wynosi zero. Natomiast
dla $n=1$ szereg $\sum_{m=1}^\infty a_{m,1}$
ma sumę $1$, a dla $n\ge2$ sumę zero.
Wobec tego
\[
\sum_{m=1}^\infty\left(\sum_{n=1}^\infty a_{m,n}\right)=0,
\qquad
\sum_{n=1}^\infty\left(\sum_{m=1}^\infty a_{m,n}\right)=1.
\]
Nie ma sprzeczności z Fubinim, ponieważ
$\sum_{m,n}|a_{m,n}|=\infty$.
W rzeczywistości $a^+$ i $a^-$ mają nieskończoną
całkę względem $\kappa\otimes\kappa$, więc
całka podpisana po całej przestrzeni nie istnieje.
\end{przyklad}

\begin{przyklad}[Bez $\sigma$-skończoności Tonelli może zawieść]
\label{ex:non-sigma-tonelli-12}
Niech $X=Y=[0,1]$, przy czym na $X$ bierzemy
$\mathcal F=\mathcal P([0,1])$ i miarę liczącą $\kappa$,
a na $Y$ miarę Lebesgue'a $\lambda_1$ ze zwykłą
$\sigma$-algebrą. Miara $\kappa$ nie jest
$\sigma$-skończona, gdyż jej zbiory skończonej
miary są skończone, a przeliczalna suma zbiorów
skończonych jest przeliczalna.
Zbiór przekątnej $D=\{(x,y):x=y\}$ jest
mierzalny w $\mathcal F\otimes\mathcal L_1$,
gdyż jest domknięty w zwykłej topologii produktu,
a $\mathcal B([0,1]^2)\subseteq\mathcal P([0,1])\otimes\mathcal L_1$:
każdy prostokąt otwarty należy do tej $\sigma$-algebry,
a otwarte podzbiory $[0,1]^2$ są przeliczalnymi sumami
takich prostokątów.
Dla każdego $x$ przekrój $D_x=\{x\}$ ma
$\lambda_1$-miarę zero, natomiast dla każdego $y$
przekrój $D^y=\{y\}$ ma miarę liczącą $1$.
Dlatego
\[
\int_X\left(\int_Y\chi_D(x,y)\,
\mathrm d\lambda_1(y)\right)\mathrm d\kappa(x)=0,
\qquad
\int_Y\left(\int_X\chi_D(x,y)\,
\mathrm d\kappa(x)\right)\mathrm d\lambda_1(y)=1.
\]
Nie można bez dodatkowych warunków konstruować
jednej miary produktowej przez dowolną z dwóch
kolejności całkowania. Ten przykład pokazuje
dokładnie miejsce, w którym używaliśmy
$\sigma$-skończoności w dowodzie.
\end{przyklad}

\section{Produkt miar Lebesgue'a i problem zupełności}

\begin{twierdzenie}[Produkt miar Lebesgue'a]
\label{thm:lebesgue-products-12}
Niech $m,n\ge1$. Miara $\lambda_m\otimes\lambda_n$
na $\mathcal L_m\otimes\mathcal L_n$ jest ograniczeniem
miary Lebesgue'a $\lambda_{m+n}$ do tej
$\sigma$-algebry. Ponadto uzupełnienie
$(\mathcal L_m\otimes\mathcal L_n,
\lambda_m\otimes\lambda_n)$ daje dokładnie
$(\mathcal L_{m+n},\lambda_{m+n})$.
W ogólności $\mathcal L_m\otimes\mathcal L_n$
nie jest zupełna.
\end{twierdzenie}
\begin{proof}
\emph{Krok 1: relacja między $\sigma$-algebrami.}
Dla $A\in\mathcal L_m$ istnieje zbiór borelowski
$B\subseteq\mathbb R^m$ taki, że
$A\mathbin\triangle B\subseteq N$ dla pewnego
borelowskiego zbioru zerowego $N$.
Analogicznie postępujemy z $C\in\mathcal L_n$.
Prostokąt $A\times C$ różni się od
borelowskiego prostokąta $B\times D$ co najwyżej
na podzbiorze
$(N\times\mathbb R^n)\cup(\mathbb R^m\times M)$,
gdzie $N$ i $M$ są borelowskie i mają miarę zero.
Zbiory te są $\lambda_{m+n}$-zerowe:
po przecięciu z ograniczonymi prostopadłościanami
wynika to z pokrywania małymi prostokątami
lub ze znanej już konstrukcji miary Lebesgue'a,
a następnie z przeliczalnej sumy po rosnących
prostopadłościanach. Stąd każdy prostokąt
$A\times C$ jest $\lambda_{m+n}$-mierzalny,
i $\mathcal L_m\otimes\mathcal L_n\subseteq\mathcal L_{m+n}$.

\emph{Krok 2: zgodność miar.}
Niech najpierw $Q\subseteq\mathbb R^m$ będzie ograniczonym
półotwartym prostopadłościanem. Dla borelowskich
$D\subseteq\mathbb R^n$ obie funkcje
$D\mapsto\lambda_{m+n}(Q\times D)$ oraz
$D\mapsto\lambda_m(Q)\lambda_n(D)$ są miarami
$\sigma$-skończonymi na $\mathcal B(\mathbb R^n)$;
zgadzają się na półotwartych prostopadłościanach
w $\mathbb R^n$, ponieważ ich produkty z $Q$ są
prostopadłościanami w $\mathbb R^{m+n}$.
Twierdzenie o jednoznaczności miar z rozdziału 2
daje równość dla każdego borelowskiego $D$.

Teraz ustalmy borelowski $D\subseteq\mathbb R^n$
o $\lambda_n(D)<\infty$. Obie funkcje
$B\mapsto\lambda_{m+n}(B\times D)$ oraz
$B\mapsto\lambda_n(D)\lambda_m(B)$ są
$\sigma$-skończonymi miarami borelowskimi na
$\mathbb R^m$ (gdy $\lambda_n(D)=0$, obie są
zerowe). Dzięki poprzedniemu akapitowi zgadzają
się na ograniczonych półotwartych prostopadłościanach
$B=Q$. Ponowne zastosowanie jednoznaczności
daje $\lambda_{m+n}(B\times D)=\lambda_m(B)\lambda_n(D)$
dla wszystkich borelowskich $B$ i $D$ ze
skończoną miarą drugiego czynnika. Dla dowolnego
borelowskiego $D$ wyczerpujemy go zbiorami
$D\cap[-k,k]^n$ i stosujemy ciągłość miary z dołu.
Tak otrzymujemy formułę także przy nieskończonych
miarach czynników, z umową $0\cdot\infty=0$.

Dla dowolnych $A,C$ jak wyżej wartość miary
prostokąta nie zmienia się po zastąpieniu go
odpowiednim prostokątem borelowskim, z wyjątkiem
przypadków nieskończonych miar, gdzie stosujemy
wyczerpanie obydwu czynników zbiorami ograniczonymi
i ciągłość miar z dołu. Zatem
$\lambda_{m+n}(A\times C)=\lambda_m(A)\lambda_n(C)$
(z konwencją dla iloczynu z zerem).
Jednoznaczność miary produktowej daje zgodność
miar na całej $\mathcal L_m\otimes\mathcal L_n$.

\emph{Krok 3: uzupełnienie.}
$\mathcal B(\mathbb R^{m+n})$ jest podzbiorem
$\mathcal L_m\otimes\mathcal L_n$, ponieważ
borelowska $\sigma$-algebra produktu euklidesowego
jest generowana przez prostokąty otwarte.
Każdy zbiór z $\mathcal L_{m+n}$ znajduje się
między dwoma zbiorami borelowskimi $B\subseteq E\subseteq C$
ze $\lambda_{m+n}(C\setminus B)=0$
(dzięki regularności miary Lebesgue'a i jej
zupełności). Zatem $E$ należy do uzupełnienia
$\mathcal L_m\otimes\mathcal L_n$.
W drugą stronę uzupełnienie pozostaje podzbiorem
$\mathcal L_{m+n}$, ponieważ ta ostatnia
$\sigma$-algebra jest zupełna. To kończy dowód.
\end{proof}

\begin{przyklad}[Niezmierzalny w produkcie podzbiór przekątnej]
\label{ex:diagonal-not-product-12}
Niech $V\subseteq[0,1]$ będzie zbiorem niemierzalnym
w sensie Lebesgue'a, np. zbiorem Vitaliego, i niech
$E=\{(x,x):x\in V\}\subseteq\mathbb R^2$.
Zbiór $E$ jest podzbiorem przekątnej
$D=\{(x,x):x\in[0,1]\}$, a $\lambda_2(D)=0$.
Zupełność $\lambda_2$ daje $E\in\mathcal L_2$.
Gdyby natomiast $E\in\mathcal L_1\otimes\mathcal L_1$,
to przeciwobraz $\Delta^{-1}(E)$ przez
$\Delta(x)=(x,x)$ byłby Lebesgue'a-mierzalny.
Rzeczywiście oba odwzorowania współrzędne
$\operatorname{pr}_j\circ\Delta$ są identycznością,
a więc $\Delta$ jest mierzalna jako odwzorowanie
do produktu przestrzeni mierzalnych. Jednak
$\Delta^{-1}(E)=V$, sprzeczność. Stąd
$\mathcal L_1\otimes\mathcal L_1\subsetneq\mathcal L_2$.
Co ciekawe, każdy pionowy i poziomy przekrój $E$
jest pusty lub jednoelementowy, a więc mierzalny:
to dodatkowo dowodzi, że nie można odwracać
lematu~\ref{lem:sections-measurable-12}.
\end{przyklad}

\begin{uwaga}[Pełne zbiory mierzalne i przekroje]
Z lematu~\ref{lem:sections-measurable-12} nie wynika,
że \emph{wszystkie} przekroje dowolnego zbioru
$E\in\mathcal L_{m+n}$ są mierzalne. Na przykład
dla niemierzalnego $V\subseteq[0,1]$ zbiór
$V\times\{0\}$ jest $\lambda_2$-mierzalny jako
podzbiór prostej poziomej miary zero, ale jego
przekrój przy $y=0$ jest równy $V$.
Wersja Fubiniego dla uzupełnionej miary mówi o
mierzalności i całkowalności przekrojów
\emph{prawie wszędzie}; do tego potrzebne jest
osobne uzasadnienie po zastąpieniu funkcji
przez mierzalnego przedstawiciela produktowego.
\end{uwaga}


\begin{twierdzenie}[Tonelli i Fubini dla uzupełnionej miary produktowej]
\label{thm:completed-product-fubini-12}
Niech $(X,\mathcal F,\mu)$ i $(Y,\mathcal G,\nu)$
będą przestrzeniami z miarami $\sigma$-skończonymi
i \emph{zupełnymi}. Niech $(\widehat{\mathcal H},
\widehat\pi)$ będzie uzupełnieniem przestrzeni
$(X\times Y,\mathcal F\otimes\mathcal G,\mu\otimes\nu)$.
Jeśli $f:X\times Y\to[0,\infty]$ jest
$\widehat{\mathcal H}$-mierzalna, to dla $\mu$-prawie
każdego $x$ przekrój $f_x$ jest $\mathcal G$-mierzalny,
a dla $\nu$-prawie każdego $y$ przekrój $f^y$ jest
$\mathcal F$-mierzalny. Po nadaniu wartości zero
na wyjątkowych mierzalnych zbiorach parametrów
funkcje powstałe przez całkowanie przekrojów są
mierzalne i spełniają wzór Tonellego z całką
względem $\widehat\pi$.

Jeśli $f$ jest rzeczywista lub zespolona,
$\widehat{\mathcal H}$-mierzalna i
$\int_{X\times Y}|f|\,\mathrm d\widehat\pi<\infty$,
to przekroje są całkowalne prawie wszędzie,
a wzór Fubiniego również zachodzi po pominięciu
zbiorów wyjątkowych.
\end{twierdzenie}
\begin{proof}
\emph{Krok 1: przedstawiciel mierzalny względem
$\mathcal F\otimes\mathcal G$.}
Z definicji uzupełnienia miary każda funkcja
$\widehat{\mathcal H}$-mierzalna o wartościach
w $\overline{\mathbb R}$ jest równa
$\widehat\pi$-prawie wszędzie pewnej funkcji
$g$ mierzalnej względem $\mathcal F\otimes\mathcal G$.
Aby to zobaczyć, dla $f\ge0$ wybieramy proste
mierzalne przybliżenia $s_n\uparrow f$, a każdy
zbiór poziomicowy każdej funkcji $s_n$ zastępujemy
zbiorem z $\mathcal F\otimes\mathcal G$ różniącym
się od niego tylko na podzbiorze mierzalnego
$\pi$-zbioru zerowego. Przeliczalna suma takich
zbiorów zerowych nadal ma $\pi$-miarę zero;
poza nią zastąpione funkcje proste mają tę samą
granicę co $f$. W przypadku funkcji ze znakiem
i zespolonych postępujemy osobno z częściami
rzeczywistymi i urojonymi. Daje to $g$ i zbiór
$N\in\mathcal F\otimes\mathcal G$ z $\pi(N)=0$
taki, że $f=g$ poza $N$.

\emph{Krok 2: przekroje zbioru zerowego.}
Tonelli dla $\chi_N$ daje
$0=\pi(N)=\int_X\nu(N_x)\,\mathrm d\mu(x)$.
Wobec tego $\nu(N_x)=0$ dla $\mu$-prawie każdego
$x$. Analogicznie $\mu(N^y)=0$ dla $\nu$-prawie
każdego $y$. Dla takich $x$ przekroje $f_x$ i $g_x$
różnią się jedynie na podzbiorze $N_x$.
Ponieważ $\nu$ jest zupełna, $f_x$ jest
$\mathcal G$-mierzalna i ma tę samą całkę co $g_x$;
podobnie z przekrojami poziomymi. Poza zbiorami
wyjątkowymi można zatem całkować $g$ zamiast $f$.

\emph{Krok 3: równość całek.}
Z równości $f=g$ $\widehat\pi$-prawie wszędzie
wynika równość ich całek względem $\widehat\pi$
w przypadku funkcji nieujemnej lub całkowalnej.
Na $g$ stosujemy odpowiednio twierdzenie
\ref{thm:tonelli-12} lub \ref{thm:fubini-12}.
Z kroku 2 całki przekrojowe $f$ i $g$ zgadzają
się prawie wszędzie, więc ich całki iterowane
są równe. To daje obie tezy.
\end{proof}

\section{Zbiory zerowe w produkcie i wyjątki przekrojowe}

\begin{wniosek}[Prawie wszystkie przekroje zbioru zerowego]
\label{cor:null-sections-12}
Niech $\mu,\nu$ będą miarami $\sigma$-skończonymi,
a $N\in\mathcal F\otimes\mathcal G$. Wówczas
$(\mu\otimes\nu)(N)=0$ wtedy i tylko wtedy, gdy
$\nu(N_x)=0$ dla $\mu$-prawie każdego $x$, a równoważnie
$\mu(N^y)=0$ dla $\nu$-prawie każdego $y$.
Nie wolno w tym stwierdzeniu zastępować
\emph{prawie każdego} przez \emph{każdego}.
\end{wniosek}
\begin{proof}
Na mocy twierdzenia~\ref{thm:section-identity-12}
obie funkcje przekrojowe są mierzalne i mają
nieujemne wartości. Z konstrukcji miary produktowej
\eqref{eq:measure-section-12} ich całki są równe
$(\mu\otimes\nu)(N)$. Całka funkcji nieujemnej
jest równa zeru dokładnie wtedy, gdy ta funkcja
znika prawie wszędzie. To daje obie równoważności.
Dla przykładu, gdy $X=Y=[0,1]$ z miarami Lebesgue'a,
zbiór $N=\{0\}\times[0,1]$ ma miarę produktową zero,
lecz $\lambda_1(N_0)=1$.
\end{proof}

\section{Więcej niż dwa czynniki i zastosowanie do prawdopodobieństwa}

\begin{wniosek}[Łączność skończonych produktów]
\label{cor:associative-product-12}
Niech $(X_j,\mathcal F_j,\mu_j)$, $j=1,2,3$,
będą przestrzeniami z miarami $\sigma$-skończonymi.
Przy naturalnej identyfikacji
$(X_1\times X_2)\times X_3=X_1\times(X_2\times X_3)$
zachodzi
$(\mu_1\otimes\mu_2)\otimes\mu_3
=\mu_1\otimes(\mu_2\otimes\mu_3)$.
W szczególności dla każdej nieujemnej funkcji
mierzalnej na iloczynie trzech przestrzeni
kolejność całkowania jest dowolna, a dla funkcji
bezwzględnie całkowalnej można stosować Fubiniego
w dowolnej kolejności.
\end{wniosek}
\begin{proof}
Oba konstrukty są miarami na tej samej
$\sigma$-algebrze produktowej, ponieważ
$\sigma$-algebra generowana przez trójwymiarowe
prostokąty nie zależy od umieszczenia nawiasów.
Dla każdego prostokąta $A_1\times A_2\times A_3$
obie miary przyjmują wartość
$\mu_1(A_1)\mu_2(A_2)\mu_3(A_3)$
(z konwencją dla zer i nieskończoności).
Miary są $\sigma$-skończone jako produkty miar
$\sigma$-skończonych. Jednoznaczność rozszerzenia
na generującym układzie prostokątów zapewnia ich
równość. Twierdzenia o całkach wynikają przez
kolejne zastosowanie Tonellego i Fubiniego.
\end{proof}

\begin{przyklad}[Niezależność dwóch zmiennych losowych]
Niech $P$ i $Q$ będą miarami probabilistycznymi
na $(X,\mathcal F)$ i $(Y,\mathcal G)$.
Na produkcie $(X\times Y,\mathcal F\otimes\mathcal G)$
rozważmy $P\otimes Q$ i odwzorowania współrzędne
$U:X\times Y\to X$, $U(x,y)=x$, oraz
$V:X\times Y\to Y$, $V(x,y)=y$.
Oba są mierzalne względem odpowiednich $\sigma$-algebr.
Dla $A\in\mathcal F$ i $B\in\mathcal G$
mamy
$(P\otimes Q)(\{U\in A,V\in B\})=P(A)Q(B)$.
To jest definicja niezależności zmiennych losowych
$U$ i $V$. Jeżeli $\varphi:X\to[0,\infty]$ oraz
$\psi:Y\to[0,\infty]$ są mierzalne, Tonelli daje
\[
\int_{X\times Y}\varphi(x)\psi(y)
\,\mathrm d(P\otimes Q)(x,y)
=\left(\int_X\varphi(x)\,\mathrm dP(x)\right)
\left(\int_Y\psi(y)\,\mathrm dQ(y)\right).
\]
Przy funkcjach zespolonych tożsamość obowiązuje,
jeżeli są całkowalne względem odpowiednich miar.
\end{przyklad}

\section{Podsumowanie i dalszy kierunek}
Najpierw udowodniliśmy mierzalność przekrojów,
a następnie \emph{bez użycia Fubiniego} wykazaliśmy
równość całek przekrojowych w prostokątach o
skończonej mierze. Twierdzenie o klasie monotonicznej
oraz wyczerpanie przez prostokąty $\sigma$-skończone
umożliwiły konstrukcję jednoznacznej miary
$\mu\otimes\nu$. Dopiero z tej konstrukcji
wyprowadziliśmy Tonellego i Fubiniego.

Warto zapamiętać rozdział ról: Tonelli pozwala
przenosić kolejność dla funkcji \emph{nieujemnej},
także o nieskończonej całce. Fubini zapewnia
pełną równość całek funkcji zmieniających znak,
lecz wymaga kontroli $\int|f|$. Przypadki
miar nie-$\sigma$-skończonych i niezupełność
$\sigma$-algebry produktowej ostrzegają przed
nadmiernym uogólnianiem tych wyników.

% END chapters/pl/measure_theory/12-rozdzial.tex


% BEGIN chapters/pl/measure_theory/13-rozdzial.tex
% ROZDZIAŁ 13 — samodzielny rozdział do projektu teoria_miary
\chapter{Zamiana zmiennych w całce Lebesgue'a}
\label{ch:zamiana-zmiennych}

\section{Miary obrazowe: co właściwie zmienia się w całce?}
Niech $(X,\mathcal F,\mu)$ będzie przestrzenią z miarą, $(Y,\mathcal G)$
przestrzenią mierzalną, a $T:X\to Y$ odwzorowaniem mierzalnym.
\emph{Miara obrazowa} $T_\#\mu$ jest określona na $\mathcal G$ wzorem
$(T_\#\mu)(B)=\mu(T^{-1}(B))$. Nazwa ta nie oznacza, że
$\mu$ jest funkcją i można ją w zwykły sposób złożyć z $T$.

\begin{twierdzenie}[Całkowanie względem miary obrazowej]
\label{thm:pushforward-integral-13}
Niech $T:(X,\mathcal F,\mu)\to(Y,\mathcal G)$ będzie mierzalne.
Jeśli $h:Y\to[0,\infty]$ jest $\mathcal G$-mierzalna, to
\begin{equation}\label{eq:pushforward-13}
\int_Y h(y)\,\mathrm d(T_\#\mu)(y)
=\int_X h(T(x))\,\mathrm d\mu(x).
\end{equation}
Ta sama równość zachodzi dla funkcji $h:Y\to\mathbb C$
całkowalnych względem $T_\#\mu$.
\end{twierdzenie}
\begin{proof}
Dla $h=\chi_B$ obie strony są równe $\mu(T^{-1}(B))$.
Z liniowości wynika teza dla nieujemnych funkcji prostych.
Dla dowolnej $h\ge0$ wybieramy funkcje proste $h_n\uparrow h$;
$h_n\circ T\uparrow h\circ T$, więc dwukrotne zastosowanie
twierdzenia Beppo Leviego daje \eqref{eq:pushforward-13}.
Jeśli $h$ jest zespolona i całkowalna, poprzedni przypadek dla
$|h|$ dowodzi całkowalności $h\circ T$. Następnie rozdzielamy
części rzeczywiste i urojone oraz ich części dodatnie i ujemne.
\end{proof}

\begin{uwaga}
Równość \eqref{eq:pushforward-13} nie wymaga, aby $T$ było
różnowartościowe, ciągłe, różniczkowalne, ani nawet aby przestrzenie
$X,Y$ były topologiczne. Natomiast wzór z wyznacznikiem jakobianu,
który udowodnimy poniżej, wykorzystuje dodatkowo lokalną strukturę
przestrzeni euklidesowej i różnowartościowość $T$.
\end{uwaga}

\begin{przyklad}[Miara obrazu a rozkład zmiennej losowej]
Na $X=(0,1)$ z miarą Lebesgue'a rozważmy $T(x)=x^2$.
Dla $0\le t\le1$ mamy
$(T_\#\lambda_1)((0,t])=\sqrt t$; miara $T_\#\lambda_1$
ma zatem gęstość $y\mapsto(2\sqrt y)^{-1}$ na $(0,1)$.
W istocie wzór
$\int_{(0,1)}h(x^2)\,\mathrm d\lambda_1(x)
=\int_{(0,1)}h(y)/(2\sqrt y)\,\mathrm d\lambda_1(y)$
poznamy jako szczególny przypadek zamiany zmiennych.
Na razie wyprowadziliśmy równość miar na przedziałach $(0,t]$;
jednoznaczność miar skończonych na generującym układzie przedziałów
uzasadnia stwierdzenie o gęstości.
\end{przyklad}

\section{Zamiana liniowa i geometryczne znaczenie wyznacznika}

\begin{lemat}[Translacje, permutacje i ścinanie]
\label{lem:shear-13}
Niech $d\ge1$, $b\in\mathbb R^d$ i $E\subset\mathbb R^d$
będzie zbiorem borelowskim. Translacja $E\mapsto E+b$ zachowuje
miarę $\lambda_d$, podobnie jak dowolna permutacja współrzędnych.
Dla $d\ge2$ miarę zachowuje także ścinanie postaci
\[
S(x_1,\ldots,x_d)=(x_1,\ldots,x_{j-1},
 x_j+a x_i,x_{j+1},\ldots,x_d),\qquad i\ne j.
\]
\end{lemat}
\begin{proof}
Niezmienniczość translacyjna została udowodniona w rozdziale 4.
Permutacje współrzędnych zachowują miarę prostokątów
półotwartych, zatem — przez jednoznaczność miary
$\sigma$-skończonej — wszystkich zbiorów borelowskich.
Dla ścinania ustalmy $d-1$ współrzędnych innych niż $x_j$.
W takim przekroju $S$ jest translacją zmiennej $x_j$
o stałą $a x_i$. Przekrój $S(E)$ jest więc translacją
odpowiedniego przekroju $E$. Ponieważ $S$ jest homeomorfizmem,
$S(E)$ jest borelowski. Twierdzenie Tonellego
\ref{thm:tonelli-12} zastosowane do funkcji charakterystycznych $\chi_E$ i $\chi_{S(E)}$
i niezmienniczość $\lambda_1$ dają równość ich miar.
\end{proof}

\begin{lemat}[Odwzorowania Lipschitzowskie zachowują zbiory zerowe]
\label{lem:lipschitz-null-13}
Niech $U\subset\mathbb R^d$ będzie otwarty, a
$F:U\to\mathbb R^d$ lokalnie Lipschitzowskie.
Jeżeli $N\subset U$ ma miarę Lebesgue'a zero, to
$\lambda_d(F(N))=0$ (w szczególności $F(N)$ jest mierzalny).
\end{lemat}
\begin{proof}
Zbiór $U$ można pokryć przeliczalnie domkniętymi sześcianami
$Q_j$ o dodatniej odległości od $\partial U$, takimi że
na otoczeniu każdego $Q_j$ funkcja $F$ jest Lipschitzowska
ze stałą $L_j<\infty$. Wystarczy rozważyć $N\cap Q_j$.
Dla ustalonego $j$ i $\varepsilon>0$ pokrywamy ten zbiór
przeliczalnie sześcianami $I_k$ o tak małych średnicach,
że leżą w odpowiednim otoczeniu $Q_j$ oraz
$\sum_k\lambda_d(I_k)<\varepsilon$.
Obraz $F(I_k)$ jest zawarty w sześcianie o boku
co najwyżej $2L_j\sqrt d$ razy bok $I_k$ (wokół
obrazu wybranego punktu $I_k$). Stąd
$\lambda_d^*(F(N\cap Q_j))\le C_d L_j^d\varepsilon$.
Przechodząc do $\varepsilon\downarrow0$, otrzymujemy
miarę zewnętrzną zero; zupełność miary Lebesgue'a
zapewnia mierzalność. Przeliczalna suma obrazów
$F(N\cap Q_j)$ jest więc zerowa.
\end{proof}

\begin{twierdzenie}[Liniowy wzór na objętość]
\label{thm:linear-change-13}
Jeżeli $A\in\mathrm{GL}(d,\mathbb R)$, $b\in\mathbb R^d$
i $E\subset\mathbb R^d$ jest borelowski, to
\begin{equation}\label{eq:linear-volume-13}
\lambda_d(AE+b)=|\det A|\lambda_d(E).
\end{equation}
Równość obowiązuje również dla dowolnego zbioru
$\lambda_d$-mierzalnego $E$, z wartościami w $[0,\infty]$.
\end{twierdzenie}
\begin{proof}
Mnożenie jednej współrzędnej przez $c\ne0$ mnoży miarę
przez $|c|$: w jednej zmiennej jest to skalowanie miary
Lebesgue'a, a w pozostałych stosujemy Tonellego.
Lemat~\ref{lem:shear-13} pokazuje, że operacje elementarne
polegające na dodaniu wielokrotności jednego wiersza do drugiego,
permutacji współrzędnych i skalowaniu pojedynczej współrzędnej
zmieniają miarę dokładnie przez moduł swego wyznacznika.
Eliminacja Gaussa przedstawia $A$ jako iloczyn skończenie
wielu takich macierzy. Kolejne zastosowanie ustalonych wzorów
i multiplikatywność wyznacznika dają \eqref{eq:linear-volume-13}
dla zbiorów borelowskich.

Macierz $A$ i jej odwrotność są globalnie Lipschitzowskie.
Z lematu~\ref{lem:lipschitz-null-13} wynika, że odwzorowanie lokalnie
Lipschitzowskie przenosi zbiory zerowe na zbiory miary zero.
Każdy zbiór Lebesgue'a-mierzalny różni się od zbioru borelowskiego
co najwyżej na podzbiorze borelowskiego zbioru zerowego;
zatem $AE+b$ jest mierzalny, a jego miara jest wyznaczona
przez część borelowską. To daje rozszerzenie na miarę zupełną.
\end{proof}

\begin{uwaga}[Wymiar ma znaczenie]
Twierdzenie~\ref{thm:linear-change-13} wymaga
macierzy kwadratowej nieosobliwej. W przypadku odwzorowań
$\mathbb R^m\to\mathbb R^n$, $m\ne n$, pojawiają się
inne wielkości (jakobiany wymiarowe, miary Hausdorffa),
które wprowadzimy w geometrycznej części skryptu.
\end{uwaga}

\section{Lokalna deformacja sześcianu}

\begin{lemat}[Oszacowanie obrazu małego sześcianu]
\label{lem:small-cube-13}
Niech $U\subset\mathbb R^d$ będzie otwarty, a
$\Phi:U\to\mathbb R^d$ odwzorowaniem klasy $C^1$.
Załóżmy, że $A=D\Phi(x_0)$ jest odwracalna.
Dla $r>0$ oznaczmy $Q_r=[-r,r]^d$ i przyjmijmy
$x_0+Q_r\subset U$. Istnieje funkcja $\delta(r)\to0$
przy $r\downarrow0$ taka, że dla dostatecznie małych $r$
\begin{equation}\label{eq:cube-est-13}
(1-\delta(r))^d|\det A|\lambda_d(Q_r)
\le \lambda_d(\Phi(x_0+Q_r))
\le(1+\delta(r))^d|\det A|\lambda_d(Q_r).
\end{equation}
W szczególności iloraz miary obrazu i objętości sześcianu
dąży do $|\det D\Phi(x_0)|$.
\end{lemat}
\begin{proof}
Wprowadzamy normę maksimum $\|z\|_\infty$ w $\mathbb R^d$
i funkcję $H(z)=A^{-1}(\Phi(x_0+z)-\Phi(x_0))$.
Zachodzi $H(0)=0$ i $DH(0)=I$.
Niech $R(z)=H(z)-z$. Z ciągłości $DH$ wynika, że
\[
\delta(r):=\sup_{z\in Q_r}\|DR(z)\|_{\infty\to\infty}
\longrightarrow0.
\]
Dla $z,w\in Q_r$ twierdzenie o przyrostach skończonych
wzdłuż odcinka $[z,w]$ daje
$\|R(z)-R(w)\|_\infty\le\delta(r)\|z-w\|_\infty$,
a po podstawieniu $w=0$ dostajemy
$\|R(z)\|_\infty\le\delta(r)r$.
Zatem $H(Q_r)\subseteq Q_{(1+\delta(r))r}$.

Dla $y\in Q_{(1-\delta(r))r}$ rozważmy
$T_y(z)=y-R(z)$ na pełnej, domkniętej przestrzeni
metrycznej $Q_r$. Funkcja ta odwzorowuje $Q_r$ w siebie,
gdyż $\|T_y(z)\|_\infty\le(1-\delta)r+\delta r=r$,
i jest zwężeniem o stałej $\delta(r)<1$.
Twierdzenie o punkcie stałym Banacha (iteracja
$z_{n+1}=T_y(z_n)$ i oszacowanie geometryczne różnic)
daje $z\in Q_r$ z $T_y(z)=z$, czyli $H(z)=y$.
Wynika stąd drugie zawieranie
$Q_{(1-\delta(r))r}\subseteq H(Q_r)$.
Po zastosowaniu $A$ i translacji otrzymujemy
\[
\Phi(x_0)+A Q_{(1-\delta)r}
\subseteq\Phi(x_0+Q_r)
\subseteq\Phi(x_0)+A Q_{(1+\delta)r}.
\]
Miary skrajnych zbiorów liczymy za pomocą
\eqref{eq:linear-volume-13}; monotoniczność miary
daje \eqref{eq:cube-est-13}.
\end{proof}

\begin{uwaga}
W dowodzie nie użyto nieudowodnionego wzoru o zamianie zmiennych,
lecz tylko przypadku liniowego i twierdzenia o punkcie stałym.
Co więcej, na zbiorze zwartym, na którym $D\Phi$ jest ciągła
i odwracalna, parametry $\delta(r)$ można dobrać
\emph{jednostajnie} względem środka sześcianu, o ile sześciany
leżą w ustalonym otwartym otoczeniu tego zbioru.
Istotne są tu jednostajna ciągłość $D\Phi$ i ograniczoność
normy $(D\Phi)^{-1}$ na rozważanym zbiorze zwartym.
\end{uwaga}

\section{Główne twierdzenie o zamianie zmiennych}

\begin{twierdzenie}[Zamiana zmiennych dla dyfeomorfizmu $C^1$]
\label{thm:change-variables-13}
Niech $U,V\subset\mathbb R^d$ będą otwarte i niech
$\Phi:U\to V$ będzie dyfeomorfizmem klasy $C^1$
(tzn. bijekcją, dla której $\Phi$ i $\Phi^{-1}$
są klasy $C^1$). Wtedy:
\begin{enumerate}[label=\textup{(\roman*)}]
\item dla każdego zbioru Lebesgue'a-mierzalnego $E\subset U$
zbiór $\Phi(E)$ jest Lebesgue'a-mierzalny i
\begin{equation}\label{eq:change-sets-13}
\lambda_d(\Phi(E))=\int_E |\det D\Phi(x)|\,
\mathrm d\lambda_d(x);
\end{equation}
\item dla każdej Lebesgue'a-mierzalnej
$f:V\to[0,\infty]$ zachodzi
\begin{equation}\label{eq:change-main-13}
\int_V f(y)\,\mathrm d\lambda_d(y)
=\int_U f(\Phi(x))|\det D\Phi(x)|\,
\mathrm d\lambda_d(x);
\end{equation}
\item ta sama równość zachodzi dla funkcji
$f:V\to\mathbb C$ całkowalnych względem $\lambda_d$.
\end{enumerate}
\end{twierdzenie}
\begin{proof}
\emph{Krok 1. Miara obrazu dla zbiorów borelowskich.}
Dla borelowskiego $E\subset U$ określamy
$\tau(E)=\lambda_d(\Phi(E))$. Jest to miara na
$\mathcal B(U)$, ponieważ $\Phi$ jest homeomorfizmem:
obrazy zbiorów borelowskich są borelowskie w $V$, a obrazy
rozłącznych zbiorów są rozłączne. Określamy drugą miarę
$\rho(E)=\int_E J(x)\,\mathrm d\lambda_d(x)$,
gdzie $J(x)=|\det D\Phi(x)|$. Funkcja $J$ jest dodatnia
i ciągła; obie miary są $\sigma$-skończone, ponieważ
$U$ jest przeliczalną sumą zwartych podzbiorów,
a obrazy tych zwartych zbiorów są zwarte.

\emph{Krok 2. Równość na sześcianach.}
Niech $Q\subset U$ będzie domkniętym sześcianem.
Dzielimy go na $n^d$ przystających, domkniętych
sześcianów $Q_{n,j}$, $1\le j\le n^d$, o rozłącznych wnętrzach.
Na mocy lematu~\ref{lem:lipschitz-null-13} obraz brzegu
każdego $Q_{n,j}$ ma miarę zero. Zatem
$\tau(Q)=\sum_j\tau(Q_{n,j})$.
Wybieramy środek $x_{n,j}$ sześcianu $Q_{n,j}$.
Ponieważ $Q\Subset U$, funkcje $D\Phi$ i
$(D\Phi)^{-1}$ są na zwartym zbiorze $Q$ odpowiednio
jednostajnie ciągłe i ograniczone; stąd w oszacowaniach
wszystkich $Q_{n,j}$ można przyjąć jeden parametr
$\delta_n\to0$, niezależny od $j$.
Oszacowanie z lematu~\ref{lem:small-cube-13},
jednostajne względem tych środków, prowadzi do
\[
(1-\delta_n)^d\sum_jJ(x_{n,j})\lambda_d(Q_{n,j})
\le\tau(Q)\le
(1+\delta_n)^d\sum_jJ(x_{n,j})\lambda_d(Q_{n,j}),
\qquad\delta_n\to0.
\]
Sumy $\sum_jJ(x_{n,j})\lambda_d(Q_{n,j})$ zbiegają
do $\int_QJ(x)\,\mathrm d\lambda_d(x)$ jako sumy Riemanna
funkcji ciągłej na zwartym sześcianie. Otrzymujemy
$\tau(Q)=\rho(Q)$. Ta sama równość dotyczy wnętrza
sześcianu i półotwartego sześcianu, bo ich brzegi
oraz obrazy tych brzegów mają miarę zero.

\emph{Krok 3. Równość miar na $\mathcal B(U)$.}
Niech $\mathcal P$ będzie rodziną wszystkich półotwartych
sześcianów diadycznych, których domknięcia zawierają się
w $U$, wraz ze zbiorem pustym. Dwa sześciany diadyczne
albo są rozłączne, albo jeden zawiera się w drugim,
Zatem $\mathcal P$ jest układem typu $\pi$:
dwa sześciany w tej samej siatce diadycznej albo są
rozłączne, albo identyczne, a sześciany w siatkach
o różnych skalach są rozłączne lub jeden zawiera się
w drugim. Dokładniej używamy półotwartych sześcianów
$\prod_{j=1}^d[k_j2^{-n},(k_j+1)2^{-n})$,
$n\in\mathbb N_0$, $k_j\in\mathbb Z$, których domknięcia
są zawarte w $U$. Rodzina ta jest przeliczalna,
pokrywa $U$ i generuje
$\mathcal B(U)$: dla każdego $x\in U$ oraz otwartego
otoczenia $W\subset U$ punktu $x$ można wybrać
sześcian diadyczny $Q$ z $x\in Q$ i $\overline Q\subset W$.
Na każdym $C\in\mathcal P$ mamy
$\tau(C)=\rho(C)<\infty$. Aby poprawnie zastosować
argument $\pi$--$\lambda$, dla ustalonego takiego $C$
rozważamy dwie \emph{skończone} miary
$\tau_C(E)=\tau(E\cap C)$ oraz
$\rho_C(E)=\rho(E\cap C)$ na $\mathcal B(U)$.
Na $\mathcal P$ są równe, ponieważ $E\cap C\in\mathcal P$,
a z twierdzenia $\pi$--$\lambda$ wynika
$\tau(E\cap C)=\rho(E\cap C)$ dla dowolnego
$E\in\mathcal B(U)$. Wreszcie numerujemy sześciany
$C_1,C_2,\dots$ pokrywające $U$ i rozłączamy je przez
$D_1=C_1$, $D_j=C_j\setminus\bigcup_{i<j}C_i$.
Dla każdego borelowskiego $E\subset U$ równość miar
na $E\cap D_j$ wynika ze skończonej addytywności
równości na $E\cap C_j$ i $E\cap C_j\cap C_i$
(lub z rozdzielenia przez wcześniejsze sześciany).
Po zsumowaniu po $j$ dostajemy
$\tau(E)=\rho(E)$. To rozpisanie jest potrzebne:
rodzina sześcianów $\mathcal P$ sama nie zawiera
na ogół rosnącego wyczerpania $U$.

\emph{Krok 4. Uzupełnienie i przejście do całek.}
Ponieważ $\Phi$ oraz $\Phi^{-1}$ są lokalnie
Lipschitzowskie, lemat~\ref{lem:lipschitz-null-13}
pokazuje, że obrazy i przeciwobrazy zbiorów
$\lambda_d$-zerowych są zerowe. Każdy zbiór Lebesgue'a-
mierzalny różni się od borelowskiego na zbiorze zerowym,
a więc tożsamość $\tau(E)=\rho(E)$ rozciąga się na
wszystkie Lebesgue'a-mierzalne $E\subset U$; daje to (i).

Wzór (ii) sprawdzamy najpierw dla funkcji $f=\chi_B$,
$B\subset V$ mierzalnego: jego lewa strona to
$\lambda_d(B)$, a prawa, na mocy (i), to
$\int_{\Phi^{-1}(B)}J(x)\,\mathrm d\lambda_d(x)$,
co jest równe $\lambda_d(B)$.
Następnie korzystamy z liniowości dla funkcji prostych
i z twierdzenia Beppo Leviego dla funkcji nieujemnych.
W przypadku funkcji zespolonej całkowalnej stosujemy
(ii) do $|f|$, aby zapewnić całkowalność prawej strony,
a następnie rozkładamy $f$ na cztery części nieujemne.
\end{proof}

\begin{wniosek}[Zamiana na mierzalnym podzbiorze]
\label{cor:change-subdomain-13}
Przy założeniach twierdzenia~\ref{thm:change-variables-13}
dla mierzalnego $E\subset U$ i nieujemnej
mierzalnej $g:\Phi(E)\to[0,\infty]$ zachodzi
$\int_{\Phi(E)}g(y)\,\mathrm d\lambda_d(y)
=\int_E g(\Phi(x))J(x)\,\mathrm d\lambda_d(x)$.
Analogicznie dla funkcji całkowalnych.
\end{wniosek}
\begin{proof}
Przedłużamy $g$ przez zero na $V\setminus\Phi(E)$
i stosujemy poprzednie twierdzenie.
\end{proof}

\begin{uwaga}[Skąd bierze się moduł?]
Gdy $\det D\Phi<0$, odwzorowanie zmienia orientację,
lecz miara Lebesgue'a jest nieujemna. Z tego powodu
w formule występuje $|\det D\Phi|$, nie sam wyznacznik.
W całkowaniu form różniczkowych bez modułu uwzględnia się
orientację i otrzymuje inną regułę.
\end{uwaga}

\section{Przykłady: współrzędne biegunowe, walcowe i sferyczne}

\begin{przyklad}[Współrzędne biegunowe]
\label{ex:polar-13}
Rozważmy $\Phi:(0,\infty)\times(0,2\pi)\to\mathbb R^2$,
$\Phi(r,\theta)=(r\cos\theta,r\sin\theta)$.
Jest to dyfeomorfizm na płaszczyznę bez dodatniej
półosi poziomej i punktu początkowego, czyli na
$\mathbb R^2\setminus\{(x,0):x\ge0\}$.
Jego macierz pochodnej ma kolumny
$(\cos\theta,\sin\theta)^T$ i
$(-r\sin\theta,r\cos\theta)^T$, stąd
$\det D\Phi(r,\theta)=r$. Usunięta półprosta ma
$\lambda_2$-miarę zero. Dlatego dla każdej
mierzalnej funkcji $f:\mathbb R^2\to[0,\infty]$
\begin{equation}\label{eq:polar-13}
\int_{\mathbb R^2}f(x,y)\,\mathrm d\lambda_2(x,y)
=\int_{(0,2\pi)}\left(\int_{(0,\infty)}
 f(r\cos\theta,r\sin\theta)r\,\mathrm d\lambda_1(r)\right)
\mathrm d\lambda_1(\theta).
\end{equation}
Kolejność całkowania uzasadnia twierdzenie Tonellego.
Dla $f=\chi_{\{x^2+y^2<R^2\}}$ otrzymujemy
$\lambda_2(B(0,R))=\int_{(0,2\pi)}\int_{(0,R)}r\,
\mathrm d\lambda_1(r)\,\mathrm d\lambda_1(\theta)=\pi R^2$.
\end{przyklad}

\begin{center}
\begin{tikzpicture}[scale=1.0,>=stealth]
\draw[->] (-0.25,0)--(3.45,0) node[right] {$x$};
\draw[->] (0,-0.15)--(0,2.55) node[above] {$y$};
\filldraw[fill=blue!13,draw=blue!55!black,thick]
(25:1.55) -- (25:2.15) arc[start angle=25,end angle=42,radius=2.15]
-- (42:1.55) arc[start angle=42,end angle=25,radius=1.55] -- cycle;
\draw[dashed] (0,0)--(25:2.6);
\draw[dashed] (0,0)--(42:2.6);
\draw[->] (0.7,0) arc[start angle=0,end angle=25,radius=0.7];
\node at (12:1.1) {$\theta$};
\node[anchor=west] at (2.2,1.85) {$\Delta r$};
\draw[->] (2.18,1.75)--(42:2.07);
\node at (60:2.6) {$\Delta\theta$};
\node at (15:2.7) {$r$};
\node[anchor=west] at (0.1,-0.35) {Mały wycinek pierścienia};
\end{tikzpicture}
\end{center}
\noindent
Dla małych przyrostów $\Delta r$ i $\Delta\theta$ pole
zaznaczonego wycinka jest w pierwszym przybliżeniu
$r\,\Delta r\,\Delta\theta$. Jest to geometryczna intuicja
za jakobianem $r$; ścisłym uzasadnieniem jest jednak
twierdzenie~\ref{thm:change-variables-13}.

\begin{przyklad}[Współrzędne walcowe]
Dla $r>0$, $\theta\in(0,2\pi)$, $z\in\mathbb R$
połóżmy $\Psi(r,\theta,z)=(r\cos\theta,r\sin\theta,z)$.
Wyznacznik jej macierzy pochodnej wynosi $r$.
Po odrzuceniu zbioru zerowej miary zawierającego
odpowiednią półpłaszczyznę otrzymujemy
$\int_{\mathbb R^3}f(x,y,z)\,\mathrm d\lambda_3(x,y,z)
=\int_{\mathbb R}\int_{(0,2\pi)}\int_{(0,\infty)}
 f(r\cos\theta,r\sin\theta,z)r\,
\mathrm d\lambda_1(r)\,\mathrm d\lambda_1(\theta)
\,\mathrm d\lambda_1(z)$
dla $f\ge0$, a także dla $f\in L^1(\mathbb R^3)$.
\end{przyklad}

\begin{przyklad}[Współrzędne sferyczne]
Dla $r>0$, $\vartheta\in(0,\pi)$, $\varphi\in(0,2\pi)$
rozważmy
$\Psi(r,\vartheta,\varphi)=
(r\sin\vartheta\cos\varphi,
 r\sin\vartheta\sin\varphi,r\cos\vartheta)$.
Bezpośrednie obliczenie wyznacznika daje
$|\det D\Psi|=r^2\sin\vartheta$.
Po odrzuceniu odpowiedniej półpłaszczyzny i osi,
które mają miarę zero, otrzymujemy
\begin{align*}
\int_{\mathbb R^3}f(x)\,\mathrm d\lambda_3(x)
&=\int_{(0,\infty)}\int_{(0,\pi)}\int_{(0,2\pi)}
 f(\Psi(r,\vartheta,\varphi))r^2\sin\vartheta\,
 \mathrm d\lambda_1(\varphi)\,\mathrm d\lambda_1(\vartheta)
 \,\mathrm d\lambda_1(r).
\end{align*}
Dla kuli $B(0,R)$ daje to $\lambda_3(B(0,R))=4\pi R^3/3$.
\end{przyklad}

\section{Obrazy miar posiadających gęstości}
\begin{twierdzenie}[Transformacja gęstości]
\label{thm:density-transform-13}
Niech $U,V,\Phi$ spełniają założenia
twierdzenia~\ref{thm:change-variables-13}.
Niech $w:U\to[0,\infty]$ będzie mierzalna i
$\mu(E)=\int_Ew(x)\,\mathrm d\lambda_d(x)$
dla mierzalnych $E\subset U$.
Wtedy dla każdego mierzalnego $B\subset V$
\[
(\Phi_\#\mu)(B)=\int_B
w(\Phi^{-1}(y))|\det D\Phi^{-1}(y)|\,
\mathrm d\lambda_d(y).
\]
\end{twierdzenie}
\begin{proof}
Z definicji miary obrazowej mamy
$(\Phi_\#\mu)(B)=\int_{\Phi^{-1}(B)}w(x)\,
\mathrm d\lambda_d(x)$. Stosujemy twierdzenie
o zamianie zmiennych do dyfeomorfizmu $\Phi^{-1}$
i funkcji $y\mapsto\chi_B(y)w(\Phi^{-1}(y))$.
Wynika stąd podana formuła. Jej prawą stronę
interpretujemy również dla nieskończonych gęstości
zgodnie z konwencją $0\cdot\infty=0$.
\end{proof}

\begin{przyklad}[Rozkład jednostajny i kwadratowanie]
Jeśli $X$ ma rozkład jednostajny na $(0,1)$ i
$Y=X^2$, to $Y$ ma gęstość
$y\mapsto\chi_{(0,1)}(y)/(2\sqrt y)$.
To dokładnie miara obrazowa rozważana na początku rozdziału.
Nie ma potrzeby zakładać, że sam $X$ jest określony
na $(0,1)$: wystarczy, aby jego rozkład był jednostajny.
\end{przyklad}

\section{Założenia, które naprawdę są potrzebne}
\begin{przyklad}[Dlaczego nie wolno pomijać różnowartościowości]
Dla $\Phi:\mathbb R\to\mathbb R$, $\Phi(x)=x^2$,
nie istnieje globalna odwrotność. Gdyby bez dodatkowych
poprawek użyć wzoru z jakobianem na całej prostej, dla
$f=\chi_{(0,1)}$ otrzymalibyśmy po prawej stronie
$\int_{(-1,1)}2|x|\,\mathrm d\lambda_1(x)=2$,
zaś po lewej $\int_{(0,1)}1\,\mathrm d\lambda_1(y)=1$.
Rozbieżność wynika z dwukrotnego liczenia punktów obrazu.
Dla nieiniektywnych odwzorowań właściwym narzędziem
jest wzór z krotnością włókien, a w większej ogólności
wzór na pole powierzchni (późniejszy rozdział).
\end{przyklad}

\begin{przyklad}[Wymagana regularność]
Z samej różnowartościowości i ciągłości $\Phi$ nie wynika
wzór z pochodną. Na przykład istnieją ściśle rosnące
homeomorfizmy przedziału, których pochodna znika
prawie wszędzie (homeomorfizmy osobliwe).
Przy takiej funkcji całka pochodnej nie daje długości obrazu.
Nie wystarczy więc samo istnienie pochodnej prawie wszędzie:
w odpowiedniej wersji uogólnionej potrzebne są dodatkowe
własności, np. absolutna ciągłość w wymiarze jeden.
\end{przyklad}

\section{Podsumowanie}
Wzór na całkę względem miary obrazowej jest czysto
mierzalny. Natomiast wyznacznik jakobianu opisuje
lokalną zmianę objętości pod dyfeomorfizmem $C^1$.
Dowód przeszedł od macierzy elementarnych do lokalnych
sześcianów, a następnie do równości miar i funkcji
mierzalnych. Formuły biegunowe i sferyczne wynikają
z jednego twierdzenia, a nie z odrębnych reguł rachunkowych.

% END chapters/pl/measure_theory/13-rozdzial.tex


% BEGIN chapters/pl/measure_theory/14-rozdzial.tex
% ROZDZIAŁ 14 — samodzielny rozdział do projektu teoria_miary
\chapter{Przestrzenie $L^p$ i ich podstawowe własności}
\label{ch:przestrzenie-lp}

\section{Dlaczego utożsamiamy funkcje równe prawie wszędzie?}
Niech $(X,\mathcal F,\mu)$ będzie dowolną przestrzenią z miarą.
Nie zakładamy, że $\mu(X)<\infty$, ani że $\mu$ jest $\sigma$-skończona.
Dla $1\le p<\infty$ rozważmy mierzalne funkcje
$f:X\to\mathbb C$ spełniające
$\int_X|f(x)|^p\,\mathrm d\mu(x)<\infty$.
Naturalna wielkość
$(\int_X|f|^p\,\mathrm d\mu(x))^{1/p}$
może znikać dla funkcji nieidentycznie zerowej:
wystarczy, że $f$ różni się od zera jedynie na zbiorze
miary zero. Aby otrzymać normę, trzeba przejść
od funkcji do ich klas równoważności.

\begin{definicja}[Równość prawie wszędzie i przestrzeń $L^p$]
\label{def:lp-14}
Dwie mierzalne funkcje zespolone $f,g$ na $X$
uznajemy za równoważne, jeśli $\mu(\{x:f(x)\ne g(x)\})=0$.
Dla $1\le p<\infty$ przez $L^p(X,\mathcal F,\mu)$, lub
krótko $L^p(\mu)$, rozumiemy przestrzeń klas
równoważności mierzalnych funkcji $f:X\to\mathbb C$,
dla których $\int_X|f(x)|^p\,\mathrm d\mu(x)<\infty$.
Wartość wyrażenia $\int_X|f|^p\,\mathrm d\mu(x)$
nie zależy od przedstawiciela klasy. Kładziemy
\begin{equation}\label{eq:p-norm-14}
\|f\|_p:=\left(\int_X|f(x)|^p\,\mathrm d\mu(x)\right)^{1/p}.
\end{equation}
Dla funkcji rzeczywistych używamy oznaczenia $L^p(\mu;\mathbb R)$.
\end{definicja}

\begin{definicja}[Kres górny istotny i przestrzeń $L^\infty$]
\label{def:linf-14}
Dla mierzalnej funkcji $f:X\to\mathbb C$ określamy
\[
\|f\|_\infty:=\inf\{M\ge0:\mu(\{|f|>M\})=0\},
\]
z umową $\inf\varnothing=\infty$.
Przestrzeń $L^\infty(\mu)$ jest przestrzenią klas
funkcji o skończonej wartości $\|f\|_\infty$.
Liczbę tę nazywamy \emph{kresem górnym istotnym} funkcji $|f|$.
\end{definicja}

\begin{lemat}[Znaczenie kresu istotnego]
\label{lem:esssup-14}
Jeśli $f$ jest mierzalna i $\|f\|_\infty<\infty$,
to $|f(x)|\le\|f\|_\infty$ dla prawie wszystkich $x$.
Ponadto $\|f\|_\infty=0$ wtedy i tylko wtedy, gdy
$f=0$ prawie wszędzie.
\end{lemat}
\begin{proof}
Niech $M=\|f\|_\infty$. Dla każdego $n$ istnieje
$M_n<M+1/n$ spełniające $\mu(\{|f|>M_n\})=0$.
Poza przeliczalną sumą tych zbiorów zachodzi
$|f|\le M_n<M+1/n$ dla wszystkich $n$, więc $|f|\le M$.
W przypadku $M=0$ otrzymujemy $f=0$ poza zbiorem zerowym.
Implikacja odwrotna wynika wprost z definicji.
\end{proof}

\begin{uwaga}[Przedstawiciele a klasy]
Zapis $f(x)$ dla $f\in L^p(\mu)$ jest skrótem:
ściśle rzecz biorąc $f$ jest klasą równoważności, więc
wartość w pojedynczym punkcie nie musi być określona.
Ilekroć piszemy $f(x)$, wybieramy przedstawiciela;
tezy o wartościach punktowych formułujemy prawie wszędzie.
Operacje algebraiczne na klasach są poprawnie określone,
ponieważ suma skończonej liczby zbiorów zerowych jest zerowa.
\end{uwaga}

\section{Elementarna nierówność Younga}

\begin{lemat}[Young dla iloczynu dwóch liczb]
\label{lem:young-numbers-14}
Niech $1<p,q<\infty$ oraz $1/p+1/q=1$.
Dla $a,b\ge0$ zachodzi
\begin{equation}\label{eq:young-14}
ab\le\frac{a^p}{p}+\frac{b^q}{q}.
\end{equation}
Równość zachodzi dokładnie wtedy, gdy $a^p=b^q$.
\end{lemat}
\begin{proof}
Dla $b>0$ funkcja $\varphi(a)=a^p/p-ab+b^q/q$
ma pochodną $\varphi'(a)=a^{p-1}-b$.
Jej minimum na $[0,\infty)$ wypada w punkcie
$a=b^{1/(p-1)}=b^{q-1}$, gdzie
$\varphi(a)=0$, ponieważ $(q-1)p=q$.
Dla $b=0$ nierówność jest natychmiastowa;
równość zachodzi wtedy jedynie przy $a=0$.
\end{proof}

\begin{uwaga}
Wykładniki $p$ i $q$ w lemacie nazywamy
\emph{sprzężonymi wykładnikami Höldera}.
Dla $p=1$ przyjmujemy formalnie $q=\infty$ i odwrotnie.
Przypadki brzegowe rozpatrujemy osobno.
\end{uwaga}

\section{Nierówność Höldera}
\begin{twierdzenie}[Hölder]
\label{thm:holder-14}
Niech $(X,\mathcal F,\mu)$ będzie przestrzenią z miarą,
$1\le p,q\le\infty$ i $1/p+1/q=1$, przy czym $1/\infty=0$.
Jeśli $f\in L^p(\mu)$ i $g\in L^q(\mu)$,
to $fg\in L^1(\mu)$ oraz
\begin{equation}\label{eq:holder-14}
\int_X|f(x)g(x)|\,\mathrm d\mu(x)
\le\|f\|_p\|g\|_q.
\end{equation}
\end{twierdzenie}
\begin{proof}
Załóżmy najpierw $1<p,q<\infty$.
Jeśli $\|f\|_p=0$ lub $\|g\|_q=0$,
to odpowiedni czynnik znika prawie wszędzie,
a więc teza jest oczywista. W przeciwnym przypadku
kładziemy $F=|f|/\|f\|_p$ oraz $G=|g|/\|g\|_q$.
Zachodzi $\int_XF^p\,\mathrm d\mu(x)=
\int_XG^q\,\mathrm d\mu(x)=1$.
Z lematu~\ref{lem:young-numbers-14} mamy punktowo
$FG\le F^p/p+G^q/q$; całkowanie daje
$\int_XFG\,\mathrm d\mu(x)\le1/p+1/q=1$.
Po pomnożeniu przez $\|f\|_p\|g\|_q$ otrzymujemy
\eqref{eq:holder-14} i zarazem całkowalność iloczynu.

Gdy $p=1$ i $q=\infty$, poza zbiorem zerowym zachodzi
$|fg|\le\|g\|_\infty|f|$ na mocy
lematu~\ref{lem:esssup-14}. Całkowanie daje nierówność;
przypadek $p=\infty$, $q=1$ jest symetryczny.
\end{proof}

\begin{wniosek}[Uogólniony Hölder]
\label{cor:generalized-holder-14}
Niech $m\ge2$, $1\le p_j\le\infty$ oraz
$\sum_{j=1}^m1/p_j=1$. Jeżeli $f_j\in L^{p_j}(\mu)$,
to $\prod_{j=1}^m f_j\in L^1(\mu)$ i
$\int_X\prod_{j=1}^m|f_j(x)|\,\mathrm d\mu(x)
\le\prod_{j=1}^m\|f_j\|_{p_j}$.
\end{wniosek}
\begin{proof}
Czynniki z $p_j=\infty$ są ograniczone prawie wszędzie;
można je oszacować przez $\|f_j\|_\infty$ i usunąć
z rozważań. Jeśli pozostał tylko jeden czynnik,
jego wykładnik musi być równy $1$ i wynik jest oczywisty.
Dla większej liczby czynników dowodzimy twierdzenia indukcyjnie.
Niech $p_m<\infty$ i $1/r=\sum_{j<m}1/p_j=1-1/p_m$.
Jeżeli $r=\infty$, wszystkie pierwsze czynniki
są ograniczone i dostajemy tezę bezpośrednio.
Jeśli $1<r<\infty$, to $p_j\ge r$ dla $j<m$,
a liczby $s_j=p_j/r\ge1$ spełniają
$\sum_{j<m}1/s_j=1$. Stosujemy krok indukcyjny
do funkcji $|f_j|^r\in L^{s_j}$:
\[
\int_X\prod_{j<m}|f_j(x)|^r\,\mathrm d\mu(x)
\le\prod_{j<m}\left(\int_X|f_j(x)|^{p_j}\,
\mathrm d\mu(x)\right)^{r/p_j}.
\]
Po wyciągnięciu pierwiastka stopnia $r$ uzyskujemy
$\|\prod_{j<m}f_j\|_r\le\prod_{j<m}\|f_j\|_{p_j}$.
Ostatnie zastosowanie Höldera do tego iloczynu i $f_m$
ze sprzężonymi wykładnikami $r,p_m$ kończy dowód.
\end{proof}

\begin{przyklad}[Oszacowanie iloczynu]
Na $X=(0,1)$ z miarą Lebesgue'a rozważmy
$f(x)=x^{-1/4}$ i $g(x)=x^{-1/3}$.
Wtedy $f\in L^2(0,1)$, $g\in L^2(0,1)$,
ponieważ $\int_{(0,1)}x^{-1/2}\,\mathrm d\lambda_1(x)<\infty$
i $\int_{(0,1)}x^{-2/3}\,\mathrm d\lambda_1(x)<\infty$.
Z Höldera wynika, że $fg=x^{-7/12}\in L^1(0,1)$.
Można to również sprawdzić bezpośrednio.
\end{przyklad}

\section{Nierówność Minkowskiego i struktura normowana}
\begin{twierdzenie}[Minkowski]
\label{thm:minkowski-14}
Dla każdej przestrzeni z miarą $(X,\mathcal F,\mu)$,
każdego $1\le p\le\infty$ oraz $f,g\in L^p(\mu)$
zachodzi $f+g\in L^p(\mu)$ oraz
\begin{equation}\label{eq:minkowski-14}
\|f+g\|_p\le\|f\|_p+\|g\|_p.
\end{equation}
\end{twierdzenie}
\begin{proof}
Dla $p=1$ z nierówności trójkąta dla wartości
bezwzględnej i liniowości całki otrzymujemy
$\int_X|f+g|\,\mathrm d\mu(x)\le\int_X|f|\,\mathrm d\mu(x)+
\int_X|g|\,\mathrm d\mu(x)$.
Dla $p=\infty$ z
$|f+g|\le\|f\|_\infty+\|g\|_\infty$ prawie wszędzie
wynika teza.

Niech $1<p<\infty$, a $q=p/(p-1)$.
Na początek mamy
$|f+g|^p\le 2^{p-1}(|f|^p+|g|^p)$,
co wynika z wypukłości $t\mapsto t^p$ na $[0,\infty)$.
Ta nierówność zapewnia $f+g\in L^p$, zanim zastosujemy Höldera.
Następnie dla $s=f+g$ korzystamy z
$|s|^p\le |f||s|^{p-1}+|g||s|^{p-1}$.
Ponieważ $(p-1)q=p$, funkcja $|s|^{p-1}$ należy
do $L^q$ i $\||s|^{p-1}\|_q=\|s\|_p^{p-1}$.
Hölder daje
\[
\|s\|_p^p
\le(\|f\|_p+\|g\|_p)\|s\|_p^{p-1}.
\]
Gdy $\|s\|_p=0$, teza jest natychmiastowa;
w pozostałym przypadku dzielimy przez $\|s\|_p^{p-1}$.
\end{proof}

\begin{wniosek}[Przestrzeń normowana]
\label{cor:lp-normed-14}
Dla $1\le p\le\infty$ wzory z
definicji~\ref{def:lp-14} i \ref{def:linf-14}
określają normę na zespolonej przestrzeni liniowej $L^p(\mu)$.
\end{wniosek}
\begin{proof}
Nieujemność jest oczywista; rozdzielność wynika
dla $p<\infty$ z faktu, że całka funkcji nieujemnej
równa zeru oznacza znikanie jej prawie wszędzie,
a dla $p=\infty$ z lematu~\ref{lem:esssup-14}.
Jednorodność $\|cf\|_p=|c|\|f\|_p$ wynika
z jednorodności całki lub kresu istotnego.
Nierówność trójkąta jest treścią twierdzenia
\ref{thm:minkowski-14}. Operacje na klasach
są niezależne od wyboru przedstawiciela.
\end{proof}

\section{Twierdzenie Riesza--Fischera: zupełność przestrzeni $L^p$}

\begin{twierdzenie}[Kryterium zbieżności szeregów w $L^p$]
\label{thm:series-lp-14}
Niech $(X,\mathcal F,\mu)$ będzie przestrzenią z miarą,
$1\le p\le\infty$ i $(u_n)_{n\ge1}\subset L^p(\mu)$.
Jeżeli $\sum_{n=1}^\infty\|u_n\|_p<\infty$,
to istnieje zbiór mierzalny $N$ o $\mu(N)=0$,
poza którym szereg $\sum_{n=1}^\infty u_n(x)$
jest bezwzględnie zbieżny. Jego suma określa element $u\in L^p(\mu)$
i $\sum_{n=1}^N u_n\to u$ w normie $L^p$.
\end{twierdzenie}
\begin{proof}
Dla $p<\infty$ kładziemy $v_N(x)=\sum_{n=1}^N|u_n(x)|$.
Nierówność Minkowskiego daje
$\|v_N\|_p\le\sum_{n=1}^N\|u_n\|_p\le M$,
gdzie $M=\sum_{n\ge1}\|u_n\|_p<\infty$.
Ponieważ $v_N\uparrow v:=\sum_{n\ge1}|u_n|$,
twierdzenie Beppo Leviego zastosowane do $v_N^p$ daje
$\int_Xv^p\,\mathrm d\mu(x)\le M^p$.
W konsekwencji $v<\infty$ prawie wszędzie.
Dla $x$ poza tym zbiorem wyjątkowym szereg
$\sum_nu_n(x)$ zbiega bezwzględnie do $u(x)$.
Wszystkie sumy częściowe są mierzalne, więc $u$
można określić jako funkcję mierzalną poza zbiorem
wyjątkowym i równą zero na nim. Mamy $|u|\le v$,
zatem $u\in L^p$.

Dla ogona $u-\sum_{n\le N}u_n$ mamy punktowo
$|u-\sum_{n\le N}u_n|\le\sum_{n>N}|u_n|$.
Stosując Minkowskiego do skończonych ogonów
i Fatou po przejściu z końcem ogona do nieskończoności,
otrzymujemy
$\|u-\sum_{n\le N}u_n\|_p\le\sum_{n>N}\|u_n\|_p\to0$.

Dla $p=\infty$ wybieramy przedstawicieli $u_n$
i jeden wspólny zbiór zerowy $N$, poza którym
$|u_n(x)|\le\|u_n\|_\infty$ dla każdego $n$.
Na $X\setminus N$ szereg zbiega bezwzględnie
oraz jednostajnie, gdyż suma $\sum_n\|u_n\|_\infty$
jest skończona. Otrzymujemy mierzalną sumę $u$,
$\|u\|_\infty\le M$ i
$\|u-\sum_{n\le N}u_n\|_\infty\le\sum_{n>N}\|u_n\|_\infty$.
\end{proof}

\begin{twierdzenie}[Riesz--Fischer: zupełność $L^p$]
\label{thm:riesz-fischer-14}
Dla każdej przestrzeni z miarą $(X,\mathcal F,\mu)$
i każdego $1\le p\le\infty$ przestrzeń normowana
$L^p(\mu)$ jest zupełna, a więc jest przestrzenią Banacha.
Dla $p=2$ jest ponadto przestrzenią Hilberta
z iloczynem skalarnym
$\langle f,g\rangle:=\int_X f(x)\overline{g(x)}\,\mathrm d\mu(x)$.
\end{twierdzenie}
\begin{proof}
Niech $(f_n)$ będzie ciągiem Cauchy'ego w $L^p$.
Indukcyjnie wybieramy rosnący ciąg indeksów $(n_k)$
taki, że $\|f_{n_{k+1}}-f_{n_k}\|_p\le2^{-k}$.
Wówczas szereg $\sum_{k\ge1}(f_{n_{k+1}}-f_{n_k})$
ma sumę norm mniejszą niż $\sum_k2^{-k}<\infty$.
Na mocy twierdzenia~\ref{thm:series-lp-14}
zbiega więc w $L^p$ do pewnego elementu $w$,
a zatem $f_{n_k}\to f:=f_{n_1}+w$ w $L^p$.
Aby wykazać zbieżność całego ciągu, ustalamy
$\varepsilon>0$ i wybieramy $N$ tak, aby
$\|f_m-f_n\|_p<\varepsilon/2$ dla $m,n\ge N$.
Wybieramy także $k$ z $n_k\ge N$ i
$\|f_{n_k}-f\|_p<\varepsilon/2$.
Wtedy dla $n\ge N$ nierówność trójkąta daje
$\|f_n-f\|_p\le\|f_n-f_{n_k}\|_p+
\|f_{n_k}-f\|_p<\varepsilon$.

Dla $p=2$ Hölder gwarantuje, że podana całka
jest skończona. Jest ona liniowa w pierwszym
argumencie, antyliniowa w drugim, hermitowska
i dodatnio określona na klasach. Wytworzona przez nią
norma to $\|f\|_2$, a udowodniona już zupełność
oznacza zupełność Hilberta.
\end{proof}

\begin{uwaga}[Zbieżność punktowa w dowodzie]
Istotnym elementem dowodu jest istnienie
\emph{podciągu} zbieżnego prawie wszędzie.
Nie wynika stąd, że dowolny ciąg zbieżny w $L^p$
zbiega punktowo prawie wszędzie; zobaczymy niżej
kontrprzykład związany z rodzajami zbieżności.
\end{uwaga}

\section{Nierówność Jensena i jej znaczenie}
\begin{twierdzenie}[Jensen dla miary probabilistycznej]
\label{thm:jensen-14}
Niech $(X,\mathcal F,P)$ będzie przestrzenią
probabilistyczną, $f\in L^1(P;\mathbb R)$,
a $\varphi:\mathbb R\to\mathbb R$ funkcją wypukłą.
Wtedy $\varphi\circ f$ jest mierzalna, jej część ujemna
jest całkowalna, a więc całka
$\int_X\varphi(f(x))\,\mathrm dP(x)\in(-\infty,\infty]$
jest dobrze określona i
\begin{equation}\label{eq:jensen-14}
\varphi\left(\int_Xf(x)\,\mathrm dP(x)\right)
\le\int_X\varphi(f(x))\,\mathrm dP(x).
\end{equation}
\end{twierdzenie}
\begin{proof}
Niech $m=\int_Xf\,\mathrm dP(x)$.
Funkcja wypukła skończona na całej prostej jest ciągła,
a więc $\varphi\circ f$ jest mierzalna.
Dla funkcji wypukłej na $\mathbb R$ w każdym punkcie $m$
istnieje prosta podpierająca: jest liczba $a\in\mathbb R$
taka, że $\varphi(t)\ge\varphi(m)+a(t-m)$
dla wszystkich $t\in\mathbb R$.
Istnienie takiego $a$ wynika z nierówności między
lewym i prawym nachyleniem siecznych funkcji wypukłej;
można wybrać dowolną liczbę między jej lewą
i prawą pochodną jednostronną w $m$.
Po podstawieniu $t=f(x)$ dostajemy
$\varphi(f(x))\ge\varphi(m)+a(f(x)-m)$.
Prawa strona jest całkowalna, więc część ujemna
$\varphi(f)$ ma skończoną całkę. Możemy całkować
nierówność, otrzymując
$\int_X\varphi(f)\,\mathrm dP(x)
\ge\varphi(m)+a(\int_Xf\,\mathrm dP(x)-m)=\varphi(m)$.
\end{proof}

\begin{przyklad}[Wartość średnia a momenty]
Dla $1\le p<\infty$ funkcja $\varphi(t)=|t|^p$
jest wypukła. Gdy $P(X)=1$ i $f\in L^p(P)$,
Jensen daje
$|\int_X f\,\mathrm dP(x)|^p\le
\int_X|f(x)|^p\,\mathrm dP(x)$ dla $f$ rzeczywistej.
Dla $f$ zespolonej analogiczna nierówność wynika
z nierówności trójkąta dla całki i Höldera
z funkcją stałą $1$ (należącą do $L^q(P)$).
\end{przyklad}

\section{Zależności między przestrzeniami $L^p$}

\begin{twierdzenie}[Włożenia przy skończonej mierze]
\label{thm:lp-embedding-finite-14}
Niech $\mu(X)<\infty$ i $1\le p<q\le\infty$.
Wtedy $L^q(\mu)\subseteq L^p(\mu)$ oraz
\begin{equation}\label{eq:finite-embedding-14}
\|f\|_p\le\mu(X)^{\,1/p-1/q}\|f\|_q
\quad(f\in L^q(\mu)),
\end{equation}
przy umowie $1/\infty=0$ i naturalnej interpretacji
w przypadku $\mu(X)=0$.
\end{twierdzenie}
\begin{proof}
Jeśli $\mu(X)=0$, wszystkie przestrzenie są
zerowe i teza jest natychmiastowa. Załóżmy więc
$0<\mu(X)<\infty$.
Dla $q<\infty$ stosujemy Höldera do
$|f|^p\cdot1$ z wykładnikami $q/p$ i $q/(q-p)$:
\[
\int_X|f|^p\,\mathrm d\mu(x)
\le\left(\int_X|f|^q\,\mathrm d\mu(x)\right)^{p/q}
\mu(X)^{1-p/q}.
\]
Podnosimy obie strony do potęgi $1/p$.
Dla $q=\infty$ nierówność
$|f|^p\le\|f\|_\infty^p$ prawie wszędzie
i całkowanie dają wynik.
\end{proof}

\begin{przyklad}[Na nieskończonej mierze włożenia zawodzą]
Na $(1,\infty)$ z miarą Lebesgue'a
funkcja $f(x)=x^{-\alpha}$ spełnia
$f\in L^p$ dokładnie wtedy, gdy $\alpha p>1$.
Na $(0,1)$ funkcja $g(x)=x^{-\beta}$
należy do $L^p$ dokładnie wtedy, gdy $\beta p<1$.
Dla $1\le p<q<\infty$ wybór
$1/q<\alpha\le1/p$ daje na $(1,\infty)$
funkcję z $L^q\setminus L^p$,
a wybór $1/q\le\beta<1/p$ daje na $(0,1)$
funkcję z $L^p\setminus L^q$.
Po przedłużeniu przez zero odpowiednie przykłady
działają na całej $\mathbb R$.
Zatem na przestrzeni nieskończonej miary
w ogólności nie zachodzi żadna z inkluzji między
różnymi $L^p$.
\end{przyklad}

\begin{twierdzenie}[Zbieżność w $L^p$ a zbieżność według miary]
\label{thm:lp-to-measure-14}
Niech $1\le p<\infty$ i niech $f_n,f\in L^p(\mu)$.
Jeżeli $\|f_n-f\|_p\to0$, to $f_n\to f$
według miary $\mu$, tj. dla każdego $\varepsilon>0$
$\mu(\{|f_n-f|>\varepsilon\})\to0$.
Dla $p=\infty$ zachodzi silniejsze stwierdzenie:
$\|f_n-f\|_\infty\to0$ implikuje
$\mu(\{|f_n-f|>\varepsilon\})=0$ dla dostatecznie
dużych $n$.
\end{twierdzenie}
\begin{proof}
Dla $p<\infty$ nierówność Czebyszewa daje
$\varepsilon^p\mu(\{|f_n-f|>\varepsilon\})
\le\int_X|f_n-f|^p\,\mathrm d\mu(x)=\|f_n-f\|_p^p$.
Prawa strona zbiega do zera.
Dla $p=\infty$ wybieramy $n$ tak duże, że
$\|f_n-f\|_\infty<\varepsilon$ i korzystamy
z definicji kresu istotnego.
\end{proof}

\begin{przyklad}[Zbieżność według miary nie wystarcza]
Na $(0,1)$ z miarą Lebesgue'a połóżmy
$f_n(x)=n^{1/p}\chi_{(0,1/n)}(x)$ dla $p<\infty$.
Wtedy $f_n\to0$ według miary i prawie wszędzie,
ale $\|f_n\|_p=1$ dla każdego $n$.
Zbieżność według miary nie pociąga więc zbieżności
w normie, nawet na przestrzeni probabilistycznej.
\end{przyklad}

\section{Interpolacja pomiędzy normami $L^p$}

\begin{twierdzenie}[Nierówność interpolacyjna]
\label{thm:lp-interpolation-14}
Niech $(X,\mathcal F,\mu)$ będzie przestrzenią z miarą,
$1\le p<r<q\le\infty$, i niech $\theta\in(0,1)$
będzie jedyną liczbą spełniającą
$1/r=\theta/p+(1-\theta)/q$, przy czym $1/\infty=0$.
Dla każdej funkcji $f\in L^p(\mu)\cap L^q(\mu)$
zachodzi $f\in L^r(\mu)$ oraz
\begin{equation}\label{eq:lp-interpolation-14}
\|f\|_r\le\|f\|_p^{\theta}\|f\|_q^{1-\theta}.
\end{equation}
W szczególności nie zakładamy $\mu(X)<\infty$.
\end{twierdzenie}
\begin{proof}
Najpierw załóżmy $q<\infty$ i obie normy po prawej
dodatnie. Wybieramy sprzężone wykładniki
$a=p/(\theta r)>1$ oraz $b=q/((1-\theta)r)>1$;
warunek na $\theta$ jest równoważny $1/a+1/b=1$.
Z twierdzenia Höldera dla czynników
$|f|^{\theta r}$ i $|f|^{(1-\theta)r}$ otrzymujemy
\[
\int_X|f(x)|^r\,\mathrm d\mu(x)
\le\left(\int_X|f(x)|^p\,\mathrm d\mu(x)\right)^{\theta r/p}
\left(\int_X|f(x)|^q\,\mathrm d\mu(x)\right)^{(1-\theta)r/q}.
\]
Podnosząc obie strony do potęgi $1/r$, dostajemy tezę.
Gdy $q=\infty$, warunek na $\theta$ daje
$\theta r=p$. Wówczas
$|f|^r\le\|f\|_\infty^{r-p}|f|^p$
prawie wszędzie, z czego po całkowaniu wynika
\eqref{eq:lp-interpolation-14}. Jeżeli któraś
z norm po prawej jest zerowa, funkcja znika
prawie wszędzie i nierówność jest oczywista.
\end{proof}

\begin{przyklad}[Interpolacja bez skończoności miary]
Na $X=\mathbb R^d$ założenie $f\in L^1\cap L^\infty$
gwarantuje $f\in L^r$ dla wszystkich $1<r<\infty$,
a dokładniej $\|f\|_r\le\|f\|_1^{1/r}
\|f\|_\infty^{1-1/r}$. Nie wolno jednak wnioskować
z samego $f\in L^q$ o przynależności do $L^p$
przy $p<q$ bez dodatkowych założeń o mierze.
\end{przyklad}

\section{Aproksymacja i gęstość wybranych klas funkcji}

\begin{twierdzenie}[Gęstość funkcji prostych]
\label{thm:simple-dense-lp-14}
Niech $(X,\mathcal F,\mu)$ będzie dowolną przestrzenią
z miarą i $1\le p<\infty$. Wtedy w $L^p(\mu)$ gęste są
klasy funkcji prostych, przyjmujących skończenie wiele
wartości i mających nośnik w zbiorze skończonej miary.
Dla $p=\infty$ funkcje proste są gęste w $L^\infty(\mu)$,
ale bez dodatkowych założeń nie można żądać, aby
ich nośniki miały skończoną miarę.
\end{twierdzenie}
\begin{proof}
Najpierw niech $f\ge0$ i $f\in L^p$ dla $p<\infty$.
Dla $n\ge1$ kładziemy $E_n=\{f>1/n\}$.
Z nierówności Czebyszewa mamy
$\mu(E_n)\le n^p\int_Xf^p\,\mathrm d\mu(x)<\infty$.
Funkcje $f\chi_{E_n}\chi_{\{f\le n\}}$
zbiegają punktowo do $f$ poza zbiorem, na którym
$f=\infty$ (zerowym, gdy $f\in L^p$), a ich
różnice z $f$ są zdominowane przez $f$.
Twierdzenie Lebesgue'a o zbieżności zmajoryzowanej
zastosowane do $p$-tych potęg modułów różnic pokazuje,
że obcięcia zbiegają do $f$ w $L^p$.

Dla ustalonych $n$ obcięta funkcja jest ograniczona
przez $n$ i ma nośnik w $E_n$ o skończonej mierze.
Aproksymujemy ją z dołu funkcjami prostymi
$o$ skoku wartości co najwyżej $1/m$; błąd
jest co najwyżej $1/m$ na $E_n$, a zatem jego
norma $L^p$ jest co najwyżej
$\mu(E_n)^{1/p}/m\to0$. W przypadku funkcji
zespolonej stosujemy aproksymację osobno do
części rzeczywistej i urojonej, a następnie
korzystamy z Minkowskiego.

Dla $p=\infty$ każdą funkcję istotnie ograniczoną
można poza zbiorem zerowym aproksymować
jednostajnie funkcjami prostymi, przez podział
ograniczonego prostokąta wartości w $\mathbb C$
na kwadraty o dowolnie małych średnicach.
Jeżeli $\mu(X)=\infty$, funkcja $1\in L^\infty(\mu)$,
ale odległość w normie $L^\infty$ między nią
a każdą funkcją zerową poza zbiorem skończonej
miary wynosi co najmniej $1$.
\end{proof}

\begin{twierdzenie}[Gęstość funkcji gładkich w przestrzeni euklidesowej]
\label{thm:cc-dense-lp-14}
Dla $d\ge1$ oraz $1\le p<\infty$ zbiór
$C_c^\infty(\mathbb R^d)$ jest gęsty w
$L^p(\mathbb R^d,\mathcal L_d,\lambda_d)$.
Nie jest gęsty w $L^\infty(\mathbb R^d)$.
\end{twierdzenie}
\begin{proof}
Przypadek $1\le p<\infty$ został wykazany
w rozdziale~\ref{ch:regularnosc-funkcji},
w twierdzeniu~\ref{thm:cc-in-lp-10}:
korzysta on z regularności miary Lebesgue'a,
aproksymacji funkcji charakterystycznych zbiorów mierzalnych
funkcjami ciągłymi o zwartym nośniku,
a następnie z wygładzania przez mollifikatory.
Przytoczenie tego wyniku nie wykorzystuje
żadnego twierdzenia z późniejszych rozdziałów.

Dla $p=\infty$ rozważmy funkcję stałą $1$.
Dla dowolnej $\varphi\in C_c^\infty(\mathbb R^d)$
poza zwartym nośnikiem $\varphi$ zachodzi $\varphi=0$,
a zbiór ten ma dodatnią miarę. Stąd
$\|1-\varphi\|_\infty\ge1$, więc gęstość nie zachodzi.
\end{proof}

\begin{twierdzenie}[Ciągłość translacji w $L^p(\mathbb R^d)$]
\label{thm:translation-lp-14}
Niech $1\le p<\infty$ i $f\in L^p(\mathbb R^d)$.
Dla $h\in\mathbb R^d$ kładziemy
$(\tau_hf)(x)=f(x-h)$.
Wówczas $\|\tau_hf-f\|_p\to0$ przy $h\to0$.
\end{twierdzenie}
\begin{proof}
Niezmienniczość miary Lebesgue'a względem translacji
daje $\|\tau_hg\|_p=\|g\|_p$.
Jeżeli $g\in C_c^\infty(\mathbb R^d)$, to jest
jednostajnie ciągła, a dla $|h|\le1$
nośniki $g(\cdot-h)$ oraz $g$ leżą we wspólnym
zbiorze zwartym $K$ o skończonej mierze.
Mamy zatem
$\|\tau_hg-g\|_p\le
\lambda_d(K)^{1/p}\sup_x|g(x-h)-g(x)|\to0$.

Dla dowolnego $f\in L^p$ i $\varepsilon>0$
wybieramy z gęstości $g\in C_c^\infty$ tak, że
$\|f-g\|_p<\varepsilon$.
Nierówność trójkąta prowadzi do
$\|\tau_hf-f\|_p\le
\|\tau_h(f-g)\|_p+\|\tau_hg-g\|_p+\|g-f\|_p
<2\varepsilon+\|\tau_hg-g\|_p$.
Biorąc granicę górną przy $h\to0$ i następnie
$\varepsilon\downarrow0$, otrzymujemy tezę.
\end{proof}

\begin{przyklad}[Dlaczego $p=\infty$ jest inny]
Dla $f=\chi_{(0,\infty)}$ na $\mathbb R$
$\|\tau_hf-f\|_\infty=1$ dla każdego $h\ne0$.
Zatem translacje nie są silnie ciągłe w
$L^\infty(\mathbb R)$.
\end{przyklad}

\section{Separowalność: kiedy istnieje przeliczalna rodzina gęsta?}

\begin{twierdzenie}[Separowalność $L^p$]
\label{thm:separable-lp-14}
Niech $(X,\mathcal F,\mu)$ będzie przestrzenią
z miarą $\sigma$-skończoną, przy czym $\mathcal F$
jest generowana przez przeliczalną rodzinę zbiorów.
Wówczas $L^p(\mu)$ jest separowalna dla
każdego $1\le p<\infty$.
\end{twierdzenie}
\begin{proof}
Wybieramy rosnące wyczerpanie $X_n\uparrow X$
zbiorami o skończonej mierze. Do przeliczalnej
rodziny generującej $\mathcal F$ dodajemy zbiory $X_n$
i oznaczamy przez $\mathcal A$ algebrę przez nią generowaną.
Ponieważ wszystkie operacje w algebrze są skończone,
$\mathcal A$ jest przeliczalna i $\sigma(\mathcal A)=\mathcal F$.

Na każdym $X_n$ rozważamy klasę zbiorów $E\in\mathcal F|_{X_n}$,
dla których dla każdego $\varepsilon>0$ istnieje
$A\in\mathcal A|_{X_n}$ z
$\mu(E\mathbin\triangle A)<\varepsilon$.
Klasa ta jest $\sigma$-algebrą: dla dwóch zbiorów
aproksymowanych zbiorami $A,B$ z algebry zachodzi
$\mu((E\cup F)\mathbin\triangle(A\cup B))\le
\mu(E\mathbin\triangle A)+\mu(F\mathbin\triangle B)$;
analogicznie traktujemy dopełnienia względem $X_n$.
Dla $E_k\uparrow E$ mamy
$\mu(E\setminus E_k)\to0$ dzięki
$\mu(X_n)<\infty$. Stąd przeliczalne sumy
aproksymujemy najpierw skończonymi, a następnie
zbiorami z algebry. Klasa zawiera $\mathcal A|_{X_n}$,
a więc jest całą $\mathcal F|_{X_n}$.

Z twierdzenia~\ref{thm:simple-dense-lp-14}
wystarczy aproksymować funkcje proste o nośniku
skończonej miary. Dzięki $X_n\uparrow X$ i twierdzeniu
Beppo Leviego o ogonach ich część poza $X_n$
można uczynić małą w $L^p$. Następnie każdą
funkcję charakterystyczną na $X_n$ aproksymujemy
w $L^p$ przez funkcje charakterystyczne zbiorów z
$\mathcal A|_{X_n}$, bo
$\|\chi_E-\chi_A\|_p
=\mu(E\mathbin\triangle A)^{1/p}$.
Wreszcie współczynniki zespolone aproksymujemy
liczbami o wymiernych częściach rzeczywistych
i urojonych. Zbiór wszystkich takich skończonych
kombinacji funkcji charakterystycznych zbiorów z przeliczalnej
algebry (z obcięciem do wybranych $X_n$) jest
przeliczalny i gęsty w $L^p$.
\end{proof}

\begin{przyklad}[Dwie przeszkody dla separowalności]
Dla $X=[0,1]$ z miarą Lebesgue'a przestrzeń
$L^\infty(X)$ nie jest separowalna: funkcje
$\chi_{(0,t)}$, $t\in(0,1)$, są parami
oddalone w normie $L^\infty$ o $1$.
Zbiór nieprzeliczalny, którego różne elementy są
od siebie oddalone co najmniej o $1$, nie może
mieścić się w przestrzeni separowalnej.

Również dla $p<\infty$ istnieją przestrzenie
nieseparowalne: jeśli $I$ jest nieprzeliczalny,
a $\kappa$ jest miarą liczącą na $\mathcal P(I)$,
to funkcje $e_i=\chi_{\{i\}}\in L^p(\kappa)$
dla $i\in I$ mają parami odległość $2^{1/p}$.
Tutaj założenie o przeliczalnym generowaniu
$\sigma$-algebry zawodzi.
\end{przyklad}

\section{Co dalej: dualność i słaba zbieżność}
Dzięki zupełności możemy traktować $L^p$ jako
przestrzenie Banacha, a $L^2$ jako przestrzeń Hilberta.
Nie przesądza to jednak ich struktury dualnej.
W następnym rozdziale zbadamy, dla jakich
wykładników i przestrzeni z miarą każdy ciągły
funkcjonał liniowy na $L^p$ można zapisać jako
$f\mapsto\int_Xf(x)g(x)\,\mathrm d\mu(x)$.
Będziemy tam precyzyjnie odróżniać
przestrzeń dualną od samego zbioru funkcji reprezentujących,
a także omówimy refleksywność i zbieżność słabą.

% END chapters/pl/measure_theory/14-rozdzial.tex


% BEGIN chapters/pl/measure_theory/15-rozdzial.tex
% ROZDZIAŁ 15 — pełny, z dowodami; nie zakłada twierdzenia Radona–Nikodyma
\chapter[Dualność i geometria przestrzeni $L^p$]{Dualność, wypukłość i słaba zbieżność w przestrzeniach $L^p$}
\label{ch:dualnosc-lp}

\section{Przestrzeń dualna i podstawowy problem reprezentacji}
Niech $(X,\mathcal F,\mu)$ będzie przestrzenią z miarą. Przestrzeń
$L^p(\mu)$ zbudowaliśmy w rozdziale~\ref{ch:przestrzenie-lp} jako
przestrzeń Banacha. Celem jest znalezienie \emph{wszystkich} ciągłych
funkcjonałów liniowych na tej przestrzeni. Należy odróżnić
wykazanie, że pewien wzór daje funkcjonał, od twierdzenia, że
każdy funkcjonał ma tę postać.

\begin{definicja}[Dualność i norma funkcjonału]
\label{def:dual-lp-15}
Dla zespolonej przestrzeni Banacha $E$ jej \emph{przestrzenią dualną}
$E^*$ nazywamy przestrzeń wszystkich ciągłych \emph{zespolono-liniowych}
odwzorowań $\Lambda:E\to\mathbb C$, z normą
$\|\Lambda\|:=\sup\{|\Lambda(f)|:\|f\|_E\le1\}$.
W całym rozdziale parowanie $L^p$ z $L^q$ zapisujemy
$\langle f,g\rangle=\int_X f(x)g(x)\,\mathrm d\mu(x)$;
\emph{nie} wstawiamy tu sprzężenia zespolonego. Jest to parowanie
dwuliniowe, a nie iloczyn skalarny na $L^2$.
\end{definicja}

\begin{lemat}[Izometryczne zanurzenie do dualnej]
\label{lem:canonical-lpq-15}
Niech $1\le p<\infty$ i $q\in(1,\infty]$ będzie wykładnikiem
sprzężonym ($1/p+1/q=1$). Jeśli $p=1$, dodatkowo
zakładamy $\sigma$-skończoność $\mu$. Dla każdego $g\in L^q(\mu)$ wzór
$\Lambda_g(f)=\int_X f(x)g(x)\,\mathrm d\mu(x)$ definiuje element
$(L^p(\mu))^*$ i zachodzi $\|\Lambda_g\|=\|g\|_q$.
Dla $1<p<\infty$ nie zakładamy $\sigma$-skończoności.
\end{lemat}
\begin{proof}
Z nierówności Höldera (twierdzenie~\ref{thm:holder-14}) mamy
$|\Lambda_g(f)|\le\|f\|_p\|g\|_q$, a więc operator jest ciągły.
Dla $1<p<\infty$, jeśli $g\ne0$, definiujemy
$f=\overline{g}|g|^{q-2}$ w punktach, gdzie $g\ne0$, i $f=0$
poza nimi. Gdy $q<2$, zapis w punktach $g=0$ zastępujemy wartością zero;
moduł tak określonej funkcji wynosi $|g|^{q-1}$.
Ponieważ $(q-1)p=q$, mamy $f\in L^p$,
$\Lambda_g(f)=\|g\|_q^q$ i $\|f\|_p=\|g\|_q^{q-1}$.
Wynika stąd nierówność odwrotna $\|\Lambda_g\|\ge\|g\|_q$.

Jeśli $p=1$, kładziemy $M=\|g\|_\infty$. Dla dowolnego
$0<a<M$ zbiór $A=\{|g|>a\}$ ma dodatnią miarę.
Z założonej $\sigma$-skończoności wynika, że istnieje
mierzalny zbiór $B\subseteq A$ o $0<\mu(B)<\infty$:
przecinamy $A$ z jednym ze zbiorów wyczerpania
$X=\bigcup_nX_n$, $\mu(X_n)<\infty$.
Wybierając $f=\chi_B\overline g/(|g|\mu(B))$ na $B$
i $f=0$ poza $B$, otrzymujemy $\|f\|_1=1$ i
$|\Lambda_g(f)|\ge a$. Przechodzimy z $a\uparrow M$.
\end{proof}

\begin{uwaga}[Dlaczego wymagane jest doprecyzowanie dla $p=1$?]
W przestrzeni z miarą, która nie jest półskończona, funkcja $g$
może mieć niezerowy kres istotny, ale całkować się do zera
z każdą funkcją $f\in L^1$. Na przykład na jednoelementowej
przestrzeni $\{*\}$ z miarą $\mu(\{*\})=\infty$ mamy
$L^1=\{0\}$ i $L^\infty\cong\mathbb C$.
Zatem ogólne twierdzenie o izometrycznym zanurzeniu z poprzedniego
lematu należy dla $p=1$ rozumieć tylko przy założeniu
$\sigma$-skończoności. Jest to istotna granica twierdzenia;
przy $1<p<\infty$ problem nie występuje.
\end{uwaga}

\section{Nierówności wypukłości w $L^p$}
W szczególności skonstruujemy reprezentację funkcjonałów
\emph{bez użycia twierdzenia Radona--Nikodyma}. Narzędziem
będzie ścisła wypukłość potęgi normy.

\begin{lemat}[Dwie nierówności dla punktu środkowego]
\label{lem:midpoint-lp-15}
Niech $f,h\in L^p(\mu)$ i $1<p<\infty$. Piszemy
$m=(f+h)/2$ oraz $v=(f-h)/2$.
\begin{enumerate}[label=\textup{(\roman*)}]
\item Jeśli $p\ge2$, to
$\|m\|_p^p+\|v\|_p^p\le(\|f\|_p^p+\|h\|_p^p)/2$.
\item Jeśli $1<p<2$ i $\|f\|_p+\|h\|_p>0$, to
\begin{equation}\label{eq:midpoint-low-15}
\frac{\|f\|_p^p+\|h\|_p^p}{2}-\|m\|_p^p
\ge \frac{p(p-1)}2\,
\frac{\|v\|_p^2}{(\|f\|_p+\|h\|_p)^{2-p}}.
\end{equation}
\end{enumerate}
\end{lemat}
\begin{proof}
Dla $p\ge2$ wystarczy wykazać dla $a,b\in\mathbb C$ nierówność
$|a+b|^p+|a-b|^p\le2^{p-1}(|a|^p+|b|^p)$.
Ponieważ $r=p/2\ge1$ i $u^r+v^r\le(u+v)^r$ dla $u,v\ge0$,
\[
|a+b|^p+|a-b|^p
\le (|a+b|^2+|a-b|^2)^{p/2}
=2^{p/2}(|a|^2+|b|^2)^{p/2}.
\]
Wypukłość $t\mapsto t^{p/2}$ daje
$(|a|^2+|b|^2)^{p/2}\le2^{p/2-1}(|a|^p+|b|^p)$.
Po podzieleniu przez $2^p$ i scałkowaniu otrzymujemy (i).

Dla $1<p<2$ identyfikujemy $\mathbb C$ z $\mathbb R^2$.
Dla $z\ne0$ hesjan funkcji $\psi(z)=|z|^p$ ma postać
$D^2\psi(z)=p|z|^{p-2}I+p(p-2)|z|^{p-4}z\otimes z$.
Najmniejsza wartość własna wynosi
$p(p-1)|z|^{p-2}$. Niech $a,b\in\mathbb C$,
$u=(a+b)/2$, $v=(a-b)/2$ i $S=|a|+|b|$.
Wzdłuż odcinka $t\mapsto u+tv$, $-1\le t\le1$,
zachodzi $|u+tv|\le S$. Dla $v\ne0$ funkcja
$t\mapsto |u+tv|^p$ ma bezwzględnie ciągłą pochodną
na $[-1,1]$. Istotnie, poza ewentualnym pojedynczym
punktem $t_0$, w którym $u+t_0v=0$, jest ona klasy $C^2$,
a jeśli taki punkt istnieje, to druga pochodna jest
w jego otoczeniu oszacowana przez wielokrotność
$|t-t_0|^{p-2}$. Jest to funkcja całkowalna, gdyż
$p-2>-1$; ponadto pierwsza pochodna jest ciągła
również w $t_0$. Można więc zastosować dwukrotnie
wzór Newtona--Leibniza na odcinku, rozbijając go
w razie potrzeby w punkcie $t_0$. Z nierówności
$D^2\psi(z)[v,v]\ge p(p-1)|z|^{p-2}|v|^2$
oraz $p-2<0$ otrzymujemy
\[
\frac{|a|^p+|b|^p}{2}-|u|^p
\ge \frac{p(p-1)}2\,|v|^2S^{p-2}.
\]
W przypadku $S=0$ obie strony rozumiemy jako zero.
Przy $S>0$ wyrażenie z prawej jest nieujemne i mierzalne.
Aby oszacować jego całkę od dołu, używamy Höldera
z wykładnikami $2/p$ i $2/(2-p)$:
\[
\int_X|v|^p\,\mathrm d\mu(x)
\le\left(\int_X |v|^2S^{p-2}\,\mathrm d\mu(x)\right)^{p/2}
\left(\int_XS^p\,\mathrm d\mu(x)\right)^{(2-p)/2}.
\]
Jeżeli pierwsza całka jest nieskończona, teza jest oczywista.
W przeciwnym razie podnosimy obie strony do potęgi $2/p$.
Minkowski zapewnia $\|S\|_p\le\|f\|_p+\|h\|_p$.
Otrzymujemy (ii).
\end{proof}

\begin{twierdzenie}[Jednostajna wypukłość]
\label{thm:uniform-convexity-15}
Dla dowolnej przestrzeni z miarą $(X,\mathcal F,\mu)$
i każdego $1<p<\infty$ przestrzeń $L^p(\mu)$ jest
\emph{jednostajnie wypukła}: dla każdego $\varepsilon>0$
istnieje $\delta>0$ taka, że jeżeli $\|f\|_p=\|h\|_p=1$
i $\|f-h\|_p\ge\varepsilon$, to
$\|(f+h)/2\|_p\le1-\delta$.
\end{twierdzenie}
\begin{proof}
Dla $p\ge2$ lemat~\ref{lem:midpoint-lp-15} daje
$\|(f+h)/2\|_p^p\le1-(\varepsilon/2)^p$.
Dla $1<p<2$ wzór \eqref{eq:midpoint-low-15} daje
$\|(f+h)/2\|_p^p\le1-c_p\varepsilon^2$,
gdzie $c_p=p(p-1)/32>0$ jest dopuszczalną stałą.
Ponieważ $\|f-h\|_p\le2$, dla $0<\varepsilon\le2$
prawa strona jest dodatnia. W obu przypadkach
możemy przyjąć $\delta=1-(1-c\varepsilon^r)^{1/p}>0$
z odpowiednimi $c,r$, jak wyżej; dla $\varepsilon>2$
warunek $\|f-h\|_p\ge\varepsilon$ jest niemożliwy.
\end{proof}

\section{Metoda wariacyjna i dualność dla $1<p<\infty$}
\begin{lemat}[Minimum energii]
\label{lem:variational-lp-15}
Niech $1<p<\infty$ i $\Lambda\in(L^p(\mu))^*$.
Funkcjonał rzeczywisty
$J(f)=\|f\|_p^p/p-\operatorname{Re}\Lambda(f)$
osiąga minimum w pewnym $u\in L^p(\mu)$.
\end{lemat}
\begin{proof}
Z oszacowania $J(f)\ge\|f\|_p^p/p-\|\Lambda\|\|f\|_p$
wynika, że $J$ jest ograniczony z dołu, a każdy ciąg
minimalizujący $(f_n)$ jest ograniczony w $L^p$:
$\|f_n\|_p\le R$ dla pewnego $R<\infty$.
Ponieważ człon liniowy zachowuje średnie, z
lematu~\ref{lem:midpoint-lp-15} wynika przy $p\ge2$
\[
\frac{J(f_n)+J(f_m)}2-J\left(\frac{f_n+f_m}{2}\right)
\ge\frac1p\left\|\frac{f_n-f_m}{2}\right\|_p^p.
\]
Dla $1<p<2$ analogiczna różnica jest co najmniej
$c(p,R)\|f_n-f_m\|_p^2$, gdzie $c(p,R)>0$;
przy $R=0$ ciąg jest zerowy. Niech $a=\inf J$.
Skoro $J((f_n+f_m)/2)\ge a$, a $J(f_n),J(f_m)\to a$,
prawe strony powyższych nierówności zmierzają do zera.
Zatem $(f_n)$ jest ciągiem Cauchy'ego w $L^p$.
Na mocy twierdzenia Riesza--Fischera
\ref{thm:riesz-fischer-14} zbiega do $u\in L^p$.
Ciągłość normy i $\Lambda$ daje $J(u)=a$.
\end{proof}

\begin{twierdzenie}[Dualność $L^p$ dla $1<p<\infty$]
\label{thm:dual-lp-15}
Niech $(X,\mathcal F,\mu)$ będzie \emph{dowolną} przestrzenią
z miarą, $1<p<\infty$, a $q=p/(p-1)$.
Wówczas każdy $\Lambda\in (L^p(\mu))^*$ ma \emph{jedyną}
reprezentację
\begin{equation}\label{eq:dual-lp-15}
\Lambda(f)=\int_X f(x)g(x)\,\mathrm d\mu(x)
\qquad(f\in L^p(\mu))
\end{equation}
przez pewne $g\in L^q(\mu)$. Zachodzi
$\|\Lambda\|=\|g\|_q$. Nie jest potrzebna
$\sigma$-skończoność ani twierdzenie Radona--Nikodyma.
\end{twierdzenie}
\begin{proof}
Z lematu~\ref{lem:variational-lp-15} istnieje minimum $u$
funkcjonału $J(f)=\|f\|_p^p/p-\operatorname{Re}\Lambda(f)$.
Dla $v\in L^p$ i $t\in\mathbb R$ funkcja $t\mapsto J(u+tv)$
ma minimum przy $t=0$. Jej pochodna w zerze wynosi
\[
\int_X |u(x)|^{p-2}\operatorname{Re}(\overline{u(x)}v(x))
\,\mathrm d\mu(x)-\operatorname{Re}\Lambda(v).
\]
Istnienie tej pochodnej uzasadniamy następująco:
dla $|t|\le1$ iloraz różnicowy $|u+tv|^p/p$ jest
oszacowany przez $(|u|+|v|)^{p-1}|v|$,
a funkcja ta jest całkowalna na mocy Höldera,
ponieważ $u,v\in L^p$. Możemy więc zastosować
twierdzenie o zbieżności zmajoryzowanej. Z minimum
wynika $\operatorname{Re}\Lambda(v)=
\operatorname{Re}\int_Xv\overline u|u|^{p-2}\,\mathrm d\mu(x)$.
Zastępując $v$ przez $iv$, otrzymujemy także
równość części urojonych. Zatem \eqref{eq:dual-lp-15}
zachodzi dla $g=\overline u|u|^{p-2}$ (przyjmujemy $g=0$
tam, gdzie $u=0$). Ponieważ $(p-1)q=p$,
$\int_X|g|^q\,\mathrm d\mu(x)=\int_X|u|^p\,\mathrm d\mu(x)<\infty$.

Lemat~\ref{lem:canonical-lpq-15} daje równość norm.
Jeśli $g_1,g_2$ reprezentują ten sam funkcjonał,
to funkcjonał $\Lambda_{g_1-g_2}$ jest zerowy,
a z izometrii wynika $\|g_1-g_2\|_q=0$.
Jest to jednoznaczność w sensie równości prawie wszędzie.
\end{proof}

\begin{przyklad}[Norma funkcjonału całkowego]
Niech $1<p<\infty$, $X=(0,1)$ z miarą Lebesgue'a
oraz $g(x)=x^{-\alpha}$ dla $0<\alpha<1/q$.
Funkcjonał $\Lambda(f)=\int_{(0,1)}f(x)x^{-\alpha}
\,\mathrm d\lambda_1(x)$ należy do $(L^p)^*$, a
$\|\Lambda\|=(1-\alpha q)^{-1/q}$.
Jeśli $\alpha q\ge1$, funkcja $g$ nie należy do $L^q$;
nie istnieje wtedy \emph{ciągłe na całym $L^p$} rozszerzenie
tego funkcjonału całkowego zachowujące tę samą postać.
\end{przyklad}

\section{Przypadek $p=1$ i szczególna rola $L^\infty$}
\begin{twierdzenie}[Dualność $L^1$ przy $\sigma$-skończoności]
\label{thm:dual-l1-15}
Jeżeli $(X,\mathcal F,\mu)$ jest $\sigma$-skończona, to
każdy $\Lambda\in(L^1(\mu))^*$ ma jednoznaczną reprezentację
$\Lambda(f)=\int_X f(x)g(x)\,\mathrm d\mu(x)$
przez $g\in L^\infty(\mu)$, a $\|\Lambda\|=\|g\|_\infty$.
\end{twierdzenie}
\begin{proof}
Wybieramy mierzalne rozłączne $E_n$, których suma jest $X$
i $0\le\mu(E_n)<\infty$. Pomijając zbiory zerowe, kładziemy
$w(x)=\sum_{n\ge1}2^{-n}(1+\mu(E_n))^{-1}\chi_{E_n}(x)$.
Mamy $w>0$ prawie wszędzie i $\int_Xw\,\mathrm d\mu(x)\le1$.
Określamy skończoną miarę $P(A)=\int_Aw\,\mathrm d\mu(x)$.
Dla $h\in L^2(P)$ funkcja $wh$ należy do $L^1(\mu)$, gdyż
$\int_X|wh|\,\mathrm d\mu(x)=\int_X|h|\,\mathrm dP(x)
\le P(X)^{1/2}\|h\|_{L^2(P)}$. Zatem
$T(h):=\Lambda(wh)$ jest ciągłym funkcjonałem na
przestrzeni Hilberta $L^2(P)$. Z twierdzenia o reprezentacji
Riesza dla przestrzeni Hilberta istnieje $g\in L^2(P)$
takie, że $T(h)=\int_X h(x)g(x)\,\mathrm dP(x)$.
Uzasadnienie wykorzystanej reprezentacji jest elementarne:
dla $T\ne0$ wybieramy $u$ prostopadły do $\ker T$ o
$T(u)=1$, po rzutowaniu ortogonalnym na domkniętą
podprzestrzeń $\ker T$; wtedy
$T(h)=\langle h,u\rangle_2/\|u\|_2^2$ przy odpowiedniej
konwencji liniowości. Dla samodzielności przypomnijmy dowód potrzebnego
przypadku twierdzenia o rzucie. Jeśli $K\subset H$ jest
domkniętą podprzestrzenią przestrzeni Hilberta,
a $h\in H$, wybieramy $z_n\in K$ tak, że
$\|h-z_n\|^2\to d^2:=\inf_{z\in K}\|h-z\|^2$.
Tożsamość równoległoboku daje
\[
\|z_n-z_m\|^2
=2\|h-z_n\|^2+2\|h-z_m\|^2
-4\left\|h-\frac{z_n+z_m}{2}\right\|^2
\le2\|h-z_n\|^2+2\|h-z_m\|^2-4d^2.
\]
Stąd $(z_n)$ jest ciągiem Cauchy'ego i ma granicę
$z\in K$. Dla każdego $k\in K$ funkcja
$t\mapsto\|h-z-tk\|^2$ ma minimum przy $t=0$;
biorąc również $tk$ z czysto urojonym współczynnikiem,
otrzymujemy $h-z\perp K$. Tak otrzymujemy potrzebny rzut
bez odwołania do zewnętrznego twierdzenia.

Dla każdego $h\in L^2(P)\cap L^1(P)$ mamy ponadto
$|T(h)|\le\|\Lambda\|\int_X|h|\,\mathrm dP(x)$.
W szczególności, podstawiając
$h=\chi_A\overline g/|g|$ (równe zero przy $g=0$),
otrzymujemy
$\int_A|g|\,\mathrm dP(x)\le\|\Lambda\|P(A)$
dla dowolnego $A\in\mathcal F$. Stąd
$|g|\le\|\Lambda\|$ $P$-prawie wszędzie; ponieważ
$w>0$ $\mu$-prawie wszędzie, także $\mu$-prawie wszędzie.

Jeśli $f\in L^1(\mu)$, to $h=f/w\in L^1(P)$
(zera $w$ pomijamy). Funkcje
$h_n=h\chi_{\{|h|\le n\}}$ należą do
$L^2(P)$, a $wh_n\to f$ w $L^1(\mu)$, co wynika
z dominacji przez $|f|$. Zatem
$\Lambda(f)=\lim_nT(h_n)=\lim_n\int_Xh_ng\,\mathrm dP(x)
=\int_Xfg\,\mathrm d\mu(x)$.
Izometrię i jednoznaczność otrzymujemy z
lematu~\ref{lem:canonical-lpq-15} w wariancie $p=1$
i z $\sigma$-skończonością.
\end{proof}

\begin{uwaga}[Dlaczego dualna $L^\infty$ jest większa?]
Nie jest prawdą, że $(L^\infty)^*=L^1$ nawet dla
miary liczącej na $\mathbb N$. Wtedy $L^\infty=\ell^\infty$.
Na podprzestrzeni $c$ ciągów zbieżnych funkcjonał
$T((x_n))=\lim_nx_n$ jest ciągły i ma normę $1$.
Twierdzenie Hahna--Banacha przedłuża go do
$\widetilde T\in(\ell^\infty)^*$ o tej samej normie.
Jeżeli byłby postaci $\widetilde T(x)=\sum_n a_nx_n$
z $(a_n)\in\ell^1$, to dla wektorów jednostkowych $e_n$
mielibyśmy $a_n=\widetilde T(e_n)=0$ (każdy $e_n\in c_0$),
lecz $\widetilde T(\chi_X)=1$. Sprzeczność pokazuje,
że taki funkcjonał nie pochodzi z $\ell^1$.
W rozdziałach o miarach ze znakiem i Radonie--Nikodymie
poznamy dodatkowo reprezentacje przez miary, lecz nie
należy utożsamiać dualnej $L^\infty$ z miarami
\emph{przeliczalnie} addytywnymi.
\end{uwaga}

\begin{wniosek}[Dokładne maksimum w nierówności Höldera]
\label{cor:holder-sharp-15}
Niech $1<p<\infty$, $q=p/(p-1)$, a
$g\in L^q(\mu)$, $g\neq0$. Istnieje
$f\in L^p(\mu)$ o $\|f\|_p=1$ takie, że
\[
\int_X f(x)g(x)\,\mathrm d\mu(x)=\|g\|_q.
\]
Można wybrać
$f(x)=\overline{g(x)}|g(x)|^{q-2}/\|g\|_q^{q-1}$
przy $g(x)\neq0$ oraz $f(x)=0$ przy $g(x)=0$.
\end{wniosek}
\begin{proof}
Mamy $|f|^p=|g|^q/\|g\|_q^q$,
ponieważ $(q-1)p=q$, zatem $\|f\|_p=1$.
Ponadto $fg=|g|^q/\|g\|_q^{q-1}$, więc
jego całka wynosi $\|g\|_q$. Te rachunki
pokazują, że stała $1$ w nierówności Höldera
jest optymalna niezależnie od przestrzeni z miarą.
\end{proof}

\section{Refleksywność, zbieżność słaba i słaba-*}
\begin{definicja}[Refleksywność i zbieżność słaba]
\label{def:weak-lp-15}
Dla przestrzeni Banacha $E$ zanurzenie kanoniczne
$J:E\to E^{**}$ jest dane przez $(Jx)(\Lambda)=\Lambda(x)$.
Przestrzeń $E$ jest \emph{refleksywna}, jeżeli $J$ jest
surjekcją. Ciąg $x_n$ \emph{zbiega słabo} do $x$,
co zapisujemy $x_n\rightharpoonup x$, jeżeli
$\Lambda(x_n)\to\Lambda(x)$ dla każdego $\Lambda\in E^*$.
Na dualnej $E^*$ mówimy o \emph{zbieżności słabej-*}
$\Lambda_n\overset{*}{\rightharpoonup}\Lambda$,
gdy $\Lambda_n(x)\to\Lambda(x)$ dla każdego $x\in E$.
Są to dwie na ogół różne topologie na $E^*$.
\end{definicja}

\begin{twierdzenie}[Refleksywność $L^p$]
\label{thm:reflexive-lp-15}
Dla dowolnej przestrzeni z miarą $(X,\mathcal F,\mu)$
i $1<p<\infty$ przestrzeń $L^p(\mu)$ jest refleksywna.
\end{twierdzenie}
\begin{proof}
Niech $q=p/(p-1)$. Z twierdzenia~\ref{thm:dual-lp-15}
izometrycznie identyfikujemy $(L^p)^*$ z $L^q$ przez
$g\mapsto\Lambda_g$, gdzie $\Lambda_g(f)=\int_Xfg\,\mathrm d\mu(x)$.
Ponieważ $1<q<\infty$, to samo twierdzenie zastosowane
z zamienionymi rolami $p,q$ utożsamia $(L^q)^*$ z $L^p$.
Pod tymi identyfikacjami kanoniczny obraz $J(f)$
działa na $g\in L^q$ jako $J(f)(g)=\int_Xfg\,\mathrm d\mu(x)$.
Każdy element $(L^q)^*$ ma tę postać, co dowodzi
surjektywności $J$.
\end{proof}

\begin{twierdzenie}[Kryterium słabej zbieżności w $L^p$]
\label{thm:weak-lp-15}
Niech $1<p<\infty$, $q=p/(p-1)$ i $f_n,f\in L^p(\mu)$.
Wtedy $f_n\rightharpoonup f$ w $L^p$ dokładnie wtedy, gdy
\[
\int_Xf_n(x)g(x)\,\mathrm d\mu(x)
\longrightarrow\int_Xf(x)g(x)\,\mathrm d\mu(x)
\quad\text{dla każdego }g\in L^q(\mu).
\]
Jeśli $f_n\to f$ w normie $L^p$, to $f_n\rightharpoonup f$.
\end{twierdzenie}
\begin{proof}
Pierwsza równoważność jest dokładnie połączeniem
definicji~\ref{def:weak-lp-15} z
reprezentacją~\ref{thm:dual-lp-15}.
Druga wynika z $|\Lambda(f_n-f)|\le\|\Lambda\|\|f_n-f\|_p$.
\end{proof}

\begin{twierdzenie}[Dolna półciągłość normy]
\label{thm:weak-lsc-15}
W dowolnej przestrzeni Banacha $E$, jeśli $x_n\rightharpoonup x$,
to $\|x\|\le\liminf_n\|x_n\|$.
\end{twierdzenie}
\begin{proof}
Dla każdego $\Lambda\in E^*$ z $\|\Lambda\|\le1$
mamy $|\Lambda(x)|=\lim_n|\Lambda(x_n)|
\le\liminf_n\|x_n\|$.
Na mocy twierdzenia Hahna--Banacha istnieje
funkcjonał normujący wektor $x$, co daje
$\|x\|\le\liminf_n\|x_n\|$.
W przypadku $E=L^p$, $1<p<\infty$, można uniknąć
Hahna--Banacha: z lematu~\ref{lem:canonical-lpq-15}
$\|x\|_p=\sup_{\|g\|_q\le1}|\int_Xxg\,\mathrm d\mu(x)|$.
\end{proof}

\begin{przyklad}[Słaba zbieżność bez zbieżności normowej]
Na $X=(0,1)$ z miarą Lebesgue'a niech
$r_n(x)=\operatorname{sign}(\sin(2^{n+1}\pi x))$,
przy dowolnym ustaleniu wartości na zbiorze miejsc zerowych.
Funkcja $r_n$ przyjmuje naprzemiennie wartości $1$ i $-1$
na kolejnych przedziałach diadycznych długości $2^{-n-1}$.
Dla każdego $1<p<\infty$ mamy $\|r_n\|_p=1$, lecz
$r_n\rightharpoonup0$ w $L^p(0,1)$. Istotnie, całka
$r_n$ po dowolnym przedziale diadycznym ustalonego
rzędu jest zerowa dla wszystkich dostatecznie dużych $n$.
W konsekwencji $\int r_ng\to0$ dla funkcji schodkowych
na przedziałach diadycznych. Ich gęstość w $L^q(0,1)$
i Hölder pozwalają przenieść zbieżność na każde $g\in L^q$.
\end{przyklad}

\begin{twierdzenie}[Podciąg słabo zbieżny: wariant separowalny]
\label{thm:weak-subseq-15}
Niech $1<p<\infty$, a $L^q(\mu)$ będzie separowalna.
Wtedy z każdego ciągu ograniczonego $(f_n)$ w $L^p(\mu)$
można wybrać podciąg słabo zbieżny do pewnego $f\in L^p(\mu)$.
\end{twierdzenie}
\begin{proof}
Niech $(g_j)_{j\ge1}$ będzie przeliczalnym gęstym
zbiorem w $L^q$. Dla każdego $j$ ciąg liczb
$\int_Xf_n g_j\,\mathrm d\mu(x)$ jest ograniczony przez
$M\|g_j\|_q$, gdzie $M=\sup_n\|f_n\|_p$.
Kolejno wybieramy podciągi, na których te liczby
są zbieżne, i przyjmujemy podciąg diagonalny
$(f_{n_k})$. Granica definiuje funkcjonał liniowy na
liniowej powłoce $(g_j)$, ograniczony przez $M\|g\|_q$.
Przedłużamy go przez ciągłość na całe $L^q$.
Z twierdzenia~\ref{thm:dual-lp-15} zastosowanego do $q$
funkcjonał ten ma postać $g\mapsto\int_Xfg\,\mathrm d\mu(x)$
dla pewnego $f\in L^p$. Gęstość i jednolite
oszacowanie $M$ zapewniają zbieżność całek
$\int_Xf_{n_k}g\to\int_Xfg$ dla wszystkich $g\in L^q$.
Kryterium~\ref{thm:weak-lp-15} daje tezę.
\end{proof}

\begin{uwaga}[Zastosowanie i ograniczenie]
W ogólnym przypadku refleksywności przestrzeni Banacha
odpowiednikiem ostatniego twierdzenia jest twierdzenie
Eberleina--\v Smuliana (zwartość słaba i zwartość ciągowa).
Nie dowodzimy go tutaj: nie jest twierdzeniem czysto
miarowym. Dowód diagonalny powyżej jest samodzielny
i wystarcza między innymi dla $L^p(\mathbb R^d)$,
$1<p<\infty$.
\end{uwaga}

\section{Zastosowanie: minimalizacja funkcjonałów całkowych}
\begin{twierdzenie}[Minimalizacja wypukłej energii]
\label{thm:minimize-lp-15}
Niech $1<p<\infty$, $g\in L^q(\mu)$ i $\alpha>0$.
Funkcjonał $\mathcal E(f)=\alpha\|f\|_p^p/p
-\operatorname{Re}\int_X f g\,\mathrm d\mu(x)$
ma dokładnie jedno minimum w $L^p(\mu)$. Jego
punkt minimalny jest dany prawie wszędzie przez
\[
f_*(x)=\alpha^{-1/(p-1)}\,\overline{g(x)}\,|g(x)|^{q-2},
\]
przy umowie $f_*(x)=0$ dla $g(x)=0$.
\end{twierdzenie}
\begin{proof}
Istnienie minimum dla $\alpha>0$ wynika z dowodu
lematu~\ref{lem:variational-lp-15} (stały dodatni
współczynnik nie zmienia argumentu). Funkcja
$z\mapsto|z|^p$ jest ściśle wypukła na $\mathbb C$,
a więc całka jest ściśle wypukła na klasach $L^p$;
stąd jednoznaczność. Warunek pierwszego rzędu daje
$\alpha|f_*|^{p-2}\overline{f_*}=g$.
Rozwiązanie tego równania w punktach $g\ne0$
ma wskazaną postać, ponieważ $q-1=1/(p-1)$;
$|f_*|^p=\alpha^{-q}|g|^q$ jest całkowalne.
To dowodzi zarówno istnienia wskazanej funkcji,
jak i spełnienia warunku minimalności (wypukłość
$\mathcal E$ sprawia, że warunek jest wystarczający).
\end{proof}

\section{Mapa zależności}
Dowiedliśmy, że dla $1<p<\infty$ kolejno zachodzą:
nierówności punktu środkowego, jednostajna wypukłość,
istnienie minimizatora, reprezentacja dualna i
refleksywność. Przypadek $p=1$ wykorzystał zamiast tego
reprezentację Hilbertowską po zmianie miary na skończoną.
Nie zakładaliśmy rozkładu Radona--Nikodyma; w rozdziale 18
zobaczymy jego szerszą wersję dla samych miar.

% END chapters/pl/measure_theory/15-rozdzial.tex


% BEGIN chapters/pl/measure_theory/16-rozdzial.tex
% ROZDZIAŁ 16 — skończone miary ze znakiem i zespolone
\chapter[Miary ze znakiem i zespolone]{Miary ze znakiem, miary zespolone i wariacja całkowita}
\label{ch:miary-ze-znakiem}

\section{Od miar dodatnich do miar przyjmujących wartości ujemne}
Dotychczas miara $\mu$ przyjmowała wartości w $[0,\infty]$.
W zastosowaniach pojawiają się różnice miar, np.
$\nu(E)=\int_E f(x)\,\mathrm d\mu(x)$ dla rzeczywistej
funkcji $f\in L^1(\mu)$, oraz miary zespolone
powstające przez całkowanie funkcji zespolonych.
W tym rozdziale zaczynamy od \emph{skończonych wartości},
aby nie tworzyć niezdefiniowanego wyrażenia $+\infty-\infty$.
O miarach ze znakiem przyjmujących jedną z wartości
nieskończonych wspomnimy osobno na końcu rozdziału;
główne twierdzenia udowadniamy dla miar skończenie wartościowych.

\begin{definicja}[Skończona miara ze znakiem]
\label{def:signed-measure-16}
Niech $(X,\mathcal F)$ będzie przestrzenią mierzalną.
\emph{Skończoną (skończenie wartościową) miarą ze znakiem}
nazywamy funkcję $\nu:\mathcal F\to\mathbb R$, dla której
$\nu(\varnothing)=0$ oraz dla dowolnych parami rozłącznych
$(E_n)_{n\ge1}\subset\mathcal F$ zachodzi
$\nu(\bigcup_{n\ge1}E_n)=\sum_{n\ge1}\nu(E_n)$.
Szereg po prawej jest wówczas zbieżny i — jako że równość
obowiązuje przy dowolnej permutacji zbiorów — bezwzględnie
zbieżny. Nie żądamy monotoniczności $\nu$.
\end{definicja}

\begin{przyklad}[Różnica skończonych miar]
Jeżeli $\mu_1,\mu_2$ są dodatnimi miarami na $\mathcal F$
i $\mu_1(X),\mu_2(X)<\infty$, to $\nu=\mu_1-\mu_2$
jest skończoną miarą ze znakiem.
Jeżeli $f\in L^1(\mu;\mathbb R)$, to
$\nu_f(E)=\int_E f\,\mathrm d\mu(x)$ jest skończoną miarą ze znakiem.
Jej przeliczalna addytywność wynika z całkowania sumy
funkcji charakterystycznych rozłącznych zbiorów, osobno dla $f^+$ i $f^-$.
\end{przyklad}

\begin{uwaga}[Ciągłość z góry i z dołu]
Dla skończonej miary ze znakiem $\nu$, jeśli
$E_n\uparrow E$, to $\nu(E_n)\to\nu(E)$, natomiast
jeśli $E_n\downarrow E$, to $\nu(E_n)\to\nu(E)$.
Pierwszy fakt otrzymujemy, rozkładając $E$ na rozłączne
przyrosty $E_1,E_2\setminus E_1,\dots$; drugi wynika
z pierwszego dla dopełnień oraz skończonej addytywności.
W odróżnieniu od miar nieujemnych nie jest potrzebne
założenie $\nu(E_1)<\infty$, bo wszystkie wartości
$\nu$ są rzeczywiste.
\end{uwaga}

\section{Dlaczego skończona miara ze znakiem jest ograniczona?}
\begin{lemat}[Ograniczoność zakresu]
\label{lem:bounded-signed-16}
Jeśli $\nu:\mathcal F\to\mathbb R$ jest skończoną
miarą ze znakiem, to $M:=\sup_{E\in\mathcal F}|\nu(E)|<\infty$.
\end{lemat}
\begin{proof}
Przypuśćmy przeciwnie, że wartości $\nu$ są nieograniczone
na podzbiorach $X$. Powiemy, że zbiór $B\in\mathcal F$
jest \emph{nieograniczony}, jeżeli
$\sup\{|\nu(A)|:A\in\mathcal F,A\subseteq B\}=\infty$.
Mamy więc nieograniczony zbiór $B_1=X$.
Wybieramy $A\subseteq B_1$ z
$|\nu(A)|>|\nu(B_1)|+2$. Ponieważ
$\nu(B_1\setminus A)=\nu(B_1)-\nu(A)$,
oba rozłączne zbiory $A$ i $B_1\setminus A$ mają
wartość $\nu$ o module większym niż $2$.
Co najmniej jeden z nich jest nieograniczony: gdyby oba
były ograniczone, dla każdego $C\subseteq B_1$ mielibyśmy
$\nu(C)=\nu(C\cap A)+\nu(C\cap(B_1\setminus A))$,
co dawałoby ograniczoność na $B_1$.
Jako $B_2$ wybieramy część nieograniczoną, a drugą
nazywamy $D_1$. Wtedy $|\nu(D_1)|>2$.
Powtarzając konstrukcję w $B_2$, otrzymujemy kolejne
nieograniczone $B_n$ i parami rozłączne $D_n$
o $|\nu(D_n)|>2$. Jednak przeliczalna addytywność
wymaga zbieżności szeregu $\sum_n\nu(D_n)$,
a więc w szczególności $\nu(D_n)\to0$.
Otrzymaliśmy sprzeczność.
\end{proof}

\section{Rozkład Hahna}
\begin{definicja}[Zbiór dodatni i ujemny]
\label{def:positive-negative-16}
Niech $\nu$ będzie skończoną miarą ze znakiem.
Zbiór $P\in\mathcal F$ nazywamy $\nu$-\emph{dodatnim},
jeżeli $\nu(A)\ge0$ dla każdego mierzalnego $A\subseteq P$.
Zbiór $N\in\mathcal F$ nazywamy $\nu$-\emph{ujemnym},
jeżeli $\nu(A)\le0$ dla każdego mierzalnego $A\subseteq N$.
Nierówność dla samego $\nu(P)$, bez kwantyfikacji po
wszystkich podzbiorach, nie wystarcza.
\end{definicja}

\begin{twierdzenie}[Rozkład Hahna]
\label{thm:hahn-16}
Dla każdej skończonej miary ze znakiem $\nu$ na
$(X,\mathcal F)$ istnieją rozłączne mierzalne zbiory
$P,N\in\mathcal F$ takie, że $X=P\cup N$,
$P$ jest $\nu$-dodatni, a $N$ jest $\nu$-ujemny.
Ponadto jeśli $(P',N')$ jest drugim takim rozkładem,
to $P\mathbin\triangle P'$ jest $\nu$-zerowy
w silnym sensie: $\nu(A)=0$ dla każdego mierzalnego
$A\subseteq P\mathbin\triangle P'$.
\end{twierdzenie}
\begin{proof}
\emph{Istnienie.} Na mocy lematu~\ref{lem:bounded-signed-16}
liczba $\alpha:=\sup\{\nu(A):A\in\mathcal F\}$ jest
skończona i $\alpha\ge0$. Dla $n\ge1$ wybieramy
$A_n\in\mathcal F$ z $\nu(A_n)>\alpha-2^{-n}$.
Dla każdego mierzalnego $B\subseteq A_n$ zachodzi
$\nu(B)\ge-2^{-n}$: w przeciwnym przypadku
$\nu(A_n\setminus B)=\nu(A_n)-\nu(B)>\alpha$,
sprzecznie z definicją $\alpha$.

Niech $F_m=\bigcup_{n\ge m}A_n$. Dla $k\ge m$,
przy dokładaniu $A_{k+1}$ do $\bigcup_{n=m}^kA_n$
dokładamy zbiór mierzalny zawarty w $A_{k+1}$,
a więc jego wartość $\nu$ jest nie mniejsza
niż $-2^{-(k+1)}$. Zatem
\[
\nu\left(\bigcup_{n=m}^kA_n\right)
\ge\nu(A_m)-\sum_{n=m+1}^k2^{-n}
>\alpha-2^{1-m}.
\]
Dla $k\to\infty$ ciągłość $\nu$ z dołu daje
$\nu(F_m)\ge\alpha-2^{1-m}$; z definicji
$\alpha$ mamy również $\nu(F_m)\le\alpha$.
Zbiory $F_m$ maleją, więc dla
$P:=\bigcap_{m\ge1}F_m$ ciągłość $\nu$ z góry daje
$\nu(P)=\lim_m\nu(F_m)=\alpha$.

Jeśli $B\subseteq P$ jest mierzalny, to
$\nu(P\setminus B)\le\alpha$ i
$\nu(B)=\nu(P)-\nu(P\setminus B)\ge0$.
Zatem $P$ jest dodatni. Jeżeli $B\subseteq N:=X\setminus P$,
to $\nu(P\cup B)=\alpha+\nu(B)\le\alpha$,
więc $\nu(B)\le0$. To dowodzi, że $N$ jest ujemny.

\emph{Jednoznaczność.} Zbiór $P\cap N'$ jest jednocześnie
podzbiorem zbioru dodatniego i ujemnego. Każdy jego
mierzalny podzbiór ma zatem miarę ze znakiem równą zero.
To samo dotyczy $N\cap P'$. Ich suma
$(P\cap N')\cup(N\cap P')=P\mathbin\triangle P'$ ma
więc wskazaną własność.
\end{proof}

\begin{przyklad}[Rozkład Hahna nie zależy tylko od znaku całkowitej masy]
Na $X=[0,1]$ z miarą Lebesgue'a połóżmy
$\nu(E)=\int_E(2x-1)\,\mathrm d\lambda_1(x)$.
Wtedy $P=[1/2,1]$ oraz $N=[0,1/2)$ tworzą rozkład
Hahna, chociaż $\nu(X)=0$. Dowolny fragment
zbioru $\{1/2\}$ można przenieść z $P$ do $N$;
nie zmienia to wartości $\nu$ na żadnym podzbiorze.
\end{przyklad}

\begin{center}
\begin{tikzpicture}[scale=0.92]
\draw[rounded corners=8pt,thick] (0,0) rectangle (9,2.55);
\fill[blue!12,rounded corners=8pt] (0.04,0.04) rectangle (4.5,2.51);
\fill[red!11,rounded corners=8pt] (4.5,0.04) rectangle (8.96,2.51);
\draw[thick] (4.5,0.05)--(4.5,2.5);
\node at (2.25,1.85) {$P$ -- zbiór dodatni};
\node at (6.75,1.85) {$N$ -- zbiór ujemny};
\node at (2.25,0.85) {$\nu(A)\ge0$ dla $A\subseteq P$};
\node at (6.75,0.85) {$\nu(B)\le0$ dla $B\subseteq N$};
\node[above] at (4.5,2.56) {$X=P\sqcup N$};
\end{tikzpicture}
\end{center}
\noindent Rysunek ilustruje kwantyfikację po \emph{wszystkich}
zbiorach mierzalnych zawartych w częściach $P$ i $N$;
nie wystarcza sam znak $\nu(P)$ oraz $\nu(N)$.

\section{Rozkład Jordana i wariacja całkowita}
\begin{definicja}[Wzajemna osobliwość]
\label{def:singular-16}
Dwie nieujemne miary $\mu_1,\mu_2$ na
$(X,\mathcal F)$ są \emph{wzajemnie osobliwe}, co zapisujemy
$\mu_1\perp\mu_2$, jeżeli istnieje $P\in\mathcal F$
taki, że $\mu_1(X\setminus P)=0$ oraz $\mu_2(P)=0$.
\end{definicja}

\begin{twierdzenie}[Rozkład Jordana]
\label{thm:jordan-16}
Niech $\nu$ będzie skończoną miarą ze znakiem.
Istnieje dokładnie jedna para skończonych miar nieujemnych
$(\nu^+,\nu^-)$ taka, że
$\nu=\nu^+-\nu^-$ i $\nu^+\perp\nu^-$.
Jeżeli $(P,N)$ jest rozkładem Hahna, to
\[
\nu^+(E)=\nu(E\cap P),\qquad
\nu^-(E)=-\nu(E\cap N)\quad(E\in\mathcal F).
\]
Ponadto $\nu^+(X)=\sup_E\nu(E)$ i
$\nu^-(X)=-\inf_E\nu(E)$.
\end{twierdzenie}
\begin{proof}
Dla rozkładu Hahna $P,N$ podane funkcje
są nieujemne i przeliczalnie addytywne, bo ich
argumenty $E\cap P$ i $E\cap N$ zachowują rozłączność.
Mamy $\nu(E)=\nu(E\cap P)+\nu(E\cap N)
=\nu^+(E)-\nu^-(E)$ oraz
$\nu^+(N)=\nu^-(P)=0$. Obie miary są skończone,
bo $\nu$ przyjmuje wartości rzeczywiste.

Jeżeli $\nu=\alpha-\beta$, gdzie $\alpha,\beta$ są
nieujemnymi miarami skończonymi i $\alpha\perp\beta$,
to istnieje $P'$ z $\alpha(X\setminus P')=0$
i $\beta(P')=0$. Każdy podzbiór $P'$ ma nieujemną
wartość $\nu$, a każdy podzbiór $X\setminus P'$ ujemną.
Jest to więc inny rozkład Hahna. Z jednoznaczności
w twierdzeniu~\ref{thm:hahn-16} zbiór
$Z=P\mathbin\triangle P'$ ma własność $\nu(C)=0$
dla każdego mierzalnego $C\subseteq Z$.
Na $P\cap N'$ mamy $\alpha=0$ i $\nu=-\beta$;
zerowość $\nu$ na podzbiorach $P\cap N'$ wymusza
$\beta(P\cap N')=0$. Na $N\cap P'$ mamy analogicznie
$\beta=0$, $\nu=\alpha$, a więc
$\alpha(N\cap P')=0$. Dla dowolnego $E\in\mathcal F$
możemy teraz policzyć
$\alpha(E)=\alpha(E\cap P')=\alpha(E\cap P\cap P')
=\nu(E\cap P\cap P')=\nu(E\cap P)=\nu^+(E)$.
Tak samo $\beta(E)=-\nu(E\cap N)=\nu^-(E)$.
Stąd jednoznaczność pary Jordana.
Wreszcie dla każdego $E$ mamy
$\nu(E)=\nu^+(E)-\nu^-(E)\le\nu^+(X)$,
a równość jest osiągana na $P$. Analogicznie
$\nu(E)\ge-\nu^-(X)$, z równością na $N$.
\end{proof}

\begin{definicja}[Wariacja całkowita]
\label{def:variation-signed-16}
Dla skończonej miary ze znakiem $\nu$ kładziemy
$|\nu|:=\nu^++\nu^-$. Jest to nieujemna miara
skończona, zwana \emph{wariacją całkowitą} $\nu$.
Liczba $\|\nu\|_{\mathrm{TV}}:=|\nu|(X)$ jest
\emph{normą w wariacji całkowitej}.
\end{definicja}

\begin{twierdzenie}[Opis wariacji przez podziały]
\label{thm:partition-variation-16}
Dla skończonej miary ze znakiem $\nu$ i $E\in\mathcal F$
zachodzi
\begin{equation}\label{eq:variation-partitions-16}
|\nu|(E)=\sup\left\{\sum_{j=1}^m|\nu(E_j)|:
E=\bigsqcup_{j=1}^mE_j,\ E_j\in\mathcal F,\ m\ge1\right\}.
\end{equation}
Ponadto $|\nu(E)|\le|\nu|(E)$ dla każdego $E$.
\end{twierdzenie}
\begin{proof}
Dla dowolnego podziału $E=\bigsqcup_jE_j$ mamy
$|\nu(E_j)|=|\nu^+(E_j)-\nu^-(E_j)|
\le\nu^+(E_j)+\nu^-(E_j)$.
Po zsumowaniu otrzymujemy oszacowanie supremum
przez $|\nu|(E)$. Podział $E=(E\cap P)\sqcup(E\cap N)$
daje sumę $\nu(E\cap P)-\nu(E\cap N)=|\nu|(E)$,
więc supremum jest osiągnięte. Ostatnia nierówność
wynika z dopuszczenia podziału jednoelementowego.
\end{proof}

\begin{wniosek}[Najmniejsza miara dominująca]
\label{cor:minimal-variation-16}
Jeśli $\rho$ jest miarą nieujemną na $(X,\mathcal F)$
i $|\nu(E)|\le\rho(E)$ dla wszystkich $E$, to
$|\nu|(E)\le\rho(E)$ dla każdego $E\in\mathcal F$.
\end{wniosek}
\begin{proof}
Dla każdego skończonego podziału
$E=\bigsqcup_jE_j$ suma $\sum_j|\nu(E_j)|$
nie przekracza $\sum_j\rho(E_j)=\rho(E)$.
Bierzemy supremum w \eqref{eq:variation-partitions-16}.
\end{proof}

\section{Miary zespolone}
\begin{definicja}[Miara zespolona]
\label{def:complex-measure-16}
Na przestrzeni mierzalnej $(X,\mathcal F)$
\emph{miarą zespoloną} nazywamy przeliczalnie addytywną
funkcję $\eta:\mathcal F\to\mathbb C$ z $\eta(\varnothing)=0$.
Nie dopuszczamy wartości $\infty$; szeregi na rozłącznych
rodzinach zbiorów są zatem bezwarunkowo zbieżne.
Jej część rzeczywista i urojona są skończonymi
miarami ze znakiem.
\end{definicja}

\begin{twierdzenie}[Wariacja miary zespolonej]
\label{thm:complex-variation-16}
Dla miary zespolonej $\eta$ definiujemy
\[
|\eta|(E)=\sup\left\{\sum_{j=1}^m|\eta(E_j)|:
E=\bigsqcup_{j=1}^mE_j,\ E_j\in\mathcal F\right\}.
\]
Jest to skończona miara nieujemna. Ponadto
$|\eta(E)|\le|\eta|(E)$ i
$|\eta|(E)\le|\operatorname{Re}\eta|(E)+
|\operatorname{Im}\eta|(E)$, gdzie po prawej występują
wariacje skończonych miar ze znakiem.
\end{twierdzenie}
\begin{proof}
Nierówność $|\eta(A)|\le|\operatorname{Re}\eta(A)|+
|\operatorname{Im}\eta(A)|$ i twierdzenie
\ref{thm:partition-variation-16} dają skończoność
$|\eta|(X)$ oraz podane oszacowanie dla każdego $E$.
Jeżeli $E,F$ są rozłączne, każdy ich skończony podział
można uszczegółowić, przecinając części z $E$ i $F$;
z nierówności trójkąta dostajemy
$|\eta|(E\cup F)\le|\eta|(E)+|\eta|(F)$.
W drugą stronę sklejamy dowolne podziały $E$ i $F$;
stąd wynika równość.

Niech teraz $E_n$ będą parami rozłączne, a
$E=\bigcup_nE_n$. Dla dowolnego skończonego
podziału $E=\bigsqcup_{j=1}^m A_j$ mamy
\[
\sum_{j=1}^m|\eta(A_j)|
=\sum_j\left|\sum_n\eta(A_j\cap E_n)\right|
\le\sum_n\sum_j|\eta(A_j\cap E_n)|
\le\sum_n|\eta|(E_n).
\]
Biorąc supremum po podziałach $E$, otrzymujemy
$|\eta|(E)\le\sum_n|\eta|(E_n)$. Z drugiej strony
skończona addytywność i monotoniczność $|\eta|$
dają $|\eta|(E)\ge\sum_{n=1}^N|\eta|(E_n)$
dla każdego $N$. Przechodząc do granicy, uzyskujemy
przeliczalną addytywność.
\end{proof}

\begin{przyklad}[Miara atomowa na $\mathbb N$]
Dla $(a_n)\in\ell^1(\mathbb N;\mathbb C)$ wzór
$\eta(A)=\sum_{n\in A}a_n$ definiuje miarę zespoloną
na $\mathcal P(\mathbb N)$. Jej wariacja jest równa
$|\eta|(A)=\sum_{n\in A}|a_n|$, a
$\|\eta\|_{\mathrm{TV}}=\sum_n|a_n|$.
Odwrotnie, każda miara zespolona na zbiorze
przeliczalnym ma tę postać: wystarczy przyjąć
$a_n=\eta(\{n\})$ i zastosować przeliczalną
addytywność oraz skończoność wariacji.
\end{przyklad}

\section{Miary o gęstości i warunek bezwzględnej całkowalności}
\begin{twierdzenie}[Wariacja miary określonej całką]
\label{thm:density-variation-16}
Niech $(X,\mathcal F,\mu)$ będzie dowolną przestrzenią
z miarą, a $f\in L^1(\mu;\mathbb C)$. Wzór
$\eta_f(E)=\int_E f(x)\,\mathrm d\mu(x)$
definiuje miarę zespoloną. Jej wariacja jest dana przez
\begin{equation}\label{eq:density-var-16}
|\eta_f|(E)=\int_E|f(x)|\,\mathrm d\mu(x)
\quad(E\in\mathcal F).
\end{equation}
Gdy $f$ jest rzeczywista, rozkład Jordana ma postać
$\eta_f^+(E)=\int_E f^+\,\mathrm d\mu(x)$ oraz
$\eta_f^-(E)=\int_E f^-\,\mathrm d\mu(x)$.
\end{twierdzenie}
\begin{proof}
Przeliczalna addytywność $\eta_f$ wynika z liniowości
całki i twierdzenia o zbieżności zmajoryzowanej dla
$\sum_{n=1}^N\chi_{E_n}f$ przy rozłącznych $E_n$.
Dla każdego skończonego podziału $E=\bigsqcup_jE_j$
\[
\sum_j\left|\int_{E_j}f\,\mathrm d\mu(x)\right|
\le\sum_j\int_{E_j}|f|\,\mathrm d\mu(x)
=\int_E|f|\,\mathrm d\mu(x).
\]
Daje to nierówność $|\eta_f|(E)\le\int_E|f|\,\mathrm d\mu(x)$.
Dla odwrotnej niech $u=\overline f/|f|$ przy $f\ne0$
i $u=0$ przy $f=0$. Wtedy $|u|\le1$ i $uf=|f|$.
Aproksymujemy $u\chi_E$ punktowo funkcjami prostymi
$u_k=\sum_jc_{k,j}\chi_{A_{k,j}}$, tak aby
$|c_{k,j}|\le1$, wszystkie $A_{k,j}\subseteq E$ i zbiory
$A_{k,j}$ były parami rozłączne.
Istnienie takich aproksymacji wynika z podziału dysku
jednostkowego na skończenie wiele zbiorów o malejących
średnicach. Mamy
\[
\left|\int_Xu_kf\,\mathrm d\mu(x)\right|
=\left|\sum_jc_{k,j}\eta_f(A_{k,j})\right|
\le\sum_j|\eta_f(A_{k,j})|\le|\eta_f|(E).
\]
Z twierdzenia o zbieżności zmajoryzowanej
$\int_Xu_kf\,\mathrm d\mu(x)\to\int_E|f|\,\mathrm d\mu(x)$.
To daje równość. Dla rzeczywistego $f$ zbiory
$\{f\ge0\}$ i $\{f<0\}$ tworzą rozkład Hahna,
a więc wzory na części Jordana wynikają z
twierdzenia~\ref{thm:jordan-16}.
\end{proof}

\begin{uwaga}[Odróżnienie warunku $\eta\ll\mu$ od istnienia gęstości]
Z wzoru \eqref{eq:density-var-16} wynika, że
$\mu(E)=0$ implikuje $\eta_f(E)=0$. W rozdziale 17
nazwiemy tę relację \emph{bezwzględną ciągłością miar}.
Nie każda miara, nawet skończona, ma gęstość względem
zadanej miary $\mu$: na przykład miara Diraca $\delta_0$
nie ma gęstości względem miary Lebesgue'a na $\mathbb R$.
Dokładne kryterium istnienia gęstości poda twierdzenie
Radona--Nikodyma w rozdziale 18.
\end{uwaga}

\section{Przestrzeń miar skończonych jako przestrzeń Banacha}
\begin{definicja}[Przestrzeń $\mathcal M(X,\mathcal F)$]
\label{def:complex-measures-space-16}
Niech $\mathcal M(X,\mathcal F)$ oznacza zespoloną
przestrzeń wektorową wszystkich miar zespolonych
na $(X,\mathcal F)$, z normą
$\|\eta\|_{\mathrm{TV}}=|\eta|(X)$.
\end{definicja}

\begin{twierdzenie}[Zupełność przestrzeni miar]
\label{thm:complete-measures-16}
Dla każdej przestrzeni mierzalnej $(X,\mathcal F)$
przestrzeń $(\mathcal M(X,\mathcal F),\|\cdot\|_{\mathrm{TV}})$
jest przestrzenią Banacha.
\end{twierdzenie}
\begin{proof}
Wariacja jest nieujemna i jednorodna. Nierówność
$\|\eta+\xi\|_{\mathrm{TV}}\le\|\eta\|_{\mathrm{TV}}+
\|\xi\|_{\mathrm{TV}}$ wynika z nierówności trójkąta
zastosowanej do sum po dowolnym podziale $X$.
Jeśli $\|\eta\|_{\mathrm{TV}}=0$, to
$|\eta(E)|\le|\eta|(E)\le|\eta|(X)=0$ dla wszystkich $E$.
Jest to więc norma.

Niech $(\eta_n)$ będzie ciągiem Cauchy'ego w tej normie.
Dla każdego $E\in\mathcal F$ z nierówności
$|\eta_n(E)-\eta_m(E)|\le\|\eta_n-\eta_m\|_{\mathrm{TV}}$
wynika istnienie granicy $\eta(E)=\lim_n\eta_n(E)$.
Zachowana jest skończona addytywność oraz
$\eta(\varnothing)=0$. Wykażemy przeliczalną addytywność.
Niech $E_j$ będą parami rozłączne i
$R_K=\bigcup_{j>K}E_j$. Dla dowolnego ustalonego $n$
\[
|\eta(R_K)|\le|\eta(R_K)-\eta_n(R_K)|+|\eta_n|(R_K).
\]
Pierwszy składnik jest nie większy niż
$\limsup_{m\to\infty}\|\eta_m-\eta_n\|_{\mathrm{TV}}$,
a drugi dąży do zera przy $K\to\infty$ z ciągłości
z góry skończonej miary $|\eta_n|$. Dobieramy najpierw
$n$ tak duże, aby pierwszy składnik był dowolnie mały,
a następnie $K$. Wtedy $\eta(R_K)\to0$,
co wraz ze skończoną addytywnością dowodzi
$\eta(\bigcup_jE_j)=\sum_j\eta(E_j)$.

Niech $\varepsilon>0$ i $n$ będzie tak duże, że
$\|\eta_m-\eta_n\|_{\mathrm{TV}}<\varepsilon$ dla $m\ge n$.
Dla każdego skończonego podziału $X=\bigsqcup_{j=1}^rA_j$
\[
\sum_j|\eta(A_j)-\eta_n(A_j)|
=\lim_{m\to\infty}\sum_j|\eta_m(A_j)-\eta_n(A_j)|
\le\varepsilon.
\]
Biorąc supremum po podziałach i korzystając z definicji
wariacji, dostajemy $\|\eta-\eta_n\|_{\mathrm{TV}}
\le\varepsilon$. Oznacza to zbieżność w normie.
\end{proof}

\begin{twierdzenie}[Oszacowanie całki względem miary zespolonej]
\label{thm:bounded-integral-complex-16}
Niech $\eta$ będzie miarą zespoloną na $(X,\mathcal F)$,
a $f:X\to\mathbb C$ funkcją ograniczoną i mierzalną.
Całka $\int_Xf\,\mathrm d\eta$, zdefiniowana przez
cztery części Jordana miar
$\operatorname{Re}\eta$ i $\operatorname{Im}\eta$, jest dobrze
określona i
$|\int_Xf\,\mathrm d\eta|\le\|f\|_{L^\infty(|\eta|)}\,|\eta|(X)$,
gdzie kres istotny liczony jest względem miary $|\eta|$.
Dokładniej, dla każdego $E\in\mathcal F$ zachodzi
$|\int_Ef\,\mathrm d\eta|\le\int_E|f|\,\mathrm d|\eta|$.
\end{twierdzenie}
\begin{proof}
Dla funkcji prostej $f=\sum_jc_j\chi_{A_j}$
o parami rozłącznych $A_j$ mamy
$|\int_Ef\,\mathrm d\eta|
=|\sum_jc_j\eta(E\cap A_j)|
\le\sum_j|c_j||\eta|(E\cap A_j)
=\int_E|f|\,\mathrm d|\eta|$.
Każdą funkcję ograniczoną mierzalną aproksymujemy
jednostajnie funkcjami prostymi. Całka po stronie
lewej jest ciągła przy tej aproksymacji, gdyż
$\eta$ jest kombinacją czterech skończonych miar
nieujemnych; również całka po prawej jest ciągła.
Przechodzimy do granicy. Ostatnie oszacowanie wynika z nierówności
$|f|\le\|f\|_{L^\infty(|\eta|)}$ prawie wszędzie
względem $|\eta|$.
\end{proof}

\begin{twierdzenie}[Dualny opis normy w wariacji całkowitej]
\label{thm:tv-dual-tests-16}
Niech $(X,\mathcal F)$ będzie przestrzenią mierzalną,
a $\eta$ skończoną miarą zespoloną na $\mathcal F$.
Wówczas
\begin{equation}\label{eq:tv-dual-tests-16}
 |\eta|(X)=\sup\left\{\left|\int_X \varphi(x)\,\mathrm d\eta(x)\right|:
 \varphi:X\to\mathbb C\ \text{mierzalna},\quad |\varphi|\le1\right\}.
\end{equation}
W supremum wystarczy dopuścić funkcje proste o module
nieprzekraczającym $1$.
\end{twierdzenie}
\begin{proof}
Z twierdzenia~\ref{thm:bounded-integral-complex-16}
każda dopuszczalna funkcja spełnia
$|\int_X\varphi(x)\,\mathrm d\eta(x)|\le|\eta|(X)$.
Dla dowolnego $\varepsilon>0$ definicja wariacji daje
skończony podział $X=\bigsqcup_{j=1}^mE_j$ z
$\sum_{j=1}^m|\eta(E_j)|>|\eta|(X)-\varepsilon$.
Ustalmy
$c_j=\overline{\eta(E_j)}/|\eta(E_j)|$, jeżeli
$\eta(E_j)\neq0$, i $c_j=0$ w przeciwnym przypadku.
Wtedy $\varphi=\sum_{j=1}^mc_j\chi_{E_j}$ jest
mierzalną funkcją prostą o module co najwyżej $1$, oraz
\[
 \int_X\varphi(x)\,\mathrm d\eta(x)
 =\sum_{j=1}^mc_j\eta(E_j)
 =\sum_{j=1}^m|\eta(E_j)|>|\eta|(X)-\varepsilon.
\]
Przechodząc z $\varepsilon\downarrow0$, otrzymujemy
odwrotną nierówność i równość supremów.
\end{proof}

\begin{uwaga}[Dlaczego jest to norma silniejsza od zbieżności na zbiorach?]
Równość~\eqref{eq:tv-dual-tests-16} pokazuje, że zbieżność
w wariacji całkowitej kontroluje całki \emph{jednostajnie}
po wszystkich mierzalnych funkcjach testowych o module
co najwyżej $1$. Sama zbieżność $\eta_n(E)\to0$ dla
każdego ustalonego zbioru $E$ nie daje takiej kontroli.
\end{uwaga}

\section{Przykłady i kontrprzykłady}
\begin{przyklad}[Miara całkowita równa zero, wariacja dodatnia]
Na $X=\{a,b\}$ z $\mathcal F=\mathcal P(X)$ połóżmy
$\nu(\{a\})=1$, $\nu(\{b\})=-1$. Wtedy
$\nu(X)=0$, ale $|\nu|(X)=2$.
Masa $\nu(X)$ nie jest więc normą miary ze znakiem.
\end{przyklad}

\begin{przyklad}[Oscylujące gęstości]
Na $X=(0,2\pi)$ z miarą Lebesgue'a definiujemy
$\nu_n(E)=\int_E\sin(nx)\,\mathrm d\lambda_1(x)$.
Wtedy $\nu_n(X)=0$ dla każdego $n$, natomiast
$\|\nu_n\|_{\mathrm{TV}}=
\int_{(0,2\pi)}|\sin(nx)|\,\mathrm d\lambda_1(x)=4$.
Zatem nawet zanikanie całkowitej masy wszystkich miar
nie kontroluje wariacji całkowitej. Poza tym dla każdej
funkcji ciągłej $f$ na $[0,2\pi]$ zachodzi
$\int f\,\mathrm d\nu(x)_n\to0$ (można to sprawdzić
najpierw dla funkcji schodkowych, a potem przez
aproksymację jednostajną). Nie zachodzi jednak
zbieżność $\nu_n\to0$ w normie wariacji całkowitej.
\end{przyklad}

\begin{uwaga}[Miary o wartościach nieskończonych]
W części literatury ``miarą ze znakiem'' nazywa się
również przeliczalnie addytywną funkcję
$\nu:\mathcal F\to[-\infty,\infty]$, która nie przyjmuje
jednocześnie wartości $+\infty$ i $-\infty$.
Dla takiej miary prawdziwe są odpowiednie wersje
rozkładów Hahna i Jordana, ale \emph{co najwyżej jedna}
z części $\nu^+$, $\nu^-$ może mieć nieskończoną masę.
Wtedy $\nu=\nu^+-\nu^-$ jest dobrze określone,
natomiast $|\nu|(X)$ nie musi być skończone.
Pełne dowody przedstawione w tym rozdziale dotyczą
skończenie wartościowych miar ze znakiem; nie wolno
bez sprawdzenia przenosić kroków opartych na
ciągłości z góry i finitości $\alpha$ na przypadek
nieskończony.
\end{uwaga}

% END chapters/pl/measure_theory/16-rozdzial.tex


% BEGIN chapters/pl/measure_theory/17-rozdzial.tex
% ROZDZIAŁ 17 — rozkład Lebesgue'a dowodzony niezależnie od RN
\chapter[Bezwzględna ciągłość i rozkład Lebesgue'a]{Bezwzględna ciągłość, osobliwość i rozkład Lebesgue'a miar}
\label{ch:ac-singular-17}

\section{Porównywanie miar: co znaczy, że jedna miara nie widzi więcej niż druga?}
Niech $(X,\mathcal F)$ będzie przestrzenią mierzalną. Dwie miary mogą
przypisywać tym samym zbiorom różne wartości, ale istotne jest również,
czy mają te same zbiory miary zero. Zagadnienie to prowadzi do dwóch
przeciwstawnych pojęć: bezwzględnej ciągłości i wzajemnej osobliwości.
Na tym etapie nie będziemy zakładać istnienia gęstości jednej miary
względem drugiej: właśnie możliwość zbudowania takich gęstości
stanowi treść następnego rozdziału.

\begin{definicja}[Bezwzględna ciągłość miary]
\label{def:ac-measure-17}
Niech $\mu$ i $\nu$ będą miarami nieujemnymi na $(X,\mathcal F)$.
Mówimy, że $\nu$ jest \emph{bezwzględnie ciągła względem $\mu$},
i piszemy $\nu\ll\mu$, jeżeli dla każdego $E\in\mathcal F$
z warunku $\mu(E)=0$ wynika $\nu(E)=0$.
Dla skończonej miary zespolonej $\eta$ przyjmujemy tę samą definicję,
równoważnie $|\eta|\ll\mu$, gdzie $|\eta|$ jest wariacją całkowitą.
\end{definicja}

\begin{uwaga}[Dlaczego warunki dla miary zespolonej są równoważne?]
Jeżeli $\mu(E)=0$, to $\mu(A)=0$ dla każdego mierzalnego $A\subseteq E$.
Stąd $\eta(A)=0$, a z definicji wariacji jako supremum sum po podziałach
wynika $|\eta|(E)=0$. Odwrotna implikacja wynika z
$|\eta(E)|\le|\eta|(E)$. Sam warunek $\eta(E)=0$ dla jednego,
ustalonego $E$ nie oznacza jednak $|\eta|(E)=0$.
\end{uwaga}

\begin{przyklad}[Miara określona gęstością]
\label{ex:ac-density-17}
Dla dowolnej miary $\mu$ i mierzalnej funkcji $f:X\to[0,\infty]$
miara $\nu_f(E)=\int_E f(x)\,\mathrm d\mu(x)$ jest bezwzględnie
ciągła względem $\mu$, nawet jeśli $\nu_f(X)=\infty$.
W przypadku zespolonego $f\in L^1(\mu)$ analogiczna własność
wynika z twierdzenia~\ref{thm:density-variation-16}.
Nie wyciągamy jeszcze wniosku odwrotnego: $\nu\ll\mu$
samo w sobie nie jest \emph{definicją} istnienia gęstości.
\end{przyklad}

\begin{przyklad}[Miary Diraca a miara Lebesgue'a]
Na $(\mathbb R,\mathcal L_1)$ oznaczmy przez $\delta_a$ miarę Diraca
w punkcie $a\in\mathbb R$. Wówczas $\delta_a\not\ll\lambda_1$,
ponieważ $\lambda_1(\{a\})=0$, lecz $\delta_a(\{a\})=1$.
Z drugiej strony $\lambda_1\not\ll\delta_a$, gdyż
$\delta_a(\mathbb R\setminus\{a\})=0$ i
$\lambda_1(\mathbb R\setminus\{a\})=\infty$.
\end{przyklad}

\begin{definicja}[Wzajemna równoważność miar]
\label{def:equivalent-measures-17}
Dwie nieujemne miary $\mu,\nu$ na $(X,\mathcal F)$ nazywamy
\emph{równoważnymi}, co zapisujemy $\mu\sim\nu$,
jeśli $\mu\ll\nu$ oraz $\nu\ll\mu$.
Znaczy to dokładnie, że mają te same mierzalne zbiory miary zero.
\end{definicja}

\begin{przyklad}[Równoważne miary nie muszą być równe]
Na $X=(0,1)$ miary $\mu=\lambda_1|_X$ i
$\nu(E)=\int_E x\,\mathrm d\lambda_1(x)$ są równoważne,
ponieważ $x>0$ na $X$. Nie są równe: $\mu(X)=1$,
zaś $\nu(X)=1/2$. Zmiana wagi może zachować wszystkie
zbiory zerowe, lecz zmienić wartości miar.
\end{przyklad}

\section{Wzajemna osobliwość i przykłady}
Pojęcie wzajemnej osobliwości wprowadziliśmy już w
definicji~\ref{def:singular-16}. Powtórzymy je w postaci dogodnej
do badania rozkładów.

Przypominamy definicję~\ref{def:singular-16}:
dwie nieujemne miary $\nu,\mu$ na $(X,\mathcal F)$ są
\emph{wzajemnie osobliwe}, $\nu\perp\mu$, gdy istnieje
$S\in\mathcal F$ takie, że $\mu(S)=0$ i
$\nu(X\setminus S)=0$. Dla miary zespolonej $\eta$
piszemy $\eta\perp\mu$, gdy $|\eta|\perp\mu$.
Jest to rozszerzenie tej samej relacji na miary zespolone,
a nie nowa definicja dla miar nieujemnych.

Zbiór $S$ w definicji jest \emph{zbiorem skupiającym}
miarę $\nu$ (można też mówić o mierzalnym nośniku). Nie należy mylić go z
\emph{nośnikiem topologicznym}: ten drugi wymaga topologii
i jest definiowany za pomocą otoczeń punktów.

\begin{lemat}[Elementarne własności osobliwości]
\label{lem:singularity-17}
Niech $\mu,\nu,\rho$ będą miarami nieujemnymi na $(X,\mathcal F)$.
Wtedy:
\begin{enumerate}[label=\textup{(\roman*)}]
\item $\nu\perp\mu$ wtedy i tylko wtedy, gdy $\mu\perp\nu$;
\item jeśli $\nu\ll\mu$ i $\nu\perp\mu$, to $\nu=0$;
\item jeśli $\nu\perp\mu$ oraz $\rho\perp\mu$, to
$\nu+\rho\perp\mu$;
\item dla dowolnego $S\in\mathcal F$ miary
$\mu|_S(E):=\mu(E\cap S)$ i
$\mu|_{X\setminus S}(E):=\mu(E\setminus S)$ są wzajemnie osobliwe.
\end{enumerate}
\end{lemat}
\begin{proof}
(i) Jeśli $\mu(S)=0$ i $\nu(X\setminus S)=0$, to
$X\setminus S$ świadczy o relacji odwrotnej.
(ii) Z $\nu\ll\mu$ wynika $\nu(S)=0$, a z osobliwości
$\nu(X\setminus S)=0$, więc $\nu(X)=0$ i $\nu=0$.
(iii) Jeśli $\nu$ jest skupiona na $S$, a $\rho$ na $T$,
to $\mu(S\cup T)=0$ i $(\nu+\rho)(X\setminus(S\cup T))=0$.
(iv) Pierwsza miara jest skupiona na $S$, a druga na $X\setminus S$.
\end{proof}

\begin{przyklad}[Miara dyskretna i miara ciągła]
Niech $(a_n)_{n\ge1}$ będzie ciągiem punktów w $\mathbb R$, a
$(c_n)_{n\ge1}$ ciągiem liczb nieujemnych o
$\sum_nc_n<\infty$. Miara $\nu=\sum_nc_n\delta_{a_n}$
jest osobliwa względem $\lambda_1$: wystarczy
$S=\{a_n:n\ge1\}$, zbiór przeliczalny, więc zerowy
względem miary Lebesgue'a. Nie ma znaczenia, czy $S$ jest
gęsty w $\mathbb R$.
\end{przyklad}

\section{Kryterium epsilon--delta i rola skończoności}
Warunek $\nu\ll\mu$ mówi jedynie o zbiorach dokładnie zerowej
miary. Gdy $\nu$ jest skończona, można go wzmocnić do
\emph{jednostajnej} kontroli małych zbiorów.

\begin{twierdzenie}[Kryterium epsilon--delta]
\label{thm:eps-delta-ac-17}
Niech $\mu$ będzie dowolną miarą nieujemną na $(X,\mathcal F)$,
a $\nu$ miarą nieujemną o $\nu(X)<\infty$.
Następujące warunki są równoważne:
\begin{enumerate}[label=\textup{(\roman*)}]
\item $\nu\ll\mu$;
\item dla każdego $\varepsilon>0$ istnieje $\delta>0$ taka,
że dla każdego $E\in\mathcal F$ zachodzi implikacja
$\mu(E)<\delta\implies\nu(E)<\varepsilon$.
\end{enumerate}
Dla skończonej miary zespolonej $\eta$ równoważność zachodzi po
zastąpieniu $\nu(E)$ przez $|\eta|(E)$.
\end{twierdzenie}
\begin{proof}
(ii)$\Rightarrow$(i). Jeśli $\mu(E)=0$, to
$\mu(E)<\delta$ dla każdej liczby $\delta>0$.
Stąd $\nu(E)<\varepsilon$ dla każdego $\varepsilon>0$,
a więc $\nu(E)=0$.

(i)$\Rightarrow$(ii). Gdyby (ii) było fałszywe, istniałoby
$\varepsilon_0>0$ takie, że dla każdego $n\ge1$
można wybrać $E_n\in\mathcal F$ z
$\mu(E_n)<2^{-n}$ oraz $\nu(E_n)\ge\varepsilon_0$.
Niech $F_m=\bigcup_{n\ge m}E_n$ i
$F=\bigcap_{m\ge1}F_m=\limsup_nE_n$.
Z przeliczalnej podaddytywności otrzymujemy
$\mu(F_m)\le\sum_{n\ge m}2^{-n}$,
zatem $\mu(F)=0$. Każdy $F_m$ zawiera $E_m$, więc
$\nu(F_m)\ge\varepsilon_0$. Ponieważ $F_m\downarrow F$
i $\nu(F_1)\le\nu(X)<\infty$, ciągłość miary z góry daje
$\nu(F)=\lim_m\nu(F_m)\ge\varepsilon_0$.
Przeczy to $\nu\ll\mu$. W przypadku zespolonym
stosujemy wynik do skończonej miary $|\eta|$.
\end{proof}

\begin{uwaga}[Dlaczego skończoności $\nu$ nie wolno pominąć?]
Na $X=(0,1)$ z miarą $\mu=\lambda_1$ określmy
$\nu(E)=0$ dla $\mu(E)=0$ oraz $\nu(E)=\infty$
dla $\mu(E)>0$. Jest to miara: dla rozłącznych
$(E_n)$ suma $\sum_n\nu(E_n)$ jest równa $0$ wtedy
i tylko wtedy, gdy wszystkie $E_n$ są zerowe; w przeciwnym
razie jest równa $\infty$. Mamy $\nu\ll\mu$,
ale nie zachodzi warunek epsilon--delta, bo przedziały
o dowolnie małej dodatniej długości mają $\nu(E)=\infty$.
\end{uwaga}

\begin{wniosek}[Ciągłość całki względem zbioru całkowania]
\label{cor:uniform-abs-cont-int-17}
Niech $f\in L^1(\mu)$ na dowolnej przestrzeni z miarą.
Wtedy dla każdego $\varepsilon>0$ istnieje $\delta>0$ taka,
że $\mu(E)<\delta$ implikuje
$\int_E|f(x)|\,\mathrm d\mu(x)<\varepsilon$.
\end{wniosek}
\begin{proof}
Miara $E\mapsto\int_E|f|\,\mathrm d\mu(x)$ jest skończona
i bezwzględnie ciągła względem $\mu$, zatem stosujemy
poprzednie twierdzenie.
\end{proof}

\section{Rozkład Lebesgue'a: przypadek miar skończonych}
\begin{twierdzenie}[Rozkład Lebesgue'a dla miar skończonych]
\label{thm:lebesgue-decomp-finite-17}
Niech $\mu$ i $\nu$ będą skończonymi miarami nieujemnymi
na tej samej przestrzeni mierzalnej $(X,\mathcal F)$.
Istnieje dokładnie jedna para miar nieujemnych
$(\nu_{\mathrm{ac}},\nu_{\mathrm{s}})$ taka, że
\begin{equation}\label{eq:decomp-finite-17}
\nu=\nu_{\mathrm{ac}}+\nu_{\mathrm{s}},\qquad
\nu_{\mathrm{ac}}\ll\mu,\qquad \nu_{\mathrm{s}}\perp\mu.
\end{equation}
Obie składowe są skończone. W dowodzie nie korzystamy
z twierdzenia Radona--Nikodyma.
\end{twierdzenie}
\begin{proof}
\emph{Konstrukcja.} Dla każdego $k\ge1$ stosujemy rozkład
Hahna (twierdzenie~\ref{thm:hahn-16}) do skończonej miary
ze znakiem $\nu-2^k\mu$. Otrzymujemy rozłączne zbiory
$P_k,N_k\in\mathcal F$ takie, że
$P_k\cup N_k=X$ oraz
\begin{equation}\label{eq:ineq-hahn-decomp-17}
\nu(A)\ge2^k\mu(A)\quad(A\subseteq P_k),\qquad
\nu(B)\le2^k\mu(B)\quad(B\subseteq N_k)
\end{equation}
dla wszystkich mierzalnych $A,B$. W szczególności
$\mu(P_k)\le2^{-k}\nu(P_k)\le2^{-k}\nu(X)$.
Definiujemy
\[
S=\limsup_{k\to\infty}P_k
  =\bigcap_{m\ge1}\bigcup_{k\ge m}P_k.
\]
Ponieważ $\sum_k\mu(P_k)<\infty$, lemat
Borela--Cantellego (lub bezpośrednia podaddytywność)
daje $\mu(S)=0$.

Dla każdego $x\in X\setminus S$ istnieje $m$ takie,
że $x\notin P_k$ dla wszystkich $k\ge m$.
W szczególności
$X\setminus S\subseteq\bigcup_{k\ge1}N_k$.
Jeśli $E\in\mathcal F$ i $\mu(E)=0$, to z drugiej
nierówności w \eqref{eq:ineq-hahn-decomp-17} wynika
$\nu(E\cap N_k)=0$ dla każdego $k$.
Zatem $\nu(E\setminus S)=0$.
Kładziemy
$\nu_{\mathrm{s}}(A)=\nu(A\cap S)$ oraz
$\nu_{\mathrm{ac}}(A)=\nu(A\setminus S)$
dla wszystkich $A\in\mathcal F$.
Obie funkcje są skończonymi miarami nieujemnymi,
a ich suma jest równa $\nu$.
Miara $\nu_{\mathrm{s}}$ jest skupiona na $S$,
przy $\mu(S)=0$, a z poprzedniego akapitu
wynika $\nu_{\mathrm{ac}}\ll\mu$.

\emph{Jednoznaczność.} Załóżmy, że
$\nu=\alpha+\sigma=\alpha'+\sigma'$,
gdzie $\alpha,\alpha'\ll\mu$ oraz
$\sigma,\sigma'\perp\mu$. Niech $T,T'$ będą
mierzalnymi zbiorami $\mu$-zerowymi skupiającymi
odpowiednio $\sigma,\sigma'$, i niech $Z=T\cup T'$.
Wówczas $\mu(Z)=0$, zatem $\alpha(Z)=\alpha'(Z)=0$.
Na $X\setminus Z$ mamy $\sigma=\sigma'=0$,
a więc dla każdego $A\in\mathcal F$
$\alpha(A\setminus Z)=\nu(A\setminus Z)=\alpha'(A\setminus Z)$.
Ponieważ obie miary $\alpha,\alpha'$ zerują się na $Z$,
wynika $\alpha=\alpha'$ i następnie $\sigma=\sigma'$.
\end{proof}

\begin{uwaga}[Dlaczego ten dowód nie jest kołowy?]
Użyliśmy wyłącznie rozkładu Hahna dla skończonych miar ze
znakiem, oszacowania szeregu $\sum_k2^{-k}$ oraz ciągłości
miary. Nie korzystaliśmy z reprezentacji przez gęstość.
Dopiero twierdzenie Radona--Nikodyma pokaże, że składową
$\nu_{\mathrm{ac}}$ można zapisać jako całkę funkcji.
\end{uwaga}

\section{Rozkład Lebesgue'a dla miar sigma-skończonych}
\begin{twierdzenie}[Rozkład Lebesgue'a w przypadku $\sigma$-skończonym]
\label{thm:lebesgue-decomp-sigfin-17}
Niech $\mu$ i $\nu$ będą miarami nieujemnymi,
$\sigma$-skończonymi na $(X,\mathcal F)$.
Wtedy istnieje dokładnie jedna para miar nieujemnych
$\nu_{\mathrm{ac}},\nu_{\mathrm{s}}$, spełniająca
$\nu=\nu_{\mathrm{ac}}+\nu_{\mathrm{s}}$,
$\nu_{\mathrm{ac}}\ll\mu$ i $\nu_{\mathrm{s}}\perp\mu$.
Obie składowe są $\sigma$-skończone.
\end{twierdzenie}
\begin{proof}
Wybieramy mierzalne pokrycia
$X=\bigcup_iA_i=\bigcup_jB_j$ takie, że
$\mu(A_i)<\infty$ i $\nu(B_j)<\infty$.
Rodzina $(A_i\cap B_j)_{i,j\ge1}$ jest przeliczalnym
pokryciem $X$ przez zbiory mające skończoną miarę
względem obu miar. Po uporządkowaniu tej rodziny
i odejmowaniu wcześniejszych zbiorów uzyskujemy
mierzalny, parami rozłączny podział $X=\bigsqcup_{n\ge1}C_n$
z $\mu(C_n),\nu(C_n)<\infty$.

Na każdym $C_n$, wyposażonym w $\sigma$-algebrę śladową,
stosujemy twierdzenie~\ref{thm:lebesgue-decomp-finite-17}.
Niech $S_n\subseteq C_n$ będzie zbiorem o
$\mu(S_n)=0$, na którym skupia się część osobliwa
ograniczenia $\nu|_{C_n}$, a na $C_n\setminus S_n$
miara $\nu|_{C_n}$ jest bezwzględnie ciągła względem
$\mu|_{C_n}$. Przyjmujemy $S=\bigcup_nS_n$.
Ponieważ $\mu(S)=0$, miara
$\nu_{\mathrm{s}}(E):=\nu(E\cap S)$ jest osobliwa.
Jeśli $\mu(E)=0$, to
$\nu((E\setminus S)\cap C_n)=0$ dla każdego $n$;
przeliczalna addytywność daje $\nu(E\setminus S)=0$.
Zatem $\nu_{\mathrm{ac}}(E):=\nu(E\setminus S)$
jest bezwzględnie ciągła względem $\mu$.
Składowe są $\sigma$-skończone, gdyż dominują je
ograniczenia $\nu$ na zbiorach $B_j$.
Jednoznaczność dowodzi się dokładnie tak jak
w przypadku skończonym, bez użycia skończoności masy.
\end{proof}

\begin{wniosek}[Rozkład miary zespolonej]
\label{cor:complex-lebesgue-decomp-17}
Niech $\mu$ będzie nieujemną miarą $\sigma$-skończoną,
a $\eta$ skończoną miarą zespoloną na $(X,\mathcal F)$.
Wtedy istnieje dokładnie jedna para miar zespolonych
$\eta_{\mathrm{ac}}$ i $\eta_{\mathrm{s}}$ taka, że
$\eta=\eta_{\mathrm{ac}}+\eta_{\mathrm{s}}$,
$\eta_{\mathrm{ac}}\ll\mu$ oraz $\eta_{\mathrm{s}}\perp\mu$.
\end{wniosek}
\begin{proof}
Stosujemy twierdzenie~\ref{thm:lebesgue-decomp-sigfin-17}
do nieujemnej skończonej miary $|\eta|$, otrzymując
zbiór $S$ o $\mu(S)=0$, taki że
$|\eta||_{X\setminus S}\ll\mu$.
Definiujemy $\eta_{\mathrm{s}}(E)=\eta(E\cap S)$
oraz $\eta_{\mathrm{ac}}(E)=\eta(E\setminus S)$.
Dla każdego $E$ wariacja ograniczenia $\eta|_S$
nie przekracza $|\eta|(E\cap S)$, co dowodzi osobliwości;
podobnie $|\eta_{\mathrm{ac}}|(E)\le|\eta|(E\setminus S)$,
co daje bezwzględną ciągłość.
Dla jednoznaczności bierzemy sumę zbiorów
$\mu$-zerowych skupiających dwie hipotetyczne części
osobliwe i powtarzamy argument z dowodu twierdzenia
\ref{thm:lebesgue-decomp-finite-17}, z uwzględnieniem
zerowości części bezwzględnie ciągłych na tych zbiorach.
\end{proof}

\section{Co może być składową osobliwą?}
\begin{przyklad}[Rozkład dyskretno-ciągły]
Na $X=(0,1)$ z miarą $\mu=\lambda_1|_X$ określmy
\[
\nu(E)=\int_E(1+x)\,\mathrm d\lambda_1(x)
       +3\delta_{1/3}(E)+2\delta_{2/3}(E).
\]
Wtedy $\nu_{\mathrm{ac}}(E)=\int_E(1+x)\,\mathrm d\lambda_1(x)$,
a $\nu_{\mathrm{s}}=3\delta_{1/3}+2\delta_{2/3}$.
Jest to rozkład Lebesgue'a: część osobliwa jest skupiona
na dwuelementowym zbiorze miary Lebesgue'a zero.
\end{przyklad}

\begin{przyklad}[Miara Cantora — osobliwość bez atomów]
\label{ex:cantor-measure-17}
Podamy konstrukcję niewymagającą wcześniejszego twierdzenia
o miarach Lebesgue'a--Stieltjesa. Dla $x\in[0,1)$ niech
\[
b_n(x)=\lfloor2^nx\rfloor-2\lfloor2^{n-1}x\rfloor\in\{0,1\},
\qquad n\ge1.
\]
Są to cyfry rozwinięcia dwójkowego, przy czym dla liczb
dwójkowo wymiernych przyjęliśmy rozwinięcie zakończone
samymi zerami. Określamy
$T(x)=\sum_{n\ge1}2b_n(x)3^{-n}$ dla $x\in[0,1)$
oraz $T(1)=1$. Każda funkcja $b_n$ jest borelowska,
a szereg definiujący $T$ jest jednostajnie zbieżny,
więc $T:[0,1]\to[0,1]$ jest borelowska. Jej wartości
należą do środkowego zbioru Cantora $C$, ponieważ ich
rozwinięcia trójkowe używają wyłącznie cyfr $0$ i $2$.

Niech $\kappa=T_*(\lambda_1|_{[0,1]})$. Jest to miara
probabilistyczna, $\kappa(C)=1$. Na poziomie $n$
konstrukcji zbioru Cantora mamy $2^n$ rozłącznych
domkniętych przedziałów podstawowych. Przeciwobraz
każdego z nich przez $T$ zawiera się w jednym przedziale
dwójkowym długości $2^{-n}$, z wyjątkiem ewentualnego
zbioru końców dyadycznych o mierze zero. Zatem
$\kappa(I)\le2^{-n}$ dla każdego przedziału podstawowego
$I$ poziomu $n$. Każdy $y\in C$ należy do jednego
takiego przedziału na każdym poziomie, skąd
$\kappa(\{y\})\le2^{-n}$ dla każdego $n$ i
$\kappa(\{y\})=0$. Dla $y\notin C$ mamy
$\kappa(\{y\})=0$, gdyż $\kappa$ jest skupiona na $C$.
Miara $\kappa$ jest więc bezatomowa. Ponieważ
$\lambda_1(C)=0$, otrzymujemy $\kappa\perp\lambda_1$.
Pokazuje to, że \emph{osobliwość} nie jest tym samym co
\emph{dyskretność}. Dystrybuanta miary $\kappa$ jest
klasyczną funkcją Cantora.
\end{przyklad}

\begin{uwaga}[Trzy składniki klasycznego rozkładu na prostej]
Dla skończonej miary borelowskiej na $\mathbb R$ część
osobliwą można dalej rozdzielić na część atomową,
skupioną na przeliczalnym zbiorze, oraz część osobliwą
bezatomową. Zatem otrzymujemy składowe: bezwzględnie
ciągłą, osobliwą atomową i osobliwą bezatomową.
To \emph{dalszy} podział względem miary Lebesgue'a,
nie część definicji rozkładu Lebesgue'a.
\end{uwaga}

\begin{twierdzenie}[Dalszy rozkład miary osobliwej na prostej]
\label{thm:atomic-cont-singular-17}
Niech $\nu$ będzie skończoną miarą borelowską na
$\mathbb R$. Zbiór
$A=\{x\in\mathbb R:\nu(\{x\})>0\}$ jest co najwyżej
przeliczalny, a miary
\[
 \nu_{\mathrm{at}}(E):=\nu(E\cap A),\qquad
 \nu_{\mathrm{na}}(E):=\nu(E\setminus A)
 \quad(E\in\mathcal B(\mathbb R))
\]
są jednoznacznie określone. Pierwsza ma postać
$\nu_{\mathrm{at}}=\sum_{a\in A}\nu(\{a\})\delta_a$,
a druga jest bezatomowa. Jeśli $\nu\perp\lambda_1$,
to również $\nu_{\mathrm{at}}\perp\lambda_1$ i
$\nu_{\mathrm{na}}\perp\lambda_1$.
\end{twierdzenie}
\begin{proof}
Dla każdego $n\ge1$ zbiór
$A_n=\{x:\nu(\{x\})\ge1/n\}$ jest skończony:
gdyby zawierał $m$ różnych punktów, mielibyśmy
$\nu(\mathbb R)\ge m/n$. Ponieważ
$A=\bigcup_{n\ge1}A_n$, jest przeliczalny.
Przeliczalna addytywność daje podany wzór na
$\nu_{\mathrm{at}}$ oraz rozkład $\nu$.
Z definicji $\nu_{\mathrm{na}}(\{x\})=0$ dla
każdego $x$. Wykażemy, że miara ta nie ma także
żadnych atomów w sensie teorii miary.

Przypuśćmy, że mierzalny zbiór $B$ jest atomem
$\nu_{\mathrm{na}}$ o dodatniej masie $b$:
każdy mierzalny podzbiór $B$ ma masę $0$ lub $b$.
Ponieważ $\mathbb R=\bigcup_{m\ge1}[-m,m]$,
pewien domknięty ograniczony przedział $I_0$
spełnia $\nu_{\mathrm{na}}(B\cap I_0)=b$.
Dzielimy $I_0$ na dwa domknięte przedziały równej
długości. Co najmniej jeden z nich, nazwany $I_1$,
ma $\nu_{\mathrm{na}}(B\cap I_1)>0$, a więc masę $b$.
Powtarzając, otrzymujemy zstępujący ciąg przedziałów
zamkniętych $I_n$ o długościach malejących do zera,
dla których $\nu_{\mathrm{na}}(B\cap I_n)=b$.
Ich przecięcie jest singletonem $\{x\}$.
Z ciągłości skończonej miary z góry mamy
$\nu_{\mathrm{na}}(B\cap\{x\})=b>0$,
sprzecznie z zerowością wszystkich singletonów.
Zatem $\nu_{\mathrm{na}}$ jest bezatomowa.

Jeśli $\nu$ skupia się na zbiorze borelowskim $S$
o $\lambda_1(S)=0$, to obie części jako ograniczenia
$\nu$ również skupiają się na $S$, więc są osobliwe.
\end{proof}

\section{Główna idea na przyszłość}
W rozdziale~\ref{ch:radon-nikodym-18} udowodnimy, że przy
$\sigma$-skończoności obie mierzone przestrzenie można porównać
jeszcze precyzyjniej: istnieje mierzalna funkcja
$f\ge0$ taka, że
$\nu_{\mathrm{ac}}(E)=\int_E f\,\mathrm d\mu(x)$.
Twierdzenie o rozkładzie Lebesgue'a mówi zatem,
\emph{którą część miary warto próbować opisywać gęstością};
twierdzenie Radona--Nikodyma daje tę gęstość i
zapewnia jej jednoznaczność prawie wszędzie.

% END chapters/pl/measure_theory/17-rozdzial.tex


% BEGIN chapters/pl/measure_theory/18-rozdzial.tex
% ROZDZIAŁ 18 — Radon-Nikodym, pełny dowód przez przestrzeń Hilberta
\chapter[Twierdzenie Radona--Nikodyma]{Twierdzenie Radona--Nikodyma i pochodne miar}
\label{ch:radon-nikodym-18}

\section{Problem reprezentacji miary przez funkcję}
Niech $(X,\mathcal F,\mu)$ będzie przestrzenią z miarą,
a $\nu$ drugą miarą nieujemną na tej samej $\sigma$-algebrze.
Chcielibyśmy znaleźć mierzalną funkcję $f\ge0$, dla której
\begin{equation}\label{eq:RN-question-18}
\nu(E)=\int_E f(x)\,\mathrm d\mu(x)
\qquad(E\in\mathcal F).
\end{equation}
Wynikałoby stąd natychmiast $\nu\ll\mu$.
Czy warunek ten jest wystarczający? Odpowiedź jest twierdząca
przy odpowiednich założeniach o skończoności miar.

\begin{definicja}[Pochodna Radona--Nikodyma]
\label{def:RN-derivative-18}
Jeśli $\nu$ i $\mu$ są miarami nieujemnymi na $(X,\mathcal F)$
i istnieje mierzalna funkcja $f:X\to[0,\infty]$
spełniająca \eqref{eq:RN-question-18}, nazywamy ją
\emph{gęstością} lub \emph{pochodną Radona--Nikodyma}
miary $\nu$ względem $\mu$ i piszemy
$f=\mathrm d\nu/\mathrm d\mu$. Jeżeli taka pochodna
jest jednoznaczna z dokładnością do równości $\mu$-prawie
wszędzie, symbol ten oznacza jej klasę równoważności.
W tej definicji symbol $\mathrm d\nu/\mathrm d\mu$
nie jest zwykłym ilorazem dwóch funkcji.
\end{definicja}

\begin{lemat}[Całkowanie względem miary o gęstości]
\label{lem:integration-density-18}
Niech $\lambda$ będzie miarą nieujemną na $(X,\mathcal F)$,
$w:X\to[0,\infty]$ funkcją mierzalną oraz
$\mu(E)=\int_Ew\,\mathrm d\lambda$ dla $E\in\mathcal F$.
Wówczas dla każdej mierzalnej $\varphi:X\to[0,\infty]$
zachodzi
\begin{equation}\label{eq:integration-density-18}
\int_X\varphi(x)\,\mathrm d\mu(x)
=\int_X\varphi(x)w(x)\,\mathrm d\lambda(x),
\end{equation}
przy konwencji $0\cdot\infty=0$ dla iloczynu
nieujemnych funkcji mierzalnych. Jeśli $\varphi$ jest
zespolona i całkowalna względem $\mu$, to analogiczna
równość zachodzi dla całek zespolonych.
\end{lemat}
\begin{proof}
Dla $\varphi=\chi_A$ jest to definicja $\mu$.
Z liniowości dla funkcji prostych nieujemnych otrzymujemy
równość dla wszystkich takich funkcji. Każdą funkcję
$\varphi\ge0$ aproksymujemy z dołu ciągiem mierzalnych
funkcji prostych $\varphi_n\uparrow\varphi$.
Dwukrotne zastosowanie twierdzenia Beppo Leviego
daje \eqref{eq:integration-density-18}. W przypadku
funkcji zespolonej stosujemy tę tożsamość do części
dodatnich i ujemnych jej części rzeczywistej i urojonej,
które są całkowalne dzięki założeniu.
\end{proof}

\section{Krok pierwszy: miary skończone}
\begin{twierdzenie}[Radon--Nikodym dla miar skończonych]
\label{thm:RN-finite-18}
Niech $\mu,\nu$ będą skończonymi miarami nieujemnymi
na $(X,\mathcal F)$ i niech $\nu\ll\mu$.
Istnieje mierzalna funkcja $f:X\to[0,\infty)$ taka, że
\begin{equation}\label{eq:RN-finite-18}
\nu(E)=\int_E f(x)\,\mathrm d\mu(x)
\quad\text{dla każdego }E\in\mathcal F.
\end{equation}
Funkcja $f$ należy do $L^1(\mu)$, jest jednoznaczna
$\mu$-prawie wszędzie i
$\int_Xf\,\mathrm d\mu(x)=\nu(X)$.
\end{twierdzenie}
\begin{proof}
\emph{Etap 1: wspólna miara dominująca.}
Definiujemy $\lambda=\mu+\nu$. Jest to miara skończona
na $\mathcal F$. Dla $u\in L^2(\lambda;\mathbb C)$
kładziemy
\[
T(u)=\int_X u(x)\,\mathrm d\nu(x).
\]
Funkcjonał ten jest dobrze określony: ponieważ
$\nu\le\lambda$ w sensie wartości na zbiorach,
$\int_X|u|^2\,\mathrm d\nu(x)\le
\int_X|u|^2\,\mathrm d\lambda$.
Nierówność Cauchy'ego--Schwarza daje
\[
|T(u)|\le\nu(X)^{1/2}
\left(\int_X|u|^2\,\mathrm d\nu(x)\right)^{1/2}
\le\nu(X)^{1/2}\|u\|_{L^2(\lambda)}.
\]
Jeśli dwa przedstawiciele $u$ są równe $\lambda$-prawie
wszędzie, to są równe $\nu$-prawie wszędzie, więc
$T$ jest funkcjonałem na klasach $L^2(\lambda)$.

\emph{Etap 2: reprezentacja Hilbertowska.}
Wykorzystujemy twierdzenie Riesza o reprezentacji
funkcjonałów na przestrzeni Hilberta $L^2(\lambda)$:
istnieje $h\in L^2(\lambda)$ takie, że
\begin{equation}\label{eq:RN-lambda-h-18}
T(u)=\int_Xu(x)h(x)\,\mathrm d\lambda(x)
\qquad(u\in L^2(\lambda)).
\end{equation}
Dla kompletności przypomnijmy uzasadnienie użytej
reprezentacji. Jeśli $T=0$, bierzemy $h=0$.
W przeciwnym razie $K=\ker T$ jest domkniętą
podprzestrzenią Hilberta; każdy $v$ ma jednoznaczny
rzut ortogonalny $v_K$ na $K$, którego istnienie wynika
z zupełności i tożsamości równoległoboku: ciąg
minimalizujący odległość od $K$ jest ciągiem Cauchy'ego.
Wybieramy $v\in L^2(\lambda)$ z $T(v)=1$ i rzutujemy
$v$ na $K^\perp$. Otrzymany wektor $u$ spełnia
$u\in K^\perp$ i $T(u)=1$, ponieważ $T$ zeruje się na $K$.
Dla dowolnego $v$ różnica $v-T(v)u$ należy do $K$,
więc przy iloczynie skalarnym liniowym w pierwszym
argumencie mamy
$\langle v,u\rangle=T(v)\|u\|_2^2$.
Zatem $T(v)=\langle v,u\rangle/\|u\|_2^2$,
a reprezentantem w zapisie \eqref{eq:RN-lambda-h-18}
jest $h=\overline u/\|u\|_2^2$.
Nie użyliśmy w tym miejscu twierdzenia o gęstościach miar.

\emph{Etap 3: oszacowanie $0\le h\le1$.}
Ponieważ $\lambda(X)<\infty$, funkcje charakterystyczne $\chi_E$
należą do $L^2(\lambda)$. Z \eqref{eq:RN-lambda-h-18}
otrzymujemy
\begin{equation}\label{eq:RN-mu-nu-h-18}
\nu(E)=\int_Eh\,\mathrm d\lambda,\qquad
\mu(E)=\int_E(1-h)\,\mathrm d\lambda.
\end{equation}
Pierwsza całka jest rzeczywista i nieujemna dla
każdego $E$. Wybierając kolejno zbiory, na których
$\operatorname{Im}h$ jest dodatnia lub ujemna,
dostajemy $\operatorname{Im}h=0$ $\lambda$-prawie wszędzie.
Podobnie z nieujemności obu całek w
\eqref{eq:RN-mu-nu-h-18} otrzymujemy
$h\ge0$ i $1-h\ge0$ prawie wszędzie.
Po modyfikacji na zbiorze $\lambda$-zerowym możemy
więc założyć, że $0\le h\le1$ wszędzie.

\emph{Etap 4: wykluczenie obszaru $h=1$.}
Niech $S=\{x:h(x)=1\}$. Druga równość
w \eqref{eq:RN-mu-nu-h-18} daje $\mu(S)=0$.
Z założenia $\nu\ll\mu$ wynika $\nu(S)=0$.
Zatem $\lambda(S)=\mu(S)+\nu(S)=0$.
Poza $S$ mianownik $1-h$ jest dodatni.
Definiujemy
\[
f(x)=\begin{cases}
\dfrac{h(x)}{1-h(x)},&h(x)<1,\\[4pt]
0,&h(x)=1.
\end{cases}
\]
Funkcja $f$ jest mierzalna i nieujemna. Na mocy
lematu~\ref{lem:integration-density-18}, zastosowanego
do $\mu(E)=\int_E(1-h)\,\mathrm d\lambda$,
\[
\int_E f\,\mathrm d\mu(x)
 =\int_E f(1-h)\,\mathrm d\lambda
 =\int_Eh\,\mathrm d\lambda
 =\nu(E).
\]
W środkowej równości ewentualne punkty $S$ nie
mają znaczenia, gdyż $\lambda(S)=0$.
Dla $E=X$ mamy $\int_Xf\,\mathrm d\mu(x)=\nu(X)<\infty$,
więc $f$ jest skończona $\mu$-prawie wszędzie
i należy do $L^1(\mu)$.

\emph{Etap 5: jednoznaczność.}
Jeżeli $f,g\in L^1(\mu)$, $f,g\ge0$, spełniają
\eqref{eq:RN-finite-18}, to dla każdego $E\in\mathcal F$
$\int_E(f-g)\,\mathrm d\mu(x)=0$.
Biorąc $E=\{f>g\}$, otrzymujemy
$\int_{\{f>g\}}(f-g)\,\mathrm d\mu(x)=0$.
Całkowana funkcja jest nieujemna, więc
$\mu(\{f>g\})=0$. Zamiana $f$ z $g$ daje
$\mu(\{g>f\})=0$. Stąd $f=g$ prawie wszędzie.
\end{proof}

\begin{uwaga}[Rola pomocniczej miary $\lambda=\mu+\nu$]
Dowód nie zakłada z góry, że $\nu$ ma gęstość względem
$\mu$. Najpierw reprezentujemy funkcjonał całkowy w
$L^2(\mu+\nu)$, a następnie uzyskujemy dwie gęstości
$h$ i $1-h$ względem tej samej miary $\lambda$.
Dopiero warunek $\nu\ll\mu$ pozwala wykluczyć
zbiór $\{h=1\}$ i podzielić te gęstości.
W ten sposób znika potencjalne błędne koło.
\end{uwaga}

\section{Przypadek sigma-skończony: twierdzenie główne}
\begin{twierdzenie}[Twierdzenie Radona--Nikodyma]
\label{thm:RN-sigfin-18}
Niech $\mu$ i $\nu$ będą miarami nieujemnymi,
$\sigma$-skończonymi na przestrzeni mierzalnej $(X,\mathcal F)$.
Jeśli $\nu\ll\mu$, to istnieje mierzalna funkcja
$f:X\to[0,\infty)$, skończona $\mu$-prawie wszędzie,
taka, że
\begin{equation}\label{eq:RN-general-18}
\nu(E)=\int_E f(x)\,\mathrm d\mu(x)
\qquad(E\in\mathcal F).
\end{equation}
Funkcja $f$ jest jednoznaczna $\mu$-prawie wszędzie.
Dla każdej mierzalnej funkcji $\varphi\ge0$ zachodzi
\begin{equation}\label{eq:RN-integral-change-18}
\int_X\varphi(x)\,\mathrm d\nu(x)
=\int_X\varphi(x)f(x)\,\mathrm d\mu(x).
\end{equation}
W szczególności, jeśli $\nu(X)<\infty$, to
$f\in L^1(\mu)$ i $\int_Xf\,\mathrm d\mu(x)=\nu(X)$.
\end{twierdzenie}
\begin{proof}
Ponieważ obie miary są $\sigma$-skończone,
można wybrać mierzalny, parami rozłączny podział
$X=\bigsqcup_{n\ge1}C_n$ taki, że
$\mu(C_n)<\infty$ i $\nu(C_n)<\infty$.
Jego konstrukcję podaliśmy w dowodzie twierdzenia
\ref{thm:lebesgue-decomp-sigfin-17}.
Dla ograniczeń obu miar do $C_n$ zachodzi
$\nu|_{C_n}\ll\mu|_{C_n}$. Z twierdzenia
\ref{thm:RN-finite-18} otrzymujemy mierzalną
$\mathcal F|_{C_n}$-funkcję $f_n\ge0$ taką, że
$\nu(E\cap C_n)=\int_{E\cap C_n}f_n\,\mathrm d\mu(x)$
dla wszystkich $E\in\mathcal F$.
Funkcja $f$, określona przez $f|_{C_n}=f_n$,
jest mierzalna, gdyż $C_n$ tworzą przeliczalny
podział mierzalny. Dla każdego $E$ twierdzenie
Beppo Leviego i przeliczalna addytywność $\nu$ dają
\[
\int_E f\,\mathrm d\mu(x)
=\sum_{n\ge1}\int_{E\cap C_n}f_n\,\mathrm d\mu(x)
=\sum_{n\ge1}\nu(E\cap C_n)=\nu(E).
\]
Z sumowalności $\int_{C_n}f_n\,\mathrm d\mu(x)=\nu(C_n)<\infty$
wynika, że $f$ jest skończona prawie wszędzie na każdym
$C_n$, a więc także na $X$.

Jeśli $f,g$ reprezentują $\nu$, to na każdym $C_n$
obie funkcje są całkowalne i z jednoznaczności w
przypadku skończonym są równe $\mu$-prawie wszędzie.
Po zsumowaniu przeliczalnej rodziny zbiorów wyjątkowych
otrzymujemy jednoznaczność na $X$.
Wzór \eqref{eq:RN-integral-change-18} jest zastosowaniem
lematu~\ref{lem:integration-density-18}.
\end{proof}

\begin{uwaga}[Co dokładnie zapewnia $\sigma$-skończoność?]
Założenie dotyczy \emph{obu} miar i wystarcza, aby rozbić
$X$ na przeliczalnie wiele części, na których
$\mu$ i $\nu$ są skończone. Nie jest ono warunkiem
koniecznym w każdej konkretnej przestrzeni;
istnieją szersze klasy miar, dla których reprezentacja
Radona--Nikodyma również działa. Bez dalszych warunków
nie można go jednak po prostu wykreślić.
\end{uwaga}

\begin{przyklad}[Brak $\sigma$-skończoności może zniszczyć wniosek]
\label{ex:RN-failure-nonsigma-18}
Niech $X=[0,1]$ z $\sigma$-algebrą borelowską,
$\mu$ miarą liczącą, a $\nu$ ograniczeniem miary
Lebesgue'a do zbiorów borelowskich. Miara $\mu$
nie jest $\sigma$-skończona, gdyż każdy jej zbiór
skończonej miary jest skończony, a przeliczalna suma
takich zbiorów nie pokrywa nieprzeliczalnego $X$.
Ponieważ jedynym zbiorem $\mu$-zerowym jest $\varnothing$,
mamy $\nu\ll\mu$. Przypuśćmy, że
$\nu(E)=\int_Ef\,\mathrm d\mu(x)$ dla pewnej
mierzalnej $f\ge0$. Dla każdego $x\in X$:
\[
0=\nu(\{x\})=\int_{\{x\}}f\,\mathrm d\mu(x)=f(x).
\]
Zatem $f=0$ wszędzie, wbrew
$\nu(X)=1$. Jest to nawet przykład z miarą $\mu$
\emph{półskończoną}: każdy zbiór nieskończonej miary
zawiera niepusty zbiór skończonej dodatniej miary.
\end{przyklad}

\section{Pochodna miary ze znakiem i miary zespolonej}
\begin{twierdzenie}[Radon--Nikodym dla miar zespolonych]
\label{thm:RN-complex-18}
Niech $\mu$ będzie nieujemną miarą $\sigma$-skończoną,
a $\eta$ skończoną miarą zespoloną na $(X,\mathcal F)$.
Jeżeli $\eta\ll\mu$, to istnieje dokładnie jedna
(z dokładnością do równości $\mu$-prawie wszędzie)
funkcja $f\in L^1(\mu;\mathbb C)$, dla której
\begin{equation}\label{eq:RN-complex-18}
\eta(E)=\int_E f(x)\,\mathrm d\mu(x)
\qquad(E\in\mathcal F).
\end{equation}
Ponadto $|\eta|(E)=\int_E|f|\,\mathrm d\mu(x)$,
a w szczególności $\|\eta\|_{\mathrm{TV}}=\|f\|_1$.
Dla rzeczywistej miary ze znakiem $\eta$ funkcję $f$
można wybrać rzeczywistą.
\end{twierdzenie}
\begin{proof}
Jeśli $\eta$ jest rzeczywistą skończoną miarą ze znakiem,
rozważamy jej rozkład Jordana
$\eta=\eta^+-\eta^-$ z rozdziału~\ref{ch:miary-ze-znakiem}.
Ponieważ $\eta\ll\mu$, dla każdego mierzalnego
zbioru $E$ o $\mu(E)=0$ również wszystkie podzbiory
mierzalne $A\subseteq E$ spełniają $\eta(A)=0$.
Z wzorów na części Jordana wynika
$\eta^+(E)=\eta^-(E)=0$. Obie miary dodatnie są
skończone, a więc $\sigma$-skończone, i są
bezwzględnie ciągłe względem $\mu$.
Stosujemy twierdzenie~\ref{thm:RN-sigfin-18} do każdej
z nich, otrzymując $f_+,f_-\in L^1(\mu)$.
Wtedy $f=f_+-f_-$ jest żądaną reprezentacją.
Dla miary zespolonej stosujemy ten argument do
$\operatorname{Re}\eta$ i $\operatorname{Im}\eta$
i otrzymujemy $f=f_1+if_2\in L^1(\mu;\mathbb C)$.
Jednoznaczność wynika z jednoznaczności części
rzeczywistych i urojonych, a wzór na wariację
jest twierdzeniem~\ref{thm:density-variation-16}.
\end{proof}

\begin{wniosek}[Pełny rozkład miary na część z gęstością i osobliwą]
\label{cor:RN-lebesgue-sum-18}
Niech $\mu,\nu$ będą nieujemnymi miarami
$\sigma$-skończonymi na $(X,\mathcal F)$.
Istnieją dokładnie jedna (z dokładnością do
$\mu$-prawie wszędzie) funkcja mierzalna
$f:X\to[0,\infty)$ i dokładnie jedna miara
$\nu_{\mathrm{s}}\perp\mu$ takie, że
\begin{equation}\label{eq:RN-Lebesgue-decomp-18}
\nu(E)=\int_E f(x)\,\mathrm d\mu(x)+\nu_{\mathrm{s}}(E)
\qquad(E\in\mathcal F).
\end{equation}
Dla skończonej miary zespolonej $\eta$ i
$\sigma$-skończonej $\mu$ analogiczna reprezentacja ma
$f\in L^1(\mu;\mathbb C)$ i osobliwą miarę zespoloną.
\end{wniosek}
\begin{proof}
Twierdzenie~\ref{thm:lebesgue-decomp-sigfin-17}
daje jednoznaczny rozkład $\nu=\nu_{\mathrm{ac}}+
\nu_{\mathrm{s}}$. Obie części są $\sigma$-skończone,
a więc \ref{thm:RN-sigfin-18} daje
$\nu_{\mathrm{ac}}(E)=\int_Ef\,\mathrm d\mu(x)$.
Jednoznaczność wynika z jednoznaczności obu składników
i pochodnej. Wersję zespoloną otrzymujemy z
wniosku~\ref{cor:complex-lebesgue-decomp-17}
i twierdzenia~\ref{thm:RN-complex-18}.
\end{proof}

\begin{twierdzenie}[Rozkład biegunowy miary zespolonej]
\label{thm:polar-complex-measure-18}
Niech $\eta$ będzie skończoną miarą zespoloną
na $(X,\mathcal F)$, a $|\eta|$ jej wariacją
całkowitą. Istnieje funkcja mierzalna
$u:X\to\mathbb C$, określona jednoznacznie
$|\eta|$-prawie wszędzie, taka, że
\begin{equation}\label{eq:polar-complex-measure-18}
 |u(x)|=1\quad |\eta|\text{-prawie wszędzie},\qquad
 \eta(E)=\int_Eu(x)\,\mathrm d|\eta|(x)
 \quad(E\in\mathcal F).
\end{equation}
Jeśli $\eta$ jest rzeczywistą miarą ze znakiem, to
$u$ można wybrać o wartościach w $\{-1,1\}$.
\end{twierdzenie}
\begin{proof}
Z definicji wariacji wynika $\eta\ll|\eta|$.
Miara $|\eta|$ jest skończona, toteż twierdzenie
\ref{thm:RN-complex-18} daje $u\in L^1(|\eta|;\mathbb C)$
spełniającą drugą równość w~\eqref{eq:polar-complex-measure-18}.
Twierdzenie~\ref{thm:density-variation-16} zastosowane
do tej reprezentacji daje
$|\eta|(E)=\int_E|u(x)|\,\mathrm d|\eta|(x)$.
Z drugiej strony $|\eta|(E)=\int_E1\,\mathrm d|\eta|(x)$.
Jednoznaczność pochodnej Radona--Nikodyma implikuje
$|u|=1$ $|\eta|$-prawie wszędzie.
W przypadku rzeczywistym gęstość jest rzeczywista,
więc warunek $|u|=1$ daje $u\in\{-1,1\}$
prawie wszędzie. Ustalając wartość $1$ na pozostałym
zbiorze zerowym, otrzymujemy żądaną funkcję mierzalną.
\end{proof}

\begin{przyklad}[Rozkład biegunowy jako znak gęstości]
Niech $\eta(E)=\int_E f(x)\,\mathrm d\mu(x)$,
gdzie $f\in L^1(\mu;\mathbb C)$. Wówczas
$|\eta|(E)=\int_E|f(x)|\,\mathrm d\mu(x)$,
a funkcję z poprzedniego twierdzenia można przyjąć
równą $u=f/|f|$ na $\{f\neq0\}$ oraz $u=1$
na $\{f=0\}$. W szczególności dla rzeczywistego $f$
jest to odpowiednik znaku funkcji.
\end{przyklad}

\section{Reguły rachunku dla pochodnych miar}
\begin{twierdzenie}[Reguła łańcuchowa]
\label{thm:RN-chain-18}
Niech $\mu,\nu,\rho$ będą nieujemnymi miarami
$\sigma$-skończonymi na $(X,\mathcal F)$
i niech $\rho\ll\nu\ll\mu$.
Oznaczmy $f=\mathrm d\nu/\mathrm d\mu$
i $g=\mathrm d\rho/\mathrm d\nu$.
Wtedy $\rho\ll\mu$ i
\begin{equation}\label{eq:RN-chain-18}
\frac{\mathrm d\rho}{\mathrm d\mu}
=g f\qquad\mu\text{-prawie wszędzie},
\end{equation}
przy czym na zbiorze $\{f=0\}$ przyjmujemy $gf=0$
niezależnie od wartości reprezentanta $g$.
\end{twierdzenie}
\begin{proof}
Bezwzględna ciągłość $\rho\ll\mu$ wynika
z przechodniości implikacji dla zbiorów zerowych.
Z lematu~\ref{lem:integration-density-18}, dla każdego
$E\in\mathcal F$,
\[
\rho(E)=\int_Eg\,\mathrm d\nu(x)
=\int_Eg f\,\mathrm d\mu(x).
\]
Ponieważ zmiana $g$ na zbiorze $\nu$-zerowym nie zmienia
całki względem $\nu$, a $\nu(\{f=0\})=0$,
przyjęta konwencja jest spójna. Jednoznaczność
pochodnej daje \eqref{eq:RN-chain-18}.
\end{proof}

\begin{twierdzenie}[Pochodna odwrotna]
\label{thm:RN-inverse-18}
Niech $\mu,\nu$ będą nieujemnymi miarami
$\sigma$-skończonymi i $\mu\sim\nu$.
Jeżeli $f=\mathrm d\nu/\mathrm d\mu$, to
$f>0$ $\mu$-prawie wszędzie i
\begin{equation}\label{eq:RN-inverse-18}
\frac{\mathrm d\mu}{\mathrm d\nu}=\frac1f
\qquad\nu\text{-prawie wszędzie}.
\end{equation}
\end{twierdzenie}
\begin{proof}
Z równoważności $\mu\ll\nu$ i wzoru
$\nu(\{f=0\})=\int_{\{f=0\}}f\,\mathrm d\mu(x)=0$
wynika $\mu(\{f=0\})=0$. Zatem $1/f$ jest
zdefiniowana i skończona prawie wszędzie.
Dla każdego $E\in\mathcal F$ lemat
\ref{lem:integration-density-18} daje
\[
\int_E\frac1f\,\mathrm d\nu(x)
=\int_E\frac1f f\,\mathrm d\mu(x)=\mu(E).
\]
Jednoznaczność pochodnej kończy dowód.
\end{proof}

\begin{wniosek}[Różnica i suma miar z gęstościami]
\label{cor:RN-linearity-18}
Niech $\mu$ będzie miarą $\sigma$-skończoną,
a $\eta_1,\eta_2$ skończonymi miarami zespolonymi,
bezwzględnie ciągłymi względem $\mu$.
Dla $a,b\in\mathbb C$ zachodzi
\[
\frac{\mathrm d(a\eta_1+b\eta_2)}{\mathrm d\mu}
=a\frac{\mathrm d\eta_1}{\mathrm d\mu}
+b\frac{\mathrm d\eta_2}{\mathrm d\mu}
\quad\mu\text{-prawie wszędzie}.
\]
\end{wniosek}
\begin{proof}
Wstawiamy reprezentacje całkowe z twierdzenia
\ref{thm:RN-complex-18}, korzystamy z liniowości całki,
a następnie z jednoznaczności gęstości.
\end{proof}

\section{Zastosowanie: warunkowa wartość oczekiwana}
Przedstawimy zastosowanie pokazujące, że twierdzenie
Radona--Nikodyma nie jest jedynie wynikiem o porównywaniu
miar. Pozwala ono skonstruować jeden z podstawowych
operatorów teorii prawdopodobieństwa.

\begin{definicja}[Warunkowa wartość oczekiwana]
\label{def:conditional-exp-18}
Niech $(\Omega,\mathcal F,\mathbb P)$ będzie przestrzenią
probabilistyczną, $\mathcal G\subseteq\mathcal F$
pod-$\sigma$-algebrą, a $Y\in L^1(\mathbb P;\mathbb R)$.
\emph{Warunkową wartością oczekiwaną} $Y$ względem
$\mathcal G$ nazywamy klasę funkcji
$Z\in L^1(\mathbb P;\mathbb R)$, która jest
$\mathcal G$-mierzalna i spełnia
\begin{equation}\label{eq:CE-def-18}
\int_A Z(\omega)\,\mathrm d\mathbb P(\omega)
=\int_A Y(\omega)\,\mathrm d\mathbb P(\omega)
\qquad(A\in\mathcal G).
\end{equation}
Piszemy $Z=\mathbb E[Y\mid\mathcal G]$.
Równość jest rozumiana z dokładnością do
$\mathbb P$-prawie wszystkich punktów. Ściślej, pracujemy
z klasami funkcji mierzalnych względem $\mathcal G$;
nie każda modyfikacja na zbiorze $\mathbb P$-zerowym musi
pozostać $\mathcal G$-mierzalna, gdy $\mathcal G$ nie jest zupełna.
\end{definicja}

\begin{twierdzenie}[Istnienie, jednoznaczność i własności podstawowe]
\label{thm:CE-existence-18}
Dla każdej przestrzeni probabilistycznej
$(\Omega,\mathcal F,\mathbb P)$, każdej
pod-$\sigma$-algebry $\mathcal G$ oraz $Y\in L^1(\mathbb P;\mathbb R)$
istnieje dokładnie jedna warunkowa wartość oczekiwana
$Z=\mathbb E[Y\mid\mathcal G]$ w sensie równości
$\mathbb P$-prawie wszędzie. Ponadto:
\begin{enumerate}[label=\textup{(\roman*)}]
\item odwzorowanie $Y\mapsto\mathbb E[Y\mid\mathcal G]$
jest liniowe i $\|Z\|_{L^1}\le\|Y\|_{L^1}$;
\item $Y\ge0$ prawie wszędzie implikuje $Z\ge0$
prawie wszędzie;
\item jeśli $Y$ jest $\mathcal G$-mierzalna,
to $Z=Y$ prawie wszędzie;
\item jeśli $\mathcal H\subseteq\mathcal G\subseteq\mathcal F$,
to $\mathbb E[\mathbb E[Y\mid\mathcal G]\mid\mathcal H]
=\mathbb E[Y\mid\mathcal H]$ prawie wszędzie.
\end{enumerate}
\end{twierdzenie}
\begin{proof}
Na przestrzeni mierzalnej $(\Omega,\mathcal G)$
określamy skończoną miarę ze znakiem
$\eta(A)=\int_A Y\,\mathrm d\mathbb P$.
Jest ona bezwzględnie ciągła względem
$\mu=\mathbb P|_{\mathcal G}$, a $\mu$ jest miarą
skończoną. Twierdzenie~\ref{thm:RN-complex-18}
(dla rzeczywistych miar) daje $\mathcal G$-mierzalną
$Z\in L^1(\mu)$, która spełnia
\eqref{eq:CE-def-18}. Jednoznaczność wynika z
jednoznaczności pochodnej Radona--Nikodyma.

Liniowość wynika z liniowości całki po $A$ i
jednoznaczności. Jeśli $Y\ge0$, to $\eta$ jest
miarą nieujemną, a więc jej pochodna $Z$ jest
nieujemna. Jeśli $Y$ jest $\mathcal G$-mierzalna,
sama spełnia \eqref{eq:CE-def-18}, co daje (iii).
Aby dowieść oszacowania $L^1$, korzystamy z
$|\eta|(\Omega)=\int_\Omega|Z|\,\mathrm d\mathbb P$
oraz z faktu, że dla każdego skończonego podziału
$\Omega=\bigsqcup_jA_j$, $A_j\in\mathcal G$,
\[
\sum_j|\eta(A_j)|
\le\sum_j\int_{A_j}|Y|\,\mathrm d\mathbb P
=\int_\Omega|Y|\,\mathrm d\mathbb P.
\]
Biorąc supremum po podziałach otrzymujemy (i).
Dla (iv) niech $A\in\mathcal H$; wówczas także
$A\in\mathcal G$, więc
\[
\int_A\mathbb E[Y\mid\mathcal G] \,\mathrm d\mathbb P
=\int_A Y\,\mathrm d\mathbb P.
\]
Zatem obie strony tożsamości (iv) spełniają
warunek definiujący wartość oczekiwaną względem
$\mathcal H$; jednoznaczność kończy dowód.
\end{proof}

\begin{twierdzenie}[Przypadek $L^2$: rzut ortogonalny]
\label{thm:CE-orthogonal-18}
W założeniach definicji~\ref{def:conditional-exp-18},
jeśli $Y\in L^2(\mathbb P;\mathbb R)$, to
$Z=\mathbb E[Y\mid\mathcal G]\in L^2(\mathbb P)$,
$\|Z\|_2\le\|Y\|_2$ i
\[
\int_\Omega(Y-Z)W\,\mathrm d\mathbb P=0
\qquad\text{dla każdego }W\in L^2(\Omega,\mathcal G,\mathbb P).
\]
Oznacza to, że $Z$ jest rzutem ortogonalnym $Y$
na domkniętą podprzestrzeń funkcji $\mathcal G$-mierzalnych.
\end{twierdzenie}
\begin{proof}
Nie wiemy jeszcze, czy $Z\in L^2$.
Niech $Z_n=\max(-n,\min(Z,n))$; jest to ograniczona
funkcja $\mathcal G$-mierzalna, więc z definicji
warunkowej wartości oczekiwanej (najpierw dla funkcji
prostych, potem przez zbieżność zmajoryzowaną) mamy
$\int_\Omega ZZ_n\,\mathrm d\mathbb P
=\int_\Omega YZ_n\,\mathrm d\mathbb P$.
Ponieważ $ZZ_n\ge Z_n^2$ punktowo,
\[
\|Z_n\|_2^2\le\int_\Omega ZZ_n\,\mathrm d\mathbb P
\le\|Y\|_2\|Z_n\|_2.
\]
Zatem $\|Z_n\|_2\le\|Y\|_2$.
Ponieważ $Z_n^2\uparrow Z^2$, Beppo Levi daje
$\|Z\|_2\le\|Y\|_2$. Definicyjna równość
$\int_A(Y-Z)\,\mathrm d\mathbb P=0$ dla $A\in\mathcal G$
rozciąga się liniowo na funkcje proste i przez
zbieżność zmajoryzowaną na ograniczone funkcje
$\mathcal G$-mierzalne. Trunkacje dowolnego
$W\in L^2(\mathcal G)$ zbliżają się do niego w $L^2$;
Cauchy--Schwarz pozwala przejść do granicy.
Podprzestrzeń $L^2(\mathcal G)$ jest domknięta:
jeśli $W_n$ są $\mathcal G$-mierzalne i zbieżne w $L^2$
do $W$, to wybieramy podciąg $(W_{n_k})$ taki, że
$\|W_{n_k}-W\|_2<2^{-2k}$. Z nierówności
Czebyszewa i lematu Borela--Cantellego wynika
$W_{n_k}\to W$ prawie wszędzie względem $\mathbb P$.
Zbiór $D$ punktów, w których ciąg $(W_{n_k})$
ma granicę skończoną, jest $\mathcal G$-mierzalny,
gdyż warunek Cauchy'ego i skończoność wartości
opisuje się przeliczalnymi działaniami mierzalnymi.
Funkcja $\widetilde W=\lim_kW_{n_k}$ na $D$
i $\widetilde W=0$ poza $D$ jest
$\mathcal G$-mierzalna i równa $W$ prawie wszędzie.
To konstruuje przedstawiciela w $L^2(\mathcal G)$,
bez zakładania zupełności $\mathcal G$.
Równanie ortogonalności identyfikuje $Z$ z rzutem.
\end{proof}

\section{Przykłady obliczania pochodnych}
\begin{przyklad}[Nieskończona gęstość przy punkcie brzegowym]
Na $X=(0,1)$ z miarą Lebesgue'a określmy
$\nu(E)=\int_E x^{-1/2}\,\mathrm d\lambda_1(x)$.
Wtedy $\nu\ll\lambda_1$, a
$\mathrm d\nu/\mathrm d\lambda_1=x^{-1/2}$.
Choć gęstość rośnie bez ograniczenia przy $0$,
jest całkowalna i $\nu(X)=2$.
\end{przyklad}

\begin{przyklad}[Miara osobliwa nie ma pochodnej reprezentującej masę]
Dla miary Cantora $\kappa$ z przykładu
\ref{ex:cantor-measure-17} zachodzi
$\kappa\perp\lambda_1$. Nie istnieje
$f\in L^1(\lambda_1)$ spełniająca
$\kappa(E)=\int_Ef\,\mathrm d\lambda_1(x)$ dla wszystkich
zbiorów borelowskich $E$, bo przy $E=C$ prawa strona
jest równa zero, a $\kappa(C)=1$.
Fakt, że pochodna klasycznej funkcji Cantora wynosi
zero prawie wszędzie, nie oznacza, że jej miara
Stieltjesa jest zerowa. Jest ona osobliwa względem
miary Lebesgue'a.
\end{przyklad}

\section{Mapa implikacji}
Przy $\sigma$-skończoności $\mu,\nu$ uzyskaliśmy
następujący ciąg rezultatów: rozkład Lebesgue'a
$\nu=\nu_{\mathrm{ac}}+\nu_{\mathrm{s}}$,
reprezentację $\nu_{\mathrm{ac}}=f\mu$,
jednoznaczność gęstości, reguły rachunkowe i
konstrukcję warunkowej wartości oczekiwanej.
Przykład~\ref{ex:RN-failure-nonsigma-18} ostrzega,
że sama implikacja między zbiorami zerowymi
nie zastępuje hipotez twierdzenia.

% END chapters/pl/measure_theory/18-rozdzial.tex


% BEGIN chapters/pl/measure_theory/19-rozdzial.tex
% ROZDZIAŁ 19 — zastosowania pochodnej Radona--Nikodyma
\chapter{Zastosowania twierdzenia Radona--Nikodyma}
\label{ch:rn-applications-19}

\section{Gęstość jako opis zmiany miary}
Twierdzenie Radona--Nikodyma z rozdziału~\ref{ch:radon-nikodym-18}
zamienia porównywanie miar w rachunek funkcyjny. Nie wystarczy jednak
bezrefleksyjnie dzielić gęstości: przy każdej zmianie miary należy
wskazać, względem której miary zachodzi równość prawie wszędzie,
a także sprawdzić, czy wyrażenia są całkowalne.

\begin{twierdzenie}[Całkowanie po zmianie miary]
\label{thm:change-measure-19}
Niech $(X,\mathcal F)$ będzie przestrzenią mierzalną,
$\mu,\nu$ nieujemnymi miarami $\sigma$-skończonymi na niej,
$\nu\ll\mu$, a $f=\mathrm d\nu/\mathrm d\mu$.
Dla każdej mierzalnej funkcji $h:X\to[0,\infty]$ zachodzi
\begin{equation}\label{eq:change-measure-19}
 \int_X h(x)\,\mathrm d\nu(x)
 =\int_X h(x)f(x)\,\mathrm d\mu(x).
\end{equation}
Jeśli $h:X\to\mathbb C$ jest mierzalna, to $h\in L^1(\nu)$
wtedy i tylko wtedy, gdy $hf\in L^1(\mu)$; wówczas
\eqref{eq:change-measure-19} pozostaje prawdziwy.
\end{twierdzenie}
\begin{proof}
Dla $h=\chi_A$ równość jest definicją pochodnej.
Przez liniowość otrzymujemy ją dla nieujemnych funkcji prostych,
a następnie przez aproksymację monotoniczną i twierdzenie Beppo
Leviego dla każdej $h\ge0$. Zastosowanie wyniku do $|h|$
daje równoważność całkowalności. Rozłożenie funkcji zespolonej
na dodatnie i ujemne części jej części rzeczywistej i urojonej
kończy dowód.
\end{proof}

\begin{przyklad}[Ważenie miary]
Niech $\mu=\lambda_1|_{(0,1)}$ i $f(x)=2x$.
Wówczas $\nu(E)=\int_E 2x\,\mathrm d\lambda_1(x)$ jest
miarą probabilistyczną. Dla $h(x)=x^2$ otrzymujemy
$\int_{(0,1)}h\,\mathrm d\nu(x)
=\int_{(0,1)}2x^3\,\mathrm d\lambda_1(x)=1/2$.
Z kolei $\mathrm d\mu/\mathrm d\nu=(2x)^{-1}$
$\nu$-prawie wszędzie. Mimo nieograniczoności tej pochodnej
przy zerze zachodzi
$\int_{(0,1)}(2x)^{-1}\,\mathrm d\nu(x)=1$.
\end{przyklad}

\section{Miary obrazowe oraz gęstości obserwacji}
Dla mierzalnego $T:X\to Y$ i miary $\mu$ na $(X,\mathcal F)$
przypominamy miarę obrazową
$T_*\mu(B)=\mu(T^{-1}(B))$, $B\in\mathcal G$,
gdzie $(Y,\mathcal G)$ jest przestrzenią mierzalną.
Nie wymagamy, aby $T$ była odwracalna.

\begin{twierdzenie}[Bezwzględna ciągłość zachowana przez obraz]
\label{thm:image-ac-19}
Niech $T:(X,\mathcal F)\to(Y,\mathcal G)$ będzie mierzalne,
a $\mu,\nu$ miarami nieujemnymi na $\mathcal F$.
Jeżeli $\nu\ll\mu$, to $T_*\nu\ll T_*\mu$.
Jeżeli ponadto $\mu,\nu$ są miarami probabilistycznymi,
$f=\mathrm d\nu/\mathrm d\mu$ i
$g=\mathrm d(T_*\nu)/\mathrm d(T_*\mu)$,
to $g\circ T$ jest wersją warunkowej wartości oczekiwanej
$\mathbb E_\mu[f\mid\sigma(T)]$, gdzie
$\sigma(T)=\{T^{-1}(B):B\in\mathcal G\}$.
\end{twierdzenie}
\begin{proof}
Jeżeli $(T_*\mu)(B)=0$, to $\mu(T^{-1}(B))=0$,
zatem $\nu(T^{-1}(B))=0$, co jest równoważne
$(T_*\nu)(B)=0$. Obie miary obrazowe są skończone,
więc pochodna $g$ istnieje. Dla każdego $B\in\mathcal G$
z twierdzenia o całkowaniu względem miary obrazowej
(definicję miary obrazowej z rozdziału~\ref{ch:struktura-przestrzeni})
uzyskujemy
\[
\int_{T^{-1}(B)}g(T(x))\,\mathrm d\mu(x)
=\int_Bg(y)\,\mathrm d(T_*\mu)(y)
=(T_*\nu)(B)
=\int_{T^{-1}(B)}f(x)\,\mathrm d\mu(x).
\]
Funkcja $g\circ T$ jest $\sigma(T)$-mierzalna. Jest również
całkowalna, ponieważ
\[
\int_X|g(T(x))|\,\mathrm d\mu(x)
=\int_Y|g(y)|\,\mathrm d(T_*\mu)(y)<\infty.
\]
Spełnia więc warunek charakteryzujący warunkową wartość oczekiwaną
z definicji~\ref{def:conditional-exp-18}.
\end{proof}

\begin{uwaga}[Kiedy informacja znika?]
W ogólności $g(T(x))\neq f(x)$. Po zastosowaniu
odwzorowania $T$ pozostaje tylko średnia warunkowa $f$
względem informacji dostępnej poprzez $T$. Jeśli $T$
jest stałe, a $\mu,\nu$ są miarami probabilistycznymi,
to $T_*\mu=T_*\nu$, choć $\mu$ i $\nu$ mogą być różne.
Nie wolno też z samego $\nu\ll\mu$ wnioskować, że dowolne
odwzorowanie $T$ zachowuje osobliwość miar.
\end{uwaga}

\section{Wariacja całkowita a testowanie różnic między rozkładami}
Dla miar probabilistycznych $P,Q$ na $(X,\mathcal F)$
wprowadzamy odległość
\begin{equation}\label{eq:tv-prob-19}
d_{\mathrm{TV}}(P,Q):=\sup_{A\in\mathcal F}|P(A)-Q(A)|.
\end{equation}
Należy odróżnić ją od normy wariacji całkowitej miary ze znakiem
$P-Q$, oznaczonej w rozdziale~\ref{ch:miary-ze-znakiem} przez
$\|P-Q\|_{\mathrm{TV}}=|P-Q|(X)$.

\begin{twierdzenie}[Formuła wariacyjna dla dwóch prawdopodobieństw]
\label{thm:total-variation-formula-19}
Niech $P,Q$ będą miarami probabilistycznymi na $(X,\mathcal F)$.
Niech $\lambda=P+Q$, $p=\mathrm dP/\mathrm d\lambda$
i $q=\mathrm dQ/\mathrm d\lambda$. Wówczas
\begin{equation}\label{eq:tv-formula-19}
 d_{\mathrm{TV}}(P,Q)
 =\frac12\int_X|p(x)-q(x)|\,\mathrm d\lambda(x)
 =\frac12\|P-Q\|_{\mathrm{TV}}
 =1-\int_X\min(p,q)\,\mathrm d\lambda.
\end{equation}
Supremum z \eqref{eq:tv-prob-19} jest osiągane na zbiorze
$A_+=\{p>q\}$. W szczególności $0\le d_{\mathrm{TV}}(P,Q)\le1$,
a równość $d_{\mathrm{TV}}(P,Q)=1$ zachodzi dokładnie wtedy,
gdy $P\perp Q$.
\end{twierdzenie}
\begin{proof}
Miara ze znakiem $\eta=P-Q$ ma gęstość $p-q$
względem $\lambda$. Zatem
$\eta(A)=\int_A(p-q)\,\mathrm d\lambda$ i
$\eta(X)=0$. Rozkład na części dodatnią i ujemną
pokazuje, że $\eta^+(X)=\eta^-(X)=|\eta|(X)/2$.
Dla dowolnego $A$ mamy
$-\eta^-(X)\le\eta(A)\le\eta^+(X)$,
a $\eta(A_+)=\eta^+(X)$; stąd pierwsze dwie równości.
Ponieważ $\int_X(p+q)\,\mathrm d\lambda=2$ i punktowo
$|p-q|=p+q-2\min(p,q)$, dostajemy ostatnią równość.
Jeżeli $d_{\mathrm{TV}}=1$, to $\min(p,q)=0$
$\lambda$-prawie wszędzie, czyli $p q=0$ prawie wszędzie.
Zbiory $\{p>0\}$ i $\{q>0\}$ są wówczas rozłączne
z dokładnością do zbioru $\lambda$-zerowego i skupiają
odpowiednio $P$ i $Q$; stąd $P\perp Q$.
Implikacja odwrotna wynika bezpośrednio z definicji osobliwości.
\end{proof}

\begin{wniosek}[Przetwarzanie danych nie zwiększa wariacji całkowitej]
\label{cor:tv-data-processing-19}
W założeniach twierdzenia~\ref{thm:image-ac-19},
bez wymogu $\nu\ll\mu$, jeśli $P,Q$ są prawdopodobieństwami,
to
$ d_{\mathrm{TV}}(T_*P,T_*Q)\le d_{\mathrm{TV}}(P,Q)$.
\end{wniosek}
\begin{proof}
Dla dowolnego $B\in\mathcal G$ zachodzi
$|(T_*P)(B)-(T_*Q)(B)|
=|P(T^{-1}B)-Q(T^{-1}B)|\le d_{\mathrm{TV}}(P,Q)$.
Bierzemy supremum po $B$.
\end{proof}

\begin{przyklad}[Zbieżność słaba nie wymusza zbieżności w wariacji]
Na $\mathbb R$ rozważmy $P_n=\delta_{1/n}$ i $P=\delta_0$.
Dla każdej funkcji ciągłej i ograniczonej $h$
zachodzi $\int_{\mathbb R}h\,\mathrm dP_n\to
\int_{\mathbb R}h\,\mathrm dP$, co jest zbieżnością wąską.
Jednak $d_{\mathrm{TV}}(P_n,P)=1$ dla każdego $n$,
ponieważ $P_n(\{1/n\})=1$ i $P(\{1/n\})=0$.
Precyzyjną teorię zbieżności miar rozwiniemy w rozdziale~22.
\end{przyklad}

\section{Lemat Neymana--Pearsona: optymalność testu progowego}
Przedstawimy rezultat stanowiący użyteczne, a zarazem nieintuicyjne
zastosowanie pochodnej Radona--Nikodyma. Z dwóch prostych hipotez
można skonstruować test optymalny bez żadnych założeń o istnieniu
gęstości względem miary Lebesgue'a.

\begin{definicja}[Test losowy i jego moc]
\label{def:test-19}
Niech $P,Q$ będą prawdopodobieństwami na $(X,\mathcal F)$.
Testem losowym jest mierzalna funkcja $\varphi:X\to[0,1]$.
Liczbę $\int_X\varphi\,\mathrm dP$ nazywamy
prawdopodobieństwem odrzucenia hipotezy $P$ pod $P$,
a $\int_X\varphi\,\mathrm dQ$ jego mocą pod $Q$.
Test ma poziom nieprzekraczający $\alpha\in[0,1]$,
gdy pierwsza z tych liczb nie przekracza $\alpha$.
\end{definicja}

\begin{twierdzenie}[Lemat Neymana--Pearsona]
\label{thm:neyman-pearson-19}
Niech $P,Q$ będą prawdopodobieństwami na $(X,\mathcal F)$,
$\lambda=P+Q$, $p=\mathrm dP/\mathrm d\lambda$
i $q=\mathrm dQ/\mathrm d\lambda$.
Dla $c\ge0$ i $\gamma\in[0,1]$ ustalmy test
\[
\varphi_*=\chi_{\{q>cp\}}+
 \gamma\chi_{\{q=cp\}}.
\]
Niech $\alpha=\int_X\varphi_*\,\mathrm dP$.
Wtedy dla każdego testu $\varphi:X\to[0,1]$ o poziomie
nieprzekraczającym $\alpha$ zachodzi
\[
 \int_X\varphi\,\mathrm dQ\le
 \int_X\varphi_*\,\mathrm dQ.
\]
Dla $c>0$ dowód nie wymaga założenia, że
$P$ i $Q$ są wzajemnie bezwzględnie ciągłe.
W razie $c=0$ twierdzenie pozostaje prawdziwe;
nie wymaga istnienia gęstości ani $P$, ani $Q$ względem
żadnej wcześniej zadanej miary odniesienia.
\end{twierdzenie}
\begin{proof}
Na zbiorze $\{q>cp\}$ mamy $\varphi_*=1$, więc
$\varphi-\varphi_*\le0$; na $\{q<cp\}$ mamy
$\varphi_*=0$, więc $\varphi-\varphi_*\ge0$.
Na zbiorze $\{q=cp\}$ iloczyn jest zerem. W każdym
przypadku zatem
$(\varphi-\varphi_*)(q-cp)\le0$ punktowo.
Całkujemy względem $\lambda$ i dostajemy
\[
\int_X(\varphi-\varphi_*)\,\mathrm dQ
\le c\int_X(\varphi-\varphi_*)\,\mathrm dP\le0,
\]
gdzie ostatnia nierówność wynika z poziomu testu
i nieujemności $c$. To dokładnie teza.
\end{proof}

\begin{uwaga}[Rola losowania na granicy]
Wprowadzony parametr $\gamma$ pozwala kontrolować poziom
testu, gdy zbiór $\{q=cp\}$ ma dodatnią miarę $P$.
W ogólności nie każdy poziom $\alpha$ da się otrzymać
bez losowania. Na zbiorze $\{p=0,q>0\}$ test progowy
zawsze ma wartość $1$: odrzucenie $P$ nic tam nie kosztuje.
\end{uwaga}

\begin{wniosek}[Konstrukcja testu optymalnego o zadanym poziomie]
\label{cor:NP-exact-alpha-19}
Niech $P,Q,\lambda,p,q$ mają znaczenie z twierdzenia
\ref{thm:neyman-pearson-19}. Dla każdego
$\alpha\in(0,1)$ istnieją $c\ge0$ i
$\gamma\in[0,1]$, dla których test
$\varphi_*=\chi_{\{q>cp\}}+\gamma\chi_{\{q=cp\}}$
ma poziom dokładnie $\alpha$, a spośród wszystkich
testów poziomu co najwyżej $\alpha$ ma największą moc.
Dla $\alpha=0$ testem optymalnym jest
$\chi_{\{p=0,\,q>0\}}$, a dla $\alpha=1$
można przyjąć test stały równy $1$.
\end{wniosek}
\begin{proof}
Niech $L=q/p$ na $\{p>0\}$ i $L=0$ na $\{p=0\}$.
Ponieważ $P(\{p=0\})=0$, mamy $L<\infty$
$P$-prawie wszędzie. Dla $0<\alpha<1$ kładziemy
\[
 c=\inf\{t\ge0:P(L>t)\le\alpha\}.
\]
Z nierówności $P(L>t)\downarrow0$ przy $t\uparrow\infty$
wynika, że zbiór pod infimum jest niepusty, a $c$ jest
skończone. Dla $t_n\downarrow c$ i $t_n>c$ z ciągłości
miary z dołu otrzymujemy $P(L>c)\le\alpha$.
Jeśli $c>0$, to z definicji infimum
$P(L>t)>\alpha$ dla $0\le t<c$;
przechodząc z $t\uparrow c$ i używając ciągłości miary
z góry, dostajemy $P(L\ge c)\ge\alpha$.
Gdy $c=0$, ostatnia nierówność wynika z
$P(L\ge0)=1$.
W konsekwencji
\[
 P(L>c)\le\alpha\le P(L\ge c).
\]
Jeśli $P(L=c)>0$, przyjmujemy
$\gamma=(\alpha-P(L>c))/P(L=c)\in[0,1]$;
w przeciwnym przypadku nierówności dają
$P(L>c)=\alpha$ i wybieramy $\gamma=0$.
Zbiory $\{L>c\}$, $\{L=c\}$ zgadzają się
$P$-prawie wszędzie z $\{q>cp\}$ i $\{q=cp\}$,
więc $\int_X\varphi_*(x)\,\mathrm dP(x)=\alpha$.
Optymalność wynika z twierdzenia
\ref{thm:neyman-pearson-19}.

Dla $\alpha=0$ dopuszczalny test $\varphi$ jest zerowy
$P$-prawie wszędzie, a więc $\varphi=0$
$\lambda$-prawie wszędzie na $\{p>0\}$.
Stąd jego moc nie przekracza
$Q(\{p=0\})$, osiąganej przez wskazany test.
Przypadek $\alpha=1$ jest natychmiastowy.
\end{proof}

\section{Wzór Bayesa jako twierdzenie o gęstościach}
Wynik sformułujemy bez założenia, że zmienna obserwowana
posiada klasyczną gęstość na $\mathbb R^d$.

\begin{twierdzenie}[Bayes dla skończonej przestrzeni parametrów]
\label{thm:bayes-finite-19}
Niech $m\ge1$, $\pi_i>0$, $\sum_{i=1}^m\pi_i=1$.
Niech $(Y,\mathcal G)$ będzie przestrzenią mierzalną,
a $P_1,\ldots,P_m$ miarami probabilistycznymi na $Y$.
Na przestrzeni $\{1,\ldots,m\}\times Y$ definiujemy
prawdopodobieństwo wspólne przez
$\mathbb P(\{i\}\times B)=\pi_iP_i(B)$.
Niech $\lambda=\sum_iP_i$, $p_i=\mathrm dP_i/\mathrm d\lambda$,
a $r=\sum_i\pi_i p_i$. Dla $r(y)>0$ określmy
\[
 \beta_i(y)=\frac{\pi_i p_i(y)}{r(y)}.
\]
Na zbiorze $\{r=0\}$ możemy przyjąć dowolne mierzalne
wartości $\beta_i\in[0,1]$. Wtedy dla obserwacji
$Y(i,y)=y$ i etykiety $\Theta(i,y)=i$ zachodzi
\begin{equation}\label{eq:bayes-finite-19}
 \beta_i(Y)=\mathbb P(\Theta=i\mid\sigma(Y))
 \quad\mathbb P\text{-prawie wszędzie}.
\end{equation}
Ponadto $\sum_i\beta_i=1$ $P_Y$-prawie wszędzie,
gdzie $P_Y=\sum_i\pi_iP_i$ jest prawem obserwacji.
\end{twierdzenie}
\begin{proof}
Dla dowolnego $B\in\mathcal G$ definicja rozkładu
wspólnego daje
$\mathbb P(\Theta=i,\,Y\in B)=\pi_iP_i(B)
=\int_B\pi_ip_i(y)\,\mathrm d\lambda(y)$.
Tymczasem $P_Y(B)=\int_Br(y)\,\mathrm d\lambda(y)$,
więc stosując twierdzenie~\ref{thm:change-measure-19}
uzyskujemy
\[
\int_{Y^{-1}(B)}\beta_i(Y)\,\mathrm d\mathbb P
=\int_B\beta_i(y)r(y)\,\mathrm d\lambda(y)
=\pi_iP_i(B).
\]
Zbiór $\{r=0\}$ ma miarę $P_Y$ równą zero, więc nie
wpływa na wzory. Funkcja $\beta_i(Y)$ jest
$\sigma(Y)$-mierzalna i ograniczona. Jest zatem
warunkową wartością oczekiwaną funkcji charakterystycznej
$\chi_{\{\Theta=i\}}$ względem $\sigma(Y)$.
Równość sumy wynika z $\sum_i\pi_ip_i=r$ na $\{r>0\}$.
\end{proof}

\begin{przyklad}[Dwie obserwacje Bernoulliego]
Niech $\Theta\in\{0,1\}$ z prawdopodobieństwami
$\mathbb P(\Theta=1)=\pi$ i $\mathbb P(\Theta=0)=1-\pi$.
Warunkowo względem $\Theta=i$ zaobserwujmy
niezależną parę $(Y_1,Y_2)$ zmiennych Bernoulliego
o parametrze $\theta_i\in(0,1)$. Dla obserwacji
$(Y_1,Y_2)=(1,0)$ wzór Bayesa daje
\[
\mathbb P(\Theta=1\mid Y_1=1,Y_2=0)
=\frac{\pi\theta_1(1-\theta_1)}
{\pi\theta_1(1-\theta_1)+(1-\pi)\theta_0(1-\theta_0)}.
\]
Nie potrzeba tu miary Lebesgue'a: wspólną miarą
dominującą jest miara licząca na czteroelementowym zbiorze.
\end{przyklad}

\section{Warunkowa wartość oczekiwana i nierówności}
Warunkową wartość oczekiwaną zdefiniowaliśmy w
rozdziale~\ref{ch:radon-nikodym-18}. Teraz dowiedziemy
jej własności przydatnych w analizie oraz prawdopodobieństwie.

\begin{lemat}[Wyciąganie czynnika mierzalnego]
\label{lem:conditional-pullout-19}
Niech $(\Omega,\mathcal F,\mathbb P)$ będzie przestrzenią
probabilistyczną, $\mathcal G\subseteq\mathcal F$,
$Y\in L^1(\mathbb P)$ oraz $H$ ograniczoną rzeczywistą
funkcją $\mathcal G$-mierzalną. Wówczas
\[
 \mathbb E[HY\mid\mathcal G]
 =H\mathbb E[Y\mid\mathcal G]
 \quad\mathbb P\text{-prawie wszędzie}.
\]
\end{lemat}
\begin{proof}
Dla $H=\chi_B$, $B\in\mathcal G$, tożsamość wynika
z warunku definiującego warunkową wartość oczekiwaną,
ponieważ dla dowolnego $A\in\mathcal G$ całkujemy po
$A\cap B$. Z liniowości wynika przypadek prostych $H$.
Dla ograniczonego $H$ istnieją proste $H_n$, mierzalne
względem $\mathcal G$, takie, że $H_n\to H$ punktowo
i $|H_n|\le\|H\|_\infty+1$. Przez zbieżność zmajoryzowaną
$H_nY\to HY$ i $H_n\mathbb E[Y\mid\mathcal G]
\to H\mathbb E[Y\mid\mathcal G]$ w $L^1$.
Całkując po każdym $A\in\mathcal G$, przechodzimy
do granicy; jednoznaczność daje tezę.
\end{proof}

\begin{twierdzenie}[Warunkowa nierówność Jensena]
\label{thm:conditional-jensen-19}
Niech $(\Omega,\mathcal F,\mathbb P)$ będzie przestrzenią
probabilistyczną, $\mathcal G\subseteq\mathcal F$,
$Y\in L^1(\mathbb P;\mathbb R)$, a
$\Phi:\mathbb R\to\mathbb R$ funkcją wypukłą, wszędzie
skończoną. Załóżmy ponadto, że $\Phi(Y)\in L^1(\mathbb P)$.
Wówczas
\[
 \Phi\bigl(\mathbb E[Y\mid\mathcal G]\bigr)
 \le\mathbb E[\Phi(Y)\mid\mathcal G]
 \quad\mathbb P\text{-prawie wszędzie}.
\]
\end{twierdzenie}
\begin{proof}
Dla każdego $t\in\mathbb Q$ wybieramy liczbę
$a_t\in\partial\Phi(t)$, gdzie
$\partial\Phi(t)$ oznacza zbiór nachyleń prostych
podpierających wykres $\Phi$ w punkcie $t$.
Istnienie takiego nachylenia wynika z monotoniczności
ilorazów różnicowych funkcji wypukłej; możemy np. przyjąć
prawostronną pochodną $\Phi'_+(t)$, która jest skończona.
Wówczas
$\Phi(x)\ge a_tx+\Phi(t)-a_tt$ dla każdego $x\in\mathbb R$.
Różnica
$\Phi(Y)-a_tY-\Phi(t)+a_tt$ jest nieujemna,
więc dodatniość i liniowość warunkowej wartości oczekiwanej
dają
\[
\mathbb E[\Phi(Y)\mid\mathcal G]
\ge a_t\mathbb E[Y\mid\mathcal G]+\Phi(t)-a_tt.
\]
Ponieważ zbiór liczb wymiernych jest przeliczalny,
można wybrać jeden wspólny zbiór pełnej miary, na którym
wszystkie nierówności obowiązują jednocześnie.
Ciągłość $\Phi$ i jej lokalna lipschitzowskość
wynikają z wypukłości i skończoności na całej prostej.
Pozwala to przejść z $t\in\mathbb Q$ do dowolnego $t$:
$\Phi(x)=\sup_{t\in\mathbb Q}
(a_tx+\Phi(t)-a_tt)$.
Podstawiając $x=\mathbb E[Y\mid\mathcal G]$
i biorąc supremum otrzymujemy tezę.
\end{proof}

\begin{wniosek}[Kontrakcja w $L^p$]
\label{cor:conditional-lp-19}
Jeśli $1\le p<\infty$, $Y\in L^p(\mathbb P;\mathbb R)$,
to $\mathbb E[Y\mid\mathcal G]\in L^p(\mathbb P)$ oraz
$\|\mathbb E[Y\mid\mathcal G]\|_p\le\|Y\|_p$.
Dla $p=\infty$ zachodzi analogiczne oszacowanie
normy istotnego supremum.
\end{wniosek}
\begin{proof}
Dla $1\le p<\infty$ stosujemy twierdzenie
\ref{thm:conditional-jensen-19} do $\Phi(t)=|t|^p$.
Całkując otrzymaną nierówność względem $\mathbb P$
i korzystając z własności zachowania całki warunkowej,
dostajemy oszacowanie normy. Jeżeli $|Y|\le M$ prawie
wszędzie, to z monotoniczności
$-M\le\mathbb E[Y\mid\mathcal G]\le M$,
co daje przypadek $p=\infty$.
\end{proof}

\section{Informacja ilościowa: entropia względna}
Dla prawdopodobieństw $Q\ll P$ określamy
\emph{entropię względną} (dywergencję Kullbacka--Leiblera)
\[
 D(Q\Vert P):=\int_X f(x)\log f(x)\,\mathrm dP(x),
 \qquad f=\frac{\mathrm dQ}{\mathrm dP},
\]
przy konwencji $0\log0=0$. Całka ma wartości
w $[0,\infty]$: z $t\log t\ge -e^{-1}$
wynika, że dla prawdopodobieństwa $P$ część ujemna
całkowanej funkcji jest całkowalna. Gdy $Q\not\ll P$,
przyjmujemy $D(Q\Vert P)=+\infty$.

\begin{twierdzenie}[Nieujemność entropii względnej]
\label{thm:gibbs-19}
Dla dowolnych prawdopodobieństw $P,Q$ na $(X,\mathcal F)$
zachodzi $D(Q\Vert P)\ge0$. Jeżeli $Q\ll P$, to
$D(Q\Vert P)=0$ wtedy i tylko wtedy, gdy $Q=P$.
\end{twierdzenie}
\begin{proof}
Dla $t\ge0$ funkcja $\psi(t)=t\log t-t+1$
(z $\psi(0)=1$) jest nieujemna i zeruje się wyłącznie
w $t=1$: na $(0,\infty)$ pochodna jest równa $\log t$,
więc minimum jest w $1$. Jeśli $Q\ll P$ i $f=\mathrm dQ/\mathrm dP$,
to $\int_Xf\,\mathrm dP=1$, a zatem
\[
 D(Q\Vert P)=\int_X\psi(f(x))\,\mathrm dP(x)\ge0.
\]
Jeśli całka jest zerowa, to $\psi(f)=0$ $P$-prawie wszędzie,
a więc $f=1$ $P$-prawie wszędzie, czyli $Q=P$.
W pozostałym przypadku definicja daje nieskończoność.
\end{proof}

\begin{twierdzenie}[Zasada przetwarzania danych dla entropii]
\label{thm:kl-data-processing-19}
Niech $P,Q$ będą prawdopodobieństwami na $(X,\mathcal F)$,
$Q\ll P$, a $T:(X,\mathcal F)\to(Y,\mathcal G)$
odwzorowaniem mierzalnym. Wówczas
\[
 D(T_*Q\Vert T_*P)\le D(Q\Vert P).
\]
\end{twierdzenie}
\begin{proof}
Niech $f=\mathrm dQ/\mathrm dP$, a
$g=\mathrm d(T_*Q)/\mathrm d(T_*P)$.
Z twierdzenia~\ref{thm:image-ac-19} wynika
$g(T)=\mathbb E_P[f\mid\sigma(T)]$ prawie wszędzie.
Dla funkcji wypukłej
$\psi(t)=t\log t-t+1$ na $[0,\infty)$
warunkowa nierówność Jensena pozostaje prawdziwa
również wtedy, gdy $\psi(f)$ nie jest całkowalna:
stosujemy ją do funkcji nieujemnej $\psi$ po odjęciu
wspierającej funkcji afinicznej, lub aproksymujemy
$\psi$ supremum przeliczalnej rodziny prostych
podpierających. Wówczas
\[
 \psi(g(T))\le\mathbb E_P[\psi(f)\mid\sigma(T)]
\]
w rozumieniu rozszerzonej warunkowej wartości oczekiwanej
funkcji nieujemnej. Jeżeli $D(Q\Vert P)=\infty$,
nierówność jest natychmiastowa. W przypadku skończonym
$\psi(f)\in L^1(P)$, a powyższą nierówność można
uzyskać wprost z dowodu twierdzenia
\ref{thm:conditional-jensen-19}, z rozszerzeniem na funkcję
$\psi$ na $[0,\infty)$; nieujemność $f$ i $g$
pozwala używać wyłącznie podpór w punktach dodatnich.
Całkując względem $P$ i stosując definicję miary obrazowej,
otrzymujemy
\[
 D(T_*Q\Vert T_*P)
=\int_Y\psi(g(y))\,\mathrm d(T_*P)(y)
\le\int_X\psi(f(x))\,\mathrm dP(x)=D(Q\Vert P).
\]
\end{proof}

\begin{uwaga}[Interpretacja]
Wzory dla wariacji całkowitej i entropii względnej
pokazują dwa różne sensy zasady, że samo przekształcenie
obserwacji nie może zwiększyć informacji o tym,
z którego z dwóch rozkładów pochodzi wynik. Są to
nierówności matematyczne; nie wymagają aparatu teorii
informacji ani wcześniejszej znajomości statystyki.
\end{uwaga}

\section{Ograniczenia i dalsze kierunki}
Pochodna Radona--Nikodyma jest określona prawie wszędzie,
a nie punktowo kanonicznie. Wartość $f(x)$ w danym punkcie
nie musi być znacząca; w zastosowaniach liczą się całki,
zbiory poziomicowe z dokładnością do zbiorów zerowych
oraz odpowiednie klasy funkcji. Warunkowe rozkłady
jako rodziny miar $P(\,\cdot\mid Y=y)$ wymagają dodatkowych
założeń, np. standardowej borelowskości przestrzeni,
czym zajmiemy się w module o dezintegracji.

% END chapters/pl/measure_theory/19-rozdzial.tex


% BEGIN chapters/pl/measure_theory/20-rozdzial.tex
% ROZDZIAŁ 20 — miary Radona i ich regularność
\chapter{Miary borelowskie, regularność i miary Radona}
\label{ch:radon-measures-20}

\section{Dlaczego miara potrzebuje topologii?}
Do tej pory wszystkie miary były określone na przestrzeniach
mierzalnych. Jeśli jednak przestrzeń ma topologię, można pytać,
czy wartości miary zbioru dają się odtworzyć z miar zbiorów
otwartych, domkniętych lub zwartych. Są to trzy różne własności:
regularność zewnętrzna, regularność wewnętrzna i ciasność.
Rozróżnianie ich będzie niezbędne w twierdzeniu reprezentacyjnym
Riesza--Markowa--Kakutaniego.

\begin{definicja}[Miara borelowska i lokalna skończoność]
\label{def:borel-locfinite-20}
Niech $(X,\tau)$ będzie przestrzenią topologiczną,
a $\mathcal B(X)$ jej $\sigma$-algebrą borelowską.
\emph{Miarą borelowską} na $X$ nazywamy nieujemną miarę
$\mu:\mathcal B(X)\to[0,\infty]$.
Miara $\mu$ jest \emph{lokalnie skończona}, jeśli dla
każdego $x\in X$ istnieje otwarte otoczenie $U_x$ punktu $x$
takie, że $\mu(U_x)<\infty$.
\end{definicja}

\begin{uwaga}[Skończoność na zbiorach zwartych]
\label{rem:compact-finiteness-20}
Na dowolnej przestrzeni topologicznej lokalna skończoność
implikuje, że $\mu(K)<\infty$ dla każdego zwartego
$K\subseteq X$: ze skończonego podpokrycia
$K\subseteq U_{x_1}\cup\cdots\cup U_{x_m}$ dostajemy
$\mu(K)\le\sum_i\mu(U_{x_i})<\infty$.
Na przestrzeni lokalnie zwartej Hausdorffa zachodzi
również implikacja odwrotna: każdy punkt posiada otwarte
otoczenie o zwartym domknięciu, więc jego miara nie
przekracza skończonej miary tego domknięcia.
Bez lokalnej zwartości należy rozróżniać te dwa warunki;
sama skończoność miary na zwartych zbiorach nie daje
jeszcze bezpośrednio otoczeń o skończonej mierze.
\end{uwaga}

\begin{definicja}[Regularność względem topologii]
\label{def:regularity-20}
Niech $X$ będzie przestrzenią Hausdorffa, a $\mu$ miarą
borelowską na $X$. Mówimy, że jest
\emph{zewnętrznie regularna}, gdy dla każdego
$E\in\mathcal B(X)$ zachodzi
\begin{equation}\label{eq:outer-regular-20}
 \mu(E)=\inf\{\mu(U):E\subseteq U,\ U\subseteq X
 \text{ otwarty}\}.
\end{equation}
Mówimy, że jest \emph{wewnętrznie regularna względem
zbiorów domkniętych}, gdy dla każdego $E\in\mathcal B(X)$
\[
 \mu(E)=\sup\{\mu(F):F\subseteq E,\ F\subseteq X
 \text{ domknięty}\}.
\]
Jest \emph{wewnętrznie regularna względem zbiorów zwartych},
gdy dla każdego $E\in\mathcal B(X)$
\begin{equation}\label{eq:inner-compact-20}
 \mu(E)=\sup\{\mu(K):K\subseteq E,\ K\subseteq X
 \text{ zwarty}\}.
\end{equation}
Przyjmujemy $\inf\varnothing=+\infty$ i
$\sup\varnothing=0$ dla podanych rodzin miar nieujemnych.
\end{definicja}

\begin{definicja}[Miara Radona --- przyjęta konwencja]
\label{def:radon-measure-20}
Niech $X$ będzie lokalnie zwartą przestrzenią Hausdorffa.
\emph{Miarą Radona} na $X$ nazywamy miarę borelowską $\mu$,
która jest skończona na zbiorach zwartych, regularna
zewnętrznie na wszystkich zbiorach borelowskich oraz
\emph{regularna wewnętrznie na zbiorach otwartych}, to znaczy
\[
 \mu(U)=\sup\{\mu(K): K\subseteq U,\ K\text{ zwarty}\}
 \qquad(U\subseteq X\text{ otwarty}).
\]
Skończoność na kompaktach jest na $X$ równoważna
lokalnej skończoności (uwaga~\ref{rem:compact-finiteness-20}).
\end{definicja}

\begin{uwaga}[Dwie konwencje regularności]
W definicji~\ref{def:radon-measure-20} regularność
wewnętrzną wymagamy na zbiorach otwartych, a nie na wszystkich
borelowskich. To rozróżnienie jest konieczne na ogólnych
przestrzeniach lokalnie zwartych Hausdorffa. Silniejsza
regularność na wszystkich zbiorach borelowskich zachodzi,
jeśli przestrzeń jest $\sigma$-zwarta lub jeśli rozważany
zbiór ma miarę $\sigma$-skończoną.
W rozdziale~\ref{ch:haar-23} podamy jawny przykład
miary Haara, dla której silniejszy warunek zawodzi.
Na lokalnie zwartych przestrzeniach metrycznych ośrodkowych
obie konwencje są równoważne.
\end{uwaga}

\begin{lemat}[Regularność na zbiorach $\sigma$-skończonych]
\label{lem:radon-sigmafinite-inner-20}
Niech $X$ będzie lokalnie zwartą przestrzenią Hausdorffa,
a $\mu$ miarą Radona według definicji
\ref{def:radon-measure-20}. Jeżeli $E\in\mathcal B(X)$
i $\mu(E)<\infty$, to
\[
 \mu(E)=\sup\{\mu(K):K\subseteq E,\ K\text{ zwarty}\}.
\]
Równość ta zachodzi także dla każdego borelowskiego
$E$ o mierze $\sigma$-skończonej.
\end{lemat}
\begin{proof}
Niech $\mu(E)<\infty$ i $\varepsilon>0$.
Z regularności zewnętrznej wybieramy otwarty
$U\supseteq E$ z $\mu(U)<\mu(E)+\varepsilon$.
Mamy $\mu(U\setminus E)<\varepsilon$.
Następnie wybieramy otwarty $W\supseteq U\setminus E$
z $\mu(W)<\mu(U\setminus E)+\varepsilon<2\varepsilon$,
a z regularności wewnętrznej na otwartym $U$ kompakt
$F\subseteq U$ z $\mu(F)>\mu(U)-\varepsilon$.
Wówczas $K=F\setminus W$ jest zwartym podzbiorem $E$ oraz
$\mu(K)\ge\mu(F)-\mu(W)>\mu(E)-3\varepsilon$.
Dla zbioru $\sigma$-skończonego przedstawiamy go jako
rosnącą przeliczalną sumę borelowskich zbiorów
$E_j$ o miarach skończonych, po czym stosujemy
ciągłość miary z dołu i pierwszy przypadek.
\end{proof}

\section{Regularność skończonych miar borelowskich na przestrzeniach metrycznych}
\begin{twierdzenie}[Aproksymacja zbiorów borelowskich]
\label{thm:finite-metric-regular-20}
Niech $(X,d)$ będzie dowolną przestrzenią metryczną,
a $\mu$ skończoną miarą borelowską na $X$.
Dla każdego $E\in\mathcal B(X)$ i $\varepsilon>0$
istnieją zbiór domknięty $F$ i zbiór otwarty $U$ takie,
że $F\subseteq E\subseteq U$ oraz
\begin{equation}\label{eq:finite-regular-error-20}
 \mu(U\setminus F)<\varepsilon.
\end{equation}
W szczególności $\mu$ jest zewnętrznie regularna
i wewnętrznie regularna względem zbiorów domkniętych.
Nie twierdzimy w tym miejscu, że $F$ można wybrać zwarty.
\end{twierdzenie}
\begin{proof}
\emph{Krok 1: zbiory otwarte.}
Dla zbioru otwartego $U\neq X$ definiujemy
$F_n=\{x\in X:d(x,X\setminus U)\ge1/n\}$.
Funkcja $x\mapsto d(x,X\setminus U)$ jest
$1$-lipschitzowska, więc $F_n$ jest domknięty.
Z otwartości $U$ wynika, że każdy $x\in U$
ma dodatnią odległość od $X\setminus U$;
stąd $F_n\uparrow U$. Ponieważ $\mu(U)<\infty$,
ciągłość miary z dołu daje
$\mu(U\setminus F_n)\to0$. Jeżeli $U=X$,
wystarczy przyjąć $F_n=X$. Zatem każdy zbiór otwarty
spełnia warunek z tezy dla dowolnego $\varepsilon>0$.

\emph{Krok 2: rodzina zbiorów aproksymowalnych.}
Oznaczmy przez $\mathcal D$ rodzinę tych zbiorów
$E\in\mathcal B(X)$, dla których warunek tezy zachodzi
dla każdego $\varepsilon>0$. Rodzina ta zawiera wszystkie
zbiory otwarte. Wykażemy, że jest $\sigma$-algebrą.
Jeżeli $F\subseteq E\subseteq U$ i
$\mu(U\setminus F)<\varepsilon$, to
$X\setminus U\subseteq X\setminus E\subseteq X\setminus F$,
a zbiór $X\setminus U$ jest domknięty,
$X\setminus F$ zaś otwarty. Różnica między nimi
jest równa $U\setminus F$, więc $X\setminus E\in\mathcal D$.

Niech teraz $E=\bigcup_{n\ge1}E_n$, gdzie $E_n\in\mathcal D$.
Dla danego $\varepsilon>0$ dobieramy
$F_n\subseteq E_n\subseteq U_n$, z $F_n$ domkniętym,
$U_n$ otwartym oraz
$\mu(U_n\setminus F_n)<\varepsilon2^{-n-3}$.
Ponieważ $\mu(E)<\infty$ i
$\bigcup_{n\le N}E_n\uparrow E$, istnieje $N$ takie,
że $\mu(E\setminus\bigcup_{n\le N}E_n)<\varepsilon/4$.
Stawiamy $F=\bigcup_{n\le N}F_n$ i $U=\bigcup_{n\ge1}U_n$.
Zbiór $F$ jest domknięty jako suma skończona,
$U$ jest otwarty, a $F\subseteq E\subseteq U$.
Ponieważ
\[
U\setminus F\subseteq
(U\setminus E)\cup(E\setminus F),
\quad
U\setminus E\subseteq\bigcup_{n\ge1}(U_n\setminus E_n),
\]
oraz
$E\setminus F\subseteq
(E\setminus\bigcup_{n\le N}E_n)
\cup\bigcup_{n\le N}(E_n\setminus F_n)$,
subaddytywność daje
$\mu(U\setminus F)<\varepsilon$.
W konsekwencji $\mathcal D$ jest $\sigma$-algebrą.

\emph{Krok 3: zakończenie.}
Ponieważ $\mathcal D$ zawiera wszystkie zbiory otwarte,
zawiera także $\sigma$-algebrę przez nie generowaną,
czyli $\mathcal B(X)$. Równości w definicjach regularności
wynikają z monotoniczności miary oraz
$\mu(E\setminus F),\mu(U\setminus E)\le\mu(U\setminus F)$.
\end{proof}

\begin{uwaga}[Dlaczego nie można automatycznie zastąpić $F$ zbiorem zwartym?]
W przestrzeni metrycznej, która nie jest lokalnie zwarta,
zbiór domknięty i ograniczony nie musi być zwarty.
Na przykład zamknięta kula jednostkowa w nieskończenie
wymiarowej przestrzeni Hilberta nie jest zwarta.
Aproksymacja przez zbiory zwarte wymaga osobnego argumentu.
\end{uwaga}

\section{Ciasność i miary na przestrzeniach polskich}
\begin{definicja}[Ciasność pojedynczej miary]
\label{def:tight-single-20}
Niech $X$ będzie przestrzenią topologiczną, a $\mu$
skończoną miarą borelowską. Mówimy, że $\mu$ jest
\emph{ciasna}, jeżeli dla każdego $\varepsilon>0$
istnieje zwarty zbiór $K\subseteq X$ taki,
że $\mu(X\setminus K)<\varepsilon$.
\end{definicja}

\begin{twierdzenie}[Ciasność skończonych miar na przestrzeniach polskich]
\label{thm:polish-tight-20}
Niech $X$ będzie przestrzenią polską, tj. przestrzenią
topologiczną z metryką zupełną i ośrodkiem przeliczalnym,
zgodną z jej topologią. Każda skończona miara borelowska
$\mu$ na $X$ jest ciasna. Ponadto dla każdego
$E\in\mathcal B(X)$ zachodzi
\[
 \mu(E)=\sup\{\mu(K):K\subseteq E,\ K\text{ zwarty}\}.
\]
\end{twierdzenie}
\begin{proof}
Wybieramy zgodną metrykę zupełną $d$ oraz przeliczalny
zbiór gęsty $D=\{x_j:j\ge1\}$. Ustalmy $\varepsilon>0$.
Dla każdego $n\ge1$ kule otwarte
$B(x_j,2^{-n})$, $j\ge1$, pokrywają $X$.
Z ciągłości miary z dołu, stosowanej do ich skończonych
sum, istnieje $m_n$ takie, że
\[
\mu\left(X\setminus\bigcup_{j=1}^{m_n}B(x_j,2^{-n})\right)
<\varepsilon 2^{-n-1}.
\]
Niech $C_n=\bigcup_{j=1}^{m_n}\overline{B}(x_j,2^{-n})$.
Zbiór $C_n$ jest domknięty i
$\mu(X\setminus C_n)<\varepsilon2^{-n-1}$.
Kładziemy $K=\bigcap_{n\ge1}C_n$.
Zbiór $K$ jest domknięty, a dla każdego $n$
zawiera się w sumie skończenie wielu kul o promieniu
$2^{-n}$. Jest zatem całkowicie ograniczony:
dla każdego $\delta>0$ wybieramy $n$ tak duże,
by $2^{1-n}<\delta$, i otrzymujemy skończone
$\delta$-pokrycie $K$ przez kule o środkach z $X$.
Ponieważ $K$ jest domknięty w zupełnym $X$, jest zupełny.
Zupełność i całkowita ograniczoność implikują zwartość:
każdy ciąg ma podciąg Cauchy'ego, a ten jest zbieżny
w $K$. Wreszcie
\[
 \mu(X\setminus K)
 \le\sum_{n\ge1}\mu(X\setminus C_n)<\varepsilon.
\]
Udowodniliśmy ciasność.

Teraz ustalmy $E\in\mathcal B(X)$. Z twierdzenia
\ref{thm:finite-metric-regular-20} wybieramy domknięty
$F\subseteq E$ taki, że $\mu(E\setminus F)<\varepsilon/2$.
Z ciasności wybieramy zwarty $K_0\subseteq X$
z $\mu(X\setminus K_0)<\varepsilon/2$.
Zbiór $K=F\cap K_0$ jest zwarty, zawiera się w $E$, a
$\mu(E\setminus K)\le\mu(E\setminus F)+\mu(X\setminus K_0)
<\varepsilon$. Monotoniczność daje żądaną równość.
\end{proof}

\begin{przyklad}[Rozkłady gaussowskie w nieskończonym wymiarze]
Twierdzenie~\ref{thm:polish-tight-20} ma zastosowanie
do każdej skończonej miary borelowskiej na separowalnej
przestrzeni Banacha z metryką normową, również w nieskończonym
wymiarze. Nie oznacza to, że domknięte kule są zwarte;
zwarty zbiór zapewniający ciasność może mieć znacznie
bardziej skomplikowaną geometrię. W szczególności
przestrzeń $\ell^2$ jest polska, więc jej skończone
miary borelowskie są ciasne.
\end{przyklad}

\section{Miary Radona na przestrzeniach lokalnie zwartych}
\begin{twierdzenie}[Lokalnie skończone miary borelowskie]
\label{thm:lch-secondcountable-radon-20}
Niech $X$ będzie lokalnie zwartą, ośrodkową
przestrzenią metryczną, a $\mu$ lokalnie skończoną
miarą borelowską na $X$. Wówczas $\mu$ jest miarą
Radona w sensie definicji~\ref{def:radon-measure-20}.
Ponadto $\mu$ jest $\sigma$-skończona.
\end{twierdzenie}
\begin{proof}
\emph{Krok 1: otwarte pokrycie skończonej miary.}
Dla każdego $x\in X$ istnieje otwarte otoczenie $W_x$
z $\mu(W_x)<\infty$. Lokalna zwartość i regularność
przestrzeni metrycznej pozwalają wybrać otwarty zbiór
$V_x$ taki, że $x\in V_x$, $\overline{V_x}$ jest zwarty
oraz $\overline{V_x}\subseteq W_x$.
W szczególności $\mu(\overline{V_x})<\infty$.
Ponieważ przestrzeń metryczna ośrodkowa jest drugiej
przeliczalności i ma własność Lindelöfa,
z pokrycia $(V_x)_{x\in X}$ wybieramy przeliczalne
podpokrycie $(V_n)_{n\ge1}$. Dla każdego $n$
$\overline{V_n}$ jest zwarty i ma skończoną miarę.
Już stąd wynika $\sigma$-skończoność $\mu$.

\emph{Krok 2: regularność zewnętrzna.}
Rozbijamy $X$ na parami rozłączne zbiory borelowskie
$B_1=V_1$, $B_n=V_n\setminus\bigcup_{j<n}V_j$.
Dla danego $E\in\mathcal B(X)$ i $\varepsilon>0$
rozważamy skończoną miarę
$\mu_n(A)=\mu(A\cap V_n)$ na przestrzeni metrycznej $X$.
Z twierdzenia~\ref{thm:finite-metric-regular-20}
wybieramy otwarty $U_n\supseteq E\cap B_n$ taki, że
$\mu_n(U_n\setminus(E\cap B_n))<\varepsilon2^{-n-1}$.
Niech $U=\bigcup_{n\ge1}(U_n\cap V_n)$.
Wówczas $U$ jest otwarty, $E\subseteq U$, ponieważ
$E\cap B_n\subseteq U_n\cap V_n$. Użycie zbiorów
$B_n$ zamiast samych $V_n$ jest istotne: są one
rozłączne i każde $E\cap B_n$ pokrywa dokładnie tę
część $E$, która jest przypisana do indeksu $n$.
Ponadto $(U_n\cap V_n)\setminus E\subseteq
U_n\setminus(E\cap B_n)$, co daje
\[
\mu(U\setminus E)
\le\sum_{n\ge1}\mu((U_n\cap V_n)\setminus(E\cap B_n))
<\varepsilon.
\]
Jeżeli $\mu(E)<\infty$, daje to
$\mu(U)<\mu(E)+\varepsilon$; jeżeli $\mu(E)=\infty$,
regularność zewnętrzna jest natychmiastowa,
ponieważ każdy otwarty nadzbiór ma miarę nieskończoną.

\emph{Krok 3: regularność wewnętrzna względem zwartych.}
Niech $E\in\mathcal B(X)$ i $\varepsilon>0$.
Dla każdego $n$ rozważamy przestrzeń zwartą
$\overline{V_n}$ z miarą skończoną
$\mu|_{\overline{V_n}}$ i metryką indukowaną.
Z twierdzenia~\ref{thm:finite-metric-regular-20}
dla zbioru $E\cap V_n$ wybieramy zbiór
$K_n$ domknięty w $\overline{V_n}$, zawarty w
$E\cap V_n$, o własności
$\mu((E\cap V_n)\setminus K_n)<\varepsilon2^{-n-1}$.
Każdy $K_n$ jest zwarty w $X$.
Dla $E$ o skończonej mierze wybieramy $N$ takie, że
$\mu(E\setminus\bigcup_{n\le N}V_n)<\varepsilon/2$.
Zbiór $K=\bigcup_{n\le N}K_n$ jest zwarty i zawiera się
w $E$, a
\[
 \mu(E\setminus K)
 \le\mu(E\setminus\bigcup_{n\le N}V_n)
 +\sum_{n\le N}\mu((E\cap V_n)\setminus K_n)
 <\varepsilon.
\]
Gdy $\mu(E)=\infty$, z ciągłości miary z dołu
$\mu(E\cap\bigcup_{n\le N}V_n)\uparrow\infty$.
Dla dowolnego $M>0$ wybieramy $N$ z miarą tej części
większą niż $M+1$ i ponawiamy aproksymację przez
$K_n$, z błędami sumującymi się do mniej niż $1$.
Dostajemy zwarty $K\subseteq E$ z $\mu(K)>M$.
To dowodzi \eqref{eq:inner-compact-20} także w przypadku
nieskończonej masy. Razem z krokami 1--2 daje tezę.
\end{proof}

\begin{uwaga}[Na jakich przestrzeniach twierdzenie obowiązuje?]
Założenie ośrodkowości i metryczności nie jest częścią
definicji miary Radona; zostało wprowadzone w powyższym
twierdzeniu, aby można było podać samodzielny dowód.
Na ogólnych lokalnie zwartych przestrzeniach Hausdorffa
regularność nie każdej miary borelowskiej jest automatyczna.
Nie wystarczy również samo założenie lokalnej skończoności.
\end{uwaga}

\begin{przyklad}[Lokalna skończoność nie wystarcza]
\label{ex:nonradon-discrete-20}
Niech $X$ będzie nieprzeliczalnym zbiorem z topologią
dyskretną. Jest to lokalnie zwarta przestrzeń metryczna,
ale nieośrodkowa. Wszystkie podzbiory są borelowskie.
Definiujemy
\[
\mu(A)=\begin{cases}0,&A\text{ jest co najwyżej przeliczalny},\\
+\infty,&A\text{ jest nieprzeliczalny}.
\end{cases}
\]
To rzeczywiście miara: w przeliczalnej rodzinie zbiorów
rozłącznych suma miar jest dodatnia nieskończona wtedy
i tylko wtedy, gdy któryś ze zbiorów jest nieprzeliczalny.
Miara jest lokalnie skończona, ponieważ każde otwarte
jednoelementowe otoczenie ma miarę zero. Tymczasem
zbiory zwarte w przestrzeni dyskretnej są skończone,
więc wszystkie mają miarę zero, choć $\mu(X)=\infty$.
Regularność wewnętrzna względem zwartych zawodzi.
\end{przyklad}

\section{Przykłady miar Radona}
\begin{przyklad}[Miara Lebesgue'a jako miara Radona]
Na $\mathbb R^d$, $d\ge1$, ograniczenie miary Lebesgue'a
$\lambda_d$ do $\mathcal B(\mathbb R^d)$ jest miarą
lokalnie skończoną: każdy punkt ma otoczenie będące
kulą o skończonej objętości. Przestrzeń $\mathbb R^d$
jest lokalnie zwarta, metryczna i ośrodkowa.
Na mocy twierdzenia~\ref{thm:lch-secondcountable-radon-20}
miara ta jest Radona. Jej uzupełnienie na
$\mathcal L_d$ zawiera także zbiory nieborelowskie;
trzeba rozróżniać tę miarę zupełną od jej ograniczenia
borelowskiego.
\end{przyklad}

\begin{przyklad}[Miary z lokalnie całkowalną gęstością]
Niech $w:\mathbb R^d\to[0,\infty]$ będzie borelowska
i lokalnie całkowalna względem $\lambda_d$,
tj. $\int_Kw\,\mathrm d\lambda_d<\infty$ dla każdego
zwartego $K\subseteq\mathbb R^d$. Wówczas
$\mu_w(E)=\int_Ew(x)\,\mathrm d\lambda_d(x)$
dla $E\in\mathcal B(\mathbb R^d)$ jest miarą Radona.
Warunek lokalnej całkowalności gwarantuje lokalną
skończoność, a regularność wynika z poprzedniego
twierdzenia. Na przykład $w(x)=\exp(|x|^2)$ daje miarę
Radona o nieskończonej masie całkowitej.
\end{przyklad}

\begin{przyklad}[Atomy gęste i nośnik całej prostej]
Niech $(q_n)_{n\ge1}$ będzie numeracją zbioru
$\mathbb Q$ bez powtórzeń. Na $\mathbb R$ określmy
$\mu=\sum_{n\ge1}2^{-n}\delta_{q_n}$.
Jest to skończona miara borelowska, a więc miara Radona.
Każdy niepusty przedział otwarty zawiera pewien $q_n$
i ma dodatnią miarę $\mu$, choć $\mu(\mathbb R\setminus
\mathbb Q)=0$. W konsekwencji nośnikiem topologicznym
$\mu$ jest cała prosta. Pełny nośnik topologiczny
nie oznacza bezwzględnej ciągłości względem miary Lebesgue'a:
w tym przykładzie $\mu\perp\lambda_1$.
\end{przyklad}

\section{Nośnik topologiczny}
\begin{definicja}[Nośnik miary]
\label{def:support-radon-20}
Dla miary borelowskiej $\mu$ na przestrzeni topologicznej
$X$ określamy jej \emph{nośnik topologiczny} przez
\[
\operatorname{supp}\mu
:=\{x\in X:\mu(U)>0\text{ dla każdego otwartego }
U\ni x\}.
\]
Jest to zbiór domknięty. Nie jest to z definicji
wybrany mierzalny zbiór skupiający miarę w sensie
rozdziału~\ref{ch:ac-singular-17}.
\end{definicja}

\begin{twierdzenie}[Nośnik na przestrzeni drugiej przeliczalności]
\label{thm:support-fullmass-20}
Niech $X$ będzie przestrzenią drugiej przeliczalności,
a $\mu$ dowolną miarą borelowską na $X$. Wtedy
$\operatorname{supp}\mu$ jest domknięty,
$\mu(X\setminus\operatorname{supp}\mu)=0$ i jest
najmniejszym domkniętym zbiorem o pełnej mierze
w sensie zawierania.
\end{twierdzenie}
\begin{proof}
Z definicji
$X\setminus\operatorname{supp}\mu$
jest sumą wszystkich otwartych zbiorów o mierze zero,
a więc jest otwarty. Niech $(B_n)$ będzie przeliczalną
bazą topologii. Jeżeli $x$ należy do tej sumy, to
istnieje otwarty $U\ni x$ o $\mu(U)=0$.
Można wybrać $B_n$ takie, że $x\in B_n\subseteq U$.
Wówczas $\mu(B_n)=0$, a więc
\[
X\setminus\operatorname{supp}\mu
=\bigcup_{\{n:\mu(B_n)=0\}}B_n.
\]
Przeliczalna podaddytywność daje miarę zero tej sumy.
Niech $F\subseteq X$ będzie domknięty i
$\mu(X\setminus F)=0$. Każdy $x\notin F$ ma
otwarte otoczenie $X\setminus F$ miary zero,
zatem $x\notin\operatorname{supp}\mu$.
W konsekwencji $\operatorname{supp}\mu\subseteq F$.
\end{proof}

\begin{uwaga}[Dlaczego założenie przeliczalności ma znaczenie?]
W przykładzie~\ref{ex:nonradon-discrete-20} każda
jednopunktowa okolica ma miarę zero, zatem
$\operatorname{supp}\mu=\varnothing$, choć
$\mu(X)=+\infty$. Problem wynika z faktu, że
nieprzeliczalna suma zbiorów zerowych nie musi być
zerowa. Pokazuje to, dlaczego nie należy dowodzić
pełności miary nośnika przez nieuzasadnione przejście
od dowolnego pokrycia do przeliczalnej sumy.
\end{uwaga}

\section{Funkcje ciągłe o zwartym nośniku}
Dla lokalnie zwartej przestrzeni Hausdorffa $X$ oznaczamy
przez $C_c(X)$ przestrzeń ciągłych funkcji
$f:X\to\mathbb R$, których \emph{nośnik funkcji}
$\operatorname{supp}f=\overline{\{x:f(x)\neq0\}}$
jest zwarty. W analogiczny sposób definiuje się
$C_c(X;\mathbb C)$. Nie należy mylić nośnika funkcji
z nośnikiem topologicznym miary.

\begin{lemat}[Funkcja odcinająca]
\label{lem:cutoff-radon-20}
Niech $X$ będzie lokalnie zwartą przestrzenią metryczną,
$K\subseteq X$ zwartym zbiorem, a $U\subseteq X$
otwartym zbiorem zawierającym $K$. Wtedy istnieje
$f\in C_c(X)$ taka, że $0\le f\le1$, $f|_K=1$
oraz $\operatorname{supp}f\subseteq U$.
\end{lemat}
\begin{proof}
Jeśli $K=\varnothing$, wystarczy przyjąć $f=0$.
W dalszej części zakładamy $K\neq\varnothing$.
Dla każdego $x\in K$ wybieramy otwarte otoczenie
$W_x$ takie, że $x\in W_x$, $\overline{W_x}$ jest
zwarty i $\overline{W_x}\subseteq U$.
Zwartość $K$ pozwala wybrać skończenie wiele takich
zbiorów. Ich suma $V$ jest otwarta, zawiera $K$,
a $\overline V$ jest zwartym podzbiorem $U$.
Jeżeli $V=X$, przestrzeń $X$ jest zwarta
i możemy przyjąć $f\equiv1$. W przeciwnym razie
określamy
\[
 f(x)=\frac{d(x,X\setminus V)}
 {d(x,X\setminus V)+d(x,K)}.
\]
Mianownik nie zeruje się nigdzie: dla $x\in K$
pierwszy składnik jest dodatni, a poza $K$ drugi
składnik jest dodatni (ponieważ $K$ jest domknięty).
Obie funkcje odległości są ciągłe, stąd $f$ jest
ciągła; na $K$ ma wartość $1$, poza $V$ jest zerowa,
a $\operatorname{supp}f\subseteq\overline V\subseteq U$.
\end{proof}

\begin{twierdzenie}[Jednoznaczność miary Radona przez całki]
\label{thm:uniqueness-cc-20}
Niech $X$ będzie lokalnie zwartą przestrzenią metryczną,
a $\mu,\nu$ miarami Radona na $X$. Jeżeli
\[
\int_X f(x)\,\mathrm d\mu(x)
=\int_Xf(x)\,\mathrm d\nu(x)
\qquad\text{dla każdego }f\in C_c(X),
\]
to $\mu=\nu$ na $\mathcal B(X)$.
\end{twierdzenie}
\begin{proof}
Dla każdego otwartego $U\subseteq X$ i dowolnej
miary Radona $\rho$ na $X$ zachodzi
\begin{equation}\label{eq:cc-open-sup-20}
\rho(U)=\sup\left\{\int_Xf\,\mathrm d\rho:
 f\in C_c(X),\ 0\le f\le1,\
 \operatorname{supp}f\subseteq U\right\}.
\end{equation}
Istotnie, każda z powyższych całek jest co najwyżej
$\rho(U)$. Z drugiej strony dla każdego zwartego
$K\subseteq U$ lemat~\ref{lem:cutoff-radon-20}
daje $f$ równe $1$ na $K$, więc
$\int_Xf\,\mathrm d\rho\ge\rho(K)$.
Z regularności wewnętrznej miary $\rho$
uzyskujemy \eqref{eq:cc-open-sup-20} także dla
$\rho(U)=+\infty$. Przy założonej równości całek
zbiory funkcji w supremum są te same dla obu miar,
a więc $\mu(U)=\nu(U)$ dla wszystkich otwartych $U$.
Wreszcie dla dowolnego zbioru borelowskiego
$E\subseteq X$ z regularności zewnętrznej
$\mu(E)=\inf_{U\supseteq E}\mu(U)
=\inf_{U\supseteq E}\nu(U)=\nu(E)$.
\end{proof}

\begin{uwaga}[Związek z twierdzeniem reprezentacyjnym]
Twierdzenie~\ref{thm:uniqueness-cc-20} daje
\emph{jednoznaczność} miary reprezentującej funkcjonał
$L(f)=\int_Xf\,\mathrm d\mu$ na $C_c(X)$.
Istnienie takiej miary dla dodatniego funkcjonału
jest przedmiotem następnego rozdziału, poświęconego
twierdzeniu Riesza--Markowa--Kakutaniego. W przestrzeni
niezwartej funkcjonał całkowania względem miary
Radona o nieskończonej masie nie musi być ograniczony
w normie supremum na całym $C_c(X)$; dlatego rozróżnimy
wtedy $C_c(X)$ i $C_0(X)$.
\end{uwaga}

\begin{twierdzenie}[Gęstość $C_c(X)$ w przestrzeniach $L^p$]
\label{thm:cc-dense-lp-radon-20}
Niech $X$ będzie lokalnie zwartą przestrzenią metryczną,
a $\mu$ miarą Radona w sensie definicji
\ref{def:radon-measure-20}. Wówczas dla każdego
$1\le p<\infty$ przestrzeń $C_c(X;\mathbb C)$ jest
gęsta w $L^p(X,\mathcal B(X),\mu;\mathbb C)$.
Analogiczny wniosek zachodzi dla funkcji rzeczywistych.
\end{twierdzenie}
\begin{proof}
Najpierw ustalmy $E\in\mathcal B(X)$ takie, że
$\mu(E)<\infty$. Dla $\varepsilon>0$ z wewnętrznej
i zewnętrznej regularności wybieramy zwarty $K$
i otwarty $U$ takie, że
$K\subseteq E\subseteq U$ oraz
$\mu(E\setminus K)<\varepsilon/2$,
$\mu(U\setminus E)<\varepsilon/2$.
Lemat~\ref{lem:cutoff-radon-20} daje $g\in C_c(X)$
taką, że $0\le g\le1$, $g|_K=1$ oraz
$\operatorname{supp}g\subseteq U$.
Na $K$ funkcje $\chi_E$ i $g$ są równe, poza $U$
obie znikają, a w pozostałej części ich różnica
ma moduł co najwyżej $1$. Zatem
\[
 \int_X|\chi_E(x)-g(x)|^p\,\mathrm d\mu(x)
 \le\mu(E\setminus K)+\mu(U\setminus E)
 <\varepsilon.
\]
Otrzymujemy aproksymację funkcji charakterystycznej
każdego zbioru o skończonej mierze.

Funkcje proste złożone ze skończenie wielu składników
$a_j\chi_{E_j}$, gdzie $\mu(E_j)<\infty$, są gęste
w $L^p(\mu)$ dla $1\le p<\infty$:
przybliżamy $f$ przez
$f_n=f\chi_{\{1/n<|f|\le n\}}$.
Funkcja $f_n$ jest ograniczona i ma nośnik
zawarty w zbiorze o skończonej mierze
(na mocy nierówności Markowa); ponadto
$\|f_n-f\|_p\to0$ z twierdzenia o zbieżności
zmajoryzowanej zastosowanego do $|f_n-f|^p$.
Dla każdego $n$ aproksymujemy $f_n$
funkcjami prostymi na tym zbiorze.
Dla ustalonej liczby składników przybliżamy
każdą $\chi_{E_j}$ odpowiednią funkcją $g_j\in C_c(X)$;
z nierówności trójkąta w $L^p$ suma
$\sum_ja_jg_j$ przybliża funkcję prostą z dowolną
dokładnością. Jest ona ciągła i ma zwarty nośnik.
To dowodzi żądanej gęstości.
\end{proof}

\begin{uwaga}[Znaczenie założenia $p<\infty$]
Gęstość w normie $L^p$ nie rozciąga się bez dodatkowych
warunków na $p=\infty$. Na przykład względem miary
Lebesgue'a na $[0,1]$ funkcji $\chi_{[0,1/2]}$
nie da się aproksymować w normie $L^\infty$
funkcjami ciągłymi: ciągłość wymusza przejście pomiędzy
wartościami bliskimi $0$ i $1$ na zbiorach dodatniej miary.
\end{uwaga}

\section{Regularność a uzupełnienie miary}
Jeżeli $\mu$ jest miarą borelowską, jej uzupełnienie
(wprowadzone w rozdziale~2) może zawierać zbiory
nieborelowskie, ale nie zmienia wartości miary
na zbiorach borelowskich. Na przykład miara Lebesgue'a
na $\mathcal L_d$ jest uzupełnieniem miary borelowskiej
$\lambda_d|_{\mathcal B(\mathbb R^d)}$. Z drugiej strony
nie każda miara Radona jest zupełna na swojej
$\sigma$-algebrze borelowskiej: jeśli borelowski zbiór
zerowy ma podzbiór nieborelowski, to uzupełnienie
ma większą dziedzinę. Jest to rozróżnienie
istotne przy formułowaniu twierdzeń o obrazach miar,
całkowaniu i odsyłaczach do $L^p$.

\section{Co przygotowaliśmy na następny rozdział?}
Przeprowadziliśmy dwa kierunki regularności:
miary skończone na przestrzeniach metrycznych są
aproksymowane od wewnątrz przez zbiory domknięte,
a na przestrzeniach polskich przez zbiory zwarte;
miary lokalnie skończone na lokalnie zwartych
przestrzeniach metrycznych ośrodkowych są miarami Radona.
Ustaliliśmy ponadto, że wszystkie całki z funkcji
$C_c(X)$ jednoznacznie wyznaczają miarę Radona.
Następny rozdział odwróci tę konstrukcję:
zacznie od dodatniego funkcjonału na funkcjach ciągłych
i skonstruuje reprezentującą go miarę.

% END chapters/pl/measure_theory/20-rozdzial.tex


% BEGIN chapters/pl/measure_theory/21-rozdzial.tex
% ROZDZIAŁ 21 — RIESZ–MARKOW–KAKUTANI
\chapter{Twierdzenie Riesza--Markowa--Kakutaniego}
\label{ch:riesz-markov-21}

\section{Sformułowanie i rola twierdzenia}
Niech $X$ będzie lokalnie zwartą przestrzenią Hausdorffa,
a $\mathbb K\in\{\mathbb R,\mathbb C\}$.
Przestrzeń $C_c(X;\mathbb K)$ i pojęcie zwartego nośnika
wprowadzono w rozdziale~\ref{ch:radon-measures-20}.
Przez $C_0(X;\mathbb K)$ oznaczamy przestrzeń funkcji ciągłych
\emph{znikających w nieskończoności}, to znaczy takich, że
dla każdego $\varepsilon>0$ zbiór
$\{x:|f(x)|\ge\varepsilon\}$ jest zwarty.
Przestrzeń $C_0(X;\mathbb K)$ z normą supremum jest przestrzenią
Banacha. W rozważanej dalej klasie lokalnie zwartych
ośrodkowych przestrzeni metrycznych udowodnimy gęstość
$C_c(X;\mathbb K)$ w $C_0(X;\mathbb K)$.

\begin{twierdzenie}[Reprezentacja dodatnich funkcjonałów]
\label{thm:riesz-positive-21}
Niech $X$ będzie \emph{lokalnie zwartą, ośrodkową przestrzenią
metryczną}, a $L:C_c(X;\mathbb R)\to\mathbb R$ funkcjonałem
liniowym takim, że $L(f)\ge0$ dla $f\ge0$.
Istnieje dokładnie jedna miara Radona $\mu$ na $X$ spełniająca
\begin{equation}\label{eq:riesz-representation-21}
 L(f)=\int_X f(x)\,\mathrm d\mu(x),\qquad f\in C_c(X;\mathbb R).
\end{equation}
Przy czym $L$ nie musi być ograniczony w normie supremum na
całej przestrzeni $C_c(X)$.
\end{twierdzenie}

\begin{uwaga}[Zakres wersji twierdzenia]
Klasyczne twierdzenie zachodzi na ogólnych lokalnie zwartych
przestrzeniach Hausdorffa, przy właściwie określonej regularności
miary. Aby wszystkie kroki dowodu, w tym konstrukcję i regularność,
wyprowadzić tutaj z wcześniejszych rozdziałów bez ukrytych
argumentów topologicznych, udowodnimy je dla lokalnie zwartych
ośrodkowych przestrzeni metrycznych. Obejmuje to w szczególności
$\mathbb R^d$, rozmaitości drugiej przeliczalności i przestrzenie
polskie lokalnie zwarte. Wersja dla dowolnych lokalnie zwartych
przestrzeni Hausdorffa wymaga dodatkowej pracy z regularnymi
miarami zewnętrznymi; nie utożsamiamy tych dwóch zakresów.
\end{uwaga}

\section{Lemat o metrycznych miarach zewnętrznych}
\begin{definicja}[Miara zewnętrzna metryczna]
\label{def:metric-outer-21}
Niech $(Z,d)$ będzie przestrzenią metryczną.
Miara zewnętrzna $m^*:\mathcal P(Z)\to[0,\infty]$
jest \emph{metryczna}, gdy dla dowolnych $A,B\subseteq Z$
o dodatniej odległości $d(A,B)>0$ zachodzi
$m^*(A\cup B)=m^*(A)+m^*(B)$.
Dla zbioru pustego lub dowolnego niepustego przyjmujemy
$d(\varnothing,B)=+\infty$.
\end{definicja}

\begin{lemat}[Borelowska mierzalność]
\label{lem:metric-outer-borel-21}
Jeżeli $m^*$ jest metryczną miarą zewnętrzną na $(Z,d)$,
to każdy zbiór borelowski w $Z$ jest mierzalny w sensie
Carathéodory'ego.
\end{lemat}
\begin{proof}
Wystarczy wykazać mierzalność zbiorów domkniętych, gdyż
rodzina zbiorów mierzalnych Carathéodory'ego jest
$\sigma$-algebrą (rozdział~\ref{ch:caratheodory}).
Niech $F\subseteq Z$ będzie domknięty i $E\subseteq Z$.
Musimy dowieść
$m^*(E)\ge m^*(E\cap F)+m^*(E\setminus F)$;
nierówność odwrotna wynika z podaddytywności.
Jeśli $m^*(E)=+\infty$, teza jest oczywista. Załóżmy
$m^*(E)<\infty$ i oznaczmy $r(x)=d(x,F)$.
Dla każdego $n\ge1$ zbiory $E\cap F$ oraz
$E\cap\{r>1/n\}$ mają dodatnią odległość, gdyż
$r$ jest $1$-lipschitzowska. Zatem
\begin{equation}\label{eq:metric-outer-sep-21}
 m^*(E)\ge m^*(E\cap F)+m^*(E\cap\{r>1/n\}).
\end{equation}
Wykażemy, że drugi składnik po prawej zbiega do
$m^*(E\setminus F)$; nie wolno tutaj korzystać z ciągłości
z dołu, ponieważ $m^*$ jest dopiero miarą zewnętrzną.
Niech $A_0=E\cap\{r>1\}$ oraz
$A_j=E\cap\{1/(j+1)<r\le1/j\}$, $j\ge1$.
Dla ustalonej parzystości indeksów dowolne dwa
\emph{różne} zbiory $A_j$ mają dodatnią odległość.
Dla każdej skończonej rodziny takich zbiorów metryczność
i indukcja dają
$\sum m^*(A_j)=m^*(\bigcup A_j)\le m^*(E)$.
W konsekwencji obie serie, po indeksach parzystych
i nieparzystych, są zbieżne; zatem
$\sum_{j\ge1}m^*(A_j)<\infty$. Ponieważ
\[
 E\cap\{0<r\le1/n\}\subseteq\bigcup_{j\ge n}A_j,
\]
miara zewnętrzna części po lewej dąży do zera.
Z podaddytywności
\[
 m^*(E\setminus F)\le
 m^*(E\cap\{r>1/n\})+
 m^*(E\cap\{0<r\le1/n\}),
\]
a z monotoniczności wynika przeciwna nierówność
$m^*(E\cap\{r>1/n\})\le m^*(E\setminus F)$.
Przechodząc do granicy w \eqref{eq:metric-outer-sep-21},
otrzymujemy żądane kryterium Carathéodory'ego.
\end{proof}

\section{Przypadek zwarty: konstrukcja od podstaw}
\begin{twierdzenie}[Reprezentacja na zwartej przestrzeni metrycznej]
\label{thm:compact-riesz-21}
Niech $K$ będzie niepustą zwartą przestrzenią metryczną,
a $L:C(K;\mathbb R)\to\mathbb R$ dodatnim funkcjonałem
liniowym. Istnieje dokładnie jedna skończona regularna
miara borelowska $\mu$ na $K$ taka, że
$L(f)=\int_Kf\,\mathrm d\mu$ dla wszystkich $f\in C(K)$.
Ponadto $\mu(K)=L(\chi_K)$ i
$\|L\|=L(\chi_K)$.
\end{twierdzenie}
\begin{proof}
\emph{Krok 1: lokalne ograniczenie funkcjonału.}
Z $-\|f\|_\infty\chi_K\le f\le\|f\|_\infty\chi_K$
wynika $|L(f)|\le L(\chi_K)\|f\|_\infty$.
Dla $f=\chi_K$ mamy równość w normie.

\emph{Krok 2: przypisanie liczb zbiorom otwartym.}
Dla zbioru otwartego $U\subseteq K$ definiujemy
\begin{equation}\label{eq:open-content-21}
 a(U)=\sup\{L(f): f\in C(K),\ 0\le f\le1,\
          \operatorname{supp}f\subseteq U\}.
\end{equation}
W szczególności $a(\varnothing)=0$,
$0\le a(U)\le L(\chi_K)$, $a(K)=L(\chi_K)$,
a $a$ jest monotoniczna. Jeżeli $U,V$ są rozłączne
i otwarte, to $a(U\cup V)\ge a(U)+a(V)$:
wybieramy funkcje $f,g$ dopuszczalne odpowiednio dla
$U,V$; ponieważ ich nośniki są rozłączne,
$0\le f+g\le1$ i $L(f+g)=L(f)+L(g)$.

\emph{Krok 3: przeliczalna podaddytywność na otwartych.}
Jeżeli $U\subseteq\bigcup_{j\ge1}U_j$ i $f$ jest
dopuszczalna w \eqref{eq:open-content-21}, to
$\operatorname{supp}f$ ma skończone pokrycie
$U_{j_1},\ldots,U_{j_N}$. Istnieją ciągłe funkcje
$\theta_i\ge0$ o nośnikach zawartych w $U_{j_i}$,
takie, że $\sum_i\theta_i=1$ na
$\operatorname{supp}f$. Aby je zbudować, używamy
funkcji odcinających równych $1$ na skończonym
pokryciu mniejszymi zbiorami i dzielimy ich sumę
w otoczeniu nośnika; poza tym otoczeniem ustawiamy
funkcje na zero. Każde $f\theta_i$ jest nieujemne,
nie przekracza $1$ i ma nośnik w $U_{j_i}$.
Wobec tego
$L(f)=\sum_i L(f\theta_i)\le\sum_{j\ge1}a(U_j)$.
Po wzięciu supremum otrzymujemy
$a(U)\le\sum_j a(U_j)$.

\emph{Krok 4: miara zewnętrzna.}
Dla dowolnego $A\subseteq K$ definiujemy
\begin{equation}\label{eq:riesz-outer-21}
 m^*(A)=\inf\{a(U):A\subseteq U,\ U\text{ otwarty}\}.
\end{equation}
Jeżeli $A\subseteq\bigcup_j A_j$, dobieramy otwarte
$U_j\supseteq A_j$ z
$a(U_j)\le m^*(A_j)+\varepsilon2^{-j}$ (przyjmujemy
osobno trywialny przypadek nieskończonej sumy).
Krok 3 daje
$m^*(A)\le a(\bigcup_jU_j)\le\sum_jm^*(A_j)+\varepsilon$.
Zatem $m^*$ jest miarą zewnętrzną.
Dla zbioru otwartego $U$ mamy $m^*(U)=a(U)$, ponieważ
$U\subseteq V$ implikuje $a(U)\le a(V)$.
Jeśli $d(A,B)=\delta>0$, to dla każdego otwartego
$W\supseteq A\cup B$ zbiory
$W_A=W\cap\{x:d(x,A)<\delta/3\}$ i
$W_B=W\cap\{x:d(x,B)<\delta/3\}$ są otwarte,
rozłączne i zawierają odpowiednio $A,B$. Stąd
$a(W)\ge a(W_A)+a(W_B)\ge m^*(A)+m^*(B)$.
Biorąc infimum po $W$ otrzymujemy metryczność $m^*$.
Lemat~\ref{lem:metric-outer-borel-21} zapewnia, że
wszystkie zbiory borelowskie są mierzalne. Ograniczenie
$\mu=m^*|_{\mathcal B(K)}$ jest więc miarą, a
$\mu(K)=L(\chi_K)<\infty$.

\emph{Krok 5: regularność.}
Regularność zewnętrzna wynika z \eqref{eq:riesz-outer-21}.
Dla otwartego $U$ i $\varepsilon>0$ dobieramy
$0\le f\le1$ z $\operatorname{supp}f\subseteq U$ i
$L(f)>\mu(U)-\varepsilon$. Przyjmijmy
$H=\operatorname{supp}f$, który jest zwarty.
Każdy otwarty $V\supseteq H$ spełnia
$L(f)\le a(V)$, zatem
$\mu(H)=\inf_{V\supseteq H}a(V)\ge L(f)$.
To dowodzi regularności wewnętrznej dla otwartych.
Dla dowolnego borelowskiego $E$ i $\varepsilon>0$
wybieramy otwarty $U\supseteq E$ z
$\mu(U\setminus E)<\varepsilon$.
Następnie wybieramy zwarty $H\subseteq U$ z
$\mu(U\setminus H)<\varepsilon$ oraz otwarty
$W\supseteq U\setminus E$ z
$\mu(W)<\mu(U\setminus E)+\varepsilon$.
Zbiór $H\setminus W$ jest zwarty, zawiera się w $E$,
a $\mu(E\setminus(H\setminus W))<3\varepsilon$.
Otrzymujemy regularność wewnętrzną względem zwartych.

\emph{Krok 6: identyfikacja całki z $L$.}
Niech $f\in C(K)$, $f\ge0$, $M=\|f\|_\infty>0$.
Dla $N\ge1$ połóżmy $\delta=M/N$ i
$U_i=\{f>i\delta\}$, $i=0,\ldots,N$.
Dla $i=0,\ldots,N-1$ zbiór
$H_i=\{f\ge(i+1)\delta\}$ jest zwarty i zawiera się
w $U_i$. Istnieje $h_i\in C(K)$, $0\le h_i\le1$,
$h_i=1$ na $H_i$, $\operatorname{supp}h_i\subseteq U_i$;
wynika to z lematu o funkcji odcinającej
\ref{lem:cutoff-radon-20}. Punktowo
$f\le\delta\chi_K+\delta\sum_{i=0}^{N-1}h_i$.
Z dodatniości $L$ dostajemy
\begin{equation}\label{eq:riesz-upper-sum-21}
 L(f)\le\delta L(\chi_K)+
           \delta\sum_{i=0}^{N-1}\mu(U_i).
\end{equation}
Z drugiej strony dla dowolnych funkcji
$g_i\in C(K)$, $0\le g_i\le1$,
$\operatorname{supp}g_i\subseteq U_i$ ($1\le i\le N$),
zachodzi $\delta\sum_{i=1}^{N}g_i\le f$ punktowo.
Zatem, biorąc niezależnie supremum po $g_i$,
\begin{equation}\label{eq:riesz-lower-sum-21}
 \delta\sum_{i=1}^{N}\mu(U_i)\le L(f).
\end{equation}
Natomiast rozcięcie poziomów funkcji daje
\[
 \delta\sum_{i=1}^{N}\mu(U_i)
 \le\int_Kf(x)\,\mathrm d\mu(x)
 \le\delta\sum_{i=0}^{N-1}\mu(U_i).
\]
Różnica między lewą i prawą sumą wynosi
$\delta\mu(U_0)\le\delta L(\chi_K)$.
Z \eqref{eq:riesz-upper-sum-21} i
\eqref{eq:riesz-lower-sum-21} wynika
$|L(f)-\int_K f\,\mathrm d\mu|\le2\delta L(\chi_K)$.
Wysyłając $N\to\infty$, otrzymujemy reprezentację
$f\ge0$. Ogólny $f$ rozkładamy na części dodatnią
i ujemną.

\emph{Krok 7: jednoznaczność.}
Dla dowolnej miary regularnej $\rho$ reprezentującej $L$
i otwartego $U$ regularność oraz funkcje odcinające
dają $\rho(U)=a(U)$ (por.
\eqref{eq:cc-open-sup-20}). Stąd z regularności
zewnętrznej $\rho=\mu$ na wszystkich zbiorach
borelowskich.
\end{proof}

\section{Przestrzenie lokalnie zwarte: lokalizacja i sklejanie}
\begin{lemat}[Wyczerpanie przez zbiory relatywnie zwarte]
\label{lem:exhaustion-21}
Niech $X$ będzie lokalnie zwartą ośrodkową przestrzenią
metryczną. Istnieją otwarte zbiory $U_n$ takie, że
$U_n\subseteq\overline U_n\subseteq U_{n+1}$,
$\overline U_n$ jest zwarty i $\bigcup_nU_n=X$.
Ponadto istnieją $\chi_n\in C_c(X)$ spełniające
$0\le\chi_n\le1$, $\chi_n=1$ na $\overline U_n$
oraz $\operatorname{supp}\chi_n\subseteq U_{n+1}$.
\end{lemat}
\begin{proof}
Lokalna zwartość daje każdemu punktowi relatywnie zwarte
otwarte otoczenie. Ośrodkowa przestrzeń metryczna jest
drugiej przeliczalności, więc z tego pokrycia wybieramy
przeliczalne podpokrycie $(V_j)_{j\ge1}$.
Indukcyjnie dobieramy relatywnie zwarte otwarte $U_n$,
obejmujące $\overline U_{n-1}\cup\overline V_1\cup\cdots
\cup\overline V_n$: skończone pokrycie tego zwartego
zbioru relatywnie zwartymi otoczeniami ma relatywnie
zwartą sumę. Dobierając otwarte otoczenie o zwartej
domkniętej powłoce, uzyskujemy żądane zawierania.
Funkcje $\chi_n$ istnieją na mocy lematu
\ref{lem:cutoff-radon-20}.
\end{proof}

\begin{proof}[Dowód twierdzenia~\ref{thm:riesz-positive-21}]
\emph{Krok 1: dodatniość lokalna i zwartość.}
Dla dowolnego zwartego $H\subseteq X$ wybieramy
$\eta\in C_c(X)$, $\eta\ge0$, $\eta=1$ na $H$.
Jeśli $\operatorname{supp}f\subseteq H$, to
$-\|f\|_\infty\eta\le f\le\|f\|_\infty\eta$,
więc $|L(f)|\le\|f\|_\infty L(\eta)$.
Wybierzmy wyczerpanie i funkcje $\chi_n$ z lematu
\ref{lem:exhaustion-21}; oznaczmy $K_n=\overline U_n$
i $J_n=\overline U_{n+1}$.
Dla $h\in C(J_n)$ funkcję $\chi_nh$ na $J_n$
przedłużamy przez zero na $X\setminus J_n$.
Jest ona ciągła, ponieważ
$\operatorname{supp}\chi_n\subseteq U_{n+1}$.
Definiujemy $L_n(h)=L(\chi_nh)$; jest to dodatni
funkcjonał na $C(J_n)$.
Z twierdzenia~\ref{thm:compact-riesz-21}
istnieje skończona regularna miara $\rho_n$
na $J_n$ reprezentująca $L_n$. Przedłużamy ją
przez zero na dopełnienie $J_n$, otrzymując miarę
borelowską na $X$ (także oznaczaną $\rho_n$).

\emph{Krok 2: zgodność miar na wyczerpaniu.}
Jeżeli $m\ge n$, to $\chi_m=\chi_n=1$ na $U_n$.
Dla dowolnego $f\in C_c(U_n)$ zachodzi zatem
\[
 \int_Xf\,\mathrm d\rho_m=L(f)
  =\int_Xf\,\mathrm d\rho_n.
\]
Ograniczenia $\rho_m|_{U_n}$ i $\rho_n|_{U_n}$
są miarami Radona (na otwartej przestrzeni $U_n$).
Jednoznaczność z twierdzenia~\ref{thm:uniqueness-cc-20}
daje $\rho_m|_{U_n}=\rho_n|_{U_n}$.
Uwaga: ograniczenie miary Radona do otwartego podzbioru
jest Radona; regularność wewnętrzna oraz zewnętrzna
wynika z regularności wyjściowej i przecięcia z podprzestrzenią.

\emph{Krok 3: sklejenie.}
Stawiamy $U_0=\varnothing$, $B_n=U_n\setminus U_{n-1}$.
Dla borelowskiego $E\subseteq X$ określamy
\begin{equation}\label{eq:gluing-measures-21}
 \mu(E)=\sum_{n\ge1}\rho_n(E\cap B_n).
\end{equation}
Ponieważ $B_n$ są parami rozłączne, a suma miar jest
miarą, $\mu$ jest miarą borelowską. Dla $E\subseteq U_m$
przeliczalna addytywność i krok 2 dają
$\mu(E)=\rho_m(E)$; dotyczy to również wszystkich
$E\in\mathcal B(U_m)$. W szczególności
$\mu(U_m)=\rho_m(U_m)<\infty$.
Zatem $\mu$ jest lokalnie skończona. Rozdział
\ref{ch:radon-measures-20} pokazuje, że na lokalnie
zwartej ośrodkowej przestrzeni metrycznej jest ona
miarą Radona.

\emph{Krok 4: reprezentacja.}
Jeśli $f\in C_c(X)$, jej nośnik mieści się w pewnym
$U_m$ (ponieważ $(U_m)$ jest rosnącym pokryciem
zwartego nośnika). Wtedy
\[
 \int_X f(x)\,\mathrm d\mu(x)
 =\int_X f(x)\,\mathrm d\rho_m(x)
 =L(\chi_m f)=L(f).
\]
Jednoznaczność wynika z twierdzenia
\ref{thm:uniqueness-cc-20}.
\end{proof}

\section{Przestrzeń $C_0(X)$ oraz miary o skończonej wariacji}
\begin{lemat}[Gęstość $C_c(X)$ w $C_0(X)$]
\label{lem:density-cc-c0-21}
Dla lokalnie zwartej ośrodkowej przestrzeni metrycznej
$X$ przestrzeń $C_c(X;\mathbb K)$ jest gęsta w
$C_0(X;\mathbb K)$ w normie supremum.
\end{lemat}
\begin{proof}
Niech $f\in C_0(X)$ i $\varepsilon>0$.
Zbiór $H=\{|f|\ge\varepsilon\}$ jest zwarty.
Wybieramy $\eta\in C_c(X)$ równą $1$ na $H$,
$0\le\eta\le1$. Wówczas $\eta f\in C_c(X)$
i $\|f-\eta f\|_\infty\le\varepsilon$.
\end{proof}

\begin{twierdzenie}[Ograniczone dodatnie funkcjonały na $C_0$]
\label{thm:riesz-c0-positive-21}
W założeniach twierdzenia~\ref{thm:riesz-positive-21}
niech $L$ będzie ograniczonym dodatnim funkcjonałem
na $C_0(X;\mathbb R)$. Wtedy jego miara reprezentująca
jest skończona i $\|L\|=\mu(X)$. Odwrotnie każda
skończona miara Radona definiuje taki funkcjonał.
\end{twierdzenie}
\begin{proof}
Dla kolejnych zbiorów $U_n$ z lematu
\ref{lem:exhaustion-21} wybieramy $\chi_n=1$
na $\overline U_n$, $0\le\chi_n\le1$.
Wówczas $\mu(U_n)\le\int_X\chi_n\,\mathrm d\mu
=L(\chi_n)\le\|L\|$. Przechodząc do granicy,
$\mu(X)\le\|L\|$. W przeciwną stronę,
$|L(f)|\le\mu(X)\|f\|_\infty$ dla $f\in C_c$;
przez gęstość nierówność zachodzi na $C_0$.
Każda skończona miara Radona definiuje przez całkę
funkcjonał liniowy ciągły i dodatni, co kończy dowód.
\end{proof}

\begin{definicja}[Skończona zespolona miara Radona]\label{idx:21:definicja:08}

Niech $X$ będzie lokalnie zwartą ośrodkową przestrzenią
metryczną. Skończoną zespoloną miarę borelowską $\nu$
nazywamy miarą Radona, jeśli jej wariacja całkowita
$|\nu|$ jest skończoną miarą Radona. Wówczas części
rzeczywista i urojona są miarami ze znakiem o skończonej
wariacji, a wszystkie całki z funkcji $C_0(X)$ istnieją.
\end{definicja}

\begin{twierdzenie}[Reprezentacja zespolona i norma wariacji]
\label{thm:riesz-complex-21}
Niech $X$ będzie lokalnie zwartą ośrodkową przestrzenią
metryczną, a $T:C_0(X;\mathbb C)\to\mathbb C$
ograniczonym funkcjonałem liniowym. Istnieje dokładnie
jedna skończona zespolona miara Radona $\nu$ taka, że
\[
 T(f)=\int_X f(x)\,\mathrm d\nu(x),\qquad f\in C_0(X;\mathbb C).
\]
Ponadto $\|T\|=|\nu|(X)$, gdzie $|\nu|$ oznacza
wariację całkowitą miary zespolonej.
\end{twierdzenie}
\begin{proof}
Dla nieujemnego $g\in C_c(X;\mathbb R)$ zdefiniujmy
\[
 M(g)=\sup\{|T(h)|:h\in C_c(X;\mathbb C),\ |h|\le g\}.
\]
Jest to skończone, bo $M(g)\le\|T\|\|g\|_\infty$.
Ma własności $M(ag)=aM(g)$ dla $a\ge0$ oraz
$M(g_1+g_2)=M(g_1)+M(g_2)$.
Dla nierówności „$\le$” rozkładamy dowolną $h$,
$|h|\le g_1+g_2$, na $h_i=h g_i/(g_1+g_2)$ na
$\{g_1+g_2>0\}$, a w pozostałych punktach przyjmujemy
$h_i=0$. Ciągłość przy zerach mianownika wynika
z $|h_i|\le g_i$; mamy $|h_i|\le g_i$ i
$h_1+h_2=h$. Dla nierówności „$\ge$” bierzemy
$h_i$ niemal realizujące supremum i mnożymy je przez
stałe fazy tak, aby wartości $T(h_i)$ były nieujemne
rzeczywiste. Wtedy $|h_1+h_2|\le g_1+g_2$ i
$T(h_1+h_2)\approx M(g_1)+M(g_2)$.
Z addytywności i dodatniej jednorodności $M$ przedłuża
się jednoznacznie do dodatniego rzeczywistego funkcjonału
na $C_c(X;\mathbb R)$: $M(f)=M(f^+)-M(f^-)$.
Ponieważ $M(g)\le\|T\|\|g\|_\infty$, przedłużony
funkcjonał $M$ jest ograniczony w normie supremum
i rozciąga się przez gęstość na $C_0(X;\mathbb R)$.
Twierdzenie~\ref{thm:riesz-c0-positive-21} daje
skończoną miarę Radona $\sigma$ taką, że
$M(g)=\int_Xg\,\mathrm d\sigma$ dla $g\ge0$.
W szczególności dla $h\in C_c(X;\mathbb C)$ zachodzi
\begin{equation}\label{eq:riesz-domination-21}
 |T(h)|\le \int_X|h(x)|\,\mathrm d\sigma(x).
\end{equation}
Przestrzeń $C_c(X;\mathbb C)$ jest gęsta w
$L^1(\sigma;\mathbb C)$: przybliżamy funkcje charakterystyczne
zbiorów borelowskich przez funkcje odcinające między
zwartym podzbiorem i otwartym nadzbiorem o małej
różnicy miar, a następnie przybliżamy funkcje proste.
Zatem $T$ jednoznacznie przedłuża się do funkcjonału
na $L^1(\sigma)$ o normie co najwyżej $1$.
Dualność $L^1$--$L^\infty$ dla skończonej miary
(\ref{ch:dualnosc-lp}) daje $w\in L^\infty(\sigma)$,
$|w|\le1$ $\sigma$-prawie wszędzie, takie że
$T(h)=\int_X h w\,\mathrm d\sigma$.
Kładziemy $\nu=w\sigma$; jest to skończona zespolona
miara Radona, ponieważ $|\nu|\le\sigma$ i
regularność przechodzi na miary zdominowane.
Reprezentacja obowiązuje dla $C_c$, a przez gęstość
także dla $C_0$.

Pozostaje identyczność norm. Nierówność
$\|T\|\le|\nu|(X)$ jest bezpośrednia. Dla odwrotnej
wybieramy mierzalną funkcję fazową $u$ o $|u|\le1$
taką, że $u\,\mathrm d\nu=\mathrm d|\nu|$.
Z regularności $|\nu|$ i aproksymacji funkcji
mierzalnych przez $C_c$ w $L^1(|\nu|)$ istnieją
$f_n\in C_c(X;\mathbb C)$, $\|f_n\|_\infty\le1$,
$f_n\to u$ w $L^1(|\nu|)$. Ograniczenie normy
zapewnia obcięcie aproksymantów do domkniętego
dysku jednostkowego, które nie zwiększa błędu
punktowego względem $u$. Stąd
$T(f_n)\to\int_Xu\,\mathrm d\nu=|\nu|(X)$,
więc $\|T\|\ge|\nu|(X)$.
Jednoznaczność $\nu$ wynika z jednoznaczności dla
części rzeczywistej i urojonej po rozdzieleniu miar
oraz z regularności; szczegółowo równość całek na
$C_c$ implikuje równość każdej części miar ze znakiem,
a to wynika z twierdzenia~\ref{thm:uniqueness-cc-20}
po zastosowaniu rozkładu Jordana.
\end{proof}

\begin{przyklad}[Funkcjonał ewaluacji]
Dla ustalonego $x_0\in X$ funkcjonał $T(f)=f(x_0)$
na $C_0(X)$ jest ciągły i ma normę $1$.
Jego miarą reprezentującą jest $\delta_{x_0}$.
\end{przyklad}

\begin{przyklad}[Miara nieskończona i funkcjonał lokalny]
Dla $X=\mathbb R$ i $L(f)=\int_{\mathbb R} f(x)
\,\mathrm d\lambda_1(x)$, $f\in C_c(\mathbb R)$,
otrzymujemy dodatni funkcjonał nieograniczony w normie
supremum na całym $C_c$: funkcje odcinające równe $1$
na $[-n,n]$ mają normę supremum $1$, ale całki
co najmniej $2n$. Dlatego $L$ nie przedłuża się
ciągle do $C_0(\mathbb R)$.
\end{przyklad}

\begin{wniosek}[Porządek dodatnich funkcjonałów i miar]
\label{cor:riesz-order-21}
Niech $X$ będzie lokalnie zwartą ośrodkową przestrzenią
metryczną, a $L_1,L_2:C_c(X;\mathbb R)\to\mathbb R$
dodatnimi funkcjonałami liniowymi, reprezentowanymi
odpowiednio przez miary Radona $\mu_1,\mu_2$.
Wówczas następujące warunki są równoważne:
\begin{enumerate}[label=(\roman*)]
\item $L_1(f)\le L_2(f)$ dla każdego nieujemnego $f\in C_c(X)$;
\item $\mu_1(E)\le\mu_2(E)$ dla każdego zbioru borelowskiego
$E\subseteq X$.
\end{enumerate}
\end{wniosek}
\begin{proof}
Implikacja (ii)$\Rightarrow$(i) wynika z monotoniczności całki.
Załóżmy (i). Dla otwartego $U\subseteq X$ zachodzi
\[
 \mu_1(U)=\sup\{L_1(f):f\in C_c(X),\ 0\le f\le1,
 \ \operatorname{supp}f\subseteq U\}.
\]
Wzór ten udowodniono w rozdziale~\ref{ch:radon-measures-20};
obejmuje także przypadek $\mu_1(U)=+\infty$.
Nierówność dla funkcjonałów daje
$\mu_1(U)\le\mu_2(U)$. Dla dowolnego borelowskiego
$E$ z regularności zewnętrznej otrzymujemy
$\mu_1(E)=\inf_{U\supseteq E}\mu_1(U)
\le\inf_{U\supseteq E}\mu_2(U)=\mu_2(E)$.
\end{proof}

\begin{uwaga}[Gdzie użyliśmy twierdzenia Tichonowa?]
W konstrukcji na zwartej przestrzeni metrycznej nie
potrzebowaliśmy twierdzenia Tichonowa: miarę zewnętrzną
zbudowaliśmy bezpośrednio z dodatniego funkcjonału.
W rozdziale~\ref{ch:haar-23}, poświęconym mierze Haara,
pojawi się natomiast zwartość produktu nieprzeliczalnie
wielu przedziałów.
Nie należy jej mylić ze zwartą metryczną przestrzenią,
dla której wystarczy argument diagonalny.
\end{uwaga}

\section{Konsekwencje: oddzielenie miar i problem momentów}
\begin{wniosek}[Miary odróżniane przez funkcje ciągłe]
\label{cor:measures-separate-21}
Niech $\mu,\nu$ będą miarami Radona na lokalnie zwartej
ośrodkowej przestrzeni metrycznej $X$. Wtedy
$\mu=\nu$ wtedy i tylko wtedy, gdy
$\int_Xf\,\mathrm d\mu=\int_Xf\,\mathrm d\nu$
dla wszystkich $f\in C_c(X;\mathbb R)$.
\end{wniosek}
\begin{proof}
Implikacja w prawo jest oczywista, a odwrotna to
\ref{thm:uniqueness-cc-20}.
\end{proof}

\begin{uwaga}[Problem momentów jako ostrzeżenie]
Jeżeli na $\mathbb R$ znamy całki
$\int_{\mathbb R}x^k\,\mathrm d\mu(x)$ dla wszystkich
$k\ge0$, nie zawsze wyznacza to jednoznacznie miarę.
Funkcje wielomianowe nie należą do $C_0(\mathbb R)$
i nie ma automatycznej gęstości wielomianów w normie
supremum na całej prostej. Na zwartym przedziale
sytuacja jest inna: twierdzenie Stone'a--Weierstrassa
przybliża dowolną funkcję ciągłą wielomianami,
więc momenty jednoznacznie wyznaczają skończoną miarę.
\end{uwaga}

\section{Podsumowanie i przygotowanie do zbieżności miar}
Reprezentacja dodatnich funkcjonałów pozwala przejść
od analizy na $C_c(X)$ do teorii miar Radona. Dla
funkcjonałów na $C_0(X)$ skończoność normy odpowiada
skończoności masy, zaś dla zespolonych funkcjonałów
norma jest wariacją całkowitą. W rozdziale~\ref{ch:weak-measures-22}
zastosujemy zwarty przypadek twierdzenia do wykazania
ciągowej zwartości przestrzeni prawdopodobieństw na
zwartej przestrzeni metrycznej.

% END chapters/pl/measure_theory/21-rozdzial.tex


% BEGIN chapters/pl/measure_theory/22-rozdzial.tex
% ROZDZIAŁ 22 — ZBIEŻNOŚĆ MIAR I PROCHOROW
\chapter[Zbieżność miar i twierdzenie Prochorowa]{Zbieżność miar, twierdzenie Portmanteau i zwartość Prochorowa}
\label{ch:weak-measures-22}

\section{Trzy topologie, których nie należy utożsamiać}
Niech $(X,d)$ będzie przestrzenią metryczną,
$\mathcal B(X)$ jej $\sigma$-algebrą borelowską,
a $\mathcal P(X)$ zbiorem miar probabilistycznych na $X$.
Przestrzeń $C_b(X)$ składa się z ograniczonych rzeczywistych
funkcji ciągłych. Jeśli $X$ jest dodatkowo lokalnie zwarta,
$C_c(X)$ oznacza przestrzeń zdefiniowaną w
rozdziale~\ref{ch:radon-measures-20}, a $C_0(X)$
przestrzeń wprowadzoną w rozdziale~\ref{ch:riesz-markov-21}.

\begin{definicja}[Zbieżność wąska]
\label{def:narrow-measures-22}
Ciąg $(\mu_n)\subseteq\mathcal P(X)$ jest
\emph{wąsko zbieżny} do $\mu\in\mathcal P(X)$,
co zapisujemy $\mu_n\Rightarrow\mu$, jeśli
\begin{equation}\label{eq:narrow-22}
 \int_X f(x)\,\mathrm d\mu_n(x)
 \longrightarrow\int_X f(x)\,\mathrm d\mu(x)
 \quad\text{dla każdej }f\in C_b(X).
\end{equation}
\end{definicja}

\begin{definicja}[Zbieżność słaba-* i mglista]
\label{def:vague-weakstar-22}
Niech $X$ będzie lokalnie zwartą ośrodkową przestrzenią
metryczną, a $\mu_n,\mu$ miarami Radona.
Nazywamy zbieżność $\mu_n\to\mu$ \emph{mglistą},
gdy $\int_Xf\,\mathrm d\mu_n\to\int_Xf\,\mathrm d\mu$
dla każdego $f\in C_c(X)$.
Jeżeli miary są skończone, zbieżność
$\sigma(C_0(X)^*,C_0(X))$ (\emph{słaba-*})
oznacza zbieżność całek dla wszystkich $f\in C_0(X)$.
W przypadku miar probabilistycznych zbieżność wąska
implikuje obie te zbieżności. Na przestrzeni niezwartej
implikacje odwrotne nie muszą zachodzić.
\end{definicja}

\begin{przyklad}[Ucieczka masy]
\label{ex:mass-escape-22}
Niech $X=\mathbb R$, $\mu_n=\delta_n$.
Dla każdego $f\in C_0(\mathbb R)$ zachodzi $f(n)\to0$,
zatem $\mu_n\to0$ słabo-* w dualnej przestrzeni
$C_0(\mathbb R)^*$. Także $\mu_n\to0$ mgliście.
Jednak $\mu_n$ nie zbiega wąsko do żadnej miary
probabilistycznej: stała funkcja $1\in C_b(\mathbb R)$
ma całkę równą $1$ względem każdej $\mu_n$.
W szczególności określenie „zbieżność słaba miar” bez
wskazania klasy funkcji testowych bywa dwuznaczne.
\end{przyklad}

\section{Twierdzenie Portmanteau}
\begin{definicja}[Zbiór ciągłości miary]
\label{def:continuity-set-22}
Niech $\mu\in\mathcal P(X)$. Zbiór borelowski
$A\subseteq X$ nazywamy \emph{zbiorem ciągłości miary}
$\mu$, jeżeli $\mu(\partial A)=0$, gdzie
$\partial A=\overline A\setminus A^\circ$ jest jego brzegiem.
\end{definicja}

\begin{twierdzenie}[Portmanteau]
\label{thm:portmanteau-22}
Niech $(X,d)$ będzie dowolną przestrzenią metryczną,
a $\mu_n,\mu\in\mathcal P(X)$. Następujące warunki
są równoważne:
\begin{enumerate}[label=(\roman*)]
\item $\mu_n\Rightarrow\mu$;
\item dla każdej funkcji ograniczonej, dolnie półciągłej
$f:X\to\mathbb R$ zachodzi
$\int_Xf\,\mathrm d\mu\le\liminf_n\int_Xf\,\mathrm d\mu_n$;
\item dla każdego otwartego $G\subseteq X$ zachodzi
$\mu(G)\le\liminf_n\mu_n(G)$;
\item dla każdego domkniętego $F\subseteq X$ zachodzi
$\limsup_n\mu_n(F)\le\mu(F)$;
\item dla każdego zbioru ciągłości $A$ miary $\mu$
zachodzi $\mu_n(A)\to\mu(A)$.
\end{enumerate}
W warunku (ii) można także dopuścić dolnie półciągłe
funkcje ograniczone z dołu, przy całkach w $(-\infty,+\infty]$.
\end{twierdzenie}
\begin{proof}
\emph{(i)$\Rightarrow$(iii).}
Dla otwartego $G$ określamy
$f_k(x)=\min\{1,kd(x,X\setminus G)\}$, gdy
$G\neq X$, oraz $f_k\equiv1$, gdy $G=X$.
Funkcje $f_k$ są ciągłe, $0\le f_k\le\chi_G$,
a $f_k\uparrow\chi_G$. Zatem dla każdego $k$
\[
\liminf_n\mu_n(G)\ge\lim_n\int_X f_k\,\mathrm d\mu_n
=\int_X f_k\,\mathrm d\mu.
\]
Z twierdzenia Beppo Leviego, po przejściu $k\to\infty$,
otrzymujemy (iii).

\emph{(iii)$\Leftrightarrow$(iv).}
Dla $F$ domkniętego zbiór $X\setminus F$ jest otwarty,
a $\mu_n(X\setminus F)=1-\mu_n(F)$ i podobnie
dla $\mu$. Pozwala to zamienić dolną granicę na górną.

\emph{(iii)$\Rightarrow$(ii).}
Najpierw niech $f$ będzie dolnie półciągła i
$0\le f\le M$. Dla każdego $t\in[0,M]$ zbiór
$G_t=\{x:f(x)>t\}$ jest otwarty. Z warunku (iii)
$\mu(G_t)\le\liminf_n\mu_n(G_t)$.
Z reprezentacji warstwowej całki
$\int_X f\,\mathrm d\mu=\int_{[0,M]}\mu(G_t)
\,\mathrm d\lambda_1(t)$, twierdzenia Tonellego
i lematu Fatou otrzymujemy
\[
\int_X f\,\mathrm d\mu
\le\int_{[0,M]}\liminf_n\mu_n(G_t)\,\mathrm d\lambda_1(t)
\le\liminf_n\int_Xf\,\mathrm d\mu_n.
\]
Jeśli $f\ge a$ i nie jest ograniczona z góry,
stosujemy poprzedni argument do
$f_k=\min\{f-a,k\}$, a następnie Beppo Leviego;
odjęta stała $a$ ma ten sam całkowy wkład dla
wszystkich miar probabilistycznych.

\emph{(ii)$\Rightarrow$(i).}
Dla dowolnej $h\in C_b(X)$ zarówno $h$, jak i $-h$
są dolnie półciągłe i ograniczone z dołu. Stąd
$\int h\,\mathrm d\mu\le\liminf_n\int h\,\mathrm d\mu_n$
oraz $-\int h\,\mathrm d\mu\le
-\limsup_n\int h\,\mathrm d\mu_n$.
Obie nierówności wymuszają zbieżność całek.

\emph{(iii), (iv)$\Rightarrow$(v).}
Dla $A$ spełniającego $\mu(\partial A)=0$ mamy
$A^\circ\subseteq A\subseteq\overline A$ oraz
$\mu(A^\circ)=\mu(A)=\mu(\overline A)$.
Wobec tego
\[
\mu(A)\le\liminf_n\mu_n(A^\circ)
\le\liminf_n\mu_n(A)
\le\limsup_n\mu_n(A)
\le\limsup_n\mu_n(\overline A)\le\mu(A).
\]

\emph{(v)$\Rightarrow$(i).}
Niech $h\in C_b(X)$, $\|h\|_\infty\le M$.
Miara obrazowa $h_*\mu$ na $[-M,M]$ ma co najwyżej
przeliczalnie wiele atomów: dla każdego $k$ liczba
atomów o masie co najmniej $1/k$ jest skończona.
Zatem można wybrać podział
$-M=t_0<\cdots<t_N=M$ o dostatecznie małych oczkach,
którego punkty wewnętrzne są nieatomowe dla $h_*\mu$.
Aby uniknąć atomów w końcach, wolno wybrać
$t_0<-M$ i $t_N>M$ oraz wymagać, by również one
były nieatomowe. Dla przedziałów $(t_{j-1},t_j]$
zbiory $A_j=\{t_{j-1}<h\le t_j\}$ mają
$\partial A_j\subseteq\{h=t_{j-1}\}\cup\{h=t_j\}$,
więc są zbiorami ciągłości miary. Z (v) wynika
$\mu_n(A_j)\to\mu(A_j)$. Aproksymując $h$
na każdym $A_j$ przez stałą wybraną z przedziału,
z błędem co najwyżej długość przedziału, otrzymujemy
\[
 \limsup_n\left|\int_Xh\,\mathrm d\mu_n-
 \int_Xh\,\mathrm d\mu\right|
 \le2\max_j(t_j-t_{j-1}).
\]
Wybierając dowolnie drobne podziały kończymy dowód.
\end{proof}

\begin{przyklad}[Dlaczego przedziały nie zawsze są zbiorami ciągłości?]
Dla $\mu_n=\delta_{1/n}$ i $\mu=\delta_0$ na
$\mathbb R$ mamy $\mu_n\Rightarrow\mu$.
Tymczasem $\mu_n(\{0\})=0$ i $\mu(\{0\})=1$.
Brzeg zbioru $\{0\}$ ma dodatnią miarę $\mu$,
więc Portmanteau nie gwarantuje zbieżności
wartości miar tego zbioru.
\end{przyklad}

\begin{wniosek}[Kiedy zbieżność mglista jest słabą-*?]
\label{cor:vague-c0-bounded-22}
Niech $X$ będzie lokalnie zwartą ośrodkową przestrzenią
metryczną i niech $\mu_n,\mu$ będą skończonymi dodatnimi
miarami Radona. Jeśli $\sup_n\mu_n(X)<\infty$, to
zbieżność mglista $\mu_n\to\mu$ jest równoważna zbieżności
$\sigma(C_0(X)^*,C_0(X))$.
Bez ograniczenia mas implikacja od zbieżności mglistej
do słabej-* może zawieść.
\end{wniosek}
\begin{proof}
Zbieżność słaba-* daje mglistą, gdyż $C_c(X)\subset C_0(X)$.
Załóżmy zbieżność mglistą oraz $C:=\sup_n\mu_n(X)<\infty$.
Niech $f\in C_0(X)$ i $\varepsilon>0$. Z gęstości
$C_c(X)$ w $C_0(X)$ (lemat~\ref{lem:density-cc-c0-21})
wybieramy $g\in C_c(X)$ takie, że
$\|f-g\|_\infty<\varepsilon$. Wówczas
\[
 \left|\int_X f\,\mathrm d\mu_n-\int_X f\,\mathrm d\mu\right|
 \le\left|\int_Xg\,\mathrm d\mu_n-\int_Xg\,\mathrm d\mu\right|
 +(C+\mu(X))\varepsilon.
\]
Pierwszy składnik dąży do zera, a $\varepsilon$ jest dowolne.

Aby wykazać konieczność kontroli mas, weźmy
$X=\mathbb R$, $\mu_n=n\delta_n$ i $\mu=0$.
Dla każdego $g\in C_c(\mathbb R)$ zachodzi
$\int g\,\mathrm d\mu_n=0$ dla dostatecznie dużych $n$,
więc zachodzi zbieżność mglista. Funkcja
$f(x)=(1+|x|)^{-1/2}$ należy do $C_0(\mathbb R)$,
a $\int f\,\mathrm d\mu_n=n/\sqrt{1+n}\to+\infty$.
Zbieżności słabej-* nie ma.
\end{proof}

\section{Zbieżność względem odległości w wariacji całkowitej}
\begin{twierdzenie}[Zbieżność w wariacji implikuje wąską]
\label{thm:tv-implies-narrow-22}
Jeśli $\mu_n,\mu\in\mathcal P(X)$ i
$d_{\mathrm{TV}}(\mu_n,\mu)\to0$, to
$\mu_n\Rightarrow\mu$.
\end{twierdzenie}
\begin{proof}
Dla dowolnej $h\in C_b(X)$, korzystając z normy
wariacji miary ze znakiem, mamy
\[
\left|\int_Xh\,\mathrm d\mu_n-\int_Xh\,\mathrm d\mu\right|
\le \|h\|_\infty\|\mu_n-\mu\|_{\mathrm{TV}}
=2\|h\|_\infty d_{\mathrm{TV}}(\mu_n,\mu)\to0.
\]
\end{proof}

\begin{uwaga}[Brak implikacji odwrotnej]
Przykład $\delta_{1/n}\Rightarrow\delta_0$ pokazuje,
że wąska zbieżność nie implikuje zbieżności w wariacji:
$d_{\mathrm{TV}}(\delta_{1/n},\delta_0)=1$.
\end{uwaga}

\section{Ciasność rodzin miar}
\begin{definicja}[Jednostajna ciasność]
\label{def:uniform-tight-22}
Niech $\mathcal F\subseteq\mathcal P(X)$.
Rodzina $\mathcal F$ jest \emph{jednostajnie ciasna},
jeżeli dla każdego $\varepsilon>0$ istnieje
zwarty zbiór $K\subseteq X$ spełniający
$\inf_{\mu\in\mathcal F}\mu(K)>1-\varepsilon$.
\end{definicja}

\begin{lemat}[Przybliżanie lipschitzowskie i przedłużanie]
\label{lem:lipschitz-extend-22}
Niech $(Y,D)$ będzie przestrzenią metryczną, $K\subseteq Y$
niepustym zbiorem zwartym, a $h:K\to\mathbb R$ funkcją
ciągłą z $\|h\|_\infty\le M$. Dla każdego $\eta>0$
istnieje ograniczona funkcja Lipschitzowska
$g:Y\to\mathbb R$ spełniająca
$\|g\|_\infty\le M$ oraz
$\sup_{x\in K}|g(x)-h(x)|<\eta$.
\end{lemat}
\begin{proof}
Funkcja $h$ jest jednostajnie ciągła na $K$.
Dla $L>0$ definiujemy
$h_L(x)=\inf_{z\in K}(h(z)+L D(x,z))$, $x\in K$.
Z nierówności trójkąta wynika
$|h_L(x)-h_L(y)|\le L D(x,y)$; ponadto
$h_L\le h$ i $-M\le h_L\le M$.
Wybierzmy $\delta>0$ takie, że
$|h(z)-h(x)|<\eta/2$ dla $D(x,z)<\delta$.
Jeżeli $L>2M/\delta$, to dla $D(x,z)\ge\delta$
mamy $h(z)+LD(x,z)\ge M\ge h(x)$.
Dla $D(x,z)<\delta$ natomiast
$h(z)+LD(x,z)\ge h(x)-\eta/2$.
Zatem $0\le h-h_L\le\eta/2$ na $K$.

Funkcję $h_L$ przedłużamy na $Y$ wzorem
$G(y)=\inf_{z\in K}(h_L(z)+L D(y,z))$.
Ponownie $G$ jest $L$-lipschitzowska. Na $K$
mamy $G=h_L$: nierówność $G\le h_L$ wynika
z podstawienia $z=y$, a odwrotna z własności
Lipschitza funkcji $h_L$. Funkcja
$g(y)=\max\{-M,\min\{M,G(y)\}\}$ jest Lipschitzowska,
ma normę co najwyżej $M$ i na $K$ jest równa $h_L$.
Wobec tego $\sup_K|g-h|\le\eta/2<\eta$.
\end{proof}

\begin{lemat}[Zwartość prawdopodobieństw na przestrzeni zwartej]
\label{lem:compact-prob-22}
Jeśli $Y$ jest zwartą przestrzenią metryczną,
to każdy ciąg w $\mathcal P(Y)$ ma podciąg wąsko
zbieżny do pewnego elementu $\mathcal P(Y)$.
\end{lemat}
\begin{proof}
Przestrzeń $C(Y)$ jest separowalna w normie supremum.
Aby to udowodnić bez odwołania do twierdzenia
Stone'a--Weierstrassa, wybierzmy przeliczalny zbiór
gęsty $(y_j)$ i rozważmy przeliczalną rodzinę
funkcji postaci
$F(x)=\min_{1\le i\le N}\{q_i+L d(x,y_{j_i})\}$,
gdzie $N,j_i\in\mathbb N$, $q_i\in\mathbb Q$,
$L\in\mathbb Q_{>0}$.
Pokażemy, że jest ona gęsta. Dla dowolnej
$h\in C(Y)$ i $\varepsilon>0$ lemat
\ref{lem:lipschitz-extend-22} (z $K=Y$) daje
funkcję $L$-lipschitzowską $v$ z
$\|h-v\|_\infty<\varepsilon/3$.
Zwiększając $L$, możemy przyjąć go wymiernym.
Wybieramy skończoną $\delta$-siatkę
$y_{j_1},\ldots,y_{j_N}$ oraz wymierne $q_i$
bliskie $v(y_{j_i})$, tak by
$|q_i-v(y_{j_i})|<\varepsilon/6$ i
$2L\delta<\varepsilon/6$.
Z lipschitzowskości wynika
$F(x)\ge v(x)-\varepsilon/6$ oraz
$F(x)\le v(x)+\varepsilon/6+2L\delta$.
W konsekwencji $\|F-h\|_\infty<\varepsilon$.
Niech $(f_j)$ będzie taką gęstą
rodziną. Dla danego ciągu $(\mu_n)$ każdy ciąg
$\int_Y f_j\,\mathrm d\mu_n$ jest ograniczona.
Metoda diagonalna dostarcza podciągu $(\mu_{n_k})$,
dla którego wszystkie te całki mają granice.
Ponieważ
$|\int h\,\mathrm d\mu_{n_k}-\int f_j\,\mathrm d\mu_{n_k}|
\le\|h-f_j\|_\infty$, te granice rozszerzają się
jednoznacznie na ciągły dodatni funkcjonał $L$ na
$C(Y)$ o normie $1$ i $L(\chi_Y)=1$.
Na mocy twierdzenia~\ref{thm:compact-riesz-21}
istnieje prawdopodobieństwo $\mu$ reprezentujące $L$.
Definicja $L$ daje $\mu_{n_k}\Rightarrow\mu$.
\end{proof}

\begin{lemat}[Włożenie separowalnej przestrzeni metrycznej
w zwartą przestrzeń metryczną]
\label{lem:hilbert-cube-22}
Dla każdej separowalnej przestrzeni metrycznej $(X,d)$
istnieje zwarta przestrzeń metryczna $Y$ i homeomorfizm
$\iota:X\to\iota(X)\subseteq Y$.
\end{lemat}
\begin{proof}
Niech $(x_j)$ będzie przeliczalnym zbiorem gęstym.
Definiujemy
\[
\iota(x)=\left(\frac{d(x,x_j)}{1+d(x,x_j)}\right)_{j\ge1}
\in[0,1]^{\mathbb N}.
\]
Przestrzeń kostki Hilberta $[0,1]^{\mathbb N}$ z metryką
$D(a,b)=\sum_{j\ge1}2^{-j}|a_j-b_j|$ jest zwarta:
metoda diagonalna zapewnia zbieżny podciąg współrzędnych,
a oszacowanie ogona sumy daje zbieżność w $D$.
Każda współrzędna $\iota$ jest ciągła, więc
$\iota$ jest ciągła. Jeśli $\iota(x)=\iota(y)$,
to $d(x,x_j)=d(y,x_j)$ dla wszystkich $j$;
przez gęstość wynika $x=y$.
Aby wykazać ciągłość odwrotności na obrazie, załóżmy,
że $\iota(x_n)\to\iota(x)$, i ustalmy $\varepsilon>0$.
Wybierzmy $x_j$ z $d(x,x_j)<\varepsilon/3$.
Dla dużych $n$ współrzędna $j$ daje
$|d(x_n,x_j)-d(x,x_j)|<\varepsilon/3$,
zatem $d(x_n,x)\le d(x_n,x_j)+d(x_j,x)<\varepsilon$.
Wystarczy przyjąć $Y=\overline{\iota(X)}$ w kostce.
\end{proof}

\begin{twierdzenie}[Prochorow: kierunek ciasność $\Rightarrow$ zwartość]
\label{thm:prokhorov-forward-22}
Niech $X$ będzie separowalną przestrzenią metryczną,
a $\mathcal F\subseteq\mathcal P(X)$ rodziną jednostajnie
ciasną. Wtedy każdy ciąg w $\mathcal F$ ma podciąg
wąsko zbieżny do prawdopodobieństwa na $X$.
\end{twierdzenie}
\begin{proof}
Niech $(\mu_n)\subseteq\mathcal F$.
Z jednostajnej ciasności wybieramy rosnący ciąg
zwartych zbiorów $K_m\subseteq X$ taki, że
\begin{equation}\label{eq:tight-compacts-22}
 \inf_n\mu_n(K_m)\ge1-2^{-m},\qquad m\ge1.
\end{equation}
Stosujemy włożenie $\iota:X\to Y$ z lematu
\ref{lem:hilbert-cube-22}. Miary
$\nu_n=\iota_*\mu_n$ są probabilistyczne na
zwartej metrycznej $Y$. Z lematu
\ref{lem:compact-prob-22} wybieramy podciąg,
ponownie oznaczany $(\nu_n)$, wąsko zbieżny
do $\nu\in\mathcal P(Y)$.
Dla każdego $m$ zbiór $\iota(K_m)$ jest zwarty,
a zatem domknięty w $Y$. Z Portmanteau
\[
\nu(\iota(K_m))\ge\limsup_n\nu_n(\iota(K_m))
\ge1-2^{-m}.
\]
Wobec rosnącości $K_m$ zbiór
$C=\bigcup_m\iota(K_m)$ jest borelowski w $Y$,
a $\nu(C)=1$.

Definiujemy teraz prawdopodobieństwo $\mu$ na $X$.
Dla $A\in\mathcal B(X)$ kładziemy
\begin{equation}\label{eq:pullback-limit-22}
\mu(A)=\nu\bigl(\iota(A\cap\bigcup_mK_m)\bigr).
\end{equation}
Ta definicja jest poprawna: $A\cap K_m$ jest
borelowski w zwartej przestrzeni $K_m$,
a $\iota|_{K_m}$ jest homeomorfizmem na zwarty
podzbiór $Y$. Zatem $\iota(A\cap K_m)$ jest
borelowski w $Y$, podobnie jak ich przeliczalna suma.
Iniektywność $\iota$ i addytywność $\nu$ dowodzą,
że $\mu$ jest miarą; ponadto $\mu(X)=\nu(C)=1$
i $\mu(K_m)=\nu(\iota(K_m))$.

Niech $f\in C_b(X)$, $M=\|f\|_\infty$.
Dla ustalonego $m$ funkcja $f\circ\iota^{-1}$
jest ciągła na zwartym $\iota(K_m)$.
Lemat~\ref{lem:lipschitz-extend-22}, zastosowany
z dowolnym $\eta>0$, daje funkcję
$g_{m,\eta}\in C_b(Y)$, $\|g_{m,\eta}\|_\infty\le M$,
która na $\iota(K_m)$ różni się od
$f\circ\iota^{-1}$ o mniej niż $\eta$.
Stąd, z \eqref{eq:tight-compacts-22},
\[
\left|\int_X f\,\mathrm d\mu_n-
\int_Y g_{m,\eta}\,\mathrm d\nu_n\right|\le2M2^{-m}+\eta,
\]
a przez analogiczny rachunek
$|\int_Xf\,\mathrm d\mu-\int_Yg_{m,\eta}\,\mathrm d\nu|
\le2M2^{-m}+\eta$.
Dla ustalonych $m,\eta$ granica całek z $g_{m,\eta}$ wynika
ze zbieżności wąskiej na $Y$; dlatego
\[
\limsup_n\left|\int_X f\,\mathrm d\mu_n-
\int_X f\,\mathrm d\mu\right|\le4M2^{-m}+2\eta.
\]
Przechodząc najpierw $\eta\to0$, a następnie
$m\to\infty$, kończymy dowód.
\end{proof}

\begin{lemat}[Jednostajna ciasność ciągu wąsko zbieżnego]
\label{lem:weak-sequence-tight-22}
Jeśli $X$ jest polską przestrzenią metryczną,
$\mu_n\Rightarrow\mu$ i $\mu_n,\mu\in\mathcal P(X)$,
to rodzina $\{\mu,\mu_1,\mu_2,\ldots\}$ jest
jednostajnie ciasna.
\end{lemat}
\begin{proof}
Niech $d$ będzie zupełną, zgodną metryką, a $(x_j)$
przeliczalnym zbiorem gęstym. Ustalmy $\varepsilon>0$.
Dla każdego $m$ przeliczalna rodzina kul otwartych
$B(x_j,2^{-m})$ pokrywa $X$. Istnieje skończona
suma tych kul $O_m$, dla której
$\mu(O_m)>1-\varepsilon2^{-m-2}$.
Z otwartej wersji Portmanteau wynika, że
$\mu_n(O_m)>1-\varepsilon2^{-m-1}$ dla wszystkich
$n$ od pewnego $N_m$.
Powiększając skończoną rodzinę kul w $O_m$,
możemy zapewnić tę samą nierówność dla
$\mu_1,\ldots,\mu_{N_m-1}$; pozostaje prawdziwa
także dla wszystkich późniejszych miar.
Niech $F_m$ będzie sumą domkniętych kul
$\overline B(x_j,2^{-m})$ o środkach użytych w $O_m$.
Kładziemy $K=\bigcap_mF_m$. Zbiór $K$ jest domknięty,
a dla każdego $m$ mieści się w skończonej sumie kul
o promieniu $2^{-m}$, jest zatem całkowicie ograniczony.
Zupełność $X$ daje zwartość $K$. Ponadto
$\mu_n(X\setminus K)\le\sum_{m\ge1}\mu_n(X\setminus F_m)
<\sum_{m\ge1}\varepsilon2^{-m-1}<\varepsilon$
dla wszystkich $n$. Podobne oszacowanie zachodzi dla $\mu$.
\end{proof}

\begin{twierdzenie}[Prochorow na przestrzeni polskiej]
\label{thm:prokhorov-22}
Niech $X$ będzie przestrzenią polską, a
$\mathcal F\subseteq\mathcal P(X)$.
Następujące warunki są równoważne:
\begin{enumerate}[label=(\roman*)]
\item $\mathcal F$ jest jednostajnie ciasna;
\item każdy ciąg w $\mathcal F$ ma podciąg wąsko
zbieżny do pewnego prawdopodobieństwa na $X$.
\end{enumerate}
Warunek (ii) oznacza względną zwartość
ciągową dla topologii zbieżności wąskiej.
\end{twierdzenie}
\begin{proof}
Implikację (i)$\Rightarrow$(ii) udowodniliśmy w twierdzeniu
\ref{thm:prokhorov-forward-22}. Dla odwrotnej
ustalmy $\varepsilon>0$, zgodną metrykę zupełną $d$
i przeliczalny zbiór gęsty $\{x_j:j\ge1\}$.
Dla każdego $m\ge1$ oznaczmy
$V_{m,k}=\bigcup_{j=1}^k B(x_j,2^{-m})$.
Dla ustalonego $m$ zbiory $V_{m,k}$ rosną do $X$.
Wykażemy, że istnieje $k_m$ takie, że
\begin{equation}\label{eq:prokhorov-uniform-balls-22}
\inf_{\rho\in\mathcal F}\rho(V_{m,k_m})
>1-\varepsilon2^{-m-1}.
\end{equation}
Gdyby tak nie było, dla każdego $k$ wybralibyśmy
$\rho_k\in\mathcal F$ z
$\rho_k(V_{m,k})\le1-\varepsilon2^{-m-1}$.
Z założenia (ii) ciąg ten ma podciąg
$\rho_{k_l}\Rightarrow\rho$ z granicą probabilistyczną.
Ponieważ $V_{m,k}\uparrow X$, istnieje $k_0$ takie,
że $\rho(V_{m,k_0})>1-\varepsilon2^{-m-1}$.
Portmanteau dla otwartego $V_{m,k_0}$ daje
$\rho_{k_l}(V_{m,k_0})>1-\varepsilon2^{-m-1}$
dla dostatecznie dużych $l$. Ponieważ $k_l\ge k_0$,
monotoniczność wymusza wówczas
$\rho_{k_l}(V_{m,k_l})>1-\varepsilon2^{-m-1}$,
sprzecznie z konstrukcją $\rho_{k_l}$.
Dowiedliśmy \eqref{eq:prokhorov-uniform-balls-22}
bez odwoływania się do metryzowalności
przestrzeni prawdopodobieństw.

Niech $F_m=\bigcup_{j\le k_m}\overline B(x_j,2^{-m})$
i $K=\bigcap_{m\ge1}F_m$. Zbiór $K$ jest domknięty,
a dla każdego $m$ mieści się w skończonej sumie kul
o promieniu $2^{-m}$, więc jest całkowicie ograniczony.
Zupełność $X$ zapewnia zwartość $K$.
Z \eqref{eq:prokhorov-uniform-balls-22} dla wszystkich
$\rho\in\mathcal F$ uzyskujemy
\[
 \rho(X\setminus K)\le\sum_{m\ge1}\rho(X\setminus F_m)
 <\sum_{m\ge1}\varepsilon2^{-m-1}<\varepsilon.
\]
Jest to jednostajna ciasność $\mathcal F$.
\end{proof}

\begin{uwaga}[Metryzowalność a dowód Prochorowa]
Na przestrzeni polskiej topologia zbieżności wąskiej
na $\mathcal P(X)$ jest metryzowalna, na przykład przez
odległość opartą na ograniczonych funkcjach Lipschitzowskich.
Nie było jednak potrzeby przyjmowania tego wyniku
w dowodzie twierdzenia~\ref{thm:prokhorov-22}:
własność wyrażona przez podciągi i argument ze
wzrastającymi skończonymi pokryciami kulowymi
dały bezpośrednią konstrukcję wspólnego zwartego zbioru.
\end{uwaga}

\begin{wniosek}[Ciasność z oszacowania momentów]
\label{cor:coercive-tight-22}
Niech $X$ będzie przestrzenią polską i
$V:X\to[0,+\infty]$ funkcją dolnie półciągłą,
której zbiory podpoziomicowe
$K_R=\{x\in X:V(x)\le R\}$ są zwarte dla każdego
$R<\infty$. Jeśli rodzina $\mathcal F\subseteq\mathcal P(X)$
spełnia
$\sup_{\rho\in\mathcal F}\int_X V(x)\,\mathrm d\rho(x)
\le C<\infty$, to jest jednostajnie ciasna.
W szczególności na $\mathbb R^d$ wystarcza jednostajne
ograniczenie dodatniego momentu rzędu $p>0$.
\end{wniosek}
\begin{proof}
Dla $R>0$ nierówność Czebyszewa daje
\[
 \rho(X\setminus K_R)=\rho(\{V>R\})
 \le R^{-1}\int_X V(x)\,\mathrm d\rho(x)\le C/R.
\]
Wybierając $R>C/\varepsilon$ otrzymujemy
$\inf_{\rho\in\mathcal F}\rho(K_R)>1-\varepsilon$.
Na $\mathbb R^d$ funkcja $V(x)=|x|^p$ jest ciągła,
a jej zbiory podpoziomicowe są zwarte.
\end{proof}

\section{Następstwa i kryteria praktyczne}
\begin{wniosek}[Zbieżność mglista bez ucieczki masy]
\label{cor:vague-narrow-22}
Niech $X$ będzie lokalnie zwartą ośrodkową
przestrzenią metryczną. Jeśli $\mu_n,\mu\in\mathcal P(X)$,
$\mu_n$ zbiega mgliście do $\mu$, to
$\mu_n\Rightarrow\mu$.
\end{wniosek}
\begin{proof}
Ustalmy $\varepsilon>0$. Z rozdziału
\ref{ch:radon-measures-20} miara $\mu$ jest ciasna.
Wybieramy zwarty $K$ o $\mu(K)>1-\varepsilon$
oraz $\chi\in C_c(X)$ z $0\le\chi\le1$,
$\chi|_K=1$. Wówczas
$\int_X\chi\,\mathrm d\mu>1-\varepsilon$,
a ze zbieżności mglistej
$\int_X\chi\,\mathrm d\mu_n>1-2\varepsilon$
dla dużych $n$. Dla $h\in C_b(X)$
funkcja $h\chi\in C_c(X)$, a więc
\[
\left|\int_Xh\,\mathrm d\mu_n-
\int_Xh\chi\,\mathrm d\mu_n\right|
\le\|h\|_\infty\int_X(1-\chi)\,\mathrm d\mu_n
<2\varepsilon\|h\|_\infty.
\]
Analogiczne oszacowanie zachodzi dla $\mu$,
z błędem $\varepsilon\|h\|_\infty$.
Zbieżność całek z $h\chi$ i dowolność $\varepsilon$
kończą dowód.
\end{proof}

\begin{przyklad}[Rozkłady jednostajnie skupione na zwartym przedziale]
Niech $X=\mathbb R$, a $\mu_n$ dowolnymi
prawdopodobieństwami spełniającymi
$\mu_n([-R,R])=1$ dla stałego $R<\infty$.
Rodzina $\{\mu_n\}$ jest jednostajnie ciasna.
Z twierdzenia Prochorowa każdy jej ciąg ma podciąg
wąsko zbieżny. Nie wynika stąd zbieżność całego ciągu:
na przykład miary $\delta_{-1},\delta_1,\delta_{-1},\ldots$
mają dwa różne punkty skupienia.
\end{przyklad}

\begin{przyklad}[Brak ciasności]
Rodzina $\{\delta_n:n\ge1\}\subseteq\mathcal P(\mathbb R)$
nie jest jednostajnie ciasna: dla każdego zwartego
$K\subseteq\mathbb R$ pewne $n$ nie należy do $K$.
Nie istnieje podciąg zbieżny wąsko do prawdopodobieństwa.
\end{przyklad}

\begin{uwaga}[Gdzie pojawi się ta teoria?]
Zbieżność wąska, ciasność i twierdzenie Prochorowa
stanowią narzędzia w teorii procesów stochastycznych,
rachunku prawdopodobieństwa, analizie wariacyjnej
oraz w zagadnieniach transportu optymalnego.
Dalszy rozdział skryptu poświęcimy mierze Haara
na grupach lokalnie zwartych, gdzie niezmienniczość
miary dostarczy kolejnego sposobu jej scharakteryzowania.
\end{uwaga}

% END chapters/pl/measure_theory/22-rozdzial.tex


% BEGIN chapters/pl/measure_theory/23-rozdzial.tex
% ROZDZIAŁ 23 — MIARA HAARA
\chapter{Miara Haara na grupach lokalnie zwartych}
\label{ch:haar-23}

\section{Grupy topologiczne, translacje i miary niezmiennicze}
\begin{definicja}[Grupa topologiczna]
\label{def:topological-group-23}
Grupę $G$, wyposażoną w topologię Hausdorffa, nazywamy
\emph{grupą topologiczną}, gdy odwzorowania $(x,y)\mapsto xy$
oraz $x\mapsto x^{-1}$ są ciągłe. Element neutralny oznaczamy $e$.
Grupa jest \emph{lokalnie zwarta}, gdy każdy punkt ma otoczenie
o zwartej domkniętej powłoce. Piszemy $xA=\{xa:a\in A\}$
i $Ax=\{ax:a\in A\}$.
\end{definicja}

\begin{uwaga}[Dlaczego lokalna zwartość?]
W przestrzeni $\mathbb R^d$ miara Lebesgue'a jest niezmiennicza
względem przesunięć. Dla grup nieprzemiennych należy rozróżniać
translacje \emph{lewe} i \emph{prawe}. Lokalna zwartość gwarantuje,
że zbiory zwarte można pokrywać skończenie wieloma przesunięciami
małego otwartego otoczenia $e$; na tej własności oprzemy konstrukcję.
\end{uwaga}

\begin{definicja}[Lewostronna i prawostronna miara Haara]
\label{def:haar-23}
Niech $G$ będzie lokalnie zwartą grupą Hausdorffa.
\emph{Lewostronną miarą Haara} nazywamy niezerową miarę Radona
$\mu$ na $\mathcal B(G)$, w sensie definicji~\ref{def:radon-measure-20},
dla której $\mu(gE)=\mu(E)$ dla wszystkich $g\in G$
i $E\in\mathcal B(G)$. Analogicznie definiujemy
\emph{prawostronną miarę Haara} przez warunek
$\mu(Eg)=\mu(E)$. Zakres wymaganej regularności
przypominamy poniżej ze względu na brak założeń przeliczalności.
\end{definicja}

\begin{twierdzenie}[Haara: istnienie i jednoznaczność]
\label{thm:haar-main-23}
Na każdej lokalnie zwartej grupie Hausdorffa istnieje
lewostronna miara Haara. Każde dwie takie miary różnią się
dodatnim czynnikiem stałym. To samo dotyczy prawostronnych miar
Haara.
\end{twierdzenie}

\begin{uwaga}[Zakres twierdzenia i regularność]
W twierdzeniu~\ref{thm:haar-main-23} \emph{nie} zakładamy ani
metryzowalności, ani $\sigma$-zwartości, ani drugiej przeliczalności.
W ogólnej przestrzeni lokalnie zwartej Hausdorffa rozróżniamy
trzy warunki dla miary borelowskiej $\mu$:
\begin{enumerate}[label=(\roman*)]
\item $\mu(K)<\infty$ dla każdego kompaktu $K$;
\item $\mu(E)=\inf\{\mu(U):E\subseteq U,\ U\text{ otwarty}\}$
      dla każdego zbioru borelowskiego $E$;
\item $\mu(U)=\sup\{\mu(K):K\subseteq U,\ K\text{ zwarty}\}$
      dla każdego zbioru otwartego $U$.
\end{enumerate}
To właśnie te trzy własności przyjmujemy w definicji miary Radona
na dowolnej grupie lokalnie zwartej. Nie wymagamy regularności
wewnętrznej na \emph{każdym} zbiorze borelowskim: bez
$\sigma$-zwartości grupy jest to założenie zbyt silne.
Dla borelowskich zbiorów o mierze $\sigma$-skończonej
regularność wewnętrzna wyniknie z (i)--(iii).
Wyjaśnimy różnicę na konkretnym przykładzie na końcu rozdziału.
\end{uwaga}

\section{Lematy topologiczne poprzedzające konstrukcję}
\begin{lemat}[Małe otoczenia, zwarte rozdrobnienia i pokrycia]
\label{lem:haar-neigh-23}
Niech $G$ będzie lokalnie zwartą grupą Hausdorffa.
\begin{enumerate}[label=(\roman*)]
\item Dla każdego otoczenia $O$ elementu $e$ istnieje symetryczne,
  relatywnie zwarte otoczenie otwarte $V$ takie, że
  $V^{-1}V\subseteq O$ oraz $\overline V\subseteq O$.
  Relatywnie zwarte otoczenia otwarte
  tworzą bazę otoczeń $e$.
\item Dla każdego kompaktu $K\subseteq G$ i otwartego otoczenia
  $V$ punktu $e$ istnieje skończone pokrycie $K$
  przez lewe przesunięcia $V$.
\item Jeśli $K,L$ są rozłącznymi kompaktami, to istnieje otoczenie
  $V$ punktu $e$ takie, że żaden zbiór $xV$, $x\in G$,
  nie przecina jednocześnie $K$ i $L$.
\item Jeśli $K\subseteq\bigcup_{j=1}^r O_j$ jest zwarty,
  gdzie $O_j$ są otwarte, to istnieją kompakty
  $K_j\subseteq O_j$ spełniające $K=\bigcup_{j=1}^rK_j$.
\end{enumerate}
\end{lemat}
\begin{proof}
Dla (i) z ciągłości mnożenia wybieramy symetryczne
otoczenie otwarte $W$ takie, że $W^2\subseteq O$.
Z lokalnej zwartości istnieje zwarty $C$ z
$e\in\operatorname{int}C$. Kładziemy
$V=W\cap\operatorname{int}C\cap(\operatorname{int}C)^{-1}$.
Wówczas $V$ jest otwarty i symetryczny,
$\overline V\subseteq C$ oraz
$V^{-1}V\subseteq W^2\subseteq O$.
Ponadto $\overline W\subseteq W^2$: jeśli $x\in\overline W$,
to otwarty zbiór $xW$ przecina $W$, a więc
$x=uw^{-1}$ dla pewnych $u,w\in W$.
Stąd także $\overline V\subseteq O$.
Dobór $O$ wewnątrz dowolnego zadanego otoczenia
daje bazę relatywnie zwartych otoczeń.

W (ii) rodzina $(xV)_{x\in K}$ jest pokryciem otwartym
kompaktu $K$, skąd istnieje podpokrycie skończone.
Dla (iii) kompakt $K^{-1}L$ nie zawiera $e$;
stosujemy (i) do otwartego otoczenia
$G\setminus K^{-1}L$. Gdyby $xv\in K$ i $xw\in L$
przy $v,w\in V$, otrzymalibyśmy
$v^{-1}w\in K^{-1}L\cap V^{-1}V$, sprzeczność.

W (iv) wykorzystamy następujący fakt o każdej przestrzeni
lokalnie zwartej Hausdorffa $X$: jeśli $x\in O$ i $O$
jest otwarty, to istnieje otwarty $W\ni x$ ze
zwartym $\overline W\subseteq O$. Dla dowodu wybieramy
zwartą okolicę $C$ punktu $x$, z $x\in\operatorname{int}C$.
Zbiór $F=C\setminus O$ jest zwartym podzbiorem
przestrzeni Hausdorffa $C$. Dla każdego $y\in F$
rozdzielamy $x$ i $y$ rozłącznymi otoczeniami w $C$.
Zwartość $F$ pozwala wybrać skończenie wiele z nich;
przecięcie otoczeń $x$ daje otwarty $Z\subseteq C$
taki, że $x\in Z$ i $\overline Z^{\,C}\cap F=\varnothing$.
Bierzemy otwarty $\widetilde Z\subseteq X$ z
$\widetilde Z\cap C=Z$ i definiujemy
$W=\widetilde Z\cap\operatorname{int}C$.
Wtedy $\overline W\subseteq C\cap\overline Z^{\,C}
\subseteq O$ oraz $\overline W$ jest zwarty.

Dla każdego $x\in K$ wybieramy zatem indeks $j(x)$
oraz takie otoczenie $W_x$, aby
$x\in W_x\subseteq\overline W_x\subseteq O_{j(x)}$.
Wybierzmy $x_1,\dots,x_s$ tak, aby zbiory $W_{x_i}$
pokrywały $K$, i połóżmy
$K_j=K\cap\bigcup_{i:\,j(x_i)=j}\overline W_{x_i}$.
Wówczas każde $K_j$ jest zwarte, zawiera się w $O_j$,
a ich suma zbiorowa jest równa $K$.
\end{proof}

\section{Zwarta przestrzeń danych bez założenia przeliczalności}
\begin{lemat}[Zwartość dowolnego produktu przedziałów]
\label{lem:haar-product-compact-23}
Niech $I$ będzie dowolnym zbiorem indeksów, a
$a_i\in[0,\infty)$ dla $i\in I$. W topologii produktowej
przestrzeń $P=\prod_{i\in I}[0,a_i]$ jest zwarta.
\end{lemat}
\begin{proof}
Dla zupełności podajemy argument ultrafiltrowy, nie wymagający
uprzedniego użycia twierdzenia Tichonowa. Korzystamy jedynie
z lematu Kuratowskiego--Zorna, czyli ze standardowej postaci
aksjomatu wyboru. Jeśli rodzina zbiorów ma własność skończonych
przecięć, rodzina wszystkich filtrów ją zawierających jest
niepusta. Każdy łańcuch tych filtrów ma ograniczenie górne,
którym jest ich suma zbiorowa. Lemat Zorna dostarcza filtru
maksymalnego $\mathcal U$. Jest to ultrafiltr: gdy $A\notin\mathcal U$, maksymalność oznacza, że
rodzina $\mathcal U\cup\{A\}$ nie może być bazą filtru
właściwego. Istnieje więc $B\in\mathcal U$ z
$B\cap A=\varnothing$, skąd $X\setminus A\in\mathcal U$
($X$ jest tutaj zbiorem, na którym zdefiniowano ultrafiltr).

Aby dowieść zwartości $P$, niech $(F_\alpha)_{\alpha\in J}$
będzie rodziną domkniętych podzbiorów $P$ o własności
skończonych przecięć. Rozszerzamy ją do ultrafiltru
$\mathcal U$ na $P$. Dla każdego $i\in I$ obraz
$\mathcal U$ przez rzutowanie $\pi_i:P\to[0,a_i]$
jest ultrafiltrem na zwartym przedziale. Zwartość przedziału
i własność skończonych przecięć zapewniają, że
$C_i:=\bigcap_{A\in\pi_i(\mathcal U)}\overline A$ jest niepusty.
Wybieramy $t_i=\min C_i$ (nie trzeba dokonywać dodatkowego
wyboru punktów). Każde otwarte otoczenie $W$ punktu $t_i$
należy do $\pi_i(\mathcal U)$: w przeciwnym razie
$[0,a_i]\setminus W$ należałoby do ultrafiltru,
co wykluczałoby $t_i\in C_i$.

Punkt $t=(t_i)_{i\in I}$ należy do domknięcia każdego
zbioru z $\mathcal U$. Istotnie, podstawowe otoczenie
$t$ ogranicza tylko skończenie wiele współrzędnych,
więc należy do $\mathcal U$ i musi przecinać każdy
jego element. W szczególności $t\in F_\alpha$ dla
każdego $\alpha$, ponieważ $F_\alpha$ jest domknięty
i należy do $\mathcal U$. Wobec tego każda rodzina
zbiorów domkniętych o własności skończonych przecięć
ma niepuste przecięcie, co dowodzi zwartości $P$.
\end{proof}

\section{Konstrukcja miary przez liczby pokrycia}
Niech odtąd $G$ będzie \emph{dowolną} lokalnie zwartą
grupą Hausdorffa. Wybieramy kompakt $K_0$ będący otoczeniem $e$
oraz otwarte otoczenie $V_0$ elementu $e$,
spełniające $V_0\subseteq K_0$. Przez $\mathcal N$
oznaczamy rodzinę wszystkich relatywnie zwartych otwartych
otoczeń $e$; jest ona skierowana przez odwrotną inkluzję.
Dla zwartego $K$ i $V\in\mathcal N$ definiujemy
\begin{equation}\label{eq:haar-cover-number-23}
 N(K,V)=\min\left\{r\in\mathbb N_0:\
 K\subseteq\bigcup_{j=1}^{r}x_jV
 \text{ dla pewnych }x_1,\dots,x_r\in G\right\}.
\end{equation}
W szczególności $N(\varnothing,V)=0$, a dla
$K\ne\varnothing$ zachodzi $1\le N(K,V)<\infty$.

\begin{lemat}[Jednolite oszacowania ilorazów]
\label{lem:haar-ratio-23}
Niech $C_K=N(K,V_0)$. Dla każdego $V\in\mathcal N$ mamy
\[
 0\le q_V(K):=\frac{N(K,V)}{N(K_0,V)}\le C_K,
 \qquad q_V(K_0)=1.
\]
Ponadto $q_V(gK)=q_V(K)$, funkcje $q_V$ są monotoniczne
i podaddytywne dla sum zwartych, zaś dla rozłącznych
kompaktów $K,L$ istnieje $W\in\mathcal N$ takie, że
$q_V(K\cup L)=q_V(K)+q_V(L)$ dla każdego $V\subseteq W$.
\end{lemat}
\begin{proof}
Zbiory $K$ można pokryć przez $C_K$ przesunięć $V_0$,
a każde $xV_0\subseteq xK_0$. Pokrywając $K_0$
przez $N(K_0,V)$ przesunięć $V$ i następnie przesuwając
to pokrycie, otrzymujemy $N(K,V)\le C_KN(K_0,V)$.
Niezmienniczość i podaddytywność wynikają z przesuwania
oraz łączenia skończonych pokryć.

Dla rozłącznych $K,L$ wybierzmy $W$ według
lematu~\ref{lem:haar-neigh-23}(iii). Gdy $V\subseteq W$,
żadne przesunięcie $V$ nie przecina zarazem $K$ i $L$.
Dowolne minimalne pokrycie $K\cup L$ można zatem
podzielić na dwie rozłączne rodziny, z których pierwsza
pokrywa $K$, a druga $L$. Stąd
$N(K\cup L,V)\ge N(K,V)+N(L,V)$.
Nierówność odwrotna wynika z łączenia pokryć.
\end{proof}

\begin{lemat}[Zawartość niezmiennicza na kompaktach]
\label{lem:haar-content-23}
Istnieje skończona funkcja nieujemna $q$ na rodzinie
wszystkich kompaktów $K\subseteq G$, taka że
\begin{enumerate}[label=(\roman*)]
\item $q(\varnothing)=0$, $q(K_0)=1$ oraz $q(gK)=q(K)$;
\item $K\subseteq L$ implikuje $q(K)\le q(L)$;
\item $q(K\cup L)\le q(K)+q(L)$;
\item jeśli $K\cap L=\varnothing$, to
      $q(K\cup L)=q(K)+q(L)$.
\end{enumerate}
\end{lemat}
\begin{proof}
Niech $\mathcal K$ będzie zbiorem wszystkich zwartych
podzbiorów $G$. Na mocy
lematu~\ref{lem:haar-product-compact-23} produkt
$P=\prod_{K\in\mathcal K}[0,C_K]$ jest zwarty,
bez względu na moc $\mathcal K$.
Każde $q_V=(q_V(K))_{K\in\mathcal K}$ należy do $P$.
Dla $W\in\mathcal N$ połóżmy
$F_W=\overline{\{q_V:V\in\mathcal N,\ V\subseteq W\}}$.
Jest to niepusty zbiór domknięty. Dla skończonej
rodziny $W_1,\dots,W_r$ mamy
$F_{W_1\cap\dots\cap W_r}\subseteq\bigcap_{j=1}^rF_{W_j}$.
Zwartość $P$ pozwala wybrać
$q\in\bigcap_{W\in\mathcal N}F_W$.

Własności (i)--(iii) opisują domknięte warunki na
skończenie wielu współrzędnych, spełnione przez
każde $q_V$, więc także przez $q$.
Jeżeli $K\cap L=\varnothing$, z poprzedniego lematu
mamy otoczenie $W$, poniżej którego wszystkie $q_V$
spełniają równość addytywności dla $K,L$.
Jej zbiór rozwiązań jest domknięty i zawiera $F_W$,
a więc również $q$. To dowodzi (iv).
\end{proof}

\begin{lemat}[Od zawartości na kompaktach do miary borelowskiej]
\label{lem:haar-content-to-measure-23}
Niech $X$ będzie lokalnie zwartą przestrzenią Hausdorffa,
a $q$ skończoną funkcją nieujemną na kompaktach,
spełniającą (i)--(iv) z poprzedniego lematu z pominięciem
niezmienniczości i normalizacji. Dla otwartego $U$
i dowolnego $A\subseteq X$ definiujemy
\begin{equation}\label{eq:haar-from-content-23}
 t(U)=\sup_{K\subseteq U,\ K\ \mathrm{zwarty}}q(K),
 \qquad
 m^*(A)=\inf_{U\supseteq A,\ U\ \mathrm{otwarty}}t(U).
\end{equation}
Wówczas $m^*$ jest miarą zewnętrzną,
wszystkie zbiory borelowskie są mierzalne w sensie
Carathéodory'ego, a $\mu=m^*|_{\mathcal B(X)}$
spełnia warunki (i)--(iii) regularności podane
na początku rozdziału. Ponadto $\mu(K)\ge q(K)$.
Jeśli homeomorfizm $T:X\to X$ zachowuje $q$ na
kompaktach, to $\mu(TA)=\mu(A)$ dla każdego
zbioru borelowskiego $A$.
\end{lemat}
\begin{proof}
Wprost z definicji $m^*(U)=t(U)$, gdy $U$ jest otwarty.
Najpierw dowodzimy podaddytywności $t$ na pokryciach
otwartych. Jeśli $K\subseteq\bigcup_{j\ge1}U_j$ jest
zwarty, możemy wybrać skończone podpokrycie.
Z lematu~\ref{lem:haar-neigh-23}(iv) (którego dowód
jest topologiczny i działa w każdej przestrzeni lokalnie
zwartej Hausdorffa) dostajemy kompakty
$K_j\subseteq U_j$ pokrywające $K$. Podaddytywność $q$
daje $q(K)\le\sum_jt(U_j)$. Biorąc supremum po $K$,
uzyskujemy
$t(\bigcup_{j\ge1}U_j)\le\sum_jt(U_j)$.
Aby uzyskać przeliczalną podaddytywność $m^*$,
dla zbiorów $A_j$ o skończonej sumie $\sum_jm^*(A_j)$
i $\varepsilon>0$ wybieramy otwarte $U_j\supseteq A_j$
z $t(U_j)<m^*(A_j)+\varepsilon2^{-j}$.
Wtedy $m^*(\bigcup_jA_j)\le\sum_jm^*(A_j)+\varepsilon$;
pozostałe własności miary zewnętrznej są natychmiastowe.

Najważniejszym krokiem jest mierzalność \emph{otwartych}
zbiorów bez odwoływania się do jakiejkolwiek metryki.
Ustalmy otwarty $O\subseteq X$, dowolny zbiór $A$
o $m^*(A)<\infty$ i $\varepsilon>0$.
Wybieramy otwarty $V\supseteq A$ z
$t(V)<m^*(A)+\varepsilon$. Ponieważ
$t(V\cap O)\le t(V)<\infty$, istnieje zwarty
$K\subseteq V\cap O$ taki, że
$q(K)>t(V\cap O)-\varepsilon$; jeśli supremum
jest zerowe, bierzemy $K=\varnothing$.
Dla dowolnego kompaktu $L\subseteq V\setminus K$
addytywność na rozłącznych kompaktach daje
$q(K)+q(L)=q(K\cup L)\le t(V)$.
Stąd $t(V\setminus K)\le t(V)-q(K)$.
Ponieważ $A\setminus O\subseteq V\setminus K$,
a $A\cap O\subseteq V\cap O$, dostajemy
\[
 m^*(A\setminus O)+m^*(A\cap O)
 \le t(V\setminus K)+t(V\cap O)
 \le t(V)+\varepsilon<m^*(A)+2\varepsilon.
\]
Przy $m^*(A)=\infty$ potrzebna nierówność
jest automatyczna. Po przejściu z $\varepsilon$ do
zera warunek Carathéodory'ego jest spełniony.
Z twierdzenia~\ref{thm:caratheodory} wynika,
że $\mu$ jest miarą na $\mathcal B(X)$.

Regularność zewnętrzna jest zawarta w definicji $m^*$.
Dla kompaktu $K$ wybieramy otwarty $W\supseteq K$
o zwartej domkniętej powłoce $C$. Wówczas
$\mu(K)\le t(W)\le q(C)<\infty$.
Dla każdego otwartego $U\supseteq K$ mamy
$t(U)\ge q(K)$, więc $\mu(K)\ge q(K)$.
Dla otwartego $O$ zachodzi
$\mu(O)=t(O)=\sup_{K\subseteq O}q(K)
\le\sup_{K\subseteq O}\mu(K)\le\mu(O)$,
co dowodzi regularności wewnętrznej na otwartych.

Na koniec, jeśli $T$ jest wskazanym homeomorfizmem,
to $T$ przekształca bijektywnie kompakty zawarte w $U$
na kompakty zawarte w $TU$, skąd
$t(TU)=t(U)$ i $m^*(TA)=m^*(A)$.
\end{proof}

\begin{proof}[Dowód istnienia w twierdzeniu~\ref{thm:haar-main-23}
bez założeń przeliczalności]
Z lematu~\ref{lem:haar-content-23} otrzymujemy
zawartość $q$. Lemat~\ref{lem:haar-content-to-measure-23}
daje borelowską miarę Radona $\mu$ niezmienniczą
względem wszystkich lewych translacji.
Ponieważ $\mu(K_0)\ge q(K_0)=1$, jest niezerowa.
Jej skończoność na kompaktach wynika z konstrukcji.
Dla miary prawej stosujemy inwersję:
$\mu_R(E)=\mu(E^{-1})$.
Jest to miara Radona, bo inwersja jest homeomorfizmem,
a $\mu_R(Eg)=\mu(g^{-1}E^{-1})=\mu(E^{-1})$.
Nie było potrzebne żadne założenie o przeliczalnej bazie
otoczeń $e$ ani o $\sigma$-zwartości $G$.
\end{proof}

\begin{wniosek}[Dodatniość miary zbiorów otwartych]
\label{cor:haar-open-positive-23}
Jeśli $\mu$ jest niezerową lewą miarą Haara
na lokalnie zwartej grupie Hausdorffa $G$, to
$\mu(U)>0$ dla każdego niepustego otwartego $U\subseteq G$.
\end{wniosek}
\begin{proof}
Jeśli $\mu(U)=0$ dla pewnego niepustego $U$, wybierzmy
$a\in U$. Dla każdego kompaktu $K\subseteq G$ skończenie
wiele zbiorów $xa^{-1}U$ pokrywa $K$, ponieważ każdy
z nich jest przesunięciem otwartego otoczenia $e$.
Z niezmienniczości wynika $\mu(K)=0$.
Regularność wewnętrzna na otwartych daje teraz
$\mu(G)=0$, wbrew niezerowości miary.
\end{proof}

\section{Jednoznaczność: argument z funkcją przybliżającą jedynkę}
\begin{lemat}[Jednostajna ciągłość translacji funkcji]
\label{lem:haar-uniform-23}
Jeśli $f\in C_c(G)$, to dla każdego $\varepsilon>0$
istnieje otoczenie $V$ elementu $e$ takie, że
$|f(yz)-f(y)|<\varepsilon$ dla wszystkich
$y\in G$ i $z\in V$.
\end{lemat}
\begin{proof}
Niech $K=\operatorname{supp}f$ i wybierzmy symetryczne
relatywnie zwarte otoczenie $C$ elementu $e$.
Zbiór $M=K\overline C$ jest zwarty. Ciągłość
$(y,z)\mapsto f(yz)-f(y)$ w każdym $(y,e)$,
$y\in M$, pozwala wybrać otoczenia $W_y$ punktu $y$
i $V_y$ elementu $e$ takie, że
$|f(uz)-f(u)|<\varepsilon$ dla $u\in W_y$, $z\in V_y$.
Ze zwartego $M$ wybieramy skończone pokrycie
$W_{y_1},\dots,W_{y_r}$; połóżmy
$V=C\cap V_{y_1}\cap\cdots\cap V_{y_r}$,
zmniejszając je w razie potrzeby do symetrycznego
otoczenia. Dla $u\in M$, $z\in V$ teza wynika
z konstrukcji. Jeśli $u\notin M$, to $f(u)=0$,
a $uz\in K$ prowadziłoby do
$u\in K V^{-1}\subseteq K\overline C=M$.
Zatem również $f(uz)=0$.
\end{proof}

\begin{lemat}[Fubini dla ciągłych jąder o zwartym nośniku]
\label{lem:haar-fubini-compact-23}
Niech $X,Y$ będą lokalnie zwartymi przestrzeniami Hausdorffa,
a $\alpha,\beta$ miarami borelowskimi skończonymi na
kompaktach. Jeśli $H\in C_c(X\times Y)$, to obie całki iterowane
\[
 \int_X\!\left(\int_Y H(x,y)\,\mathrm d\beta(y)\right)
 \mathrm d\alpha(x),\qquad
 \int_Y\!\left(\int_X H(x,y)\,\mathrm d\alpha(x)\right)
 \mathrm d\beta(y)
\]
są dobrze określone, skończone i równe, nawet gdy
$X,Y$ nie są metryzowalne ani $\alpha,\beta$
nie są $\sigma$-skończone.
\end{lemat}
\begin{proof}
Rzuty zwartego nośnika $H$ są kompaktami $K\subseteq X$
i $L\subseteq Y$; całkowanie ograniczamy do $K\times L$.
Na iloczynie tych zwartych przestrzeni każdą funkcję
ciągłą przybliża się jednostajnie przez sumy
$\sum_{j=1}^Nc_j a_j(x)b_j(y)$, gdzie
$a_j\in C(K)$, $b_j\in C(L)$. Oto bezpośrednie
uzasadnienie tej aproksymacji. Dla danego $\varepsilon>0$
z ciągłości $H$ i zwartości iloczynu można wybrać
skończone otwarte pokrycia $(U_i)$ zbioru $K$ oraz
$(V_j)$ zbioru $L$ takie, że oscylacja $H$ na każdym
$U_i\times V_j$ jest mniejsza niż $\varepsilon$.
Istnienie tych pokryć sprawdzamy bez metryzowalności:
dla ustalonego $x\in K$ i każdego $y\in L$ ciągłość
$H$ daje otoczenia $U_{x,y}\ni x$, $V_{x,y}\ni y$
takie, że $|H(x',y')-H(x,y)|<\varepsilon/4$
dla $(x',y')\in U_{x,y}\times V_{x,y}$.
Zwartość $L$ pozwala wybrać skończenie wiele
$V_{x,y}$ pokrywających $L$; przecięcie
odpowiednich $U_{x,y}$ daje otoczenie $U_x$,
na którym dla dowolnych $x',x''\in U_x$ i $y\in L$
mamy $|H(x',y)-H(x'',y)|<\varepsilon/2$.
Analogicznie dobieramy otoczenia $V_y$ z oszacowaniem
$|H(x,y')-H(x,y'')|<\varepsilon/2$ jednostajnym
względem $x\in K$. Skończone podpokrycia
$(U_i)$, $(V_j)$ dają żądane oszacowanie oscylacji
na prostokątach $U_i\times V_j$.
Na zwartych przestrzeniach
Hausdorffa istnieją ciągłe podziały jedności
$(a_i)$, $(b_j)$ podporządkowane tym pokryciom.
Wybierając $c_{ij}$ jako wartość $H$ w dowolnym
punkcie niepustego $U_i\times V_j$, otrzymujemy
\[
 \left\|H-\sum_{i,j}c_{ij}a_i\otimes b_j\right\|_\infty
 \le\varepsilon.
\]
Dla sum elementarnych obie całki iterowane są równe
z liniowości. Skończoność $\alpha(K)$ i $\beta(L)$
pozwala przejść do granicy jednostajnej, także dla
całek funkcji zespolonych. Dowodzi to tezy.
\end{proof}

\begin{twierdzenie}[Jednoznaczność miary Haara]
\label{thm:haar-unique-23}
Jeśli $\mu,\nu$ są dwiema lewostronnymi miarami Haara
na lokalnie zwartej grupie Hausdorffa $G$, to istnieje
$c>0$ takie, że $\nu=c\mu$ na $\mathcal B(G)$.
\end{twierdzenie}
\begin{proof}
Niech $f\in C_c(G)$, a $\varphi\in C_c(G)$ będzie
nieujemna, niezerowa i spełniała
$\int_G\varphi(z)\,\mathrm d\mu(z)=1$.
Funkcja $(x,y)\mapsto f(x)\varphi(y^{-1}x)$
ma nośnik zawarty w zwartym iloczynie
$\operatorname{supp}f\times
(\operatorname{supp}f)(\operatorname{supp}\varphi)^{-1}$.
Stosujemy zatem lemat~\ref{lem:haar-fubini-compact-23},
który nie wymaga $\sigma$-skończoności miary na $G$.
Otrzymujemy
\begin{align*}
 I&=\int_G\!\int_G f(x)\varphi(y^{-1}x)
       \,\mathrm d\nu(y)\,\mathrm d\mu(x)\\
  &=c_\varphi\int_Gf(x)\,\mathrm d\mu(x),\qquad
 c_\varphi:=\int_G\varphi(z^{-1})\,\mathrm d\nu(z),
\end{align*}
bo po podstawieniu $y=xz$ wewnętrzna całka nie
zależy od $x$ (lewostronna niezmienniczość $\nu$).
Zamieniając kolejność całkowania i podstawiając
$x=yz$ (lewostronna niezmienniczość $\mu$), dostajemy
\begin{equation}\label{eq:haar-uniqueness-23}
 c_\varphi\int_G f(x)\,\mathrm d\mu(x)
 =\int_G\!\int_G f(yz)\varphi(z)
       \,\mathrm d\mu(z)\,\mathrm d\nu(y).
\end{equation}
Wybierzmy funkcje $\varphi$ o nośnikach zawartych
w coraz mniejszych otoczeniach $e$, każdorazowo
unormowane przez $\mu(\varphi)=1$. Są dostępne,
ponieważ każda niepusta otwarta część grupy ma
dodatnią miarę Haara: w przeciwnym razie jej
translacje pokryłyby skończenie zwarty zbiór
o dodatniej mierze.
Na mocy lematu~\ref{lem:haar-uniform-23}
wewnętrzna całka po prawej stronie
\eqref{eq:haar-uniqueness-23} zbiega jednostajnie
na odpowiednim zwartym zbiorze do $f(y)$,
a poza nim obie funkcje znikają.
Zatem prawa strona dąży do $\int_Gf(y)\,\mathrm d\nu(y)$.
Wybierając ustaloną funkcję $f_0\in C_c(G)$,
$f_0\ge0$ i $\mu(f_0)>0$, otrzymujemy
$c_\varphi\to c:=\nu(f_0)/\mu(f_0)>0$.
Wstawienie dowolnego $f$ do
\eqref{eq:haar-uniqueness-23} daje
$\nu(f)=c\mu(f)$. Pozostaje wyjaśnić, dlaczego równość całek względem
$\mu$ i $\nu$ dla wszystkich funkcji z $C_c(G)$ oznacza
równość miar również bez założenia metryzowalności $G$.
Dla każdego zwartego $K\subseteq U$, gdzie $U$ jest otwarty,
istnieje funkcja $h\in C_c(G)$ taka, że
$0\le h\le1$, $h|_K=1$ i $\operatorname{supp}h\subseteq U$
(lokalnie zwarta przestrzeń Hausdorffa ma takie funkcje
odcinające). Wówczas $\mu(K)\le\int_Gh(x)\,\mathrm d\mu(x)
=\int_Gh(x)\,\mathrm d\nu(x)\le\nu(U)$. Z regularności
wewnętrznej $\mu(U)\le\nu(U)$. Po zamianie miar
otrzymujemy równość na otwartych, a regularność zewnętrzna
daje równość na wszystkich zbiorach borelowskich.
Stąd $\nu=c\mu$, co kończy dowód.
\end{proof}

\begin{przyklad}[Dlaczego nie wymagamy regularności wewnętrznej wszędzie]
\label{ex:haar-non-sigma-compact-23}
Niech $D$ będzie nieprzeliczalną grupą dyskretną,
a $G=\mathbb R\times D$ grupą z topologią produktową.
Rozważmy domknięty zbiór $Y=\{0\}\times D$.
Każdy kompakt $K\subseteq Y$ jest skończony, ponieważ
rzutowanie $K$ na dyskretną przestrzeń $D$ daje zbiór
zwarty, a więc skończony. Punkty mają miarę Haara zero:
na otwartej podgrupie $\mathbb R\times\{e\}$ miara Haara
jest dodatnią wielokrotnością miary Lebesgue'a, zaś
translacje przenoszą ten wniosek na pozostałe składowe.
Wobec tego $\mu(K)=0$ dla każdego kompaktu $K\subseteq Y$.

Z drugiej strony, jeśli $U\supseteq Y$ jest otwarty,
to $U\cap(\mathbb R\times\{d\})$ jest niepustym
zbiorem otwartym dla każdego $d\in D$, a więc ma
ściśle dodatnią miarę Haara na tej składowej.
Składowe są parami rozłączne i jest ich nieprzeliczalnie
wiele, zatem $\mu(U)=\infty$: gdyby $\mu(U)$ było
skończone, najwyżej przeliczalnie wiele rozłącznych
zbiorów mierzalnych mogłoby mieć dodatnią miarę.
Regularność zewnętrzna daje $\mu(Y)=\infty$, chociaż
$\sup_{K\subseteq Y,\,K\text{ zwarty}}\mu(K)=0$.
Takiego zjawiska nie ma dla zbiorów o mierze
$\sigma$-skończonej, zgodnie z lematem
\ref{lem:radon-sigmafinite-inner-20}.
\end{przyklad}

\section{Prawe miary Haara, funkcja modularna i unimodularność}
\begin{twierdzenie}[Prawa miara Haara i modularność]
\label{thm:haar-modular-23}
Niech $\mu$ będzie lewostronną miarą Haara.
Wówczas $\mu_R(E):=\mu(E^{-1})$ jest prawostronną
miarą Haara. Istnieje dokładnie jeden ciągły
homomorfizm $c:G\to(0,\infty)$ taki, że
\begin{equation}\label{eq:haar-modular-23}
 \mu(Eg)=c(g)\mu(E)\qquad(E\in\mathcal B(G),\ g\in G).
\end{equation}
Grupę nazywamy \emph{unimodularną}, gdy $c\equiv1$.
\end{twierdzenie}
\begin{proof}
Inwersja jest homeomorfizmem, toteż $\mu_R$ jest
miarą Radona. Ponadto
$\mu_R(Eg)=\mu((Eg)^{-1})=\mu(g^{-1}E^{-1})
=\mu(E^{-1})=\mu_R(E)$.
Dla ustalonego $g$ miara $E\mapsto\mu(Eg)$ jest
lewostronnie niezmiennicza, a więc na mocy
twierdzenia~\ref{thm:haar-unique-23} równa $c(g)\mu$.
Złożenie prawych przesunięć daje
$c(gh)=c(g)c(h)$, zaś $c(e)=1$.
Dla niezerowej $f\in C_c(G)$, $f\ge0$,
$c(g)$ jest ilorazem całek
\[
 c(g)=\frac{\int_G f(xg^{-1})\,\mathrm d\mu(x)}
 {\int_G f(x)\,\mathrm d\mu(x)}.
\]
Ciągłość licznika wynika z ciągłości translacji
i lokalnego wspólnego zwartego nośnika funkcji,
zatem $c$ jest ciągła.
\end{proof}

\begin{uwaga}[Dwie konwencje]
W literaturze przez funkcję modularną często oznacza
się $\Delta=c^{-1}$, a nie $c$. Dlatego zawsze
należy sprawdzić wzór definiujący. W tym skrypcie
pozostajemy przy konwencji
$\mu(Eg)=c(g)\mu(E)$.
\end{uwaga}

\begin{wniosek}[Modularność grupy Liego]
\label{cor:haar-adjoint-23}
Niech $G$ będzie skończenie wymiarową grupą Liego,
$\mathfrak g=T_eG$ jej algebrą Liego i
$\operatorname{Ad}_g:\mathfrak g\to\mathfrak g$
różniczką w $e$ sprzężenia $x\mapsto gxg^{-1}$.
Dla konwencji modularnej z~\eqref{eq:haar-modular-23}
zachodzi
\[
 c(g)=|\det(\operatorname{Ad}_{g^{-1}})|
       =|\det(\operatorname{Ad}_{g})|^{-1}.
\]
W szczególności grupa Liego jest unimodularna wtedy
i tylko wtedy, gdy $|\det\operatorname{Ad}_g|=1$
dla każdego $g\in G$.
\end{wniosek}
\begin{proof}
Niech $R_g(x)=xg$, $L_x(y)=xy$.
Przez lewą translację identyfikujemy $T_xG$ i
$T_{xg}G$ z $\mathfrak g$. W tych identyfikacjach
różniczka $R_g$ ma postać
\[
 \mathrm d(L_{(xg)^{-1}})_{xg}\circ
 \mathrm d(R_g)_x\circ\mathrm d(L_x)_e
 =\operatorname{Ad}_{g^{-1}},
\]
co sprawdzamy, podstawiając krzywą $h(t)$ wychodzącą
z $e$: $(xg)^{-1}(xh(t)g)=g^{-1}h(t)g$.
Lewa gęstość objętości na grupie Liego ma w tej
trywializacji stały współczynnik. Prawa translacja
mnoży więc jej lokalną objętość przez
$|\det\operatorname{Ad}_{g^{-1}}|$.
Podstawienie w całce daje
$\mu(Eg)=|\det\operatorname{Ad}_{g^{-1}}|\mu(E)$,
a jednoznaczność czynnika modularnego kończy dowód.
\end{proof}

\begin{przyklad}[Grupy zwarte i dyskretne]
\label{ex:haar-compact-discrete-23}
Na grupie zwartej $G$ normalizujemy miarę Haara
przez $\mu(G)=1$. Dla każdego $g$ równość
$\mu(Gg)=c(g)\mu(G)$ daje $c(g)=1$.
Na grupie dyskretnej miarą Haara jest miara licząca,
więc także $c\equiv1$. Jeżeli grupa jest skończona,
można ją unormować do prawdopodobieństwa przez
podzielenie miary liczącej przez $|G|$.
\end{przyklad}

\begin{przyklad}[$\mathbb R^d$, torus i grupy macierzowe]
\label{ex:haar-standard-23}
Na $(\mathbb R^d,+)$ lewostronną i prawostronną miarą
Haara jest miara Lebesgue'a $\lambda_d$.
Na torusie $\mathbb T^d=\mathbb R^d/\mathbb Z^d$
miarę Haara dostajemy przez przeniesienie
$\lambda_d|_{[0,1)^d}$ na iloraz; jest to miara
probabilistyczna. Na grupie
$\operatorname{GL}_d(\mathbb R)$ lewa i prawa miara Haara
mają gęstość $|\det A|^{-d}$ względem miary
Lebesgue'a na przestrzeni $M_d(\mathbb R)\cong\mathbb R^{d^2}$.
Rzeczywiście, mnożenie przez $B$ z lewej zmienia
jakobian przez $|\det B|^d$, a wyznacznik $A$
przez czynnik $\det B$; oba czynniki się skracają.
Analogicznie dla mnożenia z prawej.
\end{przyklad}

\begin{przyklad}[Nieunimodularna grupa afiniczna]
\label{ex:haar-affine-23}
Rozważmy grupę $G=\mathbb R\rtimes(0,\infty)$
z mnożeniem
$(b,a)(c,t)=(b+ac,at)$. Lewą miarą Haara jest
\[
 \mathrm d\mu(b,a)=\frac{\mathrm d\lambda_2(b,a)}{a^2},
\]
a prawą miarą Haara
$\mathrm d\mu_R(b,a)=a^{-1}\mathrm d\lambda_2(b,a)$.
Przy lewym przesunięciu przez $(c,t)$ jakobian
wynosi $t^2$, a czynnik $a^{-2}$ zmienia się
o $t^{-2}$; otrzymujemy niezmienniczość.
Przy prawym przesunięciu przez $(c,t)$ jakobian
wynosi $t$ i bezpośredni rachunek daje
$\mu(E(c,t))=t^{-1}\mu(E)$.
Jest to prosty przykład $c\not\equiv1$.
\end{przyklad}

\begin{uwaga}[Grupy Liego]
Jeżeli $G$ jest $n$-wymiarową grupą Liego, wybór
niezerowej $n$-formy w $T_eG$ i przeniesienie jej
translacjami lewymi daje lewostronną, nigdzie
nieznikającą formę najwyższego stopnia. Dowodzi to
zarazem, że \emph{każda grupa Liego jest orientowalna},
niezależnie od liczby składowych spójności. Wybierając
orientację zgodną z tą formą i całkując odpowiadającą
jej dodatnią gęstość, otrzymujemy lewą miarę Haara.
W ten sposób konstrukcja analityczna łączy się
z geometrią różniczkową.
\end{uwaga}

% END chapters/pl/measure_theory/23-rozdzial.tex


% BEGIN chapters/pl/measure_theory/24-rozdzial.tex
% ROZDZIAŁ 24 — TWIERDZENIA POKRYCIOWE
\chapter{Twierdzenia pokryciowe Vitaliego i Besicovitcha}
\label{ch:covering-24}

\section{Motywacja i terminologia}
W poprzednich rozdziałach zbiory mierzalne przybliżaliśmy
zbiorami otwartymi i zwartymi. W zastosowaniach do
różniczkowania całek i miar potrzebujemy mocniejszego narzędzia:
z dowolnej rodziny małych kul należy wybrać takie kule, których
nakładanie można kontrolować. W tym rozdziale podajemy trzy
odmienne rezultaty: lemat $5r$, twierdzenie Vitaliego i twierdzenie
Besicovitcha. Nie wolno ich utożsamiać: pierwsze dotyczy
\emph{powiększonych} kul, drugie pokrycia \emph{z dokładnością do
zbioru miary zero}, a trzecie pokrycia kulami \emph{bez powiększania},
z ograniczoną krotnością.

\begin{definicja}[Kule i krotność pokrycia]
\label{def:balls-24}
Niech $(X,d)$ będzie przestrzenią metryczną. Piszemy
$B(x,r)=\{y\in X:d(x,y)<r\}$ dla $r>0$.
Jeśli $B=B(x,r)$ i $a>0$, to $aB:=B(x,ar)$.
Rodzina $(B_i)_{i\in I}$ jest \emph{parami rozłączna}, jeżeli
$B_i\cap B_j=\varnothing$ dla $i\ne j$.
Jej \emph{krotność w punkcie} $x$ to liczba indeksów $i$ takich,
że $x\in B_i$ (dopuszczamy $+\infty$).
Funkcja charakterystyczna dowolnego zbioru $A$ jest w całym
skrypcie oznaczana przez $\chi_A$.
\end{definicja}

\begin{uwaga}
Twierdzenia sformułowane dla kul w $\mathbb R^d$ odnoszą się
zawsze do metryki euklidesowej i miary Lebesgue'a $\lambda_d$.
Przypadki metryczne bez miary rozpatrujemy oddzielnie.
Symbol $v_d:=\lambda_d(B(0,1))$ oznacza objętość kuli jednostkowej;
$\lambda_d(B(x,r))=v_dr^d$.
\end{uwaga}

\section{Lemat $5r$: wybór kul bez zakładania miary}
\begin{twierdzenie}[Lemat o pięciokrotnym powiększeniu]
\label{thm:5r-24}
Niech $(X,d)$ będzie przestrzenią metryczną, a $\mathcal B$
niepustą rodziną otwartych kul o promieniach ograniczonych
z góry przez pewną liczbę $R<\infty$.
Istnieje parami rozłączna podrodzina $\mathcal D\subseteq\mathcal B$
taka, że
\[
 \bigcup_{B\in\mathcal B}B\subseteq\bigcup_{D\in\mathcal D}5D.
\]
Jeżeli $X$ jest ośrodkowa, rodzina $\mathcal D$ jest co najwyżej
przeliczalna.
\end{twierdzenie}
\begin{proof}
Rozdzielmy kule na klasy według promienia:
$\mathcal B_k=\{B(x,r)\in\mathcal B:
2^{-k-1}R<r\le2^{-k}R\}$, $k=0,1,\ldots$.
W klasie $\mathcal B_0$ wybieramy maksymalną względem inkluzji
rodzinę parami rozłącznych kul. Indukcyjnie, po dokonaniu wyborów
z klas wcześniejszych, w $\mathcal B_k$ wybieramy maksymalną
rodzinę parami rozłącznych kul, rozłącznych również z kulami
wybranymi wcześniej. Maksymalność można uzyskać z lematu
Zorna: uporządkowujemy przez inkluzję rodziny spełniające
warunek rozłączności, a suma łańcucha pozostaje taką rodziną.
Niech $\mathcal D$ będzie sumą wszystkich wybranych rodzin.

Weźmy kulę $B=B(x,r)\in\mathcal B_k$, która nie została
wybrana. Na odpowiednim etapie była albo niedopuszczalna
wskutek przecięcia z wcześniej wybraną kulą, albo przecinała
kulę z maksymalnego wyboru w swojej klasie. W obu przypadkach
istnieje $D=B(y,s)\in\mathcal D$ taka, że $B\cap D\ne\varnothing$
i $s>2^{-k-1}R\ge r/2$. Dla każdego $z\in B$ mamy
\[
 d(z,y)\le d(z,x)+d(x,y)<r+(r+s)=2r+s<5s.
\]
Stąd $B\subseteq5D$. Teza jest oczywista także dla kul
wybranych, bo $D\subseteq5D$.

Jeżeli $X$ ma przeliczalny zbiór gęsty $Q$, to każda wybrana
kula zawiera punkt $Q$. Kule z $\mathcal D$ są rozłączne,
więc przyporządkowanie im różnych punktów $Q$ dowodzi
przeliczalności $\mathcal D$.
\end{proof}

\begin{wniosek}[Oszacowanie miary sumy kul]
\label{cor:5r-volume-24}
Niech $\mathcal B$ będzie rodziną kul euklidesowych o promieniach
ograniczonych z góry. Istnieje parami rozłączna podrodzina
$(B_i)$ taka, że
\[
 \lambda_d\Bigl(\bigcup_{B\in\mathcal B}B\Bigr)
 \le 5^d\sum_i\lambda_d(B_i).
\]
\end{wniosek}
\begin{proof}
Stosujemy twierdzenie~\ref{thm:5r-24}, przeliczalną
podaddytywność oraz równość $\lambda_d(5B_i)=5^d\lambda_d(B_i)$.
Ponieważ suma rozłącznych kul jest mierzalna, szereg po prawej
jest dobrze określony również wtedy, gdy ma wartość $+\infty$.
\end{proof}

\begin{uwaga}[Inne stałe]
W przypadku kul w $\mathbb R^d$ często formułuje się wariant
$3r$, jeżeli promienie porządkuje się dokładniej. Stała $5$
jest wygodna, gdy nie zakłada się osiągania supremum promieni;
jej konkretna wartość nie jest istotna w większości zastosowań.
Nie należy jednak z lematu $5r$ wyciągać wniosku, że same
wybrane kule pokrywają całą początkową rodzinę.
\end{uwaga}

\begin{center}
\begin{tikzpicture}[scale=0.82]
  \draw[gray!80,thick] (0,0) circle (0.65);
  \draw[gray!80,thick] (1.18,0.18) circle (0.66);
  \draw[gray!80,thick] (3.1,0) circle (0.55);
  \fill (0,0) circle (1.6pt);
  \fill (1.18,0.18) circle (1.6pt);
  \fill (3.1,0) circle (1.6pt);
  \draw[blue!70!black,thick] (0,0) circle (0.65);
  \draw[blue!70!black,thick] (3.1,0) circle (0.55);
  \draw[blue!50!black,densely dashed] (0,0) circle (1.75);
  \draw[blue!50!black,densely dashed] (3.1,0) circle (1.45);
  \node[font=\small] at (1.5,-2.2) {Wybrane kule i ich powiększenia (schemat)};
  \node[font=\scriptsize,align=center] at (1.2,1.05) {kula\\odrzucona};
\end{tikzpicture}
\end{center}
\noindent\textit{Idea lematu $5r$:} kula odrzucona przecina jedną
z wybranych kul o porównywalnym promieniu, dlatego jest
zawarta w jej pięciokrotnym powiększeniu. Okręgi
kreskowane mają na rysunku promienie poglądowe, nie
w skali $5{:}1$.

\section{Pokrycia drobne i twierdzenie Vitaliego}
\begin{definicja}[Drobne pokrycie Vitaliego]
\label{def:vitali-fine-24}
Niech $E\subseteq\mathbb R^d$ i $\mathcal V$ będzie rodziną
otwartych kul euklidesowych. Mówimy, że $\mathcal V$ jest
\emph{drobne w każdym punkcie $E$}, jeżeli dla wszystkich
$x\in E$ i $\varepsilon>0$ istnieje kula
$B(x,r)\in\mathcal V$ ze środkiem w $x$ i $0<r<\varepsilon$.
Wymaganie środka w punkcie $x$ jest tutaj częścią definicji.
\end{definicja}

\begin{twierdzenie}[Pokryciowe twierdzenie Vitaliego]
\label{thm:vitali-cover-24}
Niech $E\subseteq\mathbb R^d$ będzie ograniczony,
$O\subseteq\mathbb R^d$ otwarty i ograniczony, $E\subseteq O$,
a $\mathcal V$ rodziną kul drobną w każdym punkcie $E$.
Wówczas można wybrać skończoną lub przeliczalną rodzinę
parami rozłącznych kul $(B_j)$ z $\mathcal V$, spełniającą
$B_j\subseteq O$ oraz
\[
 \lambda_d^*\Bigl(E\setminus\bigcup_jB_j\Bigr)=0,
 \qquad \sum_j\lambda_d(B_j)\le\lambda_d(O).
\]
Zapis $\lambda_d^*$ oznacza zewnętrzną miarę Lebesgue'a;
nie zakładamy z góry mierzalności zbioru $E$.
\end{twierdzenie}
\begin{proof}
Ograniczamy rodzinę do tych kul, których domknięcia mieszczą
się w $O$. Nadal jest ona drobna w każdym punkcie $E$:
przy ustalonym $x\in E$ wybieramy promień mniejszy od
odległości $x$ od $\mathbb R^d\setminus O$.
Wszystkie takie kule mieszczą się w $O$, więc ich promienie
są jednostajnie ograniczone.

Budujemy indukcyjnie rozłączne kule. Jeżeli po wybraniu
$B_1,\ldots,B_{n-1}$ pozostaje choć jedna kula naszej rodziny
rozłączna ze wszystkimi wcześniej wybranymi, oznaczmy przez
$s_n$ supremum promieni takich kul. Wybierzmy jedną z nich,
$B_n=B(x_n,r_n)$, o promieniu $r_n>s_n/2$. Jeżeli nie ma
dopuszczalnych kul, kończymy. Konstrukcja daje
$\sum_jv_dr_j^d\le\lambda_d(O)<\infty$, a zatem
$r_j\to0$ w przypadku nieskończonego ciągu.

Niech $N$ będzie dowolną dodatnią liczbą całkowitą.
Pokażemy, że poza zbiorem
$Z=\bigcup_j\partial B_j$ zachodzi inkluzja
\begin{equation}\label{eq:vitali-tail-24}
 E\setminus\bigcup_j\overline{B_j}
 \subseteq\bigcup_{j>N}5B_j.
\end{equation}
Rzeczywiście, ustalmy $x\in E\setminus\bigcup_j\overline{B_j}$.
Jest on w dodatniej odległości od skończonej sumy
$\bigcup_{j\le N}\overline{B_j}$. Z drobności pokrycia
wybieramy kulę $B=B(x,r)\in\mathcal V$ o domknięciu
w $O$, rozłączną z $B_1,\ldots,B_N$.
Jeżeli konstrukcja zakończyła się przed krokiem $N+1$,
taka kula przeczyłaby warunkowi zakończenia. Jeżeli zaś
konstrukcja trwa nieskończenie długo, $B$ musi przecinać
którąś z późniejszych kul. Gdyby bowiem była rozłączna
ze wszystkimi $B_j$, pozostawałaby dopuszczalna na
każdym etapie. Otrzymalibyśmy $s_j\ge r$ i
$r_j>s_j/2\ge r/2$ dla wszystkich $j>N$,
sprzecznie z $\sum_jr_j^d<\infty$.
Niech zatem $j>N$ będzie pierwszym indeksem, dla
którego $B\cap B_j\ne\varnothing$.
Przed krokiem $j$ kula $B$ pozostawała dopuszczalna,
zatem $r\le s_j<2r_j$. Odległość centrów jest mniejsza
od $r+r_j$, a więc $d(x,x_j)<3r_j<5r_j$.
Otrzymaliśmy $x\in5B_j$, co daje \eqref{eq:vitali-tail-24}.

Każda sfera $\partial B_j$ ma zerową miarę Lebesgue'a
(co wynika na przykład z zawierania jej w powłoce kulistej
o dowolnie małej objętości). Zatem $\lambda_d(Z)=0$.
Ze \eqref{eq:vitali-tail-24}, podaddytywności i skalowania
miary wynika
\[
 \lambda_d^*\Bigl(E\setminus\bigcup_jB_j\Bigr)
 \le\lambda_d(Z)+5^d\sum_{j>N}\lambda_d(B_j).
\]
Po przejściu $N\to\infty$ prawa strona dąży do zera.
Suma objętości nie przekracza $\lambda_d(O)$, bo kule są
parami rozłączne i zawarte w $O$. To kończy dowód.
\end{proof}

\begin{uwaga}[Dwa niezbędne ograniczenia w sformułowaniu]
Drobność jest warunkiem \emph{dla każdego punktu} $E$,
a nie tylko prawie wszędzie. Ograniczoność $O$ zapewnia
sumowalność objętości wybranych kul. Dla nieograniczonego
zbioru można podzielić przestrzeń na rozłączne warstwy
o skończonej mierze i zastosować argument lokalnie,
trzeba jednak zadbać o rozłączność wyborów pomiędzy warstwami.
Powyższa wersja wystarcza do różniczkowania miary Lebesgue'a.
\end{uwaga}

\begin{wniosek}[Przybliżanie zbioru rozłącznymi kulami]
\label{cor:vitali-open-24}
Jeżeli $O\subseteq\mathbb R^d$ jest niepustym,
ograniczonym zbiorem otwartym, to istnieją parami rozłączne
kule $B_j\subseteq O$ takie, że
$\lambda_d(O\setminus\bigcup_jB_j)=0$.
\end{wniosek}
\begin{proof}
Stosujemy twierdzenie~\ref{thm:vitali-cover-24} do $E=O$
i rodziny wszystkich kul o domknięciu zawartym w $O$.
Ta rodzina jest drobna w każdym punkcie $O$.
\end{proof}

\section{Geometria ograniczonej krotności: lemat kątowy}
\begin{lemat}[Skończoność konfiguracji Besicovitcha]
\label{lem:besicovitch-angle-24}
Niech $B_i=B(x_i,r_i)$, $i=1,\ldots,m$, będą kulami
w $\mathbb R^d$, których część wspólna zawiera punkt $p$.
Załóżmy ponadto, że dla $i\ne j$ zachodzi
$x_i\notin B_j$ i $x_j\notin B_i$.
Wówczas $m\le N_d$, gdzie $N_d$ zależy wyłącznie od $d$;
można przyjąć $N_d=3^d$.
\end{lemat}
\begin{proof}
Odrzucamy ewentualną kulę o środku $p$: skoro $p$ leży
we wszystkich kulach, taka kula wymusiłaby $m=1$.
Dla pozostałych stawiamy $u_i=(x_i-p)/|x_i-p|\in S^{d-1}$.
Ponieważ $p\in B_i$, mamy $|x_i-p|<r_i$.
Z warunku braku centrów w innych kulach wynika
$|x_i-x_j|\ge\max\{r_i,r_j\}>
\max\{|x_i-p|,|x_j-p|\}$.
Z twierdzenia cosinusów
\[
 2\langle x_i-p,x_j-p\rangle
 =|x_i-p|^2+|x_j-p|^2-|x_i-x_j|^2
 <\min\{|x_i-p|^2,|x_j-p|^2\}.
\]
W każdym przypadku z tej
nierówności otrzymujemy
$\langle u_i,u_j\rangle<1/2$, a zatem
$|u_i-u_j|>1$. Kulki $B(u_i,1/2)$ w $\mathbb R^d$
są parami rozłączne i leżą w $B(0,3/2)$.
Porównanie objętości daje
$m v_d(1/2)^d\le v_d(3/2)^d$, skąd $m\le3^d$.
\end{proof}

\begin{twierdzenie}[Skończona postać lematu Besicovitcha]
\label{thm:besicovitch-finite-24}
Niech $x_1,\ldots,x_M$ będą punktami $\mathbb R^d$,
a $r_1,\ldots,r_M>0$. Istnieje podrodzina kul
$B(x_i,r_i)$, która pokrywa wszystkie punkty
$x_1,\ldots,x_M$ i której krotność w żadnym punkcie
$\mathbb R^d$ nie przekracza $N_d=3^d$.
\end{twierdzenie}
\begin{proof}
Porządkujemy kule nierosnąco według promienia (dla remisów
wybieramy dowolny porządek). Przechodzimy po kolei przez
ich listę i wybieramy kulę, jeśli jej środek nie należy
do sumy kul wybranych wcześniej. Każdy środek jest ostatecznie
pokryty: albo przez własną wybraną kulę, albo przez jedną
z poprzednich. Jeśli wybrano $B_i$ wcześniej niż $B_j$,
to $x_j\notin B_i$, a stąd $|x_i-x_j|\ge r_i\ge r_j$
i także $x_i\notin B_j$. Wybrane kule spełniają więc
warunek lematu~\ref{lem:besicovitch-angle-24}.
Dla dowolnego punktu $p$ kule wybrane zawierające $p$
tworzą konfigurację z tego lematu, zatem jest ich najwyżej
$N_d$.
\end{proof}

\begin{wniosek}[Oszacowanie sumy miar w rodzinie skończonej]
\label{cor:finite-besicovitch-measure-24}
W założeniach twierdzenia~\ref{thm:besicovitch-finite-24}
niech $B_{i_1},\ldots,B_{i_s}$ będzie wybraną podrodziną.
Dla dowolnej dodatniej miary borelowskiej $\rho$
na $\mathbb R^d$ zachodzi
\[
 \sum_{j=1}^s\rho(B_{i_j})
 \le N_d\,\rho\Bigl(\bigcup_{j=1}^s B_{i_j}\Bigr).
\]
Teza pozostaje prawdziwa, gdy prawa strona jest nieskończona.
\end{wniosek}
\begin{proof}
Definicja krotności i skończona addytywność całki dają
\[
 \sum_{j=1}^s\rho(B_{i_j})
 =\int_{\mathbb R^d}\sum_{j=1}^s\chi_{B_{i_j}}(x)
 \,\mathrm d\rho(x)
 \le N_d\int_{\mathbb R^d}
 \chi_{\cup_jB_{i_j}}(x)\,\mathrm d\rho(x).
\]
To ostatnie wyrażenie jest równe prawej stronie tezy.
\end{proof}

\begin{twierdzenie}[Twierdzenie pokryciowe Besicovitcha -- pełny dowód]
\label{thm:besicovitch-24}
Niech $d\ge1$, niech $E\subseteq\mathbb R^d$ będzie
zbiorem ograniczonym i niech każdemu $x\in E$
przyporządkowano promień $r_x\in(0,R]$, gdzie
$R<\infty$. Wówczas z rodziny kul $\{B(x,r_x):x\in E\}$
można wybrać co najwyżej przeliczalną podrodzinę
$(B_j)_{j\in J}$ taką, że
\begin{equation}\label{eq:besicovitch-infinite-24}
 E\subseteq\bigcup_{j\in J}B_j,\qquad
 \sum_{j\in J}\chi_{B_j}(z)\le15^d
 \quad(z\in\mathbb R^d).
\end{equation}
Jeżeli dla każdego środka $x$ zadano kilka dopuszczalnych
kul o promieniach nie większych niż $R$, wolno najpierw
wybrać jedną z nich i zastosować twierdzenie.
Nie zakładamy mierzalności zbioru $E$.
\end{twierdzenie}
\begin{proof}
\emph{Krok 1: podział na skale.}
Dla $k=0,1,\ldots$ oznaczmy $a_k=2^{-k-1}R$ i
$E_k^0=\{x\in E:a_k<r_x\le2a_k\}$. Zbiory
$E_k^0$ tworzą rozłączny podział $E$.
Będziemy kolejno wybierać skończone zbiory środków
$S_k\subseteq E_k^0$; wybranej rodzinie w skali $k$
odpowiadają kule $\{B(x,r_x):x\in S_k\}$.

Przypuśćmy, że dokonano już wyboru dla skal
$0,\ldots,k-1$, i niech $O_{k-1}$ oznacza sumę
wszystkich wcześniej wybranych kul (dla $k=0$
przyjmujemy $O_{-1}=\varnothing$).
Niech $E_k=E_k^0\setminus O_{k-1}$. Spośród środków
$E_k$ wybieramy maksymalny zbiór $S_k$ o własności
\begin{equation}\label{eq:besicovitch-separated-24}
 |x-y|\ge a_k\quad(x,y\in S_k,\ x\ne y).
\end{equation}
Taki zbiór można zbudować zachłannie, bez lematu
Zorna: zaczynamy od dowolnego środka, a następnie
dodajemy kolejny, jeśli jego odległość od każdego
wcześniej wybranego jest co najmniej $a_k$.
Ponieważ $E$ jest ograniczony, każde dwa punkty
wybranego zbioru są oddalone o co najmniej $a_k$,
a kule $B(x,a_k/2)$ o tych środkach są rozłączne
i leżą we wspólnej ograniczonej kuli. Porównanie
ich objętości ogranicza z góry liczbę kroków;
procedura kończy się więc po skończenie wielu
wyborach, jeżeli $E_k\ne\varnothing$.

Z maksymalności wynika, że każdy $x\in E_k$
albo należy do $S_k$, albo istnieje $y\in S_k$
z $|x-y|<a_k<r_y$. Toteż wszystkie środki z
$E_k$ należą do kul wybranych w skali $k$.
W konsekwencji każdy $x\in E$ zostaje pokryty:
jeżeli $x\in E_k^0$, to albo znajdował się już
w $O_{k-1}$, albo zostaje pokryty w kroku $k$.
Suma $\bigcup_{k\ge0}S_k$ jest przeliczalna,
gdyż wszystkie $S_k$ są skończone.

\emph{Krok 2: liczba kul jednej skali.}
Ustalmy $z\in\mathbb R^d$ i skalę $k$.
Jeżeli $B(x,r_x)$ jest wybraną kulą skali $k$
zawierającą $z$, to $|x-z|<r_x\le2a_k$.
Z \eqref{eq:besicovitch-separated-24} wynika, że
kule $B(x,a_k/2)$ o takich środkach są rozłączne,
a każda z nich mieści się w $B(z,5a_k/2)$.
Porównując ich objętości, otrzymujemy oszacowanie
\begin{equation}\label{eq:besicovitch-single-scale-24}
 \#\{x\in S_k:z\in B(x,r_x)\}\le5^d.
\end{equation}

\emph{Krok 3: liczba skal występujących w punkcie.}
Spośród każdej skali $k$, dla której jakaś wybrana
kula zawiera $z$, wybierzmy jedną taką kulę
$B(x_k,r_{x_k})$. Gdy $k<\ell$, środek
$x_\ell$ nie należy do żadnej kuli wybranej w
skali $k$ (tak zdefiniowano $E_\ell$), więc
$|x_k-x_\ell|\ge r_{x_k}$. Ponadto
$r_{x_k}>a_k\ge r_{x_\ell}$, a zatem
$x_k\notin B(x_\ell,r_{x_\ell})$.
Tak wybrane kule wszystkich występujących skal
mają wspólny punkt $z$ i spełniają oba warunki
wyłączenia środków z
Lematu~\ref{lem:besicovitch-angle-24}.
Dowolna \emph{skończona} liczba tych skal jest
więc ograniczona przez $3^d$. Gdyby skal było
więcej niż $3^d$, otrzymalibyśmy sprzeczność
już dla pewnej skończonej podrodziny.

Każda skala wnosi najwyżej $5^d$ kul
zawierających $z$, a skal jest najwyżej $3^d$.
Daje to drugą nierówność w
\eqref{eq:besicovitch-infinite-24} i kończy dowód.
\end{proof}

\begin{uwaga}[Dlaczego metoda działa dla rodzin nieskończonych?]
Kluczowe jest oddzielenie dwóch rodzajów nakładania.
W jednej skali wystarcza pakowanie oddalonych
środków, natomiast pomiędzy skalami obowiązuje
silniejsza reguła: środek późniejszej kuli nie leży
w żadnej kuli wcześniej wybranej. Dopiero
połączenie obu argumentów daje jednostajną
krotność $15^d$. Nie przechodzimy tu od pokryć
skończonych do granicy i nie powiększamy żadnej kuli.
Stała nie jest optymalna; istotna jest jej
zależność wyłącznie od wymiaru.
\end{uwaga}

\begin{wniosek}[Besicovitch dla dodatnich miar borelowskich]
\label{cor:besicovitch-measure-infinite-24}
Przy założeniach Twierdzenia~\ref{thm:besicovitch-24},
dla dowolnej dodatniej miary borelowskiej $\rho$
na $\mathbb R^d$ wybrana rodzina kul spełnia
\[
 \sum_{j\in J}\rho(B_j)
 \le15^d\,\rho\Bigl(\bigcup_{j\in J}B_j\Bigr).
\]
\end{wniosek}
\begin{proof}
Stosujemy twierdzenie Tonellego do
$\sum_j\chi_{B_j}$ i korzystamy z oszacowania
krotności w \eqref{eq:besicovitch-infinite-24}.
\end{proof}

\section{Zastosowania do oszacowań typu maksymalnego}
\begin{definicja}[Funkcja maksymalna Hardy'ego--Littlewooda]
\label{def:HL-max-24}
Dla $f\in L^1_{\mathrm{loc}}(\mathbb R^d)$ definiujemy
\emph{centrowaną funkcję maksymalną}
\[
 Mf(x):=\sup_{r>0}\frac1{\lambda_d(B(x,r))}
 \int_{B(x,r)}|f(y)|\,\mathrm d\lambda_d(y)
 \in[0,+\infty].
\]
Funkcja $Mf$ jest mierzalna, co wynika z faktu, że dla
ustalonego $r$ średnia po $B(x,r)$ jest funkcją ciągłą
zmiennej $x$, oraz z możliwości ograniczenia supremum
do promieni wymiernych dodatnich. Ciągłość średniej wynika
z ciągłości translacji w $L^1$ na odpowiednich zbiorach
zwartych, a ciągłość względem $r>0$ z bezwzględnej
ciągłości całki po kulach.
\end{definicja}

\begin{twierdzenie}[Nierówność słabego typu $(1,1)$]
\label{thm:HL-weak-24}
Jeśli $f\in L^1(\mathbb R^d)$, to dla każdego $t>0$
\[
 \lambda_d\bigl(\{x\in\mathbb R^d:Mf(x)>t\}\bigr)
 \le\frac{5^d}{t}\int_{\mathbb R^d}|f(x)|\,
 \mathrm d\lambda_d(x).
\]
\end{twierdzenie}
\begin{proof}
Ustalmy zwarty zbiór $K\subseteq\{Mf>t\}$.
Dla każdego $x\in K$ wybieramy promień $r_x>0$ taki,
że $\int_{B(x,r_x)}|f|\,\mathrm d\lambda_d>
 t\lambda_d(B(x,r_x))$. Promienie takich kul są
ograniczone z góry: ponieważ ich objętość jest mniejsza
od $\|f\|_1/t$, zachodzi $r_x^d<\|f\|_1/(tv_d)$.
Z lematu $5r$ wybieramy rozłączne kule $(B_j)$, których
pięciokrotne powiększenia pokrywają $K$. Wówczas
\[
 \lambda_d(K)\le5^d\sum_j\lambda_d(B_j)
 <\frac{5^d}{t}\sum_j\int_{B_j}|f|\,\mathrm d\lambda_d
 \le\frac{5^d}{t}\|f\|_1.
\]
Zbiór $\{Mf>t\}$ jest otwarty: dla dowolnego $x$
warunek średniej większej od $t$ zachowuje się przy małej
zmianie środka kuli dzięki ciągłości całki z $|f|$.
Jego miara jest supremum miar zwartych podzbiorów,
na mocy regularności miary Lebesgue'a. Biorąc supremum
po $K$, kończymy dowód.
\end{proof}

\begin{wniosek}[Kontrola zbioru, na którym średnie są duże]
\label{cor:HL-level-24}
Dla $f\in L^1(\mathbb R^d)$ oraz $a>0$ funkcja $Mf$
jest skończona prawie wszędzie, a
$\lambda_d(\{Mf>a\})\le5^d\|f\|_1/a$.
\end{wniosek}
\begin{proof}
Oszacowanie wynika z twierdzenia~\ref{thm:HL-weak-24}.
Przechodząc do granicy $a\to\infty$, dostajemy
$\lambda_d(\{Mf=+\infty\})=0$.
\end{proof}

\begin{uwaga}[Co zostawiamy do następnego rozdziału?]
W rozdziale~\ref{ch:covering-24} skonstruowaliśmy narzędzie
pokryciowe i udowodniliśmy nierówność słabego typu.
W kolejnym rozdziale użyjemy ich do dowodu twierdzenia
Lebesgue'a o różniczkowaniu całek i miar, bez zakładania
go w argumentach niniejszego rozdziału.
\end{uwaga}

% END chapters/pl/measure_theory/24-rozdzial.tex


% BEGIN chapters/pl/measure_theory/25-rozdzial.tex
% ROZDZIAŁ 25 — RÓŻNICZKOWANIE CAŁEK I MIAR
\chapter{Różniczkowanie całek i miar}
\label{ch:differentiation-25}

\section{Dlaczego różniczkowanie jest twierdzeniem o średnich?}
Niech $d\ge1$, a $\lambda_d$ oznacza miarę Lebesgue'a na
$\mathbb R^d$. Dla $x\in\mathbb R^d$ i $r>0$ piszemy
$B(x,r)=\{y:|y-x|<r\}$ oraz $v_d=\lambda_d(B(0,1))$.
Jeżeli $f\in L^1_{\mathrm{loc}}(\mathbb R^d)$, to na pierwszy
rzut oka naturalne jest oczekiwanie, że średnia $f$ po kuli
malejącej do punktu $x$ dąży do $f(x)$. Nie zachodzi to
we wszystkich punktach ani dla dowolnego przedstawiciela
funkcji. Poniższa teoria precyzyjnie określa, w jakim
sensie odtwarzamy funkcję z jej całek.

\begin{definicja}[Wartość średnia i punkt Lebesgue'a]
\label{def:lebesgue-point-25}
Dla $f\in L^1_{\mathrm{loc}}(\mathbb R^d)$ definiujemy
$A_rf(x)=\lambda_d(B(x,r))^{-1}\int_{B(x,r)}f(y)\,\mathrm d\lambda_d(y)$.
Punkt $x$ nazywamy \emph{punktem Lebesgue'a} wybranego
przedstawiciela $f$, jeżeli
\begin{equation}\label{eq:lebesgue-point-25}
 \lim_{r\downarrow0}\frac1{v_dr^d}
 \int_{B(x,r)}|f(y)-f(x)|\,\mathrm d\lambda_d(y)=0.
\end{equation}
W szczególności $A_rf(x)\to f(x)$ w takim punkcie.
Przypominamy, że wartości $f(x)$ na zbiorze miary zero
można zmieniać bez zmiany klasy elementu $L^1_{\mathrm{loc}}$.
\end{definicja}

\begin{przyklad}[Nieciągłość w zbiorze zerowym]
Jeżeli $f=\chi_{\mathbb Q^d}$, to $A_rf(x)=0$ dla
wszystkich $x,r$, ponieważ $\lambda_d(\mathbb Q^d)=0$.
Żaden punkt wymierny nie jest punktem Lebesgue'a tego
konkretnego przedstawiciela, a każdy niewymierny jest.
Równość prawie wszędzie nie może więc zostać zastąpiona
równością w każdym punkcie.
\end{przyklad}

\section{Wzmocnienie nierówności maksymalnej}
\begin{definicja}[Maksymalna funkcja miary]
\label{def:max-measure-25}
Jeżeli $\nu$ jest skończoną dodatnią miarą borelowską
na $\mathbb R^d$, określamy
$M\nu(x)=\sup_{r>0}\nu(B(x,r))/(v_dr^d)$.
Dla miary ze znakiem bądź zespolonej podstawiamy jej
wariację całkowitą $|\nu|$.
\end{definicja}

\begin{twierdzenie}[Oszacowanie słabego typu dla miar]
\label{thm:weak-measure-25}
Niech $\nu$ będzie skończoną dodatnią miarą borelowską
na $\mathbb R^d$. Dla każdego $t>0$ zachodzi
\begin{equation}\label{eq:weak-measure-25}
 \lambda_d\bigl(\{x:M\nu(x)>t\}\bigr)
 \le \frac{5^d}{t}\nu(\mathbb R^d).
\end{equation}
To samo oszacowanie obowiązuje z $\nu$ zastąpioną przez
$|\nu|$, gdy $\nu$ jest miarą ze znakiem lub zespoloną.
\end{twierdzenie}
\begin{proof}
Dla ustalonego $r>0$ funkcja $x\mapsto\nu(B(x,r))$
jest półciągła z dołu: jeśli $x_n\to x$, to dla każdego
$y\in B(x,r)$ ostatecznie $y\in B(x_n,r)$, zatem lemat
Fatou zastosowany do miary $\nu$ daje
$\nu(B(x,r))\le\liminf_n\nu(B(x_n,r))$. Stąd
$\{M\nu>t\}$ jest otwarty. Niech $K$ będzie jego
zwartym podzbiorem. Dla każdego $x\in K$ istnieje
kula $B_x=B(x,r_x)$ taka, że
$\nu(B_x)>tv_dr_x^d$. Ponieważ $\nu$ jest skończona,
promienie są ograniczone z góry. Z lematu $5r$
(twierdzenie~\ref{thm:5r-24}) wybieramy rozłączne
kule $B_j$ o pokryciu $K\subseteq\bigcup_j5B_j$.
Wtedy
\[
\lambda_d(K)\le5^d\sum_jv_dr_j^d
 <\frac{5^d}{t}\sum_j\nu(B_j)
 \le\frac{5^d}{t}\nu(\mathbb R^d).
\]
Korzystając z regularności miary Lebesgue'a na zbiorze
otwartym $\{M\nu>t\}$, przechodzimy do supremum po
jego zwartych podzbiorach. W przypadku miary zespolonej
lub ze znakiem stosujemy otrzymany wynik do $|\nu|$.
\end{proof}

\begin{uwaga}[Nie mylić dwóch operatorów maksymalnych]
Dla $f\in L^1(\mathbb R^d)$ miara
$\nu_f(A)=\int_A|f(y)|\,\mathrm d\lambda_d(y)$ spełnia
$M\nu_f=Mf$, gdzie $Mf$ zdefiniowano w
\ref{def:HL-max-24}. Operator $M\nu$ pozwala jednak
również badać miary osobliwe, nieposiadające gęstości.
\end{uwaga}

\section{Twierdzenie Lebesgue'a o różniczkowaniu całek}
\begin{lemat}[Przybliżanie lokalne funkcją ciągłą]
\label{lem:continuous-dense-25}
Dla dowolnego $f\in L^1(\mathbb R^d)$ i
$\varepsilon>0$ istnieje $\varphi\in C_c(\mathbb R^d)$
taka, że $\|f-\varphi\|_1<\varepsilon$.
\end{lemat}
\begin{proof}
Jest to szczególny przypadek twierdzenia o gęstości
$C_c^\infty(\mathbb R^d)$ w $L^1$, dowiedzionego w
rozdziale~\ref{ch:regularnosc-funkcji}. W szczególności nie
zakładamy tutaj różniczkowania całek przy konstrukcji
aproksymacji.
\end{proof}

\begin{twierdzenie}[Lebesgue'a o różniczkowaniu]
\label{thm:lebesgue-diff-25}
Jeżeli $f\in L^1_{\mathrm{loc}}(\mathbb R^d)$, to
prawie każdy $x\in\mathbb R^d$ jest punktem Lebesgue'a
$f$, tzn. spełnia \eqref{eq:lebesgue-point-25}.
W szczególności
$\lim_{r\downarrow0}A_rf(x)=f(x)$ dla prawie każdego $x$.
\end{twierdzenie}
\begin{proof}
Najpierw załóżmy $f\in L^1(\mathbb R^d)$. Ustalmy
$\varphi\in C_c(\mathbb R^d)$ i połóżmy $g=f-\varphi$.
Dla każdego $x$, dla którego $f(x)$ jest określona,
z nierówności trójkąta otrzymujemy
\begin{align*}
 \limsup_{r\downarrow0}\frac1{v_dr^d}
 \int_{B(x,r)}|f(y)-f(x)|\,\mathrm d\lambda_d(y)
 &\le Mg(x)+|g(x)|,
\end{align*}
gdzie $Mg=M(|g|\lambda_d)$. Składnik z funkcją
$\varphi$ znika wskutek jej ciągłości w $x$.
Niech $D_f(x)$ oznacza powyższy $\limsup$, a $t>0$.
Dla każdego $\varphi$ zachodzi inkluzja
$\{D_f>2t\}\subseteq\{Mg>t\}\cup\{|g|>t\}$.
Nierówność słabego typu i nierówność Czebyszewa dają
\[
 \lambda_d(\{D_f>2t\})
 \le\frac{5^d+1}{t}\|f-\varphi\|_1.
\]
Lemat~\ref{lem:continuous-dense-25} pozwala uczynić
prawą stronę dowolnie małą. Zatem
$\lambda_d(\{D_f>2t\})=0$ dla każdego $t>0$.
Ponieważ $D_f\ge0$, otrzymujemy $D_f=0$ prawie
wszędzie. (Mierzalność $D_f$ wynika z mierzalności
średnich i z wyrażenia granicy górnej przy promieniach
wymiernych, wykorzystując ciągłość całki względem
promienia.)

Jeżeli $f$ jest tylko lokalnie całkowalna, dla każdego
$n\ge1$ rozważamy $f_n=f\chi_{B(0,n+2)}\in L^1$.
Dla $x\in B(0,n)$ i $r<1$ średnie funkcji
$f_n$ i $f$ są jednakowe, a wartości obu funkcji
w $x$ są równe. Zastosowanie pierwszej części do
wszystkich $n$ i przeliczalna suma zbiorów zerowych
dają tezę.
\end{proof}

\begin{wniosek}[Wartość funkcji z otaczających ją średnich]
\label{cor:diff-basis-25}
Jeżeli $f\in L^1_{\mathrm{loc}}(\mathbb R^d)$,
to istnieje zbiór $N$ o $\lambda_d(N)=0$, taki że
$\lim_{r\downarrow0}A_rf(x)=f(x)$ dla wszystkich
$x\in\mathbb R^d\setminus N$. Zbioru wyjątkowego
nie można w ogólności usunąć.
\end{wniosek}
\begin{proof}
Pierwsza część wynika z twierdzenia
\ref{thm:lebesgue-diff-25}; druga z przykładu
$f=\chi_{\mathbb Q^d}$.
\end{proof}

\begin{przyklad}[Przedstawiciel kanoniczny w punktach Lebesgue'a]
Jeżeli $f\in L^1_{\mathrm{loc}}(\mathbb R^d)$, to
na zbiorze punktów, w których istnieje skończona granica
$\lim_{r\downarrow0}A_rf(x)$, można ją przyjąć za
wartość przedstawiciela $\widetilde f(x)$. W pozostałych
punktach przyjmujemy $0$. Otrzymujemy funkcję mierzalną
$\widetilde f=f$ prawie wszędzie. Zmiana przedstawiciela
$f$ na zbiorze zerowym nie zmienia wartości żadnej średniej,
więc $\widetilde f$ nie zależy od tego wyboru.
\end{przyklad}

\section{Różniczkowanie miar: część absolutnie ciągła i osobliwa}
\begin{lemat}[Zanik gęstości miary osobliwej]
\label{lem:singular-density-25}
Niech $\nu$ będzie skończoną dodatnią miarą borelowską
na $\mathbb R^d$ i $\nu\perp\lambda_d$. Wówczas
\begin{equation}\label{eq:singular-density-25}
 \lim_{r\downarrow0}\frac{\nu(B(x,r))}{v_dr^d}=0
 \quad\text{dla }\lambda_d\text{-prawie każdego }x.
\end{equation}
\end{lemat}
\begin{proof}
Wybierzmy borelowski $N$ taki, że $\lambda_d(N)=0$
i $\nu(\mathbb R^d\setminus N)=0$. Regularność
skończonej miary borelowskiej daje dla każdego
$\varepsilon>0$ zwarty $K\subseteq N$ taki, że
$\nu(\mathbb R^d\setminus K)<\varepsilon$.
Oznaczmy $\nu_0=\nu|_K$, $\nu_1=\nu|_{\mathbb R^d\setminus K}$.
Dla każdego $x\notin K$ odległość $\operatorname{dist}(x,K)$
jest dodatnia, zatem $\nu_0(B(x,r))=0$ dla dostatecznie
małych $r$. W konsekwencji, poza $K$,
\[
 \limsup_{r\downarrow0}\frac{\nu(B(x,r))}{v_dr^d}
 \le M\nu_1(x).
\]
Niech $t>0$. Z twierdzenia~\ref{thm:weak-measure-25}
i z $\lambda_d(K)=0$ otrzymujemy
\[
 \lambda_d\left(\left\{x:\limsup_{r\downarrow0}
 \frac{\nu(B(x,r))}{v_dr^d}>t\right\}\right)
 \le\frac{5^d\varepsilon}{t}.
\]
Ponieważ $\varepsilon$ jest dowolne, zbiór ten ma
miarę zero. Dla $t=1/n$, $n\ge1$, dostajemy
zerowanie granicy górnej prawie wszędzie. Ilorazy są
nieujemne, więc istnieje granica równa zero.
\end{proof}

\begin{twierdzenie}[Różniczkowanie miary względem Lebesgue'a]
\label{thm:measure-differentiation-25}
Niech $\nu$ będzie dodatnią miarą Radona na
$\mathbb R^d$. Jej rozkład Lebesgue'a względem
$\lambda_d$ ma postać $\nu=f\lambda_d+\nu_s$,
gdzie $f\ge0$ jest lokalnie całkowalna, a
$\nu_s\perp\lambda_d$. Wtedy
\begin{equation}\label{eq:measure-density-25}
 \lim_{r\downarrow0}\frac{\nu(B(x,r))}{v_dr^d}
 =f(x)\qquad\text{dla }\lambda_d\text{-prawie każdego }x.
\end{equation}
Dla lokalnie skończonych miar ze znakiem oraz miar
zespolonych twierdzenie pozostaje prawdziwe, jeżeli
$f$ jest odpowiednią pochodną Radona--Nikodyma.
\end{twierdzenie}
\begin{proof}
Z rozkładu Lebesgue'a i twierdzenia Radona--Nikodyma
(rozdziały~\ref{ch:ac-singular-17} i
\ref{ch:radon-nikodym-18}) otrzymujemy
$\nu(B)=\int_Bf(y)\,\mathrm d\lambda_d(y)+\nu_s(B)$.
Pierwszy składnik po podzieleniu przez $v_dr^d$
zmierza do $f(x)$ prawie wszędzie z twierdzenia
\ref{thm:lebesgue-diff-25}. Dla składnika osobliwego
lokalizujemy miarę na kulach $B(0,n)$ i stosujemy
lemat~\ref{lem:singular-density-25}; przy
$x\in B(0,n-1)$ obie miary mają takie same wartości
na dostatecznie małych kulach. Przeliczalna suma
zbiorów wyjątkowych jest zerowa. W przypadku miar ze
znakiem i zespolonych stosujemy ten sam argument do
części absolutnie ciągłej, a do wariacji części osobliwej
stosujemy lemat~\ref{lem:singular-density-25}.
\end{proof}

\begin{uwaga}[Ważne rozróżnienie]
Istnienie granicy w \eqref{eq:measure-density-25} dla
$\lambda_d$-prawie wszystkich $x$ \emph{nie} oznacza,
że $\nu_s=0$. Dla miary Cantora pochodna gęstości
względem $\lambda_1$ jest zerowa prawie wszędzie,
chociaż sama miara jest niezerowa.
\end{uwaga}

\begin{przyklad}[Miara Diraca]
Niech $\nu=\delta_0$ na $\mathbb R^d$. Jeżeli $x\ne0$,
to dla wszystkich dostatecznie małych $r$ zachodzi
$\delta_0(B(x,r))=0$. Natomiast
$\delta_0(B(0,r))/(v_dr^d)=1/(v_dr^d)\to+\infty$.
Miara osobliwa może mieć więc nieskończoną gęstość
w swoich punktach koncentracji.
\end{przyklad}

\section{Pochodna całki w jednym wymiarze}
\begin{twierdzenie}[Podstawowe twierdzenie rachunku całkowego dla całki Lebesgue'a]
\label{thm:FTC-lebesgue-25}
Niech $a<b$ i $f\in L^1((a,b),\lambda_1)$. Funkcja
$F:[a,b]\to\mathbb R$ określona wzorem
$F(x)=\int_{(a,x)}f(t)\,\mathrm d\lambda_1(t)$
jest absolutnie ciągła i spełnia
$F'(x)=f(x)$ dla prawie każdego $x\in(a,b)$.
\end{twierdzenie}
\begin{proof}
Dla skończonej rodziny parami rozłącznych przedziałów
$(a_j,b_j)\subseteq[a,b]$ mamy
$\sum_j|F(b_j)-F(a_j)|
\le\int_{\cup_j(a_j,b_j)}|f|\,\mathrm d\lambda_1$.
Bezwzględna ciągłość całki zapewnia, że dla danego
$\varepsilon>0$ prawa strona jest mniejsza od
$\varepsilon$, gdy suma długości przedziałów jest
dostatecznie mała. Jest to definicja absolutnej
ciągłości funkcji $F$.

Niech $x\in(a,b)$ będzie punktem Lebesgue'a $f$.
Dla $h>0$ tak małego, by $x+h<b$, zachodzi
\[
 \left|\frac{F(x+h)-F(x)}h-f(x)\right|
 \le\frac1h\int_{(x,x+h)}|f(t)-f(x)|\,
 \mathrm d\lambda_1(t)
 \le\frac2{2h}\int_{(x-h,x+h)}|f(t)-f(x)|\,
 \mathrm d\lambda_1(t).
\]
Ostatnia średnia zmierza do zera. Analogiczny argument
obowiązuje dla $h<0$. Prawie wszystkie punkty są
punktami Lebesgue'a na mocy twierdzenia
\ref{thm:lebesgue-diff-25}, zatem $F'=f$ prawie
wszędzie.
\end{proof}

\begin{przyklad}[Dlaczego sama pochodna prawie wszędzie nie wystarcza?]
Klasyczna funkcja Cantora $C:[0,1]\to[0,1]$ jest
ciągła i niemalejąca, a $C'=0$ prawie wszędzie, lecz
$C(1)-C(0)=1$. Nie jest absolutnie ciągła.
Dlatego z samej równości $F'=f$ prawie wszędzie nie
wynika $F(b)-F(a)=\int_{(a,b)}f\,\mathrm d\lambda_1$.
\end{przyklad}

\begin{twierdzenie}[Funkcja rozkładu a absolutna ciągłość miary]
\label{thm:cdf-derivative-25}
Niech $\nu$ będzie skończoną dodatnią miarą borelowską
na $\mathbb R$. Określmy $F:\mathbb R\to[0,\nu(\mathbb R)]$
wzorem $F(x)=\nu(( -\infty,x])$.
Jeśli $\nu=f\lambda_1+\nu_s$ jest rozkładem Lebesgue'a
miary $\nu$ względem $\lambda_1$, to
$F'(x)=f(x)$ dla $\lambda_1$-prawie każdego $x$.
Ponadto następujące warunki są równoważne:
\begin{enumerate}[label=(\roman*)]
\item $\nu\ll\lambda_1$;
\item $F$ jest absolutnie ciągła na każdym zwartym
przedziale;
\item dla wszystkich $a<b$ zachodzi
$F(b)-F(a)=\int_{(a,b)}F'(t)\,\mathrm d\lambda_1(t)$.
\end{enumerate}
W (iii) pochodna istnieje prawie wszędzie.
\end{twierdzenie}
\begin{proof}
Niech $x$ będzie punktem Lebesgue'a $f$, w którym
$\nu_s(B(x,r))/(2r)\to0$ oraz $\nu(\{x\})=0$.
Punkty takie tworzą zbiór pełnej miary Lebesgue'a:
pierwsze dwie własności wynikają z twierdzenia
\ref{thm:measure-differentiation-25}, a atomów dodatniej
masy jest co najwyżej przeliczalnie wiele.
Dla $h>0$ mamy
\[
 \frac{F(x+h)-F(x)}h
 =\frac1h\int_{(x,x+h]}f(t)\,\mathrm d\lambda_1(t)
 +\frac{\nu_s((x,x+h])}h.
\]
Pierwszy składnik dąży do $f(x)$, ponieważ średnia
wartości $|f(t)-f(x)|$ po $(x,x+h)$ jest ograniczona
przez dwa razy analogiczną średnią po $(x-h,x+h)$.
Drugi składnik, dla $h>0$, jest ograniczony przez
$\nu_s(B(x,2h))/h$ i zatem dąży do zera.
Dla $h<0$ rozpatrujemy przedział $(x+h,x]$;
ponieważ $\nu_s(\{x\})=0$, ten sam argument działa
również z lewej. Dowodzi to $F'=f$ prawie wszędzie.

Jeśli (i) zachodzi, to $\nu=f\lambda_1$ z
$f\in L^1(\mathbb R)$, a więc
$F(x)=\int_{(-\infty,x]}f(t)\,\mathrm d\lambda_1(t)$.
Bezwzględna ciągłość całki na zbiorach o małej mierze
daje (ii), zaś podstawowe twierdzenie rachunku
całkowego (twierdzenie~\ref{thm:FTC-lebesgue-25})
daje (iii).

Załóżmy (ii) i wykażmy najpierw (i), bez powoływania
się na twierdzenie o całkowaniu pochodnej dowolnej funkcji
absolutnie ciągłej. Ustalmy $n\ge1$ i zbiór borelowski
$E\subseteq[-n,n]$ z $\lambda_1(E)=0$. Funkcja $F$
jest absolutnie ciągła na $[-n-1,n+1]$, więc dla
każdego $\varepsilon>0$ istnieje $\delta>0$ takie,
że dla każdej skończonej rodziny rozłącznych przedziałów
zawartych w $[-n-1,n+1]$, o łącznej długości poniżej
$\delta$, suma przyrostów $F$ jest mniejsza od
$\varepsilon$. Pokryjmy $E$ otwartym zbiorem
$O\subset(-n-1,n+1)$ o $\lambda_1(O)<\delta$.
Składowymi $O$ są parami rozłączne przedziały otwarte
$(a_j,b_j)$. Ponieważ $F$ jest ciągła,
$\nu(\{x\})=F(x)-F(x^-)=0$ dla każdego $x$
wewnątrz $[-n-1,n+1]$, a więc
$\nu((a_j,b_j))=F(b_j)-F(a_j)$.
Stosując nierówność absolutnej ciągłości do skończonych
sum i przechodząc do granicy z przeliczalnej addytywności,
dostajemy $\nu(E)\le\nu(O)\le\varepsilon$.
Zatem $\nu(E)=0$. Ponieważ $n$ był dowolny,
$\nu\ll\lambda_1$ na całej prostej. Wykazane już
(i)$\Rightarrow$(iii) daje również (ii)$\Rightarrow$(iii).
Załóżmy wreszcie (iii). Ponieważ
$\nu((a,b])=F(b)-F(a)$, dla każdych $a<b$ mamy
$\nu((a,b])=\int_{(a,b)}F'(t)\,\mathrm d\lambda_1(t)$;
przy całce nie ma znaczenia wybór końców.
Funkcja $F'\ge0$ prawie wszędzie i jest lokalnie
całkowalna, gdyż z (iii) jej całki po przedziałach
są skończone. Obie miary są skończone na zbiorach
ograniczonych i zgadzają się na przedziałach $(a,b]$.
Jednoznaczność miary na generującym $\pi$-układzie
i przejście przez zbiory $(-n,n]$ pokazują
$\nu(E)=\int_EF'(t)\,\mathrm d\lambda_1(t)$
dla wszystkich $E\in\mathcal B(\mathbb R)$.
To daje (i).
\end{proof}

\begin{przyklad}[Pochodna nie wykrywa składnika osobliwego]
W przypadku miary Cantora funkcja rozkładu $F$ jest
ciągła, niemalejąca i spełnia $F'=0$ prawie wszędzie,
ale nie jest absolutnie ciągła. Twierdzenie
\ref{thm:cdf-derivative-25} wyjaśnia przyczynę:
pochodna odzyskuje wyłącznie część absolutnie ciągłą
miary, a nie składnik osobliwy.
\end{przyklad}

\section{Przenoszenie różniczkowania na inne rodziny zbiorów}
\begin{wniosek}[Średnie po kostkach]
\label{cor:cubes-differentiation-25}
Niech $f\in L^1_{\mathrm{loc}}(\mathbb R^d)$. Dla prawie
każdego $x$ średnia $f$ po dowolnych kostkach osiowych
$Q_r(x)=x+(-r,r)^d$ dąży do $f(x)$ przy $r\downarrow0$.
Ponadto $\lambda_d(Q_r(x))^{-1}
\int_{Q_r(x)}|f(y)-f(x)|\,\mathrm d\lambda_d(y)\to0$.
\end{wniosek}
\begin{proof}
Mamy $Q_r(x)\subseteq B(x,\sqrt d\,r)$ oraz
$\lambda_d(B(x,\sqrt d\,r))/\lambda_d(Q_r(x))
=v_dd^{d/2}/2^d$, niezależnie od $r$. Dodatnia średnia
po kostce jest więc ograniczona przez stałą razy
średnia po większej kuli, która zmierza do zera
w każdym punkcie Lebesgue'a.
\end{proof}

\begin{uwaga}[Przekroje i inne miary]
Analogiczne różniczkowanie względem \emph{dowolnej}
miary Radona $\mu$ na $\mathbb R^d$ można otrzymać
z twierdzenia pokryciowego Besicovitcha:
centrowany operator maksymalny względem $\mu$ spełnia
odpowiednie oszacowanie słabego typu $(1,1)$, a dalszy
argument opiera się na aproksymacji funkcji całkowalnych.
Pełna, nieskończona wersja twierdzenia Besicovitcha
\emph{została już udowodniona} w
Twierdzeniu~\ref{thm:besicovitch-24}.
Nie wyprowadzamy tutaj jednak osobno całego twierdzenia
o różniczkowaniu względem arbitralnej miary Radona:
wymaga ono dodatkowo sformułowania i przeprowadzenia
argumentu maksymalnego z mianownikiem $\mu(B(x,r))$,
z uwzględnieniem kul o zerowej mierze.
W ten sposób rozróżniamy udowodnione narzędzie pokryciowe
od niewypisanego w pełni wniosku o różniczkowaniu.
\end{uwaga}

% END chapters/pl/measure_theory/25-rozdzial.tex


% BEGIN chapters/pl/measure_theory/26-rozdzial.tex
% ROZDZIAŁ 26 — MIARA I WYMIAR HAUSDORFFA, TWIERDZENIE FROSTMANA
\chapter{Miara Hausdorffa i wymiar Hausdorffa}
\label{ch:hausdorff-26}

\section{Dlaczego miara Lebesgue'a nie wystarcza?}
Krzywa w $\mathbb R^2$ ma zazwyczaj zerową miarę
$\lambda_2$, choć może mieć dodatnią długość.
Zbiór Cantora ma zerową długość, ale posiada nietrywialną
strukturę w skali ułamkowego wymiaru. Celem tego
rozdziału jest skonstruowanie miar $\mathcal H^s$,
które mierzą zbiory odpowiednio do ich wymiaru, oraz
wyprowadzenie skutecznych kryteriów tego wymiaru.

\section{Konstrukcja miary Hausdorffa od podstaw}
\label{sec:hausdorff-konstrukcja-26}

W tej sekcji przeprowadzimy całą konstrukcję, łącznie z
mierzalnością zbiorów borelowskich. Nie będziemy zakładać
twierdzenia, że metryczna miara zewnętrzna automatycznie
mierzy zbiory borelowskie: udowodnimy potrzebny fakt.
Wykorzystamy jedynie elementarne własności kresów,
przeliczalnych pokryć i odległości w przestrzeni metrycznej.

\begin{definicja}[Koszt pokrycia i miara Hausdorffa]
\label{def:hausdorff-measure-26}
Niech $(X,d)$ będzie przestrzenią metryczną, $s\geq0$
i $\delta>0$. Dla $U\subseteq X$ definiujemy
$\operatorname{diam}U:=\sup\{d(x,y):x,y\in U\}$,
przyjmując $\operatorname{diam}\varnothing=0$.
Ustalamy $c_s>0$, przy czym $c_0=1$.
Dla $s>0$ koszt zbioru to $h_s(U)=c_s(\operatorname{diam}U)^s$.
Dla $s=0$ przyjmujemy jawnie
$h_0(U)=0$ dla $U=\varnothing$ oraz $h_0(U)=1$
dla $U\neq\varnothing$, niezależnie od średnicy.
Definiujemy
\begin{equation}\label{eq:hausdorff-delta-26}
 \mathcal H^s_\delta(E)
 :=\inf\left\{\sum_{j=1}^\infty h_s(U_j):
 E\subseteq\bigcup_{j=1}^\infty U_j,\quad
 \operatorname{diam}U_j\leq\delta\right\}.
\end{equation}
Infimum pustej rodziny dopuszczalnych pokryć wynosi
$+\infty$; zbiór pusty ma koszt $0$ (pokrycie pustymi
zbiorami). \emph{Miarę Hausdorffa wymiaru $s$} określamy
wzorem
\begin{equation}\label{eq:hausdorff-26}
 \mathcal H^s(E):=\lim_{\delta\downarrow0}
 \mathcal H^s_\delta(E)
 =\sup_{\delta>0}\mathcal H^s_\delta(E).
\end{equation}
Jest to dobrze określona wartość z $[0,+\infty]$,
gdyż $0<\delta'<\delta$ pociąga za sobą
$\mathcal H^s_{\delta'}(E)\geq\mathcal H^s_\delta(E)$.
Dla $s>0$ popularna normalizacja to
$c_s=\omega_s2^{-s}$, gdzie
$\omega_s=\pi^{s/2}/\Gamma(1+s/2)$;
$\Gamma$ oznacza funkcję gamma Eulera, tj.
$\Gamma(t)=\int_{(0,+\infty)}u^{t-1}e^{-u}\,\mathrm d\lambda_1(u)$
dla $t>0$. Można jednak przyjąć $c_s=1$, nie zmieniając
własności mierzalności ani wymiaru Hausdorffa.
W przypadku $s=1$ standardowy wybór daje $c_1=1$.
\end{definicja}

Ważne jest rozdzielenie dwóch operacji: najpierw
wybieramy \emph{dowolne przeliczalne} pokrycie zbiorami
o średnicy co najwyżej $\delta$ i minimalizujemy koszt;
dopiero potem wymuszamy $\delta\downarrow0$.
Nie ograniczamy pokryć do kul ani do zbiorów mierzalnych.
W przestrzeni nieośrodkowej może nie istnieć przeliczalne
pokrycie danego $E$ o małych średnicach; wtedy
odpowiednia wartość jest po prostu nieskończona.

\begin{uwaga}[Miara a zawartość Hausdorffa]
Gdy w~\eqref{eq:hausdorff-delta-26} nie narzucamy
ograniczenia średnic, otrzymujemy \emph{zawartość
Hausdorffa} $\mathcal H^s_\infty(E)$. Sama zawartość
nie musi być addytywna nawet dla rozłącznych
zbiorów odległych od siebie: jedno duże pokrycie
może być tańsze od dwóch oddzielnych.
Dlatego warunek $\delta\downarrow0$ jest istotny.
\end{uwaga}

\begin{lemat}[Zewnętrzne własności kosztu pokryć]
\label{lem:hausdorff-outer-26}
Dla każdego $s\geq0$ i $\delta>0$ funkcja
$E\mapsto\mathcal H^s_\delta(E)$ spełnia:
$\mathcal H^s_\delta(\varnothing)=0$,
monotoniczność i przeliczalną podaddytywność.
To samo dotyczy funkcji $E\mapsto\mathcal H^s(E)$.
\end{lemat}
\begin{proof}
Pusta rodzina punktów nie wymaga żadnego pokrycia,
a pokrycie złożone ze zbiorów pustych ma koszt zero.
Jeśli $A\subseteq B$, każde pokrycie $B$ pokrywa $A$,
więc infimum dla $A$ nie przekracza infimum dla $B$.

Niech $E\subseteq\bigcup_{k\ge1}E_k$. Jeśli suma
$\sum_k\mathcal H^s_\delta(E_k)$ jest nieskończona,
nierówność podaddytywności jest automatyczna.
W przeciwnym razie wszystkie składniki są skończone.
Dla $\varepsilon>0$ i każdego $k$ wybieramy dopuszczalne
pokrycie $(U_{kj})_{j\ge1}$ zbioru $E_k$ o koszcie
nieprzekraczającym
$\mathcal H^s_\delta(E_k)+\varepsilon2^{-k}$.
Rodzina wszystkich $U_{kj}$ jest przeliczalna,
ma średnice co najwyżej $\delta$ i pokrywa $E$.
W konsekwencji
$\mathcal H^s_\delta(E)\leq
\sum_k\mathcal H^s_\delta(E_k)+\varepsilon$.
Przechodząc $\varepsilon\downarrow0$, otrzymujemy
podaddytywność na ustalonej skali.

Przy przejściu do granicy \emph{nie zamieniamy}
supremum z nieskończoną sumą. Dla każdego $\delta$
z poprzedniej nierówności i z definicji supremum mamy
$\mathcal H^s_\delta(E)\leq
\sum_k\mathcal H^s_\delta(E_k)
\leq\sum_k\mathcal H^s(E_k)$.
Wzięcie supremum po $\delta>0$ kończy dowód.
\end{proof}

\begin{lemat}[Addytywność na zbiorach oddalonych]
\label{lem:hausdorff-metricity-26}
Jeżeli $A,B\subseteq X$ i
$\operatorname{dist}(A,B):=
\inf\{d(a,b):a\in A,\ b\in B\}>0$,
to $\mathcal H^s(A\cup B)=\mathcal H^s(A)+\mathcal H^s(B)$.
Wniosek pozostaje prawdziwy dla każdej skończonej
rodziny parami dodatnio oddalonych zbiorów.
\end{lemat}
\begin{proof}
Ustalmy $0<\delta<\operatorname{dist}(A,B)$.
Żaden zbiór o średnicy nieprzekraczającej $\delta$
nie może przecinać jednocześnie $A$ i $B$.
Każde pokrycie $A\cup B$ rozdzielamy zatem na
podrodzinę pokrywającą $A$ i podrodzinę pokrywającą
$B$; zbiory nieprzecinające żadnego z nich pomijamy.
Suma kosztów pierwotnego pokrycia jest co najmniej
$\mathcal H^s_\delta(A)+\mathcal H^s_\delta(B)$.
Biorąc infimum, otrzymujemy nierówność
$\mathcal H^s_\delta(A\cup B)\geq
\mathcal H^s_\delta(A)+\mathcal H^s_\delta(B)$.
Po przejściu do granicy $\delta\downarrow0$
i użyciu lematu~\ref{lem:hausdorff-outer-26}
uzyskujemy równość. Dla skończonej rodziny
stosujemy ten argument indukcyjnie: suma skończonej
liczby zbiorów odległych od pozostałych ma nadal
od nich dodatnią odległość.
\end{proof}

\begin{lemat}[Kryterium Carathéodory'ego, z dowodem]
\label{lem:hausdorff-caratheodory-selfcontained-26}
Niech $m^*:2^X\to[0,+\infty]$ będzie dowolną
miarą zewnętrzną, czyli $m^*(\varnothing)=0$,
$m^*$ jest monotoniczna i przeliczalnie
podaddytywna. Zbiór $A\subseteq X$ nazywamy
$m^*$-mierzalnym, jeśli dla każdego $E\subseteq X$
\begin{equation}\label{eq:hausdorff-caratheodory-test-26}
 m^*(E)=m^*(E\cap A)+m^*(E\setminus A).
\end{equation}
Rodzina $\mathcal M_{m^*}$ wszystkich takich zbiorów
jest $\sigma$-algebrą, a $m^*|_{\mathcal M_{m^*}}$
jest miarą zupełną.
\end{lemat}
\begin{proof}
Zbiór pusty i $X$ spełniają warunek, a zamiana $A$
na jego dopełnienie jedynie przestawia składniki.
Jeżeli $A,B$ spełniają~\eqref{eq:hausdorff-caratheodory-test-26},
to dwukrotnie rozcinając $E$ otrzymujemy
\[
 m^*(E)=m^*(E\cap A\cap B)
 +m^*(E\cap A\setminus B)+m^*(E\setminus A)
 \geq m^*(E\cap A\cap B)+m^*(E\setminus(A\cap B)).
\]
Odwrotną nierówność daje podaddytywność, więc
$A\cap B$ jest mierzalny. Dopełnienia dają
zamkniętość na skończone sumy i różnice.

Niech $(A_j)$ będzie ciągiem \emph{rozłącznych}
zbiorów mierzalnych, $B=\bigcup_{j\ge1}A_j$
i $B_N=\bigcup_{j=1}^NA_j$. Wielokrotne
zastosowanie kryterium daje dla każdego $E\subseteq X$
\[
 m^*(E)=\sum_{j=1}^Nm^*(E\cap A_j)
 +m^*(E\setminus B_N)
 \geq\sum_{j=1}^Nm^*(E\cap A_j)+m^*(E\setminus B).
\]
Przechodzimy $N\to\infty$, a następnie stosujemy
podaddytywność do $E\cap B=\bigcup_j(E\cap A_j)$.
Otrzymujemy nierówność $m^*(E)\geq
m^*(E\cap B)+m^*(E\setminus B)$;
odwrotna wynika z podaddytywności. Zatem $B$
jest mierzalny. Dla dowolnego ciągu zbiorów
mierzalnych rozłączność uzyskujemy, zastępując
$A_j$ przez $A_j\setminus\bigcup_{i<j}A_i$.
To dowodzi, że $\mathcal M_{m^*}$ jest $\sigma$-algebrą.

Dla rozłącznych mierzalnych $A_j$ stosujemy
poprzedni rachunek do $E=B$; po przejściu
do granicy otrzymujemy
$m^*(B)\geq\sum_jm^*(A_j)$.
Podaddytywność daje nierówność odwrotną,
więc mamy przeliczalną addytywność.
Wreszcie, gdy $m^*(N)=0$, dla dowolnego $A\subseteq N$
i $E\subseteq X$ zachodzi
$m^*(E\cap A)=0$ oraz
$m^*(E\setminus A)\leq m^*(E)$.
Podaddytywność wymusza~\eqref{eq:hausdorff-caratheodory-test-26}
dla $A$. Każdy podzbiór zbioru zerowego jest zatem
mierzalny i zerowy: miara jest zupełna.
\end{proof}

\begin{lemat}[Metryczna miara zewnętrzna a zbiory domknięte]
\label{lem:hausdorff-metric-open-26}
Załóżmy, że miara zewnętrzna $m^*$ spełnia
$m^*(A\cup B)=m^*(A)+m^*(B)$ zawsze, gdy
$\operatorname{dist}(A,B)>0$. Wtedy każdy zbiór
domknięty $F\subseteq X$ spełnia kryterium
Carathéodory'ego~\eqref{eq:hausdorff-caratheodory-test-26}.
W szczególności wszystkie zbiory otwarte są mierzalne.
\end{lemat}
\begin{proof}
Dla $F=\varnothing$ lub $F=X$ teza jest oczywista.
Załóżmy, że $F$ jest niepusty i właściwy, i ustalmy
$E\subseteq X$. Jeśli $m^*(E)=+\infty$,
podaddytywność już daje równość z kryterium,
więc możemy przyjąć $m^*(E)<\infty$.
Funkcja $\rho(x)=\operatorname{dist}(x,F)$
jest $1$-Lipschitzowska, $\rho^{-1}(0)=F$.
Dla $k\ge0$ wprowadzamy warstwy
$S_k=E\cap\{2^{-(k+1)}<\rho\leq2^{-k}\}$.
Warstwy $S_0,S_2,S_4,\ldots$ są parami dodatnio
oddalone: między zakresami wartości $\rho$ dwóch
kolejnych parzystych warstw leży niepusty przedział,
a $|\rho(x)-\rho(y)|\leq d(x,y)$.
To samo dotyczy warstw o nieparzystych indeksach.
Z addytywności na skończonych rodzinach oddalonych
i z monotoniczności wynika zatem
$\sum_{j\ge0}m^*(S_{2j})\leq m^*(E)$ oraz
$\sum_{j\ge0}m^*(S_{2j+1})\leq m^*(E)$.
Wobec tego szereg $\sum_km^*(S_k)$ jest zbieżny,
a jego ogony dążą do zera.

Połóżmy $E_n=E\cap\{\rho>2^{-n}\}$. Mamy
$E\setminus F=\bigcup_nE_n$ oraz
$(E\setminus F)\setminus E_n\subseteq\bigcup_{k\ge n}S_k$.
Z podaddytywności i monotoniczności dostajemy
\[
 m^*(E_n)\leq m^*(E\setminus F)
 \leq m^*(E_n)+\sum_{k\ge n}m^*(S_k).
\]
W szczególności $m^*(E_n)\to m^*(E\setminus F)$.
Ponieważ zbiory $E\cap F$ i $E_n$ mają odległość
co najmniej $2^{-n}$, metryczna addytywność daje
$m^*(E)\geq m^*(E\cap F)+m^*(E_n)$.
Przechodząc do granicy, otrzymujemy
$m^*(E)\geq m^*(E\cap F)+m^*(E\setminus F)$.
Podaddytywność daje nierówność odwrotną.
Zatem $F$ jest mierzalny, a jego dopełnienie
jest mierzalne na mocy poprzedniego lematu.
\end{proof}

\begin{twierdzenie}[Samodzielna konstrukcja miary Hausdorffa]
\label{thm:hausdorff-borel-26}
Dla dowolnej przestrzeni metrycznej $(X,d)$ i $s\geq0$
funkcja $\mathcal H^s$ z~\eqref{eq:hausdorff-26}
jest metryczną miarą zewnętrzną. Rodzina jej zbiorów
mierzalnych w sensie Carathéodory'ego jest zupełną
$\sigma$-algebrą zawierającą $\mathcal B(X)$.
Ograniczenie $\mathcal H^s$ do $\mathcal B(X)$ jest
miarą borelowską (niekoniecznie zupełną).
\end{twierdzenie}
\begin{proof}
Lemat~\ref{lem:hausdorff-outer-26} daje
wszystkie aksjomaty miary zewnętrznej,
a lemat~\ref{lem:hausdorff-metricity-26}
jej metryczną addytywność.
Lemat~\ref{lem:hausdorff-metric-open-26} mówi, że
każdy zbiór domknięty, a więc również każdy otwarty,
jest mierzalny. Rodzina zbiorów mierzalnych jest
$\sigma$-algebrą na mocy
lematu~\ref{lem:hausdorff-caratheodory-selfcontained-26};
zawiera więc $\sigma$-algebrę generowaną przez
zbiory otwarte, czyli $\mathcal B(X)$.
Ten sam lemat daje przeliczalną addytywność
i zupełność na pełnej $\sigma$-algebrze
Carathéodory'ego. Nie użyliśmy żadnego twierdzenia
o mierzalności borelowskiej metrycznych miar zewnętrznych
jako przesłanki.
\end{proof}

\begin{twierdzenie}[Regularizacja przez zbiory $G_\delta$]
\label{thm:hausdorff-gdelta-26}
Dla dowolnego $E\subseteq X$ istnieje zbiór
$G\in\mathcal B(X)$ będący przeliczalnym
przecięciem zbiorów otwartych, taki że
$E\subseteq G$ i $\mathcal H^s(G)=\mathcal H^s(E)$.
W szczególności miara Hausdorffa jest zewnętrznie
regularna względem zbiorów borelowskich.
\end{twierdzenie}
\begin{proof}
Gdy $\mathcal H^s(E)=+\infty$, wystarczy wziąć $G=X$.
Niech $M:=\mathcal H^s(E)<\infty$ i $\delta_n=2^{-n}$.
Z definicji infimum wybieramy dla każdego $n$
pokrycie $E$ zbiorami $U_{nj}$ o średnicach
nieprzekraczających $\delta_n/2$, o łącznym
koszcie co najwyżej $M+1/(2n)$.
Każdy niepusty $U_{nj}$ powiększamy do
otwartego otoczenia $V_{nj}$ o średnicy
mniejszej niż $\delta_n$, wybierając grubość
otoczenia tak małą, by jego koszt przekraczał
koszt $U_{nj}$ najwyżej o $2^{-j}/(2n)$.
Jest to możliwe, ponieważ dla $s>0$ funkcja
$t\mapsto c_st^s$ jest ciągła, a dla $s=0$
koszt każdego niepustego zbioru pozostaje równy $1$.
Połóżmy $O_n=\bigcup_jV_{nj}$ i $G=\bigcap_nO_n$.
Wtedy $E\subseteq G$, a $O_n$ jest otwarty i
\[
 \mathcal H^s_{\delta_n}(O_n)
 \leq\sum_jh_s(V_{nj})\leq M+1/n.
\]
Dla dowolnego $\delta>0$ i każdego dostatecznie
dużego $n$ (gdy $\delta_n<\delta$) mamy
$\mathcal H^s_\delta(G)\leq
\mathcal H^s_{\delta_n}(O_n)\leq M+1/n$.
Po przejściu $n\to\infty$ otrzymujemy
$\mathcal H^s_\delta(G)\leq M$.
Biorąc supremum po $\delta>0$,
dostajemy $\mathcal H^s(G)\leq M$;
monotoniczność daje równość.
\end{proof}

\begin{przyklad}[Wymiar zero: miara licząca]
\label{prz:hausdorff-counting-26}
Dla każdej przestrzeni metrycznej i zbioru $E\subseteq X$
miara $\mathcal H^0(E)$ jest liczbą punktów $E$,
z wartością $+\infty$, gdy $E$ jest nieskończony.
Istotnie, dla $E=\{x_1,\ldots,x_N\}$
wybieramy $\delta$ mniejszą niż odległość
każdej pary różnych punktów.
Każdy element pokrycia o średnicy $\leq\delta$
może wtedy zawierać najwyżej jeden punkt $E$;
pokrycie singletonami osiąga koszt $N$.
Jeżeli $E$ jest nieskończony, dla każdego $N$
wybieramy $N$ różnych punktów i stosujemy
monotoniczność: $\mathcal H^0(E)\geq N$.
\end{przyklad}

\begin{twierdzenie}[W wymiarze jeden otrzymujemy miarę Lebesgue'a]
\label{thm:hausdorff-line-26}
Dla standardowej normalizacji $c_1=1$ oraz
każdego $E\subseteq\mathbb R$ zachodzi
$\mathcal H^1(E)=\lambda_1^*(E)$.
Tutaj $\lambda_1^*(E)$ oznacza infimum
$\sum_j(b_j-a_j)$ po wszystkich przeliczalnych
pokryciach $E$ otwartymi przedziałami $(a_j,b_j)$.
Jest to konstrukcja miary zewnętrznej Lebesgue'a
na prostej, podana również w rozdziale czwartym.
\end{twierdzenie}
\begin{proof}
Każdy niepusty zbiór $U\subseteq\mathbb R$
o skończonej średnicy $a$ zawiera się w przedziale
domkniętym długości $a$ (między swoim infimum
i supremum). Dla dowolnego pokrycia
$E\subseteq\bigcup_jU_j$ dostajemy więc
$\lambda_1^*(E)\leq\sum_j\operatorname{diam}U_j$.
Po przejściu do infimum i granicy
$\lambda_1^*(E)\leq\mathcal H^1(E)$.

W drugą stronę, jeśli $\lambda_1^*(E)<\infty$,
z definicji miary zewnętrznej Lebesgue'a
wybieramy przeliczalne pokrycie $E$ otwartymi
przedziałami $I_j$ o łącznej długości
mniejszej od $\lambda_1^*(E)+\varepsilon$.
Dzielimy każdy $I_j$ na skończenie wiele
przedziałów domkniętych o długościach
nieprzekraczających $\delta$; suma długości
tych odcinków jest równa długości $I_j$.
Otrzymujemy dopuszczalne pokrycie Hausdorffa,
zatem $\mathcal H^1_\delta(E)\leq
\lambda_1^*(E)+\varepsilon$.
Ponieważ $\delta,\varepsilon>0$ były dowolne,
$\mathcal H^1(E)\leq\lambda_1^*(E)$.
W przypadku $\lambda_1^*(E)=+\infty$
pierwsza nierówność już kończy dowód.
\end{proof}

\begin{twierdzenie}[Skalowanie i odwzorowania Lipschitzowskie]
\label{thm:hausdorff-lipschitz-26}
Niech $(X,d_X)$ i $(Y,d_Y)$ będą przestrzeniami
metrycznymi, $s\ge0$, a $T:X\to Y$ będzie
odwzorowaniem $L$-Lipschitzowskim, tj.
$d_Y(Tx,Ty)\le Ld_X(x,y)$ dla wszystkich $x,y$,
gdzie $L>0$. Wówczas dla każdego $E\subseteq X$
$\mathcal H_Y^s(T(E))\le L^s\mathcal H_X^s(E)$.
W szczególności dla podobieństwa $T$ przestrzeni
euklidesowej o skali $L$ zachodzi
$\mathcal H^s(T(E))=L^s\mathcal H^s(E)$.
\end{twierdzenie}
\begin{proof}
Jeżeli $E\subseteq\bigcup_jU_j$ i
$\operatorname{diam}U_j\le\delta$, to
$T(E)\subseteq\bigcup_jT(U_j)$ i
$\operatorname{diam}T(U_j)\le L\operatorname{diam}U_j
\le L\delta$. Wobec tego
$\mathcal H^s_{L\delta}(T(E))\le
L^s\mathcal H^s_\delta(E)$.
Przechodzimy $\delta\downarrow0$. W przypadku
podobieństwa stosujemy tę samą nierówność także do
$T^{-1}$, którego stała Lipschitza wynosi $L^{-1}$.
\end{proof}

\begin{wniosek}[Punkty i zbiory przeliczalne]
\label{cor:hausdorff-singletons-auditVI}
Niech $(X,d)$ będzie przestrzenią metryczną i $s>0$.
Wtedy $\mathcal H^s(\{x\})=0$ dla każdego $x\in X$.
W konsekwencji $\mathcal H^s(A)=0$ dla każdego zbioru
przeliczalnego $A\subseteq X$. Jeśli $T:X\to Y$ jest
odwzorowaniem lokalnie Lipschitzowskim w sensie pokrycia
$X$ przeliczalnie wieloma zbiorami, na których $T$
jest Lipschitzowskie, to
$\mathcal H_Y^s(T(A))=0$ zawsze, gdy $\mathcal H_X^s(A)=0$.
\end{wniosek}
\begin{proof}
Dla każdego $\delta>0$ zbiór jednoelementowy
można pokryć nim samym, którego koszt w wymiarze
$s>0$ wynosi zero. Przeliczalna podaddytywność
miary zewnętrznej dowodzi drugiego zdania.
W ostatnim zdaniu piszemy $X=\bigcup_jX_j$,
przy czym ograniczenie $T|_{X_j}$ ma stałą
Lipschitza $L_j$. Na mocy
Twierdzenia~\ref{thm:hausdorff-lipschitz-26}
$\mathcal H_Y^s(T(A\cap X_j))\le
L_j^s\mathcal H_X^s(A\cap X_j)=0$.
Ponownie stosujemy podaddytywność.
\end{proof}

\section{Wymiar Hausdorffa jako próg zanikania miar}
\begin{lemat}[Próg zerowania]
\label{lem:hausdorff-threshold-26}
Niech $E\subseteq X$ oraz $0\le s<t$. Jeżeli
$\mathcal H^s(E)<\infty$, to $\mathcal H^t(E)=0$.
Jeżeli $\mathcal H^t(E)>0$, to $\mathcal H^s(E)=+\infty$.
\end{lemat}
\begin{proof}
Rozważmy niepusty zbiór $U_j$ w pokryciu,
którego średnica nie przekracza $\delta$.
Z definicji kosztu wynika
\[
 h_t(U_j)\le\frac{c_t}{c_s}\,\delta^{t-s}h_s(U_j).
\]
Dla $s=0$ stosujemy przy tym $h_0(U_j)=1$.
Dla zbioru pustego obie strony są zerowe.
Sumując koszty i biorąc infimum, otrzymujemy
\[
\mathcal H^t_\delta(E)\le
\frac{c_t}{c_s}\,\delta^{t-s}\mathcal H^s_\delta(E).
\]
Jeżeli $\mathcal H^s(E)<\infty$, prawa strona dąży
do zera. Druga implikacja jest kontrapozycją pierwszej.
\end{proof}

\begin{definicja}[Wymiar Hausdorffa]
\label{def:hausdorff-dimension-26}
Dla $E\subseteq X$ definiujemy
\[
 \dim_H E:=\inf\{s\ge0:\mathcal H^s(E)=0\}
 =\sup\{s\ge0:\mathcal H^s(E)=+\infty\},
\]
z konwencją $\sup\varnothing=0$ (w szczególności
$\dim_H\varnothing=0$) i $\inf\varnothing=+\infty$.
Przy wartości granicznej $s=\dim_H E$ miara
$\mathcal H^s(E)$ może być zerowa, skończona dodatnia
lub nieskończona.
\end{definicja}

\begin{wniosek}[Wymiar sumy przeliczalnej]
\label{cor:hausdorff-union-26}
Dla dowolnych zbiorów $E_n\subseteq X$ zachodzi
$\dim_H(\bigcup_{n\ge1}E_n)=\sup_n\dim_HE_n$.
W szczególności zbiory przeliczalne mają wymiar zero.
\end{wniosek}
\begin{proof}
Monotoniczność daje nierówność $\ge$.
Jeżeli $s>\sup_n\dim_HE_n$, to
$\mathcal H^s(E_n)=0$ dla każdego $n$; przeliczalna
podaddytywność daje $\mathcal H^s(\bigcup_nE_n)=0$.
Po przejściu $s$ do supremum otrzymujemy nierówność
odwrotną. Każdy pojedynczy punkt ma miarę
$\mathcal H^s$ równą zero dla $s>0$, gdyż można
pokryć go zbiorem o dowolnie małej średnicy.
\end{proof}

\begin{twierdzenie}[Wymiar przestrzeni euklidesowej]
\label{thm:euclidean-dimension-26}
Jeżeli $E\subseteq\mathbb R^d$ jest mierzalny
w sensie Lebesgue'a i $\lambda_d(E)>0$, to $\dim_HE=d$. Każdy podzbiór
$\mathbb R^d$ ma wymiar Hausdorffa co najwyżej $d$.
Co więcej, $\mathcal H^d$ i $\lambda_d$ są wzajemnie
porównywalne na zbiorach borelowskich, z dodatnimi
stałymi zależnymi tylko od $d$.
\end{twierdzenie}
\begin{proof}
Niech $A\subseteq\mathbb R^d$ będzie ograniczony.
Każdy zbiór $U$ o średnicy $\rho$ mieści się w domkniętej
kuli o promieniu $\rho$ (o środku w dowolnym punkcie $U$),
a zatem $\lambda_d^*(U)\le v_d\rho^d$.
Dla pokrycia $A\subseteq\bigcup_jU_j$ zbiorami o
średnicach $\le\delta$ mamy
$\lambda_d^*(A)\le v_d\sum_j(\operatorname{diam}U_j)^d$.
Stąd $\mathcal H^d(A)\ge(c_d/v_d)\lambda_d^*(A)$.

Dla ograniczonego $A$ wybierzmy otwarty ograniczony
zbiór $O\supseteq A$ o mierze dowolnie bliskiej
$\lambda_d^*(A)$. Z ustalonej siatki diadycznej wybieramy maksymalne
kostki półotwarte $Q$, których domknięcia mieszczą się
w $O$ i których boki nie przekraczają $\delta/\sqrt d$.
Dla każdego $x\in O$ istnieje dostatecznie mała kostka
siatki zawierająca $x$ i mająca domknięcie w $O$;
zatem domknięcia wybranych kostek pokrywają $O$.
Wnętrza maksymalnych kostek są parami rozłączne
i zawarte w $O$, toteż suma ich objętości nie
przekracza $\lambda_d(O)$. Każda domknięta kostka
ma średnicę równą $\sqrt d$ razy długość boku,
więc jej koszt $s=d$ jest nie większy od $d^{d/2}$
razy jej objętości. Z tego wynika
$\mathcal H^d_\delta(A)\le c_dd^{d/2}\lambda_d(O)$.
Przechodząc do infimum po $O$ i do granicy po $\delta$,
dostajemy porównywalność na ograniczonych zbiorach,
a następnie na całej przestrzeni przez wyczerpanie
kulami. W szczególności $\mathcal H^d(E)>0$, jeśli
$\lambda_d(E)>0$.

Dla $s>d$, pokrycie kostki o boku $L$ przez
$O((L/\delta)^d)$ małych kostek średnicy $O(\delta)$
daje $\mathcal H^s_\delta([0,L]^d)
\le C_{d,s,L}\delta^{s-d}\to0$.
Całe $\mathbb R^d$ jest przeliczalną sumą takich
kostek, więc $\dim_HE\le d$ dla każdego $E$.
Połączenie z dodatniością $\mathcal H^d(E)$ daje
$\dim_HE=d$ przy $\lambda_d(E)>0$.
\end{proof}

\begin{twierdzenie}[Kalibracja wymiaru całkowitego bez twierdzeń geometrycznych]
\label{thm:hausdorff-calibration-26}
Niech $d\ge1$ będzie liczbą całkowitą. Zapiszmy
$\widetilde{\mathcal H}^{d}$ dla miary Hausdorffa
na $\mathbb R^d$ skonstruowanej powyżej z $c_d=1$
i połóżmy $\kappa_d=\widetilde{\mathcal H}^{d}([0,1]^d)$.
Wtedy $0<\kappa_d<\infty$ oraz dla każdego
zbioru borelowskiego $E\subseteq\mathbb R^d$
zachodzi
\[
 \widetilde{\mathcal H}^{d}(E)=\kappa_d\lambda_d(E).
\]
W szczególności wybór $c_d=\kappa_d^{-1}$ daje
miarę Hausdorffa spełniającą
$\mathcal H^d|_{\mathcal B(\mathbb R^d)}=\lambda_d$.
Nie potrzebujemy do tego nierówności izodiametrycznej
ani twierdzenia o mierze Haara.
\end{twierdzenie}
\begin{proof}
Dodatniość i skończoność $\kappa_d$ wynikają z
udowodnionych wyżej dwustronnych oszacowań
$\widetilde{\mathcal H}^{d}$ przez miarę Lebesgue'a.
Definicja przez średnice pokazuje bezpośrednio,
że $\widetilde{\mathcal H}^{d}$ jest niezmiennicza
względem translacji i skaluje się przez $a^d$
przy jednokładności skali $a>0$.

Wykażmy najpierw, że brzeg sześcianu ma zerową
$\widetilde{\mathcal H}^{d}$-miarę.
Każda z jego $2d$ ścian daje się pokryć przez
$O(\delta^{-(d-1)})$ kostek o średnicy nie większej
niż $\delta$; łączny koszt w wymiarze $d$ wynosi
$O(\delta)$ i dąży do zera. Wobec tego
$\widetilde{\mathcal H}^{d}([0,1)^d)=\kappa_d$.
Z niezmienniczości i skalowania wynika, że każdy
półotwarty sześcian o boku $a$ ma miarę
$\kappa_da^d$.

Niech $R=\prod_{i=1}^d[a_i,b_i)$ będzie
prostokątem o wymiernych końcach.
Wybieramy dodatnią liczbę wymierną $h$, która
dzieli długość każdego boku $b_i-a_i$.
Podział na skończenie wiele półotwartych sześcianów
o boku $h$ jest rozłączny, zatem przeliczalna
addytywność daje
$\widetilde{\mathcal H}^{d}(R)
=\kappa_d\prod_{i=1}^d(b_i-a_i)
=\kappa_d\lambda_d(R)$.

Rodzina takich prostokątów jest układem $\pi$
generującym $\mathcal B(\mathbb R^d)$.
Obie miary są skończone na prostokątach ograniczonych,
a $\mathbb R^d$ wyczerpuje rosnący ciąg prostokątów
o wymiernych końcach. Stosując twierdzenie
$\pi$--$\lambda$ o jednoznaczności miar (którego
dowód podaliśmy w rozdziale pierwszym) do
każdego z tych prostokątów, otrzymujemy równość
na wszystkich zbiorach borelowskich.
\end{proof}

\begin{uwaga}[Konwencja normalizacji w dalszych rozdziałach]
\label{rem:hausdorff-convention-26}
Od tego miejsca w całym skrypcie, dla każdego całkowitego
$k\ge1$, symbol $\mathcal H^k$ oznacza miarę Hausdorffa
zdefiniowaną w Definicji~\ref{def:hausdorff-measure-26}
z wyborem $c_k=\kappa_k^{-1}$ ustalonym w
Twierdzeniu~\ref{thm:hausdorff-calibration-26}.
W szczególności $\mathcal H^k$ na $\mathbb R^k$
zgadza się z $\lambda_k$; to samo dotyczy każdej
izometrycznie osadzonej $k$-płaszczyzny w $\mathbb R^N$.
Dla $k=0$ zachowujemy $\mathcal H^0$ jako miarę liczącą.
Dla niecałkowitych $s>0$ pozostajemy przy dowolnej,
raz ustalonej dodatniej stałej $c_s$; jej wybór nie
zmienia wymiaru Hausdorffa. W konsekwencji wzory
na pole powierzchni, koobszar i obwód w rozdziałach
\ref{ch:area-coarea-28}--\ref{ch:geometric-measure-29}
nie zawierają niejawnych czynników normalizacyjnych.
\end{uwaga}

\begin{uwaga}[Stała geometryczna a stała kalibracyjna]
Często stosuje się jawny współczynnik
$c_d=\omega_d2^{-d}$, gdzie $\omega_d$ jest
objętością euklidesowej kuli jednostkowej.
Wykazanie, że ten współczynnik jest równy
$\kappa_d^{-1}$, wymaga jeszcze nierówności
izodiametrycznej. Nie korzystamy z tej
identyczności: powyższa kalibracja jest
zdefiniowana w pełni wewnętrznie i wystarcza
do wzorów geometrycznych w rozdziale 28.
\end{uwaga}

\section{Zasada rozkładu masy}
\begin{twierdzenie}[Zasada rozkładu masy]
\label{thm:mass-distribution-26}
Niech $E\subseteq\mathbb R^d$ będzie zbiorem borelowskim,
$s>0$, a $\mu$ skończoną dodatnią miarą borelowską
spełniającą $\mu(E)>0$ i $\mu(\mathbb R^d\setminus E)=0$.
Załóżmy, że istnieje $C>0$ takie, iż
$\mu(B(x,r))\le Cr^s$ dla wszystkich
$x\in\mathbb R^d$ i $r>0$. Wtedy
$\mathcal H^s(E)\ge(c_s/(2^sC))\mu(E)>0$,
a więc $\dim_HE\ge s$.
\end{twierdzenie}
\begin{proof}
Niech $E\subseteq\bigcup_jU_j$, gdzie wszystkie
$U_j$ mają średnicę dodatnią nieprzekraczającą $\delta$.
Jeśli $U_j\cap E\ne\varnothing$, wybierzmy
$x_j\in U_j\cap E$. Dla dodatniej średnicy
$U_j\subseteq B(x_j,2\operatorname{diam}U_j)$.
Nie musimy zakładać mierzalności $U_j$: zbiory
$E\cap U_j$ są zawarte w wymienionych kulach,
a funkcja $\mu$ jest przeliczalnie podaddytywna
na pokryciach przez kule borelowskie. Zbiory pokrycia
o zerowej średnicy są co najwyżej singletonami,
a $\mu$ nie ma atomów (wynika to z oszacowania
miary kul), więc ich łączna miara zewnętrzna
wynosi zero. Ostatecznie
$\mu(E)\le C2^s\sum_j(\operatorname{diam}U_j)^s$.
Przeliczalna podaddytywność i przejście do infimum dają
$\mu(E)\le(C2^s/c_s)\mathcal H^s_\delta(E)$
po przejściu do infimum po pokryciach. W granicy
$\delta\downarrow0$ otrzymujemy tezę.
\end{proof}

\section{Twierdzenie Frostmana — konstrukcja miary}
\begin{lemat}[Dodatnia zawartość zbioru zwartego]
\label{lem:positive-content-26}
Niech $K\subseteq\mathbb R^d$ będzie zwarty i $s>0$.
Jeżeli $\mathcal H^s(K)>0$, to
$\mathcal H^s_\infty(K)>0$.
\end{lemat}
\begin{proof}
Załóżmy przeciwnie, że $\mathcal H^s_\infty(K)=0$.
Dla każdego $\varepsilon>0$ można pokryć $K$
zbiorami $U_j$ takimi, że
$\sum_j(\operatorname{diam}U_j)^s<\varepsilon/c_s$.
Każda średnica w tym pokryciu jest mniejsza niż
$(\varepsilon/c_s)^{1/s}$. Wybierając
$\varepsilon$ dowolnie małe, otrzymujemy
$\mathcal H^s_\delta(K)=0$ dla każdego $\delta>0$,
a więc $\mathcal H^s(K)=0$, sprzeczność.
\end{proof}

\begin{twierdzenie}[Frostmana, wersja dla zbiorów zwartych]
\label{thm:frostman-26}
Niech $K\subseteq\mathbb R^d$ będzie niepustym
zbiorem zwartym i $s>0$. Następujące warunki
są równoważne:
\begin{enumerate}[label=(\roman*)]
 \item $\mathcal H^s(K)>0$;
 \item istnieje niezerowa skończona dodatnia miara
 borelowska $\mu$, skoncentrowana na $K$, oraz stała
 $C>0$ takie, że
 $\mu(B(x,r))\le Cr^s$ dla wszystkich
 $x\in\mathbb R^d$ i $r>0$.
\end{enumerate}
Ponadto w (ii) miarę można znormalizować do
prawdopodobieństwa na $K$.
\end{twierdzenie}
\begin{proof}
Implikacja (ii)$\Rightarrow$(i) wynika z zasady
rozkładu masy~\ref{thm:mass-distribution-26}.
Pokażemy konstruktywnie implikację odwrotną, nie
zakładając twierdzenia o przepływach w grafach.

\emph{Krok 1: drzewo kostek.} Wybierzmy półotwartą
kostkę osiową $Q_0$ o boku $L>0$, tak aby
$K\subset\operatorname{int}Q_0$. Dzielimy ją na
$2^d$ kostek potomnych o połowie boku i powtarzamy
tę operację do poziomu $n$. Każda kostka na poziomie
$j$ ma bok $L2^{-j}$. Zachowujemy także kostki
nieprzecinające $K$, ale nadajemy im zerową pojemność.
Każdej kostce $Q$ przypisujemy pojemność
\[
 w(Q)=\begin{cases}
 (\sqrt d\,\operatorname{side}(Q))^s,&Q\cap K\ne\varnothing,\\
 0,&Q\cap K=\varnothing.
 \end{cases}
\]
Zdefiniujmy liczby $a_n(Q)$ rekurencyjnie od poziomu
$n$ do korzenia: dla liści $a_n(Q)=w(Q)$, a dla
pozostałych węzłów
\begin{equation}\label{eq:frostman-recursion-26}
 a_n(Q)=\min\left\{w(Q),\sum_{Q'\text{ potomek bezpośredni }Q}
 a_n(Q')\right\}.
\end{equation}
Indukcja od liści pokazuje, że $a_n(Q)$ jest minimalną
sumą pojemności kostek z poziomów od poziomu $Q$
do $n$, których suma pokrywa $K\cap Q$ w strukturze
drzewa. Istotnie, wybieramy albo samą kostkę $Q$
(koszt $w(Q)$), albo pokrywamy oddzielnie każdy z jej
$2^d$ potomków (koszt sumy minimów). Każde takie
pokrycie wyznacza pokrycie zbioru $K$ przez domknięcia
wybranych kostek, których średnice nie przekraczają
wartości użytych w kosztach. Stąd, na mocy lematu
\ref{lem:positive-content-26},
\begin{equation}\label{eq:frostman-positive-26}
 a_n(Q_0)\ge\mathcal H^s_\infty(K)/c_s=:m_0>0.
\end{equation}
Równocześnie $a_n(Q_0)\le w(Q_0)$.

\emph{Krok 2: rozdzielanie masy po drzewie.} Nadajemy
korzeniowi masę $m_n(Q_0)=a_n(Q_0)$.
Jeżeli masa $m_n(Q)$ jest już określona i
$m_n(Q)\le a_n(Q)$, to ponieważ
$\sum_{Q'}a_n(Q')\ge a_n(Q)\ge m_n(Q)$,
można rozdzielić $m_n(Q)$ między potomków $Q'$
tak, by $0\le m_n(Q')\le a_n(Q')$ oraz
$\sum_{Q'}m_n(Q')=m_n(Q)$. Można to zrobić
kolejno, napełniając potomków do ich pojemności.
Na poziomie $n$ wybieramy dla każdej kostki
$Q$ o dodatniej masie punkt $x_Q\in K\cap Q$
i definiujemy miarę atomową
$\mu_n=\sum_{Q\text{ liść}}m_n(Q)\delta_{x_Q}$.
Z konstrukcji otrzymujemy
\begin{equation}\label{eq:frostman-dyadic-bound-26}
 m_0\le\mu_n(K)\le w(Q_0),\qquad
 \mu_n(Q)\le w(Q)
 \quad\text{dla każdej kostki }Q\text{ o poziomie }\le n.
\end{equation}

\emph{Krok 3: oszacowanie kul.} Niech $0<r<L/2$
i wybierzmy poziom $j$ taki, że
$r\le\ell_j=L2^{-j}<2r$ (przy $r$ wystarczająco
małym taki $j\ge0$ istnieje). Kula $B(x,r)$
przecina co najwyżej $4^d$ półotwartych kostek
diadycznych poziomu $j$: w każdym kierunku jej
rzut ma długość $2r\le2\ell_j$, więc przecina
co najwyżej cztery przedziały siatki.
Dla każdego $n\ge j$ z
\eqref{eq:frostman-dyadic-bound-26} wynika
\[
 \mu_n(B(x,r))\le4^d(\sqrt d\,\ell_j)^s
 \le4^d(2\sqrt d)^s r^s.
\]
Dla większych $r$ analogiczne oszacowanie wynika z
$\mu_n(K)\le w(Q_0)$ po powiększeniu stałej.
Otrzymaliśmy dla wszystkich $n$ i $x,r$ wspólne
oszacowanie $\mu_n(B(x,r))\le C_{d,s,L}r^s$.

\emph{Krok 4: granica słaba.} Miary $\mu_n$ są
skoncentrowane na zwartym $K$ i mają jednostajnie
ograniczone masy. Twierdzenie o zwartości miar na
przestrzeni zwartej, dowiedzione w rozdziale
\ref{ch:weak-measures-22}, pozwala wybrać podciąg
zbieżny wąsko do skończonej miary $\mu$ na $K$.
Testując funkcję stałą $1$, otrzymujemy
$\mu(K)=\lim\mu_{n_k}(K)\ge m_0>0$.
Dla otwartej kuli $B(x,r)$ twierdzenie Portmanteau
daje
$\mu(B(x,r))\le\liminf_k\mu_{n_k}(B(x,r))
\le C_{d,s,L}r^s$. Jest to (ii).
Normalizując przez $\mu(K)>0$, otrzymujemy miarę
prawdopodobieństwa. Cała konstrukcja została
przeprowadzona wewnątrz zwartego $K$.
\end{proof}

\begin{uwaga}[Zakres twierdzenia]
Dowiedliśmy pełnej wersji dla \emph{zbiorów zwartych}.
Rozszerzenie na dowolne zbiory analityczne lub borelowskie
korzysta z dodatkowego twierdzenia o pojemnościach
(regularności od wewnątrz w sensie pojemności).
Nie wolno automatycznie zastępować w powyższym dowodzie
zwartości przez samą mierzalność.
\end{uwaga}

\section{Zbiór Cantora i przykłady wymiarów ułamkowych}
\begin{twierdzenie}[Wymiar klasycznego zbioru Cantora]
\label{thm:cantor-dimension-26}
Niech $C\subseteq[0,1]$ będzie klasycznym zbiorem
Cantora uzyskanym przez kolejne usuwanie otwartych
środkowych trzecich części przedziałów. Wtedy
\[
 \dim_HC=\frac{\log2}{\log3}.
\]
\end{twierdzenie}
\begin{proof}
Niech $s=\log2/\log3$, zatem $2\cdot3^{-s}=1$.
Po $n$ krokach zbiór $C$ zawiera się w sumie
$2^n$ domkniętych przedziałów długości $3^{-n}$.
Dla $t>s$ otrzymujemy
$\mathcal H^t_{3^{-n}}(C)\le
c_t2^n3^{-nt}=c_t(2\cdot3^{-t})^n\to0$,
stąd $\dim_HC\le s$.

Dla nierówności odwrotnej skonstruujmy miarę
$\mu$ przez przypisanie każdemu przedziałowi
poziomu $n$ masy $2^{-n}$, zgodnie z konstrukcją
Cantora z rozdziału~\ref{ch:lebesgue}.
Pokażemy, że $\mu(B(x,r))\le Cr^s$.
Dla $0<r<1/3$ wybierzmy $n$ tak, aby
$3^{-(n+1)}<r\le3^{-n}$. Przedziały poziomu $n$
mają długość $\ell=3^{-n}$ i odległości między
różnymi przedziałami są co najmniej $\ell$.
Przedział $B(x,r)$ ma długość $2r\le2\ell$,
więc przecina co najwyżej cztery przedziały poziomu
$n$. W konsekwencji
$\mu(B(x,r))\le4\cdot2^{-n}
=4\cdot3^{-ns}\le4\cdot3^sr^s$.
Dla $r\ge1/3$ powiększamy stałą.
Zasada rozkładu masy~\ref{thm:mass-distribution-26}
daje $\dim_HC\ge s$.
\end{proof}

\begin{center}
\begin{tikzpicture}[x=10cm,y=0.8cm]
\foreach \i/\y in {0/0,1/-1,2/-2,3/-3}{
 \ifnum\i=0
  \draw[line width=2pt,blue!65!black] (0,\y)--(1,\y);
 \fi
 \ifnum\i=1
  \draw[line width=2pt,blue!65!black] (0,\y)--(1/3,\y);
  \draw[line width=2pt,blue!65!black] (2/3,\y)--(1,\y);
 \fi
 \ifnum\i=2
  \foreach \a in {0,2/9,2/3,8/9}{\draw[line width=2pt,blue!65!black] (\a,\y)--({\a+1/9},\y);}
 \fi
 \ifnum\i=3
  \foreach \a in {0,2/27,2/9,8/27,2/3,20/27,8/9,26/27}{\draw[line width=2pt,blue!65!black] (\a,\y)--({\a+1/27},\y);}
 \fi
 \node[anchor=east,font=\scriptsize] at (-0.025,\y) {etap \i};
}
\end{tikzpicture}
\end{center}
\noindent\textit{Pierwsze etapy konstrukcji zbioru Cantora.}
Każde przejście zastępuje przedział dwoma
przedziałami o jednej trzeciej jego długości.

\begin{przyklad}[Gruby zbiór Cantora]
Istnieje domknięty, nigdzie gęsty zbiór
$F\subset[0,1]$ o $\lambda_1(F)>0$.
Na mocy twierdzenia~\ref{thm:euclidean-dimension-26}
ma wymiar Hausdorffa równy $1$, choć ma puste wnętrze.
Wymiar Hausdorffa nie charakteryzuje topologicznego
wymiaru ani posiadania wnętrza.
\end{przyklad}

\begin{wniosek}[Wykres funkcji Lipschitzowskiej]
\label{cor:graph-dimension-26}
Jeżeli $A\subseteq\mathbb R^m$ i
$f:A\to\mathbb R^k$ jest Lipschitzowska, to
$\dim_H\{(x,f(x)):x\in A\}=\dim_HA$.
\end{wniosek}
\begin{proof}
Odwzorowanie $x\mapsto(x,f(x))$ jest Lipschitzowskie,
więc nie zwiększa wymiaru Hausdorffa na mocy
twierdzenia~\ref{thm:hausdorff-lipschitz-26}.
Rzut $(x,y)\mapsto x$ jest również Lipschitzowski
i odwzorowuje wykres na $A$. Stosując tę samą
nierówność w drugą stronę, otrzymujemy równość.
\end{proof}

\section{Energie potencjałowe i wymiar}
\begin{definicja}[$t$-energia miary]
\label{def:energy-26}
Niech $t>0$, a $\mu$ będzie skończoną dodatnią miarą
borelowską na $\mathbb R^d$. Jej $t$-energią nazywamy
\[
 I_t(\mu)=\int_{\mathbb R^d}\!\int_{\mathbb R^d}
 |x-y|^{-t}\,\mathrm d\mu(y)\,\mathrm d\mu(x),
\]
przyjmując $|0|^{-t}=+\infty$. Całka należy do
$[0,+\infty]$; nie zakładamy jej skończoności.
\end{definicja}

\begin{twierdzenie}[Oszacowanie energii przez warunek Frostmana]
\label{thm:frostman-energy-26}
Jeżeli $\mu$ jest skończoną dodatnią miarą o
$\mu(B(x,r))\le Cr^s$ dla wszystkich $x,r$,
gdzie $s>0$, to $I_t(\mu)<\infty$ dla każdego
$0<t<s$.
\end{twierdzenie}
\begin{proof}
Z tożsamości warstwowej dla nieujemnych funkcji
$|x-y|^{-t}=t\int_{|x-y|}^{\infty}r^{-t-1}\,\mathrm dr$
i z twierdzenia Tonellego otrzymujemy
\[
 \int_{\mathbb R^d}|x-y|^{-t}\,\mathrm d\mu(y)
 =t\int_0^\infty r^{-t-1}\mu(B(x,r))\,\mathrm dr.
\]
W przedziale $(0,1)$ całka jest ograniczona przez
$Ct\int_0^1r^{s-t-1}\,\mathrm dr=Ct/(s-t)$.
W przedziale $(1,\infty)$ nie przekracza
$t\mu(\mathbb R^d)\int_1^\infty r^{-t-1}\,\mathrm dr
=\mu(\mathbb R^d)$. Oszacowanie jest jednostajne
względem $x$. Całkowanie po $x$ dowodzi tezy.
\end{proof}

\begin{uwaga}[Dalsza perspektywa]
W rozdziałach dotyczących geometrycznej teorii miary
pojawią się miary prostowalne, wzór na pole powierzchni
i wzór koobszaru. Miara Hausdorffa nadaje im naturalną
skalę. Związek energii z pojemnością stanowi natomiast
punkt wyjścia do teorii potencjału, której pełne
rozwinięcie pozostawiamy do osobnego opracowania.
\end{uwaga}

% END chapters/pl/measure_theory/26-rozdzial.tex


% BEGIN chapters/pl/measure_theory/27-rozdzial.tex
% ROZDZIAŁ 27 -- ujęcie miarowe; definicje geometryczne w osobnym skrypcie.
\chapter[Miary na rozmaitościach]{Miary na rozmaitościach: ujęcie miarowo-teoretyczne}
\label{ch:miary-rozmaitosci-27}

\section{Zakres rozdziału i ustalone obiekty}
W opracowaniu \emph{Geometria i topologia różniczkowa}, w rozdziale
\emph{Tensor metryczny, długość i miara}, skonstruowano już
miarę riemannowską, wyprowadzono jej niezależność od mapy,
wprowadzono metrykę indukowaną i obliczono całki po gładkich
podrozmaitościach. W rozdziale o formach udowodniono twierdzenie
Stokesa. Nie powtarzamy tych dowodów ani przykładów
sfery i paraboloidy. Zamiast tego badamy miarę riemannowską
jako obiekt teorii miary: regularność, zbiory zerowe,
gęstości Radona--Nikodyma, transport miar i produkty.

Przez $(M^n,g)$ rozumiemy w tym rozdziale niepustą, gładką,
Hausdorffowską i drugą przeliczalną rozmaitość bez brzegu,
wymiaru $n\ge1$, wyposażoną w gładką metrykę riemannowską.
Nie wymagamy orientowalności. Symbol $\mathcal B(M)$ oznacza
$\sigma$-algebrę borelowską, a $\mu_g$ oznacza miarę
borelowską określoną lokalnie przez
\begin{equation}\label{eq:g-density-27}
  \mu_g(A\cap U)=\int_{x(A\cap U)}
  \sqrt{\det G(x^{-1}(u))}\,\mathrm d\lambda_n(u),
\end{equation}
gdzie $(U,x)$ jest mapą, $A\in\mathcal B(M)$,
a $G=(g_{ij})$ jest macierzą metryki w tej mapie.
Jest to przypomnienie definicji, nie nowa konstrukcja.
Jej uzupełnienie będziemy oznaczać przez
$(\overline{\mathcal B(M)}^{\,\mu_g},\overline\mu_g)$.

\section{Regularność i $\sigma$-skończoność}
\begin{twierdzenie}[Miara objętościowa jest miarą Radona]
\label{thm:riemann-radon-27}
Niech $(M^n,g)$ spełnia powyższe założenia. Wówczas
$\mu_g$ jest lokalnie skończoną, $\sigma$-skończoną
miarą Radona na $\mathcal B(M)$. W szczególności
$\mu_g(K)<\infty$ dla każdego zwartego $K\subseteq M$.
\end{twierdzenie}
\begin{proof}
Weźmy punkt $p\in M$ i mapę $(U,x)$ wokół $p$.
Wybieramy kulę euklidesową $B$ tak, aby
$\overline B\subset x(U)$, a jej przeciwobraz $K=x^{-1}(\overline B)$
był zwartym otoczeniem punktu $p$. Gęstość
$\rho(u)=\sqrt{\det G(x^{-1}(u))}$ jest dodatnia
i ciągła, zatem na $\overline B$ jest ograniczona
z góry i z dołu przez dodatnie stałe. Z
\eqref{eq:g-density-27} wynika $\mu_g(K)<\infty$.
Dzięki zwartości dowolny zwarty $L\subset M$ pokrywamy
skończenie wieloma takimi otoczeniami, skąd
$\mu_g(L)<\infty$.

Rozmaitość $M$ jest lokalnie zwarta i ma przeliczalną
bazę, więc jest $\sigma$-zwarta: wybieramy
przeliczalnie wiele otoczeń o zwartych domknięciach,
które ją pokrywają. W konsekwencji $\mu_g$ jest
$\sigma$-skończona. Ponadto $M$ jest przestrzenią
metryzowalną i lokalnie zwartą. Stosujemy
udowodnione wcześniej twierdzenie o regularności
lokalnie skończonych miar borelowskich na takich
przestrzeniach (Rozdział~\ref{ch:radon-measures-20}),
aby otrzymać regularność wewnętrzną przez zbiory
zwarte i zewnętrzną przez zbiory otwarte.
\end{proof}

\begin{uwaga}[Miara borelowska a jej uzupełnienie]
Nie należy utożsamiać $\mathcal B(M)$ z jej
uzupełnieniem. Podzbiory borelowskiego zbioru
miary zero mogą nie być borelowskie, ale należą do
$\overline{\mathcal B(M)}^{\,\mu_g}$.
Definicja miary Radona dotyczy w tym rozdziale
miary borelowskiej; do całkowania funkcji
mierzalnych względem jej uzupełnienia używamy
standardowego rozszerzenia całki.
\end{uwaga}

\section{Zbiory o zerowej objętości}
\begin{twierdzenie}[Lokalne kryterium zbiorów zerowych]
\label{thm:null-local-27}
Niech $A\in\mathcal B(M)$, a
$\{(U_j,x_j)\}_{j\ge1}$ będzie przeliczalnym atlasem
pokrywającym $M$. Wówczas
\[
 \mu_g(A)=0\quad\Longleftrightarrow\quad
 \lambda_n\bigl(x_j(A\cap U_j)\bigr)=0
 \quad\text{dla każdego }j\ge1.
\]
To samo kryterium działa dla zbiorów mierzalnych
w uzupełnionej $\sigma$-algebrze, gdy ich obrazy
rozumiemy jako zbiory mierzalne w sensie Lebesgue'a.
\end{twierdzenie}
\begin{proof}
Na każdej mapie gęstość $\rho_j$ jest dodatnia
i ciągła. Dla zbioru borelowskiego $E\subseteq x_j(U_j)$
mamy $\int_E\rho_j\,\mathrm d\lambda_n=0$
wtedy i tylko wtedy, gdy $\lambda_n(E)=0$:
wystarcza wyczerpać $x_j(U_j)$ zbiorami zwartymi,
na których $\rho_j$ ma dodatnie minimum.
Z \eqref{eq:g-density-27} otrzymujemy kryterium
lokalne, a przeliczalność atlasu pozwala
zsumować je na całym $M$. W wersji zupełnej
zbiory zerowe można zastąpić borelowskimi
zbiorami zerowymi, które je zawierają.
\end{proof}

\begin{wniosek}[Obrazy immersji niższego wymiaru]
\label{cor:submanifold-null-27}
Niech $N^k$ będzie drugą przeliczalną gładką
rozmaitością, $0\le k<n$, a $F:N\to M^n$
immersją klasy $C^1$. Jeśli $S=F(N)$ jest
zbiorem borelowskim w $M$, to $\mu_g(S)=0$.
W szczególności teza zachodzi dla każdej osadzonej
podrozmaitości gładkiej wymiaru $k<n$.
\end{wniosek}
\begin{proof}
Twierdzenie o stałej randze pozwala dla każdego
$p\in N$ wybrać otoczenie $V_p\subseteq N$,
na którym $F$ jest osadzeniem, a $F(V_p)$
ma w odpowiedniej mapie celu postać kawałka
$\mathbb R^k\times\{0\}^{n-k}$.
Taki zbiór ma zerową $n$-wymiarową miarę
Lebesgue'a z twierdzenia Tonellego, więc
ma zerową $\mu_g$-miarę w każdej mapie celu.
Ponieważ $N$ ma przeliczalną bazę, z pokrycia
$(V_p)_{p\in N}$ wybieramy podpokrycie
przeliczalne $(V_j)$. Zatem
$S=\bigcup_jF(V_j)$, a przeliczalna
podaddytywność miary zewnętrznej daje
$\mu_g^*(S)=0$. Symbol $\mu_g^*$ oznacza tu miarę
zewnętrzną wyznaczoną przez $\mu_g$, czyli infimum
$\mu_g(B)$ po borelowskich $B\supseteq S$. Założenie borelowskości $S$
pozwala zapisać wniosek jako $\mu_g(S)=0$.
Argument nie zakłada, że \emph{cały} obraz immersji
jest lokalnie osadzoną podrozmaitością w otoczeniu
każdego punktu obrazu: przy samoprzecięciach lub
nagromadzeniach może to być nieprawdą.
\end{proof}

\begin{uwaga}[Dwie różne miary na tym samym zbiorze]
Wniosek nie oznacza, że długość krzywej lub pole
powierzchni znika. Dla osadzonej $k$-wymiarowej
podrozmaitości $S$, miara $\mu_{g|_S}$ może być
dodatnia, choć $\mu_g(S)=0$, gdy $k<n$.
To porównanie prowadzi do miar Hausdorffa;
identyfikację $\mu_{g|_S}$ z odpowiednio
unormowaną $\mathcal H^k$ uzasadnimy
w Rozdziale~\ref{ch:area-coarea-28}.
\end{uwaga}

\section{Miary z gęstością i rozkład Lebesgue'a}
\begin{twierdzenie}[Miary absolutnie ciągłe względem objętości]
\label{thm:rn-manifold-27}
Niech $(M,g)$ będzie jak wyżej i niech
$\nu$ będzie dodatnią, $\sigma$-skończoną
miarą na $\mathcal B(M)$, przy czym
$\nu\ll\mu_g$. Wtedy istnieje nieujemna funkcja
borelowska $h$, jednoznaczna $\mu_g$-prawie
wszędzie, taka że dla każdego $A\in\mathcal B(M)$
\[
 \nu(A)=\int_A h(p)\,\mathrm d\mu_g(p).
\]
Jeżeli dodatkowo $\nu$ jest lokalnie skończona,
to $h\in L^1_{\mathrm{loc}}(M,\mu_g)$,
tj. $\int_K h\,\mathrm d\mu_g<\infty$
dla każdego zwartego $K\subseteq M$.
\end{twierdzenie}
\begin{proof}
Miara $\mu_g$ jest $\sigma$-skończona na mocy
Twierdzenia~\ref{thm:riemann-radon-27}.
Stosujemy twierdzenie Radona--Nikodyma
(Rozdział~\ref{ch:radon-nikodym-18})
do pary $(\nu,\mu_g)$. Otrzymujemy $h$ mierzalną
względem uzupełnienia; można wybrać przedstawiciela
borelowskiego, modyfikując ją na zbiorze
$\mu_g$-miary zero. Jeżeli $\nu$ jest lokalnie
skończona, każdy zwarty $K$ można przykryć
skończenie wieloma otoczeniami o skończonej
$\nu$-mierze, więc $\nu(K)<\infty$.
Wstawienie $A=K$ daje lokalną całkowalność.
\end{proof}

\begin{przyklad}[Gęstość w mapie]
Jeżeli w mapie $x:U\to V\subseteq\mathbb R^n$
miara $\nu$ ma gęstość $q(u)$ względem $\lambda_n$,
to gęstość względem $\mu_g$ wynosi
$h(p)=q(x(p))/\sqrt{\det G(p)}$
na $U$, w sensie równości prawie wszędzie.
Wyrażenie $q$ zmienia się po zmianie mapy;
$h$ jest funkcją skalarną na rozmaitości.
\end{przyklad}

\begin{twierdzenie}[Rozkład na część objętościową i osobliwą]
\label{thm:leb-manifold-27}
Jeżeli $\nu$ jest dodatnią, $\sigma$-skończoną
miarą na $\mathcal B(M)$, to ma jednoznaczny rozkład
$\nu=h\mu_g+\nu_{\mathrm s}$, gdzie
$h\ge0$ jest borelowska, a
$\nu_{\mathrm s}\perp\mu_g$.
\end{twierdzenie}
\begin{proof}
Jest to zastosowanie twierdzenia o rozkładzie
Lebesgue'a (Rozdział~\ref{ch:ac-singular-17})
i twierdzenia Radona--Nikodyma do części
absolutnie ciągłej. Jednoznaczność obu składników
wynika z tych twierdzeń.
\end{proof}

\section{Przenoszenie objętości przez dyfeomorfizm}
\begin{definicja}[Jakobian objętościowy]
\label{def:volume-jacobian-27}
Niech $(M^n,g)$ i $(N^n,h)$ będą rozmaitościami
riemannowskimi, a $F:M\to N$ dyfeomorfizmem $C^1$.
Dla $p\in M$ wybieramy ortonormalne bazy
w $T_pM$ i $T_{F(p)}N$, a przez $A_p$ oznaczamy
macierz $\mathrm dF_p$ w tych bazach. Definiujemy
$J_{g,h}F(p)=|\det A_p|$. Wartość nie zależy
od wyboru baz. W lokalnych współrzędnych $u$ i $v$
wyraża się wzorem
\[
 J_{g,h}F(p)=|\det D(v\circ F\circ u^{-1})(u(p))|
 \frac{\sqrt{\det H(v(F(p)))}}{\sqrt{\det G(u(p))}}.
\]
\end{definicja}

\begin{twierdzenie}[Zamiana zmiennych na rozmaitościach]
\label{thm:manifold-change-27}
Dla powyższego dyfeomorfizmu i każdej funkcji
borelowskiej $f:N\to[0,+\infty]$ zachodzi
\begin{equation}\label{eq:manifold-change-27}
 \int_N f(q)\,\mathrm d\mu_h(q)=
 \int_M f(F(p))J_{g,h}F(p)\,\mathrm d\mu_g(p).
\end{equation}
Ta sama równość zachodzi dla funkcji ze znakiem,
gdy $f\in L^1(N,\mu_h)$.
W szczególności $(F_*\mu_g)\ll\mu_h$ i
\[
 \frac{\mathrm d(F_*\mu_g)}{\mathrm d\mu_h}(q)
 =\frac{1}{J_{g,h}F(F^{-1}(q))}
 \quad\text{dla }\mu_h\text{-prawie wszystkich }q.
\]
\end{twierdzenie}
\begin{proof}
Najpierw załóżmy, że $f$ ma nośnik zawarty
w mapie celu, której przeciwobraz przez $F$
mieści się w mapie źródła. Z definicji
$\mu_g,\mu_h$ oraz twierdzenia o zamianie
zmiennych w $\mathbb R^n$ otrzymujemy
\eqref{eq:manifold-change-27}, po skróceniu
wyznaczników macierzy metryk.

W ogólności wybieramy przeliczalne pokrycie
mapami $V_j$ przestrzeni $N$ i rozłączny podział
borelowski $B_1=V_1$, $B_j=V_j\setminus
\bigcup_{i<j}V_i$. Każdy $B_j$ rozbijamy dalej
na przeliczalną rodzinę części, na których
$F^{-1}$ ma wspólną mapę źródła. Dla
$f\chi_{B_j}$ stosujemy argument lokalny,
a następnie sumujemy po rozłącznych częściach,
używając twierdzenia Tonellego.

Aby obliczyć miarę obrazową, przyjmujemy
$f=\chi_A$ w równaniu zastosowanym do $F^{-1}$.
Ponieważ
$J_{h,g}(F^{-1})(q)=[J_{g,h}F(F^{-1}q)]^{-1}$,
dostajemy podaną pochodną Radona--Nikodyma.
Dla funkcji całkowalnych stosujemy równość
oddzielnie do części dodatniej i ujemnej.
\end{proof}

\begin{przyklad}[Konforemna zmiana metryki]
Niech $\varphi\in C^\infty(M;\mathbb R)$ i
$g'=e^{2\varphi}g$. Macierze spełniają
$G'=e^{2\varphi}G$, więc
$\det G'=e^{2n\varphi}\det G$. W konsekwencji
\[
 \mathrm d\mu_{g'}(p)=e^{n\varphi(p)}\,\mathrm d\mu_g(p).
\]
Nie oznacza to, że krzywizna obu metryk jest
taka sama: wzór dotyczy tylko miary objętościowej.
\end{przyklad}

\section{Miara na produkcie rozmaitości}
\begin{twierdzenie}[Miara produktu riemannowskiego]
\label{thm:product-riemann-27}
Niech $(M^n,g)$ i $(N^m,h)$ będą rozmaitościami
jak wyżej. Na $M\times N$ rozważmy metrykę
$g\oplus h$. Wtedy na $\mathcal B(M\times N)$
zachodzi $\mu_{g\oplus h}=\mu_g\otimes\mu_h$.
Dla nieujemnej funkcji borelowskiej $f$ mamy
\[
 \int_{M\times N} f(p,q)\,\mathrm d\mu_{g\oplus h}(p,q)
 =\int_M\!\left(\int_N f(p,q)\,\mathrm d\mu_h(q)\right)
   \mathrm d\mu_g(p).
\]
\end{twierdzenie}
\begin{proof}
W mapach produktowych macierz metryki jest
blokowo diagonalna $\operatorname{diag}(G,H)$.
Stąd $\sqrt{\det(G\oplus H)}=
\sqrt{\det G}\sqrt{\det H}$. Miara lokalna
jest produktem odpowiednich miar lokalnych
na mapach. Równość na całym produkcie wynika
z przeliczalnego rozbicia na mapy i jednoznaczności
miary produktowej dla miar $\sigma$-skończonych.
Ostatnia równość jest twierdzeniem Tonellego.
\end{proof}

\section{Od gładkich podrozmaitości do zbiorów prostowalnych}
W geometrii rozważaliśmy miarę indukowaną
na \emph{gładkich} podrozmaitościach. Obraz
odwzorowania Lipschitzowskiego nie musi być
gładką podrozmaitością, lecz nadal może
posiadać sensowną miarę $k$-wymiarową.
W następnym rozdziale pokażemy, jak
$\mathcal H^k$ mierzy takie obrazy i dlaczego
przy odwzorowaniach nieiniektywnych konieczny
jest czynnik krotności. Wprowadzone tam wzory
są ściśle miarowe: nie potrzebują globalnej
orientacji ani globalnych parametryzacji.

% END chapters/pl/measure_theory/27-rozdzial.tex


% BEGIN chapters/pl/measure_theory/28-rozdzial.tex
% ROZDZIAŁ 28 -- wzory miarowo-geometryczne, nie powtórzenie całek po podrozmaitościach.
\chapter{Wzór na pole powierzchni, krotności i koobszar}
\label{ch:area-coarea-28}

\section{Od wzoru parametrycznego do wzoru miarowego}
W skrypcie \emph{Geometria i topologia różniczkowa}
udowodniono wzór na całkę po kawałku gładkiej
podrozmaitości opisanej parametryzacją, która
pokrywa ten kawałek dokładnie raz. Tutaj rozważamy
\emph{odwzorowania nieiniektywne} oraz ich włókna.
Nowością są krotności obrazów, całkowanie po
przeciwobrazach punktów i formuły dla
odwzorowań Lipschitzowskich.

W tym rozdziale $\lambda_k$ oznacza miarę Lebesgue'a
na $\mathbb R^k$, a $\mathcal H^k$ -- miarę
Hausdorffa z Rozdziału~\ref{ch:hausdorff-26}.
Stosujemy konwencję z
Uwagi~\ref{rem:hausdorff-convention-26}: dla wymiaru
całkowitego $k\ge1$ wybieramy $c_k=\kappa_k^{-1}$,
tak aby $\mathcal H^k|_{\mathbb R^k}=\lambda_k$;
$\mathcal H^0$ jest miarą liczącą.
Dzięki temu poniższe wzory nie zakładają
nieudowodnionej nierówności izodiametrycznej.
Przy innej stałej normalizującej ich prawe strony
należy odpowiednio przeskalować.

\begin{definicja}[Jakobian wymiaru $k$]
\label{def:k-jacobian-28}
Niech $1\le k\le N$. Dla dowolnego odwzorowania
liniowego $A:\mathbb R^k\to\mathbb R^N$ definiujemy
\[
 J_k(A)=\sqrt{\det(A^{\mathsf T}A)}.
\]
Jeżeli $U\subseteq\mathbb R^k$ jest otwarty,
a $F:U\to\mathbb R^N$ jest różniczkowalne
w punkcie $x$, to kładziemy $J_kF(x)=J_k(DF(x))$.
W szczególności dotyczy to wszystkich punktów odwzorowania
$C^1$ oraz prawie wszystkich punktów odwzorowania
lokalnie Lipschitzowskiego. Jakobian jest zerowy,
gdy $\operatorname{rank}DF(x)<k$; dla $k=N$
jest równy $|\det DF(x)|$.
\end{definicja}

\begin{definicja}[Krotność odwzorowania]
\label{def:multiplicity-28}
Dla dowolnego odwzorowania $F:X\to Y$
i zbioru $E\subseteq X$ definiujemy
\[
 N(F,E,y)=\#\{x\in E:F(x)=y\}\in\mathbb N_0\cup\{+\infty\}.
\]
Symbol $\#$ oznacza moc zbioru, a nie operację
podnoszenia indeksów za pomocą metryki.
Pojawia się tu również wartość $+\infty$,
gdy włókno ma nieskończenie wiele punktów.
\end{definicja}

\section{Miara Hausdorffa gładkiego wykresu}
\begin{lemat}[Małe zniekształcenie wykresu]
\label{lem:graph-hausdorff-28}
Niech $V\subseteq\mathbb R^k$ będzie otwarty,
a $h:V\to\mathbb R^{N-k}$ funkcją $C^1$.
Niech $\Gamma(z)=(z,h(z))$ oraz
$\nu(A)=\mathcal H^k(\Gamma(A))$
dla borelowskich $A\subseteq V$. Wtedy $\nu$
jest miarą Radona na $V$ i dla każdego
$z_0\in V$ zachodzi
\[
 \lim_{r\downarrow0}
 \frac{\nu(B(z_0,r))}{\lambda_k(B(z_0,r))}
 =\sqrt{\det\bigl(I_k+(Dh(z_0))^{\mathsf T}Dh(z_0)\bigr)}.
\]
\end{lemat}
\begin{proof}
Na każdej kuli domkniętej $K\subset V$ wykres $\Gamma$
jest Lipschitzowski, ponieważ $Dh$ jest ograniczona
na $K$. Zatem, na mocy
Twierdzenia~\ref{thm:hausdorff-lipschitz-26},
$\mathcal H^k(\Gamma(K))<\infty$, gdyż
$\mathcal H^k(K)=\lambda_k(K)$ jest skończona.
Z drugiej strony rzut ortogonalny $\pi$ na
pierwszych $k$ współrzędnych jest odwrotnością
$\Gamma$ na obrazie i jest $1$-Lipschitzowski.
Ponieważ $\Gamma$ jest homeomorfizmem $V$ na
$\Gamma(V)$, przeciwobrazy zbiorów borelowskich
są borelowskie, a $\nu$ jest miarą borelowską
lokalnie skończoną. Stosujemy twierdzenie o
regularności miar lokalnie skończonych na lokalnie
zwartych ośrodkowych przestrzeniach metrycznych
z Rozdziału~\ref{ch:radon-measures-20}.
Otrzymujemy, że $\nu$ jest miarą Radona,
bez zakładania z góry regularności miary
Hausdorffa na całym wykresie.

Ustalmy $z_0\in V$ i liniowe odwzorowanie
$L(z)=(z,h(z_0)+Dh(z_0)(z-z_0))$.
Na małej kuli $B(z_0,r)$ z ciągłości $Dh$ wynika
$|[h(u)-h(v)]-Dh(z_0)(u-v)|\le\varepsilon_r|u-v|$,
gdzie $\varepsilon_r\to0$.
Ponieważ $|L(u)-L(v)|\ge|u-v|$,
odwzorowanie $\Gamma\circ L^{-1}$ między
odpowiednimi fragmentami wykresów jest
dwulipschitzowskie ze stałymi należącymi
do przedziału $[1-\varepsilon_r,1+\varepsilon_r]$
po ewentualnym zwiększeniu $\varepsilon_r$.
Z oszacowania miary Hausdorffa przez stałą
Lipschitza otrzymujemy
\[
 (1-\varepsilon_r)^k\mathcal H^k(L(B(z_0,r)))
 \le\nu(B(z_0,r))\le
 (1+\varepsilon_r)^k\mathcal H^k(L(B(z_0,r))).
\]
Liniowy obraz przestrzeni $\mathbb R^k$ jest
$k$-płaszczyzną euklidesową. Diagonalizacja
macierzy Grama pokazuje, że
$\mathcal H^k(L(A))=J_kL\,\lambda_k(A)$
dla każdego borelowskiego $A$, gdzie
$J_kL=\sqrt{\det(I_k+Dh(z_0)^{\mathsf T}Dh(z_0))}$.
Dzielimy przez $\lambda_k(B(z_0,r))$ i
przechodzimy do granicy.
\end{proof}

\begin{twierdzenie}[Miara na gładkim wykresie]
\label{thm:graph-area-28}
Przy założeniach Lematu~\ref{lem:graph-hausdorff-28}
dla każdej nieujemnej funkcji borelowskiej
$\psi$ na $\Gamma(V)$ mamy
\[
 \int_{\Gamma(V)}\psi(y)\,\mathrm d\mathcal H^k(y)
 =\int_V\psi(\Gamma(z))
 \sqrt{\det(I_k+Dh(z)^{\mathsf T}Dh(z))}
 \,\mathrm d\lambda_k(z).
\]
To samo zachodzi dla funkcji całkowalnych.
\end{twierdzenie}
\begin{proof}
Miara $\nu$ jest lokalnie skończona i
$\nu\ll\lambda_k$: obraz zbioru o zerowej
mierze Lebesgue'a ma zerową miarę Hausdorffa,
co wynika z lokalnej Lipschitzowskości $\Gamma$.
Istnieje zatem gęstość $a=\mathrm d\nu/\mathrm d\lambda_k$.
Twierdzenie Lebesgue'a o różniczkowaniu miar
(Rozdział~\ref{ch:differentiation-25})
identyfikuje $a(z)$ prawie wszędzie z granicą
ilorazów mas małych kul. Poprzedni lemat
pokazuje, że granica ta jest równa $J_k\Gamma(z)$.
Otrzymujemy więc równość miar
$\nu=J_k\Gamma\,\lambda_k$, a następnie
równość całek.
\end{proof}

\section{Wzór na pole powierzchni dla immersji}
\begin{twierdzenie}[Wzór z krotnością dla immersji $C^1$]
\label{thm:immersion-area-28}
Niech $1\le k\le N$, $U\subseteq\mathbb R^k$
będzie otwarty, a $F:U\to\mathbb R^N$
immersją klasy $C^1$, tj.
$\operatorname{rank}DF(x)=k$ dla każdego $x\in U$.
Niech $E\subseteq U$ będzie borelowski,
a $f:E\to[0,+\infty]$ borelowska.
Wówczas funkcja
$y\mapsto\sum_{x\in E\cap F^{-1}(\{y\})}f(x)$
jest $\mathcal H^k$-mierzalna i
\begin{equation}\label{eq:area-immersion-28}
 \int_E f(x)J_kF(x)\,\mathrm d\lambda_k(x)
 =\int_{\mathbb R^N}\left(
 \sum_{x\in E\cap F^{-1}(\{y\})}f(x)\right)
 \,\mathrm d\mathcal H^k(y).
\end{equation}
Przy $f\equiv1$ na $E$ (czyli $f=\chi_E$ na $U$)
prawa strona redukuje się do
\[
 \int_{\mathbb R^N} N(F,E,y)\,\mathrm d\mathcal H^k(y).
\]
\end{twierdzenie}
\begin{proof}
Twierdzenie o stałej randze zapewnia, że każdy
$x\in U$ ma otoczenie $U_x$, na którym $F$
jest osadzeniem na gładką podrozmaitość.
Przeliczalność bazy pozwala wybrać przeliczalne
pokrycie takimi otoczeniami $U_j$. Tworzymy
rozłączny podział borelowski
$B_1=U_1$, $B_j=U_j\setminus\bigcup_{i<j}U_i$.

Na każdym $U_j$ obraz $F(U_j)$ jest gładką
podrozmaitością w pewnym otoczeniu, a po
lokalnym przekształceniu ortogonalnym i zmianie
współrzędnych można go przedstawić jako wykres.
Twierdzenie~\ref{thm:graph-area-28}, zastosowane
w tych współrzędnych, oraz klasyczny wzór
podstawienia na dziedzinie $\mathbb R^k$
dają dla borelowskich $A\subseteq B_j$
\[
 \int_A f(x)J_kF(x)\,\mathrm d\lambda_k(x)
 =\int_{F(A)}f((F|_{U_j})^{-1}(y))
 \,\mathrm d\mathcal H^k(y).
\]
Różniczka zmiany parametrów wprowadza dokładnie
ten czynnik, który przekształca macierz Grama
$F$ w macierz Grama wykresu; jest to tożsamość
$J_k(F\circ\phi)=J_kF(\phi)|\det D\phi|$
dla lokalnego dyfeomorfizmu $\phi$.

Każdy $F(B_j\cap E)$ jest zbiorem borelowskim
w lokalnie osadzonym wykresie, a zatem również
w $\mathbb R^N$. Odpowiednie funkcje na obrazie
są borelowskie. Sumujemy równości po $j$;
po lewej używamy rozłączności $B_j$, a po prawej
twierdzenia Tonellego. Dla każdego $y$ suma
wkładów z kolejnych $B_j$ jest dokładnie sumą
po całym włóknie $E\cap F^{-1}(\{y\})$.
\end{proof}

\begin{przyklad}[Dlaczego trzeba liczyć krotności?]
Niech $F:(0,4\pi)\to\mathbb R^2$,
$F(t)=(\cos t,\sin t)$. Wówczas $J_1F=1$,
więc lewa strona \eqref{eq:area-immersion-28}
dla $f\equiv1$ wynosi $4\pi$.
Prawie każdy punkt okręgu ma dwa przeciwobrazy
w $(0,4\pi)$, zatem prawa strona wynosi
$2\mathcal H^1(S^1)=4\pi$.
Bez krotności otrzymalibyśmy wynik dwa razy za mały.
\end{przyklad}

\section{Jakobian normalny i całkowanie po włóknach}
\begin{definicja}[Jakobian normalny]
\label{def:normal-jacobian-28}
Niech $n\ge m\ge1$. Dla dowolnego odwzorowania
liniowego $A:\mathbb R^n\to\mathbb R^m$ definiujemy
\[
 \operatorname{Jac}_m(A)=\sqrt{\det(AA^{\mathsf T})}.
\]
Dla odwzorowania $F:U\to\mathbb R^m$ różniczkowalnego
w punkcie $x$ przyjmujemy
$\operatorname{Jac}_mF(x)=\operatorname{Jac}_m(DF(x))$.
Definicja obejmuje zatem odwzorowania $C^1$ i punkty
różniczkowalności odwzorowań Lipschitzowskich.
Jeżeli $\operatorname{rank}DF(x)<m$, to
$\operatorname{Jac}_mF(x)=0$; dla pełnej rangi
jest dodatni. Gdy $m=1$, jest równy $|\nabla F(x)|$.
\end{definicja}

\begin{lemat}[Tożsamość liniowo-algebraiczna koobszaru]
\label{lem:linear-coarea-28}
Niech $A:\mathbb R^n\to\mathbb R^m$ będzie
surjektywnym odwzorowaniem liniowym, a
$G:\mathbb R^m\times\mathbb R^{n-m}\to\mathbb R^n$
odwracalnym odwzorowaniem liniowym, dla którego
$AG(y,z)=y$. Wówczas
\[
 \sqrt{\det(AA^{\mathsf T})}\,|\det DG|
 =\sqrt{\det((D_zG)^{\mathsf T}D_zG)}.
\]
\end{lemat}
\begin{proof}
Wybieramy ortonormalne rozkłady
$\mathbb R^n=(\ker A)^\perp\oplus\ker A$
i zapisujemy macierze względem tych baz.
Ograniczenie $A_0=A|_{(\ker A)^\perp}$
jest izomorfizmem, a macierz $G$ ma blokową postać
$\begin{psmallmatrix} A_0^{-1}&0\\B&C\end{psmallmatrix}$,
gdzie $C$ jest odwracalną macierzą kwadratową
rozmiaru $n-m$. Stąd
$|\det DG|=|\det A_0|^{-1}|\det C|$.
Ponieważ $AA^{\mathsf T}=A_0A_0^{\mathsf T}$,
lewy pierwiastek jest równy $|\det A_0|$.
Kolumny $D_zG$ leżą w $\ker A$, a ich macierz
w wybranej ortonormalnej bazie jest $C$,
więc prawy pierwiastek wynosi $|\det C|$.
Dla $n=m$ przyjmujemy wyznacznik pustej
macierzy równy $1$ i dowód pozostaje poprawny.
\end{proof}

\begin{twierdzenie}[Wzór koobszaru dla submersji]
\label{thm:coarea-submersion-28}
Niech $n\ge m\ge1$, $U\subseteq\mathbb R^n$
będzie otwarty, a $F:U\to\mathbb R^m$
submersją klasy $C^1$, tj.
$\operatorname{rank}DF(x)=m$ dla każdego $x\in U$.
Dla każdej borelowskiej funkcji $g:U\to[0,+\infty]$
zachodzi
\begin{equation}\label{eq:coarea-smooth-28}
 \int_U g(x)\operatorname{Jac}_mF(x)\,\mathrm d\lambda_n(x)
 =\int_{\mathbb R^m}\left(
 \int_{F^{-1}(\{y\})}g(x)\,\mathrm d\mathcal H^{n-m}(x)
 \right)\mathrm d\lambda_m(y).
\end{equation}
Wewnętrzna całka dotyczy miary Hausdorffa
na włóknie, które jest gładką
$(n-m)$-wymiarową podrozmaitością (a dla
$n=m$ jest zliczaniem punktów). Jeżeli jedna ze stron
jest skończona dla $|g|$, równość zachodzi także
dla rzeczywistych i zespolonych funkcji całkowalnych.
\end{twierdzenie}
\begin{proof}
Gdy $n=m$, odwzorowanie $F$ jest lokalnym
dyfeomorfizmem. Dziedzinę dzielimy na przeliczalnie
wiele części borelowskich, na których $F$ jest
różnowartościowe. Zwykła zamiana zmiennych i
Tonelli dają tezę, ponieważ
$\mathcal H^0$ jest miarą liczącą. W dalszym
dowodzie przyjmujemy zatem $n>m$.

\emph{Krok 1: lokalne współrzędne włóknowe.}
Na mocy twierdzenia o funkcji uwikłanej, dla
każdego $x_0\in U$ istnieje otoczenie $V$ i
dyfeomorfizm $G:W\to V$, gdzie
$W\subseteq\mathbb R^m\times\mathbb R^{n-m}$,
taki że $F(G(y,z))=y$ dla $(y,z)\in W$.
Wybieramy $g$ o nośniku zawartym w takim $V$.

\emph{Krok 2: zamiana zmiennych i Fubini.}
Przez klasyczną zamianę zmiennych dostajemy
\[
 \int_V g(x)\operatorname{Jac}_mF(x)\,\mathrm d\lambda_n(x)
 =\int_W g(G(y,z))\operatorname{Jac}_mF(G(y,z))
 |\det DG(y,z)|\,\mathrm d\lambda_n(y,z).
\]
Ponieważ $DF(G)D_{(y,z)}G=(I_m\ \ 0)$,
Lemat~\ref{lem:linear-coarea-28} zastosowany
punktowo do różniczek daje
$\operatorname{Jac}_mF(G)|\det DG|=
J_{n-m}(z\mapsto G(y,z))$.
Przez Tonellego całka jest więc równa
\[
 \int_{\mathbb R^m}\left(\int_{W_y}
 g(G(y,z))J_{n-m}(z\mapsto G(y,z))
 \,\mathrm d\lambda_{n-m}(z)\right)\mathrm d\lambda_m(y),
\]
gdzie $W_y=\{z:(y,z)\in W\}$.
Odwzorowanie $z\mapsto G(y,z)$ jest gładkim
osadzeniem w $F^{-1}(\{y\})$. Stosujemy
Twierdzenie~\ref{thm:immersion-area-28} na włóknie
(bez krotności) i otrzymujemy prawą stronę
\eqref{eq:coarea-smooth-28} dla $g$ wspieranej w $V$.

\emph{Krok 3: sklejenie.}
Wybieramy przeliczalne pokrycie $U$ otoczeniami
$V_j$ jak w kroku 1 i rozłączny podział borelowski
$B_1=V_1$, $B_j=V_j\setminus\bigcup_{i<j}V_i$.
Zastępujemy $g$ przez $g\chi_{B_j}$, sumujemy
równości lokalne po $j$ i stosujemy Tonellego
po obu stronach. Dla $n=m$ wewnętrzne całkowanie
względem $\mathcal H^0$ oznacza sumowanie po
włóknach; argument z lokalnymi odwrotnościami
pozostaje poprawny.
\end{proof}

\begin{wniosek}[Różniczkowanie poziomic]
\label{cor:coarea-level-28}
Dla $U\subseteq\mathbb R^n$ otwartego i
funkcji $F\in C^1(U;\mathbb R)$
bez punktów krytycznych oraz dowolnej
nieujemnej funkcji borelowskiej $g$ zachodzi
\[
 \int_U g(x)|\nabla F(x)|\,\mathrm d\lambda_n(x)
 =\int_{\mathbb R}\left(\int_{F^{-1}(\{t\})}
 g(x)\,\mathrm d\mathcal H^{n-1}(x)\right)
 \,\mathrm d\lambda_1(t).
\]
\end{wniosek}

\section{Zastosowanie: rozkład miary obrazowej}
\begin{twierdzenie}[Gęstość obrazu przez submersję]
\label{thm:pushforward-coarea-28}
Niech $F:U\to\mathbb R^m$ spełnia założenia
Twierdzenia~\ref{thm:coarea-submersion-28},
a $w:U\to[0,+\infty]$ będzie borelowska.
Niech $\nu=w\lambda_n$ i $\eta=F_*\nu$.
Wówczas $\eta\ll\lambda_m$, a jej gęstość jest
równa
\[
 q(y)=\int_{F^{-1}(\{y\})}
 \frac{w(x)}{\operatorname{Jac}_mF(x)}
 \,\mathrm d\mathcal H^{n-m}(x)
\]
dla $\lambda_m$-prawie wszystkich $y$,
z możliwością przyjmowania wartości $+\infty$.
\end{twierdzenie}
\begin{proof}
Dla borelowskiego $A\subseteq\mathbb R^m$
stosujemy wzór koobszaru do funkcji
\[
 g(x)=\frac{\chi_{F^{-1}(A)}(x)w(x)}
 {\operatorname{Jac}_mF(x)}.
\]
Mianownik jest dodatni na całym $U$, więc
jest to funkcja borelowska nieujemna.
Wynik ma postać
$\eta(A)=\int_A q(y)\,\mathrm d\lambda_m(y)$.
To zarazem dowodzi absolutnej ciągłości
i identyfikuje pochodną Radona--Nikodyma.
\end{proof}

\begin{przyklad}[Poziomice promienia w pierścieniu]
Niech $U=\{x\in\mathbb R^2:a<|x|<b\}$,
gdzie $0<a<b$, i $F(x)=|x|^2$.
Wtedy $|\nabla F(x)|=2|x|$. Coarea daje
\[
 \int_U 2|x|\,\mathrm d\lambda_2(x)
 =\int_{a^2}^{b^2}\mathcal H^1(F^{-1}(\{t\}))
 \,\mathrm d\lambda_1(t)
 =\int_{a^2}^{b^2}2\pi\sqrt t\,\mathrm d\lambda_1(t)
 =\frac{4\pi}{3}(b^3-a^3).
\]
Jest to ilustracja całkowania po \emph{włóknach},
a nie kolejny sposób obliczania powierzchni sfery.
\end{przyklad}


\section{Twierdzenie Rademachera: skąd bierze się różniczka prawie wszędzie?}
\label{sec:rademacher-28}

Wzory geometryczne wymagają jakobianu, czyli różniczki. W przypadku
odwzorowań klasy $C^1$ istnieje ona wszędzie, natomiast odwzorowanie
Lipschitzowskie może nie być różniczkowalne w pewnych punktach.
Twierdzenie Rademachera dowodzi, że taki zbiór wyjątkowy ma zerową
miarę Lebesgue'a. Samo istnienie różniczki prawie wszędzie nie dowodzi
jednak ogólnych wzorów area i coarea: do ich uzasadnienia potrzebne są
jeszcze twierdzenia o pokrywaniu i mierzeniu obrazów zbiorów krytycznych.

\begin{definicja}[Różniczkowalność Frécheta]
\label{def:frechet-rademacher-28}
Niech $U\subseteq\mathbb R^n$ będzie otwarty,
$f:U\to\mathbb R^m$ i $x\in U$. Mówimy, że $f$ jest
\emph{różniczkowalna w sensie Frécheta} w $x$, jeżeli istnieje
liniowe $A:\mathbb R^n\to\mathbb R^m$ takie, że
\[
 \lim_{h\to0}\frac{|f(x+h)-f(x)-Ah|}{|h|}=0.
\]
Takie $A$ jest jednoznaczne i oznaczamy je przez $Df(x)$.
Jeżeli dla każdego $x\in K\Subset U$ istnieją $L_K$ oraz
otoczenie $V_K$ zbioru $K$, na którym
$|f(x)-f(y)|\le L_K|x-y|$, nazywamy $f$
\emph{lokalnie Lipschitzowskim}. Równoważnie każdy punkt
$U$ ma otoczenie, na którym obowiązuje taka nierówność
z pewną stałą.
\end{definicja}

\begin{lemat}[Jednowymiarowa różniczkowalność]
\label{lem:rademacher-line-28}
Niech $a<b$, a $f:[a,b]\to\mathbb R$ będzie funkcją
Lipschitzowską ze stałą $L$. Wtedy istnieje
$g\in L^\infty((a,b),\lambda_1)$ takie, że
$|g|\le L$ prawie wszędzie oraz
\[
 f(t)-f(a)=\int_{(a,t)}g(s)\,\mathrm d\lambda_1(s)
 \quad (t\in[a,b]).
\]
W konsekwencji $f'=g$ prawie wszędzie.
\end{lemat}
\begin{proof}
Ponieważ $f$ jest Lipschitzowska, jest absolutnie ciągła.
Dla przejrzystości odzyskamy gęstość bez powoływania się na
odwrotne podstawowe twierdzenie rachunku całkowego.
Połóżmy $p(t)=f(t)+Lt$. Dla $s<t$ mamy
$p(t)-p(s)\ge -L(t-s)+L(t-s)=0$, więc $p$ jest
niemalejąca i ciągła. Jej przyrosty na półotwartych
przedziałach określają dodatnią miarę borelowską
$\mu_p((s,t])=p(t)-p(s)$; istnienie tej miary wynika
z konstrukcji Lebesgue'a--Stieltjesa przez rozszerzenie
Carathéodory'ego (ciągłość $p$ zapewnia ciągłość
premiary z góry na przedziałach). Dla każdego przedziału
$[s,t]$ zachodzi
$0\le\mu_p((s,t])\le2L(t-s)$.
Twierdzenie o jednoznaczności miar na przedziałach daje
$\mu_p\le2L\lambda_1$. Twierdzenie Radona--Nikodyma
dostarcza więc $q\in L^\infty$ o wartościach w $[0,2L]$
takiego, że $\mu_p=q\lambda_1$. Zatem
$f(t)-f(a)=\int_{(a,t)}(q(s)-L)\,\mathrm d\lambda_1(s)$.
Dla $g=q-L$ zachodzi $|g|\le L$ prawie wszędzie;
Twierdzenie~\ref{thm:FTC-lebesgue-25} kończy dowód.
\end{proof}

\begin{lemat}[Słabe i kierunkowe pochodne funkcji Lipschitzowskiej]
\label{lem:rademacher-direction-28}
Niech $f:U\to\mathbb R$ będzie lokalnie Lipschitzowska.
Istnieją mierzalne funkcje $g_i\in L^\infty_{\mathrm{loc}}(U)$,
$i=1,\ldots,n$, takie, że dla $\varphi\in C_c^1(U)$
\[
 \int_U f(x)\partial_i\varphi(x)\,\mathrm d\lambda_n(x)
 =-\int_U g_i(x)\varphi(x)\,\mathrm d\lambda_n(x).
\]
Dla każdego ustalonego $v\in\mathbb R^n$ pochodna
kierunkowa $\partial_v f(x)=\lim_{t\to0}
(f(x+tv)-f(x))/t$ istnieje dla prawie wszystkich
$x\in U$ i jest tam równa $\sum_i v_i g_i(x)$.
\end{lemat}
\begin{proof}
Przypuśćmy najpierw, że $f$ jest Lipschitzowska na
prostopadłościanie $Q\Subset U$. Z Lematu
\ref{lem:rademacher-line-28} i Fubiniego otrzymujemy
istnienie $\partial_i f$ prawie wszędzie na $Q$ oraz
$|\partial_i f|\le L$. Jest to funkcja mierzalna,
ponieważ jest granicą mierzalnych ilorazów różnicowych
w punktach istnienia granicy. Niech $\varphi\in C_c^1(Q)$.
Na prawie każdej prostej równoległej do osi $i$ funkcja
$f$ jest absolutnie ciągła. Jednowymiarowe całkowanie
przez części i Fubini dają
$\int_Q f\partial_i\varphi\,\mathrm d\lambda_n
=-\int_Q (\partial_i f)\varphi\,\mathrm d\lambda_n$.
Wyczerpując $U$ przeliczalnie wieloma takimi
prostopadłościanami, otrzymujemy globalną funkcję
$g_i$ z lokalnym ograniczeniem.

Dla dowolnego ustalonego $v\ne0$ rozkładamy
$\mathbb R^n=v^\perp\oplus\mathbb Rv$.
Po odpowiedniej ortogonalnej zmianie współrzędnych
stosujemy jednowymiarowy lemat do przekrojów
w kierunku $v$ i twierdzenie Fubiniego.
Wynika stąd, że $\partial_v f$ istnieje prawie
wszędzie oraz że jej działanie dystrybucyjne jest
równe $\partial_v f=\sum_i v_i g_i$.
Jednoznaczność funkcji lokalnie całkowalnej
reprezentującej dystrybucję dowodzi ostatniej równości.
\end{proof}

\begin{twierdzenie}[Rademacher]
\label{thm:rademacher-28}
Niech $n,m\ge1$, $U\subseteq\mathbb R^n$ będzie
otwarty, a $f:U\to\mathbb R^m$ lokalnie Lipschitzowska.
Wtedy istnieje zbiór borelowski $N\subseteq U$ o
$\lambda_n(N)=0$ taki, że $f$ jest różniczkowalna
w sensie Frécheta w każdym punkcie $U\setminus N$.
Ponadto w każdym punkcie różniczkowalności
$\|Df(x)\|\le\operatorname{Lip}(f;V)$ dla
każdego otoczenia $V$ punktu $x$, na którym $f$
jest Lipschitzowska.
\end{twierdzenie}
\begin{proof}
Wystarczy przypadek $m=1$, ponieważ stosujemy go
do wszystkich współrzędnych funkcji wektorowej.
Niech $g_1,\ldots,g_n$ będą słabymi pochodnymi z
Lematu~\ref{lem:rademacher-direction-28}. Rozważmy
przeliczalny zbiór kierunków $D=\mathbb Q^n\setminus\{0\}$.
Dla każdego $v\in D$ na zbiorze pełnej miary
istnieje $\partial_v f(x)=\sum_i v_i g_i(x)$.
Przecięcie tych przeliczalnie wielu zbiorów nadal
ma pełną miarę. Ustalmy punkt $x$ tego przecięcia,
w którym wszystkie $g_i(x)$ są skończone, i liniowe
odwzorowanie $A_xh=\sum_i g_i(x)h_i$.

Pokażemy, że $f(x+h)=f(x)+A_xh+o(|h|)$.
Ustalmy $\varepsilon>0$. Z gęstości racjonalnych
kierunków na sferze można wybrać \emph{skończenie}
wiele wektorów jednostkowych
$q_1,\ldots,q_M\in\{v/|v|:v\in D\}$, tak aby
każdy $u\in S^{n-1}$ był w odległości mniejszej
od $\delta$ od któregoś $q_j$, gdzie
$\delta>0$ wybierzemy później. Dla $h=tu$, $t=|h|$,
wybieramy taki $q_j$. Jeżeli $L$ jest stałą
Lipschitza $f$ w otoczeniu $x$, to
\begin{align*}
 |f(x+h)-f(x)-A_xh|
 &\le |f(x+tu)-f(x+tq_j)|\\
 &\quad+|f(x+tq_j)-f(x)-tA_xq_j|
       +|tA_x(q_j-u)|\\
 &\le t(L+\|A_x\|)\delta+o_j(t).
\end{align*}
Ponieważ lista kierunków $q_j$ jest skończona,
po podzieleniu przez $t$ ilorazy $o_j(t)/t$
zbiegają do zera \emph{jednostajnie względem $j$}.
Wybieramy $\delta$ tak, aby
$(L+\|A_x\|)\delta<\varepsilon$, a potem
dostatecznie małe $t$. Wynika stąd
$\limsup_{h\to0}|f(x+h)-f(x)-A_xh|/|h|\le\varepsilon$.
Dowolność $\varepsilon$ daje różniczkowalność
Frécheta z $Df(x)=A_x$.

Pozostaje uzasadnić borelowskość zbioru wyjątkowego.
Dla $A\in\mathbb Q^{m\times n}$, $r>0$ i $x\in U$
połóżmy
\[
 R_{A,r}(x)=\sup_{\substack{h\in\mathbb Q^n,\;0<|h|<r\\
                                    x+h\in U}}
 \frac{|f(x+h)-f(x)-Ah|}{|h|},
\]
przy czym supremum pustego zbioru przyjmujemy równe
zeru. Dla ustalonych $A,r$ funkcja $R_{A,r}$ jest
borelowska: supremum liczymy po przeliczalnie wielu
borelowskich ilorazach na otwartych dziedzinach
$\{x\in U:x+h\in U\}$.
Punkt $x$ jest różniczkowalnościowy wtedy i tylko wtedy,
gdy dla każdego $q\ge1$ istnieją
$A\in\mathbb Q^{m\times n}$ i $j\ge1$ takie, że
$R_{A,1/j}(x)<1/q$. Istotnie, jeżeli $Df(x)$ istnieje,
przybliżamy go macierzą wymierną. W drugą stronę
wybrane macierze $A_q$ spełniają
$\|A_q-A_p\|\le1/q+1/p$:
wspólne małe przyrosty wymierne, po podzieleniu przez
ich długość i przejściu do granicy, oszacowują normę
różnicy tych macierzy. Stąd $A_q\to A$, a z nierówności
resztowej wynika $f(x+h)=f(x)+Ah+o(|h|)$, najpierw
dla wymiernych $h$, potem dla wszystkich $h$ przez
ciągłość $f$. To daje przeliczalny opis borelowski
zbioru różniczkowalności. Jego dopełnienie ma zerową
miarę na mocy poprzedniego argumentu.
Oszacowanie normy różniczki wynika z przejścia do
granicy w nierówności Lipschitza.
\end{proof}

\begin{uwaga}[Dlaczego pochodne cząstkowe nie wystarczają?]
Na przykład $f(x,y)=|x|$ ma pochodne cząstkowe
w punktach z $x\ne0$, ale nie ma różniczki na prostej
$x=0$; zbiór ten jest zerowy dla $\lambda_2$.
Ponadto istnienie wszystkich pochodnych kierunkowych
w \emph{jednym} punkcie nie zapewnia różniczkowalności
Frécheta: dla $f(x,y)=\sqrt{x^2+y^2}$ w $(0,0)$
istnieją jednostronne pochodne kierunkowe, które nie
tworzą odwzorowania liniowego. Istotna w dowodzie
Rademachera jest zgodność pochodnych kierunkowych
z jedną \emph{słabą} różniczką, a następnie
argument skończonej siatki kierunków.
\end{uwaga}

\begin{przyklad}[Wszystkie pochodne cząstkowe w punkcie, ale brak różniczki]
Niech $f:\mathbb R^2\to\mathbb R$ będzie zadana przez
$f(0,0)=0$ oraz $f(x,y)=xy/\sqrt{x^2+y^2}$
dla $(x,y)\ne(0,0)$. Poza początkiem ma ona
ograniczony gradient: jest jednorodna stopnia $1$,
a jej pochodne pierwszego rzędu są ciągłe na
okręgu jednostkowym i jednorodne stopnia $0$.
W związku z tym $f$ jest globalnie Lipschitzowska
(dla odcinka przechodzącego przez $0$ dzielimy go
na dwa odcinki). Obie pochodne cząstkowe w $(0,0)$
istnieją i są równe $0$, ponieważ $f(x,0)=f(0,y)=0$.
Jednak dla $h=(t,t)$, $t>0$, mamy
$f(t,t)=t/\sqrt2$, a więc
$|f(h)-f(0)-0\cdot h|/|h|=1/2$.
Nie istnieje zatem różniczka Frécheta w $(0,0)$.
Pokazuje to, dlaczego sam argument z istnienia
pochodnych cząstkowych prawie wszędzie jest
niewystarczający dla twierdzenia Rademachera.
\end{przyklad}

\begin{przyklad}[Różniczkowalność prawie wszędzie a brak gładkości]
Niech $F:\mathbb R^2\to\mathbb R^2$,
$F(x,y)=(x,y+|x|)$. Jest globalnie Lipschitzowskie,
choć nie jest klasy $C^1$. Dla $x>0$ i $x<0$
ma odpowiednio macierze $\begin{psmallmatrix}1&0\\1&1\end{psmallmatrix}$
i $\begin{psmallmatrix}1&0\\-1&1\end{psmallmatrix}$,
a na prostej $x=0$ nie ma różniczki. Mimo to
$|\det DF|=1$ prawie wszędzie. Wzór zamiany
zmiennych jest poprawny, ale jego uzasadnienie
w klasie Lipschitzowskiej nie może ograniczać się
do zwykłego twierdzenia o dyfeomorfizmach $C^1$.
\end{przyklad}

\begin{uwaga}[Związek z miarami wektorowymi]
Twierdzenie Rademachera wykorzystuje skończony wymiar
\emph{przestrzeni wartości}. Dla dowolnej przestrzeni
Banacha $X$ nie musi być prawdą, że każda funkcja
Lipschitzowska $[0,1]\to X$ ma pochodną prawie wszędzie.
Dokładną granicę wyznacza \emph{własność Radona--Nikodyma},
która zostanie zdefiniowana i scharakteryzowana w
Rozdziale~\ref{ch:vector-measures-30}.
\end{uwaga}

\section{Aproksymacja typu Luzina--Whitneya w klasie $C^1$}
Aby przejść od wersji gładkich do odwzorowań Lipschitzowskich,
nie wystarczy różniczkowalność prawie wszędzie. Potrzebujemy
możliwości zastąpienia odwzorowania na dużych zbiorach
\emph{dokładnie} funkcją klasy $C^1$, wraz z jej różniczką.
Poniżej dowodzimy odpowiedniego faktu, nie przyjmując
jako gotowego twierdzenia o przedłużaniu Whitneya.

\begin{lemat}[Przedłużenie zgodnych danych pierwszego rzędu]
\label{lem:whitney-c1-28}
Niech $K\subseteq\mathbb R^n$ będzie niepustym zbiorem zwartym,
$a:K\to\mathbb R^N$ i
$A:K\to\mathcal L(\mathbb R^n,\mathbb R^N)$ funkcjami ciągłymi.
Załóżmy, że
\begin{equation}\label{eq:whitney-condition-28}
 \lim_{r\downarrow0}\ \sup_{\substack{x,y\in K,\;0<|x-y|\le r}}
 \frac{|a(y)-a(x)-A(x)(y-x)|}{|y-x|}=0,
\end{equation}
przy czym supremum pustego zbioru jest równe zeru.
Istnieje funkcja $\widetilde a\in C^1(\mathbb R^n;\mathbb R^N)$
taka, że $\widetilde a|_K=a$ i
$D\widetilde a(x)=A(x)$ dla $x\in K$.
\end{lemat}
\begin{proof}
Wprowadzamy dwa moduły ciągłości: $\omega(t)$ jest
supremum z \eqref{eq:whitney-condition-28} z $r=t$,
a $\alpha(t)=\sup\{\|A(x)-A(y)\|:x,y\in K,
|x-y|\le t\}$. Obydwa dążą do zera dla $t\downarrow0$.

\emph{Krok 1: kostki otwartego dopełnienia.}
Zbiór $\mathbb R^n\setminus K$ rozkładamy na kostki
diadyczne $Q$ o rozłącznych wnętrzach, dobierając
kostki maksymalne (względem włączenia w rodzinie
kostek diadycznych) spośród tych, dla których
$20\sqrt n\,\ell(Q)\le\operatorname{dist}(Q,K)$.
Tutaj $\ell(Q)$ oznacza długość boku kostki.
Każdy punkt dopełnienia należy do jakiejś kostki:
dostatecznie mała kostka diadyczna zawierająca ten
punkt spełnia nierówność. Maksymalność zapewnia
również oszacowanie odwrotne
$\operatorname{dist}(Q,K)\le C_n\ell(Q)$:
gdyby odległość była dużo większa, kostka macierzysta
(o dwukrotnie większym boku) też spełniałaby warunek.
Zatem długość boku jest porównywalna z odległością
kostki od $K$. W szczególności kostki o przecinających
się nieznacznie powiększonych wnętrzach mają
porównywalne boki, a liczba takich kostek w otoczeniu
dowolnego punktu jest ograniczona przez stałą zależną
wyłącznie od $n$ (wynika to z rozłączności wnętrz
i porównania objętości).

Dla każdej kostki wybieramy $p_Q\in K$ minimalizujący
odległość od jej środka. Wybieramy $C^\infty$-funkcje
$\psi_Q\ge0$, równe $1$ na $Q$, o nośnikach zawartych
w kostkach współśrodkowych o boku $(9/8)\ell(Q)$,
i spełniające $\|D\psi_Q\|\le C_n/\ell(Q)$.
Można je otrzymać jako iloczyny przeskalowanych
jednowymiarowych funkcji odcinających.
Suma $\sum_Q\psi_Q$ jest lokalnie skończona i nie
mniejsza niż $1$ na $\mathbb R^n\setminus K$.
Funkcje $\varphi_Q=\psi_Q/(\sum_R\psi_R)$ tworzą
więc gładki podział jedności, dla którego
$\|D\varphi_Q\|\le C_n/\ell(Q)$.

\emph{Krok 2: sklejanie przybliżeń afinicznych.}
Dla $x\notin K$ definiujemy
\[
 P_Q(x)=a(p_Q)+A(p_Q)(x-p_Q),\qquad
 \widetilde a(x)=\sum_Q\varphi_Q(x)P_Q(x),
\]
a dla $x\in K$ przyjmujemy $\widetilde a(x)=a(x)$.
Niech $x\notin K$, $\delta=\operatorname{dist}(x,K)$,
a $p\in K$ będzie punktem najbliższym $x$.
Dla każdej kostki, której $\varphi_Q(x)\ne0$,
mamy $\ell(Q)\asymp\delta$ i $|p_Q-p|\le C_n\delta$.
Warunek \eqref{eq:whitney-condition-28} daje
\[
 |P_Q(x)-[a(p)+A(p)(x-p)]|
 \le C_n\delta\,[\omega(C_n\delta)+\alpha(C_n\delta)],
\]
zaś $\|A(p_Q)-A(p)\|\le\alpha(C_n\delta)$.
Ponieważ $\sum_QD\varphi_Q=0$, różniczkując
sumę na dopełnieniu, otrzymujemy
\begin{align*}
 |\widetilde a(x)-a(p)-A(p)(x-p)|
 &\le C_n\delta[\omega(C_n\delta)+\alpha(C_n\delta)],\\
 \|D\widetilde a(x)-A(p)\|
 &\le C_n[\omega(C_n\delta)+\alpha(C_n\delta)].
\end{align*}
Stałe mogą również zależeć od wyboru funkcji
odcinających, ale nie od $x$.

Gdy $x\to p_0\in K$ od zewnątrz, mamy $p\to p_0$
i $\delta\to0$; druga nierówność zapewnia
$D\widetilde a(x)\to A(p_0)$. Z pierwszej oraz
\eqref{eq:whitney-condition-28} wynika
$\widetilde a(x)-a(p_0)-A(p_0)(x-p_0)=o(|x-p_0|)$.
Ten sam warunek bezpośrednio daje taki rozwój,
gdy $x\in K$. Zatem $\widetilde a$ jest
różniczkowalna także na $K$, jej różniczka jest
tam równa $A$ i jest ciągła na całej przestrzeni.
\end{proof}

\begin{lemat}[Zwarte fragmenty klasy $C^1$]
\label{lem:lusin-whitney-28}
Niech $U\subseteq\mathbb R^n$ będzie otwarty,
$F:U\to\mathbb R^N$ lokalnie Lipschitzowski,
$Q\subseteq U$ zwarty i $\varepsilon>0$.
Istnieją zbiór zwarty $K\subseteq Q$ oraz
$F_K\in C^1(\mathbb R^n;\mathbb R^N)$ takie, że
\[
 \lambda_n(Q\setminus K)<\varepsilon,\qquad
 F_K(x)=F(x),\quad DF_K(x)=DF(x)\quad(x\in K).
\]
Różniczkę $DF$ rozumiemy tylko w punktach,
w których istnieje; wszystkie punkty $K$ są
takimi punktami. Ponadto $U$ jest sumą przeliczalnej
rodziny zbiorów zwartych $K_j$, z których każdy
ma opisane wyżej przedłużenie, a
$\lambda_n(U\setminus\bigcup_jK_j)=0$.
\end{lemat}
\begin{proof}
Niech $D\subseteq U$ będzie borelowskim zbiorem
punktów różniczkowalności $F$. Na mocy
Twierdzenia~\ref{thm:rademacher-28} dopełnienie
$D$ jest zerowe. Funkcja $x\mapsto DF(x)$ jest
borelowska na $D$, ponieważ jej współczynniki
są granicami ilorazów różnicowych w kierunkach
wektorów bazowych.

Dla $j,l\ge1$ oznaczmy przez $H_{j,l}$ zbiór
punktów $x\in D\cap Q$, dla których
\[
 |F(x+h)-F(x)-DF(x)h|\le j^{-1}|h|
\]
dla każdego $h\in\mathbb Q^n$ spełniającego
$0<|h|<l^{-1}$ i $x+h\in U$. Jest to zbiór
borelowski, bo warunek sprawdzamy dla
przeliczalnej liczby przyrostów. Z różniczkowalności
wynika $\bigcup_lH_{j,l}=D\cap Q$.
Dla każdego $j$ wybierzmy zatem $l_j$ tak duże,
aby $\lambda_n((D\cap Q)\setminus H_{j,l_j})
<\varepsilon2^{-j-2}$. Zbiór
$G=\bigcap_{j\ge1}H_{j,l_j}$ spełnia
$\lambda_n(Q\setminus G)<\varepsilon/2$.
Z twierdzenia Łuzina, zastosowanego do mierzalnej
funkcji macierzowej $DF$ na $G$, i z wewnętrznej
regularności miary Lebesgue'a otrzymujemy zwarty
$K\subseteq G$, dla którego
$\lambda_n(G\setminus K)<\varepsilon/2$ oraz
$DF|_K$ jest ciągła. Jeśli $K=\varnothing$,
przyjmujemy $F_K\equiv0$ i pierwsza teza jest
trywialna; w dalszym ciągu zakładamy $K\ne\varnothing$.

Warunek w definicji $H_{j,l_j}$, przez ciągłość
$F$ względem przyrostu $h$, pozostaje prawdziwy
również dla rzeczywistych $h$ o tej samej
ograniczonej długości. Stąd dla każdego $j$ i
$x,y\in K$ o $0<|x-y|<l_j^{-1}$ zachodzi
\[
 |F(y)-F(x)-DF(x)(y-x)|\le j^{-1}|x-y|.
\]
Zatem dane $a=F|_K$ i $A=DF|_K$ spełniają
założenia Lematu~\ref{lem:whitney-c1-28}.
Jego zastosowanie daje żądane przedłużenie.

Aby uzyskać ostatnią tezę, wybieramy przeliczalne
pokrycie $U$ zwartymi kostkami $Q_s\Subset U$.
Dla każdego $s$ i $j$ wykonujemy konstrukcję
z $\varepsilon=1/j$ na $Q_s$, uzyskując
$K_{s,j}$. Ponieważ
$\lambda_n(Q_s\setminus\bigcup_jK_{s,j})\le1/j$
dla każdego $j$, miara ta jest zerowa.
Przeliczalna suma pozostałych zbiorów zerowych
ma miarę zero, co kończy dowód.
\end{proof}

\section{Dwa oszacowania dla zbiorów krytycznych}
\begin{lemat}[Obraz zbioru o niedostatecznej randze]
\label{lem:critical-area-null-28}
Niech $1\le k\le N$, $V\subseteq\mathbb R^k$ będzie
otwarty, $G\in C^1(V;\mathbb R^N)$, a
$K\subseteq V$ zwarty. Jeżeli
$\operatorname{rank}DG(x)<k$ dla wszystkich $x\in K$,
to $\mathcal H^k(G(K))=0$.
\end{lemat}
\begin{proof}
Niech $M$ ogranicza normę $DG$ na pewnym
zwartym otoczeniu $K$. Z jednostajnej ciągłości
$DG$ wynika, że dla każdego $\eta>0$ i dostatecznie
małej kostki $Q$ o boku $r$, przecinającej $K$,
po wybraniu $x_Q\in K\cap Q$ mamy
\[
 |G(x)-G(x_Q)-DG(x_Q)(x-x_Q)|\le\eta\sqrt k\,r
 \quad(x\in K\cap Q).
\]
Obraz $G(K\cap Q)$ leży więc w otoczeniu o
grzubości $C_k\eta r$ zbioru zawartego
w afinicznej podprzestrzeni wymiaru co najwyżej
$k-1$ i promienia $C_kMr$. Taki zbiór można
pokryć najwyżej $C_{k,N,M}\eta^{-(k-1)}$
kulami promienia $C_k\eta r$ (siatka w owej
podprzestrzeni); suma ich $k$-tych potęg
średnic nie przekracza $C_{k,N,M}\eta r^k$.

Ponieważ $K$ mieści się w ustalonej ograniczonej
kostce, liczba kostek siatki o boku $r$
przecinających $K$ jest nie większa niż $Cr^{-k}$.
Otrzymujemy pokrycie $G(K)$ zbiorami o średnicy
nie większej niż $C\eta r$, którego suma
$k$-tych potęg średnic jest nie większa niż
$C_{k,N,M}\eta$. Gdy $r\downarrow0$, dopuszczalne
średnice także dążą do zera, więc
$\mathcal H^k(G(K))\le C_{k,N,M}\eta$.
Przechodząc $\eta\downarrow0$, dostajemy tezę.
\end{proof}

\begin{lemat}[Zbiory zerowe i ich włókna]
\label{lem:null-fibres-lipschitz-28}
Niech $n\ge m\ge1$, $F:\mathbb R^n\to\mathbb R^m$
będzie lokalnie Lipschitzowska, a
$Z\subseteq\mathbb R^n$ zbiorem o zerowej zewnętrznej
mierze Lebesgue'a. Wówczas
\[
 \mathcal H^{n-m}(Z\cap F^{-1}(\{y\}))=0
 \quad\text{dla }\lambda_m\text{-prawie wszystkich }y.
\]
\end{lemat}
\begin{proof}
Przez przeliczalne pokrycie dziedziny otwartymi
kostkami, na których $F$ ma wspólną stałą
Lipschitza $L$, wystarczy rozpatrzyć część $Z$
zawartą w jednej ograniczonej takiej kostce.
Piszemy $d=n-m$. Ustalmy $\delta,\varepsilon>0$.
Z warunku $\lambda_n^*(Z)=0$ wynika istnienie
przeliczalnego pokrycia $Z$ kostkami $Q_i$
o średnicach mniejszych niż $\delta$ i takiego, że
$\sum_i(\operatorname{diam}Q_i)^n<\varepsilon$.
Kostki możemy wybrać wewnątrz otoczenia,
na którym działa stała $L$.
Niech $z_i$ będzie środkiem $Q_i$ i niech
$V_i=B(F(z_i),L\operatorname{diam}Q_i)$.
Jeśli $y\in F(Z\cap Q_i)$, to $y\in V_i$.
Z definicji miary Hausdorffa z ograniczeniem
średnic uzyskujemy punktowo
\[
 \mathcal H^d_\delta(Z\cap F^{-1}(\{y\}))
 \le C_{n,m}\sum_i(\operatorname{diam}Q_i)^d\chi_{V_i}(y).
\]
Dla $d=0$ koszt niepustego zbioru w pokryciu
wynosi $1$. Prawa strona jest borelowska,
a całkując ją, otrzymujemy
\[
 \int_{\mathbb R^m}C_{n,m}\sum_i
 (\operatorname{diam}Q_i)^d\chi_{V_i}(y)
 \,\mathrm d\lambda_m(y)
 \le C_{n,m}L^m\varepsilon.
\]
Zatem dla dowolnego $t>0$ zbiór, na którym
$\mathcal H^d_\delta$ rozpatrywanego włókna
przekracza $t$, ma zewnętrzną miarę nie większą
niż $C_{n,m}L^m\varepsilon/t$. Ponieważ
$\varepsilon$ jest dowolne, miara ta jest zerowa.
Stosując to do $\delta=1/j$ i sumując
przeliczalnie zbiory wyjątkowe, dostajemy
$\mathcal H^d=0$ na prawie wszystkich włóknach.
\end{proof}

\begin{lemat}[Włókna zwartego zbioru krytycznego]
\label{lem:critical-coarea-null-28}
Niech $n\ge m\ge1$, $V\subseteq\mathbb R^n$
będzie otwarty, $G\in C^1(V;\mathbb R^m)$,
a $K\subseteq V$ zwarty. Jeżeli
$\operatorname{rank}DG(x)<m$ dla wszystkich $x\in K$,
to
\[
 \mathcal H^{n-m}(K\cap G^{-1}(\{y\}))=0
 \quad\text{dla }\lambda_m\text{-prawie wszystkich }y.
\]
\end{lemat}
\begin{proof}
Piszemy $d=n-m$. Najpierw zauważmy, że dla
$\delta>0$ funkcja
$y\mapsto\mathcal H^d_\delta(K\cap G^{-1}(\{y\}))$
jest borelowska. Rzeczywiście, dla $t>0$ warunek,
że jej wartość jest mniejsza niż $t$, można
wyrazić istnieniem skończonego pokrycia włókna
otwartymi zbiorami będącymi skończonymi
sumami kul o wymiernych środkach i promieniach,
średnicach mniejszych niż $\delta$ i sumie kosztów
mniejszej niż $t$. Dla ustalonego takiego pokrycia
$\mathcal U$ zbiór parametrów $y$, dla których
całe włókno mieści się w $\mathcal U$, jest
$\mathbb R^m\setminus G(K\setminus\mathcal U)$,
a więc jest otwarty, gdyż $K\setminus\mathcal U$
jest zwarty. Przeliczalna suma tych otwartych
zbiorów daje borelowskość. W konsekwencji
również funkcja dla $\delta\downarrow0$ jest
borelowska jako granica monotoniczna.

Ustalmy $\eta>0$. Pokryjmy $K$ kostkami siatki
o boku $r$, dla których wybieramy
$x_Q\in K\cap Q$. Gdy $r$ jest dostatecznie
mały, jednostajna ciągłość $DG$ daje
\[
 |G(x)-G(x_Q)-DG(x_Q)(x-x_Q)|\le C_n\eta r
 \quad(x\in K\cap Q).
\]
Ponieważ ranga $DG(x_Q)$ jest mniejsza niż $m$,
zbiór $G(K\cap Q)$ zawiera się w
$C_n\eta r$-otoczeniu afinicznej podprzestrzeni
wymiaru nie większego niż $m-1$, w obrębie
kuli o promieniu $C_nMr$, gdzie $M$ jest
lokalnym ograniczeniem $\|DG\|$.
Porównanie objętości po rozłożeniu
$\mathbb R^m$ na tę podprzestrzeń i jej
ortogonalne dopełnienie daje
\[
 \lambda_m(G(K\cap Q))\le C_{n,m,M}\eta r^m.
\]
Dla każdego $y$ zbiór $K\cap G^{-1}(\{y\})$
można pokryć fragmentami leżącymi w tych
kostkach $Q$, dla których $y\in G(K\cap Q)$.
Każdy fragment ma średnicę nie większą niż
$\sqrt n\,r$, więc dla $\delta_r=2\sqrt n\,r$
\[
 \mathcal H^d_{\delta_r}(K\cap G^{-1}(\{y\}))
 \le C_{n,m}\sum_Q r^d\chi_{G(K\cap Q)}(y).
\]
W sumie występuje najwyżej $Cr^{-n}$ kostek.
Całkując i stosując oszacowanie ich obrazów,
dostajemy
\[
 \int_{\mathbb R^m}\mathcal H^d_{\delta_r}
 (K\cap G^{-1}(\{y\}))\,\mathrm d\lambda_m(y)
 \le C_{n,m,M}\eta.
\]
Przechodząc $r\downarrow0$ (wzdłuż malejącego
ciągu) i korzystając z monotoniczności
$\mathcal H^d_\delta$ oraz twierdzenia o
zbieżności monotonicznej, otrzymujemy
$\int\mathcal H^d(K\cap G^{-1}(\{y\}))
\,\mathrm d\lambda_m(y)\le C\eta$.
Dowolność $\eta$ kończy dowód.
\end{proof}


\section{Wzory area i coarea dla odwzorowań Lipschitzowskich}
W poniższych wzorach $DF$ oznacza różniczkę z twierdzenia
Rademachera, istniejącą prawie wszędzie. Jakobian $J_kF$
oraz jakobian normalny $\operatorname{Jac}_mF$ obliczamy
w punktach różniczkowalności, podstawiając $DF(x)$
do funkcji jakobianowych z
Definicji~\ref{def:k-jacobian-28} i
Definicji~\ref{def:normal-jacobian-28},
a na pozostałym zbiorze zerowym przypisujemy im wartość $0$.
Tak określone funkcje są mierzalne, a wybór wartości
na zbiorze wyjątkowym nie zmienia całek.
\begin{twierdzenie}[Wzory area i coarea dla odwzorowań Lipschitzowskich]
\label{thm:lipschitz-area-coarea-28}
Niech $E\subseteq\mathbb R^k$ będzie mierzalny
w sensie Lebesgue'a, a $F:\mathbb R^k\to\mathbb R^N$,
$k\le N$, będzie lokalnie Lipschitzowski.
Wówczas różniczka $DF(x)$ istnieje dla prawie
wszystkich $x$ (twierdzenie Rademachera)
i dla każdej mierzalnej $f\ge0$ zachodzi
\[
 \int_E f(x)J_kF(x)\,\mathrm d\lambda_k(x)
 =\int_{\mathbb R^N}\left(\sum_{x\in E\cap F^{-1}(\{y\})}
 f(x)\right)\mathrm d\mathcal H^k(y).
\]
Ponadto jeżeli $n\ge m\ge1$ i
$F:\mathbb R^n\to\mathbb R^m$ jest lokalnie
Lipschitzowski, to dla każdej mierzalnej
funkcji $g\ge0$ zachodzi
\[
 \int_{\mathbb R^n}g(x)\operatorname{Jac}_mF(x)
 \,\mathrm d\lambda_n(x)
 =\int_{\mathbb R^m}\left(
 \int_{F^{-1}(\{y\})}g(x)
 \,\mathrm d\mathcal H^{n-m}(x)\right)
 \,\mathrm d\lambda_m(y).
\]
W przypadku $n=m$ wewnętrzna całka po włóknie
jest sumą względem miary liczącej.
\end{twierdzenie}

\begin{lemat}[Wzór area dla dowolnego odwzorowania $C^1$]
\label{lem:c1-general-area-28}
Niech $1\le k\le N$, $V\subseteq\mathbb R^k$ będzie
otwarty, $G\in C^1(V;\mathbb R^N)$, a $A\subseteq V$
borelowski. Dla każdej nieujemnej funkcji borelowskiej
$f:A\to[0,+\infty]$ zachodzi
\[
 \int_Af(x)J_kG(x)\,\mathrm d\lambda_k(x)
 =\int_{\mathbb R^N}\left(
 \sum_{x\in A\cap G^{-1}(\{y\})}f(x)\right)
 \,\mathrm d\mathcal H^k(y).
\]
\end{lemat}
\begin{proof}
Zbiór $V_{\mathrm{reg}}=\{x\in V:\operatorname{rank}DG(x)=k\}$
jest otwarty, a $G|_{V_{\mathrm{reg}}}$ jest immersją.
Stosujemy do niego
Twierdzenie~\ref{thm:immersion-area-28}.
Jego dopełnienie w $V$ jest względnie domkniętym
zbiorem punktów krytycznych, który można wyczerpać
przeliczalną rodziną zwartych podzbiorów.
Lemat~\ref{lem:critical-area-null-28} dowodzi, że obraz
tych punktów jest $\mathcal H^k$-zerowy.
Na części krytycznej $J_kG=0$, a suma po włóknach
może być różna od zera tylko na wspomnianym
zbiorze zerowym. Równości po zsumowaniu obu
części są więc prawdziwe.
\end{proof}

\begin{proof}[Dowód ogólnych wzorów area i coarea z
Twierdzenia~\ref{thm:lipschitz-area-coarea-28}]
\emph{Krok 1: wzór area.}
Z Lematu~\ref{lem:lusin-whitney-28} otrzymujemy
zwarte zbiory $K_j\subseteq\mathbb R^k$ oraz
odwzorowania $F_j\in C^1(\mathbb R^k;\mathbb R^N)$
takie, że $F_j=F$ i $DF_j=DF$ na $K_j$, a
$Z=\mathbb R^k\setminus\bigcup_jK_j$ jest
$\lambda_k$-zerowy. Połóżmy
$A_1=K_1$ i
$A_j=K_j\setminus\bigcup_{i<j}K_i$ dla $j\ge2$.
Są to parami rozłączne zbiory borelowskie,
których suma jest dopełnieniem $Z$.

Dla $E$ i $f$ borelowskich stosujemy
Lemat~\ref{lem:c1-general-area-28} do
$F_j$ na $A_j\cap E$. Daje on
\[
 \int_{A_j\cap E}f(x)J_kF(x)\,\mathrm d\lambda_k(x)
 =\int_{\mathbb R^N}
 \sum_{x\in A_j\cap E\cap F^{-1}(\{y\})}f(x)
 \,\mathrm d\mathcal H^k(y).
\]
Po zsumowaniu względem $j$ stosujemy Tonellego.
Zbiór $Z$ nie wnosi nic do lewej strony.
Ponieważ $F$ jest lokalnie Lipschitzowskie,
$\mathcal H^k(F(Z))=0$: dla każdego zwartego
fragmentu dziedziny wynika to z nierówności
$\mathcal H^k(F(A))\le L^k\mathcal H^k(A)$
z Rozdziału~\ref{ch:hausdorff-26}.
Zatem również prawa strona nie zmienia się
po dołączeniu punktów z $Z$.

Jeżeli $E$ jest tylko mierzalny w sensie
Lebesgue'a, a $f$ jest mierzalna, wybieramy
borelowskie przedstawiciele $E$ i $f$ różniące
się od zadanych jedynie na zbiorze
$\lambda_k$-zerowym. Jego obraz przez $F$
jest $\mathcal H^k$-zerowy, więc ta modyfikacja
nie zmienia żadnej strony równości. Zapewnia
to również mierzalność prawej strony względem
uzupełnionej miary $\mathcal H^k$.

\emph{Krok 2: wzór coarea na częściach pełnej rangi.}
Przyjmijmy na razie, że $g$ jest borelowska.
Dla odwzorowania $F:\mathbb R^n\to\mathbb R^m$
wybieramy analogiczne $K_j$ i $F_j\in C^1$,
z rozłącznymi zbiorami $A_j\subseteq K_j$.
Niech
$V_j=\{x:\operatorname{rank}DF_j(x)=m\}$.
Jest to zbiór otwarty, a $F_j|_{V_j}$ jest
submersją. Twierdzenie~\ref{thm:coarea-submersion-28}
zastosowane do funkcji
$g\chi_{A_j\cap V_j}$ daje
\begin{align*}
 &\int_{A_j\cap V_j}g(x)\operatorname{Jac}_mF(x)
 \,\mathrm d\lambda_n(x)\\
 &\quad=\int_{\mathbb R^m}\left(
 \int_{A_j\cap V_j\cap F^{-1}(\{y\})}
 g(x)\,\mathrm d\mathcal H^{n-m}(x)\right)
 \,\mathrm d\lambda_m(y).
\end{align*}

\emph{Krok 3: części krytyczne i zbiór wyjątkowy.}
Zbiór
$C_j=K_j\cap\{x:\operatorname{rank}DF_j(x)<m\}$
jest zwarty. Na nim $\operatorname{Jac}_mF=0$,
a Lemat~\ref{lem:critical-coarea-null-28}
pokazuje, że
$\mathcal H^{n-m}(C_j\cap F^{-1}(\{y\}))=0$
dla prawie wszystkich $y$. Na pozostałym zbiorze
$Z=\mathbb R^n\setminus\bigcup_jK_j$ jakobian
również nie wnosi nic do całki, bo
$\lambda_n(Z)=0$. Z kolei
Lemat~\ref{lem:null-fibres-lipschitz-28} dowodzi,
że $Z$ jest $\mathcal H^{n-m}$-zerowy na
prawie wszystkich włóknach. Stąd po zsumowaniu
równości z kroku 2 otrzymujemy pełny wzór coarea.

W przypadku funkcji $g$ jedynie mierzalnej
w sensie Lebesgue'a zastępujemy ją wersją
borelowską poza zbiorem zerowym. Lemat o
zerowych włóknach gwarantuje, że taka zmiana
nie wpływa na prawą stronę dla prawie wszystkich
$y$; jej mierzalność wynika z rozpatrzenia
części borelowskich i uzupełnienia miary.
Nie trzeba w żadnym kroku zakładać, że odwzorowanie
$F$ jest globalnie różnowartościowe lub że ma
stałą rangę.
\end{proof}

\begin{uwaga}[Zakres uzyskanego dowodu]
W porównaniu z wersjami gładkimi nowym elementem
jest jednoczesna kontrola zbiorów krytycznych i ich
włókien. Lemat~\ref{lem:whitney-c1-28} wyprowadza
potrzebne przedłużenie $C^1$, a lematy
\ref{lem:critical-area-null-28}--\ref{lem:critical-coarea-null-28}
wyprowadzają oszacowania bez powoływania się na
same wzory area/coarea. Dzięki temu dowód nie jest
błędnym kołem. Nadal poza zakresem pozostają
analogiczne wzory na ogólnych przestrzeniach
metrycznych oraz dla odwzorowań, które nie są
lokalnie Lipschitzowskie.
\end{uwaga}

\section{Postać na rozmaitościach}
Niech $(M^n,g)$ i $(N^m,h)$ będą rozmaitościami
riemannowskimi, $n\ge m$, a $F:M\to N$
submersją $C^1$. Różniczkę i jej sprzężenie
$\mathrm dF_p^*$ wyznaczają metryki $g,h$.
Jakobian normalny definiujemy jako
$\operatorname{Jac}_{g,h}F(p)=
\sqrt{\det(\mathrm dF_p\mathrm dF_p^*)}$,
gdzie wyznacznik liczymy w bazie ortonormalnej
$T_{F(p)}N$. Dla nieujemnej borelowskiej $\varphi$
analogiczny lokalny rachunek daje
\begin{equation}\label{eq:coarea-manifolds-28}
 \int_M\varphi(p)\operatorname{Jac}_{g,h}F(p)
 \,\mathrm d\mu_g(p)
 =\int_N\left(\int_{F^{-1}(\{q\})}\varphi(p)
 \,\mathrm d\mu_{g|_{F^{-1}(\{q\})}}(p)\right)
 \,\mathrm d\mu_h(q).
\end{equation}
Dowód redukuje się w mapach do
Twierdzenia~\ref{thm:coarea-submersion-28}:
zmiana objętości w mapach źródła, celu i włókna
prowadzi do tej samej tożsamości wyznaczników
co w Lemacie~\ref{lem:linear-coarea-28}, zapisanej
w bazach ortonormalnych. Lokalne równości
łączymy przeliczalnym podziałem borelowskim.
Nie wykorzystujemy orientowalności.

\begin{wniosek}[Miara riemannowska i Hausdorffa]
\label{cor:riemann-hausdorff-28}
Na gładkiej rozmaitości riemannowskiej $(M^n,g)$
miara $\mu_g$ jest równa $n$-wymiarowej mierze
Hausdorffa zbudowanej z odległości riemannowskiej,
przy konwencji kalibracji z
Uwagi~\ref{rem:hausdorff-convention-26}. Na osadzonej
gładkiej podrozmaitości $S^k\subseteq M$ zachodzi
analogiczna identyfikacja miary indukowanej
$\mu_{g|_S}$ z $\mathcal H^k$ obliczaną
z odległości ograniczonej z $M$ na $S$.
\end{wniosek}
\begin{proof}
Wokół każdego punktu wybieramy mapę normalną
lub dowolną mapę, w której metryka ma ciągłe
współczynniki. Na dostatecznie małym otoczeniu
metryczne odległości są, z błędem mnożnikowym
$1+\varepsilon$, porównywalne z odległościami
określonymi przez stały iloczyn skalarny
$g_p$. Wynika to z ciągłości $g$ i
lokalnego porównania długości krzywych
(zmniejszając otoczenie tak, aby minimalizujące
krzywe nie opuszczały większego otoczenia).
Miara Hausdorffa tych samych zbiorów różni się
wtedy co najwyżej o czynnik $(1+\varepsilon)^n$
od miary Hausdorffa dla metryki stałej.
Ta ostatnia ma gęstość $\sqrt{\det G(p)}$
względem $\lambda_n$. Przez lokalną
regularność miar i przejście $\varepsilon\downarrow0$
otrzymujemy równość miar w mapie. Przeliczalne
pokrycie mapami daje ją globalnie.
Dla $S$ powtarzamy argument dla metryki
indukowanej; odległość wewnętrzna na $S$
i zewnętrzna odległość ograniczona z $M$
są lokalnie asymptotycznie równoważne,
więc ich miary Hausdorffa wymiaru $k$ są równe.
\end{proof}

\begin{uwaga}[Co przechodzi do geometrycznej teorii miary?]
Następny rozdział zastosuje miary Hausdorffa,
formuły area i coarea oraz twierdzenia o
różniczkowaniu miar do zbiorów prostowalnych
i zbiorów o skończonym obwodzie. Te obiekty
nie muszą być gładkimi podrozmaitościami,
a zatem stanowią właściwe rozszerzenie
geometrii różniczkowej.
\end{uwaga}

\begin{uwaga}[Dżety przybliżone a wzory miarowo-geometryczne]
Pochodna $DF(x)$ z twierdzenia Rademachera jest częścią
pierwszego dżetu odwzorowania $F$ w punktach różniczkowalności.
Przy osłabieniu pojęcia granicy do granicy przybliżonej otrzymujemy
\emph{dżety przybliżone}, określone miarowo prawie wszędzie.
W rozdziale~\ref{ch:dzety} rozwijamy geometryczną teorię wiązek
dżetów, omawiamy twierdzenie Whitneya o przedłużaniu danych
dżetowych oraz jego miarową wersję Whitneya--Luzina.
Jest to uzupełnienie obecnego rozdziału, a nie ponowny dowód
wzorów area i coarea.
\end{uwaga}

% END chapters/pl/measure_theory/28-rozdzial.tex


% BEGIN chapters/pl/measure_theory/29-rozdzial.tex
% ROZDZIAŁ 29 -- podstawy geometrycznej teorii miary
\chapter[Prostowalność i zbiory o skończonym obwodzie]{Zbiory prostowalne, funkcje BV i zbiory o skończonym obwodzie}
\label{ch:geometric-measure-29}

\section{Cel i zakres teorii}
W geometrii różniczkowej badaliśmy gładkie podrozmaitości, ich
metryki indukowane oraz normalne. W tym rozdziale będziemy mierzyć
również brzegi zbiorów, które mają narożniki, pęknięcia lub
nieregularności i nie są podrozmaitościami. Zamiast zakładać
istnienie klasycznej normalnej w każdym punkcie, zdefiniujemy ją
\emph{miarowo}. Przygotujemy w ten sposób język geometrycznej
teorii miary, teorii wariacyjnej i równań różniczkowych cząstkowych.

Pracujemy na $\mathbb R^n$, gdzie $n\ge1$, z miarą Lebesgue'a
$\lambda_n$ i miarami Hausdorffa $\mathcal H^s$
z Rozdziału~\ref{ch:hausdorff-26}; dla wymiarów
całkowitych stosujemy dokładnie kalibrację
z Uwagi~\ref{rem:hausdorff-convention-26}.
W szczególności $\mathcal H^{n-1}$ ma tę samą
normalizację co we wzorze area z
Rozdziału~\ref{ch:area-coarea-28}, a gdy $n=1$,
$\mathcal H^0$ jest miarą liczącą. Oznaczenie $\chi_E$ oznacza
funkcję charakterystyczną $E$. Wszystkie przestrzenie $L^1$
są przestrzeniami klas funkcji równych prawie wszędzie.
Stosujemy przy tym konwencję dystrybucyjną
$\langle D_i u,\varphi\rangle=-\int_{\mathbb R^n}u(x)
\partial_i\varphi(x)\,\mathrm d\lambda_n(x)$ dla
$\varphi\in C_c^1(\mathbb R^n)$.

\section{Prostowalność: co pozostaje z gładkiej powierzchni?}
\begin{definicja}[Zbiory przeliczalnie prostowalne]
\label{def:rectifiable-29}
Niech $1\le k\le n$. Zbiór $E\subseteq\mathbb R^n$ nazywamy
\emph{przeliczalnie $k$-prostowalnym}, jeśli istnieją
zbiory $A_j\subseteq\mathbb R^k$ oraz odwzorowania
Lipschitzowskie $f_j:A_j\to\mathbb R^n$, $j\in\mathbb N$,
takie że
\[
 \mathcal H^k\!\left(E\setminus\bigcup_{j=1}^\infty f_j(A_j)\right)=0.
\]
Mówimy, że $E$ jest \emph{$k$-prostowalny o skończonej mierze},
jeżeli dodatkowo $\mathcal H^k(E)<\infty$.
Zbiór $P\subseteq\mathbb R^n$ jest \emph{czysto
$k$-nieprostowalny}, gdy
$\mathcal H^k(P\cap f(A))=0$ dla każdego
$A\subseteq\mathbb R^k$ i każdego odwzorowania
Lipschitzowskiego $f:A\to\mathbb R^n$.
\end{definicja}

Dla $k=0$ zbiór nazywamy przeliczalnie $0$-prostowalnym,
jeśli jest przeliczalny; wtedy $\mathcal H^0$ jest miarą liczącą.

W tej definicji nie utożsamiamy prostowalności z topologiczną
gładkością. Funkcja Lipschitzowska może mieć punkty załamania,
a zbiór prostowalny może być przeliczalną sumą rozłącznych
fragmentów powierzchni. Warunek z miarą Hausdorffa pozwala
pominąć tylko zbiór $k$-wymiarowo zerowy, nie dowolny zbiór
mały w sensie miary $n$-wymiarowej.

\begin{przyklad}[Wykres z narożnikiem]
\label{ex:lip-graph-29}
Zbiór $E=\{(t,|t|):-1\le t\le1\}\subseteq\mathbb R^2$
jest $1$-prostowalny: to obraz przedziału przez odwzorowanie
$F(t)=(t,|t|)$, którego stała Lipschitza jest nie większa
niż $\sqrt2$. W punkcie $(0,0)$ nie jest on jednowymiarową
podrozmaitością klasy $C^1$. Mimo to $\mathcal H^1(E)$
jest skończona i dodatnia: rzut $E\to[-1,1]$ jest
$1$-Lipschitzowski, a odwzorowanie $F$ jest
$\sqrt2$-Lipschitzowskie, więc
$2\le\mathcal H^1(E)\le2\sqrt2$.
Złożenie dwóch odcinków daje nawet
$\mathcal H^1(E)=2\sqrt2$ na podstawie addytywności
miary na fragmentach przecinających się tylko w jednym punkcie.
\end{przyklad}

\begin{lemat}[Podstawowe własności prostowalności]
\label{lem:rectifiable-union-29}
Niech $1\le k\le n$. Wtedy:
\begin{enumerate}[label=(\roman*)]
\item przeliczalna suma zbiorów przeliczalnie
$k$-prostowalnych jest przeliczalnie $k$-prostowalna;
\item każdy podzbiór zbioru przeliczalnie
$k$-prostowalnego jest przeliczalnie $k$-prostowalny;
\item obraz zbioru przeliczalnie $k$-prostowalnego
przez odwzorowanie Lipschitzowskie określone na jego
otoczeniu jest przeliczalnie $k$-prostowalny.
\end{enumerate}
\end{lemat}
\begin{proof}
W (i) przestawiamy pary indeksów w jedną rodzinę
przeliczalną. W (ii) korzystamy z monotoniczności
miary zewnętrznej Hausdorffa: zbiór wyjątkowy
nowego podzbioru mieści się w zbiorze wyjątkowym
pierwotnego pokrycia. W (iii) ograniczamy każde $f_j$ do zbioru
$A_j' = \{x\in A_j:f_j(x)\in E\}$, tak aby złożenie
$G\circ f_j|_{A_j'}$ było określone. Jeśli $G$ ma stałą
Lipschitza $L$ w otoczeniu rozważanego zbioru,
to te złożenia są Lipschitzowskie, a obraz
wyjątkowego zbioru ma zerową miarę $\mathcal H^k$,
bo $\mathcal H^k(G(A))\le L^k\mathcal H^k(A)$.
Dla odwzorowania lokalnie Lipschitzowskiego rozbijamy
dziedzinę na przeliczalnie wiele obszarów, na których
jest ono Lipschitzowskie.
\end{proof}

\begin{uwaga}[Rozkład prostowalny -- granica argumentu elementarnego]
Przy założeniu $\mathcal H^k(E)<\infty$ istnieje rozkład
$E=R\,\dot\cup\,P$ (z dokładnością do zbiorów zerowych),
w którym $R$ jest przeliczalnie $k$-prostowalny,
a $P$ jest czysto $k$-nieprostowalny.
Jego dowód można oprzeć na maksymalizacji miary części
prostowalnej, ale w pełnej ogólności potrzebna jest
\emph{mierzalność obrazów borelowskich przez odwzorowania
Lipschitzowskie} (teoria zbiorów analitycznych).
Nie przedstawiamy tego argumentu jako samodzielnego dowodu:
w dalszej części wystarczą definicja prostowalności,
Lemat~\ref{lem:rectifiable-union-29} oraz jawne przykłady
wykresów i brzegów kawałkami gładkich.
\end{uwaga}

\section{Dystrybucyjna pochodna funkcji całkowalnej lokalnie}
\begin{definicja}[Wariacja funkcji]
\label{def:BV-variation-29}
Niech $\Omega\subseteq\mathbb R^n$ będzie otwarty,
a $u\in L^1_{\mathrm{loc}}(\Omega)$. Definiujemy
\[
 |Du|(\Omega)=\sup\left\{
 \int_\Omega u(x)\operatorname{div}\varphi(x)
   \,\mathrm d\lambda_n(x):
 \varphi\in C_c^1(\Omega;\mathbb R^n),\quad
 |\varphi(x)|\le1\ \text{dla wszystkich }x\in\Omega
 \right\}.
\]
Przez $BV(\Omega)$ oznaczamy przestrzeń
funkcji $u\in L^1(\Omega)$, dla których
$|Du|(\Omega)<\infty$. Normę definiujemy wzorem
$\|u\|_{BV}=\|u\|_{L^1}+|Du|(\Omega)$.
W zastosowaniach do obszarów nieograniczonych
wygodna jest również lokalna klasa $BV_{\mathrm{loc}}(\Omega)$:
na każdym $V\Subset\Omega$ funkcja ma skończoną
wariację $|Du|(V)$.
\end{definicja}

W definicji nie jest wymagane istnienie klasycznej
pochodnej $u$. Symbol $Du$ na początku oznacza
funkcjonał dystrybucyjny; dopiero niżej udowodnimy,
że przy skończonej wariacji reprezentuje go
wektorowa miara Radona.

\begin{lemat}[Wariacja gładkiej funkcji]
\label{lem:BV-smooth-29}
Jeżeli $u\in C^1(\Omega)\cap L^1_{\mathrm{loc}}(\Omega)$,
to
\[
 |Du|(\Omega)=\int_\Omega|\nabla u(x)|
 \,\mathrm d\lambda_n(x)\in[0,+\infty].
\]
\end{lemat}
\begin{proof}
Całkowanie przez części dla testów o zwartym nośniku
daje $\int_\Omega u\operatorname{div}\varphi\,\mathrm d\lambda_n=-\int_\Omega\nabla u\cdot\varphi\,\mathrm d\lambda_n$.
Z nierówności Cauchy'ego--Schwarza wynika
$|Du|(\Omega)\le\int_\Omega|\nabla u|\,\mathrm d\lambda_n$.
Dla nierówności odwrotnej ustalmy $K\Subset\Omega$
i $\varepsilon>0$. Na $K$ pole
$-\nabla u/(|\nabla u|+\varepsilon)$ jest ciągłe i ma normę
nie większą niż $1$. Aproksymujemy je jednostajnie na
$K$ przez pole $C_c^1(\Omega;\mathbb R^n)$ o normie
nieprzekraczającej $1$, mnożąc przez funkcję odcinającą
i ewentualnie przez czynnik bliski $1$.
Z definicji wariacji dostajemy
$|Du|(\Omega)\ge\int_K|\nabla u|^2/(|\nabla u|+\varepsilon)$.
Przez $\varepsilon\downarrow0$ i wyczerpanie $\Omega$
zbiorami zwartymi otrzymujemy przeciwną nierówność.
\end{proof}

Przez miarę wektorową o wartościach w $\mathbb R^n$
rozumiemy w tym miejscu przeliczalnie addytywną
w normie funkcję zbioru; ogólna definicja i teoria
wariacji takich miar znajduje się w
Rozdziale~\ref{ch:vector-measures-30}.

\begin{twierdzenie}[Pochodna $Du$ jako miara wektorowa]
\label{thm:BV-measure-29}
Niech $u\in L^1_{\mathrm{loc}}(\Omega)$ oraz
$|Du|(\Omega)<\infty$. Wtedy istnieje jedyna
skończona miara wektorowa $Du=(D_1u,\ldots,D_nu)$,
której współrzędne są skończonymi miarami ze znakiem,
i
\[
 \int_\Omega u(x)\operatorname{div}\varphi(x)
 \,\mathrm d\lambda_n(x)
 =-\int_\Omega\varphi(x)\cdot\mathrm dDu(x)
 \qquad(\varphi\in C_c^1(\Omega;\mathbb R^n)).
\]
Jej wariacja całkowita spełnia
$|Du|(\Omega)=\sup_{|\varphi|\le1}|
\int_\Omega \varphi\cdot\mathrm dDu|$,
gdzie supremum przebiega testy z definicji.
Jeżeli $u\in BV_{\mathrm{loc}}(\Omega)$,
analogicznie otrzymujemy lokalnie skończoną
wektorową miarę Radona.
\end{twierdzenie}
\begin{proof}
Oznaczmy przez $V=|Du|(\Omega)$ wartość supremum.
Dla każdego $i$ funkcjonał
$T_i(\psi)=-\int_\Omega u\partial_i\psi\,\mathrm d\lambda_n$
spełnia $|T_i(\psi)|\le V\|\psi\|_\infty$,
gdyż testujemy pole $\varphi=\psi e_i/\|\psi\|_\infty$
(zera nie wymagają osobnego przypadku). Funkcjonał ten
przedłuża się ciągle z $C_c^1(\Omega)$ do $C_0(\Omega)$:
przestrzeń testów jest gęsta dla normy supremum,
co wynika z konstrukcji funkcji odcinających i wygładzania.
Twierdzenie Riesza--Markowa--Kakutaniego
(Rozdział~\ref{ch:riesz-markov-21})
daje miarę ze znakiem $D_i u$. Złóżmy je w miarę wektorową
$Du=(D_iu)_i$, uzyskując podaną równość.

Aby określić dokładną wariację, przyjmijmy
$\mu=\sum_{i=1}^n|D_i u|$. Jest to skończona dodatnia
miara. Twierdzenie Radona--Nikodyma daje
$D_i u=a_i\mu$, z funkcjami $a_i$ mierzalnymi.
Definiujemy miarę $|Du|=|a|\mu$.
Dla dowolnego pola testowego $|\varphi|\le1$ mamy
$\left|\int_\Omega\varphi\cdot\mathrm dDu\right|
\le\int_\Omega|a|\,\mathrm d\mu$.
Dla nierówności odwrotnej aproksymujemy
$a/|a|$ (przyjmując zero, gdy $a=0$) w $L^1(|Du|)$
przez pola $C_c^1$ z normą nieprzekraczającą $1$:
korzystamy z regularności $|Du|$ i przybliżania
funkcji mierzalnych ciągłymi na dużym zbiorze zwartym
oraz wygładzania i odcinania. Wtedy całki zbiegają
do $\int_\Omega|a|\,\mathrm d\mu$, co dowodzi równości.
Lokalna wersja wynika przez wykonywanie konstrukcji
na wyczerpaniu $\Omega$ i jednoznaczność miar na
przecięciach.
\end{proof}

\begin{twierdzenie}[Zupełność przestrzeni $BV$]
\label{thm:BV-banach-29}
Niech $\Omega\subseteq\mathbb R^n$ będzie otwarty.
Przestrzeń $BV(\Omega)$ z normą
$\|u\|_{BV}=\|u\|_{L^1(\Omega)}+|Du|(\Omega)$
jest przestrzenią Banacha.
\end{twierdzenie}
\begin{proof}
Niech $(u_j)$ będzie ciągiem Cauchy'ego w normie
$BV$. W szczególności jest to ciąg Cauchy'ego w
$L^1(\Omega)$, więc istnieje $u\in L^1(\Omega)$,
$u_j\to u$ w $L^1$. Z twierdzenia o reprezentacji
pochodnych mamy skończone miary wektorowe $Du_j$.
Ponieważ
$|D(u_j-u_\ell)|(\Omega)\le\|u_j-u_\ell\|_{BV}$,
ciąg tych miar jest Cauchy'ego w normie wariacji.
Przestrzeń skończonych miar wektorowych jest zupełna:
dla miar o wartościach w $\mathbb R^n$ wynika to z
zupełności przestrzeni skończonych miar ze znakiem
w każdej współrzędnej (Rozdział~\ref{ch:miary-ze-znakiem})
oraz nierówności
$\sum_i|\mu_i|(\Omega)/\sqrt n\le|\mu|(\Omega)
\le\sum_i|\mu_i|(\Omega)$.
Istnieje więc skończona miara wektorowa $\eta$
taka, że $|Du_j-\eta|(\Omega)\to0$.

Dla pola $\varphi\in C_c^1(\Omega;\mathbb R^n)$
przechodzimy do granicy w równości
$\int_\Omega u_j\operatorname{div}\varphi\,\mathrm d\lambda_n
=-\int_\Omega\varphi\cdot\mathrm dDu_j$.
Zbieżność $L^1$ po lewej i w wariacji po prawej daje
$\int_\Omega u\operatorname{div}\varphi\,\mathrm d\lambda_n
=-\int_\Omega\varphi\cdot\mathrm d\eta$, zatem $Du=\eta$.
W szczególności $u\in BV(\Omega)$, a
$\|u_j-u\|_{BV}=\|u_j-u\|_{L^1}+|Du_j-\eta|(\Omega)\to0$.
\end{proof}

\section{Zbiory o skończonym obwodzie}
\begin{definicja}[Obwód w sensie De Giorgiego]
\label{def:finite-perimeter-29}
Niech $E\subseteq\mathbb R^n$ będzie zbiorem mierzalnym
w sensie Lebesgue'a, a $\Omega\subseteq\mathbb R^n$ otwartym.
Definiujemy \emph{obwód $E$ w $\Omega$} przez
\[
 P(E;\Omega):=|D\chi_E|(\Omega)
 =\sup\left\{\int_E\operatorname{div}\varphi(x)
 \,\mathrm d\lambda_n(x):
 \varphi\in C_c^1(\Omega;\mathbb R^n),\quad|\varphi|\le1
 \right\}.
\]
Zbiór $E$ ma \emph{lokalnie skończony obwód},
gdy $P(E;V)<\infty$ dla wszystkich $V\Subset\mathbb R^n$.
Jeśli $P(E;\mathbb R^n)<\infty$, mówimy o skończonym
obwodzie globalnym. Nie wymagamy $\lambda_n(E)<\infty$:
półprzestrzeń ma lokalnie skończony obwód, ale
nieskończony obwód globalny przy $n\ge2$.
\end{definicja}

\begin{twierdzenie}[Granica klasycznej geometrii]
\label{thm:smooth-perimeter-29}
Niech $E\subseteq\mathbb R^n$ będzie ograniczonym
obszarem z brzegiem klasy $C^1$ i niech
$\nu_E$ będzie jednostkowym polem normalnych
zewnętrznych na $\partial E$. Wtedy dla każdego
otwartego $\Omega\subseteq\mathbb R^n$
\[
 D\chi_E=-\nu_E\,\mathcal H^{n-1}\llcorner\partial E,
 \qquad P(E;\Omega)=\mathcal H^{n-1}(\partial E\cap\Omega).
\]
Zapis $\mu\llcorner A$ oznacza ograniczenie miary
$\mu$ do zbioru $A$: $(\mu\llcorner A)(B)=\mu(B\cap A)$.
\end{twierdzenie}
\begin{proof}
Dla pola $\varphi\in C_c^1(\Omega;\mathbb R^n)$
przedłużonego przez zero poza $\Omega$ twierdzenie
o dywergencji dla brzegu $C^1$ (uzyskiwane lokalnie
przez zapis brzegu jako wykresu $C^1$, całkowanie
w kierunku normalnym i podział jedności) daje
$\int_E\operatorname{div}\varphi
\,\mathrm d\lambda_n
=\int_{\partial E}\varphi\cdot\nu_E
\,\mathrm d\mathcal H^{n-1}$.
Przez definicję $D\chi_E$ oznacza to
$D\chi_E=-\nu_E\mathcal H^{n-1}\llcorner\partial E$.
Ponieważ $|\nu_E|=1$, wariacja miary wektorowej
jest dokładnie $\mathcal H^{n-1}\llcorner\partial E$.
Nie liczymy w tym miejscu powierzchni gładkiej
od początku: jej miarę skonstruowano w
Rozdziale~\ref{ch:area-coarea-28}.
\end{proof}

\begin{przyklad}[Kostka ma skończony obwód]
\label{ex:cube-perimeter-29}
Niech $Q=(0,a)^n$, $a>0$. Całkując kolejno po
współrzędnych i stosując zasadnicze twierdzenie
rachunku całkowego, otrzymujemy
\[
 \int_Q\operatorname{div}\varphi(x)\,\mathrm d\lambda_n(x)
 =\sum_{i=1}^n\int_{(0,a)^{n-1}}
 [\varphi_i(x_1,\ldots,a,\ldots,x_n)-
  \varphi_i(x_1,\ldots,0,\ldots,x_n)]
 \,\mathrm d\lambda_{n-1}(x_{\widehat i}).
\]
Wyraz $x_{\widehat i}$ oznacza wszystkie współrzędne
poza $i$-tą. Tożsamość identyfikuje
$D\chi_Q$ z miarą powierzchniową $2n$ ścian,
pomnożoną przez odpowiednie normalne. Ich
przecięcia mają zerową $\mathcal H^{n-1}$-miarę,
więc $P(Q;\mathbb R^n)=2na^{n-1}$.
Nie wymagaliśmy gładkości brzegu w narożnikach.
\end{przyklad}

\begin{center}
\begin{tikzpicture}[scale=1.02,>=stealth]
  \fill[blue!8] (0,0) rectangle (3.5,2.4);
  \draw[blue!65!black,thick] (0,0) rectangle (3.5,2.4);
  \node at (1.75,1.2) {$E$};
  \foreach \x in {0.65,1.75,2.85} {
    \draw[->,very thick,red!75!black] (\x,2.4)--(\x,2.99);
    \draw[->,very thick,red!75!black] (\x,0)--(\x,-0.59);
  }
  \foreach \y in {0.6,1.2,1.8} {
    \draw[->,very thick,red!75!black] (0,\y)--(-0.6,\y);
    \draw[->,very thick,red!75!black] (3.5,\y)--(4.1,\y);
  }
  \node[anchor=west,font=\small] at (4.25,2.1) {$\nu_E$ --- normalna zewnętrzna};
  \node[anchor=west,font=\small] at (4.25,1.65) {$D\chi_E=-\nu_E\,\mathcal H^1|_{\partial E}$};
  \node[anchor=west,font=\small] at (4.25,1.2) {$P(E)=2(a+b)$ dla boków $a,b$};
  \node[anchor=west,font=\small] at (4.25,.75) {narożniki nie dają atomów};
\end{tikzpicture}
\end{center}
\noindent\emph{Rysunek:} miarowo-teoretyczna normalna i obwód
prostokąta w $\mathbb R^2$. Jego cztery wierzchołki
są zbiorami $\mathcal H^1$-miary zero i nie wnoszą
osobnej masy do obwodu.

\begin{lemat}[Pochodna i dopełnienie]
\label{lem:perimeter-complement-29}
Dla dowolnego zbioru mierzalnego $E$ oraz dowolnego
otwartego $\Omega$ mamy
$D\chi_{E^c}=-D\chi_E$ jako dystrybucje na $\Omega$
oraz $P(E^c;\Omega)=P(E;\Omega)$.
Jeżeli $\lambda_n(E\triangle F)=0$,
to $P(E;\Omega)=P(F;\Omega)$.
\end{lemat}
\begin{proof}
Ponieważ $\chi_{E^c}=1-\chi_E$ prawie wszędzie
i $\int_{\mathbb R^n}\operatorname{div}\varphi
\,\mathrm d\lambda_n=0$ dla testów zwartych,
w definicji wariacji zmienia się tylko znak.
Dla zbiorów równych prawie wszędzie ich funkcje
charakterystyczne są równe w $L^1_{\mathrm{loc}}$,
więc dają tę samą dystrybucję pochodnej.
\end{proof}

\section{Zwartość i półciągłość: dwa pełne dowody}
\begin{twierdzenie}[Dolna półciągłość obwodu]
\label{thm:perimeter-lsc-29}
Niech $\Omega\subseteq\mathbb R^n$ będzie otwarty.
Jeśli $\chi_{E_j}\to\chi_E$ w
$L^1_{\mathrm{loc}}(\Omega)$, to
\[
 P(E;\Omega)\le\liminf_{j\to\infty}P(E_j;\Omega).
\]
Analogicznie, jeśli $u_j\to u$ w
$L^1_{\mathrm{loc}}(\Omega)$, to
$|Du|(\Omega)\le\liminf_j|Du_j|(\Omega)$.
\end{twierdzenie}
\begin{proof}
Ustalmy pole $\varphi\in C_c^1(\Omega;\mathbb R^n)$,
$|\varphi|\le1$. Ponieważ $\operatorname{div}\varphi$
jest ograniczona i ma zwarty nośnik,
$\int_\Omega u_j\operatorname{div}\varphi
\,\mathrm d\lambda_n\to\int_\Omega u\operatorname{div}\varphi
\,\mathrm d\lambda_n$.
Każda z całek po lewej jest nie większa niż
$|Du_j|(\Omega)$. Zatem całka dla $u$
jest nie większa niż $\liminf_j|Du_j|(\Omega)$.
Bierzemy supremum po wszystkich testach,
uzyskując tezę. Przypadek zbiorów otrzymujemy
po podstawieniu $u_j=\chi_{E_j}$.
\end{proof}

\begin{lemat}[Oszacowanie przesunięcia]
\label{lem:BV-translation-29}
Jeżeli $u\in BV(\mathbb R^n)$, to dla każdego
$h\in\mathbb R^n$
\[
 \int_{\mathbb R^n}|u(x+h)-u(x)|
 \,\mathrm d\lambda_n(x)\le|h|\,|Du|(\mathbb R^n).
\]
W szczególności dla mierzalnego zbioru $E$ o
skończonej mierze i skończonym obwodzie zachodzi
$\lambda_n((E+h)\triangle E)\le|h|P(E;\mathbb R^n)$.
\end{lemat}
\begin{proof}
Wybierzmy nieujemny gładki mollifikator
$\rho\in C_c^\infty(\mathbb R^n)$ z
$\int_{\mathbb R^n}\rho\,\mathrm d\lambda_n=1$ i
$\rho_\varepsilon(x)=\varepsilon^{-n}\rho(x/\varepsilon)$.
Niech $u_\varepsilon=u*\rho_\varepsilon$.
Z twierdzenia o reprezentacji pochodnej jako miary
oraz Fubiniego wynika
$\nabla u_\varepsilon=(Du)*\rho_\varepsilon$.
Stąd
$\int_{\mathbb R^n}|\nabla u_\varepsilon|\,\mathrm d\lambda_n
\le|Du|(\mathbb R^n)$ na mocy nierówności trójkąta
pod całką i tego, że $\rho_\varepsilon$ ma masę $1$.
Dla gładkiego $u_\varepsilon$ twierdzenie o wartości
średniej na odcinku od $x$ do $x+h$ daje
$|u_\varepsilon(x+h)-u_\varepsilon(x)|
\le\int_0^1|\nabla u_\varepsilon(x+th)|\,|h|\,\mathrm dt$.
Całkując po $x$, stosując Tonellego i niezmienniczość
miary Lebesgue'a przez translację, otrzymujemy
oszacowanie z tezy dla $u_\varepsilon$.
Ponieważ $u_\varepsilon\to u$ w $L^1$,
możemy przejść do granicy w obu stronach nierówności.
Dla $u=\chi_E$ norma $L^1$ różnicy przesunięć
jest miarą różnicy symetrycznej.
\end{proof}

\begin{twierdzenie}[Zwartość rodzin zbiorów o ograniczonym obwodzie]
\label{thm:perimeter-compact-29}
Niech $Q\subset\mathbb R^n$ będzie ograniczonym
prostopadłościanem, a $E_j\subseteq Q$ zbiorami
mierzalnymi takimi, że
$\sup_j P(E_j;\mathbb R^n)<\infty$.
Wówczas istnieje podciąg $E_{j_\ell}$ oraz
mierzalny $E\subseteq\overline Q$, dla których
$\chi_{E_{j_\ell}}\to\chi_E$ w $L^1(\mathbb R^n)$,
a $P(E;\mathbb R^n)\le\liminf_\ell P(E_{j_\ell};\mathbb R^n)$.
\end{twierdzenie}
\begin{proof}
Niech $P_0=\sup_jP(E_j;\mathbb R^n)$.
Wybierzmy duży domknięty sześcian $S$, którego wnętrze
zawiera $\overline Q$. Dla $m\ge1$ podzielmy $S$
na skończenie wiele jednakowych półotwartych
kostek o boku $\delta_m\downarrow0$.
Każdy $u_j=\chi_{E_j}$ zastępujemy funkcją
$A_m u_j$, stałą na każdej kostce i równą
średniej $u_j$ na tej kostce, oraz równą $0$ poza $S$.
Współczynniki $A_m u_j$ należą do $[0,1]$.

Oszacujmy błąd przybliżenia. Dla kostki $C$
i jej średniej $u_C$ nierówność trójkąta i Fubini dają
\[
 \int_C|u_j(x)-u_C|\,\mathrm d\lambda_n(x)
 \le\frac1{\lambda_n(C)}\int_C\int_C
 |u_j(x)-u_j(y)|\,\mathrm d\lambda_n(y)
 \,\mathrm d\lambda_n(x).
\]
Po podstawieniu $y=x+h$ i rozszerzeniu dziedziny
całkowania w $x$ do całej $\mathbb R^n$ otrzymujemy
\[
 \sum_C\int_C|u_j-u_C|\,\mathrm d\lambda_n
 \le\delta_m^{-n}\int_{[-\delta_m,\delta_m]^n}
 \int_{\mathbb R^n}|u_j(x+h)-u_j(x)|
 \,\mathrm d\lambda_n(x)\,\mathrm d\lambda_n(h)
 \le 2^n\sqrt n\,\delta_m P_0.
\]
Ostatnia nierówność wynika z Lematu
\ref{lem:BV-translation-29}, ponieważ $|h|\le\sqrt n\,\delta_m$
na całym $[-\delta_m,\delta_m]^n$, którego objętość
wynosi $(2\delta_m)^n$. Zatem
$\sup_j\|u_j-A_mu_j\|_{L^1}\to0$.

Dla ustalonego $m$ rodzina $A_m u_j$ należy do
skończeniewymiarowego sześcianu $[0,1]^{N_m}$,
który jest zwarty. Wybierając kolejno podciągi dla
$m=1,2,\ldots$ i przechodząc do przekątnej,
uzyskujemy podciąg $u_{j_\ell}$, dla którego
$A_m u_{j_\ell}$ jest zbieżny w $L^1$
przy każdym ustalonym $m$. Jednostajna kontrola
błędu aproksymacji implikuje, że $u_{j_\ell}$
jest ciągiem Cauchy'ego w $L^1$. Zupełność
$L^1$ daje granicę $u$.

Wybieramy dalszy podciąg zbieżny prawie wszędzie.
Ponieważ wszystkie $u_j$ przyjmują tylko wartości
$0,1$, także $u$ ma te wartości prawie wszędzie,
czyli $u=\chi_E$ dla pewnego mierzalnego $E$.
Ponadto $u=0$ prawie wszędzie poza $\overline Q$.
Nierówność półciągłości z
Twierdzenia~\ref{thm:perimeter-lsc-29} kończy dowód.
\end{proof}

\section{Normalna miarowo-teoretyczna i brzeg zredukowany}
\begin{definicja}[Normalna z pochodnej miary]
\label{def:normal-BV-29}
Niech $E$ będzie zbiorem lokalnie skończonego obwodu.
Miara wektorowa $D\chi_E$ ma rozkład biegunowy
$D\chi_E=-\nu_E|D\chi_E|$, gdzie
$\nu_E:\mathbb R^n\to\mathbb R^n$ jest
$|D\chi_E|$-mierzalna i
$|\nu_E(x)|=1$ dla $|D\chi_E|$-prawie wszystkich $x$.
Funkcję $\nu_E$ nazywamy \emph{miarowo-teoretyczną
normalną zewnętrzną}.
\end{definicja}

\begin{proof}[Uzasadnienie istnienia normalnej]
Z Twierdzenia~\ref{thm:BV-measure-29} dla lokalnie
skończonej miary $\mu=\sum_i|D_i\chi_E|$
otrzymujemy pochodne Radona--Nikodyma
$D_i\chi_E=a_i\mu$. Miara
$|D\chi_E|=|a|\mu$, a $\nu_E=-a/|a|$ na
$\{|a|>0\}$, daje wymagany rozkład i normę $1$.
Zbiór, na którym $a=0$, ma zerową miarę wariacji.
\end{proof}

\begin{definicja}[Punkty gęstości i brzeg miarowy]
\label{def:measure-boundary-29}
Dla zbioru mierzalnego $E\subseteq\mathbb R^n$
niech $E^{(1)}$ i $E^{(0)}$ oznaczają zbiory punktów,
w których granica
\[
 \theta(E,x)=\lim_{r\downarrow0}
 \frac{\lambda_n(E\cap B(x,r))}{\lambda_n(B(x,r))}
\]
istnieje i wynosi odpowiednio $1$ lub $0$.
\emph{Brzegiem miarowym} nazywamy
$\partial^{\mathrm m}E=\mathbb R^n\setminus(E^{(0)}\cup E^{(1)})$.
Punkty, w których powyższa granica nie istnieje,
również należą do $\partial^{\mathrm m}E$.
\end{definicja}

\begin{lemat}[Brzeg miarowy jest objętościowo mały]
\label{lem:essential-boundary-null-29}
Dla każdego mierzalnego $E\subseteq\mathbb R^n$
zachodzi $\lambda_n(\partial^{\mathrm m}E)=0$.
\end{lemat}
\begin{proof}
Twierdzenie Lebesgue'a o różniczkowaniu stosujemy
do $\chi_E\in L^1_{\mathrm{loc}}(\mathbb R^n)$.
Dla prawie każdego $x$ średnie funkcji $\chi_E$
na kulach $B(x,r)$ zbiegają do $\chi_E(x)$,
a więc do $0$ lub $1$. Takie punkty należą
do $E^{(0)}\cup E^{(1)}$.
\end{proof}

\begin{definicja}[Brzeg zredukowany]
\label{def:reduced-boundary-29}
Niech $E$ będzie zbiorem lokalnie skończonego obwodu.
\emph{Brzeg zredukowany} $\partial^*E$ składa się
z tych $x\in\mathbb R^n$, dla których
$|D\chi_E|(B(x,r))>0$ dla każdego $r>0$
oraz istnieje granica
\[
 \nu_E^*(x)=-\lim_{r\downarrow0}
 \frac{D\chi_E(B(x,r))}{|D\chi_E|(B(x,r))}
 \quad\text{o normie }1.
\]
To definicja punktowa, nie tylko klasy funkcji
wyznaczonej prawie wszędzie.
\end{definicja}

\begin{twierdzenie}[Twierdzenie Gaussa--Greena w postaci miarowej]
\label{thm:gauss-green-BV-29}
Niech $E$ będzie zbiorem lokalnie skończonego
obwodu, a $\varphi\in C_c^1(\mathbb R^n;\mathbb R^n)$.
Wtedy
\begin{equation}\label{eq:gauss-green-BV-29}
 \int_E\operatorname{div}\varphi(x)
 \,\mathrm d\lambda_n(x)
 =\int_{\mathbb R^n}\varphi(x)\cdot\nu_E(x)
 \,\mathrm d|D\chi_E|(x).
\end{equation}
\end{twierdzenie}
\begin{proof}
Definicja dystrybucyjnej pochodnej daje
$\int_E\operatorname{div}\varphi\,\mathrm d\lambda_n=-\int_{\mathbb R^n}\varphi\cdot\mathrm dD\chi_E$.
Podstawiamy rozkład polarny z
Definicji~\ref{def:normal-BV-29} i otrzymujemy
prawą stronę. Nie trzeba zakładać gładkości brzegu.
\end{proof}

\subsection*{Narzędzia do dowodu twierdzenia De Giorgiego}
Główny trud polega na wyprowadzeniu z samej miary $D\chi_E$
zarówno lokalnej płaskości brzegu, jak i jego miary powierzchniowej.
Przedstawimy ten argument w kilku krokach. W tej części zakładamy
$n\ge2$; przypadek $n=1$ omówimy na końcu dowodu.
Przez $\omega_j$ oznaczamy objętość kuli jednostkowej
w $\mathbb R^j$.

\begin{lemat}[Nierówności Sobolewa i względna nierówność izoperymetryczna]
\label{lem:dg-relative-isoperimetric}
Dla każdego $n\ge2$ istnieje $C_n>0$ taka, że
\begin{enumerate}[label=(\roman*)]
\item jeśli $v\in C_c^1(\mathbb R^n)$, to
 $\|v\|_{L^{n/(n-1)}(\mathbb R^n)}\le
 C_n\int_{\mathbb R^n}|\nabla v(x)|\,\mathrm d\lambda_n(x)$;
\item dla każdej kuli $B=B(x,r)$ i $v\in BV(B)$ zachodzi
 $\|v-v_B\|_{L^{n/(n-1)}(B)}\le C_n|Dv|(B)$,
 gdzie $v_B=\lambda_n(B)^{-1}\int_Bv\,\mathrm d\lambda_n$;
\item dla mierzalnego $E$ o lokalnie skończonym obwodzie
 i dowolnej kuli $B=B(x,r)$ zachodzi
 \[
 \min\{\lambda_n(E\cap B),\lambda_n(B\setminus E)\}^{(n-1)/n}
 \le C_n P(E;B).
 \]
 Ponadto, jeżeli $A\subset\mathbb R^n$ jest mierzalny,
 ma skończoną miarę i skończony obwód, to
 $\lambda_n(A)^{(n-1)/n}\le C_nP(A;\mathbb R^n)$.
\end{enumerate}
\end{lemat}
\begin{proof}
Dla nieujemnych całkowalnych funkcji $g_i$ zależnych od
wszystkich współrzędnych poza $i$-tą zachodzi nierówność
\[
 \int_{\mathbb R^n}\prod_{i=1}^n
 g_i(x_{\widehat i})^{1/(n-1)}\,\mathrm d\lambda_n(x)
 \le\prod_{i=1}^n
 \left(\int_{\mathbb R^{n-1}}g_i(y)\,
 \mathrm d\lambda_{n-1}(y)\right)^{1/(n-1)}.
\]
Dla $n=2$ nierówność jest równością Fubiniego.
W kroku indukcyjnym oznaczamy
$h_i(x'_{\widehat i})=\int_{\mathbb R}
 g_i(x_{\widehat i})\,\mathrm d\lambda_1(x_n)$
przy $i<n$, gdzie $x'=(x_1,\ldots,x_{n-1})$.
Nierówność Höldera dla całkowania po $x_n$ daje
\[
 I\le\int_{\mathbb R^{n-1}}g_n(x')^{1/(n-1)}
 \prod_{i<n}h_i(x'_{\widehat i})^{1/(n-1)}
 \,\mathrm d\lambda_{n-1}(x').
\]
Stosując ponownie Höldera, z wykładnikami
$n-1$ oraz $(n-1)/(n-2)$, a następnie
założenie indukcyjne dla $(n-1)$ funkcji $h_i$,
otrzymujemy następujący rachunek. Oznaczmy
$G_i=\int_{\mathbb R^{n-1}}g_i(y)\,
\mathrm d\lambda_{n-1}(y)$ oraz
\[
 K=\int_{\mathbb R^{n-1}}\prod_{i<n}
 h_i(x'_{\widehat i})^{1/(n-2)}
 \,\mathrm d\lambda_{n-1}(x').
\]
Wtedy
\[
 I\le G_n^{1/(n-1)}K^{(n-2)/(n-1)}
 \le\prod_{i=1}^n G_i^{1/(n-1)}.
\]
To zamyka indukcję.
Dla $v\in C_c^1$ kładziemy
$g_i(x_{\widehat i})=\int_{\mathbb R}
|\partial_i v(x)|\,\mathrm d\lambda_1(x_i)$.
Z zasadniczego twierdzenia rachunku całkowego wynika
$|v(x)|\le g_i(x_{\widehat i})$ dla każdego $i$.
Mnożymy te nierówności i stosujemy poprzednie
oszacowanie, otrzymując
$\|v\|_{n/(n-1)}\le\prod_i\|\partial_i v\|_1^{1/n}
\le C_n\|\nabla v\|_1$. To dowodzi (i).

Dla (ii) wystarczy rozważyć kulę jednostkową,
ponieważ oba człony nierówności skalują się jak
$r^{n-1}$. Najpierw dla $v\in C^1(\overline B)$
pokażemy nierówność Poincarégo w $L^1$.
Ponieważ kula jest wypukła,
$|v(x)-v(y)|\le|x-y|\int_0^1
|\nabla v((1-t)x+ty)|\,\mathrm dt$.
Całkujemy po $x,y\in B$. Dla $t\le1/2$ zmiana
zmiennej $x\mapsto(1-t)x+ty$ ma jakobian odwrotny
nie większy niż $2^n$; dla $t\ge1/2$ wykonujemy
analogiczną zmianę względem $y$. Zatem
$\|v-v_B\|_1\le C_n\int_B|\nabla v(x)|\,\mathrm d\lambda_n(x)$.

Aby podnieść wykładnik, przedłużamy $v-v_B$
na nieco większą kulę, odbijając funkcję radialnie
w pierścieniu $1<|x|<3/2$ za pomocą
$(r,\theta)\mapsto(2-r,\theta)$, a następnie
mnożąc przez gładką funkcję odcinającą równą $1$
na $B$. Ponieważ odwzorowanie odbijające i jego
odwrotność mają ograniczone pochodne na tym pierścieniu,
dla otrzymanego przedłużenia $V\in W^{1,1}_c(\mathbb R^n)$
zachodzi
$\|\nabla V\|_1\le C_n(\|\nabla v\|_{L^1(B)}+
\|v-v_B\|_{L^1(B)})\le C_n\|\nabla v\|_{L^1(B)}$.
Stosujemy (i) do wygładzeń $V$ i przechodzimy do
granicy w $W^{1,1}$. Otrzymujemy (ii) dla funkcji
gładkich. Dla $v\in BV(B(x,r))$ ustalamy najpierw
$0<\rho<r$. Wygładzamy $v$ splotem o nośniku
promienia mniejszego niż $r-\rho$; na $B(x,\rho)$
całka z modułu gradientu wygładzenia jest
ograniczona przez $|Dv|(B(x,\rho+\varepsilon))$.
Stosujemy nierówność na $B(x,\rho)$, przechodzimy
$\varepsilon\downarrow0$ w $L^1$, korzystamy z
dolnej półciągłości normy $L^{n/(n-1)}$, a
następnie pozwalamy $\rho\uparrow r$. W ten
sposób otrzymujemy dokładnie $|Dv|(B(x,r))$,
bez przedłużania $v$ poza jej dziedzinę.

W (iii) bierzemy $v=\chi_E$, a w (ii)
$v_B=\theta:=\lambda_n(E\cap B)/\lambda_n(B)$.
Obliczenie normy $\chi_E-\theta$ daje
$\|\chi_E-\theta\|_{n/(n-1)}\ge
 c_n\min\{\lambda_n(E\cap B),\lambda_n(B\setminus E)\}^{(n-1)/n}$.
Dla ostatniej nierówności stosujemy (i) do
mollifikacji $\chi_A$ w całej przestrzeni.
Zbieżność w $L^1$ i oszacowanie całki z gradientu
przez $|D\chi_A|(\mathbb R^n)$ kończą dowód.
\end{proof}

\begin{lemat}[Przecięcie zbioru skończonego obwodu z kulą]
\label{lem:dg-ball-cut}
Niech $E\subseteq\mathbb R^n$ ma lokalnie skończony obwód,
$x\in\mathbb R^n$, $\mu=|D\chi_E|$.
Dla prawie wszystkich $r>0$ zachodzi równość miar
\[
 D\chi_{E\cap B(x,r)}
 =\chi_{B(x,r)}D\chi_E
 -\frac{y-x}{r}\,\chi_E(y)
 \mathcal H^{n-1}\llcorner\partial B(x,r).
\]
W ostatnim członie $\chi_E$ oznacza przedstawiciela
określonego dla prawie wszystkich promieni przez
polarne twierdzenie Fubiniego. Ponadto, przy
$v(r)=\lambda_n(E\cap B(x,r))$,
$w(r)=\lambda_n(B(x,r)\setminus E)$,
\[
 v'(r)=\mathcal H^{n-1}(E\cap\partial B(x,r)),\qquad
 w'(r)=\mathcal H^{n-1}(E^c\cap\partial B(x,r))
\]
dla prawie wszystkich $r>0$.
\end{lemat}
\begin{proof}
Dla pola testowego $\varphi$ i gładkiej
funkcji $\eta$ na $\mathbb R$ o wartościach
$0\le\eta\le1$, przechodzącej z $1$ do $0$,
wykorzystujemy tożsamość dystrybucyjną
$D(\chi_E\eta(|\cdot-x|))=
\eta(|\cdot-x|)D\chi_E+
\chi_E\eta'(|\cdot-x|)(y-x)|y-x|^{-1}\lambda_n$.
Otrzymujemy ją najpierw dla mollifikacji $\chi_E$,
a następnie przechodzimy do granicy w $L^1_{\rm loc}$
i słabo względem miar pochodnych. Podstawmy
$\eta((|y-x|-r)/\varepsilon)$ i pozwólmy
$\varepsilon\downarrow0$. Pierwszy składnik
zbiega do $\chi_{B(x,r)}D\chi_E$ dla wszystkich
promieni poza przeliczalnym zbiorem, na których
$\mu(\partial B(x,r))>0$. W drugim składniku
współrzędne polarne i twierdzenie Lebesgue'a
o różniczkowaniu jednowymiarowym względem $r$
dają granicę powierzchniową dla prawie wszystkich
promieni. To dowodzi wzoru na pochodną.
Ostatnie dwie równości wynikają bezpośrednio
z polarnego wzoru całkowego dla miary Lebesgue'a
i zasadniczego twierdzenia rachunku całkowego.
\end{proof}

\begin{lemat}[Oszacowania gęstości na brzegu zredukowanym]
\label{lem:dg-density}
Niech $E$ ma lokalnie skończony obwód,
$\mu=|D\chi_E|$, a $x\in\partial^*E$.
Istnieją liczby $r_x>0$ i $0<c_n<C_n<\infty$,
zależne wyłącznie od wymiaru poza $r_x$, takie że
\[
 c_nr^{n-1}\le\mu(B(x,r))\le C_nr^{n-1},\qquad
 \lambda_n(E\cap B(x,r))\ge c_nr^n,\quad
 \lambda_n(E^c\cap B(x,r))\ge c_nr^n
\]
dla każdego $0<r<r_x$.
\end{lemat}
\begin{proof}
Niech $\nu_0=\nu_E^*(x)$ oraz
$D\chi_E=-\nu_E\mu$. Z definicji brzegu
zredukowanego otrzymujemy
\[
 \frac{1}{\mu(B(x,r))}\int_{B(x,r)}\nu_E(y)
 \,\mathrm d\mu(y)\longrightarrow\nu_0.
\]
Ponieważ $|\nu_E|=|\nu_0|=1$,
\[
 \frac1{\mu(B(x,r))}\int_{B(x,r)}
 |\nu_E(y)-\nu_0|^2\,\mathrm d\mu(y)
 =2-2\nu_0\cdot\frac{\int_{B(x,r)}\nu_E(y)\,\mathrm d\mu(y)}
 {\mu(B(x,r))}\longrightarrow0.
\]
W szczególności, dla małych $r$ mamy
$\int_{B(x,r)}\nu_0\cdot\nu_E(y)\,\mathrm d\mu(y)
\ge\frac12\mu(B(x,r))$.

Całkowita masa każdej składowej miary
$D\chi_{E\cap B(x,r)}$ jest zerowa, ponieważ
$E\cap B(x,r)$ jest ograniczony i pole stałe
można testować funkcją odcinającą równą $1$
na jego otoczeniu. Z Lematu~\ref{lem:dg-ball-cut}
po pomnożeniu skalarnym przez $\nu_0$ otrzymujemy
dla prawie wszystkich małych $r$
$\mu(B(x,r))/2\le v'(r)$.
Stosując to samo rozumowanie do $E^c$ oraz
$D\chi_{E^c}=-D\chi_E$, otrzymujemy
$\mu(B(x,r))/2\le w'(r)$.
Ponieważ $v'(r),w'(r)\le n\omega_n r^{n-1}$,
daje to górne oszacowanie dla $\mu(B(x,r))$
(przez monotoniczność także dla wszystkich
małych $r$).

Dla prawie wszystkich małych $r$ nierówność
izoperymetryczna zastosowana do $E\cap B(x,r)$
i wzór na jego obwód dają
\[
 v(r)^{(n-1)/n}\le C_n
 \bigl(\mu(B(x,r))+v'(r)\bigr)
 \le C'_n v'(r).
\]
Analogicznie $w(r)^{(n-1)/n}\le C'_nw'(r)$.
Z warunku $\mu(B(x,r))>0$ wynika, że
$v(r),w(r)>0$ dla każdego $r>0$:
gdyby którakolwiek z tych objętości była zerowa,
pochodna dystrybucyjna $\chi_E$ zerowałaby się
w tej kuli. Dzieląc nierówności przez
$v^{(n-1)/n}$ i $w^{(n-1)/n}$ oraz całkując
od $\varepsilon$ do $r$, a następnie przechodząc
$\varepsilon\downarrow0$, uzyskujemy
$v(r),w(r)\ge c_nr^n$.
Względna nierówność izoperymetryczna z
Lematu~\ref{lem:dg-relative-isoperimetric}
daje teraz $\mu(B(x,r))\ge c_nr^{n-1}$.
\end{proof}

\begin{lemat}[Półprzestrzeń jako jedyna granica po przeskalowaniu]
\label{lem:dg-blowup}
Niech $E$, $\mu$ i $x\in\partial^*E$ będą jak
w Lemacie~\ref{lem:dg-density}, $\nu=\nu_E^*(x)$
i $E_{x,r}=\{z:x+rz\in E\}$. Wtedy
\[
 \chi_{E_{x,r}}\longrightarrow\chi_{H_\nu}
 \quad\text{w }L^1_{\rm loc}(\mathbb R^n),\qquad
 H_\nu=\{z:z\cdot\nu<0\}.
\]
Ponadto miary $\mu_{x,r}(A)=r^{1-n}\mu(x+rA)$
zbiegają słabo na zbiorach zwartych, gdy
$r\downarrow0$, do
$\mathcal H^{n-1}\llcorner\nu^\perp$, a więc
\[
 \lim_{r\downarrow0}
 \frac{\mu(B(x,r))}{\omega_{n-1}r^{n-1}}=1.
\]
\end{lemat}
\begin{proof}
Weźmy dowolny ciąg $r_j\downarrow0$.
Lemat~\ref{lem:dg-density} zapewnia jednostajne
oszacowania $\mu_{x,r_j}(B(0,R))\le C_nR^{n-1}$
dla każdego ustalonego $R$, gdy $j$ jest duże.
Dla funkcji $\chi_{E_{x,r_j}}$ daje to lokalne
ograniczenie norm $BV$. Wybieramy funkcję
odcinającą równą $1$ na $B(0,R)$ i o nośniku
w $B(0,R+1)$. Iloczyny z funkcjami
charakterystycznymi mają jednostajnie
ograniczone normy $L^1$ i wariacje.
Stosujemy do nich argument skończeniewymiarowych
średnich na kostkach z dowodu
Twierdzenia~\ref{thm:perimeter-compact-29}:
oszacowanie błędu używa wyłącznie wariacji,
nie dwuwartościowości funkcji. Przechodząc
przekątniowo względem $R\in\mathbb N$,
otrzymujemy podciąg zbieżny w $L^1_{\rm loc}$.
Wybrany dalej podciąg zbiega prawie wszędzie,
a zatem jego granica ma wartości w $\{0,1\}$
prawie wszędzie; jest nią $\chi_F$ dla pewnego
zbioru mierzalnego $F$.
Równocześnie, po wybraniu dalszego podciągu,
miary $\mu_{x,r_j}$ zbiegają słabo na zwartych
zbiorach do dodatniej miary Radona $\sigma$.

Z oszacowania w dowodzie Lematu~\ref{lem:dg-density}
oraz górnego ograniczenia $\mu(B(x,r))\le Cr^{n-1}$
wynika dla każdego $R>0$
\[
 \int_{B(0,R)}|\nu_E(x+r_jz)-\nu|\,
 \mathrm d\mu_{x,r_j}(z)\longrightarrow0.
\]
W języku dystrybucji otrzymujemy więc
$D\chi_F=-\nu\sigma$. Dla każdego $w\perp\nu$
pochodna kierunkowa $D_w\chi_F$ jest zerowa.
Z twierdzenia Fubiniego zastosowanego po obrocie
współrzędnych wynika stąd, że
$\chi_F(z)=h(z\cdot\nu)$ prawie wszędzie.
Ponieważ $D_\nu\chi_F=-\sigma$ jest miarą
niedodatnią, jednowymiarowa funkcja $h$ jest
nierosnąca w sensie dystrybucyjnym.
Wartości $h$ należą do $\{0,1\}$ prawie wszędzie,
więc $F$ jest albo półprzestrzenią
$\{z:z\cdot\nu<a\}$, albo zbiorem pustym,
albo całą przestrzenią, z dokładnością do miary zero.

Dolne oszacowanie z Lematu~\ref{lem:dg-density}
pokazuje, że $\sigma(\overline{B(0,1/2)})>0$:
dla zwartych kul stosujemy nierówność górną
w twierdzeniu o słabej zbieżności miar do
$\mu_{x,r_j}(\overline{B(0,1/2)})\ge c_n2^{1-n}$.
Zatem $F$ nie jest ani pusty, ani pełny.
Co więcej, w każdej kuli $B(0,R)$ granica $F$
i jej dopełnienie mają dodatnią objętość:
obydwie objętości dla $E_{x,r_j}$ są
nie mniejsze niż $c_nR^n$. Stąd $a=0$.
Otrzymujemy $F=H_\nu$.

Miara $D\chi_{H_\nu}$ wynosi
$-\nu\,\mathcal H^{n-1}\llcorner\nu^\perp$,
co wynika z jednowymiarowego całkowania
przez części i Fubiniego. Ponieważ
$D\chi_F=-\nu\sigma$, mamy
$\sigma=\mathcal H^{n-1}\llcorner\nu^\perp$.
Każdy podciąg dopuszcza dalszy podciąg o tej
samej granicy, więc granice dla $r\downarrow0$
są jednoznaczne. Słaba zbieżność miar i fakt,
że hiperpłaszczyzna $\nu^\perp$ nie nadaje masy
sferze $\partial B(0,1)$, dają na koniec
$\mu_{x,r}(B(0,1))\to\omega_{n-1}$.
\end{proof}

\begin{lemat}[Prostowalność brzegu zredukowanego]
\label{lem:dg-rectifiability}
Dla zbioru $E\subseteq\mathbb R^n$ o lokalnie
skończonym obwodzie zbiór $\partial^*E$ jest
przeliczalnie $(n-1)$-prostowalny.
\end{lemat}
\begin{proof}
Niech $S=\partial^*E$. Ustalmy $0<\alpha<1/4$
i skończony zbiór kierunków $e_1,\ldots,e_N$
taki, że każda jednostkowa normalna ma odległość
mniejszą niż $\alpha/8$ od któregoś $e_i$.
Z Lematu~\ref{lem:dg-blowup} każdy $x\in S$
spełnia, dla dostatecznie małych $r$,
\[
 \lambda_n\bigl((E\triangle(x+H_{\nu_E^*(x)}))
 \cap B(x,2r)\bigr)\le\delta r^n,
\]
gdzie $\delta>0$ wybierzemy wyłącznie w
zależności od $n$ i $\alpha$.
Podzielmy zatem $S$ na przeliczalne zbiory
$S_{i,k,m}$, na których $|x|\le k$,
$|\nu_E^*(x)-e_i|<\alpha/8$ oraz powyższe
oszacowanie zachodzi dla każdego $0<r<1/m$.
Nie pomijamy żadnego punktu $S$.

Pokażemy, że jeśli $x,y\in S_{i,k,m}$ i
$0<|x-y|<1/(2m)$, to
$|(y-x)\cdot e_i|<\alpha|y-x|$.
Gdyby nie, połóżmy $r=|x-y|$ i
$\rho=\alpha r/16$. Z bliskości normalnych
do $e_i$ wynika, że kula $B(y,\rho)$ leży
ściśle po jednej stronie hiperpłaszczyzny
$x+\nu_E^*(x)^\perp$. Z oszacowania w punkcie
$x$ zbiór $E$ zajmuje w tej kuli objętość
różniącą się od $0$ albo od pełnej objętości
o co najwyżej $\delta r^n$.
Z oszacowania w punkcie $y$ przy skali $\rho/2$
wynika natomiast, że w $B(y,\rho)$ funkcja
$\chi_E$ jest bliska charakterystycznej
półprzestrzeni przechodzącej przez $y$,
więc zajmuje objętość różniącą się od
$\lambda_n(B(y,\rho))/2$ o co najwyżej
$C_n\delta\rho^n$. Dla
$\delta$ mniejszego od stałej zależnej od
$n$ i $\alpha$ sprzeczność jest nieunikniona.

W konsekwencji rzut ortogonalny
$\pi_i:\mathbb R^n\to e_i^\perp$ jest
różnowartościowy na każdym fragmencie
$S_{i,k,m}$ o średnicy mniejszej niż $1/(2m)$.
Jego odwrotność ma postać
$z\mapsto z+h(z)e_i$ z funkcją $h$
Lipschitzowską o stałej co najwyżej
$\alpha/\sqrt{1-\alpha^2}$.
Rozbijamy $S_{i,k,m}$ na przeliczalnie wiele
takich fragmentów, a funkcję $h$ przedłużamy
na całe $e_i^\perp$ wzorem McShane'a
$\widetilde h(z)=\inf_{w\in\pi_i(S_{i,k,m})}
(h(w)+L|z-w|)$, stosowanym osobno na każdym
fragmencie. Każdy fragment leży więc na
wykresie Lipschitzowskim. Takich wykresów jest
przeliczalnie wiele, co dowodzi tezy.
\end{proof}

\begin{lemat}[Gęstość miary na wykresach prostowalnych]
\label{lem:dg-density-identification}
Niech $S\subseteq\mathbb R^n$ będzie przeliczalnie
$(n-1)$-prostowalny, $\mu$ dodatnią miarą Radona
skupioną na $S$, a dla każdego $x\in S$ niech
$\lim_{r\downarrow0}\mu(B(x,r))/
(\omega_{n-1}r^{n-1})=1$.
Wtedy $\mu=\mathcal H^{n-1}\llcorner S$.
\end{lemat}
\begin{proof}
Najpierw zauważmy, że $S$ jest zbiorem borelowskim:
warunki z Definicji~\ref{def:reduced-boundary-29}
można sprawdzać na promieniach wymiernych,
a funkcje $x\mapsto\mu(B(x,r))$ i
$x\mapsto D\chi_E(B(x,r))$ są borelowskie.
Z prostowalności wybieramy rozłączne zbiory
borelowskie $S_j\subseteq\Gamma_j$, gdzie $\Gamma_j$
jest pełnym wykresem funkcji Lipschitzowskiej,
takie że $S\setminus\bigcup_jS_j$ ma zerową
$\mathcal H^{n-1}$-miarę. Grafy pełne są
domknięte, więc taki rozkład uzyskujemy
przez kolejne odejmowanie wcześniejszych grafów.

Dla $M\in\mathbb N$ niech $S_M$ oznacza
zbiór punktów $x\in S$, dla których
\[
 \frac12\omega_{n-1}r^{n-1}
 \le\mu(B(x,r))\le2\omega_{n-1}r^{n-1}
 \qquad(0<r<1/M).
\]
Z założenia $S=\bigcup_MS_M$.
Na każdym $S_M$ oba oszacowania są jednostajne.
Z górnego i z definicji miary Hausdorffa wynika,
że $\mu(A)\le C_n\mathcal H^{n-1}(A)$
dla dowolnego borelowskiego $A\subseteq S_M$:
pokrywamy $A$ zbiorami małej średnicy,
wybieramy punkt $A$ w każdym niepustym
fragmencie i zastępujemy fragment kulą o
porównywalnej średnicy, po czym sumujemy
oszacowania na masy kul. W szczególności
$\mu(S\setminus\bigcup_jS_j)=0$.

Pokażemy również odwrotną absolutną ciągłość
na każdym $S_j$. Niech
$K\subseteq S_M\cap S_j$ będzie zwarty.
Dla $\delta<1/M$ pokrywamy $K$ kulami
o środkach w $K$ i promieniach mniejszych
od $\delta$. Lemat selekcji kul z
Rozdziału~\ref{ch:covering-24} pozwala wybrać
parami rozłączne kule $B(x_i,r_i)$, których
pięciokrotne powiększenia pokrywają $K$.
Definicja miary Hausdorffa i dolne oszacowanie
masy dają
\[
 \mathcal H^{n-1}_{10\delta}(K)
 \le C_n\sum_i r_i^{n-1}
 \le C'_n\sum_i\mu(B(x_i,r_i))
 \le C'_n\mu(K_\delta),
\]
gdzie $K_\delta=\{z:\operatorname{dist}(z,K)<\delta\}$.
Przez $\delta\downarrow0$ i skończoność
$\mu$ w pewnym otoczeniu zwartego $K$
otrzymujemy $\mathcal H^{n-1}(K)\le C'_n\mu(K)$.
Ponieważ $\mathcal H^{n-1}\llcorner\Gamma_j$
jest miarą Radona (wykres Lipschitzowski
jest lokalnie obrazem bi-Lipschitzowskim),
regularność przenosi to oszacowanie na borelowskie
podzbiory $S_M\cap S_j$.
Po zsumowaniu po $M$ stwierdzamy, że
$\mu\llcorner S_j\ll\mathcal H^{n-1}\llcorner S_j$
oraz $\mathcal H^{n-1}\llcorner S_j\ll\mu\llcorner S_j$.

Dla $\mathcal H^{n-1}$-prawie każdego $x\in S_j$
twierdzenie Rademachera, wzór area i
różniczkowanie funkcji charakterystycznej
$\chi_{S_j}$ po przeniesieniu na parametr
wykresu dają
\[
 \lim_{r\downarrow0}
 \frac{\mathcal H^{n-1}(S_j\cap B(x,r))}
 {\omega_{n-1}r^{n-1}}=1.
\]
Dla $\mu$-prawie każdego $x\in S_j$,
a zatem również dla $\mathcal H^{n-1}$-prawie
każdego takiego $x$ dzięki wzajemnej
absolutnej ciągłości, różniczkowanie miary
$\mu$ daje
$\mu(B(x,r)\cap S_j)/\mu(B(x,r))\to1$.
Łącząc te dwa rezultaty z założeniem o gęstości
$\mu$, uzyskujemy
\[
 \frac{\mu(B(x,r)\cap S_j)}
 {\mathcal H^{n-1}(S_j\cap B(x,r))}
 \longrightarrow1
 \quad\text{dla }\mathcal H^{n-1}\text{-prawie wszystkich }x\in S_j.
\]
Ponieważ na wykresie $S_j$ miary są lokalnie
skończone i wzajemnie absolutnie ciągłe,
twierdzenie o różniczkowaniu miar identyfikuje
ten iloraz z pochodną Radona--Nikodyma.
Otrzymujemy $\mu\llcorner S_j=
\mathcal H^{n-1}\llcorner S_j$.
Sumując po $j$ i uwzględniając zerowy zbiór
wyjątkowy, dostajemy tezę.
\end{proof}

\begin{twierdzenie}[Twierdzenie strukturalne De Giorgiego]
\label{thm:de-giorgi-structure-29}
Niech $E\subseteq\mathbb R^n$ będzie mierzalny
i lokalnie skończonego obwodu. Wtedy
$\partial^*E$ jest przeliczalnie $(n-1)$-prostowalny,
$|D\chi_E|=\mathcal H^{n-1}\llcorner\partial^*E$
oraz dla każdego $x\in\partial^*E$ zachodzi
\[
 \lim_{r\downarrow0}\int_{B(0,R)}
 \bigl|\chi_E(x+rz)-\chi_{H_{\nu_E^*(x)}}(z)\bigr|
 \,\mathrm d\lambda_n(z)=0\qquad(R>0),
\]
gdzie $H_\nu=\{z:z\cdot\nu<0\}$.
W szczególności, dla każdego pola
$\varphi\in C_c^1(\mathbb R^n;\mathbb R^n)$
\[
 \int_E\operatorname{div}\varphi(x)\,\mathrm d\lambda_n(x)
 =\int_{\partial^*E}\varphi(x)\cdot\nu_E^*(x)
 \,\mathrm d\mathcal H^{n-1}(x).
\]
\end{twierdzenie}
\begin{proof}
Dla $n\ge2$ przyjmijmy $\mu=|D\chi_E|$.
Z twierdzenia o różniczkowaniu miar zastosowanego
do funkcji polarnej $\nu_E\in L^1_{\rm loc}(\mu)$
wynika
$\mu(\mathbb R^n\setminus\partial^*E)=0$:
w prawie każdym punkcie względem $\mu$
średnie $\nu_E$ na małych kulach zbiegają do
wartości o normie $1$. Punkty poza nośnikiem
$\mu$ również nie spełniają definicji
brzegu zredukowanego.
Lemat~\ref{lem:dg-blowup} dowodzi granicy
po przeskalowaniu i istnienia dokładnej
gęstości $(n-1)$-wymiarowej równej $1$.
Lemat~\ref{lem:dg-rectifiability} daje
przeliczalną prostowalność, a
Lemat~\ref{lem:dg-density-identification}
identyfikuje miarę wariacji z miarą Hausdorffa.
Wzór całkowy wynika z
Twierdzenia~\ref{thm:gauss-green-BV-29}.

Jeśli $n=1$, funkcja $\chi_E$ ma na każdym
przedziale zwartym skończoną wariację
jednowymiarową, czyli jej przedstawiciel
ma skończenie wiele skoków o wielkości $1$.
Istotnie, każda zmiana wartości między
$0$ i $1$ wnosi co najmniej $1$ do wariacji,
więc nieskończenie wiele zmian na przedziale
przeczyłoby jej skończoności. Punkty skoków
tworzą lokalnie skończony zbiór, w każdym z nich
przeskalowane $E$ jest półprostą, a
$|D\chi_E|$ jest miarą liczącą na tych punktach,
czyli $\mathcal H^0\llcorner\partial^*E$.
\end{proof}

\begin{center}
\begin{tikzpicture}[>=stealth,scale=0.84]
  \begin{scope}[shift={(0,0)}]
    \clip (0,0) circle (1.35);
    \fill[blue!13] (-1.5,-1.5)--(-0.10,-1.5)--(-0.12,-1.3)--
      (0.08,-0.9)--(-0.10,-0.5)--(0.13,0)--(-0.08,0.45)--
      (0.09,0.9)--(-0.02,1.5)--(-1.5,1.5)--cycle;
    \draw[blue!75!black,very thick] plot[smooth] coordinates
      {(-0.12,-1.35) (0.08,-0.9) (-0.10,-0.5) (0.13,0)
       (-0.08,0.45) (0.09,0.9) (-0.02,1.35)};
  \end{scope}
  \draw (0,0) circle (1.35);
  \node at (0,-1.68) {$E_{x,r}\cap B(0,1)$};
  \draw[->,thick] (1.75,0)--(3.3,0)
    node[pos=.5,above] {$r\downarrow0$};
  \begin{scope}[shift={(5.15,0)}]
    \clip (0,0) circle (1.35);
    \fill[blue!13] (-1.5,-1.5) rectangle (0,1.5);
    \draw[blue!75!black,very thick] (0,-1.4)--(0,1.4);
    \draw[->,red!70!black,thick] (0,0)--(0.9,0)
      node[above,font=\scriptsize] {$\nu_E^*(x)$};
  \end{scope}
  \draw (5.15,0) circle (1.35);
  \node at (5.15,-1.68) {$H_{\nu_E^*(x)}\cap B(0,1)$};
\end{tikzpicture}

\small\emph{Schemat granicy De Giorgiego:}
po przeskalowaniu przy punkcie brzegu zredukowanego
zbiór staje się półprzestrzenią w sensie zbieżności
$L^1_{\mathrm{loc}}$; rysunek nie oznacza zbieżności
brzegów w odległości Hausdorffa.
\end{center}

\begin{uwaga}[Brzeg topologiczny a zredukowany]
Twierdzenie nie orzeka, że $\partial E=\partial^*E$.
Zmiana $E$ na zbiorze $\lambda_n$-zerowym może
zupełnie zmienić jego brzeg topologiczny,
lecz nie zmienia $D\chi_E$, obwodu ani wzoru
Gaussa--Greena. W szczególności właściwą
miarą powierzchniową dla ogólnych zbiorów
skończonego obwodu jest
$\mathcal H^{n-1}\llcorner\partial^*E$.
\end{uwaga}

\begin{lemat}[Mierzalność obwodów poziomic]
\label{lem:BV-level-meas-auditVI}
Niech $\Omega\subseteq\mathbb R^n$ będzie otwarty,
a $u:\Omega\to\mathbb R$ mierzalna. Wówczas funkcja
$t\mapsto P(\{u>t\};\Omega)$ o wartościach
w $[0,+\infty]$ jest borelowska. Analogicznie
$t\mapsto P(\{u<-t\};\Omega)$ jest borelowska.
Nie zakładamy skończoności obwodów.
\end{lemat}
\begin{proof}
Wybierzmy przeliczalne wyczerpanie $\Omega$
zbiorami zwartymi $K_j\subset\operatorname{int}K_{j+1}$.
Dla każdego $j$ istnieje przeliczalna rodzina
pól $C_c^1(\Omega;\mathbb R^n)$ o nośnikach
w ustalonym otoczeniu $K_j$, gęsta w normie $C^1$
wśród pól testowych wspartych w $K_j$.
Można ją skonstruować z mollifikowanych funkcji
schodkowych na wymiernych siatkach, z wymiernymi
współczynnikami; następnie skalujemy pola, aby
ich normy punktowe nie przekraczały $1$.
Po złączeniu tych rodzin otrzymujemy ciąg
$(\varphi_j)$, dla którego
\[
 P(E;\Omega)=\sup_j\int_E
 \operatorname{div}\varphi_j(x)\,\mathrm d\lambda_n(x)
\]
dla każdego zbioru mierzalnego $E$; każdą całkę
wolno tu zastąpić jej wartością z przeciwnym znakiem,
bo rodzina pól jest symetryczna. Uzasadnieniem
równości jest aproksymacja w normie $C^1$
na wspólnym zwartym nośniku i ograniczoność
$\lambda_n$ na tym nośniku.
Dla ustalonego $j$ odwzorowanie
$t\mapsto\int_\Omega\chi_{\{u>t\}}(x)
\operatorname{div}\varphi_j(x)\,\mathrm d\lambda_n(x)$
jest borelowskie: funkcja
$(t,x)\mapsto\chi_{\{u(x)>t\}}\operatorname{div}\varphi_j(x)$
jest mierzalna i jest zdominowana przez
całkowalną funkcję $|\operatorname{div}\varphi_j|$.
Co więcej, powstała funkcja $t$ jest borelowska:
rozcinamy $\operatorname{div}\varphi_j$ na części dodatnią
i ujemną; każda z odpowiadających całek jest
niemalejąca przy $t\downarrow-\infty$ (czyli
nierosnąca jako funkcja $t$), więc jest borelowska.
Supremum przeliczalnie wielu funkcji
borelowskich jest borelowskie. Analogiczny
argument działa dla $\{u<-t\}$.
\end{proof}

\section{Współobszar, funkcje BV i warstwowanie}
\begin{twierdzenie}[Jedna nierówność formuły koobszaru dla $BV$]
\label{thm:BV-layercake-29}
Niech $\Omega\subseteq\mathbb R^n$ będzie otwarty,
$u\in L^1(\Omega)$, a dla prawie wszystkich $t>0$
zbiory $E_t^+=\{u>t\}$ i $E_t^-=\{u<-t\}$
mają skończony obwód w $\Omega$. Jeśli
\[
 \int_0^\infty\!P(E_t^+;\Omega)\,\mathrm d\lambda_1(t)
 +\int_0^\infty\!P(E_t^-;\Omega)\,\mathrm d\lambda_1(t)<\infty,
\]
to $u\in BV(\Omega)$ i
\[
 |Du|(\Omega)\le\int_0^\infty
 [P(E_t^+;\Omega)+P(E_t^-;\Omega)]\,
 \mathrm d\lambda_1(t).
\]
\end{twierdzenie}
\begin{proof}
Warstwowanie funkcji daje punktowo
$u(x)=\int_0^\infty[\chi_{\{u>t\}}(x)
-\chi_{\{u<-t\}}(x)]\,\mathrm dt$
poza zbiorem, na którym $u$ nie jest skończona.
Równość ma sens w $L^1(\Omega)$, gdyż
$\int_0^\infty\lambda_n(\{|u|>t\})\,\mathrm dt
=\int_\Omega|u(x)|\,\mathrm d\lambda_n(x)<\infty$.
Dla testowego pola $\varphi$ o $|\varphi|\le1$
możemy zastosować Fubiniego i otrzymujemy
\[
 \int_\Omega u(x)\operatorname{div}\varphi(x)
 \,\mathrm d\lambda_n(x)
 =\int_0^\infty\!\left(
 \int_{E_t^+}\operatorname{div}\varphi(x)
 \,\mathrm d\lambda_n(x)
 -\int_{E_t^-}\operatorname{div}\varphi(x)
 \,\mathrm d\lambda_n(x)\right)\mathrm dt.
\]
Z definicji obwodu wartość wyrażenia w nawiasie
jest nie większa niż
$P(E_t^+;\Omega)+P(E_t^-;\Omega)$.
Biorąc supremum po testach, otrzymujemy tezę.
\end{proof}

\section{Pełna formuła koobszaru dla funkcji $BV$}
Związek między obwodem zbiorów poziomic a wariacją
funkcji nie wymaga twierdzenia strukturalnego De Giorgiego.
Wystarczą ścisła aproksymacja w $BV$, gładki wzór
koobszaru z Rozdziału~\ref{ch:area-coarea-28} i dolna
półciągłość wariacji.

\begin{lemat}[Ścisła aproksymacja gładka w $BV$]
\label{lem:BV-strict-approx}
Niech $\Omega\subseteq\mathbb R^n$ będzie otwarty,
a $u\in BV(\Omega)$. Istnieją funkcje
$u_j\in C^\infty(\Omega)\cap W^{1,1}(\Omega)$ takie, że
\[
 \|u_j-u\|_{L^1(\Omega)}\longrightarrow0,
 \qquad
 \int_\Omega|\nabla u_j(x)|\,\mathrm d\lambda_n(x)
 \longrightarrow |Du|(\Omega).
\]
Nie wymaga się regularności brzegu $\Omega$.
\end{lemat}
\begin{proof}
Wybieramy lokalnie skończoną gładką partycję
jedności $(\psi_i)_{i\ge1}$ na $\Omega$,
$\psi_i\ge0$, $\psi_i\in C_c^\infty(\Omega)$,
$\sum_i\psi_i=1$. Można ją skonstruować z
przeliczalnego pokrycia $\Omega$ kulami o zwartych
domknięciach w $\Omega$: lokalnie skończone
funkcje odcinające normalizujemy ich sumą.
Niech $\rho_\varepsilon$ będzie nieujemnym
gładkim mollifikatorem o masie $1$.
Dla ustalonego $j$ wybieramy dodatnie
$\varepsilon_{i,j}$ tak małe, aby nośniki
$\rho_{\varepsilon_{i,j}}*(\psi_i u)$ mieściły
się w ustalonych lokalnie skończenie wielu
otoczeniach $\operatorname{supp}\psi_i$,
a ponadto
\begin{align*}
 \|\rho_{\varepsilon_{i,j}}*(\psi_i u)-\psi_i u\|_1
 &<2^{-i-j},\\
 \|\rho_{\varepsilon_{i,j}}*(u\nabla\psi_i)
              -u\nabla\psi_i\|_1
 &<2^{-i-j}.
\end{align*}
Wszystkie normy dotyczą $\mathbb R^n$ po
przedłużeniu funkcji o zwartej podporze przez $0$.
Możliwość wyboru wynika z ciągłości mollifikacji
w $L^1$; obydwa wyrażenia są całkowalne, ponieważ
$u\in L^1(\Omega)$, a $\psi_i$ ma zwarty nośnik.
Definiujemy
$u_j=\sum_i\rho_{\varepsilon_{i,j}}*(\psi_i u)$.
Suma jest lokalnie skończona i należy do
$C^\infty(\Omega)\cap L^1(\Omega)$, a z pierwszego
oszacowania
$\|u_j-u\|_1\le2^{-j}$.

Reguła Leibniza w sensie miar daje
$D(\psi_i u)=\psi_i Du+u\nabla\psi_i\,\lambda_n$:
wystarczy w definicji pochodnej dystrybucyjnej
podstawić pole testowe pomnożone przez $\psi_i$
i rozwinąć klasyczną pochodną iloczynu.
Dlatego
\[
 \nabla u_j
 =\sum_i\rho_{\varepsilon_{i,j}}*(\psi_i Du)
 +\sum_i\bigl[\rho_{\varepsilon_{i,j}}*(u\nabla\psi_i)
                  -u\nabla\psi_i\bigr],
\]
gdyż $\sum_i\nabla\psi_i=0$ lokalnie,
a więc $\sum_i u\nabla\psi_i=0$ prawie wszędzie.
Dla pierwszej sumy stosujemy nierówność trójkąta,
Tonellego i nieujemność $\psi_i$:
\[
 \int_\Omega\left|
 \sum_i\rho_{\varepsilon_{i,j}}*(\psi_i Du)(x)
 \right|\,\mathrm d\lambda_n(x)
 \le\sum_i\int_\Omega\psi_i(x)\,\mathrm d|Du|(x)
 =|Du|(\Omega).
\]
Druga suma ma normę $L^1$ co najwyżej $2^{-j}$.
Zatem $\int_\Omega|\nabla u_j(x)|\,\mathrm d\lambda_n(x)
\le|Du|(\Omega)+2^{-j}$, co jednocześnie dowodzi,
że $u_j\in W^{1,1}(\Omega)$.
Nierówność odwrotna dla granicy wynika z
dolnej półciągłości wariacji z
Twierdzenia~\ref{thm:perimeter-lsc-29}:
$|Du|(\Omega)\le\liminf_j\int_\Omega|\nabla u_j|$.
\end{proof}

\begin{lemat}[Obwody poziomic funkcji gładkiej]
\label{lem:BV-smooth-levels}
Jeżeli $v\in C^1(\Omega)$ i
$\int_\Omega|\nabla v(x)|\,\mathrm d\lambda_n(x)<\infty$,
to dla prawie każdego $t\in\mathbb R$
zbiór $\{v>t\}$ ma skończony obwód oraz
\[
 \int_{-\infty}^{\infty}P(\{v>t\};\Omega)
 \,\mathrm d\lambda_1(t)
 \le\int_\Omega|\nabla v(x)|\,\mathrm d\lambda_n(x).
\]
\end{lemat}
\begin{proof}
Wybierzmy gładką, niemalejącą funkcję
$\theta:\mathbb R\to[0,1]$, równą $0$ na
$(-\infty,0]$ i $1$ na $[1,\infty)$, i niech
$\theta_\varepsilon(s)=\theta(s/\varepsilon)$.
Dla pola $\varphi\in C_c^1(\Omega;\mathbb R^n)$
całkowanie przez części daje
\[
 \int_\Omega\theta_\varepsilon(v(x)-t)
 \operatorname{div}\varphi(x)\,\mathrm d\lambda_n(x)
 =-\int_\Omega\theta'_\varepsilon(v(x)-t)
 \nabla v(x)\cdot\varphi(x)\,\mathrm d\lambda_n(x).
\]
Dla prawie każdego $t$ zbiór $\{v=t\}$ ma
zerową miarę Lebesgue'a (możliwe wyjątki
są co najwyżej przeliczalne na wyczerpaniu
$\Omega$ zwartymi podzbiorami). Lewa strona
zbiega więc do $\int_{\{v>t\}}\operatorname{div}\varphi$.
Z kolei wzór koobszaru dla odwzorowania
$v:\Omega\to\mathbb R$ pozwala zapisać
prawą stronę jako
\[
 -\int_{\mathbb R}\theta'_\varepsilon(s-t)
 \left(\int_{\{v=s,\,\nabla v\ne0\}}
 \frac{\nabla v(x)}{|\nabla v(x)|}\cdot\varphi(x)
 \,\mathrm d\mathcal H^{n-1}(x)\right)
 \,\mathrm d\lambda_1(s).
\]
Funkcja w nawiasie należy do $L^1(\mathbb R)$,
ponieważ jej moduł jest ograniczony przez
$\|\varphi\|_\infty\mathcal H^{n-1}
(\{v=s,\nabla v\ne0\}\cap\operatorname{supp}\varphi)$,
a całka tej ostatniej wielkości jest skończona
na mocy koobszaru i całkowalności $|\nabla v|$.
Twierdzenie Lebesgue'a o różniczkowaniu daje
zatem dla prawie wszystkich $t$
\[
 \int_{\{v>t\}}\operatorname{div}\varphi(x)
 \,\mathrm d\lambda_n(x)
 =-\int_{\{v=t,\,\nabla v\ne0\}}
 \frac{\nabla v(x)}{|\nabla v(x)|}\cdot\varphi(x)
 \,\mathrm d\mathcal H^{n-1}(x).
\]
Wybierając przeliczalny zbiór pól testowych
gęsty w normie $C^1$ na zwartych nośnikach
i przechodząc do granicy, uzyskujemy powyższą
równość dla wszystkich pól testowych na
wspólnym zbiorze poziomów pełnej miary.
W szczególności
$P(\{v>t\};\Omega)\le
\mathcal H^{n-1}(\{v=t,\nabla v\ne0\})$.
Całkujemy po $t$ i ponownie stosujemy koobszar,
otrzymując tezę.
\end{proof}

\begin{twierdzenie}[Pełna formuła koobszaru dla $BV$]
\label{thm:BV-coarea-full-29}
Niech $\Omega\subseteq\mathbb R^n$ będzie otwarty.
Jeżeli $u\in BV(\Omega)$, to dla prawie wszystkich
$t\in\mathbb R$ zbiór $\{x\in\Omega:u(x)>t\}$
ma skończony obwód w $\Omega$ oraz
\begin{equation}\label{eq:BV-coarea-full-29}
 |Du|(\Omega)=\int_{-\infty}^{\infty}
 P(\{u>t\};\Omega)\,\mathrm d\lambda_1(t).
\end{equation}
W szczególności dla każdego $u\in BV_{\mathrm{loc}}(\Omega)$
i otwartego $V\subseteq\Omega$ zachodzi równość,
być może z dwiema stronami równymi $+\infty$,
\[
 |Du|(V)=\int_{-\infty}^{\infty}
 P(\{u>t\};V)\,\mathrm d\lambda_1(t).
\]
\end{twierdzenie}
\begin{proof}
Najpierw niech $u\in BV(\Omega)$.
Wybierzmy z Lematu~\ref{lem:BV-strict-approx}
ciąg $u_j\in C^\infty(\Omega)\cap W^{1,1}(\Omega)$
zbiegający do $u$ w $L^1$, ze zbieżnością
całek z gradientu do $|Du|(\Omega)$.
Przechodząc do podciągu możemy założyć,
że $\sum_j\|u_j-u\|_{L^1}<\infty$.
Tożsamość liczbowa
$\int_{\mathbb R}|\chi_{\{a>t\}}-
\chi_{\{b>t\}}|\,\mathrm d\lambda_1(t)=|a-b|$
i twierdzenie Tonellego dają
\[
 \int_{\mathbb R}\|\chi_{\{u_j>t\}}-
 \chi_{\{u>t\}}\|_{L^1(\Omega)}
 \,\mathrm d\lambda_1(t)=\|u_j-u\|_{L^1(\Omega)}.
\]
Stąd dla prawie wszystkich $t$ zachodzi
$\chi_{\{u_j>t\}}\to\chi_{\{u>t\}}$
w $L^1(\Omega)$. Półciągłość obwodu,
Lemat~\ref{lem:BV-level-meas-auditVI},
nierówność Fatou i Lemat~\ref{lem:BV-smooth-levels}
dają
\begin{align*}
 \int_{\mathbb R}P(\{u>t\};\Omega)\,
 \mathrm d\lambda_1(t)
 &\le\liminf_j\int_{\mathbb R}
 P(\{u_j>t\};\Omega)\,\mathrm d\lambda_1(t)\\
 &\le\liminf_j\int_\Omega|\nabla u_j(x)|
 \,\mathrm d\lambda_n(x)=|Du|(\Omega).
\end{align*}
W szczególności prawie wszystkie poziomice
mają skończony obwód.

Nierówność w przeciwną stronę wynika z
Twierdzenia~\ref{thm:BV-layercake-29}.
Istotnie, dla $t>0$ używamy poziomic $\{u>t\}$,
a dla $t<0$ ich dopełnień, które są równe
$\{u<-|t|\}$ z dokładnością do zbioru
$\lambda_n$-zerowego, poza co najwyżej
przeliczalną liczbą poziomów $t$.
Wynika to z faktu, że na przestrzeni o mierze
$\sigma$-skończonej tylko przeliczalnie
wiele rozłącznych poziomic $\{u=t\}$ może
mieć dodatnią miarę. Obwód dopełnienia jest
taki sam, zatem
$\int_{\mathbb R}P(\{u>t\};\Omega)\,\mathrm d\lambda_1(t)$
można zapisać jako sumę obydwu całek
z Twierdzenia~\ref{thm:BV-layercake-29}.
To kończy dowód równości.

Dla $u\in BV_{\mathrm{loc}}(\Omega)$ i otwartego
$V\subseteq\Omega$ wybieramy wyczerpanie $V$
przez rosnące zbiory otwarte $V_j$ o zwartych
domknięciach w $\Omega$. Na każdym $V_j$
funkcja $u$ jest całkowalna i ma skończoną
wariację; stosujemy właśnie udowodnioną równość.
Ponieważ $|Du|(V_j)\uparrow|Du|(V)$
i $P(\{u>t\};V_j)\uparrow P(\{u>t\};V)$,
twierdzenie o zbieżności monotonicznej pozwala
przejść do granicy.
\end{proof}

\begin{wniosek}[Wzór koobszaru z wagą]
\label{cor:BV-coarea-weighted-29}
Niech $u\in BV_{\mathrm{loc}}(\Omega)$.
Dla każdej borelowskiej $g:\Omega\to[0,+\infty]$
zachodzi
\[
 \int_\Omega g(x)\,\mathrm d|Du|(x)
 =\int_{\mathbb R}\left(\int_\Omega g(x)
 \,\mathrm d|D\chi_{\{u>t\}}|(x)\right)
 \,\mathrm d\lambda_1(t).
\]
\end{wniosek}
\begin{proof}
Dla zbioru otwartego $V\subseteq\Omega$
równość obu stron, z $g=\chi_V$, to druga
część Twierdzenia~\ref{thm:BV-coarea-full-29}.
Zatem miara $|Du|$ i miara zadana prawą
stroną są równe na zbiorach otwartych.
Rzeczywiście, ta druga jest miarą: addytywność
wynika z twierdzenia Tonellego i addytywności
miar $|D\chi_{\{u>t\}}|$. Mierzalność względem
$t$ dla zbiorów otwartych wynika z
Lematu~\ref{lem:BV-level-meas-auditVI}.
Dla pozostałych zbiorów borelowskich ustalamy
wyczerpanie $\Omega$ zbiorami otwartymi o zwartych
domknięciach i ograniczamy się do takich zbiorów.
Dla prawie wszystkich $t$ obwód na każdym z nich
jest skończony. Dlatego twierdzenie o klasie
monotonicznej, zastosowane do zbiorów otwartych
i różnic takich zbiorów, daje mierzalność
$t\mapsto|D\chi_{\{u>t\}}|(A)$ dla każdego
zbioru borelowskiego $A$. Przez wyczerpanie
otrzymujemy to dla całej $\Omega$.
Obie miary są regularne i zgadzają się na
zbiorach otwartych, zatem są równe na
zbiorach borelowskich. Aproksymacja
nieujemnej funkcji borelowskiej funkcjami
prostymi kończy dowód.
\end{proof}

\begin{przyklad}[Funkcja ze skokiem i formuła koobszaru]
Niech $\Omega=(-2,2)$ i $u=\chi_{(0,1)}$.
Wtedy $Du=\delta_0-\delta_1$ jako miara
ze znakiem, zatem $|Du|(\Omega)=2$.
Dla $0<t<1$ mamy $\{u>t\}=(0,1)$ i
$P(\{u>t\};\Omega)=2$. Dla pozostałych
$t$, poza nieistotnymi poziomami $0$ i $1$,
obwód poziomicy jest zerowy: zbiór jest
pusty albo równy $\Omega$. Stąd
$\int_{\mathbb R}P(\{u>t\};\Omega)
\,\mathrm d\lambda_1(t)=\int_0^1 2\,
\mathrm d\lambda_1(t)=2$, zgodnie z
Twierdzeniem~\ref{thm:BV-coarea-full-29}.
\end{przyklad}

\begin{przyklad}[Gęsta zmiana reprezentanta nie zmienia brzegu zredukowanego]
Dla $E=(0,1)^n\setminus\mathbb Q^n$
zachodzi $\chi_E=\chi_{(0,1)^n}$ prawie wszędzie.
Stąd $|D\chi_E|$ jest miarą powierzchniową
na ścianach sześcianu, a jego brzeg zredukowany
zgadza się z sumą wnętrz ścian, z dokładnością
do zbioru $\mathcal H^{n-1}$-zerowego.
Natomiast brzeg topologiczny $E$ zawiera
cały domknięty sześcian, ponieważ zarówno
$E$, jak i jego dopełnienie są gęste
wewnątrz $[0,1]^n$. To pokazuje, że
brzeg topologiczny nie może zastępować
$\partial^*E$ w twierdzeniu De Giorgiego.
\end{przyklad}

\section{Elementy teorii prądów}
\begin{definicja}[Prąd i jego brzeg]
\label{def:current-29}
Niech $\Omega\subseteq\mathbb R^n$ będzie otwarty.
\emph{Prądem wymiaru $k$} nazywamy ciągły
funkcjonał liniowy na przestrzeni $\mathcal D^k(\Omega)$
gładkich $k$-form o zwartym nośniku, wyposażonej
w zwykłą topologię funkcji testowych.
Jeśli $T$ jest prądem wymiaru $k\ge1$, jego
\emph{brzeg} definiujemy przez
$\partial T(\omega)=T(\mathrm d\omega)$
dla $\omega\in\mathcal D^{k-1}(\Omega)$.
\end{definicja}

\begin{lemat}[Brzeg brzegu znika]
\label{lem:current-boundary-square-29}
Dla każdego prądu $T$ wymiaru $k\ge2$ zachodzi
$\partial(\partial T)=0$.
\end{lemat}
\begin{proof}
Dla $\omega\in\mathcal D^{k-2}(\Omega)$ mamy
$\partial(\partial T)(\omega)=
\partial T(\mathrm d\omega)=T(\mathrm d^2\omega)=0$,
gdyż z definicji różniczki zewnętrznej
$\mathrm d^2=0$.
\end{proof}

\begin{przyklad}[Prąd zadany przez powierzchnię]
Niech $S$ będzie zorientowaną zwartą gładką
$k$-wymiarową podrozmaitością z brzegiem
w $\Omega$. Wówczas $T_S(\omega)=\int_S\omega$
jest prądem wymiaru $k$. Z twierdzenia Stokesa
wynika $\partial T_S=T_{\partial S}$. Właśnie
w tej postaci klasyczny rachunek na rozmaitościach
staje się szczególnym przypadkiem teorii prądów.
\end{przyklad}

\section{Podsumowanie zależności i dalsze kierunki}
W rozdziale otrzymaliśmy reprezentację pochodnych
$BV$ przez miary, półciągłość i zwartość obwodu,
miarową formułę Gaussa--Greena, twierdzenie strukturalne
De Giorgiego wraz z geometrycznym opisem brzegu
zredukowanego oraz pełną formułę koobszaru dla $BV$.
Dowód twierdzenia strukturalnego opierał się na
nierównościach izoperymetrycznych, zwartości w $BV$,
granicach zbiorów po przeskalowaniu i prostowalności
wyprowadzonej z lokalnej płaskości.
Dowód formuły koobszaru był niezależny od
tej struktury: wykorzystał ścisłą aproksymację
gładką i wzór koobszaru dla odwzorowań Lipschitzowskich.
Dalsze zagadnienia to m.in. dokładne przedstawiciele
funkcji $BV$, twierdzenia śladowe, łańcuchowe
oraz teoria powierzchni minimalnych.

\begin{uwaga}[Miary wektorowe w dalszym ciągu skryptu]
Dla $u\in BV(\Omega)$ miara $Du$ ma skończoną wariację
i jest szczególnym przypadkiem miary o wartościach
w przestrzeni Banacha, omawianej w
Rozdziale~\ref{ch:vector-measures-30}.
W przypadku $u\in BV_{\mathrm{loc}}(\Omega)$ jej wariacja
jest tylko lokalnie skończona; nie musimy wtedy otrzymać
skończonego wektora $Du(\Omega)$.
Rozdział 30 interpretuje rozkład polarny poprzez RNP
przestrzeni $\mathbb R^n$. Nie należy jednak utożsamiać
tego rozkładu z istnieniem gęstości $Du$ względem
miary Lebesgue'a.
\end{uwaga}

% END chapters/pl/measure_theory/29-rozdzial.tex


% BEGIN chapters/pl/measure_theory/30-rozdzial.tex
% ROZDZIAL 30: miary wektorowe i wlasnosc Radona-Nikodyma
\chapter[Miary wektorowe i własność Radona--Nikodyma]{Miary wektorowe, całka Bochnera i własność Radona--Nikodyma}
\label{ch:vector-measures-30}

\section{Od miar skalarnych do miar w przestrzeni Banacha}
\label{sec:vector-intro-30}
W Rozdziale~\ref{ch:geometric-measure-29} pojawiła się
pochodna $Du$ o wartościach w $\mathbb R^n$,
a w Rozdziale~\ref{ch:area-coarea-28} różniczka
odwzorowania między przestrzeniami skończenie wymiarowymi.
W pierwszym przypadku dla $u\in BV(\Omega)$ mamy
skończoną wariację miary $Du$, natomiast dla
$u\in BV_{\mathrm{loc}}(\Omega)$ tylko wariację
lokalnie skończoną. Definicja miary wektorowej
poniżej wymaga, aby wartość na każdym zbiorze
mierzalnym należała do $X$; nie należy bez dodatkowych
założeń podstawiać do niej lokalnie skończonej
miary $Du$ na całym $\Omega$. Teraz pozwalamy, aby miara przyjmowała
wartości w \emph{dowolnej} rzeczywistej lub zespolonej przestrzeni
Banacha $X$. Szczególnie ważne okaże się pytanie, kiedy istnieje
funkcja $g:\Omega\to X$ taka, że $\nu(A)=\int_Ag\,\mathrm d\mu$.
W przypadku skalarnym odpowiedź daje twierdzenie Radona--Nikodyma;
w przypadku Banachowskim zależy ona od geometrii $X$.

\begin{definicja}[Miara wektorowa]
\label{def:vector-measure-30}
Niech $(\Omega,\mathcal A)$ będzie przestrzenią mierzalną,
a $X$ przestrzenią Banacha. \emph{Miarą wektorową o wartościach w $X$}
nazywamy odwzorowanie $\nu:\mathcal A\to X$, dla którego
$\nu(\varnothing)=0$ i dla każdej rodziny parami rozłącznych
$A_j\in\mathcal A$ zachodzi
\[
 \nu\left(\bigcup_{j=1}^{\infty}A_j\right)
 =\lim_{N\to\infty}\sum_{j=1}^{N}\nu(A_j)
 \quad\text{w normie }X.
\]
Przez $\|\nu(A)\|$ oznaczamy normę \emph{wektora} $\nu(A)$.
Nie należy jej mylić z wariacją miary na zbiorze $A$.
\end{definicja}

\begin{przyklad}[Dwie podstawowe konstrukcje]
\label{ex:basic-vector-measures-30}
Jeżeli $\mu$ jest dodatnią miarą skończoną, a $x_0\in X$,
to $\nu(A)=\mu(A)x_0$ jest miarą wektorową.
Jeśli $(\Omega,\mathcal A,\mu)$ jest przestrzenią o mierze
skończonej, a $g:\Omega\to X$ jest całkowalna w sensie
Bochnera (definicja poniżej), to
$\nu_g(A)=\int_Ag\,\mathrm d\mu$ jest miarą wektorową.
Zbieżność w normie sum po rozłącznych zbiorach
wynika z oszacowania
$\|\int_Bg\,\mathrm d\mu\|\le\int_B\|g\|\,\mathrm d\mu$
i bezwzględnej ciągłości całki skalarnej.
\end{przyklad}

\section{Wariacja całkowita i półwariacja}
\label{sec:vector-variation-30}
\begin{definicja}[Wariacja i półwariacja]
\label{def:vector-variation-30}
Niech $\nu:\mathcal A\to X$ będzie miarą wektorową.
Dla $A\in\mathcal A$ definiujemy \emph{wariację całkowitą}
\[
 |\nu|(A)=\sup\left\{\sum_{j=1}^{N}\|\nu(A_j)\|:
 A=\dot\bigcup_{j=1}^{N}A_j,\quad A_j\in\mathcal A,\ N\ge1\right\}.
\]
Jeżeli $|\nu|(\Omega)<\infty$, mówimy, że $\nu$ ma
\emph{skończoną wariację}. Dla $x^*\in X^*$ odwzorowanie
$A\mapsto x^*(\nu(A))$ jest miarą skalarną.
\emph{Półwariacją} nazywamy
$\|\nu\|_{\mathrm{sv}}(A)
=\sup_{\|x^*\|\le1}|x^*\nu|(A)$.
\end{definicja}

\begin{twierdzenie}[Wariacja jest miarą dodatnią]
\label{thm:vector-variation-measure-30}
Dla każdej miary wektorowej $\nu:\mathcal A\to X$
funkcja $|\nu|:\mathcal A\to[0,+\infty]$ jest miarą.
Ponadto dla każdego $A\in\mathcal A$ zachodzi
$\|\nu(A)\|\le\|\nu\|_{\mathrm{sv}}(A)\le|\nu|(A)$.
\end{twierdzenie}
\begin{proof}
Monotoniczność wariacji wynika z dopisania zbioru
$B\setminus A$ do każdego podziału $A\subseteq B$.
Niech $A,B$ będą rozłączne. Łącząc podziały $A$ i $B$,
otrzymujemy $|\nu|(A\cup B)\ge|\nu|(A)+|\nu|(B)$.
W drugą stronę dowolny skończony podział
$A\cup B=\dot\bigcup_i C_i$ rozcinamy przez $A$ i $B$.
Trójkąt daje
$\sum_i\|\nu(C_i)\|
\le\sum_i\|\nu(C_i\cap A)\|+\sum_i\|\nu(C_i\cap B)\|
\le|\nu|(A)+|\nu|(B)$.
Zatem wariacja jest skończenie addytywna, także wtedy,
gdy przyjmuje wartość nieskończoną.

Jeśli $A=\dot\bigcup_{j=1}^{\infty}A_j$,
to z monotoniczności i skończonej addytywności
$|\nu|(A)\ge\sum_j|\nu|(A_j)$.
Dla nierówności odwrotnej wybierzmy dowolny
skończony podział $A=\dot\bigcup_i C_i$.
Z przeliczalnej addytywności $\nu$ w normie wynika
\[
 \|\nu(C_i)\|\le\sum_j\|\nu(C_i\cap A_j)\|.
\]
Sumujemy po skończenie wielu $i$ i zamieniamy kolejność
sum nieujemnych, uzyskując
$\sum_i\|\nu(C_i)\|\le\sum_j|\nu|(A_j)$.
Supremum po podziałach daje nierówność odwrotną.
Dla dowolnego $x^*$ z $\|x^*\|\le1$ mamy
$|x^*\nu|(A)\le|\nu|(A)$. Supremum po $x^*$
daje drugą nierówność z tezy, zaś pierwsza wynika
z twierdzenia Hahna--Banacha:
$\|\nu(A)\|=\sup_{\|x^*\|\le1}|x^*(\nu(A))|$.
\end{proof}

\begin{przyklad}[Skończona półwariacja nie implikuje skończonej wariacji]
\label{ex:semivariation-not-variation-30}
Na $\Omega=\mathbb N$ z $\mathcal A=2^{\mathbb N}$
niech $X=\ell^2$ i $\nu(A)=\sum_{n\in A}e_n/n$,
gdzie $(e_n)$ jest bazą kanoniczną. Szereg po dowolnym
zbiorze zbiega bezwarunkowo w $\ell^2$, więc $\nu$
jest miarą wektorową. Dla podziału na singletony
$|\nu|(\mathbb N)\ge\sum_{n=1}^{N}1/n$; zatem
$|\nu|(\mathbb N)=+\infty$.
Natomiast dla $x^*=(a_n)\in\ell^2$ z $\|x^*\|\le1$
\[
 |x^*\nu|(\mathbb N)=\sum_{n=1}^{\infty}\frac{|a_n|}{n}
 \le\left(\sum_n|a_n|^2\right)^{1/2}
 \left(\sum_n\frac1{n^2}\right)^{1/2}.
\]
Równość w kresie górnym wynika z nierówności
Cauchy'ego--Schwarza i właściwego doboru $a_n$.
Półwariacja jest skończona, mimo że wariacja nie jest.
\end{przyklad}

\begin{lemat}[Bezwzględna ciągłość wariacji]
\label{lem:vector-variation-ac-30}
Jeśli $\nu$ ma skończoną wariację i $\nu\ll\mu$
dla dodatniej miary $\mu$, to $|\nu|\ll\mu$.
\end{lemat}
\begin{proof}
Jeśli $\mu(A)=0$, to $\nu(B)=0$ dla każdego
mierzalnego $B\subseteq A$. Wszystkie sumy definiujące
$|\nu|(A)$ są więc zerowe.
\end{proof}

\section{Całka Bochnera i miary z gęstością}
\label{sec:bochner-30}
\begin{definicja}[Silna mierzalność i całka Bochnera]
\label{def:bochner-30}
Niech $(\Omega,\mathcal A,\mu)$ będzie przestrzenią
z dodatnią miarą, a $X$ przestrzenią Banacha.
Funkcja $g:\Omega\to X$ jest \emph{silnie mierzalna},
jeżeli jest granicą punktową prawie wszędzie ciągu
funkcji prostych o wartościach w $X$.
Jest \emph{całkowalna w sensie Bochnera}, gdy dodatkowo
$\int_\Omega\|g(\omega)\|\,\mathrm d\mu(\omega)<\infty$.
Dla funkcji prostej $s=\sum_{j=1}^{N}x_j\chi_{A_j}$
o rozłącznych zbiorach $A_j$ o skończonej mierze
kładziemy $\int_\Omega s\,\mathrm d\mu=\sum_j\mu(A_j)x_j$.
Dla ogólnej $g$ całkę określa granica tych sum
w normie $X$ dla aproksymacji zbieżnej w
$L^1(\mu;X)$. Przestrzeń tych klas funkcji
oznaczamy $L^1(\mu;X)$.
\end{definicja}

\begin{twierdzenie}[Konstrukcja i oszacowanie całki Bochnera]
\label{thm:bochner-basic-30}
Przestrzeń $L^1(\mu;X)$ jest zupełna, a wzór
$g\mapsto\int_\Omega g\,\mathrm d\mu$ wyznacza
ciągły operator liniowy $L^1(\mu;X)\to X$
o normie co najwyżej $1$. Dla każdego $g\in L^1(\mu;X)$
zachodzi
\[
 \left\|\int_\Omega g\,\mathrm d\mu\right\|
 \le\int_\Omega\|g\|\,\mathrm d\mu.
\]
\end{twierdzenie}
\begin{proof}
Dla funkcji prostych nierówność jest zwykłą nierównością
trójkąta. Niech $g$ będzie silnie mierzalna i
$\int_\Omega\|g\|\,\mathrm d\mu<\infty$.
Z definicji silnej mierzalności wybierzmy proste
$s_n\to g$ prawie wszędzie. Możemy dobrać je tak, by
$\|s_n(\omega)\|\le2\|g(\omega)\|$ prawie wszędzie:
w punkcie, w którym pierwotna prosta jest zbyt duża,
zastępujemy jej wartość przez $0$. Otrzymany zbiór
wycięcia jest mierzalny, a po podziale na skończenie
wiele wartości prostej nadal otrzymujemy funkcję prostą.
Twierdzenie Lebesgue'a o zbieżności zmajoryzowanej
zastosowane do norm daje $\|s_n-g\|_{L^1}\to0$.
Wybór przybliżeń z nośnikami o skończonej mierze
wymaga jeszcze uzasadnienia, gdy $\mu(\Omega)=\infty$.
Dla $k\ge1$ zbiór $\{\|g\|>1/k\}$ ma skończoną
miarę z nierówności Markowa. Ich suma pokrywa
$\{\|g\|>0\}$, a z dominowanej zbieżności wynika,
że $g\chi_{\{\|g\|>1/k\}}\to g$ w $L^1$.
Na każdym z tych zbiorów o skończonej mierze
powyższe funkcje proste można obcinać bez zmiany
argumentu dominacji. Zatem funkcje proste
o całkowalnych normach są gęste także bez
założenia $\sigma$-skończoności całej miary $\mu$.

Dla ciągu prostych Cauchy'ego w $L^1$ wybieramy
podciąg $s_{n_k}$ taki, że
$\sum_k\|s_{n_{k+1}}-s_{n_k}\|_{L^1}<\infty$.
Twierdzenie Beppo Leviego daje
$\sum_k\|s_{n_{k+1}}(\omega)-s_{n_k}(\omega)\|<\infty$
prawie wszędzie. Zupełność $X$ zapewnia punktową
zbieżność do silnie mierzalnej $g$; nierówność Fatou
daje $s_{n_k}\to g$ w $L^1$, a w konsekwencji zbiega
cały ciąg. To dowodzi zupełności. Nierówność dla całki
prostej i gęstość jednoznacznie przedłużają całkę
i oszacowanie na całe $L^1(\mu;X)$.
\end{proof}

\begin{twierdzenie}[Wariacja miary o gęstości Bochnera]
\label{thm:variation-density-30}
Niech $g\in L^1(\mu;X)$, a
$\nu_g(A)=\int_Ag\,\mathrm d\mu$ dla $A\in\mathcal A$.
Wtedy $\nu_g$ jest miarą wektorową o skończonej wariacji
oraz
\[
 |\nu_g|(A)=\int_A\|g(\omega)\|\,\mathrm d\mu(\omega)
 \quad(A\in\mathcal A).
\]
\end{twierdzenie}
\begin{proof}
Przeliczalna addytywność wynika z
$\|\nu_g(\bigcup_{j>N}A_j)\|
\le\int_{\bigcup_{j>N}A_j}\|g\|\,\mathrm d\mu\to0$.
Dla każdego podziału $A=\dot\bigcup_jA_j$ mamy
$\sum_j\|\nu_g(A_j)\|\le\int_A\|g\|\,\mathrm d\mu$.
W drugą stronę dla funkcji prostej o rozłącznych
nośnikach $g=\sum_jx_j\chi_{B_j}$ podział $A$ przez
$B_j$ daje równość. Dla dowolnego $g$ i prostej $s$
wykorzystujemy oszacowanie
$\big||\nu_g|(A)-|\nu_s|(A)\big|
\le|\nu_{g-s}|(A)\le\int_A\|g-s\|\,\mathrm d\mu$.
Zbieżność $s\to g$ w $L^1$ kończy dowód.
\end{proof}

\begin{lemat}[Różniczkowanie całki Bochnera]
\label{lem:bochner-differentiation-30}
Jeśli $g\in L^1((a,b),\lambda_1;X)$, to
\[
 \lim_{h\to0}\left\|
 \frac1h\int_{(t,t+h)}g(s)\,\mathrm d\lambda_1(s)-g(t)
 \right\|=0
\]
dla prawie każdego $t\in(a,b)$; dla $h<0$ całka
po przedziale jest rozumiana ze znakiem, tzn.
$\int_{(t,t+h)}=-\int_{(t+h,t)}$.
\end{lemat}
\begin{proof}
Wystarczy wykazać, że prawie każdy punkt jest
punktem Lebesgue'a funkcji $X$-wartościowej, czyli
średnia z $\|g(s)-g(t)\|$ na małych symetrycznych
przedziałach dąży do zera. Funkcje ciągłe
$\phi:[a,b]\to X$ są gęste w $L^1((a,b);X)$:
najpierw aproksymujemy $g$ funkcjami prostymi,
a funkcje charakterystyczne zbiorów mierzalnych
aproksymujemy w $L^1$ funkcjami ciągłymi,
korzystając z regularności miary Lebesgue'a.
Niech $\phi_j$ zbiega do $g$ w $L^1$ tak szybko, aby
$\sum_j2^j\|g-\phi_j\|_{L^1}<\infty$.
Dla $a_j(t)=\|g(t)-\phi_j(t)\|$ mamy
$\sum_j2^j a_j(t)<\infty$ dla prawie wszystkich $t$.
Słaba nierówność maksymalna z Rozdziału
\ref{ch:differentiation-25} daje dla ustalonego
$\varepsilon>0$ oszacowanie
\[
 \lambda_1\!\left(\left\{t:
 \sup_{r>0}\frac1{2r}\int_{(t-r,t+r)}a_j(s)
 \,\mathrm d\lambda_1(s)>\varepsilon\right\}\right)
 \le C\varepsilon^{-1}\|g-\phi_j\|_{L^1}.
\]
Przez $j\to\infty$ odpowiednie zbiory wyjątkowe
mają miarę dążącą do zera. Jednostajna ciągłość
$\phi_j$ i nierówność trójkąta pokazują więc, że
$\limsup_{r\downarrow0}\frac1{2r}\int_{(t-r,t+r)}
\|g(s)-g(t)\|\,\mathrm d\lambda_1(s)\le2\varepsilon$
poza zbiorem o dowolnie małej mierze.
Przy $\varepsilon\downarrow0$ otrzymujemy tezę.
Średnią po przedziale jednostronnym ograniczamy
przez dwukrotność średniej symetrycznej.
\end{proof}

\section{Własność Radona--Nikodyma przestrzeni Banacha}
\label{sec:rnp-30}
\begin{definicja}[Własność Radona--Nikodyma]
\label{def:RNP-30}
Przestrzeń Banacha $X$ ma \emph{własność Radona--Nikodyma}
(w skrócie RNP), jeśli dla każdej przestrzeni z miarą
skończoną $(\Omega,\mathcal A,\mu)$ oraz każdej
$X$-wartościowej miary wektorowej $\nu$ o skończonej
wariacji, spełniającej $\nu\ll\mu$, istnieje
$g\in L^1(\mu;X)$ takie, że
\begin{equation}\label{eq:RNP-30}
 \nu(A)=\int_Ag(\omega)\,\mathrm d\mu(\omega)
 \qquad(A\in\mathcal A).
\end{equation}
Funkcja $g$ jest wyznaczona jednoznacznie
$\mu$-prawie wszędzie.
\end{definicja}
\begin{proof}[Jednoznaczność]
Niech $g,h$ mają takie same całki po wszystkich
$A\in\mathcal A$. Funkcja $k=g-h$ ma istotnie
separowalny zbiór wartości. Istnieje przeliczalna
rodzina funkcjonałów $x_j^*\in X^*$ o normach co
najwyżej $1$, która normuje tę separowalną przestrzeń:
$\|x\|=\sup_j|x_j^*(x)|$ dla jej elementów
(konstruujemy ją przez wybranie funkcjonałów
normujących punkty przeliczalnego zbioru gęstego,
a następnie używamy ciągłości).
Dla każdego $j$ skalarny Radon--Nikodym i
$\int_A x_j^*(k)\,\mathrm d\mu=0$ dla każdego $A$
dają $x_j^*(k)=0$ prawie wszędzie. Po usunięciu
przeliczalnej sumy zbiorów zerowych otrzymujemy
$\|k\|=0$. Stąd $g=h$ prawie wszędzie.
\end{proof}

\begin{przyklad}[Przestrzenie skończenie wymiarowe]
Każda przestrzeń Banacha $X$ o skończonym wymiarze
ma RNP. Wybierając bazę $e_1,\ldots,e_d$,
zapisujemy $\nu=\sum_j\nu_je_j$.
Twierdzenie Radona--Nikodyma dla miar skalarnych
daje $\nu_j=g_j\mu$. Funkcja
$g=\sum_jg_je_j$ jest silnie mierzalna i całkowalna,
gdyż $\|g\|\le C\sum_j|g_j|$ dla pewnej stałej $C$.
Wtedy $\nu=g\mu$. Ta argumentacja nie działa
bezpośrednio przy nieskończenie wielu współrzędnych.
\end{przyklad}

\begin{lemat}[Przybliżanie miary przez skończone podziały]
\label{lem:rnp-partition-30}
Niech $(\Omega,\mathcal A,\mu)$ będzie przestrzenią
o mierze skończonej, a $\nu:\mathcal A\to X$ miarą
wektorową spełniającą $\|\nu(A)\|\le\mu(A)$
dla każdego $A\in\mathcal A$.
Dla każdego skończonego podziału mierzalnego
$\pi=\{A_1,\ldots,A_N\}$ definiujemy
\[
 g_\pi=\sum_{j:\mu(A_j)>0}
 \frac{\nu(A_j)}{\mu(A_j)}\chi_{A_j}
 \quad\in L^1(\mu;X).
\]
Wówczas $\|g_\pi\|\le1$ prawie wszędzie,
a następujące warunki są równoważne:
\begin{enumerate}[label=(\roman*)]
 \item istnieje $g\in L^1(\mu;X)$ z
       $\nu(A)=\int_Ag\,\mathrm d\mu$ dla każdego $A$;
 \item sieć $(g_\pi)$, uporządkowana przez
       rozdrobnienie skończonych podziałów,
       jest siecią Cauchy'ego w $L^1(\mu;X)$.
\end{enumerate}
Jeśli $\rho$ rozdrabnia $\pi$, to
$g_\pi=\mathbb E_\pi g_\rho$, gdzie
$\mathbb E_\pi$ oznacza uśrednienie na komórkach
$\pi$. Operator $\mathbb E_\pi$ jest kontrakcją
na $L^1(\mu;X)$.
\end{lemat}
\begin{proof}
Na komórce $A_j$ uśrednienie funkcji $h$ jest równe
$\mu(A_j)^{-1}\int_{A_j}h\,\mathrm d\mu$
(na komórkach zerowych przyjmujemy $0$).
Z nierówności trójkąta dla całki Bochnera mamy
$\|\mathbb E_\pi h\|_{L^1}\le\|h\|_{L^1}$.
Jeżeli $\rho$ rozdrabnia $\pi$, to
$\int_{A_j}g_\rho\,\mathrm d\mu=\nu(A_j)$,
co daje równanie z tezy.

Załóżmy (i). Wtedy $g_\pi=\mathbb E_\pi g$.
Dla $\varepsilon>0$ wybierzmy funkcję prostą
$s$ z $\|g-s\|_{L^1}<\varepsilon$ i podział
$\pi_0$, na którego komórkach $s$ jest stała.
Dla każdego $\pi$ rozdrabniającego $\pi_0$
zachodzi $\mathbb E_\pi s=s$, skąd
\[
 \|g_\pi-g\|_{L^1}
 \le\|\mathbb E_\pi(g-s)\|_{L^1}+\|s-g\|_{L^1}
 <2\varepsilon.
\]
Zatem sieć jest zbieżna, a więc Cauchy'ego.

Jeśli zachodzi (ii), zupełność $L^1(\mu;X)$
daje granicę $g$. Ustalmy $A\in\mathcal A$.
Dla każdego podziału $\pi$ rozdrabniającego
$\{A,\Omega\setminus A\}$ mamy
$\int_Ag_\pi\,\mathrm d\mu=\nu(A)$.
Przechodząc do granicy w $L^1$ otrzymujemy
$\int_Ag\,\mathrm d\mu=\nu(A)$.
\end{proof}

\begin{lemat}[Całka po krzywej z pochodną Banachowską]
\label{lem:FTC-Banach-30}
Niech $F:[a,b]\to X$ będzie Lipschitzowska i
różniczkowalna w normie $X$ dla prawie wszystkich
$t\in(a,b)$. Wtedy jej pochodna $F'$ ma silnie
mierzalnego przedstawiciela, jest istotnie ograniczona,
a dla wszystkich $a\le s<t\le b$ zachodzi
\[
 F(t)-F(s)=\int_{(s,t)}F'(r)\,\mathrm d\lambda_1(r).
\]
\end{lemat}
\begin{proof}
W punktach różniczkowalności $F'$ jest granicą
ilorazów różnicowych z dodatnimi wymiernymi
przyrostami. Każdy iloraz jest funkcją ciągłą na
odpowiednim przedziale, o wartościach w separowalnej
podprzestrzeni rozpiętej przez $F([a,b])$.
Przedłużamy $F'$ przez zero na zbiorze wyjątkowym;
punktowa granica funkcji silnie mierzalnych jest
silnie mierzalna. Ponadto $\|F'\|\le\operatorname{Lip}(F)$.
Niech $x^*\in X^*$. Funkcja skalarna
$x^*\circ F$ jest Lipschitzowska, a prawie wszędzie
$(x^*\circ F)'=x^*(F')$. Jednowymiarowe
Twierdzenie~\ref{lem:rademacher-line-28} daje
\[
 x^*(F(t)-F(s))
 =\int_{(s,t)}x^*(F'(r))\,\mathrm d\lambda_1(r)
 =x^*\!\left(\int_{(s,t)}F'(r)\,\mathrm d\lambda_1(r)\right).
\]
Ponieważ zachodzi to dla każdego $x^*$,
Hahn--Banach dowodzi równości wektorów.
\end{proof}

\begin{lemat}[Dzielenie miary bezatomowej]
\label{lem:nonatomic-divisibility-30}
Niech $\mu$ będzie skończoną miarą bezatomową.
Dla każdego mierzalnego $A$ i każdej liczby
$0\le t\le\mu(A)$ istnieje mierzalny $B\subseteq A$
taki, że $\mu(B)=t$. W konsekwencji dla $\eta>0$
każdy mierzalny $A$ można podzielić na skończenie
wiele zbiorów miary co najwyżej $\eta$.
\end{lemat}
\begin{proof}
Najpierw zauważmy, że każdy zbiór dodatniej miary
zawiera podzbiór o dowolnie małej dodatniej mierze.
Istotnie, jeśli $E$ ma dodatnią miarę, z braku atomów
istnieje $D\subset E$ z $0<\mu(D)<\mu(E)$.
Mniejszy z $D$ i $E\setminus D$ ma dodatnią miarę
nie większą niż $\mu(E)/2$. Iterując otrzymujemy
zbiory dodatniej miary mniejszej niż zadane $\eta$.

Dla $0<t<\mu(A)$ rozważmy częściowo uporządkowaną
przez inkluzję modulo zbiory zerowe rodzinę
$\mathcal C=\{B\subseteq A:B\in\mathcal A,\ \mu(B)\le t\}$.
Każdy łańcuch ma kres górny: jeśli $s$ jest kresem
górnym miar zbiorów łańcucha, wybieramy rosnący
(w porządku modulo zbiory zerowe) ciąg $B_j$ o
$\mu(B_j)\uparrow s$, a jego suma $B=\bigcup_jB_j$
stanowi ograniczenie górne całego łańcucha modulo
zbiory zerowe. Rzeczywiście, każdy zbiór łańcucha
o mierze mniejszej od $s$ mieści się modulo zbiory
zerowe w pewnym $B_j$, a każdy zbiór o mierze $s$
jest modulo zbiory zerowe równy $B$.
Z lematu Zorna istnieje element maksymalny $B$.
Jeśli $\mu(B)<t$, to $\mu(A\setminus B)>0$;
wybierając w $A\setminus B$ zbiór $C$ o mierze
$0<\mu(C)<t-\mu(B)$, otrzymalibyśmy większy
element $B\cup C\in\mathcal C$, sprzeczność.
Stąd $\mu(B)=t$. Przypadki $t=0$ i $t=\mu(A)$
są bezpośrednie. Dzieląc kolejno $A$ na części
miary $\eta$ aż do uzyskania reszty o mierze
co najwyżej $\eta$, dostajemy skończony podział.
\end{proof}

\begin{twierdzenie}[RNP a różniczkowalność krzywych Lipschitzowskich]
\label{thm:rnp-lipschitz-equiv-30}
Dla rzeczywistej lub zespolonej przestrzeni Banacha $X$
następujące warunki są równoważne:
\begin{enumerate}[label=(\roman*)]
 \item $X$ ma własność Radona--Nikodyma z
       Definicji~\ref{def:RNP-30};
 \item każda funkcja Lipschitzowska
       $F:[0,1]\to X$ jest różniczkowalna w normie $X$
       dla $\lambda_1$-prawie każdego $t\in(0,1)$.
\end{enumerate}
Wtedy każda taka $F$ spełnia wzór
$F(t)=F(0)+\int_{(0,t)}F'(s)\,\mathrm d\lambda_1(s)$.
\end{twierdzenie}
\begin{proof}
\emph{Kierunek (i)$\Rightarrow$(ii).}
Niech $F:[0,1]\to X$ ma stałą Lipschitza $L$.
Na algebrze skończonych sum rozłącznych przedziałów
półotwartych definiujemy
$\nu((s,t])=F(t)-F(s)$. Definicja jest skończenie
addytywna, ponieważ przy dzieleniu przedziału przyrosty
się teleskopują. Z nierówności Lipschitza wynika
$\|\nu(A)\|\le L\lambda_1(A)$ dla takich zbiorów.
Każdy zbiór mierzalny $A\subseteq[0,1]$ można
przybliżyć w mierze Lebesgue'a skończonymi sumami
przedziałów półotwartych. Przybliżające wartości
$\nu(A_j)$ są ciągiem Cauchy'ego w $X$:
$\|\nu(A_j)-\nu(A_k)\|\le L\lambda_1(A_j\triangle A_k)$.
Ich granica nie zależy od aproksymacji i określa
rozszerzenie $\nu$ na $\sigma$-algebrę Lebesgue'a z
$\|\nu(A)\|\le L\lambda_1(A)$. Dla rozłącznych
$A_j$ to oszacowanie kontroluje normę
$\nu(\bigcup_{j>N}A_j)$, zatem rozszerzenie jest
przeliczalnie addytywne. Ma skończoną wariację:
$|\nu|(A)\le L\lambda_1(A)$.

Na mocy RNP istnieje $g\in L^1(\lambda_1;X)$ takie,
że $\nu(A)=\int_Ag\,\mathrm d\lambda_1$.
Dla $t\in[0,1]$ otrzymujemy
$F(t)=F(0)+\int_{(0,t)}g\,\mathrm d\lambda_1$.
Lemat~\ref{lem:bochner-differentiation-30} daje
$F'=g$ prawie wszędzie.

\emph{Kierunek (ii)$\Rightarrow$(i): redukcja do części bezatomowej.}
Niech $\nu:\mathcal A\to X$ ma skończoną wariację
oraz $\nu\ll\mu_0$ dla skończonej miary $\mu_0$.
Wybierzmy miarę $\mu=|\nu|$. Jest skończona i
spełnia $\|\nu(A)\|\le\mu(A)$.
Jeśli $\nu=g\mu$ dla pewnego $g\in L^1(\mu;X)$,
to na mocy Lematu~\ref{lem:vector-variation-ac-30}
mamy $\mu\ll\mu_0$. Skalarne twierdzenie
Radona--Nikodyma daje $w=\mathrm d\mu/\mathrm d\mu_0$;
wtedy $gw\in L^1(\mu_0;X)$ reprezentuje $\nu$
względem $\mu_0$. Wystarczy zatem zbudować
gęstość względem $\mu=|\nu|$.

Atomy dodatniej miary skończonej $\mu$ można
wybrać parami rozłącznie, a jest ich co najwyżej
przeliczalnie wiele, licząc z dokładnością do zbiorów
zerowych. Istotnie, dwa atomy o dodatnim przecięciu
są równe modulo zbiór zerowy. Dla każdego $k$
istnieje co najwyżej skończenie wiele rozłącznych
atomów miary większej niż $1/k$; ich sumy, po
przejściu po wszystkich $k$, obejmują część atomową.
Po odjęciu tej części pozostaje miara bezatomowa.
Na atomie $A$ funkcja
$\nu$ przyjmuje wartości proporcjonalne do $\mu$,
ponieważ każda mierzalna część atomu ma miarę $0$
lub $\mu(A)$; gęstość wynosi tam
$\nu(A)/\mu(A)$. Suma takich części daje funkcję
$g_{\mathrm{at}}\in L^1(\mu;X)$, gdyż
$\sum_A\|\nu(A)\|\le\sum_A\mu(A)\le\mu(\Omega)$.
Pozostaje część bezatomowa. Zakładamy więc dalej,
że $\mu$ jest bezatomowa i $\mu(\Omega)=M>0$.

\emph{Przypuszczenie o braku gęstości i wybór podziałów.}
Załóżmy przeciwnie, że $\nu$ nie ma gęstości
względem $\mu$. Z Lematu~\ref{lem:rnp-partition-30}
wynika, że sieć funkcji $g_\pi$ nie jest Cauchy'ego.
Istnieje $\varepsilon>0$ takie, że dla każdego
podziału $\pi_0$ znajdują się jego rozdrobnienia
$\alpha,\beta$ z
$\|g_\alpha-g_\beta\|_{L^1}\ge\varepsilon$.
Niech $\rho$ będzie wspólnym rozdrobnieniem
$\alpha$ i $\beta$. Ponieważ uśrednienia są
kontrakcjami i $\mathbb E_\alpha g_\rho=g_\alpha$,
$\mathbb E_\beta g_\rho=g_\beta$,
$\mathbb E_\alpha g_{\pi_0}=g_{\pi_0}$ oraz
$\mathbb E_\beta g_{\pi_0}=g_{\pi_0}$,
otrzymujemy
$\varepsilon\le2\|g_\rho-g_{\pi_0}\|_{L^1}$.
Na mocy Lematu~\ref{lem:nonatomic-divisibility-30}
każdy mierzalny zbiór można rozciąć na skończenie
wiele części o arbitralnie małej mierze. Jeśli dalej rozdrobnimy
$\rho$ do $\pi$, to ponownie przez kontrakcję
$\|g_\pi-g_{\pi_0}\|_{L^1}
\ge\|g_\rho-g_{\pi_0}\|_{L^1}$.
Możemy więc zbudować zagnieżdżone podziały
$\pi_1\prec\pi_2\prec\cdots$ takie, że
\begin{equation}\label{eq:rnp-partition-jumps-30}
 \max_{A\in\pi_n}\mu(A)\longrightarrow0,
 \qquad
 \|g_{\pi_{n+1}}-g_{\pi_n}\|_{L^1}
 \ge\varepsilon/2\quad(n\ge1).
\end{equation}
Zbiory o mierze zero pomijamy.

\emph{Kodowanie podziałów jako przedziałów.}
Każdej komórce $A\in\pi_1$ przypisujemy
półotwarty przedział $I_A\subseteq[0,M]$ o
długości $\mu(A)$, tak aby przedziały te
w kolejności wypełniały $[0,M]$.
Jeśli komórka $A$ jest podzielona na
$A_1,\ldots,A_r$ przez $\pi_2$, dzielimy
$I_A$ na kolejne podprzedziały długości
$\mu(A_1),\ldots,\mu(A_r)$. Powtarzamy to
rekurencyjnie. W ten sposób dostajemy zagnieżdżone
podziały przedziału, których maksymalna długość
komórki zmierza do zera. Niech $h_n:[0,M]\to X$
będzie funkcją stałą na $I_A$ dla $A\in\pi_n$
o wartości $\nu(A)/\mu(A)$. Z konstrukcji
\[
 \|h_n\|\le1\ \text{prawie wszędzie},\qquad
 \int_{I_A}h_n(t)\,\mathrm d\lambda_1(t)=\nu(A).
\]
Ponadto
\[
 \|h_{n+1}-h_n\|_{L^1([0,M];X)}
 =\|g_{\pi_{n+1}}-g_{\pi_n}\|_{L^1(\mu;X)}.
\]
Ostatnia równość wynika z porównania stałych
wartości na odpowiednich komórkach wspólnych
rozdrobnień, mających te same miary.

Niech $D$ będzie zbiorem końców wszystkich
przedziałów z wszystkich poziomów. Jest on
gęsty w $[0,M]$. Na $D$ definiujemy $F(0)=0$
i określamy $F(t)$ jako sumę $\nu(A)$ po
komórkach leżących na lewo od $t$, na dowolnym
poziomie, na którym $t$ jest końcem przedziału.
Definicja jest spójna dzięki skończonej addytywności
$\nu$; dla $s<t$ z $D$ mamy
$\|F(t)-F(s)\|\le\mu(\text{suma odpowiednich komórek})
=t-s$. Zatem $F$ przedłuża się jednoznacznie
do $1$-Lipschitzowskiej funkcji na $[0,M]$.
Dla każdej komórki $I_A=[c,d]$ zachodzi
\begin{equation}\label{eq:rnp-interval-average-30}
 F(d)-F(c)=\nu(A)
 =\int_{I_A}h_n(t)\,\mathrm d\lambda_1(t).
\end{equation}

\emph{Różniczkowanie wymusza sprzeczność.}
Warunek (ii), po liniowym przeskalowaniu $[0,M]$
do $[0,1]$, zapewnia istnienie $F'$ prawie wszędzie.
Lemat~\ref{lem:FTC-Banach-30} daje
$F(d)-F(c)=\int_{(c,d)}F'(t)\,\mathrm d\lambda_1(t)$.
Porównanie z \eqref{eq:rnp-interval-average-30}
pokazuje, że na każdej komórce $I_A$ funkcja $h_n$
jest średnią $F'$; inaczej mówiąc,
$h_n=\mathbb E_{\mathcal I_n}F'$,
gdzie $\mathcal I_n$ oznacza skończony podział
$[0,M]$ na opisane przedziały.
Ponieważ ich maksymalna długość dąży do zera,
$h_n\to F'$ w $L^1([0,M];X)$. Istotnie, dla ciągłej
$\phi:[0,M]\to X$ średnie po przedziałach różnią się
od $\phi$ o nie więcej niż jej moduł ciągłości
przy maksymalnej długości komórki. Dla ogólnej
$F'\in L^1$ przybliżamy ją funkcją ciągłą i
korzystamy z kontrakcji operatorów uśredniających.
Otrzymana zbieżność przeczy nierównościom
\eqref{eq:rnp-partition-jumps-30}. Zatem $\nu$
ma gęstość względem $|\nu|$, a stąd, jak wykazano
na początku, względem $\mu_0$.
\end{proof}

\section{Przykłady: przestrzenie z RNP i bez RNP}
\label{sec:rnp-examples-30}
\begin{twierdzenie}[Przestrzenie Hilberta mają RNP]
\label{thm:hilbert-rnp-30}
Każda rzeczywista lub zespolona przestrzeń Hilberta $H$
ma własność Radona--Nikodyma.
\end{twierdzenie}
\begin{proof}
Niech $\nu:\mathcal A\to H$ ma skończoną wariację
i niech $\mu=|\nu|$. Dla podziału $\pi$ funkcja
$g_\pi$ z Lematu~\ref{lem:rnp-partition-30} spełnia
$\|g_\pi\|_{L^\infty}\le1$, zatem jest ograniczona
w przestrzeni Hilberta $L^2(\mu;H)$:
$\|g_\pi\|_{L^2}^2\le\mu(\Omega)$.
W domkniętej kuli tej przestrzeni istnieje
słabo zbieżna podsieć $(g_{\pi_\alpha})$,
gdyż kula przestrzeni Hilberta jest słabo zwarta.
Niech $g\in L^2(\mu;H)$ będzie jej granicą.
Dla każdego $A\in\mathcal A$ i $h\in H$ funkcja
$\chi_A h$ należy do $L^2(\mu;H)$. Jeśli
$\pi_\alpha$ rozdrabnia $\{A,\Omega\setminus A\}$,
mamy $\int_Ag_{\pi_\alpha}\,\mathrm d\mu=\nu(A)$.
Zatem słaba zbieżność w $L^2$ daje
\[
 \left\langle\int_Ag\,\mathrm d\mu,h\right\rangle_H
 =\lim_\alpha\left\langle
 \int_Ag_{\pi_\alpha}\,\mathrm d\mu,h\right\rangle_H
 =\langle\nu(A),h\rangle_H.
\]
Dowolność $h$ dowodzi $\nu(A)=\int_Ag\,\mathrm d\mu$.
Ponieważ $\mu$ jest skończona, $L^2(\mu;H)$
zanurza się w $L^1(\mu;H)$.
Dla pierwotnej miary dominującej $\mu_0$
przechodzimy z $\mu=|\nu|$ do $\mu_0$ dokładnie jak
w dowodzie Twierdzenia~\ref{thm:rnp-lipschitz-equiv-30}.
\end{proof}

\begin{przyklad}[Dlaczego przestrzeń $L^1$ nie ma RNP]
\label{ex:L1-no-RNP-30}
Niech $X=L^1([0,1],\lambda_1)$ oraz
$F:[0,1]\to X$ będzie określona przez
$F(t)=\chi_{[0,t]}$ (jako klasa funkcji w zmiennej $s$).
Dla $t>u$ mamy $\|F(t)-F(u)\|_{L^1}=t-u$,
a więc $F$ jest $1$-Lipschitzowska.
W żadnym punkcie $t\in(0,1)$ nie istnieje jednak
pochodna w normie $L^1$. Istotnie, dla $h>0$
iloraz różnicowy to
$h^{-1}\chi_{(t,t+h]}$, którego norma wynosi $1$.
Gdyby istniała granica $v\in L^1$, to zbieżność
w $L^1$ pewnego podciągu dawałaby zbieżność
prawie wszędzie do $v$. Dla każdego $s\ne t$
ilorazy różnicowe są jednak ostatecznie równe $0$.
Zatem $v=0$ prawie wszędzie, sprzecznie z
$\|h^{-1}\chi_{(t,t+h]}\|_{L^1}=1$.
Z Twierdzenia~\ref{thm:rnp-lipschitz-equiv-30}
wynika, że $L^1([0,1])$ nie ma RNP.

Równoważnie, miara wektorowa
$\nu(A)=\chi_A\in L^1([0,1])$ spełnia
$\|\nu(A)\|_{L^1}=\lambda_1(A)$ i
$|\nu|=\lambda_1$, ale \emph{nie} ma gęstości
Bochnera $g$ względem $\lambda_1$.
Gdyby miała, jej funkcja pierwotna $F(t)=\nu([0,t])$
musiałaby być różniczkowalna prawie wszędzie.
\end{przyklad}

\begin{przyklad}[Przestrzeń funkcji ciągłych bez RNP]
Niech $X=C([0,1])$ z normą supremum i
$F(t)(s)=|s-t|$ dla $s,t\in[0,1]$.
Wtedy $\|F(t)-F(u)\|_\infty\le|t-u|$.
Dla ustalonego $t\in(0,1)$ skalarne ilorazy różnicowe
przy $h\downarrow0$ dążą punktowo (dla $s\ne t$)
do funkcji równej $+1$ dla $s<t$ oraz $-1$ dla $s>t$.
Ta funkcja ma skok, więc nie może być jednostajną
granicą funkcji ciągłych na $[0,1]$.
Nie istnieje zatem pochodna $F'(t)$ w normie supremum,
a $C([0,1])$ nie ma RNP.
\end{przyklad}

\begin{uwaga}[Nie należy mylić dwóch twierdzeń Rademachera]
Twierdzenie~\ref{thm:rademacher-28} dowodzi
różniczkowalności odwzorowań Lipschitzowskich
$U\subseteq\mathbb R^n\to\mathbb R^m$ prawie
wszędzie, bez dodatkowych założeń. Dla ogólnej
przestrzeni Banacha $X$ nawet odwzorowanie
$[0,1]\to X$ może być nigdzie nieróżniczkowalne.
Twierdzenie~\ref{thm:rnp-lipschitz-equiv-30}
wykazuje, że dokładną przeszkodą jest brak RNP.
W szczególności nie wolno wnioskować o RNP
z samego istnienia skalarnej pochodnej
Radona--Nikodyma dla $x^*\nu$ przy każdym
$x^*\in X^*$: te skalarne gęstości nie muszą
składać się w silnie mierzalną gęstość $X$-wartościową.
\end{uwaga}

\section{Miary wektorowe w analizie i geometrii}
\label{sec:vector-applications-30}
\begin{przyklad}[Miara gradientowa funkcji $BV$]
Jeśli $u\in BV(\Omega)$, to z
Twierdzenia~\ref{thm:BV-measure-29} dystrybucyjna
pochodna $Du=(D_1u,\ldots,D_nu)$ jest skończoną
miarą wektorową o wartościach w $\mathbb R^n$.
Dla $u\in BV_{\mathrm{loc}}(\Omega)$ stosujemy
poniższe wnioski o gęstościach lokalnie, po ograniczeniu
do części o skończonej wariacji.
Z RNP tej przestrzeni i twierdzenia
Radona--Nikodyma otrzymujemy gęstość
$\sigma_u=\mathrm dDu/\mathrm d|Du|$ o normie $1$
$|Du|$-prawie wszędzie. Dla funkcji gładkiej
$Du=\nabla u\,\lambda_n$, natomiast dla funkcji
charakterystycznej zbioru o skończonym obwodzie
miara $D\chi_E$ może być osobliwa względem
$\lambda_n$ i skoncentrowana na brzegu miarowym.
To wyjaśnia, dlaczego naturalną miarą dominującą
jest $|Du|$, a nie zawsze $\lambda_n$.
\end{przyklad}

\begin{wniosek}[Przestrzenie Banacha i całki z krzywych]\label{idx:30:wniosek:15}

Jeśli $X$ ma RNP i $F:[a,b]\to X$ jest
Lipschitzowska, to istnieje $g\in L^\infty((a,b);X)$
takie, że $F(t)-F(a)=\int_{(a,t)}g\,\mathrm d\lambda_1$
i $F'=g$ prawie wszędzie.
\end{wniosek}
\begin{proof}
Po przeskalowaniu przedziału stosujemy
Twierdzenie~\ref{thm:rnp-lipschitz-equiv-30} i
Lemat~\ref{lem:FTC-Banach-30}.
\end{proof}

\section{Granice zakresu i dalsze kierunki}
Teoria miar wektorowych ma również zagadnienia,
których nie potrzebowaliśmy do dowodu RNP:
całkę Pettisa, miary o skończonej półwariacji,
operatory słabo zwarte oraz twierdzenie
Bartle'a--Dunforda--Schwartza.
Są one naturalną kontynuacją w analizie funkcjonalnej.
W tym rozdziale udowodniliśmy od podstaw
charakteryzację RNP przez różniczkowalność krzywych
Lipschitzowskich. Nie zakładaliśmy ani twierdzenia
o dentowalności zbiorów wypukłych, ani ogólnego
twierdzenia o zbieżności martyngałów Banachowskich.

% END chapters/pl/measure_theory/30-rozdzial.tex


% BEGIN chapters/pl/measure_theory/31-dzety-geometria-i-teoria-miary.tex
% Rozdział współdzielony tematycznie z geometrią różniczkową; wersja audytowana 2026-11-19.
% Samodzielny rozdział do włączenia przez \input{rozdzial_dzety.tex}.
% Wymaga: amsmath, amssymb, amsthm, mathtools, tikz, hyperref.
% Środowiska: definicja, twierdzenie, lemat, wniosek, uwaga, przyklad, proof.
\chapter{Dżety: od geometrii różniczkowej do teorii miary}
\label{ch:dzety}

\section{Motywacja: ile informacji niesie punkt?}
W rozdziale~\ref{ch:area-coarea-28} wyprowadziliśmy twierdzenie
Rademachera i wzory miarowo-geometryczne. Teraz interpretujemy
różniczkowalność prawie wszędzie językiem dżetów, zachowując też
pełne ujęcie geometryczne, formy kontaktowe i twierdzenia Whitneya.
Wartość $f(p)$ opisuje jedynie punkt obrazu, pochodna $Df(p)$ ---
przybliżenie liniowe, a kolejne pochodne --- coraz dokładniejsze
przybliżenia lokalne. \emph{Dżet rzędu $k$} jest sposobem zebrania
tej informacji w jeden obiekt. Rozróżniamy trzy sytuacje:
\begin{enumerate}
\item w geometrii gładkiej dżet jest klasą równoważności
kiełków odwzorowań lub przekrojów;
\item w równaniach różniczkowych współrzędne dżetowe
traktuje się jako formalnie niezależne, lecz rozwiązania
muszą spełniać warunki holonomiczności;
\item w teorii miary dżet może istnieć tylko prawie wszędzie,
a zwykłą granicę zastępuje się \emph{granicą przybliżoną}.
\end{enumerate}
Te trzy punkty widzenia spotykają się w twierdzeniu Whitneya:
\emph{kiedy formalnie zadane pochodne dają się zrealizować przez
prawdziwą funkcję gładką?} Wersja Whitneya--Luzina mówi zaś,
że wiele funkcji mierzalnych staje się gładkimi po pominięciu
zbioru o dowolnie małej mierze.

\section{Kiełki, klasy równoważności i wielomiany Taylora}
\begin{definicja}[Kiełek odwzorowania]
\label{def:dzety-kielek}
Niech $M,N$ będą gładkimi rozmaitościami i $p\in M$.
Dwa gładkie odwzorowania $f:U\to N$ i $g:V\to N$,
gdzie $p\in U\cap V$, mają ten sam \emph{kiełek w $p$},
jeżeli pokrywają się na pewnym otoczeniu $p$ zawartym w $U\cap V$.
Zbiór kiełków oznaczamy przez $\mathcal C^\infty_p(M,N)$.
\end{definicja}

\begin{definicja}[$k$-dżet odwzorowania]
\label{def:dzety-klasa}
Niech $k\in\mathbb N_0$, $f,g$ będą gładkie w otoczeniu $p\in M$,
a $f(p)=g(p)=q\in N$. Piszemy $f\sim_{p,k}g$, jeżeli
po wybraniu mapy $x$ wokół $p$ i mapy $y$ wokół $q$ zachodzi
\[
 \partial^\alpha (y\circ f\circ x^{-1})(x(p))
 =\partial^\alpha (y\circ g\circ x^{-1})(x(p)),
 \qquad |\alpha|\le k.
\]
Klasę równoważności oznaczamy przez $j_p^k f$ i nazywamy
\emph{$k$-dżetem $f$ w $p$}. Dla $k=0$ jest to po prostu para
$(p,f(p))$. Wartości pochodnych liczymy składowo.
\end{definicja}

\begin{lemat}[Niezależność od map]
\label{lem:dzety-mapy}
Relacja w Definicji~\ref{def:dzety-klasa} nie zależy od wyboru
map źródła i celu.
\end{lemat}
\begin{proof}
Niech $\varphi,\psi$ będą zmianami współrzędnych.
Dla odwzorowania współrzędnego $F$ nowy zapis ma postać
$\psi\circ F\circ\varphi^{-1}$. Reguła łańcuchowa mówi,
że jego pochodna rzędu $r\le k$ jest skończoną sumą wyrażeń
zawierających pochodne $F$ rzędów od $1$ do $r$, pochodne
$\psi,\varphi^{-1}$ i wartości funkcji w rozpatrywanym punkcie.
Jeżeli odpowiednie dane dla $F$ i $G$ są równe do rzędu $k$,
to równe są także wszystkie nowe pochodne. Zmieniając mapy
w przeciwnym kierunku, otrzymujemy równoważność.
\end{proof}

\begin{przyklad}[Różne funkcje, ten sam dżet]
\label{ex:dzety-taylor}
Dla $f(x)=\sin x$ i $g(x)=x$ mamy $j_0^2f=j_0^2g$,
ale $j_0^3f\ne j_0^3g$. Natomiast $f(x)=x+x^2$
i $g(x)=x-x^2$ mają identyczne $1$-dżety w zerze,
ale różne $2$-dżety.
\end{przyklad}

\begin{uwaga}[Nie należy utożsamiać drugiej pochodnej z tensorem]
\label{rem:dzety-hess}
Pierwsza pochodna funkcji skalarnych definiuje kanoniczny
kowektor $\mathrm df_p\in T_p^*M$. Dla drugiej pochodnej jest inaczej.
Współrzędna $x$ na $\mathbb R$ i współrzędna $y=x+x^2$ wokół zera
dałyby dla funkcji $f(x)=x$ odpowiednio $f''(0)=0$
oraz $\frac{\mathrm d^2 f}{\mathrm dy^2}(0)=-2$.
Oznacza to, że $2$-dżet jest obiektem niezmienniczym,
ale same współrzędne drugich pochodnych nie są tensorami.
Tensor Hessego powstaje kanonicznie dopiero po określeniu
koneksji (lub w punkcie krytycznym funkcji).
\end{uwaga}

\begin{definicja}[Algebra dżetów funkcji skalarnych]
\label{def:dzety-algebra}
Niech $\mathcal E_p$ oznacza pierścień kiełków gładkich
funkcji rzeczywistych na $M$ w punkcie $p$ i
$\mathfrak m_p=\{[f]:f(p)=0\}$ jego ideał maksymalny.
Definiujemy $\mathcal J_p^k=\mathcal E_p/\mathfrak m_p^{k+1}$.
\end{definicja}

\begin{twierdzenie}[Algebra i wielomiany Taylora]
\label{thm:dzety-algebra}
Dla $M=\mathbb R^n$ i $p=0$ odwzorowanie
\[
 [f]\longmapsto \sum_{|\alpha|\le k}
 \frac{\partial^\alpha f(0)}{\alpha!}X^\alpha
\]
indukuje izomorfizm algebr
$\mathcal J_0^k\cong\mathbb R[X_1,\ldots,X_n]/(X_1,\ldots,X_n)^{k+1}$.
W szczególności $\dim\mathcal J_0^k=\binom{n+k}{k}$.
\end{twierdzenie}
\begin{proof}
Wielomian Taylora rzędu $k$ iloczynu jest iloczynem
wielomianów Taylora po odrzuceniu wyrazów stopnia większego niż $k$,
co wynika z reguły Leibniza. Jądro odwzorowania składa się
z kiełków o pochodnych rzędów $0,\ldots,k$ równych zeru.
Lemat Hadamarda, otrzymywany przez całkowanie
$f(x)-f(0)=\sum_i x_i\int_0^1 \partial_i f(tx)\,\mathrm d\lambda_1(t)$,
daje kolejno dla takiego kiełka rozkład
$f(x)=\sum_{|\alpha|=k+1}x^\alpha h_\alpha(x)$
z funkcjami gładkimi $h_\alpha$ w otoczeniu zera.
To znaczy, że jądrem jest $\mathfrak m_0^{k+1}$.
Każdy wielomian stopnia co najwyżej $k$ sam jest gładki,
więc homomorfizm jest surjektywny. Liczba jednomianów
o $|\alpha|\le k$ wynosi $\binom{n+k}{k}$.
\end{proof}

\begin{przyklad}[Liczby dualne]
\label{ex:dzety-dualne}
Dla $n=k=1$ dostajemy
$\mathcal J_0^1\cong\mathbb R[\varepsilon]/(\varepsilon^2)$.
Dżet funkcji $f$ ma zapis $f(0)+f'(0)\varepsilon$;
reguła mnożenia automatycznie daje wzór na pochodną iloczynu.
To algebraiczne źródło arytmetyki liczb dualnych używanej
w automatycznym różniczkowaniu.
\end{przyklad}

\section{Wiązki dżetów i ich naturalne odwzorowania}
\begin{definicja}[Wiązka dżetów odwzorowań]
\label{def:dzety-bundle}
Dla gładkich rozmaitości $M,N$ oznaczamy przez $J^k(M,N)$
zbiór wszystkich $k$-dżetów lokalnych gładkich odwzorowań
$M\to N$. Naturalne odwzorowania to
$\pi_0^k(j_p^k f)=(p,f(p))$ oraz dla $0\le r\le k$
$\pi_r^k(j_p^k f)=j_p^r f$.
Dla ustalonego $f$ odwzorowanie $j^kf:p\mapsto j_p^kf$
nazywamy \emph{przedłużeniem dżetowym}.
\end{definicja}

W lokalnych współrzędnych na $M$ i $N$, o wymiarach $n,m$,
przestrzeń $J^k(M,N)$ ma współrzędne $(x^i,u^a,u^a_\alpha)$,
$1\le|\alpha|\le k$, gdzie $u^a_\alpha$ odpowiada pochodnej
$\partial^\alpha u^a$. Jej wymiar wynosi
$n+m\binom{n+k}{k}$. Przestrzeń ta jest rozmaitością gładką,
ponieważ zmiany współrzędnych są otrzymywane z reguły
łańcuchowej i zależą gładko od skończonej liczby pochodnych.

\begin{twierdzenie}[Pierwsze dżety i struktura afiniczna]
\label{thm:dzety-affine}
Dla dowolnych $(p,q)\in M\times N$ włókno
$(\pi_0^1)^{-1}(p,q)$ utożsamia się kanonicznie z
$\operatorname{Hom}(T_pM,T_qN)$. Dla $k\ge2$ włókna
$\pi_{k-1}^k:J^k(M,N)\to J^{k-1}(M,N)$ są przestrzeniami
afinicznymi modelowanymi na
$S^kT_p^*M\otimes T_qN$, gdzie $(p,q)$ jest
wartością rozpatrywanego $(k-1)$-dżetu.
\end{twierdzenie}
\begin{proof}
Wartość i różniczka $f$ w $p$ jednoznacznie wyznaczają
$1$-dżet, a dowolne odwzorowanie liniowe
$T_pM\to T_qN$ jest różniczką odpowiedniej lokalnej funkcji
w mapach. Dla $k\ge2$ dwa przedłużenia tego samego
$(k-1)$-dżetu różnią się tylko pochodną rzędu $k$.
W regule zmiany współrzędnych wszystkie wyrazy
nieliniowe w pochodnych rzędu mniejszego niż $k$ znikają
przy odejmowaniu obu dżetów. Pozostaje tensorowe
przekształcenie symetrycznej $k$-formy o wartościach
w $T_qN$. Pozwala to odejmować dwa przedłużenia,
ale nie wybiera wyróżnionego zera w takim włóknie.
\end{proof}

\begin{uwaga}
Nie należy nazywać całej $J^k(M,N)$ wiązką wektorową
nad $M\times N$ dla $k\ge2$: zmiany współrzędnych
zawierają wyrazy niższych rzędów. Różnica między
\emph{wiązka wektorowa} a \emph{wiązka afiniczna}
jest w tym miejscu istotna.
\end{uwaga}

\section{Dżety przekrojów, holonomiczność i formy kontaktowe}
Niech $\pi:E\to M$ będzie gładką wiązką lokalnie trywialną.
Dżety lokalnych przekrojów $s$ tworzą $J^kE$.
Na włóknie lokalne współrzędne mają postać
$(x^i,u^a,u^a_\alpha)$ z $|\alpha|\le k$.
Każde $j^ks$ jest przekrojem $J^kE\to M$.

\begin{definicja}[Przekrój holonomiczny]
\label{def:dzety-holonomic}
Przekrój ciągły $\sigma:U\to J^kE$ nazywamy
\emph{holonomicznym}, jeżeli istnieje przekrój
$s:U\to E$ klasy $C^k$ taki, że $\sigma=j^ks$.
Gdy ponadto $\sigma$ jest klasy $C^1$, to $s$ jest
klasy $C^{k+1}$; wynika to kolejno ze współrzędnych
$u^a_\alpha\circ\sigma$, które są pochodnymi $s$.
\end{definicja}

Na $J^kE$ występują \emph{formy kontaktowe}
\[
 \theta^a_\alpha=\mathrm d u^a_\alpha
 -\sum_{i=1}^n u^a_{\alpha+e_i}\,\mathrm d x^i,
 \qquad |\alpha|\le k-1.
\]
Dla $\alpha=0$ przyjmujemy $u_0^a=u^a$.

\begin{twierdzenie}[Kryterium holonomiczności]
\label{thm:dzety-contact}
Niech $U\subseteq M$ będzie otwarty, $k\ge1$,
a $\sigma:U\to J^kE$ będzie przekrojem klasy $C^1$.
Wtedy $\sigma=j^ks$ dla pewnego przekroju
$s:U\to E$ klasy $C^{k+1}$ wtedy i tylko wtedy, gdy
$\sigma^*\theta^a_\alpha=0$ dla wszystkich form
kontaktowych, w dowolnych lokalnych trywializacjach.
\end{twierdzenie}
\begin{proof}
Dla $\sigma=j^ks$ zachodzi
$u^a_\alpha\circ\sigma=\partial^\alpha s^a$,
więc $\sigma^*\theta^a_\alpha=0$ wynika z różniczkowania.
Odwrotnie, niech $s=\pi_0^k\circ\sigma$, tj. składowa
zerowego rzędu. Równości $\sigma^*\theta^a_0=0$ dają
$\partial_i s^a=u^a_{e_i}\circ\sigma$. Indukcyjnie,
$\sigma^*\theta^a_\alpha=0$ implikuje
$\partial_i(u^a_\alpha\circ\sigma)
=u^a_{\alpha+e_i}\circ\sigma$. Otrzymujemy
$u^a_\alpha\circ\sigma=\partial^\alpha s^a$
dla $|\alpha|\le k$. Ponieważ wszystkie współrzędne $u^a_\alpha\circ\sigma$
są klasy $C^1$, otrzymujemy ponadto $s\in C^{k+1}$.
W szczególności regularność $C^1$ pola dżetów nie jest
tożsama z regularnością $C^1$ przekroju wyjściowego.
Konstrukcja jest zgodna na nakładaniu się map.
\end{proof}

\begin{przyklad}[Formalne dane nie muszą być pochodnymi]
W trywialnej wiązce $\mathbb R\times\mathbb R\to\mathbb R$
przekrój $\sigma(x)=(x,u(x),p(x))=(x,0,1)$
przestrzeni $J^1E$ ma zerową składową wartości i
stałą składową pochodnej równą $1$. Nie pochodzi od
żadnej funkcji $u$, ponieważ
$\sigma^*(\mathrm du-p\,\mathrm dx)=-\mathrm dx\ne0$.
\end{przyklad}

\begin{center}
\begin{tikzpicture}[>=stealth, every node/.style={font=\small}]
 \node (jk) at (0,1) {$J^kE$};
 \node (jm) at (2.2,1) {$J^{k-1}E$};
 \node (je) at (3.8,1) {$\cdots$};
 \node (e) at (5.2,1) {$E$};
 \node (m) at (6.6,1) {$M$};
 \node (u) at (0,-.25) {$U\subset M$};
 \draw[->] (jk) -- node[above] {$\pi_{k-1}^k$} (jm);
 \draw[->] (jm) -- (je);
 \draw[->] (je) -- (e);
 \draw[->] (e) -- node[above] {$\pi$} (m);
 \draw[->] (u) -- node[left] {$j^k s$} (jk);
\end{tikzpicture}

{\small Rzutowania w wieży dżetów odrzucają pochodne najwyższego rzędu.}
\end{center}

\section{Równania różniczkowe jako podzbiory przestrzeni dżetów}
\begin{definicja}[Równanie różniczkowe rzędu $k$]
\label{def:dzety-pde}
Niech $\pi:E\to M$ będzie gładką wiązką.
Układ równań różniczkowych rzędu nie większego niż $k$
możemy opisać \emph{lokalnie} zbiorem zer
$\mathcal R\cap V=\{z\in V:F(z)=0\}$, gdzie
$V\subseteq J^kE$ jest otwarty, a
$F:V\to\mathbb R^r$ jest gładkie. Nie każda globalna
podrozmaitość równań musi być zbiorem zer jednego
globalnie określonego odwzorowania do $\mathbb R^r$.
Przekrój $s$ jest rozwiązaniem, gdy $j^ks(x)\in\mathcal R$
dla każdego $x$ z jego obszaru określenia.
\end{definicja}

\begin{przyklad}[Hamilton--Jacobi i równanie Laplace'a]
Dla rzeczywistej funkcji $u$ na $\mathbb R^n$
równanie Hamiltona--Jacobiego
$H(x,\nabla u(x))=0$ jest podzbiorem $J^1E$
opisanym przez $H(x,p)=0$. Równanie Laplace'a
$\sum_i\partial_{ii}u=0$ jest hiperpłaszczyzną
w każdym włóknie $J^2E\to J^1E$ o równaniu
$\sum_i u_{2e_i}=0$. W obu przypadkach sam wybór
punktu z $\mathcal R$ nie wystarcza: rozwiązanie
musi ponadto wyznaczać przekrój holonomiczny.
\end{przyklad}

\begin{definicja}[Symbol główny operatora liniowego]
\label{def:dzety-symbol}
Dla skalarnego liniowego operatora
$Lu=\sum_{|\alpha|\le k}a_\alpha(x)\partial^\alpha u$
o gładkich współczynnikach, na otwartym $U\subset\mathbb R^n$,
jego symbol główny jest wielomianem
$\sigma_k(L)(x,\xi)=\sum_{|\alpha|=k}a_\alpha(x)\xi^\alpha$.
\end{definicja}

\begin{uwaga}
Symbol główny jest zależny wyłącznie od najwyższej części
dżetu; pod zmianą map jego współczynniki przekształcają się
zgodnie z tensorową częścią struktury afinicznej
Twierdzenia~\ref{thm:dzety-affine}. Dla Laplasjanu
$\sigma_2(\Delta)(x,\xi)=|\xi|^2$.
\end{uwaga}

\begin{przyklad}[Równanie Eulera--Lagrange'a]
Niech $\Omega\subset\mathbb R^n$ będzie ograniczonym obszarem,
$u\in C^2(\Omega)$, a $L(x,z,p)$ będzie funkcją klasy $C^2$.
Zakładamy, że $L(x,u(x),\nabla u(x))$ jest całkowalna na $\Omega$.
Funkcjonał zależący od $1$-dżetu ma postać
\[
 \mathcal A(u)=\int_\Omega L(x,u(x),\nabla u(x))
 \,\mathrm d\lambda_n(x).
\]
Dla $\varphi\in C_c^\infty(\Omega)$ i zróżniczkowanego
pod całką wyrażenia (na zwartym nośniku $\varphi$) mamy
\[
 \frac{\mathrm d}{\mathrm dt}\mathcal A(u+t\varphi)\big|_{t=0}
 =\int_\Omega\left(L_z\varphi+
 \sum_i L_{p_i}\partial_i\varphi\right)\mathrm d\lambda_n(x).
\]
Jeżeli $u$ jest punktem stacjonarnym, czyli powyższa
pierwsza wariacja znika dla wszystkich
$\varphi\in C_c^\infty(\Omega)$, to po całkowaniu
przez części i zastosowaniu lematu podstawowego rachunku
wariacyjnego otrzymujemy równanie
$L_z-\sum_i\partial_i L_{p_i}=0$ w $\Omega$.
Jest ono zazwyczaj równaniem drugiego rzędu, a zatem
jego geometria znajduje się w $J^2E$, mimo że
sam funkcjonał zależy tylko od $J^1E$.
\end{przyklad}

\section{Od gładkości do miary: gęstość i granice przybliżone}
Odtąd pracujemy na $\mathbb R^n$. Przez $\lambda_n$ oznaczamy
miarę Lebesgue'a, a przez $B(x,r)$ kulę otwartą.

\begin{definicja}[Punkt gęstości i granica przybliżona]
\label{def:dzety-ap-limit}
Dla zbioru mierzalnego $A\subset\mathbb R^n$ punkt $x$
ma gęstość $1$ względem $A$, jeżeli
\[
 \lim_{r\downarrow0}
 \frac{\lambda_n(A\cap B(x,r))}{\lambda_n(B(x,r))}=1.
\]
Jeżeli $v$ jest funkcją mierzalną w otoczeniu $x$,
to piszemy $\operatorname{ap}\!\lim_{y\to x}v(y)=L$,
gdy dla każdego $\varepsilon>0$ zbiór
$\{y:|v(y)-L|>\varepsilon\}$ ma w $x$ gęstość $0$.
W przypadku funkcji określonej tylko na zbiorze mierzalnym
$A$ zakładamy, że $x$ jest punktem gęstości $1$ zbioru $A$
i liczymy gęstość zbiorów wyjątkowych w $A\cap B(x,r)$.
\end{definicja}

\begin{definicja}[Przybliżony $k$-dżet]
\label{def:dzety-approx-jet}
Niech $A\subset\mathbb R^n$ będzie mierzalny,
$x\in A$ punktem jego gęstości $1$, a
$f:A\to\mathbb R^m$ funkcją mierzalną.
Mówimy, że $f$ ma \emph{przybliżony dżet Peana rzędu $k\ge1$}
w $x$, jeżeli istnieje wielomian wektorowy
$P_x(h)=\sum_{|\alpha|\le k}a_\alpha(x)h^\alpha/\alpha!$,
przy $a_0(x)=f(x)$, taki że
\[
 \operatorname{ap}\!\lim_{\substack{y\to x\\y\in A}}
 \frac{f(y)-P_x(y-x)}{|y-x|^k}=0.
\]
Dla $k=1$ współczynniki liniowe wyznaczają
\emph{różniczkę przybliżoną} $D_{\mathrm{ap}}f(x)$.
W tym sformułowaniu wartość $f(x)$ jest częścią danych:
zmiana wartości funkcji w samym punkcie $x$ może więc
zniszczyć istnienie dżetu w tym punkcie. Natomiast
\emph{prawie wszędzie} pojęcie to zależy wyłącznie od
klasy funkcji modulo zbiory miary zero: dwa równe prawie
wszędzie przedstawiciele mają te same dżety w punktach
gęstości zbioru, na którym są równe.
\end{definicja}

\begin{lemat}[Jednoznaczność przybliżonego dżetu]
\label{lem:dzety-approx-unique}
Jeżeli $x$ jest punktem gęstości $1$ zbioru $A$,
to wielomian $P_x$ spełniający
Definicję~\ref{def:dzety-approx-jet} jest jednoznaczny.
\end{lemat}
\begin{proof}
Niech $Q=P_x-\widetilde P_x$ będzie różnicą dwóch takich
wielomianów. Wtedy $Q(0)=0$ i
$\operatorname{ap}\!\lim_{h\to0}Q(h)/|h|^k=0$.
Gdyby $Q\ne0$, jego pierwszy niezerowy jednorodny składnik
$Q_j$, $j\le k$, byłby niezerowy na pewnym otwartym
podzbiorze sfery jednostkowej. Na stożku o dodatnim kącie
bryłowym zachodziłoby
$|Q(h)|/|h|^k\ge c|h|^{j-k}$ dla dostatecznie małych $h$.
Stożek ma dodatnią, niezależną od promienia, względną miarę
w każdej kuli, co przeczy granicy przybliżonej.
\end{proof}

\begin{przyklad}[Funkcja nieciągła z przybliżoną pochodną]
\label{ex:dzety-approx}
Niech $A=\bigcup_{j\ge1}(2^{-j},2^{-j}+2^{-3j})$
i $f(x)=\chi_A(x)$, gdzie $\chi_A$ oznacza funkcję
charakterystyczną zbioru $A$ i $f(0)=0$. Wtedy $A$ ma
w zerze gęstość $0$, zatem $f$ ma w zerze przybliżoną
granicę równą $0$. Ponadto
$\lambda_1(A\cap(-r,r))=O(r^3)$ dla $r\downarrow0$,
a więc również $D_{\mathrm{ap}}f(0)=0$.
Zwykła ciągłość (a tym bardziej różniczkowalność) nie zachodzi,
ponieważ w każdym otoczeniu zera są punkty $f=1$.
\end{przyklad}

\begin{przyklad}[Dlaczego trzeba ustalić przedstawiciela?]
Niech $f=\chi_{\mathbb Q}$ na $\mathbb R$. W punktach
niewymiernych funkcja ma przybliżony dżet każdego rzędu
równy zeru, ponieważ zbiór $\mathbb Q$ ma miarę zero.
W punktach wymiernych $f(x)=1$, więc nie ma nawet
przybliżonej granicy równej własnej wartości $f(x)$.
Przedstawiciel zerowy tej samej klasy w $L^0$ ma natomiast
dżet zerowy w każdym punkcie. Rozróżnienie to nie ma
znaczenia dla twierdzeń prawie wszędzie.
\end{przyklad}

\begin{uwaga}[Gdzie tkwi istota mierzalności?]
Równość $D_{\mathrm{ap}}f(x)=A$ nie oznacza
zbieżności po \emph{wszystkich} punktach $y\to x$.
Wystarcza odrzucanie zbioru o gęstości zerowej,
który może być gęsty topologicznie. Dlatego dżety
przybliżone nie są dżetami kiełków w sensie
Definicji~\ref{def:dzety-klasa}.
\end{uwaga}

\section{Rademacher, Sobolew i $BV$}
\begin{twierdzenie}[Rademacher i mierzalny pierwszy dżet]
\label{thm:dzety-rademacher}
Niech $U\subset\mathbb R^n$ będzie otwarty, a
$f:U\to\mathbb R^m$ lokalnie Lipschitzowskie.
Wtedy $f$ jest różniczkowalne w $\lambda_n$-prawie każdym
punkcie $U$. W tych punktach istnieje przybliżony
$1$-dżet i $D_{\mathrm{ap}}f(x)=Df(x)$.
Ponadto $x\mapsto Df(x)$ ma mierzalny przedstawiciel.
\end{twierdzenie}
\begin{proof}
Istnienie zwykłej różniczki prawie wszędzie jest treścią
twierdzenia Rademachera. Zwykła granica implikuje
przybliżoną, więc równość dżetów jest natychmiastowa.
Zbiór różniczkowalności jest borelowski. Istotnie,
na zbiorze istnienia wszystkich pochodnych cząstkowych
ich wartości są granicami ilorazów różnicowych z
przyrostami $1/j$, a więc są funkcjami borelowskimi.
Warunek
$|f(x+h)-f(x)-A(x)h|\le |h|/m$
dla wszystkich dostatecznie małych niezerowych $h$
można, dzięki ciągłości $f$, sprawdzać na przeliczalnym
zbiorze przyrostów wymiernych. Kwantyfikatory
$\forall m\in\mathbb N\,\exists j\in\mathbb N\,
\forall h\in\mathbb Q^n\setminus\{0\}$ z $|h|<1/j$
i $x+h\in U$ opisują zatem borelowski podzbiór $S$,
który jest dokładnie zbiorem różniczkowalności.
Wynika to z ciągłości wyrażenia względem przyrostu $h$
poza $h=0$. Przedłużenie macierzy
$Df$ przez zero na dopełnienie tego zbioru jest mierzalne.
\end{proof}

\begin{przyklad}[Punkt o dwóch geometrycznych dżetach]
Funkcja $f(x)=|x|$ na $\mathbb R$ ma zwykłą różniczkę
$Df(x)=\operatorname{sgn}(x)$ dla $x\ne0$.
W zerze nie ma ani zwykłego, ani przybliżonego
$1$-dżetu: dodatnia i ujemna półprosta mają dodatnią
gęstość i dawałyby sprzeczne pochodne. Brak dżetu w jednym
punkcie nie narusza twierdzenia Rademachera.
\end{przyklad}

\begin{twierdzenie}[Dżet przybliżony funkcji Sobolewa]
\label{thm:dzety-sobolev}
Niech $U\subset\mathbb R^n$ będzie otwarty,
$1\le p\le\infty$, a $u\in W_{\mathrm{loc}}^{1,p}(U)$.
Wówczas istnieje przedstawiciel $u$, który ma
przybliżony $1$-dżet w $\lambda_n$-prawie każdym
punkcie $U$, przy czym
$D_{\mathrm{ap}}u(x)=\nabla u(x)$ prawie wszędzie.
Po prawej stronie występuje gradient słaby.
\end{twierdzenie}
\begin{proof}
Korzystamy z następującego twierdzenia o precyzyjnych
przedstawicielach Sobolewa: funkcja klasy $W^{1,1}_{\rm loc}$
jest przybliżenie różniczkowalna prawie wszędzie, a jej
różniczka przybliżona pokrywa się z gradientem słabym.
Standardowy dowód tego wyniku opiera się na różniczkowaniu
całek i przybliżeniu funkcji Sobolewa przez funkcje gładkie.
Dla $p>1$ można go uzyskać także z oszacowań funkcji
maksymalnej Hardy'ego--Littlewooda; przypadek $p=1$
wymaga wersji słabego typu $(1,1)$. Ponieważ
$W^{1,p}_{\rm loc}\subset W^{1,1}_{\rm loc}$,
otrzymujemy tezę. Podany rezultat o przedstawicielach
Sobolewa jest w tym dowodzie \emph{jawnie użytym
wcześniejszym twierdzeniem}, a nie wynikiem wyprowadzonym
z samej definicji dżetu.
\end{proof}

\begin{uwaga}[Funkcje $BV$ i punkty skoku]
\label{rem:dzety-bv}
Dla $u\in BV_{\rm loc}(U)$ istnieje przybliżona
różniczka w prawie każdym punkcie względem $\lambda_n$,
lecz pochodna dystrybucyjna jest miarą wektorową, która
może mieć część osobliwą. W pełnym twierdzeniu
strukturalnym dla $BV$ rozkłada się ją na
\[
 Du=\nabla u\,\lambda_n
 +(u^+-u^-)\nu_u\,\mathcal H^{n-1}\llcorner J_u+D^cu,
\]
gdzie $J_u$ jest przybliżonym zbiorem skoków,
$u^\pm$ jednostronnymi śladami, $\nu_u$ normalną,
a $D^cu$ częścią Cantora. Jest to odrębny, głębszy wynik:
nie należy wyprowadzać go wyłącznie z istnienia
przybliżonego $1$-dżetu prawie wszędzie.
\end{uwaga}

\begin{przyklad}[Skok a przybliżony dżet]
Dla $E=\{x\in\mathbb R^n:x_1>0\}$ funkcja
$u=\chi_E$ jest lokalnie stała poza hiperpłaszczyzną
$x_1=0$, więc $D_{\mathrm{ap}}u=0$ prawie wszędzie.
Jednak $Du=e_1\,\mathcal H^{n-1}\llcorner\{x_1=0\}$
(z orientacją odpowiadającą przejściu z $0$ do $1$).
Zatem przybliżony dżet pierwszego rzędu nie opisuje
całej miary pochodnej funkcji $BV$.
\end{przyklad}

\section{Jakobiany i wzory area/coarea odczytywane z dżetów}
Niech $f:U\subseteq\mathbb R^k\to\mathbb R^N$,
$k\le N$, będzie lokalnie Lipschitzowskie.
W punkcie różniczkowalności określamy
$J_kf(x)=\sqrt{\det(Df(x)^{\mathsf T}Df(x))}$.
Jest to funkcja samego \emph{liniowego składnika $1$-dżetu}.
Dla mierzalnego $A\subset U$ i nieujemnej
mierzalnej funkcji $h:U\to[0,\infty]$ wzór area ma postać
\[
 \int_A h(x)J_kf(x)\,\mathrm d\lambda_k(x)
 =\int_{\mathbb R^N}\sum_{x\in A\cap f^{-1}(\{y\})}
 h(x)\,\mathrm d\mathcal H^k(y).
\]
Dla lokalnie Lipschitzowskiego $F:V\subset\mathbb R^n\to
\mathbb R^m$, $n\ge m$, gdzie $V$ jest otwarty,
i nieujemnej mierzalnej $g:V\to[0,\infty]$, wzór
koobszaru brzmi
\[
 \int_V g(x)\sqrt{\det(DF(x)DF(x)^{\mathsf T})}
 \,\mathrm d\lambda_n(x)
 =\int_{\mathbb R^m}\left(
 \int_{F^{-1}(\{y\})}g(x)\,\mathrm d\mathcal H^{n-m}(x)
 \right)\mathrm d\lambda_m(y).
\]
W obu wzorach miary Hausdorffa są znormalizowane tak,
że $\mathcal H^j|_{\mathbb R^j}=\lambda_j$ dla całkowitych
$j\ge1$ oraz $\mathcal H^0$ jest miarą liczącą.
Dla zbiorów mierzalnych w sensie Lebesgue'a konieczne
są właściwe sformułowania twierdzeń, uwzględniające
ukończenie miar i mierzalność funkcji krotnościowej.
W szczególności sama teoria dżetów nie rozwiązuje
kwestii mierzalności włókien i całek po nich.

\begin{uwaga}
Teoria dżetów wyjaśnia, \emph{jakie dane} wyznaczają
lokalny współczynnik geometryczny. Sam wzór area/coarea
wymaga jednak niezależnie twierdzeń o mierzalności,
różniczkowaniu i krotnościach włókien. Nie jest więc
czysto formalną konsekwencją teorii dżetów.
\end{uwaga}

\section{Pola dżetów i twierdzenie Whitneya}
\begin{definicja}[Pole dżetów Whitneya klasy $C^k$]
\label{def:dzety-whitney-field}
Niech $K\subset\mathbb R^n$ będzie domknięty,
$k\in\mathbb N_0$, a dla każdego $|\alpha|\le k$
niech $a_\alpha:K\to\mathbb R$ będzie funkcją.
Kładziemy $P_x(z)=\sum_{|\alpha|\le k}
 a_\alpha(x)(z-x)^\alpha/\alpha!$.
Rodzinę $(P_x)_{x\in K}$ nazywamy
\emph{polem dżetów Whitneya klasy $C^k$}, jeżeli
wszystkie $a_\alpha$ są ciągłe i dla każdego zwartego
$K_0\subset K$ oraz każdego $|\alpha|\le k$ zachodzi
\[
 \lim_{\delta\downarrow0}
 \sup_{\substack{x,y\in K_0,\;0<|x-y|<\delta}}
 \frac{|D^\alpha P_x(y)-a_\alpha(y)|}
 {|x-y|^{k-|\alpha|}}=0,
\]
przy umowie, że supremum po zbiorze pustym wynosi $0$.
\end{definicja}

\begin{uwaga}
Dla $|\alpha|=k$ mianownik jest równy $1$:
warunek mówi o ciągłości najwyższych pochodnych.
Dla $k=1$ warunki na zwartym $K$ sprowadzają się do
ciągłości $a_0,v=(a_{e_1},\ldots,a_{e_n})$ i
$|a_0(y)-a_0(x)-v(x)\cdot(y-x)|=o(|x-y|)$
jednostajnie dla $x,y\in K$.
\end{uwaga}

\begin{lemat}[Kostki Whitneya i rozkład jedności]
\label{lem:dzety-whitney-cubes}
Niech $K\subset\mathbb R^n$ będzie niepustym zbiorem
domkniętym, różnym od całej przestrzeni, a
$\Omega=\mathbb R^n\setminus K$. Istnieje przeliczalna
rodzina domkniętych kostek diadycznych $Q$, o rozłącznych
wnętrzach pokrywających $\Omega$, taka że
$\ell(Q)\asymp\operatorname{dist}(Q,K)$. Istnieje
również gładki rozkład jedności $(\varphi_Q)$ na $\Omega$
o nośnikach w nieznacznie powiększonych kostkach $Q^*$,
takich że
$|D^\alpha\varphi_Q|\le C_\alpha\ell(Q)^{-|\alpha|}$.
W każdym punkcie co najwyżej $N(n)$ funkcji
$\varphi_Q$ jest niezerowych.
\end{lemat}
\begin{proof}
Rozważmy siatkę wszystkich kostek diadycznych i wybierzmy
maksymalne kostki $Q$ spełniające
$8\sqrt n\,\ell(Q)\le\operatorname{dist}(Q,K)$.
Każdy punkt $z\in\Omega$ leży w takiej kostce, ponieważ
w dostatecznie małej skali nierówność jest spełniona,
a w dostatecznie dużej skali przestaje być spełniona.
Maksymalne kostki mają rozłączne wnętrza.
Nierówność wyboru daje dolne oszacowanie odległości.
Gdyby odległość była nieporównywalnie większa od boku,
to również kostka rodzicielska spełniałaby warunek,
wbrew maksymalności; daje to oszacowanie górne.
Kostki o przecinających się niewielkich powiększeniach
mają porównywalne boki, bo ich odległości od $K$
są porównywalne. Dlatego, po ustaleniu współczynnika
powiększenia dostatecznie bliskiego $1$, ich krotność
jest ograniczona liczbą zależną tylko od wymiaru.
Dla każdej kostki wybieramy funkcję $\eta_Q\in C_c^\infty(Q^*)$
równą $1$ na $Q$, z
$|D^\alpha\eta_Q|\le C_\alpha\ell(Q)^{-|\alpha|}$;
uzyskujemy ją przez przeskalowanie ustalonej funkcji
obcięcia na kostce jednostkowej. Suma
$S=\sum_Q\eta_Q$ jest lokalnie skończona,
$1\le S\le N(n)$. Funkcje
$\varphi_Q=\eta_Q/S$ mają wymagane własności
na mocy reguły Leibniza i oszacowania nakładania się
kostek.
\end{proof}

\begin{lemat}[Sklejanie ciągłych pochodnych przez zbiór domknięty]
\label{lem:dzety-glue}
Niech $K\subset\mathbb R^n$ będzie domknięty. Załóżmy, że
$H:\mathbb R^n\to\mathbb R$ jest ciągła, klasy $C^1$ poza $K$,
a $V:\mathbb R^n\to\mathbb R^n$ jest ciągła i spełnia
$V=\nabla H$ poza $K$. Jeżeli dla każdego $a\in K$
zachodzi $H(a+h)=H(a)+V(a)\cdot h+o(|h|)$, to
$H\in C^1(\mathbb R^n)$ i $\nabla H=V$ wszędzie.
\end{lemat}
\begin{proof}
Poza $K$ teza jest znana. W punkcie $a\in K$ podany
warunek jest definicją różniczkowalności i daje
$\nabla H(a)=V(a)$. Ponieważ $V$ jest ciągła na całej
przestrzeni, otrzymany gradient jest ciągły, czyli
$H\in C^1$. Nie trzeba różniczkować po krzywych
zawartych w $K$.
\end{proof}

\begin{twierdzenie}[Twierdzenie Whitneya o przedłużaniu]
\label{thm:dzety-whitney-extension}
Niech $K\subset\mathbb R^n$ będzie domknięty,
$k\in\mathbb N_0$, a $(a_\alpha)_{|\alpha|\le k}$
polem dżetów Whitneya klasy $C^k$ na $K$.
Wtedy istnieje $F\in C^k(\mathbb R^n)$ spełniająca
$D^\alpha F(x)=a_\alpha(x)$ dla wszystkich
$x\in K$ i $|\alpha|\le k$.
Warunki Definicji~\ref{def:dzety-whitney-field} są
zarazem konieczne dla takich przedłużeń.
\end{twierdzenie}

\begin{proof}
Konieczność wynika z wzoru Taylora z resztą Peana,
stosowanego do wszystkich $D^\alpha F$, i z jednostajnej
ciągłości pochodnych na zbiorach zwartych.
Przedstawimy konstrukcję wystarczalności; dla $K=\varnothing$
można przyjąć $F=0$.

\emph{Krok 1: kostki Whitneya.}
Jeżeli $K=\mathbb R^n$, przyjmujemy $F=a_0$
i teza wynika wprost z warunków Whitneya.
W przeciwnym przypadku Lemat~\ref{lem:dzety-whitney-cubes}
dostarcza rozkładu zbioru
$\Omega=\mathbb R^n\setminus K$
na kostki $Q$, wraz z rozkładem jedności
$(\varphi_Q)$ i oszacowaniami ich pochodnych.

\emph{Krok 2: sklejanie wielomianów.}
Dla każdej kostki wybierzmy $b_Q\in K$ takie, że
$|b_Q-Q|\le2\operatorname{dist}(Q,K)$, i połóżmy
$F(z)=\sum_Q\varphi_Q(z)P_{b_Q}(z)$ dla $z\in\Omega$,
a na $K$ niech $F(x)=a_0(x)$. Suma jest lokalnie skończona
i funkcja jest gładka poza $K$.

\emph{Krok 3: oszacowanie w pobliżu brzegu.}
Ustalmy zwarty fragment $K_0\subset K$ i niewielką
jego otoczkę. Z warunków Whitneya istnieje funkcja
$\omega(t)\downarrow0$ dla $t\downarrow0$ taka, że
jeżeli $a,b$ należą do tej otoczki i $|a-b|\le t$, to
\[
 |D^\beta P_a(b)-a_\beta(b)|
 \le\omega(t)|a-b|^{k-|\beta|},\qquad |\beta|\le k.
\]
Z rozwinięcia wielomianów i tego oszacowania wynika,
że dla $|z-a|+|z-b|\le Ct$ zachodzi
$|D^\beta(P_a-P_b)(z)|\le C\omega(Ct)t^{k-|\beta|}$.
Niech teraz $z\in\Omega$, $r=\operatorname{dist}(z,K)$,
a $a\in K$ spełnia $|z-a|\le2r$. Wszystkie kostki
$Q$ z $\varphi_Q(z)\ne0$ mają $\ell(Q)\asymp r$
i $|b_Q-a|\le Cr$. Różniczkując
$F(z)-P_a(z)=\sum_Q\varphi_Q(z)(P_{b_Q}(z)-P_a(z))$
i korzystając z ograniczenia liczby nakładających się
kostek otrzymujemy dla $|\beta|\le k$
\[
 |D^\beta F(z)-D^\beta P_a(z)|
 \le C\omega(Cr)r^{k-|\beta|}.
\]
W szczególności pochodne do rzędu $k$ mają przy $K$
granice zadane przez współczynniki $a_\beta$.

\emph{Krok 4: regularność na $K$.}
Dla $x\in K$ i $z\to x$ porównujemy $P_a$ z $P_x$,
gdzie $a\in K$ jest jak wyżej położonym punktem najbliższym
$z$. Z poprzednich oszacowań i warunków Whitneya wynika
$F(z)-P_x(z)=o(|z-x|^k)$; analogiczne rozwinięcia
obowiązują dla przedłużeń $D^\beta F$ rzędu do
$k-|\beta|$. Dla $k=0$ wystarcza uzyskana ciągłość $F$ na $K$
i na $\Omega$. Dla $k\ge1$ doprecyzujemy indukcję.
Dla
$|\beta|\le k$ na $\Omega$ mamy funkcję
$D^\beta F$, której ciągłym przedłużeniem na $K$
jest $a_\beta$ (wynika to z oszacowania poprzedniego
kroku i ciągłości danych). Dla $|\beta|<k$ warunki
Whitneya zapewniają w każdym $a\in K$ rozwinięcie
tego przedłużenia pierwszego rzędu z gradientem
$(a_{\beta+e_i})_{i=1}^n$. Stosując
Lemat~\ref{lem:dzety-glue} kolejno do pochodnych
rzędów $0,1,\ldots,k-1$, wykazujemy
$F\in C^k(\mathbb R^n)$ oraz $D^\alpha F|_K=a_\alpha$.
W ten sposób przejście przez $K$ nie wymaga
nieuzasadnionego różniczkowania funkcji na samym zbiorze
$K$, który może mieć puste wnętrze.
\end{proof}

\begin{przyklad}[Dlaczego sama ciągłość dżetów nie wystarcza?]
Niech $K=\{0\}\cup\{1/j:j\ge1\}\subset\mathbb R$,
$a_0(0)=0$, $a_0(1/j)=1/j$ i $a_1\equiv0$.
Obie funkcje $a_0,a_1$ są ciągłe na $K$, ale
$|a_0(1/j)-a_0(0)-a_1(0)/j|/(1/j)=1$.
Nie istnieje więc funkcja $C^1$ o tych wartościach i
pochodnych. Jest to konkretne naruszenie warunku Whitneya.
\end{przyklad}

\section{Twierdzenie Whitneya--Luzina: od dżetów przybliżonych do $C^1$}
\begin{lemat}[Mierzalność różniczki przybliżonej]
\label{lem:dzety-ap-meas}
Niech $U\subset\mathbb R^n$ będzie otwarty,
a $f:U\to\mathbb R^m$ mierzalna. Na mierzalnym
zbiorze punktów, w których $f$ jest przybliżenie
różniczkowalna, odwzorowanie $D_{\mathrm{ap}}f$
ma mierzalny przedstawiciel.
\end{lemat}
\begin{proof}
Zmienimy wartości na zbiorze miary zero tylko wtedy,
gdy jest to potrzebne do zdefiniowania funkcji na całym $U$.
Dla $M\in\mathbb N$ i $r>0$ rozważmy macierz
\[
 T_{M,r}(x)=\frac{n}{\lambda_n(B(0,r))}
 \int_{B(0,r)}\frac{f(x+h)-f(x)}{|h|}
 \frac{h^{\mathsf T}}{|h|}
 \chi_{\{|f(x+h)-f(x)|\le M|h|\}}
 \,\mathrm d\lambda_n(h),
\]
gdzie dla $x+h\notin U$ całkowaną funkcję przyjmujemy za zero.
Jest to odwzorowanie mierzalne, ponieważ całkowana funkcja
jest mierzalna i ograniczona przez $M$.
Jeżeli $f$ ma różniczkę przybliżoną $A$ w $x$ i
$M>\|A\|+1$, to na kulach malejących do zera
iloraz różnicowy jest w mierze bliski
$A(h/|h|)$; obcięcie nie usuwa prawie wszystkich
kierunków. Ponieważ całkowana funkcja jest ograniczona,
własność zbieżności w mierze daje
$\lim_{r\downarrow0}T_{M,r}(x)=A$.
Wykorzystaliśmy tożsamość
$\lambda_n(B(0,r))^{-1}\int_{B(0,r)}
 hh^{\mathsf T}/|h|^2\,\mathrm d\lambda_n(h)=I_n/n$.
Dla ciągu $r=1/j$ granice macierzy $T_{M,1/j}$ są
mierzalne na zbiorach swojego istnienia. W każdym punkcie
przybliżonej różniczkowalności wszystkie dostatecznie
duże wartości $M$ dają tę samą granicę $A$.
Wybór pierwszego stabilnego ogona tych granic jest
operacją na przeliczalnej rodzinie funkcji mierzalnych,
a więc daje mierzalny przedstawiciel $A$.
\end{proof}

\begin{twierdzenie}[Whitney--Luzin dla pierwszego dżetu]
\label{thm:dzety-whitney-lusin}
Niech $U\subset\mathbb R^n$ będzie otwarty,
a $f:U\to\mathbb R^m$ funkcją mierzalną,
przybliżenie różniczkowalną w $\lambda_n$-prawie każdym
punkcie $U$. Dla każdego zbioru mierzalnego $A\subset U$
o skończonej mierze i każdego $\varepsilon>0$ istnieją
zbiór zwarty $K\subset A$ i odwzorowanie
$F\in C^1(\mathbb R^n,\mathbb R^m)$ takie, że
\[
 \lambda_n(A\setminus K)<\varepsilon,
 \qquad F(x)=f(x),\quad DF(x)=D_{\mathrm{ap}}f(x)
 \quad (x\in K).
\]
Odwrotnie, jeżeli dla każdego mierzalnego $A\subset U$
o skończonej mierze i każdego $\varepsilon>0$ istnieją
$K,F$ spełniające powyższą równość wartości (bez wymogu
dotyczącego pochodnych), to $f$ jest przybliżenie
różniczkowalna prawie wszędzie na $U$.
\end{twierdzenie}

\begin{proof}
Wystarczy dowodzić dla funkcji rzeczywistej: dla wektorowych
stosujemy ten sam wybór zwartego zbioru i wspólne
oszacowania normy w $\mathbb R^m$.

\emph{Krok 1: duży zbiór z ciągłymi danymi.}
Niech $D$ będzie pełnomiarowym podzbiorem $U$, gdzie
istnieje różniczka przybliżona $G(x)=D_{\mathrm{ap}}f(x)$.
Dla sformułowania „prawie wszędzie” wybieramy mierzalny
zbiór $D$ o pełnej mierze, na którym dżet istnieje;
na nim działa Lemat~\ref{lem:dzety-ap-meas}.
Funkcje $f,G$ są więc mierzalne na $D$. Z twierdzenia Luzina i regularności
miary Lebesgue'a wybieramy zwarty $K_0\subset A\cap D$,
$\lambda_n(A\setminus K_0)<\varepsilon/3$, tak że
$f|_{K_0}$ i $G|_{K_0}$ są ciągłe.
W szczególności $G|_{K_0}$ jest jednostajnie ciągła.
Ponieważ $K_0\Subset U$, ma dodatnią odległość od
$\mathbb R^n\setminus U$.

\emph{Krok 2: jednostajna kontrola zbiorów złych.}
Dla $j\in\mathbb N$ i $r>0$ oznaczmy przez
$\vartheta_{j,r}(x)$ iloraz
\[
 \frac{\lambda_n\bigl(\{z\in B(x,r)\cap U:
 |f(z)-f(x)-G(x)(z-x)|>j^{-1}|z-x|\}\bigr)}
 {\lambda_n(B(x,r))}.
\]
Z przybliżonej różniczkowalności wynika
$\vartheta_{j,r}(x)\to0$ przy $r\downarrow0$ dla
każdego $x\in D$. Dla ustalonego $j$ funkcje
$H_{j,l}(x)=\sup_{0<r<1/l}\vartheta_{j,r}(x)$ są
mierzalne: dla ustalonego $x$ licznik zależy ciągle
od $r$ (sfery mają miarę zero), toteż supremum można
brać po promieniach wymiernych; mierzalność całki
zależącej od $x$ daje twierdzenie Tonellego.
Mamy $H_{j,l}(x)\downarrow0$.

Wybierzmy stałą $\delta=8^{-n}$. Z twierdzenia Jegorowa
stosowanego kolejno dla $j=1,2,\ldots$ wybieramy
mierzalne podzbiory $E_j\subset K_0$ tak, że
$\lambda_n(K_0\setminus E_j)<\varepsilon/(3\cdot 2^j)$
oraz $H_{j,l}\to0$ jednostajnie na $E_j$.
Niech $E=\bigcap_j E_j$. Następnie wybierzmy zwarty
$K\subset E$ z $\lambda_n(E\setminus K)<\varepsilon/3$.
Wtedy $\lambda_n(A\setminus K)<\varepsilon$.
Dla każdego $j$ istnieje $\rho_j>0$ taka, że
$\vartheta_{j,r}(x)<\delta$ jednocześnie dla
$x\in K$ i $0<r<\rho_j$.

\emph{Krok 3: z kontroli w mierze do warunku Whitneya.}
Niech $x,y\in K$, $d=|x-y|>0$ i
$2d<\rho_j$; wymagamy także, aby $B(x,2d),B(y,2d)$
były zawarte w $U$. Przecięcie tych kul zawiera
$B((x+y)/2,3d/2)$, więc jego miara jest większa
niż suma dwóch górnych oszacowań
$\delta\lambda_n(B(0,2d))$ na zbiory złe
(dla $\delta=8^{-n}$). Istnieje zatem punkt $z$
w przecięciu, dla którego jednocześnie
\[
 |f(z)-f(x)-G(x)(z-x)|\le 2d/j,
 \quad |f(z)-f(y)-G(y)(z-y)|\le 2d/j.
\]
Odejmując równości otrzymujemy
\[
 |f(y)-f(x)-G(x)(y-x)|
 \le4d/j+2d\,\|G(x)-G(y)\|.
\]
Jednostajna ciągłość $G$ na $K$ i dowolność $j$
dowodzą, że iloraz lewej strony przez $|x-y|$
dąży jednostajnie do zera. Tak więc $(f|_K,G|_K)$
jest polem dżetów Whitneya klasy $C^1$.

\emph{Krok 4: przedłużenie.}
Twierdzenie~\ref{thm:dzety-whitney-extension}
daje $F\in C^1(\mathbb R^n)$ o danych wartościach
i różniczkach na $K$.

\emph{Odwrotność.}
Dla każdego wyczerpującego ciągu zbiorów
$A_l=U\cap B(0,l)$ i każdego $l$ możemy wybrać
zbiory zwarte $K_{l,j}\subset A_l$, na których
$f=F_{l,j}\in C^1$, z
$\lambda_n(A_l\setminus K_{l,j})<1/j$.
Prawie każdy punkt $A_l$ należy do przynajmniej
jednego z tych zbiorów i jest jego punktem gęstości.
W takim punkcie $f$ ma przybliżoną różniczkę
$DF_{l,j}$, ponieważ obie funkcje zgadzają się na
zbiorze gęstości $1$. Biorąc sumę po $l$,
otrzymujemy przybliżoną różniczkowalność prawie wszędzie.
\end{proof}

\begin{wniosek}[Luzinowskie wygładzenie odwzorowania Lipschitzowskiego]
\label{cor:dzety-lusin-lipschitz}
Niech $U\subset\mathbb R^n$ będzie otwarty,
$f:U\to\mathbb R^m$ lokalnie Lipschitzowskie,
a $A\subset U$ mierzalny o skończonej mierze.
Dla każdego $\varepsilon>0$ istnieje
$F\in C^1(\mathbb R^n,\mathbb R^m)$ takie, że
\[
 \lambda_n(\{x\in A:F(x)\ne f(x)\})<\varepsilon,
\]
a na pewnym zwartym $K\subset A$ o
$\lambda_n(A\setminus K)<\varepsilon$ zachodzi
$DF=Df$.
\end{wniosek}
\begin{proof}
Twierdzenie Rademachera daje przybliżoną
różniczkowalność prawie wszędzie.
Stosujemy Twierdzenie~\ref{thm:dzety-whitney-lusin}.
\end{proof}

\begin{uwaga}[Granice twierdzenia]
Twierdzenie~\ref{thm:dzety-whitney-lusin} nie mówi,
że $f$ jest funkcją $C^1$ po zmianie na zbiorze miary zero.
Wolno odrzucić zbiór o \emph{dowolnie małej}, lecz niekoniecznie
zerowej mierze, zależny od $\varepsilon$.
Na przykład $f(x)=|x|$ na $(-1,1)$ nie ma globalnego
przedstawiciela $C^1$ równego $f$ prawie wszędzie.
\end{uwaga}

\section{Dżety wyższych rzędów i zgodność danych}
Dla $k\ge2$ sam formalny zapis danych
$(a_\alpha(x))_{|\alpha|\le k}$ nie gwarantuje,
że są one pochodnymi jednej funkcji. W teorii miary
istotne jest, czy można znaleźć duży zbiór, na którym
dane spełniają \emph{wszystkie} warunki Whitneya.

\begin{twierdzenie}[Luzinowskie przedłużenie zgodnych dżetów rzędu $k$]
\label{thm:dzety-lusin-higher}
Niech $U\subset\mathbb R^n$ będzie otwarty, $k\ge1$,
a $a_\alpha:U\to\mathbb R$, $|\alpha|\le k$,
będą funkcjami mierzalnymi. Załóżmy, że dla
$\lambda_n$-prawie każdego $x\in U$ i każdego
$|\alpha|\le k$ zachodzi zwykła, względna do
pełnomiarowego zbioru $D\subset U$, ekspansja
\[
 a_\alpha(y)=\sum_{|\beta|\le k-|\alpha|}
 \frac{a_{\alpha+\beta}(x)}{\beta!}(y-x)^\beta
 +o(|y-x|^{k-|\alpha|})\quad(D\ni y\to x).
\]
Wówczas dla każdego mierzalnego $A\subset U$
o skończonej mierze i każdego $\varepsilon>0$
istnieją $K\subset A$ zwarty i
$F\in C^k(\mathbb R^n)$ takie, że
$\lambda_n(A\setminus K)<\varepsilon$ oraz
$D^\alpha F|_K=a_\alpha|_K$ dla wszystkich
$|\alpha|\le k$.
\end{twierdzenie}
\begin{proof}
Twierdzenie Luzina zastosowane do skończonej rodziny
$a_\alpha$ daje zwarty $K_0\subset A\cap D$ o dowolnie
małym dopełnieniu i ciągłości wszystkich współczynników.
Dla każdej wielopochodnej $\alpha$ i promienia $r$
weźmy supremum po $y\in K_0$ z $0<|x-y|<r$
ilorazu reszty Taylora przez $|x-y|^{k-|\alpha|}$.
Jest to funkcja mierzalna względem $x\in K_0$:
możemy zastąpić wszystkie $y$ przeliczalnym zbiorem
gęstym w $K_0$, ponieważ na $K_0$ współczynniki
są ciągłe. Z hipotezy supremum dąży do zera punktowo przy
$r\downarrow0$, ponieważ $K_0\subset D$ i każda
reszta jest oceniana wzdłuż całego zbioru $D$.
To właśnie tu korzystamy z silniejszej, względnej
wersji założenia, a nie jedynie ze zbieżności przybliżonej. Twierdzenie Jegorowa
pozwala przejść do zwartego $K\subset K_0$ z
arbitralnie małą utratą miary, na którym zbieżność
jest jednostajna dla skończenie wielu $\alpha$.
Właśnie to są warunki pola Whitneya z
Definicji~\ref{def:dzety-whitney-field}. Zastosowanie
Twierdzenia~\ref{thm:dzety-whitney-extension}
kończy dowód.
\end{proof}

\begin{uwaga}[Ważne rozróżnienie]
Warunek z Twierdzenia~\ref{thm:dzety-lusin-higher}
jest mocniejszy od samego istnienia przybliżonego
wielomianu Peana rzędu $k$ dla jednej funkcji.
Dla ogólnego twierdzenia o przybliżonych dżetach
wyższych rzędów potrzebny jest dodatkowy krok
przechodzący od informacji w mierze do zgodności
wszystkich pochodnych na wybranym zbiorze.
Nie wolno pomijać go, stosując twierdzenie Whitneya.
\end{uwaga}

\section{Dodatkowe przykłady i powiązania}
\begin{przyklad}[Dżet holomorficzny]
Niech $U\subset\mathbb C^n$ będzie obszarem,
$p\in U$ i $f\in\mathcal O(U)$. Jego holomorficzny
$k$-dżet zapisujemy w lokalnych współrzędnych jako
\[
 j_p^kf\equiv\sum_{|\alpha|\le k}
 \frac{\partial_z^\alpha f(p)}{\alpha!}(z-p)^\alpha
 \pmod{\mathfrak m_p^{k+1}}.
\]
Warunki interpolacji wartości i pochodnych
(np. w zadaniach typu Carathéodory'ego--Fejéra)
są warunkami na holomorficzne dżety.
W przeciwieństwie do dowolnych gładkich funkcji,
holomorficzne dżety podlegają globalnym ograniczeniom
wynikającym z istnienia funkcji holomorficznej na
ustalonym obszarze. Twierdzenie Whitneya w klasie $C^k$
nie zapewnia holomorficznego przedłużenia.
\end{przyklad}

\begin{przyklad}[Dżety a metoda elementów skończonych]
Przybliżenia typu Hermite'a w jednym wymiarze
łączą zadane wartości $u(x_i)$ i pochodne $u'(x_i)$
na końcach przedziałów. Są więc interpolacją
$1$-dżetów na skończonym zbiorze węzłów.
Warunki zgodności na węzłach decydują,
czy powstaje aproksymacja klasy $C^1$.
\end{przyklad}

\begin{przyklad}[Dżety i geometryczne oszacowania rzędu]
Dwie gładkie funkcje o identycznym $k$-dżecie w $p$
mogą różnić się globalnie, ale ich różnica
spełnia $f(x)-g(x)=o(|x-p|^k)$ przy $x\to p$.
Jeżeli mają również równe pochodne rzędu $k+1$,
to ich $(k+1)$-dżety są równe i otrzymujemy
odpowiednią poprawę lokalnego rzędu reszty.
\end{przyklad}

\section{Mapa zależności i zakres udowodnionych rezultatów}
\begin{center}
\begin{tikzpicture}[>=stealth, every node/.style={font=\small,align=center},
 box/.style={draw,rounded corners=2pt,inner sep=5pt,text width=3.8cm}]
 \node[box] (j) at (-2.4,2.8) {Dżety gładkie\\wiązki i formy kontaktowe};
 \node[box] (a) at (2.4,2.8) {Gęstość Lebesgue'a\\dżety przybliżone};
 \node[box] (w) at (-2.4,0.7) {Warunki Whitneya\\przedłużenie $C^k$};
 \node[box] (l) at (2.4,0.7) {Whitney--Luzin\\wybór dużego $K$};
 \node[box] (c) at (0,-1.45) {Zastosowania: Sobolew, $BV$,\\area/coarea, PDE};
 \draw[->] (j) -- (w);
 \draw[->] (a) -- (l);
 \draw[->] (w) -- (l);
 \draw[->] (l) -- (c);
\end{tikzpicture}
\end{center}

\begin{uwaga}[Co zostało udowodnione, a co przyjęte?]
W tym rozdziale szczegółowo wyprowadzono algebraiczny opis
dżetów, strukturę afiniczną, kryterium holonomiczności,
jednoznaczność dżetów przybliżonych, twierdzenie Whitneya
metodą kostek i twierdzenie Whitneya--Luzina dla $C^1$.
Osobno przyjęto znane twierdzenia Rademachera, Luzina,
Jegorowa i Sobolewa o przybliżonej różniczkowalności
oraz ogólne wzory area/coarea i rozkład $BV$.
Ich przytoczenie pokazuje zastosowania, nie stanowi
nowego dowodu tych wyników.
\end{uwaga}


\section*{Wskazówki bibliograficzne}
Konstrukcja wiązek dżetów i form kontaktowych jest standardowo
przedstawiana w monografii D.~J.~Saundersa,
\emph{The Geometry of Jet Bundles}. Dla teorii przedłużeń
podstawową pracą jest H.~Whitney,
\emph{Analytic extensions of differentiable functions defined
in closed sets} (1934). Przybliżona różniczkowalność,
twierdzenia typu Luzina i ich związek z teorią Sobolewa i $BV$
są omawiane w książce L.~C.~Evansa i R.~F.~Gariepy'ego,
\emph{Measure Theory and Fine Properties of Functions}.
W tekście rozdzielono twierdzenia dowodzone tutaj od
rezultatów zewnętrznych (Rademacher, własności Sobolewa,
wzory area i koobszaru). W konstrukcji Whitneya wykorzystuje
się standardowy rozkład na kostki i rozkład jedności;
podano ich własności oraz argument sklejania na zbiorze
domkniętym. Dowody powinny być czytane z zachowaniem
wskazanych założeń regularności.

% END chapters/pl/measure_theory/31-dzety-geometria-i-teoria-miary.tex

\backmatter

% BEGIN chapters/pl/measure_theory/supplements/indeks_definicji_i_twierdzen.tex
% Indeks generowany na podstawie \begin{definicja|twierdzenie|lemat|wniosek}.
\chapter*{Indeks definicji i twierdzeń}
\markboth{INDEKS DEFINICJI I TWIERDZEŃ}{INDEKS DEFINICJI I TWIERDZEŃ}
\phantomsection\addcontentsline{toc}{chapter}{Indeks definicji i twierdzeń}
Indeks podaje nazwę rezultatu, jego numer i stronę. Nazwy i numery są aktywnymi odsyłaczami.
Uwzględniono także lematy i wnioski.

\section*{Definicje}
\begingroup\small
\setlength\LTleft{0pt}\setlength\LTright{0pt}
\begin{longtable}{@{}p{0.72\linewidth}p{0.135\linewidth}p{0.08\linewidth}@{}}
\textbf{Nazwa} & \textbf{Nr} & \textbf{Str.}\\\hline
\endfirsthead
\multicolumn{3}{@{}l}{\emph{Definicje -- ciąg dalszy}}\\[2pt]
\textbf{Nazwa} & \textbf{Nr} & \textbf{Str.}\\\hline
\endhead
\hyperref[def:atom]{Atom i miara bezatomowa} & \ref*{def:atom} & \pageref*{def:atom} \\
\hyperref[def:ac-measure-17]{Bezwzględna ciągłość miary} & \ref*{def:ac-measure-17} & \pageref*{def:ac-measure-17} \\
\hyperref[def:reduced-boundary-29]{Brzeg zredukowany} & \ref*{def:reduced-boundary-29} & \pageref*{def:reduced-boundary-29} \\
\hyperref[def:signed-integral-over-set-7]{Całka funkcji całkowalnej po zbiorze} & \ref*{def:signed-integral-over-set-7} & \pageref*{def:signed-integral-over-set-7} \\
\hyperref[def:integral-simple]{Całka funkcji prostej} & \ref*{def:integral-simple} & \pageref*{def:integral-simple} \\
\hyperref[def:complex-integral]{Całka funkcji zespolonej} & \ref*{def:complex-integral} & \pageref*{def:complex-integral} \\
\hyperref[def:integral-positive]{Całka Lebesgue'a funkcji nieujemnej} & \ref*{def:integral-positive} & \pageref*{def:integral-positive} \\
\hyperref[def:integral-over-set]{Całka po podzbiorze} & \ref*{def:integral-over-set} & \pageref*{def:integral-over-set} \\
\hyperref[def:tight-single-20]{Ciasność pojedynczej miary} & \ref*{def:tight-single-20} & \pageref*{def:tight-single-20} \\
\hyperref[def:cauchy-measure-9]{Ciąg Cauchy'ego według miary} & \ref*{def:cauchy-measure-9} & \pageref*{def:cauchy-measure-9} \\
\hyperref[def:positive-negative-parts]{Części dodatnia i ujemna} & \ref*{def:positive-negative-parts} & \pageref*{def:positive-negative-parts} \\
\hyperref[def:vitali-fine-24]{Drobne pokrycie Vitaliego} & \ref*{def:vitali-fine-24} & \pageref*{def:vitali-fine-24} \\
\hyperref[def:dual-lp-15]{Dualność i norma funkcjonału} & \ref*{def:dual-lp-15} & \pageref*{def:dual-lp-15} \\
\hyperref[def:HL-max-24]{Funkcja maksymalna Hardy'ego--Littlewooda} & \ref*{def:HL-max-24} & \pageref*{def:HL-max-24} \\
\hyperref[def:simple-function]{Funkcja prosta} & \ref*{def:simple-function} & \pageref*{def:simple-function} \\
\hyperref[idx:06:definicja:11]{Funkcja zespolona mierzalna} & \ref*{idx:06:definicja:11} & \pageref*{idx:06:definicja:11} \\
\hyperref[def:limsup-sets]{Granica górna i dolna zbiorów} & \ref*{def:limsup-sets} & \pageref*{def:limsup-sets} \\
\hyperref[def:topological-group-23]{Grupa topologiczna} & \ref*{def:topological-group-23} & \pageref*{def:topological-group-23} \\
\hyperref[def:normal-jacobian-28]{Jakobian normalny} & \ref*{def:normal-jacobian-28} & \pageref*{def:normal-jacobian-28} \\
\hyperref[def:volume-jacobian-27]{Jakobian objętościowy} & \ref*{def:volume-jacobian-27} & \pageref*{def:volume-jacobian-27} \\
\hyperref[def:k-jacobian-28]{Jakobian wymiaru $k$} & \ref*{def:k-jacobian-28} & \pageref*{def:k-jacobian-28} \\
\hyperref[def:uniform-integrability]{Jednostajna całkowalność} & \ref*{def:uniform-integrability} & \pageref*{def:uniform-integrability} \\
\hyperref[def:uniform-tight-22]{Jednostajna ciasność} & \ref*{def:uniform-tight-22} & \pageref*{def:uniform-tight-22} \\
\hyperref[def:monotone-class]{Klasa monotoniczna} & \ref*{def:monotone-class} & \pageref*{def:monotone-class} \\
\hyperref[def:measure-finiteness]{Klasy miar} & \ref*{def:measure-finiteness} & \pageref*{def:measure-finiteness} \\
\hyperref[def:hausdorff-measure-26]{Koszt pokrycia i miara Hausdorffa} & \ref*{def:hausdorff-measure-26} & \pageref*{def:hausdorff-measure-26} \\
\hyperref[def:linf-14]{Kres górny istotny i przestrzeń $L^\infty$} & \ref*{def:linf-14} & \pageref*{def:linf-14} \\
\hyperref[def:multiplicity-28]{Krotność odwzorowania} & \ref*{def:multiplicity-28} & \pageref*{def:multiplicity-28} \\
\hyperref[def:balls-24]{Kule i krotność pokrycia} & \ref*{def:balls-24} & \pageref*{def:balls-24} \\
\hyperref[def:haar-23]{Lewostronna i prawostronna miara Haara} & \ref*{def:haar-23} & \pageref*{def:haar-23} \\
\hyperref[def:local-majorant-11]{Lokalna dominacja całkowalna} & \ref*{def:local-majorant-11} & \pageref*{def:local-majorant-11} \\
\hyperref[def:max-measure-25]{Maksymalna funkcja miary} & \ref*{def:max-measure-25} & \pageref*{def:max-measure-25} \\
\hyperref[def:measure]{Miara} & \ref*{def:measure} & \pageref*{def:measure} \\
\hyperref[def:borel-locfinite-20]{Miara borelowska i lokalna skończoność} & \ref*{def:borel-locfinite-20} & \pageref*{def:borel-locfinite-20} \\
\hyperref[def:pushforward]{Miara obrazowa} & \ref*{def:pushforward} & \pageref*{def:pushforward} \\
\hyperref[def:radon-measure-20]{Miara Radona --- przyjęta konwencja} & \ref*{def:radon-measure-20} & \pageref*{def:radon-measure-20} \\
\hyperref[def:restriction]{Miara śladowa i ograniczenie} & \ref*{def:restriction} & \pageref*{def:restriction} \\
\hyperref[def:vector-measure-30]{Miara wektorowa} & \ref*{def:vector-measure-30} & \pageref*{def:vector-measure-30} \\
\hyperref[def:complex-measure-16]{Miara zespolona} & \ref*{def:complex-measure-16} & \pageref*{def:complex-measure-16} \\
\hyperref[def:outer-measure]{Miara zewnętrzna} & \ref*{def:outer-measure} & \pageref*{def:outer-measure} \\
\hyperref[def:lebesgue-outer]{Miara zewnętrzna Lebesgue'a} & \ref*{def:lebesgue-outer} & \pageref*{def:lebesgue-outer} \\
\hyperref[def:metric-outer-21]{Miara zewnętrzna metryczna} & \ref*{def:metric-outer-21} & \pageref*{def:metric-outer-21} \\
\hyperref[def:complete-measure]{Miara zupełna} & \ref*{def:complete-measure} & \pageref*{def:complete-measure} \\
\hyperref[def:extended-measurable]{Mierzalność funkcji rozszerzonej} & \ref*{def:extended-measurable} & \pageref*{def:extended-measurable} \\
\hyperref[def:normal-BV-29]{Normalna z pochodnej miary} & \ref*{def:normal-BV-29} & \pageref*{def:normal-BV-29} \\
\hyperref[def:support-radon-20]{Nośnik miary} & \ref*{def:support-radon-20} & \pageref*{def:support-radon-20} \\
\hyperref[def:finite-perimeter-29]{Obwód w sensie De Giorgiego} & \ref*{def:finite-perimeter-29} & \pageref*{def:finite-perimeter-29} \\
\hyperref[def:measurable-map]{Odwzorowanie mierzalne} & \ref*{def:measurable-map} & \pageref*{def:measurable-map} \\
\hyperref[def:algebry]{Pierścień, algebra, $\sigma$-algebra} & \ref*{def:algebry} & \pageref*{def:algebry} \\
\hyperref[def:five-convergences]{Pięć rodzajów zbieżności} & \ref*{def:five-convergences} & \pageref*{def:five-convergences} \\
\hyperref[def:RN-derivative-18]{Pochodna Radona--Nikodyma} & \ref*{def:RN-derivative-18} & \pageref*{def:RN-derivative-18} \\
\hyperref[def:current-29]{Prąd i jego brzeg} & \ref*{def:current-29} & \pageref*{def:current-29} \\
\hyperref[def:premeasure]{Premiara} & \ref*{def:premeasure} & \pageref*{def:premeasure} \\
\hyperref[def:halfopen-box]{Prostopadłościan półotwarty} & \ref*{def:halfopen-box} & \pageref*{def:halfopen-box} \\
\hyperref[def:complex-measures-space-16]{Przestrzeń $\mathcal M(X,\mathcal F)$} & \ref*{def:complex-measures-space-16} & \pageref*{def:complex-measures-space-16} \\
\hyperref[def:standard-borel]{Przestrzeń polska i standardowa borelowska} & \ref*{def:standard-borel} & \pageref*{def:standard-borel} \\
\hyperref[def:measure-boundary-29]{Punkty gęstości i brzeg miarowy} & \ref*{def:measure-boundary-29} & \pageref*{def:measure-boundary-29} \\
\hyperref[def:weak-lp-15]{Refleksywność i zbieżność słaba} & \ref*{def:weak-lp-15} & \pageref*{def:weak-lp-15} \\
\hyperref[def:regularity-20]{Regularność względem topologii} & \ref*{def:regularity-20} & \pageref*{def:regularity-20} \\
\hyperref[def:ae-functions]{Równość i zbieżność prawie wszędzie} & \ref*{def:ae-functions} & \pageref*{def:ae-functions} \\
\hyperref[def:lp-14]{Równość prawie wszędzie i przestrzeń $L^p$} & \ref*{def:lp-14} & \pageref*{def:lp-14} \\
\hyperref[def:frechet-rademacher-28]{Różniczkowalność Frécheta} & \ref*{def:frechet-rademacher-28} & \pageref*{def:frechet-rademacher-28} \\
\hyperref[def:borel]{$\sigma$-algebra borelowska} & \ref*{def:borel} & \pageref*{def:borel} \\
\hyperref[def:sigma-generated]{$\sigma$-algebra generowana} & \ref*{def:sigma-generated} & \pageref*{def:sigma-generated} \\
\hyperref[def:sigma-of-function]{$\sigma$-algebra generowana przez odwzorowanie} & \ref*{def:sigma-of-function} & \pageref*{def:sigma-of-function} \\
\hyperref[def:product-sigma]{$\sigma$-algebra produktowa} & \ref*{def:product-sigma} & \pageref*{def:product-sigma} \\
\hyperref[idx:01:definicja:04]{$\sigma$-algebra śladowa} & \ref*{idx:01:definicja:04} & \pageref*{idx:01:definicja:04} \\
\hyperref[idx:03:definicja:09]{\texorpdfstring{$\sigma$}{sigma}-skończona premiara} & \ref*{idx:03:definicja:09} & \pageref*{idx:03:definicja:09} \\
\hyperref[def:bochner-30]{Silna mierzalność i całka Bochnera} & \ref*{def:bochner-30} & \pageref*{def:bochner-30} \\
\hyperref[def:signed-measure-16]{Skończona miara ze znakiem} & \ref*{def:signed-measure-16} & \pageref*{def:signed-measure-16} \\
\hyperref[idx:21:definicja:08]{Skończona zespolona miara Radona} & \ref*{idx:21:definicja:08} & \pageref*{idx:21:definicja:08} \\
\hyperref[def:cantor-set]{Środkowotrzeci zbiór Cantora} & \ref*{def:cantor-set} & \pageref*{def:cantor-set} \\
\hyperref[def:test-19]{Test losowy i jego moc} & \ref*{def:test-19} & \pageref*{def:test-19} \\
\hyperref[def:energy-26]{$t$-energia miary} & \ref*{def:energy-26} & \pageref*{def:energy-26} \\
\hyperref[def:pi-lambda]{Układ $\pi$ i układ Dynkina} & \ref*{def:pi-lambda} & \pageref*{def:pi-lambda} \\
\hyperref[def:variation-signed-16]{Wariacja całkowita} & \ref*{def:variation-signed-16} & \pageref*{def:variation-signed-16} \\
\hyperref[def:BV-variation-29]{Wariacja funkcji} & \ref*{def:BV-variation-29} & \pageref*{def:BV-variation-29} \\
\hyperref[def:vector-variation-30]{Wariacja i półwariacja} & \ref*{def:vector-variation-30} & \pageref*{def:vector-variation-30} \\
\hyperref[def:lebesgue-point-25]{Wartość średnia i punkt Lebesgue'a} & \ref*{def:lebesgue-point-25} & \pageref*{def:lebesgue-point-25} \\
\hyperref[def:conditional-exp-18]{Warunkowa wartość oczekiwana} & \ref*{def:conditional-exp-18} & \pageref*{def:conditional-exp-18} \\
\hyperref[def:almost-everywhere]{Własność prawie wszędzie} & \ref*{def:almost-everywhere} & \pageref*{def:almost-everywhere} \\
\hyperref[def:RNP-30]{Własność Radona--Nikodyma} & \ref*{def:RNP-30} & \pageref*{def:RNP-30} \\
\hyperref[def:hausdorff-dimension-26]{Wymiar Hausdorffa} & \ref*{def:hausdorff-dimension-26} & \pageref*{def:hausdorff-dimension-26} \\
\hyperref[def:singular-16]{Wzajemna osobliwość} & \ref*{def:singular-16} & \pageref*{def:singular-16} \\
\hyperref[def:equivalent-measures-17]{Wzajemna równoważność miar} & \ref*{def:equivalent-measures-17} & \pageref*{def:equivalent-measures-17} \\
\hyperref[def:vague-weakstar-22]{Zbieżność słaba-* i mglista} & \ref*{def:vague-weakstar-22} & \pageref*{def:vague-weakstar-22} \\
\hyperref[def:narrow-measures-22]{Zbieżność wąska} & \ref*{def:narrow-measures-22} & \pageref*{def:narrow-measures-22} \\
\hyperref[def:convergence-measure]{Zbieżność według miary} & \ref*{def:convergence-measure} & \pageref*{def:convergence-measure} \\
\hyperref[def:lp-convergence-9]{Zbieżność w średniej rzędu $p$} & \ref*{def:lp-convergence-9} & \pageref*{def:lp-convergence-9} \\
\hyperref[def:rectifiable-29]{Zbiory przeliczalnie prostowalne} & \ref*{def:rectifiable-29} & \pageref*{def:rectifiable-29} \\
\hyperref[def:continuity-set-22]{Zbiór ciągłości miary} & \ref*{def:continuity-set-22} & \pageref*{def:continuity-set-22} \\
\hyperref[def:positive-negative-16]{Zbiór dodatni i ujemny} & \ref*{def:positive-negative-16} & \pageref*{def:positive-negative-16} \\
\hyperref[def:caratheodory-measurable]{Zbiór mierzalny w sensie Carathéodory'ego} & \ref*{def:caratheodory-measurable} & \pageref*{def:caratheodory-measurable} \\
\end{longtable}
\endgroup

\section*{Twierdzenia}
\begingroup\small
\setlength\LTleft{0pt}\setlength\LTright{0pt}
\begin{longtable}{@{}p{0.72\linewidth}p{0.135\linewidth}p{0.08\linewidth}@{}}
\textbf{Nazwa} & \textbf{Nr} & \textbf{Str.}\\\hline
\endfirsthead
\multicolumn{3}{@{}l}{\emph{Twierdzenia -- ciąg dalszy}}\\[2pt]
\textbf{Nazwa} & \textbf{Nr} & \textbf{Str.}\\\hline
\endhead
\hyperref[thm:positive-additivity]{Addytywność całki nieujemnej} & \ref*{thm:positive-additivity} & \pageref*{thm:positive-additivity} \\
\hyperref[thm:cc-approx-in-measure-10]{Aproksymacja funkcji o nośniku mierzalnym skończonej miary} & \ref*{thm:cc-approx-in-measure-10} & \pageref*{thm:cc-approx-in-measure-10} \\
\hyperref[thm:simple-real-approx]{Aproksymacja funkcji o wartościach rzeczywistych} & \ref*{thm:simple-real-approx} & \pageref*{thm:simple-real-approx} \\
\hyperref[thm:lebesgue-outer-regular]{Aproksymacja otwartymi zbiorami} & \ref*{thm:lebesgue-outer-regular} & \pageref*{thm:lebesgue-outer-regular} \\
\hyperref[thm:lebesgue-inner-regular]{Aproksymacja zbiorami zwartymi} & \ref*{thm:lebesgue-inner-regular} & \pageref*{thm:lebesgue-inner-regular} \\
\hyperref[thm:finite-metric-regular-20]{Aproksymacja zbiorów borelowskich} & \ref*{thm:finite-metric-regular-20} & \pageref*{thm:finite-metric-regular-20} \\
\hyperref[thm:atom-sigma-finite]{Atomy miar $\sigma$-skończonych mają skończoną masę} & \ref*{thm:atom-sigma-finite} & \pageref*{thm:atom-sigma-finite} \\
\hyperref[thm:bayes-finite-19]{Bayes dla skończonej przestrzeni parametrów} & \ref*{thm:bayes-finite-19} & \pageref*{thm:bayes-finite-19} \\
\hyperref[thm:image-ac-19]{Bezwzględna ciągłość zachowana przez obraz} & \ref*{thm:image-ac-19} & \pageref*{thm:image-ac-19} \\
\hyperref[thm:lebesgue-completion]{Borelowskie aproksymacje i uzupełnienie} & \ref*{thm:lebesgue-completion} & \pageref*{thm:lebesgue-completion} \\
\hyperref[thm:borel-Rd]{Borelowskie prostokąty} & \ref*{thm:borel-Rd} & \pageref*{thm:borel-Rd} \\
\hyperref[thm:moving-domains-11]{Całkowanie po zależnej od parametru dziedzinie} & \ref*{thm:moving-domains-11} & \pageref*{thm:moving-domains-11} \\
\hyperref[thm:change-measure-19]{Całkowanie po zmianie miary} & \ref*{thm:change-measure-19} & \pageref*{thm:change-measure-19} \\
\hyperref[thm:sum-absolutely-integrable]{Całkowanie szeregów absolutnie sumowalnych} & \ref*{thm:sum-absolutely-integrable} & \pageref*{thm:sum-absolutely-integrable} \\
\hyperref[thm:pushforward-integral-13]{Całkowanie względem miary obrazowej} & \ref*{thm:pushforward-integral-13} & \pageref*{thm:pushforward-integral-13} \\
\hyperref[thm:product-projections]{Charakteryzacja przez współrzędne} & \ref*{thm:product-projections} & \pageref*{thm:product-projections} \\
\hyperref[thm:polish-tight-20]{Ciasność skończonych miar na przestrzeniach polskich} & \ref*{thm:polish-tight-20} & \pageref*{thm:polish-tight-20} \\
\hyperref[thm:parameter-continuity]{Ciągłość pod znakiem całki} & \ref*{thm:parameter-continuity} & \pageref*{thm:parameter-continuity} \\
\hyperref[thm:translation-lp-14]{Ciągłość translacji w $L^p(\mathbb R^d)$} & \ref*{thm:translation-lp-14} & \pageref*{thm:translation-lp-14} \\
\hyperref[thm:measure-continuity]{Ciągłość z dołu i z góry} & \ref*{thm:measure-continuity} & \pageref*{thm:measure-continuity} \\
\hyperref[thm:atomic-cont-singular-17]{Dalszy rozkład miary osobliwej na prostej} & \ref*{thm:atomic-cont-singular-17} & \pageref*{thm:atomic-cont-singular-17} \\
\hyperref[thm:weak-lsc-15]{Dolna półciągłość normy} & \ref*{thm:weak-lsc-15} & \pageref*{thm:weak-lsc-15} \\
\hyperref[thm:perimeter-lsc-29]{Dolna półciągłość obwodu} & \ref*{thm:perimeter-lsc-29} & \pageref*{thm:perimeter-lsc-29} \\
\hyperref[thm:dual-lp-15]{Dualność $L^p$ dla $1<p<\infty$} & \ref*{thm:dual-lp-15} & \pageref*{thm:dual-lp-15} \\
\hyperref[thm:dual-l1-15]{Dualność $L^1$ przy $\sigma$-skończoności} & \ref*{thm:dual-l1-15} & \pageref*{thm:dual-l1-15} \\
\hyperref[thm:tv-dual-tests-16]{Dualny opis normy w wariacji całkowitej} & \ref*{thm:tv-dual-tests-16} & \pageref*{thm:tv-dual-tests-16} \\
\hyperref[thm:ui-two-criteria-8]{Dwa równoważne opisy jednostajnej całkowalności} & \ref*{thm:ui-two-criteria-8} & \pageref*{thm:ui-two-criteria-8} \\
\hyperref[thm:measurable-algebra]{Działania algebraiczne} & \ref*{thm:measurable-algebra} & \pageref*{thm:measurable-algebra} \\
\hyperref[thm:total-variation-formula-19]{Formuła wariacyjna dla dwóch prawdopodobieństw} & \ref*{thm:total-variation-formula-19} & \pageref*{thm:total-variation-formula-19} \\
\hyperref[thm:frostman-26]{Frostmana, wersja dla zbiorów zwartych} & \ref*{thm:frostman-26} & \pageref*{thm:frostman-26} \\
\hyperref[thm:fubini-12]{Fubini} & \ref*{thm:fubini-12} & \pageref*{thm:fubini-12} \\
\hyperref[thm:cdf-derivative-25]{Funkcja rozkładu a absolutna ciągłość miary} & \ref*{thm:cdf-derivative-25} & \pageref*{thm:cdf-derivative-25} \\
\hyperref[thm:measurable-real]{Funkcje rzeczywiste i działania} & \ref*{thm:measurable-real} & \pageref*{thm:measurable-real} \\
\hyperref[thm:pushforward-composition]{Funktorialność miar obrazowych} & \ref*{thm:pushforward-composition} & \pageref*{thm:pushforward-composition} \\
\hyperref[thm:cc-dense-lp-radon-20]{Gęstość $C_c(X)$ w przestrzeniach $L^p$} & \ref*{thm:cc-dense-lp-radon-20} & \pageref*{thm:cc-dense-lp-radon-20} \\
\hyperref[thm:cc-in-lp-10]{Gęstość funkcji gładkich} & \ref*{thm:cc-in-lp-10} & \pageref*{thm:cc-in-lp-10} \\
\hyperref[thm:cc-dense-lp-14]{Gęstość funkcji gładkich w przestrzeni euklidesowej} & \ref*{thm:cc-dense-lp-14} & \pageref*{thm:cc-dense-lp-14} \\
\hyperref[thm:simple-dense-lp-14]{Gęstość funkcji prostych} & \ref*{thm:simple-dense-lp-14} & \pageref*{thm:simple-dense-lp-14} \\
\hyperref[thm:pushforward-coarea-28]{Gęstość obrazu przez submersję} & \ref*{thm:pushforward-coarea-28} & \pageref*{thm:pushforward-coarea-28} \\
\hyperref[thm:smooth-perimeter-29]{Granica klasycznej geometrii} & \ref*{thm:smooth-perimeter-29} & \pageref*{thm:smooth-perimeter-29} \\
\hyperref[thm:haar-main-23]{Haara: istnienie i jednoznaczność} & \ref*{thm:haar-main-23} & \pageref*{thm:haar-main-23} \\
\hyperref[thm:laplace-holo-11]{Holomorficzność transformaty Laplace'a} & \ref*{thm:laplace-holo-11} & \pageref*{thm:laplace-holo-11} \\
\hyperref[thm:holder-14]{Hölder} & \ref*{thm:holder-14} & \pageref*{thm:holder-14} \\
\hyperref[thm:finiteness-implications]{Implikacje} & \ref*{thm:finiteness-implications} & \pageref*{thm:finiteness-implications} \\
\hyperref[thm:lebesgue-exists]{Istnienie i podstawowe własności} & \ref*{thm:lebesgue-exists} & \pageref*{thm:lebesgue-exists} \\
\hyperref[thm:vitali]{Istnienie zbioru niemierzalnego} & \ref*{thm:vitali} & \pageref*{thm:vitali} \\
\hyperref[thm:CE-existence-18]{Istnienie, jednoznaczność i własności podstawowe} & \ref*{thm:CE-existence-18} & \pageref*{thm:CE-existence-18} \\
\hyperref[thm:BV-layercake-29]{Jedna nierówność formuły koobszaru dla $BV$} & \ref*{thm:BV-layercake-29} & \pageref*{thm:BV-layercake-29} \\
\hyperref[thm:uniform-convexity-15]{Jednostajna wypukłość} & \ref*{thm:uniform-convexity-15} & \pageref*{thm:uniform-convexity-15} \\
\hyperref[thm:haar-unique-23]{Jednoznaczność miary Haara} & \ref*{thm:haar-unique-23} & \pageref*{thm:haar-unique-23} \\
\hyperref[thm:uniqueness-cc-20]{Jednoznaczność miary Radona przez całki} & \ref*{thm:uniqueness-cc-20} & \pageref*{thm:uniqueness-cc-20} \\
\hyperref[thm:uniqueness-finite]{Jednoznaczność miar skończonych} & \ref*{thm:uniqueness-finite} & \pageref*{thm:uniqueness-finite} \\
\hyperref[thm:uniqueness-sigma-finite]{Jednoznaczność przy wspólnym wyczerpaniu} & \ref*{thm:uniqueness-sigma-finite} & \pageref*{thm:uniqueness-sigma-finite} \\
\hyperref[thm:extension-uniqueness]{Jednoznaczność rozszerzenia} & \ref*{thm:extension-uniqueness} & \pageref*{thm:extension-uniqueness} \\
\hyperref[thm:egorov-9]{Jegorow} & \ref*{thm:egorov-9} & \pageref*{thm:egorov-9} \\
\hyperref[thm:jensen-14]{Jensen dla miary probabilistycznej} & \ref*{thm:jensen-14} & \pageref*{thm:jensen-14} \\
\hyperref[thm:hausdorff-calibration-26]{Kalibracja wymiaru całkowitego bez twierdzeń geometrycznych} & \ref*{thm:hausdorff-calibration-26} & \pageref*{thm:hausdorff-calibration-26} \\
\hyperref[thm:bochner-basic-30]{Konstrukcja i oszacowanie całki Bochnera} & \ref*{thm:bochner-basic-30} & \pageref*{thm:bochner-basic-30} \\
\hyperref[thm:product-measure-12]{Konstrukcja miary produktowej} & \ref*{thm:product-measure-12} & \pageref*{thm:product-measure-12} \\
\hyperref[thm:outer-from-cover]{Konstrukcja przez pokrycia} & \ref*{thm:outer-from-cover} & \pageref*{thm:outer-from-cover} \\
\hyperref[thm:local-to-global-measure-9]{Kontrola ucieczki masy przy zbieżności lokalnej} & \ref*{thm:local-to-global-measure-9} & \pageref*{thm:local-to-global-measure-9} \\
\hyperref[thm:pointwise-limits]{Kresy i granice} & \ref*{thm:pointwise-limits} & \pageref*{thm:pointwise-limits} \\
\hyperref[thm:integrability-criterion]{Kryterium absolutnej całkowalności} & \ref*{thm:integrability-criterion} & \pageref*{thm:integrability-criterion} \\
\hyperref[thm:premeasure-continuity]{Kryterium ciągłości premiary skończonej} & \ref*{thm:premeasure-continuity} & \pageref*{thm:premeasure-continuity} \\
\hyperref[thm:eps-delta-ac-17]{Kryterium epsilon--delta} & \ref*{thm:eps-delta-ac-17} & \pageref*{thm:eps-delta-ac-17} \\
\hyperref[thm:tails-11]{Kryterium jednostajnej kontroli ogonów} & \ref*{thm:tails-11} & \pageref*{thm:tails-11} \\
\hyperref[thm:riemann-criterion]{Kryterium Lebesgue'a całkowalności Riemanna} & \ref*{thm:riemann-criterion} & \pageref*{thm:riemann-criterion} \\
\hyperref[thm:measurable-generator]{Kryterium przez generatory} & \ref*{thm:measurable-generator} & \pageref*{thm:measurable-generator} \\
\hyperref[thm:weak-lp-15]{Kryterium słabej zbieżności w $L^p$} & \ref*{thm:weak-lp-15} & \pageref*{thm:weak-lp-15} \\
\hyperref[thm:param-vitali-11]{Kryterium Vitaliego dla parametru} & \ref*{thm:param-vitali-11} & \pageref*{thm:param-vitali-11} \\
\hyperref[thm:series-lp-14]{Kryterium zbieżności szeregów w $L^p$} & \ref*{thm:series-lp-14} & \pageref*{thm:series-lp-14} \\
\hyperref[thm:lebesgue-diff-25]{Lebesgue'a o różniczkowaniu} & \ref*{thm:lebesgue-diff-25} & \pageref*{thm:lebesgue-diff-25} \\
\hyperref[thm:borel-cantelli-1]{Lemat Borela--Cantellego, część pierwsza} & \ref*{thm:borel-cantelli-1} & \pageref*{thm:borel-cantelli-1} \\
\hyperref[thm:fatou]{Lemat Fatou} & \ref*{thm:fatou} & \pageref*{thm:fatou} \\
\hyperref[thm:fatou-integrable-lower-8]{Lemat Fatou z całkowalnym ograniczeniem od dołu} & \ref*{thm:fatou-integrable-lower-8} & \pageref*{thm:fatou-integrable-lower-8} \\
\hyperref[thm:neyman-pearson-19]{Lemat Neymana--Pearsona} & \ref*{thm:neyman-pearson-19} & \pageref*{thm:neyman-pearson-19} \\
\hyperref[thm:5r-24]{Lemat o pięciokrotnym powiększeniu} & \ref*{thm:5r-24} & \pageref*{thm:5r-24} \\
\hyperref[thm:scheffe-9]{Lemat Scheffégo} & \ref*{thm:scheffe-9} & \pageref*{thm:scheffe-9} \\
\hyperref[thm:integral-linear-real]{Liniowość całki funkcji całkowalnych} & \ref*{thm:integral-linear-real} & \pageref*{thm:integral-linear-real} \\
\hyperref[thm:integral-complex-linear]{Liniowość zespolona i oszacowanie modułu} & \ref*{thm:integral-complex-linear} & \pageref*{thm:integral-complex-linear} \\
\hyperref[thm:linear-change-13]{Liniowy wzór na objętość} & \ref*{thm:linear-change-13} & \pageref*{thm:linear-change-13} \\
\hyperref[thm:local-convergence-measure]{Lokalna zbieżność według miary} & \ref*{thm:local-convergence-measure} & \pageref*{thm:local-convergence-measure} \\
\hyperref[thm:local-continuity-11]{Lokalne kryterium ciągłości} & \ref*{thm:local-continuity-11} & \pageref*{thm:local-continuity-11} \\
\hyperref[thm:null-local-27]{Lokalne kryterium zbiorów zerowych} & \ref*{thm:null-local-27} & \pageref*{thm:null-local-27} \\
\hyperref[thm:lch-secondcountable-radon-20]{Lokalnie skończone miary borelowskie} & \ref*{thm:lch-secondcountable-radon-20} & \pageref*{thm:lch-secondcountable-radon-20} \\
\hyperref[thm:lusin-10]{Łuzin} & \ref*{thm:lusin-10} & \pageref*{thm:lusin-10} \\
\hyperref[thm:measure-metric-9]{Metryka zbieżności według miary} & \ref*{thm:measure-metric-9} & \pageref*{thm:measure-metric-9} \\
\hyperref[thm:graph-area-28]{Miara na gładkim wykresie} & \ref*{thm:graph-area-28} & \pageref*{thm:graph-area-28} \\
\hyperref[thm:riemann-radon-27]{Miara objętościowa jest miarą Radona} & \ref*{thm:riemann-radon-27} & \pageref*{thm:riemann-radon-27} \\
\hyperref[thm:product-riemann-27]{Miara produktu riemannowskiego} & \ref*{thm:product-riemann-27} & \pageref*{thm:product-riemann-27} \\
\hyperref[thm:rn-manifold-27]{Miary absolutnie ciągłe względem objętości} & \ref*{thm:rn-manifold-27} & \pageref*{thm:rn-manifold-27} \\
\hyperref[thm:restriction-meas-5]{Mierzalność ograniczenia} & \ref*{thm:restriction-meas-5} & \pageref*{thm:restriction-meas-5} \\
\hyperref[thm:minimize-lp-15]{Minimalizacja wypukłej energii} & \ref*{thm:minimize-lp-15} & \pageref*{thm:minimize-lp-15} \\
\hyperref[thm:initial-sigma]{Minimalna struktura mierzalna} & \ref*{thm:initial-sigma} & \pageref*{thm:initial-sigma} \\
\hyperref[thm:minkowski-14]{Minkowski} & \ref*{thm:minkowski-14} & \pageref*{thm:minkowski-14} \\
\hyperref[thm:null-modifications]{Modyfikacje na zbiorach zerowych} & \ref*{thm:null-modifications} & \pageref*{thm:null-modifications} \\
\hyperref[thm:simple-up]{Monotoniczna aproksymacja funkcji nieujemnej} & \ref*{thm:simple-up} & \pageref*{thm:simple-up} \\
\hyperref[thm:fatou-for-sets]{Nierówności dla granic zbiorów} & \ref*{thm:fatou-for-sets} & \pageref*{thm:fatou-for-sets} \\
\hyperref[thm:lp-implies-measure-9]{Nierówność Czebyszewa i zbieżność według miary} & \ref*{thm:lp-implies-measure-9} & \pageref*{thm:lp-implies-measure-9} \\
\hyperref[thm:lp-interpolation-14]{Nierówność interpolacyjna} & \ref*{thm:lp-interpolation-14} & \pageref*{thm:lp-interpolation-14} \\
\hyperref[thm:markov-nonnegative]{Nierówność Markowa dla całki} & \ref*{thm:markov-nonnegative} & \pageref*{thm:markov-nonnegative} \\
\hyperref[thm:HL-weak-24]{Nierówność słabego typu $(1,1)$} & \ref*{thm:HL-weak-24} & \pageref*{thm:HL-weak-24} \\
\hyperref[thm:gibbs-19]{Nieujemność entropii względnej} & \ref*{thm:gibbs-19} & \pageref*{thm:gibbs-19} \\
\hyperref[thm:ae-integral-invariance]{Niezmienniczość całki przy zmianie na zbiorze zerowym} & \ref*{thm:ae-integral-invariance} & \pageref*{thm:ae-integral-invariance} \\
\hyperref[thm:support-fullmass-20]{Nośnik na przestrzeni drugiej przeliczalności} & \ref*{thm:support-fullmass-20} & \pageref*{thm:support-fullmass-20} \\
\hyperref[thm:riesz-c0-positive-21]{Ograniczone dodatnie funkcjonały na $C_0$} & \ref*{thm:riesz-c0-positive-21} & \pageref*{thm:riesz-c0-positive-21} \\
\hyperref[thm:extended-operations]{Operacje na funkcjach rozszerzonych} & \ref*{thm:extended-operations} & \pageref*{thm:extended-operations} \\
\hyperref[thm:partition-variation-16]{Opis wariacji przez podziały} & \ref*{thm:partition-variation-16} & \pageref*{thm:partition-variation-16} \\
\hyperref[thm:bounded-integral-complex-16]{Oszacowanie całki względem miary zespolonej} & \ref*{thm:bounded-integral-complex-16} & \pageref*{thm:bounded-integral-complex-16} \\
\hyperref[thm:frostman-energy-26]{Oszacowanie energii przez warunek Frostmana} & \ref*{thm:frostman-energy-26} & \pageref*{thm:frostman-energy-26} \\
\hyperref[thm:weak-measure-25]{Oszacowanie słabego typu dla miar} & \ref*{thm:weak-measure-25} & \pageref*{thm:weak-measure-25} \\
\hyperref[thm:mct]{O zbieżności monotonicznej; Beppo Levi} & \ref*{thm:mct} & \pageref*{thm:mct} \\
\hyperref[thm:dct]{O zbieżności zmajoryzowanej; Lebesgue} & \ref*{thm:dct} & \pageref*{thm:dct} \\
\hyperref[thm:BV-coarea-full-29]{Pełna formuła koobszaru dla $BV$} & \ref*{thm:BV-coarea-full-29} & \pageref*{thm:BV-coarea-full-29} \\
\hyperref[thm:local-differentiation-11]{Pierwsza pochodna — wersja lokalna} & \ref*{thm:local-differentiation-11} & \pageref*{thm:local-differentiation-11} \\
\hyperref[thm:nonnegative-first-properties]{Pierwsze własności całki nieujemnej} & \ref*{thm:nonnegative-first-properties} & \pageref*{thm:nonnegative-first-properties} \\
\hyperref[thm:BV-measure-29]{Pochodna $Du$ jako miara wektorowa} & \ref*{thm:BV-measure-29} & \pageref*{thm:BV-measure-29} \\
\hyperref[thm:RN-inverse-18]{Pochodna odwrotna} & \ref*{thm:RN-inverse-18} & \pageref*{thm:RN-inverse-18} \\
\hyperref[thm:higher-differentiation-11]{Pochodne wyższego rzędu} & \ref*{thm:higher-differentiation-11} & \pageref*{thm:higher-differentiation-11} \\
\hyperref[thm:weak-subseq-15]{Podciąg słabo zbieżny: wariant separowalny} & \ref*{thm:weak-subseq-15} & \pageref*{thm:weak-subseq-15} \\
\hyperref[thm:measure-subsequence]{Podciąg zbieżny prawie wszędzie} & \ref*{thm:measure-subsequence} & \pageref*{thm:measure-subsequence} \\
\hyperref[thm:conv-implications-9]{Podstawowe implikacje} & \ref*{thm:conv-implications-9} & \pageref*{thm:conv-implications-9} \\
\hyperref[thm:FTC-lebesgue-25]{Podstawowe twierdzenie rachunku całkowego dla całki Lebesgue'a} & \ref*{thm:FTC-lebesgue-25} & \pageref*{thm:FTC-lebesgue-25} \\
\hyperref[thm:basic-measure]{Podstawowe własności miary} & \ref*{thm:basic-measure} & \pageref*{thm:basic-measure} \\
\hyperref[thm:cantor-properties]{Podstawowe własności zbioru Cantora} & \ref*{thm:cantor-properties} & \pageref*{thm:cantor-properties} \\
\hyperref[thm:vitali-cover-24]{Pokryciowe twierdzenie Vitaliego} & \ref*{thm:vitali-cover-24} & \pageref*{thm:vitali-cover-24} \\
\hyperref[thm:portmanteau-22]{Portmanteau} & \ref*{thm:portmanteau-22} & \pageref*{thm:portmanteau-22} \\
\hyperref[thm:haar-modular-23]{Prawa miara Haara i modularność} & \ref*{thm:haar-modular-23} & \pageref*{thm:haar-modular-23} \\
\hyperref[thm:ae-smooth-approx-10]{Prawie wszędzie gładka aproksymacja funkcji mierzalnej} & \ref*{thm:ae-smooth-approx-10} & \pageref*{thm:ae-smooth-approx-10} \\
\hyperref[thm:ae-to-measure-finite]{Prawie wszędzie implikuje zbieżność według miary przy mierze skończonej} & \ref*{thm:ae-to-measure-finite} & \pageref*{thm:ae-to-measure-finite} \\
\hyperref[thm:volume-premeasure]{Premiara na ograniczonym prostopadłościanie} & \ref*{thm:volume-premeasure} & \pageref*{thm:volume-premeasure} \\
\hyperref[thm:prokhorov-22]{Prochorow na przestrzeni polskiej} & \ref*{thm:prokhorov-22} & \pageref*{thm:prokhorov-22} \\
\hyperref[thm:prokhorov-forward-22]{Prochorow: kierunek ciasność $\Rightarrow$ zwartość} & \ref*{thm:prokhorov-forward-22} & \pageref*{thm:prokhorov-forward-22} \\
\hyperref[thm:borel-product]{Produkt borelowski przestrzeni drugoprzeliczalnych} & \ref*{thm:borel-product} & \pageref*{thm:borel-product} \\
\hyperref[thm:lebesgue-products-12]{Produkt miar Lebesgue'a} & \ref*{thm:lebesgue-products-12} & \pageref*{thm:lebesgue-products-12} \\
\hyperref[thm:compact-extension-10]{Przedłużenie z kompaktu w $\mathbb R^d$} & \ref*{thm:compact-extension-10} & \pageref*{thm:compact-extension-10} \\
\hyperref[thm:intersection-sigma]{Przekroje $\sigma$-algebr} & \ref*{thm:intersection-sigma} & \pageref*{thm:intersection-sigma} \\
\hyperref[thm:borel-generators-R]{Przeliczalne generatory na $\mathbb R$} & \ref*{thm:borel-generators-R} & \pageref*{thm:borel-generators-R} \\
\hyperref[thm:standard-countably-generated]{Przeliczalne generowanie standardowych przestrzeni} & \ref*{thm:standard-countably-generated} & \pageref*{thm:standard-countably-generated} \\
\hyperref[thm:hilbert-rnp-30]{Przestrzenie Hilberta mają RNP} & \ref*{thm:hilbert-rnp-30} & \pageref*{thm:hilbert-rnp-30} \\
\hyperref[thm:CE-orthogonal-18]{Przypadek $L^2$: rzut ortogonalny} & \ref*{thm:CE-orthogonal-18} & \pageref*{thm:CE-orthogonal-18} \\
\hyperref[thm:simple-properties]{Rachunek dla funkcji prostych} & \ref*{thm:simple-properties} & \pageref*{thm:simple-properties} \\
\hyperref[thm:rademacher-28]{Rademacher} & \ref*{thm:rademacher-28} & \pageref*{thm:rademacher-28} \\
\hyperref[thm:RN-finite-18]{Radon--Nikodym dla miar skończonych} & \ref*{thm:RN-finite-18} & \pageref*{thm:RN-finite-18} \\
\hyperref[thm:RN-complex-18]{Radon--Nikodym dla miar zespolonych} & \ref*{thm:RN-complex-18} & \pageref*{thm:RN-complex-18} \\
\hyperref[thm:reflexive-lp-15]{Refleksywność $L^p$} & \ref*{thm:reflexive-lp-15} & \pageref*{thm:reflexive-lp-15} \\
\hyperref[thm:hausdorff-gdelta-26]{Regularizacja przez zbiory $G_\delta$} & \ref*{thm:hausdorff-gdelta-26} & \pageref*{thm:hausdorff-gdelta-26} \\
\hyperref[thm:RN-chain-18]{Reguła łańcuchowa} & \ref*{thm:RN-chain-18} & \pageref*{thm:RN-chain-18} \\
\hyperref[thm:riesz-positive-21]{Reprezentacja dodatnich funkcjonałów} & \ref*{thm:riesz-positive-21} & \pageref*{thm:riesz-positive-21} \\
\hyperref[thm:compact-riesz-21]{Reprezentacja na zwartej przestrzeni metrycznej} & \ref*{thm:compact-riesz-21} & \pageref*{thm:compact-riesz-21} \\
\hyperref[thm:riesz-complex-21]{Reprezentacja zespolona i norma wariacji} & \ref*{thm:riesz-complex-21} & \pageref*{thm:riesz-complex-21} \\
\hyperref[thm:riesz-fischer-14]{Riesz--Fischer: zupełność $L^p$} & \ref*{thm:riesz-fischer-14} & \pageref*{thm:riesz-fischer-14} \\
\hyperref[thm:rnp-lipschitz-equiv-30]{RNP a różniczkowalność krzywych Lipschitzowskich} & \ref*{thm:rnp-lipschitz-equiv-30} & \pageref*{thm:rnp-lipschitz-equiv-30} \\
\hyperref[thm:polar-complex-measure-18]{Rozkład biegunowy miary zespolonej} & \ref*{thm:polar-complex-measure-18} & \pageref*{thm:polar-complex-measure-18} \\
\hyperref[thm:hahn-16]{Rozkład Hahna} & \ref*{thm:hahn-16} & \pageref*{thm:hahn-16} \\
\hyperref[thm:jordan-16]{Rozkład Jordana} & \ref*{thm:jordan-16} & \pageref*{thm:jordan-16} \\
\hyperref[thm:lebesgue-decomp-finite-17]{Rozkład Lebesgue'a dla miar skończonych} & \ref*{thm:lebesgue-decomp-finite-17} & \pageref*{thm:lebesgue-decomp-finite-17} \\
\hyperref[thm:lebesgue-decomp-sigfin-17]{Rozkład Lebesgue'a w przypadku $\sigma$-skończonym} & \ref*{thm:lebesgue-decomp-sigfin-17} & \pageref*{thm:lebesgue-decomp-sigfin-17} \\
\hyperref[thm:leb-manifold-27]{Rozkład na część objętościową i osobliwą} & \ref*{thm:leb-manifold-27} & \pageref*{thm:leb-manifold-27} \\
\hyperref[thm:extension-caratheodory]{Rozszerzenie Carathéodory'ego} & \ref*{thm:extension-caratheodory} & \pageref*{thm:extension-caratheodory} \\
\hyperref[thm:measure-differentiation-25]{Różniczkowanie miary względem Lebesgue'a} & \ref*{thm:measure-differentiation-25} & \pageref*{thm:measure-differentiation-25} \\
\hyperref[thm:differentiating-integrals]{Różniczkowanie pod znakiem całki} & \ref*{thm:differentiating-integrals} & \pageref*{thm:differentiating-integrals} \\
\hyperref[thm:hausdorff-borel-26]{Samodzielna konstrukcja miary Hausdorffa} & \ref*{thm:hausdorff-borel-26} & \pageref*{thm:hausdorff-borel-26} \\
\hyperref[thm:separable-lp-14]{Separowalność $L^p$} & \ref*{thm:separable-lp-14} & \pageref*{thm:separable-lp-14} \\
\hyperref[thm:hausdorff-lipschitz-26]{Skalowanie i odwzorowania Lipschitzowskie} & \ref*{thm:hausdorff-lipschitz-26} & \pageref*{thm:hausdorff-lipschitz-26} \\
\hyperref[thm:besicovitch-finite-24]{Skończona postać lematu Besicovitcha} & \ref*{thm:besicovitch-finite-24} & \pageref*{thm:besicovitch-finite-24} \\
\hyperref[thm:measure-operations]{Sumy i mnożenie przez stałą} & \ref*{thm:measure-operations} & \pageref*{thm:measure-operations} \\
\hyperref[thm:sum-nonnegative-integrals]{Szeregi nieujemnych funkcji mierzalnych} & \ref*{thm:sum-nonnegative-integrals} & \pageref*{thm:sum-nonnegative-integrals} \\
\hyperref[thm:extended-tests]{Testy mierzalności na półprostych} & \ref*{thm:extended-tests} & \pageref*{thm:extended-tests} \\
\hyperref[thm:complex-test]{Test na częściach rzeczywistej i urojonej} & \ref*{thm:complex-test} & \pageref*{thm:complex-test} \\
\hyperref[thm:tonelli-12]{Tonelli} & \ref*{thm:tonelli-12} & \pageref*{thm:tonelli-12} \\
\hyperref[thm:completed-product-fubini-12]{Tonelli i Fubini dla uzupełnionej miary produktowej} & \ref*{thm:completed-product-fubini-12} & \pageref*{thm:completed-product-fubini-12} \\
\hyperref[thm:section-identity-12]{Tożsamość przekrojowa przy $\sigma$-skończoności} & \ref*{thm:section-identity-12} & \pageref*{thm:section-identity-12} \\
\hyperref[thm:density-transform-13]{Transformacja gęstości} & \ref*{thm:density-transform-13} & \pageref*{thm:density-transform-13} \\
\hyperref[thm:invariance-lebesgue]{Translacje i jednolite skalowanie} & \ref*{thm:invariance-lebesgue} & \pageref*{thm:invariance-lebesgue} \\
\hyperref[thm:caratheodory]{Twierdzenie Carathéodory'ego} & \ref*{thm:caratheodory} & \pageref*{thm:caratheodory} \\
\hyperref[thm:gauss-green-BV-29]{Twierdzenie Gaussa--Greena w postaci miarowej} & \ref*{thm:gauss-green-BV-29} & \pageref*{thm:gauss-green-BV-29} \\
\hyperref[thm:monotone-class]{Twierdzenie o klasie monotonicznej} & \ref*{thm:monotone-class} & \pageref*{thm:monotone-class} \\
\hyperref[thm:pi-lambda]{Twierdzenie $\pi$--$\lambda$ Dynkina} & \ref*{thm:pi-lambda} & \pageref*{thm:pi-lambda} \\
\hyperref[thm:besicovitch-24]{Twierdzenie pokryciowe Besicovitcha -- pełny dowód} & \ref*{thm:besicovitch-24} & \pageref*{thm:besicovitch-24} \\
\hyperref[thm:RN-sigfin-18]{Twierdzenie Radona--Nikodyma} & \ref*{thm:RN-sigfin-18} & \pageref*{thm:RN-sigfin-18} \\
\hyperref[thm:de-giorgi-structure-29]{Twierdzenie strukturalne De Giorgiego} & \ref*{thm:de-giorgi-structure-29} & \pageref*{thm:de-giorgi-structure-29} \\
\hyperref[thm:completion]{Uzupełnienie przestrzeni z miarą} & \ref*{thm:completion} & \pageref*{thm:completion} \\
\hyperref[thm:caratheodory-completion]{Uzupełnienie rozszerzenia} & \ref*{thm:caratheodory-completion} & \pageref*{thm:caratheodory-completion} \\
\hyperref[thm:vector-variation-measure-30]{Wariacja jest miarą dodatnią} & \ref*{thm:vector-variation-measure-30} & \pageref*{thm:vector-variation-measure-30} \\
\hyperref[thm:density-variation-16]{Wariacja miary określonej całką} & \ref*{thm:density-variation-16} & \pageref*{thm:density-variation-16} \\
\hyperref[thm:variation-density-30]{Wariacja miary o gęstości Bochnera} & \ref*{thm:variation-density-30} & \pageref*{thm:variation-density-30} \\
\hyperref[thm:complex-variation-16]{Wariacja miary zespolonej} & \ref*{thm:complex-variation-16} & \pageref*{thm:complex-variation-16} \\
\hyperref[thm:conditional-jensen-19]{Warunkowa nierówność Jensena} & \ref*{thm:conditional-jensen-19} & \pageref*{thm:conditional-jensen-19} \\
\hyperref[thm:lp-embedding-finite-14]{Włożenia przy skończonej mierze} & \ref*{thm:lp-embedding-finite-14} & \pageref*{thm:lp-embedding-finite-14} \\
\hyperref[thm:cantor-dimension-26]{Wymiar klasycznego zbioru Cantora} & \ref*{thm:cantor-dimension-26} & \pageref*{thm:cantor-dimension-26} \\
\hyperref[thm:euclidean-dimension-26]{Wymiar przestrzeni euklidesowej} & \ref*{thm:euclidean-dimension-26} & \pageref*{thm:euclidean-dimension-26} \\
\hyperref[thm:lipschitz-area-coarea-28]{Wzory area i coarea dla odwzorowań Lipschitzowskich} & \ref*{thm:lipschitz-area-coarea-28} & \pageref*{thm:lipschitz-area-coarea-28} \\
\hyperref[thm:coarea-submersion-28]{Wzór koobszaru dla submersji} & \ref*{thm:coarea-submersion-28} & \pageref*{thm:coarea-submersion-28} \\
\hyperref[thm:layer-cake-8]{Wzór warstwowy na całkę nieujemną} & \ref*{thm:layer-cake-8} & \pageref*{thm:layer-cake-8} \\
\hyperref[thm:immersion-area-28]{Wzór z krotnością dla immersji $C^1$} & \ref*{thm:immersion-area-28} & \pageref*{thm:immersion-area-28} \\
\hyperref[thm:hausdorff-line-26]{W wymiarze jeden otrzymujemy miarę Lebesgue'a} & \ref*{thm:hausdorff-line-26} & \pageref*{thm:hausdorff-line-26} \\
\hyperref[thm:change-variables-13]{Zamiana zmiennych dla dyfeomorfizmu $C^1$} & \ref*{thm:change-variables-13} & \pageref*{thm:change-variables-13} \\
\hyperref[thm:manifold-change-27]{Zamiana zmiennych na rozmaitościach} & \ref*{thm:manifold-change-27} & \pageref*{thm:manifold-change-27} \\
\hyperref[thm:kl-data-processing-19]{Zasada przetwarzania danych dla entropii} & \ref*{thm:kl-data-processing-19} & \pageref*{thm:kl-data-processing-19} \\
\hyperref[thm:mass-distribution-26]{Zasada rozkładu masy} & \ref*{thm:mass-distribution-26} & \pageref*{thm:mass-distribution-26} \\
\hyperref[thm:uniform-parameter-11]{Zbieżność jednostajna po scałkowaniu} & \ref*{thm:uniform-parameter-11} & \pageref*{thm:uniform-parameter-11} \\
\hyperref[thm:mct-ae]{Zbieżność monotoniczna prawie wszędzie} & \ref*{thm:mct-ae} & \pageref*{thm:mct-ae} \\
\hyperref[thm:vitali-l1]{Zbieżność Vitaliego na przestrzeni o mierze skończonej} & \ref*{thm:vitali-l1} & \pageref*{thm:vitali-l1} \\
\hyperref[thm:lp-to-measure-14]{Zbieżność w $L^p$ a zbieżność według miary} & \ref*{thm:lp-to-measure-14} & \pageref*{thm:lp-to-measure-14} \\
\hyperref[thm:tv-implies-narrow-22]{Zbieżność w wariacji implikuje wąską} & \ref*{thm:tv-implies-narrow-22} & \pageref*{thm:tv-implies-narrow-22} \\
\hyperref[thm:borel-trace]{Zbiory borelowskie podprzestrzeni} & \ref*{thm:borel-trace} & \pageref*{thm:borel-trace} \\
\hyperref[thm:riemann-to-lebesgue]{Zgodność całki Riemanna i Lebesgue'a} & \ref*{thm:riemann-to-lebesgue} & \pageref*{thm:riemann-to-lebesgue} \\
\hyperref[thm:composition-measurable]{Złożenia i funkcje ciągłe} & \ref*{thm:composition-measurable} & \pageref*{thm:composition-measurable} \\
\hyperref[thm:BV-banach-29]{Zupełność przestrzeni $BV$} & \ref*{thm:BV-banach-29} & \pageref*{thm:BV-banach-29} \\
\hyperref[thm:complete-measures-16]{Zupełność przestrzeni miar} & \ref*{thm:complete-measures-16} & \pageref*{thm:complete-measures-16} \\
\hyperref[thm:cauchy-measure-complete-9]{Zupełność zbieżności według miary} & \ref*{thm:cauchy-measure-complete-9} & \pageref*{thm:cauchy-measure-complete-9} \\
\hyperref[thm:perimeter-compact-29]{Zwartość rodzin zbiorów o ograniczonym obwodzie} & \ref*{thm:perimeter-compact-29} & \pageref*{thm:perimeter-compact-29} \\
\end{longtable}
\endgroup

\section*{Lematy}
\begingroup\small
\setlength\LTleft{0pt}\setlength\LTright{0pt}
\begin{longtable}{@{}p{0.72\linewidth}p{0.135\linewidth}p{0.08\linewidth}@{}}
\textbf{Nazwa} & \textbf{Nr} & \textbf{Str.}\\\hline
\endfirsthead
\multicolumn{3}{@{}l}{\emph{Lematy -- ciąg dalszy}}\\[2pt]
\textbf{Nazwa} & \textbf{Nr} & \textbf{Str.}\\\hline
\endhead
\hyperref[lem:hausdorff-metricity-26]{Addytywność na zbiorach oddalonych} & \ref*{lem:hausdorff-metricity-26} & \pageref*{lem:hausdorff-metricity-26} \\
\hyperref[lem:outer-approx-algebra]{Aproksymacja od góry} & \ref*{lem:outer-approx-algebra} & \pageref*{lem:outer-approx-algebra} \\
\hyperref[lem:vector-variation-ac-30]{Bezwzględna ciągłość wariacji} & \ref*{lem:vector-variation-ac-30} & \pageref*{lem:vector-variation-ac-30} \\
\hyperref[lem:metric-outer-borel-21]{Borelowska mierzalność} & \ref*{lem:metric-outer-borel-21} & \pageref*{lem:metric-outer-borel-21} \\
\hyperref[lem:current-boundary-square-29]{Brzeg brzegu znika} & \ref*{lem:current-boundary-square-29} & \pageref*{lem:current-boundary-square-29} \\
\hyperref[lem:essential-boundary-null-29]{Brzeg miarowy jest objętościowo mały} & \ref*{lem:essential-boundary-null-29} & \pageref*{lem:essential-boundary-null-29} \\
\hyperref[lem:FTC-Banach-30]{Całka po krzywej z pochodną Banachowską} & \ref*{lem:FTC-Banach-30} & \pageref*{lem:FTC-Banach-30} \\
\hyperref[lem:integration-density-18]{Całkowanie względem miary o gęstości} & \ref*{lem:integration-density-18} & \pageref*{lem:integration-density-18} \\
\hyperref[lem:positive-content-26]{Dodatnia zawartość zbioru zwartego} & \ref*{lem:positive-content-26} & \pageref*{lem:positive-content-26} \\
\hyperref[lem:midpoint-lp-15]{Dwie nierówności dla punktu środkowego} & \ref*{lem:midpoint-lp-15} & \pageref*{lem:midpoint-lp-15} \\
\hyperref[lem:nonatomic-divisibility-30]{Dzielenie miary bezatomowej} & \ref*{lem:nonatomic-divisibility-30} & \pageref*{lem:nonatomic-divisibility-30} \\
\hyperref[lem:singularity-17]{Elementarne własności osobliwości} & \ref*{lem:singularity-17} & \pageref*{lem:singularity-17} \\
\hyperref[lem:haar-fubini-compact-23]{Fubini dla ciągłych jąder o zwartym nośniku} & \ref*{lem:haar-fubini-compact-23} & \pageref*{lem:haar-fubini-compact-23} \\
\hyperref[lem:cutoff-radon-20]{Funkcja odcinająca} & \ref*{lem:cutoff-radon-20} & \pageref*{lem:cutoff-radon-20} \\
\hyperref[lem:density-cc-c0-21]{Gęstość $C_c(X)$ w $C_0(X)$} & \ref*{lem:density-cc-c0-21} & \pageref*{lem:density-cc-c0-21} \\
\hyperref[lem:dg-density-identification]{Gęstość miary na wykresach prostowalnych} & \ref*{lem:dg-density-identification} & \pageref*{lem:dg-density-identification} \\
\hyperref[lem:smooth-cutoff-10]{Gładka funkcja odcinająca} & \ref*{lem:smooth-cutoff-10} & \pageref*{lem:smooth-cutoff-10} \\
\hyperref[lem:canonical-lpq-15]{Izometryczne zanurzenie do dualnej} & \ref*{lem:canonical-lpq-15} & \pageref*{lem:canonical-lpq-15} \\
\hyperref[lem:haar-ratio-23]{Jednolite oszacowania ilorazów} & \ref*{lem:haar-ratio-23} & \pageref*{lem:haar-ratio-23} \\
\hyperref[lem:weak-sequence-tight-22]{Jednostajna ciasność ciągu wąsko zbieżnego} & \ref*{lem:weak-sequence-tight-22} & \pageref*{lem:weak-sequence-tight-22} \\
\hyperref[lem:haar-uniform-23]{Jednostajna ciągłość translacji funkcji} & \ref*{lem:haar-uniform-23} & \pageref*{lem:haar-uniform-23} \\
\hyperref[lem:rademacher-line-28]{Jednowymiarowa różniczkowalność} & \ref*{lem:rademacher-line-28} & \pageref*{lem:rademacher-line-28} \\
\hyperref[lem:uniform-absolute-continuity]{Kontrola całek na małych zbiorach} & \ref*{lem:uniform-absolute-continuity} & \pageref*{lem:uniform-absolute-continuity} \\
\hyperref[lem:hausdorff-caratheodory-selfcontained-26]{Kryterium Carathéodory'ego, z dowodem} & \ref*{lem:hausdorff-caratheodory-selfcontained-26} & \pageref*{lem:hausdorff-caratheodory-selfcontained-26} \\
\hyperref[lem:dynkin-intersection]{Lemat z rozdziału 1} & \ref*{lem:dynkin-intersection} & \pageref*{lem:dynkin-intersection} \\
\hyperref[lem:lipschitz-approx-compact-10]{Lipschitzowska aproksymacja na zbiorze zwartym} & \ref*{lem:lipschitz-approx-compact-10} & \pageref*{lem:lipschitz-approx-compact-10} \\
\hyperref[lem:lusin-simple-10]{Łuzin dla funkcji prostych} & \ref*{lem:lusin-simple-10} & \pageref*{lem:lusin-simple-10} \\
\hyperref[lem:haar-neigh-23]{Małe otoczenia, zwarte rozdrobnienia i pokrycia} & \ref*{lem:haar-neigh-23} & \pageref*{lem:haar-neigh-23} \\
\hyperref[lem:graph-hausdorff-28]{Małe zniekształcenie wykresu} & \ref*{lem:graph-hausdorff-28} & \pageref*{lem:graph-hausdorff-28} \\
\hyperref[lem:hausdorff-metric-open-26]{Metryczna miara zewnętrzna a zbiory domknięte} & \ref*{lem:hausdorff-metric-open-26} & \pageref*{lem:hausdorff-metric-open-26} \\
\hyperref[lem:BV-level-meas-auditVI]{Mierzalność obwodów poziomic} & \ref*{lem:BV-level-meas-auditVI} & \pageref*{lem:BV-level-meas-auditVI} \\
\hyperref[lem:sections-measurable-12]{Mierzalność wszystkich przekrojów} & \ref*{lem:sections-measurable-12} & \pageref*{lem:sections-measurable-12} \\
\hyperref[lem:variational-lp-15]{Minimum energii} & \ref*{lem:variational-lp-15} & \pageref*{lem:variational-lp-15} \\
\hyperref[lem:dg-relative-isoperimetric]{Nierówności Sobolewa i względna nierówność izoperymetryczna} & \ref*{lem:dg-relative-isoperimetric} & \pageref*{lem:dg-relative-isoperimetric} \\
\hyperref[lem:rect-cover-ineq]{Nierówność objętościowa dla pokryć} & \ref*{lem:rect-cover-ineq} & \pageref*{lem:rect-cover-ineq} \\
\hyperref[lem:simple-refinement]{Niezależność od reprezentacji} & \ref*{lem:simple-refinement} & \pageref*{lem:simple-refinement} \\
\hyperref[lem:critical-area-null-28]{Obraz zbioru o niedostatecznej randze} & \ref*{lem:critical-area-null-28} & \pageref*{lem:critical-area-null-28} \\
\hyperref[lem:BV-smooth-levels]{Obwody poziomic funkcji gładkiej} & \ref*{lem:BV-smooth-levels} & \pageref*{lem:BV-smooth-levels} \\
\hyperref[lem:lipschitz-null-13]{Odwzorowania Lipschitzowskie zachowują zbiory zerowe} & \ref*{lem:lipschitz-null-13} & \pageref*{lem:lipschitz-null-13} \\
\hyperref[lem:haar-content-to-measure-23]{Od zawartości na kompaktach do miary borelowskiej} & \ref*{lem:haar-content-to-measure-23} & \pageref*{lem:haar-content-to-measure-23} \\
\hyperref[lem:bounded-signed-16]{Ograniczoność zakresu} & \ref*{lem:bounded-signed-16} & \pageref*{lem:bounded-signed-16} \\
\hyperref[lem:dg-density]{Oszacowania gęstości na brzegu zredukowanym} & \ref*{lem:dg-density} & \pageref*{lem:dg-density} \\
\hyperref[lem:small-cube-13]{Oszacowanie obrazu małego sześcianu} & \ref*{lem:small-cube-13} & \pageref*{lem:small-cube-13} \\
\hyperref[lem:BV-translation-29]{Oszacowanie przesunięcia} & \ref*{lem:BV-translation-29} & \pageref*{lem:BV-translation-29} \\
\hyperref[lem:perimeter-complement-29]{Pochodna i dopełnienie} & \ref*{lem:perimeter-complement-29} & \pageref*{lem:perimeter-complement-29} \\
\hyperref[lem:premeasure-cover]{Podaddytywność premiary na pokryciach} & \ref*{lem:premeasure-cover} & \pageref*{lem:premeasure-cover} \\
\hyperref[lem:rectifiable-union-29]{Podstawowe własności prostowalności} & \ref*{lem:rectifiable-union-29} & \pageref*{lem:rectifiable-union-29} \\
\hyperref[lem:dg-blowup]{Półprzestrzeń jako jedyna granica po przeskalowaniu} & \ref*{lem:dg-blowup} & \pageref*{lem:dg-blowup} \\
\hyperref[lem:dg-rectifiability]{Prostowalność brzegu zredukowanego} & \ref*{lem:dg-rectifiability} & \pageref*{lem:dg-rectifiability} \\
\hyperref[lem:hausdorff-threshold-26]{Próg zerowania} & \ref*{lem:hausdorff-threshold-26} & \pageref*{lem:hausdorff-threshold-26} \\
\hyperref[lem:dg-ball-cut]{Przecięcie zbioru skończonego obwodu z kulą} & \ref*{lem:dg-ball-cut} & \pageref*{lem:dg-ball-cut} \\
\hyperref[lem:whitney-c1-28]{Przedłużenie zgodnych danych pierwszego rzędu} & \ref*{lem:whitney-c1-28} & \pageref*{lem:whitney-c1-28} \\
\hyperref[lem:finite-rectangle-sections-12]{Przekroje na prostokącie skończonej miary} & \ref*{lem:finite-rectangle-sections-12} & \pageref*{lem:finite-rectangle-sections-12} \\
\hyperref[lem:lipschitz-extend-22]{Przybliżanie lipschitzowskie i przedłużanie} & \ref*{lem:lipschitz-extend-22} & \pageref*{lem:lipschitz-extend-22} \\
\hyperref[lem:continuous-dense-25]{Przybliżanie lokalne funkcją ciągłą} & \ref*{lem:continuous-dense-25} & \pageref*{lem:continuous-dense-25} \\
\hyperref[lem:rnp-partition-30]{Przybliżanie miary przez skończone podziały} & \ref*{lem:rnp-partition-30} & \pageref*{lem:rnp-partition-30} \\
\hyperref[lem:compact-inside-10]{Przybliżenie mierzalnego zbioru przez zbiór zwarty} & \ref*{lem:compact-inside-10} & \pageref*{lem:compact-inside-10} \\
\hyperref[lem:radon-sigmafinite-inner-20]{Regularność na zbiorach $\sigma$-skończonych} & \ref*{lem:radon-sigmafinite-inner-20} & \pageref*{lem:radon-sigmafinite-inner-20} \\
\hyperref[lem:bochner-differentiation-30]{Różniczkowanie całki Bochnera} & \ref*{lem:bochner-differentiation-30} & \pageref*{lem:bochner-differentiation-30} \\
\hyperref[lem:besicovitch-angle-24]{Skończoność konfiguracji Besicovitcha} & \ref*{lem:besicovitch-angle-24} & \pageref*{lem:besicovitch-angle-24} \\
\hyperref[lem:rect-grid]{Skończony podział na siatkę} & \ref*{lem:rect-grid} & \pageref*{lem:rect-grid} \\
\hyperref[lem:rademacher-direction-28]{Słabe i kierunkowe pochodne funkcji Lipschitzowskiej} & \ref*{lem:rademacher-direction-28} & \pageref*{lem:rademacher-direction-28} \\
\hyperref[lem:BV-strict-approx]{Ścisła aproksymacja gładka w $BV$} & \ref*{lem:BV-strict-approx} & \pageref*{lem:BV-strict-approx} \\
\hyperref[lem:trace-generation]{Ślad i generowanie} & \ref*{lem:trace-generation} & \pageref*{lem:trace-generation} \\
\hyperref[lem:linear-coarea-28]{Tożsamość liniowo-algebraiczna koobszaru} & \ref*{lem:linear-coarea-28} & \pageref*{lem:linear-coarea-28} \\
\hyperref[lem:shear-13]{Translacje, permutacje i ścinanie} & \ref*{lem:shear-13} & \pageref*{lem:shear-13} \\
\hyperref[lem:BV-smooth-29]{Wariacja gładkiej funkcji} & \ref*{lem:BV-smooth-29} & \pageref*{lem:BV-smooth-29} \\
\hyperref[lem:hilbert-cube-22]{Włożenie separowalnej przestrzeni metrycznej w zwartą przestrzeń metryczną} & \ref*{lem:hilbert-cube-22} & \pageref*{lem:hilbert-cube-22} \\
\hyperref[lem:critical-coarea-null-28]{Włókna zwartego zbioru krytycznego} & \ref*{lem:critical-coarea-null-28} & \pageref*{lem:critical-coarea-null-28} \\
\hyperref[lem:conditional-pullout-19]{Wyciąganie czynnika mierzalnego} & \ref*{lem:conditional-pullout-19} & \pageref*{lem:conditional-pullout-19} \\
\hyperref[lem:exhaustion-21]{Wyczerpanie przez zbiory relatywnie zwarte} & \ref*{lem:exhaustion-21} & \pageref*{lem:exhaustion-21} \\
\hyperref[lem:mollify-cc-10]{Wygładzanie funkcji ciągłej o zwartym nośniku} & \ref*{lem:mollify-cc-10} & \pageref*{lem:mollify-cc-10} \\
\hyperref[lem:c1-general-area-28]{Wzór area dla dowolnego odwzorowania $C^1$} & \ref*{lem:c1-general-area-28} & \pageref*{lem:c1-general-area-28} \\
\hyperref[lem:young-numbers-14]{Young dla iloczynu dwóch liczb} & \ref*{lem:young-numbers-14} & \pageref*{lem:young-numbers-14} \\
\hyperref[lem:singular-density-25]{Zanik gęstości miary osobliwej} & \ref*{lem:singular-density-25} & \pageref*{lem:singular-density-25} \\
\hyperref[lem:haar-content-23]{Zawartość niezmiennicza na kompaktach} & \ref*{lem:haar-content-23} & \pageref*{lem:haar-content-23} \\
\hyperref[lem:simple-exhaustion]{Zbieżność całek prostej funkcji na zbiorach rosnących} & \ref*{lem:simple-exhaustion} & \pageref*{lem:simple-exhaustion} \\
\hyperref[lem:hyperplane-null]{Zbiory zawarte w hiperpłaszczyźnie} & \ref*{lem:hyperplane-null} & \pageref*{lem:hyperplane-null} \\
\hyperref[lem:null-fibres-lipschitz-28]{Zbiory zerowe i ich włókna} & \ref*{lem:null-fibres-lipschitz-28} & \pageref*{lem:null-fibres-lipschitz-28} \\
\hyperref[lem:hausdorff-outer-26]{Zewnętrzne własności kosztu pokryć} & \ref*{lem:hausdorff-outer-26} & \pageref*{lem:hausdorff-outer-26} \\
\hyperref[lem:restriction-integral]{Zgodność z ograniczeniem miary} & \ref*{lem:restriction-integral} & \pageref*{lem:restriction-integral} \\
\hyperref[lem:esssup-14]{Znaczenie kresu istotnego} & \ref*{lem:esssup-14} & \pageref*{lem:esssup-14} \\
\hyperref[lem:lusin-whitney-28]{Zwarte fragmenty klasy $C^1$} & \ref*{lem:lusin-whitney-28} & \pageref*{lem:lusin-whitney-28} \\
\hyperref[lem:haar-product-compact-23]{Zwartość dowolnego produktu przedziałów} & \ref*{lem:haar-product-compact-23} & \pageref*{lem:haar-product-compact-23} \\
\hyperref[lem:compact-prob-22]{Zwartość prawdopodobieństw na przestrzeni zwartej} & \ref*{lem:compact-prob-22} & \pageref*{lem:compact-prob-22} \\
\end{longtable}
\endgroup

\section*{Wnioski}
\begingroup\small
\setlength\LTleft{0pt}\setlength\LTright{0pt}
\begin{longtable}{@{}p{0.72\linewidth}p{0.135\linewidth}p{0.08\linewidth}@{}}
\textbf{Nazwa} & \textbf{Nr} & \textbf{Str.}\\\hline
\endfirsthead
\multicolumn{3}{@{}l}{\emph{Wnioski -- ciąg dalszy}}\\[2pt]
\textbf{Nazwa} & \textbf{Nr} & \textbf{Str.}\\\hline
\endhead
\hyperref[cor:smooth-approx-in-measure-10]{Aproksymacja w mierze funkcjami gładkimi} & \ref*{cor:smooth-approx-in-measure-10} & \pageref*{cor:smooth-approx-in-measure-10} \\
\hyperref[cor:besicovitch-measure-infinite-24]{Besicovitch dla dodatnich miar borelowskich} & \ref*{cor:besicovitch-measure-infinite-24} & \pageref*{cor:besicovitch-measure-infinite-24} \\
\hyperref[cor:borel-representative-10]{Borelowska reprezentacja funkcji Lebesgue'a} & \ref*{cor:borel-representative-10} & \pageref*{cor:borel-representative-10} \\
\hyperref[cor:coordinate-measurable]{Borelowskie odwzorowania wektorowe} & \ref*{cor:coordinate-measurable} & \pageref*{cor:coordinate-measurable} \\
\hyperref[cor:meas-simple-characterization]{Charakteryzacja mierzalności przez funkcje proste} & \ref*{cor:meas-simple-characterization} & \pageref*{cor:meas-simple-characterization} \\
\hyperref[cor:coercive-tight-22]{Ciasność z oszacowania momentów} & \ref*{cor:coercive-tight-22} & \pageref*{cor:coercive-tight-22} \\
\hyperref[cor:uniform-abs-cont-int-17]{Ciągłość całki względem zbioru całkowania} & \ref*{cor:uniform-abs-cont-int-17} & \pageref*{cor:uniform-abs-cont-int-17} \\
\hyperref[cor:l1-param-11]{Ciągłość w normie $L^1$} & \ref*{cor:l1-param-11} & \pageref*{cor:l1-param-11} \\
\hyperref[cor:haar-open-positive-23]{Dodatniość miary zbiorów otwartych} & \ref*{cor:haar-open-positive-23} & \pageref*{cor:haar-open-positive-23} \\
\hyperref[cor:holder-sharp-15]{Dokładne maksimum w nierówności Höldera} & \ref*{cor:holder-sharp-15} & \pageref*{cor:holder-sharp-15} \\
\hyperref[cor:lusin-finite-ae-10]{Funkcje skończone prawie wszędzie} & \ref*{cor:lusin-finite-ae-10} & \pageref*{cor:lusin-finite-ae-10} \\
\hyperref[cor:lebesgue-regularity-characterization]{Geometryczna charakterystyka mierzalności} & \ref*{cor:lebesgue-regularity-characterization} & \pageref*{cor:lebesgue-regularity-characterization} \\
\hyperref[cor:countable-lusin-10]{Jednoczesne twierdzenie Łuzina dla przeliczalnej rodziny} & \ref*{cor:countable-lusin-10} & \pageref*{cor:countable-lusin-10} \\
\hyperref[cor:measure-limit-unique]{Jednoznaczność granicy według miary} & \ref*{cor:measure-limit-unique} & \pageref*{cor:measure-limit-unique} \\
\hyperref[cor:vague-c0-bounded-22]{Kiedy zbieżność mglista jest słabą-*?} & \ref*{cor:vague-c0-bounded-22} & \pageref*{cor:vague-c0-bounded-22} \\
\hyperref[cor:NP-exact-alpha-19]{Konstrukcja testu optymalnego o zadanym poziomie} & \ref*{cor:NP-exact-alpha-19} & \pageref*{cor:NP-exact-alpha-19} \\
\hyperref[cor:simple-approx-integral]{Konstruktywne obliczanie całki} & \ref*{cor:simple-approx-integral} & \pageref*{cor:simple-approx-integral} \\
\hyperref[cor:conditional-lp-19]{Kontrakcja w $L^p$} & \ref*{cor:conditional-lp-19} & \pageref*{cor:conditional-lp-19} \\
\hyperref[cor:HL-level-24]{Kontrola zbioru, na którym średnie są duże} & \ref*{cor:HL-level-24} & \pageref*{cor:HL-level-24} \\
\hyperref[cor:tonelli-integrability-12]{Kryterium nieujemnej całkowalności} & \ref*{cor:tonelli-integrability-12} & \pageref*{cor:tonelli-integrability-12} \\
\hyperref[cor:practical-fubini-12]{Kryterium praktyczne} & \ref*{cor:practical-fubini-12} & \pageref*{cor:practical-fubini-12} \\
\hyperref[cor:every-subseq-9]{Kryterium za pomocą wszystkich podciągów} & \ref*{cor:every-subseq-9} & \pageref*{cor:every-subseq-9} \\
\hyperref[cor:nonborel-lebesgue]{Lebesgue'owskie zbiory nieborelowskie} & \ref*{cor:nonborel-lebesgue} & \pageref*{cor:nonborel-lebesgue} \\
\hyperref[cor:associative-product-12]{Łączność skończonych produktów} & \ref*{cor:associative-product-12} & \pageref*{cor:associative-product-12} \\
\hyperref[cor:riemann-hausdorff-28]{Miara riemannowska i Hausdorffa} & \ref*{cor:riemann-hausdorff-28} & \pageref*{cor:riemann-hausdorff-28} \\
\hyperref[cor:measures-separate-21]{Miary odróżniane przez funkcje ciągłe} & \ref*{cor:measures-separate-21} & \pageref*{cor:measures-separate-21} \\
\hyperref[cor:haar-adjoint-23]{Modularność grupy Liego} & \ref*{cor:haar-adjoint-23} & \pageref*{cor:haar-adjoint-23} \\
\hyperref[idx:06:wniosek:06]{Monotoniczne ciągi funkcji} & \ref*{idx:06:wniosek:06} & \pageref*{idx:06:wniosek:06} \\
\hyperref[cor:minimal-variation-16]{Najmniejsza miara dominująca} & \ref*{cor:minimal-variation-16} & \pageref*{cor:minimal-variation-16} \\
\hyperref[cor:integral-triangle]{Nierówność trójkąta i struktura $L^1$} & \ref*{cor:integral-triangle} & \pageref*{cor:integral-triangle} \\
\hyperref[cor:no-all-subsets-invariant]{Nie istnieje miara długości wszystkich podzbiorów} & \ref*{cor:no-all-subsets-invariant} & \pageref*{cor:no-all-subsets-invariant} \\
\hyperref[cor:submanifold-null-27]{Obrazy immersji niższego wymiaru} & \ref*{cor:submanifold-null-27} & \pageref*{cor:submanifold-null-27} \\
\hyperref[cor:reverse-fatou]{Odwrotny lemat Fatou przy dominacji} & \ref*{cor:reverse-fatou} & \pageref*{cor:reverse-fatou} \\
\hyperref[cor:5r-volume-24]{Oszacowanie miary sumy kul} & \ref*{cor:5r-volume-24} & \pageref*{cor:5r-volume-24} \\
\hyperref[cor:finite-besicovitch-measure-24]{Oszacowanie sumy miar w rodzinie skończonej} & \ref*{cor:finite-besicovitch-measure-24} & \pageref*{cor:finite-besicovitch-measure-24} \\
\hyperref[cor:RN-lebesgue-sum-18]{Pełny rozkład miary na część z gęstością i osobliwą} & \ref*{cor:RN-lebesgue-sum-18} & \pageref*{cor:RN-lebesgue-sum-18} \\
\hyperref[cor:riesz-order-21]{Porządek dodatnich funkcjonałów i miar} & \ref*{cor:riesz-order-21} & \pageref*{cor:riesz-order-21} \\
\hyperref[cor:null-sections-12]{Prawie wszystkie przekroje zbioru zerowego} & \ref*{cor:null-sections-12} & \pageref*{cor:null-sections-12} \\
\hyperref[cor:real-borel-tests]{Prawostronny test mierzalności} & \ref*{cor:real-borel-tests} & \pageref*{cor:real-borel-tests} \\
\hyperref[cor:null-union]{Przeliczalne sumy zbiorów miary zero} & \ref*{cor:null-union} & \pageref*{cor:null-union} \\
\hyperref[cor:countably-many-disjoint-atoms]{Przeliczalność rozłącznych atomów} & \ref*{cor:countably-many-disjoint-atoms} & \pageref*{cor:countably-many-disjoint-atoms} \\
\hyperref[idx:30:wniosek:15]{Przestrzenie Banacha i całki z krzywych} & \ref*{idx:30:wniosek:15} & \pageref*{idx:30:wniosek:15} \\
\hyperref[cor:lp-normed-14]{Przestrzeń normowana} & \ref*{cor:lp-normed-14} & \pageref*{cor:lp-normed-14} \\
\hyperref[cor:tv-data-processing-19]{Przetwarzanie danych nie zwiększa wariacji całkowitej} & \ref*{cor:tv-data-processing-19} & \pageref*{cor:tv-data-processing-19} \\
\hyperref[cor:indicator-smooth-10]{Przybliżanie funkcji charakterystycznej} & \ref*{cor:indicator-smooth-10} & \pageref*{cor:indicator-smooth-10} \\
\hyperref[cor:vitali-open-24]{Przybliżanie zbioru rozłącznymi kulami} & \ref*{cor:vitali-open-24} & \pageref*{cor:vitali-open-24} \\
\hyperref[cor:hausdorff-singletons-auditVI]{Punkty i zbiory przeliczalne} & \ref*{cor:hausdorff-singletons-auditVI} & \pageref*{cor:hausdorff-singletons-auditVI} \\
\hyperref[cor:integral-countable-partition]{Rozkład całki na rozłączne zbiory} & \ref*{cor:integral-countable-partition} & \pageref*{cor:integral-countable-partition} \\
\hyperref[cor:complex-lebesgue-decomp-17]{Rozkład miary zespolonej} & \ref*{cor:complex-lebesgue-decomp-17} & \pageref*{cor:complex-lebesgue-decomp-17} \\
\hyperref[cor:RN-linearity-18]{Różnica i suma miar z gęstościami} & \ref*{cor:RN-linearity-18} & \pageref*{cor:RN-linearity-18} \\
\hyperref[cor:coarea-level-28]{Różniczkowanie poziomic} & \ref*{cor:coarea-level-28} & \pageref*{cor:coarea-level-28} \\
\hyperref[cor:cubes-differentiation-25]{Średnie po kostkach} & \ref*{cor:cubes-differentiation-25} & \pageref*{cor:cubes-differentiation-25} \\
\hyperref[cor:bct]{Twierdzenie o zbieżności ograniczonej} & \ref*{cor:bct} & \pageref*{cor:bct} \\
\hyperref[idx:08:wniosek:20]{Uogólniona formuła rozkładu na poziomice} & \ref*{idx:08:wniosek:20} & \pageref*{idx:08:wniosek:20} \\
\hyperref[cor:generalized-holder-14]{Uogólniony Hölder} & \ref*{cor:generalized-holder-14} & \pageref*{cor:generalized-holder-14} \\
\hyperref[cor:diff-basis-25]{Wartość funkcji z otaczających ją średnich} & \ref*{cor:diff-basis-25} & \pageref*{cor:diff-basis-25} \\
\hyperref[idx:09:wniosek:14]{Wersja dla gęstości prawdopodobieństwa} & \ref*{idx:09:wniosek:14} & \pageref*{idx:09:wniosek:14} \\
\hyperref[cor:local-egorov-9]{Wersja lokalna dla przestrzeni $\sigma$-skończonej} & \ref*{cor:local-egorov-9} & \pageref*{cor:local-egorov-9} \\
\hyperref[cor:lusin-extension-10]{Wersja Łuzina przez funkcję ciągłą na całej przestrzeni} & \ref*{cor:lusin-extension-10} & \pageref*{cor:lusin-extension-10} \\
\hyperref[idx:06:wniosek:15]{Wersje mierzalne funkcji określonych p.w.} & \ref*{idx:06:wniosek:15} & \pageref*{idx:06:wniosek:15} \\
\hyperref[cor:graph-dimension-26]{Wykres funkcji Lipschitzowskiej} & \ref*{cor:graph-dimension-26} & \pageref*{cor:graph-dimension-26} \\
\hyperref[cor:hausdorff-union-26]{Wymiar sumy przeliczalnej} & \ref*{cor:hausdorff-union-26} & \pageref*{cor:hausdorff-union-26} \\
\hyperref[cor:BV-coarea-weighted-29]{Wzór koobszaru z wagą} & \ref*{cor:BV-coarea-weighted-29} & \pageref*{cor:BV-coarea-weighted-29} \\
\hyperref[cor:change-subdomain-13]{Zamiana na mierzalnym podzbiorze} & \ref*{cor:change-subdomain-13} & \pageref*{cor:change-subdomain-13} \\
\hyperref[cor:vague-narrow-22]{Zbieżność mglista bez ucieczki masy} & \ref*{cor:vague-narrow-22} & \pageref*{cor:vague-narrow-22} \\
\hyperref[cor:dct-l1]{Zbieżność zmajoryzowana w $L^1$} & \ref*{cor:dct-l1} & \pageref*{cor:dct-l1} \\
\hyperref[cor:measure-metric-complete-9]{Zupełność metryki zbieżności według miary} & \ref*{cor:measure-metric-complete-9} & \pageref*{cor:measure-metric-complete-9} \\
\end{longtable}
\endgroup


% Dżety dodane jako rozdział 31 — bez przenumerowania indeksu historycznego.

% END chapters/pl/measure_theory/supplements/indeks_definicji_i_twierdzen.tex

\end{document}
